
\documentclass[11pt]{article}
\usepackage[T1]{fontenc}
\usepackage{amsmath,amssymb,amsthm}
\usepackage[margin=1in]{geometry}
\usepackage{enumitem}
\usepackage[colorlinks=true,linkcolor=blue,urlcolor=blue]{hyperref}

\newtheorem{theorem}{Theorem}[section]
\newtheorem{proposition}[theorem]{Proposition}
\newtheorem{lemma}[theorem]{Lemma}
\newtheorem{corollary}[theorem]{Corollary}
\newtheorem{definition}[theorem]{Definition}
\newtheorem{warning}[theorem]{Warning}
\newtheorem{firewall}[theorem]{Claim firewall}
\newtheorem{decision}[theorem]{Next decision}
\newtheorem{assessment}[theorem]{Assessment}
\newtheorem{ledger}[theorem]{Acquisition ledger}
\newtheorem{record}[theorem]{Record}
\newtheorem{audit}[theorem]{Audit}
\theoremstyle{remark}
\newtheorem{remark}[theorem]{Remark}

\title{Renewal Standard Model--Einstein Continuum Research Notes\\
Fifteenth Cycle, V1}
\author{}
\date{10 September 2026}

\begin{document}
\maketitle
\tableofcontents

\section{Context, lineage, and purpose}
\label{sec:v15v1}

These notes continue the attempt to construct one rigorous
Standard Model--Einstein continuum from the attached Grand Tensor
regulator and its related Standard Model, Yang--Mills, and
Navier--Stokes programmes.  The intended endpoint is a single
continuum state with compatible gauge, Higgs, fermion, ghost, and
gravitational sectors; exact physical Ward and BRST identities;
Osterwalder--Schrader reconstruction; a local covariant observable
net; a controlled Einstein source equation; and the required
massive and massless spectral structure.

That endpoint has not been reached.  The earlier notes contain
exact finite-cutoff constructions, conditional continuum theorems,
spectral compilers, and counterexamples to insufficient proof
routes.  This compact restart summarizes those results without
strengthening their hypotheses, records the present bottleneck,
and then continues the proof programme with a distance adapted to
the unbounded quartic Higgs field.

The completed fourteenth cycle is frozen as
\[
 \texttt{Renewal\_SM\_Einstein\_Continuum\_Research\_Notes\_V100\_f14.tex}.
                                                                  \tag{1.1}
\]
Its validated record is
\[
\begin{array}{c|c}
 \text{lines}&34825\\
 \text{bytes}&1166228\\
 \text{SHA--256}&
 \texttt{C6608FD46EE95038605B57FE09CC584946F88DF6D480C7A53443385B94E64BD9}.
\end{array}                                                       \tag{1.2}
\]
Future versions append to this V1 and replace only the immediate
active predecessor.  The newest Yang--Mills and Navier--Stokes
notes are inspected at every fifth version.  When an active line
reaches V100, it is frozen with the next unused suffix and another
compact V1 begins.

\begin{record}[Primary source provenance]
\label{rec:v15v1-sources}
The attached Grand Tensor source used in the latest audit is
\begin{center}
\path{Renewal_Grand_Tensor_Monograph_V067.tex},
\end{center}
with \(2804854\) bytes, \(68938\) lines, and SHA--256
\[
 \texttt{FACD058EC8BB210AD8B296C9D7F4D78A229D08D7A2AE378808700655CB8BC0FE}.
                                                                  \tag{1.3}
\]
The attached Einstein--Standard-Model programme source is
\begin{center}
\path{Renewal_Einstein_Standard_Model_Physical_Source_Metric_Duhamel_Ancestry_and_Quantum_Continuum_Programme_V30.tex},
\end{center}
with \(1262823\) bytes, \(32301\) lines, and SHA--256
\[
 \texttt{E0AE0E0990DE472021A1E90A9C10D41C3F747F28D97A5B1D8F1281852E3A0E69}.
                                                                  \tag{1.4}
\]
Statements inside those documents are treated as mathematical
source claims to audit, not as instructions governing these notes.
\end{record}

\section{Major mathematical results acquired so far}
\label{sec:v15v1-summary}

This is a status map.  Every result below retains the hypotheses
and claim firewall of its archived proof.

\subsection{Finite regulator and quantum-field structure}

\begin{enumerate}[leftmargin=2.2em]
\item Finite renewal regulators, common action records, Schur and
      Feshbach reductions, determinant and Pfaffian lines, shell
      normalizations, and source writers have been separated into
      explicit algebraic, analytic, and provenance rows.
\item Standard Model representation arithmetic cancels the
      universal local anomaly polynomial for the declared
      representation.  Regulator remainders, determinant-line
      topology, torsion, and global anomalies remain independent
      checks.
\item Finite cochain Hodge decompositions split regulated anomalies
      into harmonic obstructions and exact counterterms.  A harmonic
      class in a proved integral lattice vanishes when it is below
      the corresponding systole.
\item Exact or vanishing-defect Ward identities pass through
      normalized complex shell martingales and convergent joint
      states on declared local cores.  Closed BRST derivations pass
      to graph limits under the recorded domain and predual
      convergence hypotheses.
\item Relative trace-class and Hilbert--Schmidt control transports
      determinant moduli and phases on zero-free paths.  Local phase
      forms, cycle holonomy, real Pfaffian orientation, and
      scalarization of line bundles have been kept distinct.
\item A joint classical--quantum state has a trace-class
      conditional-density disintegration when a compact-resolvent
      energy moment prevents quantum mass from escaping.
      Compatible local restrictions construct a quasi-local state.
\item All-order positive moment data are necessary for the
      interacting state; no fixed finite correlation order
      determines it.  Dense test-function and graph seminorm
      estimates construct operator-valued distributions on a common
      invariant domain.
\item Reflection positivity on the physical algebra reconstructs
      the OS Hilbert space, time semigroup, and Hamiltonian under the
      archived domain hypotheses.  Dense covariance and locality
      extend only when uniform identity and causal-margin controls
      are also supplied.
\end{enumerate}

\subsection{Gravity and the ordered Einstein interface}

\begin{enumerate}[leftmargin=2.2em]
\item The unmodified Euclidean Einstein--Hilbert action is unbounded
      below along bounded fixed-volume conformal oscillations.
      Coercive real references and complex contours therefore carry
      separate normalization, divisor-avoidance, homology, reality,
      and positivity requirements.
\item Projectively consistent Gaussian gravitational laws and
      summable-defect corrections have quantitative cofinal limits.
      Uniform regularity moments and an action floor give quotient
      tightness and Wasserstein subsequences without requiring a
      global gauge section.
\item A countable distance-coordinate ledger reconstructs the
      metric quotient topology and its Borel laws.  A determining
      observable ledger upgrades subsequential equality to
      full-sequence convergence.
\item For a noncentral coframe \(E\), two-sided inverse floors give
      controlled derivatives and landing of
      \(G=E^*\eta E\), \(E^{-1}\), and their jets.
\item Ordered Cartan calculus gives exact curvature and torsion
      identities.  Finite-stage Bianchi failures factor through the
      actual nilpotence and Leibniz defects.
\item The coframe-induced connection \(E^{-1}\dd E\) is metric
      compatible but flat.  Every metric-compatible connection is
      \(E^{-1}\dd E+G^{-1}H\) with \(H^*=-H\); torsion-freeness is a
      linear Cartan selector equation.
\item At a central orthonormal anchor the selector is explicitly
      invertible.  A Neumann margin transports uniqueness and
      cofinal landing to a noncommutative neighborhood.
\item The symmetric ordered Einstein word has a divergence defect
      which factors into Ricci, scalar, metricity, and uncontracted
      Bianchi defects.  A common limiting stress tensor and support
      of the joint state on the resulting ordered Einstein equation
      remain unproved.
\end{enumerate}

\subsection{Literal RPESM regulator and the spectral boundary}

\begin{enumerate}[leftmargin=2.2em]
\item The Grand Tensor constructs an existential, positive,
      reflection-compatible, exactly projective RPESM regulator
      with compact gauge variables, an unbounded quartically
      controlled Higgs field, fermion and ghost records, and a
      stationary Markov--Renewal slab.
\item Its global refresh adds the exact centered spectral shift
      \(\rho\), lifts relative-sector zero modes, and cancels from
      any honest comparison with a refresh-free physical
      Hamiltonian.  It is therefore not a local mass gap.
\item The exact conditional shell skeleton is
      \(P_{n+1}=J_nP_nJ_n^*\).  It has block spectrum
      \(P_n\oplus0\), inherits the old Dirichlet gap, and annihilates
      every centered new-detail observable.
\item The resulting zero transfer sector cannot equal
      \(e^{-aH}\) for a finite self-adjoint Hamiltonian on the full
      detail space.  Scalar laziness cannot retain both finite
      nonzero coarse and detail energies on one continuum clock.
\item A measure-preserving conditional-fiber trivialization gives
      an injective faithful old transfer.  Adding a detail generator
      produces the exact product gap
      \(\min\{\gamma_{\rm old},\gamma_{\rm det}\}\).  A
      noncommuting symmetric split has an explicit
      Friedrichs-angle gap and commutator debit.
\item Conditional Moser transport constructs the needed
      trivialization on uniformly convex Euclidean fibers and gives
      a mixed-derivative influence bound.  Compact gauge topology,
      permitted nonconvex Higgs terms, elliptic nonlocality, and
      cutoff holonomy prevent its unqualified use on the full shell.
\item Every local conditional-expectation update in the source has
      exact conditional fiber gap one.  Its global refresh-free
      Dirichlet form is the sum of the conditional variances.
      Sectorwise irreducibility and approximate tensorization are
      the missing assembly rows.
\item Uniform total-variation influence fails maximally for an
      allowed unbounded quartic linear channel.  A high-probability
      bounded-field set alone cannot yield a spectral gap; an
      independent tail-return or Lyapunov estimate is necessary.
\end{enumerate}

\section{Current dependency boundary}
\label{sec:v15v1-boundary}

The full continuum goal currently depends on several interacting
rows:
\[
\begin{array}{p{0.29\linewidth}|p{0.61\linewidth}}
\textbf{row}&\textbf{missing evidence}\\ \hline
\text{refresh-free dynamics}
 &constraint-preserving local blocks, sector irreducibility, and a
  cutoff-uniform local-to-global estimate\\
\text{unbounded Higgs control}
 &a weighted influence contraction or a Lyapunov return inequality
  which remains uniform outside compact field sets\\
\text{compact gauge transport}
 &a finite physical chart atlas with covariant overlap maps and
  uniform local constants\\
\text{fermionic physical measure}
 &selected chiral crossing positivity, determinant/Pfaffian phase
  and orientation, and common-cutoff compatibility\\
\text{OS Hamiltonian limit}
 &tightness of low-energy states, norm-resolvent or equivalent
  isolated-vacuum control, locality, and the physical clock\\
\text{Einstein source}
 &uniform stress/contact bounds, Ward descent, metric-jet landing,
  and support on one ordered Einstein equation\\
\text{continuum identification}
 &one common regulator populating all rows, with cutoff removal and
  the required massive and massless spectral sectors.
\end{array}                                                     \tag{1.5}
\]
The present note attacks the second row.  Total variation is too
strong because widely separated external fields can make
conditional laws nearly disjoint.  Their locations can still vary
regularly in transport distance.

\section{Quartic conditional laws have finite weighted influence}
\label{sec:v15v1-quartic}

For \(\lambda>0\), \(b\in\mathbb R\), and \(t\in\mathbb R\), let
\[
 \nu_t(\dd h)
 =Z_t^{-1}\exp\{-\lambda h^4+bh^2-th\}\,\dd h.                \tag{1.6}
\]

\begin{lemma}[Uniform variance under every linear tilt]
\label{lem:v15v1-uniform-variance}
The constant
\[
 C_{\lambda,b}:=\sup_{t\in\mathbb R}
 \operatorname{Var}_{\nu_t}(h)                               \tag{1.7}
\]
is finite.
\end{lemma}

\begin{proof}
On every compact \(t\)-interval, the density and all polynomial
moments are dominated by
\[
 C_T\exp\{-(\lambda/2)h^4\}.
\]
Dominated convergence therefore makes
\(t\mapsto\operatorname{Var}_{\nu_t}(h)\) continuous.

It remains to control \(|t|\to\infty\).  Put
\[
 a=\left(\frac{|t|}{4\lambda}\right)^{1/3},
 \qquad h=aq,\qquad s=\operatorname{sgn}t.
\]
Up to normalization, the rescaled exponent is
\[
 -a^4\lambda(q^4+4sq)+ba^2q^2.                              \tag{1.8}
\]
The function \(q\mapsto q^4+4sq\) has the unique minimizer
\(-s\), with second derivative \(12\) there, and grows quartically
away from that point.  Splitting the \(q\)-integrals into a fixed
neighborhood of \(-s\) and its complement gives the standard
Laplace bound
\[
 \mathbb E_{\nu_t}\bigl[(q-\mathbb E q)^2\bigr]=O(a^{-4}).
\]
The lower-order term \(ba^2q^2\) is uniform on the local
\(a^{-2}\) fluctuation scale and does not alter that bound.
Consequently
\[
 \operatorname{Var}_{\nu_t}(h)
 =a^2\operatorname{Var}_{\nu_t}(q)=O(a^{-2})\longrightarrow0.
\]
Continuity on compact intervals and decay at both ends prove
\((1.7)\).
\end{proof}

The previous cycle proved that
\(\sup_{t,s}\|\nu_t-\nu_s\|_{\rm TV}=1\).  The next theorem shows
that the same family is uniformly Lipschitz in \(W_1\).

\begin{theorem}[Exact one-dimensional Wasserstein response]
\label{thm:v15v1-w1-response}
For all \(s,t\in\mathbb R\),
\[
 \boxed{
 W_1(\nu_s,\nu_t)
 =\left|\int_s^t\operatorname{Var}_{\nu_r}(h)\,\dd r\right|
 \le C_{\lambda,b}|t-s|.}                                    \tag{1.9}
\]
\end{theorem}

\begin{proof}
If \(t>s\), the likelihood ratio
\[
 \frac{\dd\nu_t}{\dd\nu_s}(h)
 =\frac{Z_s}{Z_t}e^{-(t-s)h}
\]
is decreasing.  Hence \(\nu_t\) is stochastically smaller than
\(\nu_s\).  For stochastically ordered probability laws on
\(\mathbb R\), the monotone quantile coupling is optimal for
\(|h-k|\) and
\[
 W_1(\nu_s,\nu_t)
 =\mathbb E_{\nu_s}h-\mathbb E_{\nu_t}h.                     \tag{1.10}
\]
Differentiation under the integral, justified by quartic
domination, gives
\[
 \frac{\dd}{\dd r}\mathbb E_{\nu_r}h
 =-\operatorname{Var}_{\nu_r}(h).                            \tag{1.11}
\]
Integrating \((1.11)\) proves the identity in \((1.9)\); the
uniform variance bound proves the inequality.  The case \(s>t\)
follows by symmetry.
\end{proof}

\section{Weighted path coupling for the quartic heat-bath chain}
\label{sec:v15v1-path-coupling}

Let \(I\) be a finite set of Higgs blocks and consider
\[
 \mu(\dd x)=Z^{-1}\exp\left\{
 -\sum_{i\in I}(\lambda_ix_i^4-b_ix_i^2)
 -\sum_{\{i,j\}}J_{ij}x_ix_j
 \right\}\prod_{i\in I}\dd x_i,                              \tag{1.12}
\]
where \(\lambda_i>0\), \(b_i\in\mathbb R\), and
\(J_{ij}=J_{ji}\).  Its conditional law at block \(i\) is
\[
 \nu_{i,t_i(x)},\qquad
 t_i(x)=\sum_{j\ne i}J_{ij}x_j,                              \tag{1.13}
\]
with variance constant
\[
 C_i=\sup_t\operatorname{Var}_{\nu_{i,t}}(h)<\infty.
\]
For update rates \(r_i>0\), define the continuous-time heat-bath
generator
\[
 (Lf)(x)=\sum_{i\in I}r_i
 \left[
 \int f(x^{i\to h})\nu_{i,t_i(x)}(\dd h)-f(x)
 \right].                                                    \tag{1.14}
\]
This generator is reversible with respect to \(\mu\).

For weights \(w_i>0\), use the transport metric
\[
 d_w(x,y)=\sum_{i\in I}w_i|x_i-y_i|                           \tag{1.15}
\]
and define
\[
 \kappa_w
 =\min_{j\in I}
 \left[
 r_j-\frac1{w_j}
 \sum_{i\ne j}r_iw_iC_i|J_{ij}|
 \right].                                                    \tag{1.16}
\]

\begin{theorem}[Quartic heat-bath Wasserstein contraction]
\label{thm:v15v1-path-contraction}
If \(\kappa_w>0\), the heat-bath semigroup \(P_\tau=e^{\tau L}\)
satisfies
\[
 \boxed{
 W_{1,d_w}(\delta_xP_\tau,\delta_yP_\tau)
 \le e^{-\kappa_w\tau}d_w(x,y)}
 \qquad(\tau\ge0).                                           \tag{1.17}
\]
More generally, for probability laws \(\alpha,\beta\) with finite
first \(d_w\)-moment,
\[
 W_{1,d_w}(\alpha P_\tau,\beta P_\tau)
 \le e^{-\kappa_w\tau}W_{1,d_w}(\alpha,\beta).                \tag{1.18}
\]
In particular, \(\mu\) is the unique invariant law with finite
first moment and convergence to it is exponential in this
transport distance.
\end{theorem}

\begin{proof}
Couple the two chains by using the same rate-\(r_i\) clock at site
\(i\).  At a ring, couple the two new conditional values by their
monotone quantiles.  Theorem~\ref{thm:v15v1-w1-response} gives
\[
 \begin{split}
 \mathbb E|X_i'-Y_i'|
 &=
 W_1\bigl(\nu_{i,t_i(X)},\nu_{i,t_i(Y)}\bigr)\\
 &\le C_i|t_i(X)-t_i(Y)|\\
 &\le C_i\sum_{j\ne i}|J_{ij}|\,|X_j-Y_j|.                  \tag{1.19}
 \end{split}
\]
The coupled generator applied to \(d_w\) is therefore at most
\[
 \begin{split}
 &\sum_i r_i\left[
 -w_i|X_i-Y_i|
 +w_iC_i\sum_{j\ne i}|J_{ij}|\,|X_j-Y_j|
 \right]\\
 &=
 \sum_j\left[
 -r_jw_j+\sum_{i\ne j}r_iw_iC_i|J_{ij}|
 \right]|X_j-Y_j|\\
 &\le-\kappa_wd_w(X,Y).                                      \tag{1.20}
 \end{split}
\]
Gronwall's inequality proves \((1.17)\).  Taking an initial optimal
coupling and integrating proves \((1.18)\).  Since the quartic law
has finite first moments and is invariant by the heat-bath
identity, contraction gives uniqueness and convergence.
\end{proof}

\begin{assessment}[What V1 gains and what it does not]
\label{ass:v15v1-gain}
Theorem~\ref{thm:v15v1-path-contraction} remains informative in
the regime where the total-variation coefficient is one.  It gives
a finite, explicit dependence matrix
\[
 {\cal C}_{ij}=C_i|J_{ij}|                                   \tag{1.21}
\]
and reduces transport contraction to the weighted inequality
\((1.16)\).

This is not yet an \(L^2(\mu)\) spectral-gap theorem for the
continuous unbounded heat-bath chain.  A Wasserstein contraction
controls Lipschitz observables; an upgrade to all square-integrable
observables requires a gradient regularization, conditional
Poincare, or variance-tensorization argument.  Nor does the scalar
prototype include compact gauge variables, gauge constraints,
multi-component Higgs fields, fermionic weights, or cutoff scaling.
Those are the next rows rather than hidden assumptions.
\end{assessment}

\begin{decision}[V2 transport-to-variance upgrade]
\label{dec:v15v1-next}
Use the exact unit conditional heat-bath gaps from the frozen V99
proof together with the weighted matrix \((1.21)\).  Derive a
variance tensorization theorem under an explicit weighted
interdependence norm, or prove by counterexample that the
Wasserstein row alone is insufficient.  Track the dependence of
\(C_i\) on the quartic and quadratic coefficients and on the
lattice scaling.  Only after this \(L^2\) bridge is valid should the
compact gauge charts and constraint-preserving blocks be added.
\end{decision}

\begin{firewall}[Claim boundary after fifteenth-cycle V1]
\label{fire:v15v1}
V1 proves uniform \(W_1\) response of every one-dimensional
quartic conditional law and exponential weighted path coupling for
a finite scalar quartic heat-bath system under \(\kappa_w>0\).
It does not prove a cutoff-uniform RPESM spectral gap, a continuum
Hamiltonian, or the full Standard Model--Einstein continuum.
\end{firewall}


\section{Fifteenth-cycle V2: the exact two-block \(L^2\) bridge}
\label{sec:v15v2}

V1 obtained weighted \(W_1\) contraction for scalar quartic
heat-bath dynamics but did not infer an \(L^2\) gap.  The two-block
case admits an exact answer in terms of maximal correlation.

Let \((X,Y)\) have law \(\mu\), put
\[
 {\cal H}=L^2(\mu),\qquad
 P_Xf=\mathbb E[f\mid X],\qquad
 P_Yf=\mathbb E[f\mid Y],
\]
and let \(E_0f=\mathbb E_\mu f\).  The generator which updates one
coordinate at unit rate is
\[
 A=(I-P_X)+(I-P_Y).                                          \tag{2.1}
\]
Here \(P_X\) is the heat-bath update of \(Y\), while \(P_Y\)
updates \(X\).

\begin{definition}[Maximal correlation]
\label{def:v15v2-maxcorr}
The Hirschfeld--Gebelein--Renyi maximal correlation is
\[
 c(\mu)=
 \sup\left\{
 |\mathbb E[f(X)g(Y)]|:
 \mathbb Ef=\mathbb Eg=0,\ 
 \mathbb E|f|^2=\mathbb E|g|^2=1
 \right\}.                                                   \tag{2.2}
\]
Equivalently,
\[
 c(\mu)=
 \bigl\|(P_X-E_0)(P_Y-E_0)\bigr\|.                           \tag{2.3}
\]
\end{definition}

\begin{theorem}[Exact two-block heat-bath gap]
\label{thm:v15v2-two-projection}
Assume the common \(X\)- and \(Y\)-measurable sigma algebra is
trivial modulo \(\mu\).  Then
\[
 \boxed{\operatorname{gap}A=1-c(\mu).}                       \tag{2.4}
\]
For update rates \(r_X,r_Y>0\), the corresponding generator obeys
the lower bound
\[
 \operatorname{gap}\bigl(r_X(I-P_X)+r_Y(I-P_Y)\bigr)
 \ge\min\{r_X,r_Y\}\,[1-c(\mu)].                             \tag{2.5}
\]
\end{theorem}

\begin{proof}
On the centered space \({\cal H}_0=(I-E_0){\cal H}\), the
restrictions of \(P_X\) and \(P_Y\) are orthogonal projections
onto the centered \(X\)- and \(Y\)-measurable subspaces.  Their
Friedrichs cosine is \((2.3)\).  The two-projection canonical
decomposition gives, on every nontrivial two-dimensional angle
block with cosine \(s\),
\[
 \operatorname{spec}(2I-P_X-P_Y)=\{1-s,1+s\}.                \tag{2.6}
\]
The remaining orthogonal blocks have values \(1\) or \(2\).
Taking the lower spectral edge over the angle blocks, including an
approximating sequence if the supremum is not attained, gives
\((2.4)\).  Alternatively, the lower bound follows directly from
\[
 \|(I-P_X)f\|^2+\|(I-P_Y)f\|^2
 \ge(1-c(\mu))\|f\|^2,
\]
and the canonical decomposition proves sharpness.  Multiplying the
two nonnegative terms by their rates proves \((2.5)\).
\end{proof}

We apply this identity to the coupled quartic pair
\[
 \mu_J(\dd x,\dd y)
 =Z_J^{-1}\exp\{
 -\lambda x^4+bx^2-\rho y^4+dy^2-Jxy
 \}\,\dd x\,\dd y,                                           \tag{2.7}
\]
where \(\lambda,\rho>0\) and \(b,d,J\in\mathbb R\).

\begin{lemma}[Hilbert--Schmidt conditional operator]
\label{lem:v15v2-hs}
Let \(p(x,y)\) be the density in \((2.7)\), with marginals
\(p_X,p_Y\).  The conditional expectation operator
\[
 (Kg)(x)=\mathbb E[g(Y)\mid X=x]
\]
from \(L^2(p_Y\dd y)\) to \(L^2(p_X\dd x)\) is
Hilbert--Schmidt.  Its squared Hilbert--Schmidt norm is
\[
 \|K\|_{\rm HS}^2
 =\int_{\mathbb R^2}
   \frac{p(x,y)^2}{p_X(x)p_Y(y)}\,\dd x\,\dd y<\infty.        \tag{2.8}
\]
\end{lemma}

\begin{proof}
Put
\[
 V_X(x)=\lambda x^4-bx^2,\qquad
 V_Y(y)=\rho y^4-dy^2.
\]
For fixed \(x\), restrict the marginal integral defining \(p_X(x)\)
to \(|y|\le(1+|x|)^{-1}\).  On this interval,
\(|Jxy|\le|J|\), while \(V_Y\) is bounded above uniformly.
Consequently there is \(a_X>0\) such that
\[
 p_X(x)\ge
 a_X\frac{e^{-V_X(x)}}{1+|x|}.                               \tag{2.9}
\]
The symmetric argument gives
\[
 p_Y(y)\ge
 a_Y\frac{e^{-V_Y(y)}}{1+|y|}.                               \tag{2.10}
\]
Equations \((2.9)\)--\((2.10)\) imply
\[
 \frac{p(x,y)^2}{p_X(x)p_Y(y)}
 \le C(1+|x|)(1+|y|)
 e^{-V_X(x)-V_Y(y)-2Jxy}.                                   \tag{2.11}
\]
The positive quartic terms dominate all quadratic and bilinear
terms, so the right-hand side is integrable.  Formula \((2.8)\)
is the usual kernel formula for the conditional expectation
operator.
\end{proof}

\begin{theorem}[Strict \(L^2\) gap for every finite quartic pair]
\label{thm:v15v2-quartic-pair}
For every choice of parameters in \((2.7)\),
\[
 c(\mu_J)<1
\]
and hence the two-coordinate refresh-free heat-bath generator has
the strictly positive gap
\[
 \boxed{\operatorname{gap}A=1-c(\mu_J)>0.}                   \tag{2.12}
\]
This includes the nonconvex double-well regime and requires no
global refresh.
\end{theorem}

\begin{proof}
Lemma~\ref{lem:v15v2-hs} makes the conditional operator compact.
Its restriction from the centered \(Y\)-space to the centered
\(X\)-space is also compact and has norm \(c(\mu_J)\).
If that norm were one, compactness would give centered unit
vectors \(f(X)\) and \(g(Y)\) for which
\[
 |\mathbb E[f(X)g(Y)]|=1.
\]
Equality in Cauchy--Schwarz would imply
\(f(X)=\zeta g(Y)\) almost surely for some \(|\zeta|=1\).
The density in \((2.7)\) is strictly positive at every
\((x,y)\).  Fubini's theorem then forces \(f(x)\) and \(g(y)\) to
be almost-everywhere constant, contradicting their centered unit
norms.  Therefore \(c(\mu_J)<1\), and
Theorem~\ref{thm:v15v2-two-projection} proves \((2.12)\).
\end{proof}

\begin{assessment}[Finite positivity versus a cofinal floor]
\label{ass:v15v2-uniformity}
Theorem~\ref{thm:v15v2-quartic-pair} closes the existence of a
finite two-block \(L^2\) gap.  It does not give a useful numerical
lower bound on \(1-c(\mu_J)\), nor does it prove that the gap stays
positive when:
\begin{enumerate}
\item external linear tilts from the remaining lattice sites become
      unbounded;
\item the number of blocks grows;
\item the lattice coefficients and physical clock scale with the
      cutoff;
\item compact gauge and constraint variables are included.
\end{enumerate}
Compactness at each fixed regulator cannot replace a
regulator-uniform maximal-correlation bound.

The result also explains the logical relation to V1.  Wasserstein
response controls conditional means and Lipschitz observables.
The \(L^2\) gap is controlled by the full conditional expectation
operator on every centered square-integrable observable.  The two
constants are related in special regular families but are not
identical.
\end{assessment}

\begin{decision}[V3 tilted pair and multi-block compiler]
\label{dec:v15v2-next}
Add arbitrary external linear tilts \(sx+ty\) to \((2.7)\).
Determine whether the centered conditional operator has a
uniform norm
\[
 \sup_{s,t\in\mathbb R}c(\mu_{J;s,t})<1.                     \tag{2.13}
\]
If it does, prove the limit at large tilts and compactness on the
remaining parameter set.  Then assemble a multi-block
maximal-correlation matrix and derive a volume-uniform heat-bath
gap under an explicit weighted row or column condition.  If
\((2.13)\) fails, retain the actual concentrating sequence and
return to a Lyapunov decomposition.
\end{decision}

\begin{firewall}[Claim boundary after V2]
\label{fire:v15v2}
Every finite scalar quartic pair with bilinear coupling has a
strictly positive two-block heat-bath gap.  No uniform value,
many-block tensorization, compact-gauge extension, or RPESM
continuum gap has yet been proved.
\end{firewall}

\section*{Version V2 append-only record}
\addcontentsline{toc}{section}{Version V2 append-only record}
The complete fifteenth-cycle V1 source was retained before this
continuation.  It had \(19638\) bytes, \(497\) lines, and SHA--256
\[
 \texttt{7C6F71B2EB84EB588A1D6E1A882F3A34C9CF6E7CBA35E4CB453CEE6AC601B4CB}.
\]
V2 proved the exact two-projection heat-bath gap formula and strict
finite \(L^2\) positivity for every bilinearly coupled scalar
quartic pair.


\section{Fifteenth-cycle V3: tilt-uniform maximal correlation for a quartic pair}
\label{sec:v15v3}

V2 reduced the two-block \(L^2\) gap to maximal correlation but
left its behavior under unbounded external fields open.  The
required uniformity follows in a quantitative small-coupling
region from one-dimensional Poincare estimates.

We use the convention that \(C_P(\pi)\) is the least constant in
\[
 \operatorname{Var}_\pi f
 \le C_P(\pi)\int|f'|^2\,\dd\pi.                              \tag{3.1}
\]

\begin{lemma}[Curvature-at-infinity gives tilt-uniform Poincare]
\label{lem:v15v3-curvature-infinity}
Let \(\{U_\theta\}_{\theta\in\Theta}\subset C^2(\mathbb R)\)
satisfy, for fixed \(K,m,R>0\),
\[
 U_\theta''(x)\ge-K\quad(x\in\mathbb R),\qquad
 U_\theta''(x)\ge m\quad(|x|\ge R).                           \tag{3.2}
\]
Assume \(e^{-U_\theta(x)-\ell x}\) is integrable for all
\(\theta,\ell\).  Then there is a finite constant
\(C_*(K,m,R)\) such that
\[
 \sup_{\theta\in\Theta,\ \ell\in\mathbb R}
 C_P\!\left(
 \frac{e^{-U_\theta(x)-\ell x}\dd x}
      {\int e^{-U_\theta(u)-\ell u}\dd u}
 \right)
 \le C_*(K,m,R).                                             \tag{3.3}
\]
\end{lemma}

\begin{proof}
Choose a smooth function \(\chi\), depending only on \(K,m,R\),
whose second derivative raises every curvature profile obeying
\((3.2)\) to a fixed positive floor \(m_0>0\), while
\[
 \operatorname{osc}\chi
 :=\sup\chi-\inf\chi\le D(K,m,R).                             \tag{3.4}
\]
One explicit construction takes \(\chi''=K+m_0\) on
\([-R,R]\), returns \(\chi'\) to zero in two collars where the
available exterior curvature \(m\) absorbs the negative return
curvature, and then keeps \(\chi\) affine-constant outside the
collars.  Enlarging the collars if necessary makes
\[
 (U_\theta+\ell x+\chi)''\ge m_0
\]
uniformly.  The size and oscillation of \(\chi\) depend only on
the three constants in \((3.2)\).

Bakry--Emery in one dimension gives Poincare constant at most
\(m_0^{-1}\) for the probability density proportional to
\(e^{-U_\theta-\ell x-\chi}\).  The two normalized densities
differ by the bounded factor \(e^\chi\).  The
Holley--Stroock comparison therefore gives
\[
 C_P(e^{-U_\theta-\ell x})
 \le e^{\operatorname{osc}\chi}m_0^{-1}.
\]
This proves \((3.3)\).
\end{proof}

\begin{corollary}[Uniform one-site quartic constants]
\label{cor:v15v3-one-site}
For \(\lambda>0\) and \(b\in\mathbb R\),
\[
 C_{\lambda,b}^{P}
 :=\sup_{t\in\mathbb R}C_P(\nu_t)<\infty,                    \tag{3.5}
\]
where \(\nu_t\) is the quartic law \((1.6)\).  In particular,
\[
 \sup_t\operatorname{Var}_{\nu_t}(h)
 \le C_{\lambda,b}^{P}.                                      \tag{3.6}
\]
\end{corollary}

\begin{proof}
The untilted potential
\(U(h)=\lambda h^4-bh^2\) has
\[
 U''(h)=12\lambda h^2-2b.
\]
It has a uniform lower bound and a positive curvature floor
outside a fixed compact interval.  Apply
Lemma~\ref{lem:v15v3-curvature-infinity}.  Equation \((3.6)\)
follows from \((3.1)\) with \(f(h)=h\).
\end{proof}

Add arbitrary external tilts to the V2 pair:
\[
 \mu_{s,t}(\dd x,\dd y)
 =Z_{s,t}^{-1}\exp\{
 -\lambda x^4+bx^2-\rho y^4+dy^2
 -Jxy-sx-ty
 \}\,\dd x\,\dd y.                                           \tag{3.7}
\]
Let
\[
 C_Y=C_{\rho,d}^{P},\qquad C_X=C_{\lambda,b}^{P}.             \tag{3.8}
\]

The \(X\)-marginal has density proportional to
\[
 e^{-U_t^{X}(x)-sx},\qquad
 U_t^{X}(x)
 =\lambda x^4-bx^2-\log Z_Y(t+Jx),                            \tag{3.9}
\]
where
\[
 Z_Y(q)=\int_{\mathbb R}
 e^{-\rho y^4+dy^2-qy}\,\dd y.
\]
Differentiation gives
\[
 (U_t^{X})''(x)
 =12\lambda x^2-2b
 -J^2\operatorname{Var}_{\nu^{Y}_{t+Jx}}(Y).                 \tag{3.10}
\]
By \((3.6)\),
\[
 (U_t^{X})''(x)
 \ge12\lambda x^2-2|b|-J^2C_Y.                               \tag{3.11}
\]
Lemma~\ref{lem:v15v3-curvature-infinity} therefore supplies a
finite constant \(\overline C_X\), depending only on
\(\lambda,b,J,C_Y\), such that
\[
 \sup_{s,t}C_P((\mu_{s,t})_X)\le\overline C_X.                \tag{3.12}
\]
Symmetrically there is \(\overline C_Y<\infty\), depending only
on \(\rho,d,J,C_X\), with
\[
 \sup_{s,t}C_P((\mu_{s,t})_Y)\le\overline C_Y.                \tag{3.13}
\]

\begin{theorem}[Tilt-uniform maximal-correlation bound]
\label{thm:v15v3-uniform-correlation}
For every \(s,t\in\mathbb R\),
\[
 \boxed{
 c(\mu_{s,t})
 \le |J|\min\left\{
 \sqrt{\overline C_XC_Y},
 \sqrt{\overline C_YC_X}
 \right\}.}                                                  \tag{3.14}
\]
Consequently, if
\[
 \chi_J:=
 |J|\min\left\{
 \sqrt{\overline C_XC_Y},
 \sqrt{\overline C_YC_X}
 \right\}<1,                                                 \tag{3.15}
\]
then the two-block heat-bath generator has the tilt-uniform gap
\[
 \boxed{\inf_{s,t}
 \operatorname{gap}A_{s,t}\ge1-\chi_J>0.}                    \tag{3.16}
\]
\end{theorem}

\begin{proof}
Let \(g(Y)\) be centered and square integrable under the
\(Y\)-marginal, and set
\[
 (Kg)(x)=\mathbb E_{\mu_{s,t}}[g(Y)\mid X=x].
\]
Only the conditional tilt \(t+Jx\) depends on \(x\).
Differentiation under the quartically dominated conditional
integral gives
\[
 (Kg)'(x)=-J\operatorname{Cov}(g(Y),Y\mid X=x).               \tag{3.17}
\]
Conditional Cauchy--Schwarz and \((3.6)\) yield
\[
 |(Kg)'(x)|^2
 \le J^2C_Y\operatorname{Var}(g(Y)\mid X=x).                  \tag{3.18}
\]
Integrating and using \((3.12)\),
\[
 \begin{split}
 \operatorname{Var}_X(Kg)
 &\le\overline C_X\mathbb E_X|(Kg)'(X)|^2\\
 &\le J^2\overline C_XC_Y
       \mathbb E_X\operatorname{Var}(g(Y)\mid X)\\
 &\le J^2\overline C_XC_Y\operatorname{Var}_Y(g).             \tag{3.19}
 \end{split}
\]
Thus the centered conditional-expectation operator has norm at
most \(|J|\sqrt{\overline C_XC_Y}\).  Interchanging \(X\) and
\(Y\) proves the second bound in \((3.14)\).
The exact two-projection formula \((2.4)\) then proves
\((3.16)\).
\end{proof}

\begin{assessment}[Scope of the uniformity]
\label{ass:v15v3-scope}
Equation \((3.16)\) is uniform over every external linear field
generated by the rest of a lattice.  This removes the specific
unbounded-tilt obstruction left by V2 in the explicit region
\(\chi_J<1\).

The theorem is still a two-block result.  Its constant
\(\overline C_X\) is finite and structurally explicit through the
convexification proof, but has not been optimized.  More
importantly, a many-site system cannot be treated by applying
\((3.16)\) independently to every edge: overlapping conditional
expectations share variables, and their correlation losses must be
assembled without counting the same variance repeatedly.
\end{assessment}

\begin{decision}[V4 bipartite block assembly]
\label{dec:v15v3-next}
On a bipartite interaction graph, group all even sites into one
block and all odd sites into the other.  Conditional on one
parity, the opposite scalar quartic sites are independent.  Use
tensorized conditional Poincare inequalities and the matrix of
cross couplings to bound the parity-to-parity conditional
expectation operator.  Derive a volume-uniform maximal-correlation
constant in a weighted operator norm, and hence an exact
two-parity heat-bath gap.  Record separately the locality cost of
updating an entire parity at once and the comparison needed to
return to single-site local clocks.
\end{decision}

\begin{firewall}[Claim boundary after V3]
\label{fire:v15v3}
A strictly positive two-block \(L^2\) gap is now uniform over
arbitrary external linear tilts under the explicit small-coupling
condition \((3.15)\).  No many-block, gauge-constrained, cofinal,
or continuum RPESM gap has yet been proved.
\end{firewall}

\section*{Version V3 append-only record}
\addcontentsline{toc}{section}{Version V3 append-only record}
The complete V2 source was retained before this continuation.  It
had \(27246\) bytes, \(722\) lines, and SHA--256
\[
 \texttt{EFEFAA9809F45D5E2AC826ED3C9F21D3ACD6D58AE3D728F694658E7369FB5457}.
\]
V3 proved tilt-uniform one-site Poincare control, bounded the
tilted pair's maximal correlation, and obtained an explicit
small-coupling heat-bath gap uniform over all external fields.


\section{Fifteenth-cycle V4: bipartite quartic assembly and a local refresh-free gap}
\label{sec:v15v4}

V3 controlled one coupled pair uniformly over all external linear
fields.  On a bipartite graph the full many-site conditional laws
retain exactly the product structure needed to assemble those
one-site inequalities without a volume-dependent bounded
perturbation.

Let \(G=(E\sqcup O,{\cal A})\) be a finite bipartite graph.  For
\(i\in E\) and \(j\in O\), let
\[
 V_i(x)=\lambda_ix^4-b_ix^2,\qquad
 W_j(y)=\rho_jy^4-d_jy^2,                                   \tag{4.1}
\]
with \(\lambda_i,\rho_j>0\), and let \(J=(J_{ij})\) vanish off
the edge set.  Include arbitrary linear fields \(s_i,t_j\), and
define
\[
 \begin{split}
 \mu(\dd x,\dd y)=Z^{-1}\exp\biggl\{&
 -\sum_{i\in E}\bigl(V_i(x_i)+s_ix_i\bigr)
 -\sum_{j\in O}\bigl(W_j(y_j)+t_jy_j\bigr)\\
 &-\sum_{i\in E,j\in O}J_{ij}x_iy_j
 \biggr\}\,\dd x\,\dd y.                                    \tag{4.2}
 \end{split}
\]
Conditional on \(Y=y\), the \(X_i\) are independent tilted
quartic laws; conditional on \(X=x\), the \(Y_j\) are independent
tilted quartic laws.

Let
\[
 C_E=\operatorname{diag}(C_i^E)_{i\in E},\qquad
 C_O=\operatorname{diag}(C_j^O)_{j\in O},                    \tag{4.3}
\]
where the tilt-uniform one-site Poincare constants from V3 obey
\[
 \operatorname{Var}(f(X_i)\mid Y)
 \le C_i^E\,\mathbb E(|f'(X_i)|^2\mid Y),
\]
and symmetrically for \(Y_j\).  Put
\[
 q=\|C_E^{1/2}JC_O^{1/2}\|_{2\to2}.                          \tag{4.4}
\]

Let
\[
 P_Ef=\mathbb E[f\mid Y],\qquad
 P_Of=\mathbb E[f\mid X],\qquad
 A_{\rm par}=(I-P_E)+(I-P_O).                                \tag{4.5}
\]

\begin{lemma}[Compact parity conditional operator]
\label{lem:v15v4-compact}
For every finite graph \(G\), the centered conditional operator
\[
 K:L^2_0(\mu_O)\longrightarrow L^2_0(\mu_E),
 \qquad
 Kg(x)=\mathbb E[g(Y)\mid X=x]
\]
is compact, and its norm \(c=\|K\|\) satisfies \(c<1\).
\end{lemma}

\begin{proof}
The proof of Lemma~\ref{lem:v15v2-hs} extends verbatim in finite
dimension.  Restricting each marginal integral to an inverse-linear
ball in the opposite variables gives polynomial lower bounds times
the uncoupled quartic density.  Hence the squared density ratio
\[
 \int
 \frac{p(x,y)^2}{p_E(x)p_O(y)}\,\dd x\,\dd y                 \tag{4.6}
\]
is bounded by a polynomial times a coercive quartic exponential
and is finite.  Thus \(K\) is Hilbert--Schmidt.

If its centered norm were one, compactness and equality in
Cauchy--Schwarz would give nonconstant centered functions
\(f(X)=g(Y)\) almost surely.  The density \((4.2)\) is strictly
positive on the full Euclidean product.  Fubini's theorem forces
both functions to be constant, a contradiction.
\end{proof}

\begin{theorem}[Parity maximal-correlation matrix bound]
\label{thm:v15v4-parity-correlation}
For the law \((4.2)\),
\[
 \boxed{c\le q.}                                             \tag{4.7}
\]
Consequently, if \(q<1\),
\[
 \boxed{\operatorname{gap}A_{\rm par}\ge1-q.}                \tag{4.8}
\]
The bound is uniform over all external fields \(s_i,t_j\).
\end{theorem}

\begin{proof}
If \(c=0\), there is nothing to prove.  Otherwise compactness gives
a centered unit vector \(g(Y)\) with
\[
 K^*Kg=c^2g,\qquad h(X):=Kg,\qquad \|h\|_2=c.                \tag{4.9}
\]
For fixed \(x\), differentiation of the product conditional law
gives
\[
 \partial_{x_i}h(x)
 =-\sum_{j\in O}J_{ij}
   \operatorname{Cov}(g(Y),Y_j\mid X=x).                     \tag{4.10}
\]
Write \(a_j(x)=\operatorname{Cov}(g(Y),Y_j\mid X=x)\).
For every vector \(v\),
\[
 \begin{split}
 |v^{\mathsf T}a(x)|^2
 &\le\operatorname{Var}(g(Y)\mid X=x)
      \operatorname{Var}\left(\sum_jv_jY_j\mid X=x\right)\\
 &\le\operatorname{Var}(g(Y)\mid X=x)\,
      v^{\mathsf T}C_Ov.                                    \tag{4.11}
 \end{split}
\]
The last step uses conditional independence and the one-site
Poincare inequality with the linear test function.  Duality in
\((4.11)\) gives
\[
 a(x)^{\mathsf T}C_O^{-1}a(x)
 \le\operatorname{Var}(g(Y)\mid X=x).                        \tag{4.12}
\]
Since \(\nabla h=-Ja\),
\[
 \nabla h^{\mathsf T}C_E\nabla h
 \le q^2\operatorname{Var}(g(Y)\mid X=x).                    \tag{4.13}
\]

Conditional on \(Y\), the \(X_i\) are independent, so their
Poincare inequalities tensorize:
\[
 \mathbb E\operatorname{Var}(h(X)\mid Y)
 \le\mathbb E[\nabla h^{\mathsf T}C_E\nabla h].              \tag{4.14}
\]
The singular-vector identities \((4.9)\) and total variance give
\[
 \begin{split}
 \mathbb E\operatorname{Var}(h(X)\mid Y)
 &=\|h\|_2^2-\|K^*h\|_2^2
   =c^2-c^4,\\
 \mathbb E\operatorname{Var}(g(Y)\mid X)
 &=1-\|Kg\|_2^2
   =1-c^2.                                                   \tag{4.15}
 \end{split}
\]
Combining \((4.13)\)--\((4.15)\) yields
\[
 c^2(1-c^2)\le q^2(1-c^2).
\]
Lemma~\ref{lem:v15v4-compact} gives \(c<1\), so division proves
\((4.7)\).  The exact two-projection gap formula \((2.4)\) proves
\((4.8)\).  The constants in \((4.3)\) are tilt-uniform, so the
external fields do not enter \(q\).
\end{proof}

The parity update changes an extensive set of sites at once.  The
next theorem transfers its gap to the sum of truly local
single-site heat baths.  Let
\[
 P_i f=\mathbb E[f\mid\hbox{all coordinates except }x_i],
 \quad
 P_j f=\mathbb E[f\mid\hbox{all coordinates except }y_j],
\]
and set
\[
 A_{\rm loc}
 =\sum_{i\in E}r_i(I-P_i)+\sum_{j\in O}r_j(I-P_j),
 \qquad r_*=\min_{k\in E\cup O}r_k.                          \tag{4.16}
\]

\begin{theorem}[Volume-uniform local heat-bath gap]
\label{thm:v15v4-local-gap}
If \(q<1\), then
\[
 \boxed{\operatorname{gap}A_{\rm loc}
 \ge r_*(1-q).}                                              \tag{4.17}
\]
If the degrees, one-site quartic parameters, update rates, and
weighted operator norm in \((4.4)\) remain uniform along a
sequence of bipartite lattices, the lower bound is independent of
the number of sites and uses no global refresh.
\end{theorem}

\begin{proof}
Conditional on \(Y\), the \(X_i\) are independent.  The
Efron--Stein inequality for that product conditional law gives
\[
 \mathbb E\operatorname{Var}(f\mid Y)
 \le\sum_{i\in E}
 \mathbb E\operatorname{Var}
 (f\mid\hbox{all coordinates except }x_i).                  \tag{4.18}
\]
The analogous inequality holds for the odd parity.  Therefore, in
quadratic-form order,
\[
 A_{\rm loc}\succeq r_*A_{\rm par}.                          \tag{4.19}
\]
Equation \((4.17)\) follows from \((4.8)\).
\end{proof}

\begin{corollary}[A simple bounded-degree criterion]
\label{cor:v15v4-degree}
If \(C_i^E\le C_E^*\), \(C_j^O\le C_O^*\), and the unweighted
coupling matrix satisfies \(\|J\|_{2\to2}\le J_*\), then
\[
 q\le J_*\sqrt{C_E^*C_O^*}.                                  \tag{4.20}
\]
For a regular nearest-neighbor bipartite graph with edge coupling
magnitude at most \(j_*\) and degree \(D\), one may use
\(J_*\le Dj_*\).  Hence
\[
 Dj_*\sqrt{C_E^*C_O^*}<1                                    \tag{4.21}
\]
is a volume-independent sufficient condition for \((4.17)\).
\end{corollary}

\begin{assessment}[First many-site refresh-free theorem]
\label{ass:v15v4-result}
V4 proves a genuine many-site, unbounded-field, refresh-free
spectral gap for a scalar bipartite quartic Gibbs law in the
small-coupling region \(q<1\).  The proof neither truncates the
Higgs field nor uses total variation, and the single-site
generator is spatially local.

The literal RPESM action is more general than \((4.2)\).  It
contains multi-component Higgs fields, compact gauge links,
geometry and route variables, higher local interactions, physical
constraints, and coefficient scaling with the lattice.  Before
\((4.17)\) can be imported, one must show that its conditional
block law has:
\begin{enumerate}
\item a bipartite or finitely colored conditional product
      decomposition;
\item uniform block Poincare matrices;
\item a cross-derivative operator with norm below one;
\item constraint-preserving local updates and the correct physical
      clock.
\end{enumerate}
\end{assessment}

\begin{decision}[V5 compulsory audit]
\label{dec:v15v4-next}
At V5 inspect the newest YM and NS notes for finite chart covers,
block Poincare estimates, colored decompositions, or physical
clock scalings.  Then continue in V6 by extending
Theorem~\ref{thm:v15v4-local-gap} from scalar sites to
finite-dimensional Higgs blocks and compact gauge-link blocks.
Use independently certified compact charts with overlap, as
suggested by the YM V30 atlas method, and do not impose one global
convex coordinate on the gauge group.
\end{decision}

\begin{firewall}[Claim boundary after V4]
\label{fire:v15v4}
The local gap \((4.17)\) is proved for the finite scalar bipartite
quartic model under \(q<1\).  The source has not yet been shown to
lie in this region, and no gauge-constrained or continuum RPESM
gap follows without the four additional rows above.
\end{firewall}

\section*{Version V4 append-only record}
\addcontentsline{toc}{section}{Version V4 append-only record}
The complete V3 source was retained before this continuation.  It
had \(35315\) bytes, \(971\) lines, and SHA--256
\[
 \texttt{2323CF681F818D5C8B9C0A4BE891614BEDA5318DB86CD8E48DD93D4655DA96FC}.
\]
V4 bounded parity maximal correlation by a cross-coupling
operator norm and transferred the resulting two-parity gap to the
spatially local single-site heat-bath generator.


\section{Fifteenth-cycle V5: compulsory YM/NS audit after the scalar local-gap theorem}
\label{sec:v15v5}

V4 produced a volume-uniform refresh-free gap for a scalar
bipartite quartic model under the cross-coupling condition
\(q<1\).  This checkpoint tests whether the companion programmes
supply any of the block constants needed to extend that theorem to
the literal RPESM shell.

\begin{record}[V5 companion checkpoint]
\label{rec:v15v5-companions}
The newest Yang--Mills note is
\begin{center}
\path{Renewal_Yang_Mills_Mass_Gap_Research_Notes_V32.tex},
\end{center}
with \(261923\) bytes, \(7139\) lines, and SHA--256
\[
 \texttt{D9FD37946B54B96970086B85E9ADA651E17C1EB78EAD86E6DEFE68657AB5A33B}.
                                                                  \tag{5.1}
\]
The Navier--Stokes note remains
\begin{center}
\path{Renewal_NS_Stability_Research_Notes_V26.tex},
\end{center}
with \(278283\) bytes, \(5319\) lines, and SHA--256
\[
 \texttt{82C887F8C2C69E4F28FAD7C3DC8DE5B9099CB5A4BF431EBA1FFF3E91935978AD}.
                                                                  \tag{5.2}
\]
\end{record}

Yang--Mills V31 and V32 add two independently certified rational
splittings to the V28--V30 construction.  Four radius-scale tubes
now form an exact connected corridor along the physical diagonal,
and the V32 corridor crosses the diagonal face of a quarter cell.
Every consecutive overlap is proved by rational inequalities.
The notes explicitly leave the off-diagonal covering problem open.

\begin{assessment}[What the YM corridor contributes]
\label{ass:v15v5-ym}
The corridor strengthens the method imported at the fourteenth
cycle close:
\begin{enumerate}
\item a failed global splitting can be replaced by locally constant
      certified splittings;
\item pairwise overlap edges suffice to make the accepted region
      connected;
\item coverage must still be proved in every independent direction;
      a diagonal crossing does not cover a cell.
\end{enumerate}
This is relevant to a future compact-gauge atlas.  However the YM
objects are momentum-space Wilson-pencil matrices.  V32 contains
no gauge-fiber coordinate map, conditional Poincare inequality,
colored Gibbs decomposition, heat-bath update, or physical clock
comparison.  Its rational radii and matrix floors cannot enter the
constant \(q\) of \((4.4)\).
\end{assessment}

Navier--Stokes V26 remains a deterministic Oseen-kernel derivative
calculation.  It has no common configuration measure or block
generator with \((4.2)\).

\begin{theorem}[V5 audit verdict]
\label{thm:v15v5-verdict}
At checkpoints \((5.1)\)--\((5.2)\), neither companion supplies:
\begin{enumerate}
\item a vector-Higgs conditional Poincare matrix;
\item a compact-gauge block Poincare constant;
\item a finite coloring of the literal RPESM interaction
      hypergraph;
\item a constraint-preserving heat-bath block cover;
\item a lattice-to-physical clock scaling.
\end{enumerate}
The scalar theorem V4 is therefore not numerically modified.
The finite-atlas overlap rule is retained as a construction method
for the compact sector.
\end{theorem}

\begin{proof}
The new YM conclusions are exact local Loewner inequalities for a
finite matrix pencil on momentum tubes.  NS V26 concerns
derivatives of a continuum kernel.  Neither source defines any of
the conditional laws, gradient matrices, or Markov forms listed
above, so their conclusions cannot populate those rows.
\end{proof}

\begin{decision}[V6 block-valued extension]
\label{dec:v15v5-next}
Continue with a finite-dimensional block theorem.  Replace scalar
quartic sites by Euclidean Higgs blocks with tilt-uniform block
Poincare matrices, and compact link blocks with intrinsic
Poincare constants.  Prove the parity maximal-correlation estimate
using cross-gradient operator norms.  Keep the theorem coordinate
free on compact blocks so that a later finite chart atlas is
needed only to verify constants, not to define the generator.
Then state the exact colored-block condition required by the
literal source.
\end{decision}

\begin{firewall}[Claim boundary after V5]
\label{fire:v15v5}
The audit imports no companion numerical constant and makes no
full-cell YM or RPESM gap claim.  The next extension rests on new
block Poincare and cross-gradient hypotheses.
\end{firewall}

\section*{Version V5 append-only record}
\addcontentsline{toc}{section}{Version V5 append-only record}
The complete V4 source was retained before this continuation.  It
had \(44703\) bytes, \(1248\) lines, and SHA--256
\[
 \texttt{1E2E666938DF83647EB7ABD8F7CA1A8EC0E4C61F3E681444AE13F2783899B418}.
\]
V5 performed the compulsory companion audit, retained the exact
finite-atlas overlap method, and found no block-gap or clock
constant applicable to the RPESM shell.


\section{Fifteenth-cycle V6: intrinsic block Poincare assembly for Higgs and compact gauge variables}
\label{sec:v15v6}

V4 used scalar derivatives only for the maximal-correlation
estimate.  The same proof is intrinsic on finite-dimensional
Riemannian blocks and therefore does not require one global
coordinate chart on a compact gauge group.

Let \(G=(E\sqcup O,{\cal A})\) be bipartite.  At each vertex \(k\)
let \(M_k\) be a complete finite-dimensional Riemannian manifold
with reference measure \(m_k\).  Euclidean Higgs blocks and compact
Lie-group link blocks are both allowed.  Consider
\[
 \mu(\dd z)=Z^{-1}\exp\left\{
 -\sum_kV_k(z_k)
 -\sum_{i\in E,j\in O}\Phi_{ij}(z_i,z_j)
 \right\}\prod_km_k(\dd z_k),                                \tag{6.1}
\]
where the density is positive and normalizable.  There are no
same-parity interaction terms.

Assume every one-block conditional law has a uniform intrinsic
Poincare constant:
\[
 \operatorname{Var}(f(Z_i)\mid Z_O)
 \le C_i^E\,\mathbb E(|\nabla_if|^2\mid Z_O),                 \tag{6.2}
\]
\[
 \operatorname{Var}(g(Z_j)\mid Z_E)
 \le C_j^O\,\mathbb E(|\nabla_jg|^2\mid Z_E).                 \tag{6.3}
\]
Let \({\cal T}_E=\bigoplus_{i\in E}TM_i\) and
\({\cal T}_O=\bigoplus_{j\in O}TM_j\).  At a configuration
\((x,y)\), define the weighted mixed-Hessian operator
\[
 {\mathbb B}(x,y):{\cal T}_{E,x}\longrightarrow{\cal T}_{O,y},
\qquad
 {\mathbb B}_{ji}(x,y)
 =\sqrt{C_j^OC_i^E}\,
   \nabla_{y_j}\nabla_{x_i}\Phi_{ij}(x_i,y_j),                \tag{6.4}
\]
using the Riemannian metrics to identify covectors and vectors.
Put
\[
 q=\operatorname*{ess\,sup}_{(x,y)}
   \|{\mathbb B}(x,y)\|_{2\to2}.                             \tag{6.5}
\]

\begin{theorem}[Intrinsic parity maximal-correlation bound]
\label{thm:v15v6-intrinsic-correlation}
Assume the centered parity conditional-expectation operator is
compact and the common parity-measurable sigma algebra is trivial.
Then its maximal correlation \(c\) satisfies
\[
 \boxed{c\le q.}                                             \tag{6.6}
\]
If \(q<1\), the two-parity heat-bath generator has gap at least
\(1-q\).
\end{theorem}

\begin{proof}
Let
\[
 K:L^2_0(\mu_O)\to L^2_0(\mu_E),
 \qquad Kg(x)=\mathbb E[g(Y)\mid X=x].
\]
As in V4, compactness and trivial intersection give \(c=\|K\|<1\)
and a centered unit singular vector \(g\) such that
\[
 K^*Kg=c^2g,\qquad h=Kg,\qquad\|h\|_2=c.                    \tag{6.7}
\]
For \(v=(v_i)_{i\in E}\in{\cal T}_{E,x}\), differentiation of the
conditional Gibbs density gives
\[
 \dd h(x)[v]
 =-\operatorname{Cov}\left(
 g(Y),
 \sum_{i,j}\dd_{x_i}\Phi_{ij}(x_i,Y_j)[v_i]
 \,\middle|\,X=x\right).                                    \tag{6.8}
\]
Apply this with \(v_i=\sqrt{C_i^E}\,u_i\).  Conditional on \(X\),
the odd blocks are independent.  Tensorizing \((6.3)\) bounds the
conditional variance of the score in \((6.8)\) by
\[
 \|{\mathbb B}(x,Y)u\|_{{\cal T}_O}^2
 \le q^2\|u\|_{{\cal T}_E}^2                                \tag{6.9}
\]
after integration in the conditional odd law.  Conditional
Cauchy--Schwarz and duality therefore give
\[
 \sum_{i\in E}C_i^E|\nabla_ih(x)|^2
 \le q^2\operatorname{Var}(g(Y)\mid X=x).                    \tag{6.10}
\]

Conditional on \(Y\), the even blocks are independent, so
tensorization of \((6.2)\) yields
\[
 \mathbb E\operatorname{Var}(h(X)\mid Y)
 \le q^2\mathbb E\operatorname{Var}(g(Y)\mid X).             \tag{6.11}
\]
The singular identities give the left and right sides before the
factor \(q^2\) as
\[
 c^2(1-c^2),\qquad 1-c^2,
\]
respectively.  Since \(c<1\), division gives \(c\le q\).  The
two-projection formula \((2.4)\) proves the gap statement.
\end{proof}

\begin{remark}[Essential-supremum version]
\label{rem:v15v6-average}
Equation \((6.5)\) is a sufficient pointwise bound.  The proof only
uses the conditional quadratic estimate \((6.10)\).  A sharper
version may replace the essential supremum by any measurable
positive operator bound which remains valid after conditional
integration.  This is the natural location for field-dependent
weights or a Lyapunov correction.
\end{remark}

Let \(P_k\) update only block \(k\), let \(r_k>0\), and define
\[
 A_{\rm loc}=\sum_kr_k(I-P_k),\qquad r_*=\min_kr_k.            \tag{6.12}
\]

\begin{theorem}[Intrinsic local block gap]
\label{thm:v15v6-local-block-gap}
Under the hypotheses of
Theorem~\ref{thm:v15v6-intrinsic-correlation}, if \(q<1\), then
\[
 \boxed{\operatorname{gap}A_{\rm loc}\ge r_*(1-q).}          \tag{6.13}
\]
This statement is invariant under changes of local gauge
coordinates.
\end{theorem}

\begin{proof}
Given the odd parity, the even blocks are independent.  Conditional
Efron--Stein gives
\[
 \mathbb E\operatorname{Var}(f\mid Z_O)
 \le\sum_{i\in E}
 \mathbb E\operatorname{Var}
 (f\mid\hbox{all blocks except }i).                          \tag{6.14}
\]
The same inequality holds with the parities reversed.  Thus the
local form dominates \(r_*\) times the parity form.  Apply the
preceding theorem.  All quantities in the proof are intrinsic
gradients, Hessians, conditional expectations, and Riemannian
norms, so coordinate changes do not alter the conclusion.
\end{proof}

\begin{corollary}[Finite-color interaction hypergraphs]
\label{cor:v15v6-color}
Suppose the interaction hypergraph admits \(m<\infty\) colors such
that blocks of one color are conditionally independent given the
others.  If consecutive color classes can be grouped into two
super-parities satisfying the hypotheses above with
\(q\le q_*<1\), then the local generator has the lower bound
\[
 \operatorname{gap}A_{\rm loc}
 \ge r_*(1-q_*)                                              \tag{6.15}
\]
for that grouping.  If no two-grouping preserves conditional
independence, an \(m\)-projection angle theorem is additionally
required.
\end{corollary}

\begin{proof}
A valid grouping is exactly the bipartite theorem applied to the
product of all blocks in each super-parity.  Conditional
Efron--Stein then returns to the original local blocks.
\end{proof}

\begin{proposition}[When compactness and strictness are automatic]
\label{prop:v15v6-compactness}
At a fixed finite regulator, suppose every noncompact block has a
coercive positive quartic tail, every remaining block is compact,
and all interaction terms have lower degree in the noncompact
variables.  If the joint density is smooth and strictly positive,
then the parity conditional operator is Hilbert--Schmidt and its
maximal correlation is strictly below one.
\end{proposition}

\begin{proof}
On compact factors there is no tail integral.  On Euclidean
factors, the marginal lower-bound argument of
Lemma~\ref{lem:v15v2-hs} loses only polynomial factors.
Coercive quartic tails dominate those factors and every
lower-degree mixed term, proving square integrability of the
density ratio.  Hilbert--Schmidt compactness follows.  Equality of
maximal correlation would produce nonconstant functions of the two
parities which agree almost surely.  Strict positivity and Fubini
force them to be constant.
\end{proof}

\begin{assessment}[Literal source obstruction]
\label{ass:v15v6-source}
The intrinsic theorem avoids a false global gauge chart, but it
does not make the constants in \((6.2)\)--\((6.5)\) automatic.
In a gauge--Higgs interaction such as
\[
 |H_x-U_{xy}H_y|^2,                                         \tag{6.16}
\]
the oscillation and gauge Hessian can grow like
\(|H_x||H_y|\).  Since the Higgs field is unbounded, a compact
gauge block can acquire arbitrarily high barriers or curvature.
Neither compactness of the group nor positivity of Haar measure
gives a conditional Poincare constant uniform in those external
Higgs values.  The mixed-Hessian essential supremum in \((6.5)\)
can also be infinite.

Thus V6 is a complete block compiler, and
Proposition~\ref{prop:v15v6-compactness} gives a positive gap at
each fixed finite regulator, but the literal source still needs a
weighted estimate which couples the gauge constant to the quartic
Higgs Lyapunov function.
\end{assessment}

\begin{decision}[V7 field-weighted block estimate]
\label{dec:v15v6-next}
Replace the essential supremum in \((6.5)\) by a conditional
quadratic estimate weighted by
\[
 W(H)=1+\sum_x|H_x|^2
 \quad\hbox{or}\quad
 W(H)=\exp\!\left(\varepsilon\sum_x|H_x|^2\right).
\]
Derive the gauge-block Poincare and mixed-Hessian terms on the
bounded-\(W\) region, and seek a quartic heat-bath drift which
returns the complementary region.  The combination must be a
proved form inequality, not a high-probability assertion.  If
\((6.16)\) generates multiple compact-group wells with an
exponentially small conditional gap, record that barrier as the
next upstream obstruction.
\end{decision}

\begin{firewall}[Claim boundary after V6]
\label{fire:v15v6}
The scalar V4 theorem now has a coordinate-free block extension
covering Euclidean and compact manifolds under uniform conditional
Poincare and mixed-Hessian hypotheses.  Those hypotheses have not
been populated uniformly for the literal unbounded gauge--Higgs
shell, so no RPESM continuum gap is claimed.
\end{firewall}

\section*{Version V6 append-only record}
\addcontentsline{toc}{section}{Version V6 append-only record}
The complete V5 source was retained before this continuation.  It
had \(49476\) bytes, \(1366\) lines, and SHA--256
\[
 \texttt{C4F78EE2295A6257CC7CFE1D4D66C2356DACE74E80575E6BEF88B5FCF661947A}.
\]
V6 proved the intrinsic parity-correlation and local block-gap
compiler for Euclidean Higgs and compact gauge manifolds, then
isolated the unbounded gauge--Higgs barrier in its uniform
constants.


\section{Fifteenth-cycle V7: compact-link amplitudes, a single-well theorem, and a multiwell obstruction}
\label{sec:v15v7}

V6 showed that an unbounded Higgs amplitude can make the
coefficients of a compact gauge-link conditional arbitrarily
large.  A bounded-density comparison then loses uniformity.
Large amplitude is not itself fatal: the decisive issue is whether
the compact conditional develops separated wells.

We test this on \(U(1)\).  For \(a\ge0\), let
\[
 \mu_a(\dd\theta)
 =Z_a^{-1}e^{a\cos\theta}\,\dd\theta,
 \qquad \theta\in\mathbb R/(2\pi\mathbb Z),                  \tag{7.1}
\]
and let \(\lambda_1(a)\) be the Poincare gap of
\[
 {\cal E}_a(f,f)=\int|f'(\theta)|^2\,\mu_a(\dd\theta).        \tag{7.2}
\]

\begin{theorem}[Uniform single-harmonic compact-link gap]
\label{thm:v15v7-von-mises}
There is a numerical constant \(\lambda_*>0\) such that
\[
 \boxed{\inf_{a\ge0}\lambda_1(a)\ge\lambda_*.}                \tag{7.3}
\]
Moreover,
\[
 \lambda_1(a)\asymp a\qquad(a\to\infty).                     \tag{7.4}
\]
Thus a gauge--Higgs coefficient proportional to
\(|H_x||H_y|\) does not destroy the link Poincare gap when the
conditional link density is exactly the one-well first harmonic
\((7.1)\).
\end{theorem}

\begin{proof}
For every finite \(a\), the density in \((7.1)\) is smooth and
strictly positive on the compact circle, so the gap is positive.
The associated closed forms depend continuously on \(a\) on every
compact parameter interval; the min--max principle makes
\(a\mapsto\lambda_1(a)\) continuous there.

It remains to control large \(a\).  On \(|\theta|\le\pi/3\),
\[
 a(1-\cos\theta)\asymp a\theta^2,\qquad
 ( -a\cos\theta)''=a\cos\theta\ge a/2.                       \tag{7.5}
\]
Outside this arc, the normalized mass is bounded by
\(C\sqrt a\,e^{-a/2}\).  Apply the one-dimensional Hardy
criterion on each of the two monotone half-circles, with the
minimum at \(\theta=0\).  The two Hardy products are
\[
 \sup_{0<r<\pi}
 \mu_a([r,\pi])
 \int_0^r\frac{\dd s}{p_a(s)}
 \quad\hbox{and their reflected copies},                    \tag{7.6}
\]
where \(p_a=Z_a^{-1}e^{a\cos s}\).  Splitting at
\(r=\pi/3\), the Gaussian estimates from \((7.5)\) bound
\((7.6)\) above by \(C/a\): in the core this is the ordinary
Gaussian Hardy product, while on the tail the decreasing mass
cancels the increasing reciprocal-density integral.  The matching
lower test \(f(\theta)=\theta\) on the core gives a constant
multiple of \(1/a\) for the Poincare constant.  Hence
\[
 c_1a\le\lambda_1(a)\le c_2a
\]
for all sufficiently large \(a\).  Continuity and positivity on
the remaining compact interval give the uniform floor
\((7.3)\).
\end{proof}

A source action can contain more than the first character.  For
\(m\ge2\), define
\[
 \mu_{a,m}(\dd\theta)
 =Z_{a,m}^{-1}e^{a\cos(m\theta)}\,\dd\theta.                  \tag{7.7}
\]
It has \(m\) equivalent wells.

\begin{theorem}[Multiwell compact-link soft mode]
\label{thm:v15v7-multiwell}
For every integer \(m\ge2\), there is \(C_m<\infty\) such that
\[
 \boxed{\lambda_1(a,m)
 \le C_m a\,e^{-2a}}
 \qquad(a\ge1),                                              \tag{7.8}
\]
where \(\lambda_1(a,m)\) is the Poincare gap of \(\mu_{a,m}\).
In particular,
\[
 \inf_{a\ge0}\lambda_1(a,m)=0.                               \tag{7.9}
\]
\end{theorem}

\begin{proof}
Choose a smooth periodic function \(f_m\) which is constant with
two different values on neighborhoods of two adjacent maxima of
\(\cos(m\theta)\), has mean zero by the \(2\pi/m\) symmetry, and
changes only in fixed small neighborhoods of the intervening
minima, where \(\cos(m\theta)\le-1+c_m\varepsilon^2\).
Its variance under \(\mu_{a,m}\) stays bounded below, because each
well has asymptotic mass \(1/m\).

Laplace's estimate at the \(m\) maxima gives
\[
 Z_{a,m}\asymp a^{-1/2}e^a.                                 \tag{7.10}
\]
Choose transition neighborhoods of width \(a^{-1/2}\), and scale
the transition profile so \(|f_m'|=O(\sqrt a)\) there.
On those neighborhoods,
\[
 e^{a\cos(m\theta)}\le e^{-a+C_m}.
\]
Their normalized Dirichlet contribution is therefore at most
\[
 C_m\,
 (a)(a^{-1/2})\,
 \frac{e^{-a}}{a^{-1/2}e^a}
 =C_m a\,e^{-2a}.                                            \tag{7.11}
\]
Thus, after enlarging \(C_m\), the estimate in \((7.8)\) holds for every \(a\ge1\).  The Rayleigh principle
proves the theorem.
\end{proof}

\begin{remark}[Strength of the exponential bound]
\label{rem:v15v7-bound}
The precise polynomial prefactor in \((7.8)\) is immaterial for
the present obstruction.  The two-well test proves an
exponentially vanishing gap.  A sharper saddle calculation can
replace the displayed prefactor by the corresponding
Eyring--Kramers value without changing any subsequent claim.
\end{remark}

\begin{corollary}[Gauge--Higgs conditional dichotomy]
\label{cor:v15v7-dichotomy}
Suppose a \(U(1)\) link conditional at fixed neighboring fields
has density proportional to
\[
 \exp\{a(H)\cos(\theta-\theta_0(H))\}.                        \tag{7.12}
\]
Then its intrinsic Poincare gap has a positive lower bound
independent of the unbounded amplitude \(a(H)\).
If instead the conditional action admits a term
\(a(H)\cos(m\theta)\), \(m\ge2\), with
\(a(H)\to\infty\) along physical Higgs configurations and the
lower harmonics do not select one well by a comparable amount,
then no uniform conditional link gap follows; the symmetric model
has the soft sequence \((7.8)\).
\end{corollary}

\begin{assessment}[Source row that must be audited]
\label{ass:v15v7-source}
The source RPESM definition specifies a finite-range common action
with compact gauge and route variables and nonzero mixed
couplings.  It does not present the complete one-link conditional
character expansion or prove that every large-amplitude
conditional is single-well modulo its physical stabilizer.
Therefore:
\begin{enumerate}
\item compactness of the gauge group is insufficient;
\item the crude oscillation comparison is unnecessarily weak on a
      first-harmonic branch;
\item higher characters or competing plaquette, gauge-fixing,
      counterterm, and modulus factors can create a genuine
      exponentially small barrier;
\item the physical quotient must remove exact gauge copies before
      a multiwell configuration is called a tunneling mode.
\end{enumerate}
The first missing datum is the literal conditional character and
stabilizer table, not another abstract Poincare lemma.
\end{assessment}

\begin{decision}[V8 literal compact-link audit]
\label{dec:v15v7-next}
Read the explicit bosonic, gauge-fixing, counterterm, and positive
control-modulus definitions used by the RPESM construction.
For each compact link block, list the characters which enter its
conditional density and quotient exact stabilizer copies.
Classify the remaining large-amplitude maxima as:
\begin{enumerate}
\item one physical well, where Theorem
      \ref{thm:v15v7-von-mises} suggests a uniform block gap;
\item a connected stabilizer manifold, requiring a tangential
      Poincare estimate;
\item separated physical wells, producing a tunneling debit.
\end{enumerate}
If the source leaves the action arbitrary within a class broad
enough to contain both \((7.1)\) and \((7.7)\), record
underdetermination and state the additional single-well hypothesis
needed by the V6 compiler.
\end{decision}

\begin{firewall}[Claim boundary after V7]
\label{fire:v15v7}
V7 proves that unbounded link amplitude is compatible with a
uniform compact-block gap on the one-well \(U(1)\) branch, and
that a symmetric multiwell harmonic has an exponentially vanishing
gap.  The literal source character content and physical quotient
have not yet selected between these cases.
\end{firewall}

\section*{Version V7 append-only record}
\addcontentsline{toc}{section}{Version V7 append-only record}
The complete V6 source was retained before this continuation.  It
had \(59059\) bytes, \(1624\) lines, and SHA--256
\[
 \texttt{FA3FE62B6CDF9FE63249AB355F0F4AAF400642FA35A1056FDE93F715DE006801}.
\]
V7 established the uniform one-well compact-link gap and the
multiwell exponential soft-mode obstruction, reducing the source
question to its literal conditional character and stabilizer data.




\section{Fifteenth-cycle V8: literal compact-link source audit and higher-character birth under shell marginalization}
\label{sec:v15v8}

V7 reduced the compact-link gap question to the complete
conditional character table.  The Grand Tensor contains three
different action levels which must not be conflated.

\begin{record}[Three source action levels]
\label{rec:v15v8-action-levels}
The explicit classical SM-II action \({\rm(CA.12)}\) contains a
principal-chart curvature square, a covariant Higgs edge,
the symmetry-breaking quartic, and the Dirac--Yukawa bilinear.
It is an exact local classical action, but its plaquette curvature
uses the principal logarithm chart.

The protected literal finite shell \({\rm(FS.8)}\) instead uses
the defining-representation Wilson terms
\[
 \sum_{a=3,2}\beta_a\sum_p
 \left(1-\frac1{d_a}\operatorname{Re}\operatorname{tr}U_p^{(a)}\right)
 +\beta_1\sum_p(1-\cos\alpha_p),                             \tag{8.1}
\]
together with
\[
 \kappa_H\sum_{e:x\to y}
 \|H_y-\rho_H(U_e)H_x\|^2
 +\sum_x(m_H^2\|H_x\|^2+\lambda_H\|H_x\|^4).                \tag{8.2}
\]
At this level \(H\) is restricted to a fixed compact ball, the
chiral operator is restricted to a protected singular-value chart,
and the positive field cylinder also contains
\(\|\sigma_n(U,H)\|^2\).  Its field-gap theorem assumes, rather
than derives, the uniform Dobrushin margin
\[
 \sup_n\sup_i\sum_{j\ne i}c_{ij}^{(n)}<1.                    \tag{8.3}
\]

The later projective RPESM cutoff \({\rm(RP.5)}\) restores an
unbounded quartically controlled Higgs field and declares
\[
 S_{B,n}=S_{g,n}+S_{{\rm YM},n}+S_{H,n}
          +S_{{\rm gf},n}+S_{{\rm med},n}+S_{{\rm ct},n}.     \tag{8.4}
\]
It requires finite range, reflection evenness, quartic domination,
and nonzero mixed couplings, but gives no complete one-link
Peter--Weyl table for \((8.4)\), the positive control-section
modulus, or the exact shell tails.
\end{record}

\begin{proposition}[Bare defining-representation dependence]
\label{prop:v15v8-bare-head}
Freeze every field except one link \(U_e\) in the bare action
\((8.1)\)--\((8.2)\).  Its plaquette and Higgs-edge dependence is
a finite sum of defining-representation matrix coefficients:
\[
 \operatorname{Re}\operatorname{tr}(A_pU_e),\qquad
 \operatorname{Re}\langle h_y,\rho_H(U_e)h_x\rangle,         \tag{8.5}
\]
and their inversely oriented conjugates.  In the \(U(1)\) factor,
a sum of the plaquette terms alone reduces to one first harmonic
\[
 A\cos\theta+B\sin\theta
 =R\cos(\theta-\theta_0).                                    \tag{8.6}
\]
\end{proposition}

\begin{proof}
In a plaquette containing \(e\), cyclically multiply all frozen
links into a matrix \(A_p\).  The holonomy trace is then
\(\operatorname{tr}(A_pU_e)\) or its conjugate.  Expanding the
squared covariant Higgs difference leaves a constant plus twice
the real matrix coefficient in \((8.5)\).  In the abelian factor,
addition of first sines and cosines gives \((8.6)\).
\end{proof}

Proposition~\ref{prop:v15v8-bare-head} does not classify the full
conditional cylinder.  The determinant norm, gauge-fixing,
mediator, counterterm, and exact shell-tail factors multiply the
bare weight and can carry higher representations.  The following
finite example shows that higher characters arise even when the
fine interaction is first harmonic.

\begin{theorem}[Exact higher-character birth from one hidden sign]
\label{thm:v15v8-hidden-sign}
Let \(\sigma\in\{-1,1\}\) have counting reference law and let
\(\theta\in\mathbb R/(2\pi\mathbb Z)\).  Give the fine pair the
positive reflection-even weight
\[
 w_a(\theta,\sigma)=e^{a\sigma\cos\theta},
 \qquad a>0.                                                 \tag{8.7}
\]
Eliminating \(\sigma\) gives
\[
 \boxed{
 w_a^{\rm eff}(\theta)
 =\frac12\sum_{\sigma=\pm1}w_a(\theta,\sigma)
 =\cosh(a\cos\theta).}                                       \tag{8.8}
\]
This effective weight is \(\pi\)-periodic and its Fourier
expansion contains only even characters, with a strictly positive
\(\cos(2\theta)\) coefficient.  Its two maxima at
\(\theta=0,\pi\) are separated by the minima
\(\theta=\pi/2,3\pi/2\).

Let \(\lambda_{\rm eff}(a)\) be the Poincare gap of the normalized
law proportional to \((8.8)\).  There is \(C<\infty\) such that
\[
 \boxed{\lambda_{\rm eff}(a)\le Ca^{3/2}e^{-a}}
 \qquad(a\ge1).                                              \tag{8.9}
\]
Thus exact marginalization of a first-harmonic fine interaction
can generate a physical two-well compact-link soft mode.
\end{theorem}

\begin{proof}
Equation \((8.8)\) is direct summation.  Expanding
\[
 \cosh(a\cos\theta)
 =1+\frac{a^2}{2}\cos^2\theta+\cdots
\]
and using
\(\cos^2\theta=(1+\cos2\theta)/2\) shows that the second
character is present with positive coefficient; the identity
\(w(\theta+\pi)=w(\theta)\) excludes odd characters.

The normalizing integral has the Laplace behavior
\[
 Z_a^{\rm eff}\asymp a^{-1/2}e^a,                            \tag{8.10}
\]
because there are two nondegenerate maxima.  Choose a mean-zero
smooth function which is \(+1\) near \(0\), \(-1\) near \(\pi\),
and changes only in intervals of width \(a^{-1}\) about
\(\pi/2\) and \(3\pi/2\).  Its variance stays bounded below.
On the transition intervals,
\(\cosh(a\cos\theta)\le C\), while the squared derivative is
\(O(a^2)\).  Hence its normalized Dirichlet energy is at most
\[
 C\,
 (a^2)(a^{-1})
 \frac1{a^{-1/2}e^a}
 =Ca^{3/2}e^{-a}.                                            \tag{8.11}
\]
The Rayleigh principle proves \((8.9)\).
\end{proof}

\begin{proposition}[The two abelian wells are not gauge copies]
\label{prop:v15v8-physical-wells}
If \(\theta\) is a \(U(1)\) plaquette holonomy, the configurations
\(\theta=0\) and \(\theta=\pi\) are invariant under vertex gauge
transformations and have different fundamental Wilson values.
They are therefore distinct gauge-orbit points.
\end{proposition}

\begin{proof}
A vertex gauge transformation cancels around a closed abelian
plaquette, so its holonomy is unchanged.  The fundamental Wilson
characters are \(\cos0=1\) and \(\cos\pi=-1\).
\end{proof}

\begin{theorem}[V8 literal-source verdict]
\label{thm:v15v8-verdict}
The source proves neither the uniform one-well branch of V7 nor a
multiwell obstruction for the complete projective RPESM link
conditional.  More precisely:
\begin{enumerate}
\item the bare protected action \((8.1)\)--\((8.2)\) has
      defining-representation link dependence;
\item its published uniform field gap is conditional on
      \((8.3)\) and uses a compact Higgs screen;
\item the unbounded RPESM action \((8.4)\) does not specify the
      full conditional character table or populate \((8.3)\);
\item exact shell elimination can generate even higher characters
      and the soft mode \((8.9)\).
\end{enumerate}
Therefore no gauge-link Poincare constant may be imported from the
protected shell into the unbounded projective shell without a new
single-well, stabilizer, or tunnelling estimate.
\end{theorem}

\begin{proof}
Items one through three are the literal content of
\({\rm(FS.8)}\), \({\rm(FS.20)}\)--\({\rm(FS.24)}\), and
\({\rm(RP.5)}\).  Item four is
Theorem~\ref{thm:v15v8-hidden-sign}.  The compact-screen and
unbounded-shell hypotheses are different, and the assumed
Dobrushin row cannot be transferred across that difference.
\end{proof}

\begin{assessment}[Upstream consequence]
\label{ass:v15v8-consequence}
The direct bare action is more specific than V7's first source
reading suggested, but the complete quantum/projective cylinder is
less specific than the bare head.  The correct input to the V6
block theorem is not the character content of the classical action
alone.  It is the character and stabilizer content of the
conditional logarithm of the full positive cylinder after
determinant modulus, gauge fixing, mediator variables, and the
exact shell tail have been assembled.
\end{assessment}

\begin{decision}[V9 robust single-well and retained-detail alternatives]
\label{dec:v15v8-next}
Develop two conditional routes:
\begin{enumerate}
\item prove a bounded-remainder comparison for a one-well link
      density
      \(e^{a\cos(\theta-\theta_0)-R(\theta)}\), uniform in
      \(a\), when \(\operatorname{osc}R\) is uniformly bounded;
\item for a higher-character shell term, retain the detail or
      sector label instead of marginalizing it and compute the
      resulting two-block maximal correlation.
\end{enumerate}
The first route supplies a usable sufficient hypothesis for the
V6 compiler.  The second determines whether retaining ancestry
removes the tunnelling mode or merely moves it into a slow sector
transition.
\end{decision}

\begin{firewall}[Claim boundary after V8]
\label{fire:v15v8}
The protected bare Wilson--Higgs head and the full unbounded
projective cylinder have been separated.  Higher-character birth
and its two-well soft mode are proved by an exact finite shell
example, but the literal RPESM coefficients have not been shown to
realize that example.
\end{firewall}

\section*{Version V8 append-only record}
\addcontentsline{toc}{section}{Version V8 append-only record}
The complete V7 source was retained before this continuation.  It
had \(67268\) bytes, \(1841\) lines, and SHA--256
\[
 \texttt{88E0830FFD84B74A0B7FEB27BF118DA9D2EAF2DDC834700CCA32C2A2B08DA946}.
\]
V8 audited the three distinct source action levels, proved
higher-character birth under exact sign elimination, and showed
why the protected compact-screen gap does not populate the
unbounded projective RPESM link gap.





\section{Fifteenth-cycle V9: robust one-well comparison and the retained-sign slow mode}
\label{sec:v15v9}

V8 separated the bare first-character head from higher characters
generated by the complete cylinder.  We now prove a usable
sufficient condition for the one-well branch and test whether
retaining the hidden ancestry variable removes the two-well
barrier.

\begin{theorem}[Bounded-remainder stability of the one-well link gap]
\label{thm:v15v9-bounded-remainder}
Let
\[
 \mu_{a,R}(\dd\theta)
 =Z_{a,R}^{-1}
 \exp\{a\cos(\theta-\theta_0)-R(\theta)\}\,\dd\theta,         \tag{9.1}
\]
where \(a\ge0\), \(\theta_0\in\mathbb R/(2\pi\mathbb Z)\), and
\(R\) is measurable with
\[
 \operatorname{osc}R
 :=\operatorname*{ess\,sup}R-
   \operatorname*{ess\,inf}R\le M.                           \tag{9.2}
\]
If \(\lambda_*\) is the uniform von Mises floor from
Theorem~\ref{thm:v15v7-von-mises}, then
\[
 \boxed{\operatorname{gap}(\mu_{a,R})
 \ge e^{-M}\lambda_*}                                       \tag{9.3}
\]
for every \(a,\theta_0,R\) satisfying \((9.2)\).
\end{theorem}

\begin{proof}
Rotation removes \(\theta_0\).  Relative to the von Mises law
\(\mu_a\), the density ratio of \(\mu_{a,R}\) is a normalized
multiple of \(e^{-R}\).  The ratio of its essential supremum to
its essential infimum is at most \(e^M\).  The
Holley--Stroock Poincare comparison gives
\[
 C_P(\mu_{a,R})\le e^M C_P(\mu_a).
\]
Taking reciprocals and using
\(\operatorname{gap}(\mu_a)\ge\lambda_*\) proves \((9.3)\).
\end{proof}

\begin{corollary}[A sufficient conditional-character row]
\label{cor:v15v9-character-row}
A family of \(U(1)\) link conditionals has a uniform Poincare gap
if, after quotienting exact stabilizer copies, its logarithmic
density can be written
\[
 a(\zeta)\cos(\theta-\theta_0(\zeta))-R_\zeta(\theta),
 \qquad
 \sup_\zeta\operatorname{osc}R_\zeta\le M<\infty.             \tag{9.4}
\]
The leading amplitude \(a(\zeta)\) may be unbounded.
\end{corollary}

The hypothesis controls the total residual potential, not each
Fourier coefficient separately.  It therefore permits infinitely
many higher characters when their assembled oscillation stays
bounded.

We next retain the sign eliminated in V8.  Let
\[
 \widehat\mu_a(\dd\theta,\dd\sigma)
 =\widehat Z_a^{-1}e^{a\sigma\cos\theta}\,
   \dd\theta\,(\delta_{-1}+\delta_1)(\dd\sigma).              \tag{9.5}
\]
Let \(P_\theta\) condition on \(\sigma\), so \(I-P_\theta\)
updates the angle, and let \(P_\sigma\) condition on \(\theta\),
so \(I-P_\sigma\) updates the sign.  Define
\[
 A_a=(I-P_\theta)+(I-P_\sigma).                              \tag{9.6}
\]

\begin{theorem}[Retaining the shell sign does not remove tunnelling]
\label{thm:v15v9-retained-sign}
There is \(C<\infty\) such that
\[
 \boxed{\operatorname{gap}A_a
 \le C a^{-1/2}e^{-a}}
 \qquad(a\ge1).                                              \tag{9.7}
\]
Thus the two wells of the marginalized angle law become the two
joint wells
\[
 (\sigma,\theta)=(1,0),\qquad(-1,\pi),                       \tag{9.8}
\]
and the slow mode survives in the sign heat bath.
\end{theorem}

\begin{proof}
Use the centered test function \(f(\theta,\sigma)=\sigma\).
Symmetry gives \(\operatorname{Var}_{\widehat\mu_a}f=1\).
The angle update does not change \(f\), while
\[
 \mathbb E[\sigma\mid\theta]=\tanh(a\cos\theta),
 \qquad
 \operatorname{Var}(\sigma\mid\theta)
 =\operatorname{sech}^2(a\cos\theta).                        \tag{9.9}
\]
Consequently
\[
 \langle f,A_af\rangle
 =\mathbb E_{\widehat\mu_a}
   \operatorname{sech}^2(a\cos\theta)
 =\frac{\displaystyle
   \int_0^{2\pi}\frac{\dd\theta}{\cosh(a\cos\theta)}}
  {\displaystyle
   \int_0^{2\pi}\cosh(a\cos\theta)\,\dd\theta}.               \tag{9.10}
\]
The denominator is comparable to \(a^{-1/2}e^a\).  Split the
numerator into four intervals centered at the two zeros of
\(\cos\theta\) and their complements.  On a fixed neighborhood
of each zero, the change of variable
\(u=a\cos\theta\) gives an \(O(a^{-1})\) integral because
\(\int_{\mathbb R}\operatorname{sech}u\,\dd u<\infty\).
The complement is exponentially smaller.  Hence the ratio in
\((9.10)\) is at most \(Ca^{-1/2}e^{-a}\).  The Rayleigh
principle proves \((9.7)\).
\end{proof}

\begin{corollary}[Maximal correlation tends to one]
\label{cor:v15v9-correlation}
Let \(c_a\) be the maximal correlation between \(\sigma\) and
\(\theta\) under \(\widehat\mu_a\).  Then
\[
 c_a^2\ge
 \operatorname{Var}\bigl(\mathbb E[\sigma\mid\theta]\bigr)
 =1-\mathbb E\operatorname{sech}^2(a\cos\theta),             \tag{9.11}
\]
so
\[
 1-c_a\le Ca^{-1/2}e^{-a}.                                  \tag{9.12}
\]
\end{corollary}

\begin{proof}
The variable \(\sigma\) is centered with unit variance.  Its
conditional expectation is an admissible image in the definition
of maximal correlation.  Apply \((9.10)\).
\end{proof}

There is a formal way to erase the barrier: update
\((\theta,\sigma)\) jointly from its full conditional law.  The
corresponding conditional-expectation block again has exact fiber
gap one.  Whether this is a physical local move depends on the
support of the retained sign.  If the sign is one cell variable,
the pair is a local enlarged block.  If it records an extensive
shell, topological sector, determinant orientation, or long
Wilson loop, the joint update is nonlocal and cannot be inserted
as a physical gap mechanism.

\begin{theorem}[Exact block-support criterion]
\label{thm:v15v9-support}
Let \(\sigma\) be a retained detail variable and let
\({\cal S}(\sigma)\) be the smallest physical cell set on which it
is measurable.  A joint heat-bath update of \((U_e,\sigma)\) is a
bounded-range block update only if
\[
 \sup_n\operatorname{diam}
 \bigl(\{e\}\cup{\cal S}_n(\sigma)\bigr)<\infty               \tag{9.13}
\]
in lattice units.  If this diameter diverges, its unit fiber gap
is a nonlocal shell refresh and must be excluded from the physical
local gap.
\end{theorem}

\begin{proof}
A block update may directly change every variable in its
measurable support.  Bounded-range locality requires a common
diameter bound on that support.  Conversely, under \((9.13)\) the
pair lies in a uniformly bounded block; its conditional
expectation is a local projection with unit fiber gap.
\end{proof}

\begin{assessment}[The two viable compact-link routes]
\label{ass:v15v9-routes}
The V6 block compiler can now accept either:
\begin{enumerate}
\item a quotient one-well decomposition with the uniform remainder
      bound \((9.4)\); or
\item a uniformly bounded enlarged local block which retains and
      jointly updates every detail responsible for the competing
      wells.
\end{enumerate}
Retaining ancestry without a joint local update is insufficient by
Theorem~\ref{thm:v15v9-retained-sign}.  The literal RPESM source
does not yet specify an \(M\) for \((9.4)\) or the support
diameters in \((9.13)\).
\end{assessment}

\begin{decision}[V10 compulsory audit and V11 cluster route]
\label{dec:v15v9-next}
At V10 inspect the newest companion notes for a bounded character
remainder, a local retained-sector support, a cluster heat-bath
cover, or a physical locality radius.  Then continue in V11 by
proving a fractional-coverage gap for enlarged gauge--detail
blocks satisfying \((9.13)\), with an explicit overlap and
cross-block influence debit.
\end{decision}

\begin{firewall}[Claim boundary after V9]
\label{fire:v15v9}
Bounded character remainders preserve the one-well compact-link
gap uniformly, while a retained hidden sign has an exponentially
small two-block heat-bath gap.  No literal RPESM remainder bound
or local ancestry-support certificate has yet been supplied.
\end{firewall}

\section*{Version V9 append-only record}
\addcontentsline{toc}{section}{Version V9 append-only record}
The complete V8 source was retained before this continuation.  It
had \(76819\) bytes, \(2084\) lines, and SHA--256
\[
 \texttt{55C50A210D2000D272FDF38CFF3CC44A0818211AEE783E15A5EEA7BE6C66A512}.
\]
V9 proved a uniform bounded-remainder one-well theorem, showed
that retained sign ancestry still has an exponentially slow mode,
and isolated the bounded-support condition for a legitimate joint
local update.


\section{Fifteenth-cycle V10: compulsory companion audit before the cluster compiler}
\label{sec:v15v10}

V9 identified two legitimate compact-link routes: a uniformly
bounded one-well remainder, or a joint update containing every
local detail responsible for competing wells.  This checkpoint
tests the newest companions for a cluster cover or locality
constant.

\begin{record}[V10 companion checkpoint]
\label{rec:v15v10-companions}
The newest Yang--Mills note is
\begin{center}
\path{Renewal_Yang_Mills_Mass_Gap_Research_Notes_V33.tex},
\end{center}
with \(270628\) bytes, \(7390\) lines, and SHA--256
\[
 \texttt{76F776C9F0CFCD187A77631E0CD9D10634910657D8FFCC2501C131ED1B640793}.
                                                                  \tag{10.1}
\]
The Navier--Stokes note remains
\begin{center}
\path{Renewal_NS_Stability_Research_Notes_V26.tex},
\end{center}
with \(278283\) bytes, \(5319\) lines, and SHA--256
\[
 \texttt{82C887F8C2C69E4F28FAD7C3DC8DE5B9099CB5A4BF431EBA1FFF3E91935978AD}.
                                                                  \tag{10.2}
\]
\end{record}

Yang--Mills V33 constructs three independent radius-\(1/8\)
full-active-circle charts obtained by moving separately in each
transverse coordinate.  Each chart has its own exact splitting and
matrix ledger, and each overlaps the V28 anchor tube.  This creates
a genuinely three-dimensional branch seed; the note explicitly
does not infer coverage of the phase cube from those four charts.

\begin{assessment}[Transferable atlas lesson]
\label{ass:v15v10-ym}
The applicable method is stronger than the earlier diagonal
corridor:
\begin{enumerate}
\item certify every independent coordinate direction;
\item build separate local data rather than identify unequal
      coordinate roles;
\item record an overlap graph;
\item test the higher-dimensional cells between the vertices.
\end{enumerate}
A cluster cover for the RPESM state space must follow the same
logic.  Covering each isolated interaction motif and checking
pairwise overlaps does not automatically cover configurations
where several motifs are simultaneously active.

YM V33 nevertheless contains no conditional probability kernel,
retained-detail support radius, fractional block coverage,
Poincare constant, or physical update clock.  Its momentum-space
radii and Wilson-pencil floors are not imported.
\end{assessment}

Navier--Stokes V26 contains no new common object with the
gauge--detail heat-bath problem.

\begin{theorem}[V10 audit verdict]
\label{thm:v15v10-verdict}
At checkpoints \((10.1)\)--\((10.2)\), neither companion supplies:
\begin{enumerate}
\item a bounded character remainder of the form \((9.4)\);
\item a retained-sector support satisfying \((9.13)\);
\item a constraint-preserving cluster heat-bath cover;
\item a post-constraint block influence matrix;
\item a lattice-to-physical clock comparison.
\end{enumerate}
The only imported information is the independently certified
multi-directional atlas rule.
\end{theorem}

\begin{proof}
The new YM statements are exact matrix inequalities on overlapping
momentum tubes and explicitly leave their full cell uncovered.
NS V26 remains a deterministic Oseen-kernel estimate.  Neither
defines any object in the five requested rows.
\end{proof}

\begin{decision}[V11 enlarged local clusters]
\label{dec:v15v10-next}
Use the already proved V82 overlapping-block theorem rather than
repeating it.  Specialize it to uniformly bounded gauge--detail
clusters.  Prove that the internal multiwell energy cancels from
the likelihood-ratio oscillation between two boundary conditions,
so only boundary interactions enter the inter-cluster debit.
Then state the fractional coverage, sector-kernel, and locality
conditions under which the cluster generator has a refresh-free
gap.
\end{decision}

\begin{firewall}[Claim boundary after V10]
\label{fire:v15v10}
The mandatory audit imports no numerical constant.  A physical
cluster gap remains conditional on a new RPESM motif cover and
boundary-influence estimate.
\end{firewall}

\section*{Version V10 append-only record}
\addcontentsline{toc}{section}{Version V10 append-only record}
The complete V9 source was retained before this continuation.  It
had \(84983\) bytes, \(2309\) lines, and SHA--256
\[
 \texttt{347E8BA797EA7F8F66866C7185CE770174F1B7906DD8974912B5EC4822784FCD}.
\]
V10 performed the compulsory companion audit, retained the
multi-directional atlas rule, and found no cluster-gap or locality
constant applicable to the RPESM shell.


\section{Fifteenth-cycle V11: local motif clusters and boundary-only influence}
\label{sec:v15v11}

V9 showed that jointly updating a link with its retained detail can
remove an internal tunnelling barrier only when the joint support
has uniformly bounded diameter.  We now combine that observation
with the overlapping-block compiler proved in archived V82.

Fix one physical constraint sector \(\Omega_s\) with positive law
\[
 \mu_s(\dd z)=Z_s^{-1}e^{-{\cal H}(z)}\lambda_s(\dd z).
\]
Let \(B\) be a finite cluster of field coordinates.  Split the
action relative to \(B\) as
\[
 {\cal H}(u,\xi)
 =U_B(u)+W_{\partial B}(u,\xi)+H_{B^c}(\xi),                 \tag{11.1}
\]
where \(u=z_B\), \(\xi=z_{B^c}\), \(U_B\) contains every
interaction supported wholly inside \(B\), and
\(W_{\partial B}\) contains precisely the interactions meeting
both \(B\) and \(B^c\).  The conditional block law is
\[
 K_B^\xi(\dd u)
 =Z_B(\xi)^{-1}
 e^{-U_B(u)-W_{\partial B}(u,\xi)}\lambda_B^\xi(\dd u).
                                                                  \tag{11.2}
\]

\begin{theorem}[Internal-barrier cancellation]
\label{thm:v15v11-barrier-cancel}
Assume the conditional reference \(\lambda_B^\xi\) is unchanged
when the outside block \(C\) is varied, or include its
log-density variation in \(W_{\partial B}\).  For outside
configurations \(\xi,\xi'\),
\[
 \begin{split}
 \operatorname{osc}_u
 \log\frac{\dd K_B^\xi}{\dd K_B^{\xi'}}(u)
 =
 \operatorname{osc}_u\bigl[
 W_{\partial B}(u,\xi)-W_{\partial B}(u,\xi')
 \bigr].                                                     \tag{11.3}
 \end{split}
\]
In particular, every coefficient and barrier contained entirely
in \(U_B\) cancels exactly.

If \(\xi,\xi'\) differ only in cluster \(C\), put
\[
 \Delta_{BC}^{\partial}
 =\sup_{\xi\sim_C\xi'}
 \operatorname{osc}_u\bigl[
 W_{\partial B}(u,\xi)-W_{\partial B}(u,\xi')
 \bigr].                                                     \tag{11.4}
\]
Then the block influence satisfies
\[
 \boxed{
 \sup_{\xi\sim_C\xi'}
 \|K_B^\xi-K_B^{\xi'}\|_{\rm TV}
 \le\tanh(\Delta_{BC}^{\partial}/4).}                        \tag{11.5}
\]
\end{theorem}

\begin{proof}
The logarithmic likelihood ratio is
\[
 \begin{split}
 \log\frac{\dd K_B^\xi}{\dd K_B^{\xi'}}(u)
 ={}&-W_{\partial B}(u,\xi)
     +W_{\partial B}(u,\xi')\\
 &-\log Z_B(\xi)+\log Z_B(\xi').                            \tag{11.6}
 \end{split}
\]
The internal action \(U_B\) cancels pointwise, and the last two
terms are constant in \(u\), so taking oscillation proves
\((11.3)\).  The likelihood-ratio oscillation lemma proved in V81
then gives \((11.5)\).
\end{proof}

\begin{corollary}[The hidden-sign motif becomes a unit-gap block]
\label{cor:v15v11-sign-cluster}
For the V9 motif
\[
 U_B(\theta,\sigma)=-a\sigma\cos\theta,\qquad
 B=\{\theta,\sigma\},                                       \tag{11.7}
\]
the joint conditional-expectation update \(I-{\cal E}_B\) has
exact fiber gap one for every \(a\ge0\).  Its influence on another
cluster contains no factor \(e^a\) and no internal tunnelling
debit; it is bounded solely by the boundary interaction in
\((11.4)\).
\end{corollary}

\begin{proof}
Conditional expectation is an orthogonal projection, so its
nonzero local spectrum is \(\{1\}\), as proved in V99.  The
coefficient \(a\) occurs only in \(U_B\), and
Theorem~\ref{thm:v15v11-barrier-cancel} removes it from the
boundary likelihood-ratio oscillation.
\end{proof}

Let \({\mathfrak B}_n\) be a family of clusters at cutoff \(n\),
with update rates \(r_B\), and define
\[
 A_n^{\rm cl}
 =\sum_{B\in{\mathfrak B}_n}r_B(I-{\cal E}_B).               \tag{11.8}
\]
Use the rate-weighted fractional coverage number
\(\kappa_{*,n}\) and the post-constraint overlapping-block
influence number \(\widehat\alpha_n\) of the archived V82
compiler.  Thus \(\kappa_{*,n}\) is the certified lower coverage
of every active coordinate by the admissible block rates, and
\(\widehat\alpha_n\) is computed after forbidden moves and fixed
sector labels have been removed.  V82 proves the approximate
tensorization estimate
\[
 \operatorname{Var}_{\mu_{s,n}}f
 \le
 \frac{1}{\kappa_{*,n}(1-\widehat\alpha_n)}
 \sum_{B\in{\mathfrak B}_n}r_B
 \mathbb E\operatorname{Var}
 (f\mid{\cal F}_{B^c})                                      \tag{11.9}
\]
whenever \(\kappa_{*,n}>0\) and
\(\widehat\alpha_n<1\).

\begin{theorem}[Uniformly local cluster-gap compiler]
\label{thm:v15v11-cluster-gap}
Suppose, uniformly in \(n\):
\begin{enumerate}
\item every \(B\in{\mathfrak B}_n\) is constraint preserving and
      has lattice diameter at most \(R<\infty\);
\item every retained detail responsible for a competing compact
      well is contained, together with its coupled physical link,
      in at least one cluster \(B\);
\item the common block-invariant sigma algebra is trivial within
      the selected sector;
\item
      \[
       \kappa_{*,n}\ge\kappa_*>0,\qquad
       \widehat\alpha_n\le\alpha_*<1.                        \tag{11.10}
      \]
\end{enumerate}
Then \(A_n^{\rm cl}\) is an \(R\)-local, refresh-free generator
with
\[
 \boxed{
 \operatorname{gap}_{\Omega_{s,n}}A_n^{\rm cl}
 \ge\kappa_*(1-\alpha_*).}                                  \tag{11.11}
\]
Internal multiwell amplitudes captured by the clusters do not
enter this bound.  Only their boundary interactions can enter
\(\widehat\alpha_n\).
\end{theorem}

\begin{proof}
The block form of \(A_n^{\rm cl}\) is the right-hand side of
\((11.9)\).  Multiplying \((11.9)\) by
\(\kappa_{*,n}(1-\widehat\alpha_n)\) proves the gap.  The common
invariant-algebra hypothesis makes the sector vacuum simple.
The diameter bound makes every direct update \(R\)-local.
Theorem~\ref{thm:v15v11-barrier-cancel} shows that a completely
captured motif contributes only through interactions crossing the
cluster boundary.
\end{proof}

\begin{proposition}[Uncaptured labels remain exact zero modes]
\label{prop:v15v11-uncaptured}
Let \(q_n\) be a nonconstant physical or ancestry label which is
unchanged by every cluster update.  Then every centered function
of \(q_n\) lies in \(\ker A_n^{\rm cl}\), so the unrestricted
gap is zero.  It can be removed only by:
\begin{enumerate}
\item conditioning on a selected physical sector;
\item adding a bounded-diameter constraint-preserving block which
      changes the label; or
\item proving that the label is OS-null in the physical quotient.
\end{enumerate}
\end{proposition}

\begin{proof}
For \(f=f(q_n)\), every conditional projection satisfies
\({\cal E}_Bf=f\).  Hence every summand in \((11.8)\) vanishes.
The three listed operations are exactly the ways the function can
cease to be a nonzero physical invariant.
\end{proof}

\begin{assessment}[What V11 removes]
\label{ass:v15v11-removes}
V11 removes the exponential internal barrier from the local gap
when its complete ancestry motif fits inside a physical cluster.
This is stronger than applying a bounded potential comparison to
the original single-link update.  It also makes the remaining
tasks finite and geometric:
\begin{enumerate}
\item enumerate every interaction or retained-detail motif;
\item assign it to a bounded cluster;
\item verify coverage and sector irreducibility;
\item estimate only the interactions crossing cluster boundaries.
\end{enumerate}

The theorem is still conditional for the RPESM source.  Its action
is declared finite range, but the determinant modulus, exact shell
counterterm, and quasilocal Wilsonian tail may have supports larger
than one bare plaquette.  Their support and decay must enter the
motif ledger before a uniform \(R\), \(\kappa_*\), or
\(\alpha_*\) is claimed.
\end{assessment}

\begin{decision}[V12 literal motif and tail ledger]
\label{dec:v15v11-next}
Extract the interaction supports of the explicit bare
plaquette--Higgs action, the determinant/Pfaffian control modulus,
the gauge-fixing quartet, mediator variables, and the exact shell
counterterm.  Separate:
\begin{enumerate}
\item strictly bounded motifs eligible for V11 clusters;
\item quasilocal tails with a summable boundary-influence norm;
\item global line, sector, or orientation labels which cannot be
      updated locally.
\end{enumerate}
Derive a finite coloring and coverage number for the bounded
motifs.  If the source gives no decay constant for the tail,
record that precise missing row and do not replace it by finite
range of the bare head.
\end{decision}

\begin{firewall}[Claim boundary after V11]
\label{fire:v15v11}
The enlarged-cluster theorem is rigorous and refresh free.
Application to the literal RPESM sequence still requires a
uniform motif cover, tail influence, sector kernel, and physical
clock.  None is inferred from bare-action locality alone.
\end{firewall}

\section*{Version V11 append-only record}
\addcontentsline{toc}{section}{Version V11 append-only record}
The complete V10 source was retained before this continuation.  It
had \(89506\) bytes, \(2424\) lines, and SHA--256
\[
 \texttt{9AB5C6EDC4E3A0BDC18C2306258C0BC7E31AD843115C071D4D2D2F1C544AEC3A}.
\]
V11 proved exact cancellation of internal tunnelling barriers from
cluster boundary influence and combined that fact with the
archived fractional-coverage theorem to obtain a uniformly local
refresh-free cluster-gap compiler.


\section{Fifteenth-cycle V12: literal motif ledger, finite coloring, and the missing tail norm}
\label{sec:v15v12}

V11 reduced the cluster construction to a support ledger and a
boundary-influence estimate.  The source separates a local
classical head from the determinant/control and Wilsonian tail
rows.  We record those rows without identifying their different
notions of locality.

\begin{ledger}[Literal RPESM motif classes]
\label{led:v15v12-motifs}
The source supplies the following classes.

\begin{enumerate}
\item \emph{Strictly bounded bosonic head.}
      In \({\rm(FS.8)}\), a Higgs potential is supported at one
      vertex, a covariant Higgs difference is supported on one
      link and its two endpoint fields, and a Wilson term is
      supported on the links in one plaquette.  The classical
      \({\rm(CA.12)}\) action has the same vertex--edge--plaquette
      support on its principal chart.
\item \emph{Gauge-fixing quartet.}
      On the protected positive Faddeev--Popov chart,
      \({\rm(FS.31)}\) proves that the complete normalized
      bosonic gauge-fixing and ghost integrals multiply to one.
      After the quartet is integrated as declared, it contributes
      no residual positive-density motif.  Keeping only one member
      of the quartet would invalidate this cancellation.
\item \emph{Fermionic and control modulus.}
      The protected finite cylinder contains
      \(\|\operatorname{Det}D_{\chi,n}\|^2\).  The projective
      RPESM cylinder instead uses the positive control-section
      modulus
      \(\|\sigma_n^+\|^2|p_n^+|^2\), while retaining the physical
      determinant and Pfaffian lines as marked data.  These scalar
      moduli are functions of complete finite operators; finite
      matrix size does not make their spatial support uniformly
      local.
\item \emph{Mediator, route, geometry, and counterterm head.}
      Formula \({\rm(RP.5)}\) declares these interactions finite
      range at each cutoff, but does not list one cutoff-uniform
      support radius and incidence number for every term.
\item \emph{Exact shell counterterm and Wilsonian tail.}
      Exact marginalization introduces
      \(+\log Z_{n+1/n}(x)\).  The source proves the nonlinear
      log-Laplace cocycle and states that the correct Wilsonian
      object is a finite declared head plus a source-minimal
      quasilocal tail.  It does not provide a uniform exponential
      or summable boundary-influence norm for the unbounded RPESM
      sequence.
\item \emph{Global marked labels.}
      Determinant/Pfaffian orientation, divisor order, topological
      sector, global line holonomy, and superselection labels are
      not ordinary local positive-action motifs.  They must be
      fixed by sector selection, transported as line data, or
      proved OS-null; otherwise Proposition
      \ref{prop:v15v11-uncaptured} applies.
\end{enumerate}
\end{ledger}

We now quantify what follows from the first class alone.

\begin{theorem}[Finite motif coloring and coverage]
\label{thm:v15v12-coloring}
Let \({\cal A}\) be a family of bounded interaction motifs on a
finite field-coordinate set \(\Lambda\).  Suppose
\[
 |A|\le s\quad(A\in{\cal A}),\qquad
 \#\{A\in{\cal A}:v\in A\}\le m\quad(v\in\Lambda).            \tag{12.1}
\]
Add a singleton motif for any coordinate not already covered.
Then the intersection graph of the motifs has maximum degree at
most
\[
 \Delta_{\cal A}\le s(m-1),                                  \tag{12.2}
\]
and hence admits a proper coloring with
\[
 \boxed{\chi({\cal A})\le s(m-1)+1.}                         \tag{12.3}
\]
If each motif is updated at rate at least \(r_0>0\), the
rate-weighted fractional coverage of the associated cluster family
satisfies
\[
 \boxed{\kappa_*\ge r_0.}                                    \tag{12.4}
\]
Motifs of one color are disjoint and can be scheduled without
simultaneous coordinate writes; a parallel conditional heat bath
additionally requires conditional independence.
\end{theorem}

\begin{proof}
Fix \(A\in{\cal A}\).  Each of its at most \(s\) coordinates
belongs to at most \(m-1\) other motifs, so \(A\) has at most
\(s(m-1)\) neighbors in the intersection graph.  Greedy coloring
gives \((12.3)\).  Every coordinate lies in at least one cluster
whose rate is at least \(r_0\); using that cluster with unit
fraction in the coverage linear program gives \((12.4)\).
Disjoint motifs share no coordinate, which proves the scheduling statement.  Commutation of their conditional projections is not inferred from support disjointness alone.
\end{proof}

\begin{corollary}[Bare plaquette--Higgs head]
\label{cor:v15v12-bare-cover}
Assume the finite cell complexes have:
\begin{enumerate}
\item plaquette boundary length at most \(\ell_*\);
\item vertex and link incidence bounded by \(d_*\);
\item update rates bounded below by \(r_0\).
\end{enumerate}
The bare action \((8.1)\)--\((8.2)\) has a cluster cover with
\[
 s\le\max\{3,\ell_*\},\qquad m\le m(d_*,\ell_*),              \tag{12.5}
\]
a cutoff-independent finite color count from \((12.3)\), and
coverage at least \(r_0\).  This conclusion concerns the bare
head only.
\end{corollary}

\begin{proof}
The site potential uses one coordinate block.  A Higgs edge uses
the link and its two endpoints.  A plaquette uses at most
\(\ell_*\) links.  Bounded cell incidence gives a finite number
\(m(d_*,\ell_*)\) of such motifs through any coordinate.  Apply
Theorem~\ref{thm:v15v12-coloring}.
\end{proof}

A quantitative quasilocal tail can be added to this cover.  Let
\(\Delta_{BC}^{\rm tail}\) be the likelihood-ratio oscillation
\((11.4)\) contributed by the tail when the exterior cluster
\(C\) changes.  On a cluster graph with distance \(d(B,C)\),
define
\[
 \eta_\omega^{\rm tail}
 :=\sup_B\sum_C
 e^{\omega d(B,C)}
 \frac{\Delta_{BC}^{\rm tail}}4,
 \qquad \omega>0.                                           \tag{12.6}
\]
For \(\omega=0\), use the same definition without the exponential
weight.

\begin{theorem}[Weighted tail-to-influence bound]
\label{thm:v15v12-tail-bound}
If \(\eta_\omega^{\rm tail}<\infty\), then the total-variation
influence of tail interactions beyond cluster distance \(R\)
obeys
\[
 \boxed{
 \sup_B\sum_{C:d(B,C)>R}c_{BC}^{\rm tail}
 \le e^{-\omega R}\eta_\omega^{\rm tail}.}                  \tag{12.7}
\]
If a bounded head has post-constraint row constant
\(\alpha_{\rm head}<1\), then a sufficient full-action condition
is
\[
 \alpha_{\rm head}+\eta_0^{\rm tail}<1,                      \tag{12.8}
\]
or, after absorbing all tail terms of range at most \(R\) into
enlarged clusters,
\[
 \alpha_{\rm head}^{(R)}
 +e^{-\omega R}\eta_\omega^{\rm tail}<1.                    \tag{12.9}
\]
Under the remaining hypotheses of V11, either inequality supplies
the influence row in the refresh-free cluster-gap theorem.
\end{theorem}

\begin{proof}
The likelihood-ratio lemma gives
\[
 c_{BC}^{\rm tail}
 \le\tanh(\Delta_{BC}^{\rm tail}/4)
 \le\Delta_{BC}^{\rm tail}/4.                               \tag{12.10}
\]
For \(d(B,C)>R\),
\[
 \frac{\Delta_{BC}^{\rm tail}}4
 \le e^{-\omega R}
 e^{\omega d(B,C)}
 \frac{\Delta_{BC}^{\rm tail}}4.
\]
Summing and taking the supremum proves \((12.7)\).  Influence
matrices add under addition of logarithmic potentials, which gives
\((12.8)\).  Moving every range-\(R\) term into the internal
cluster action invokes the exact cancellation theorem V11; the
remaining row is bounded by \((12.7)\), proving \((12.9)\).
\end{proof}

\begin{proposition}[Finite range does not control a cofinal tail]
\label{prop:v15v12-finite-range-no-tail}
There are interaction families for which every cutoff action has
finite range, yet no uniform \(\eta_\omega^{\rm tail}\) exists.
For example, on a growing one-dimensional lattice let
\[
 \Psi_n(x)=\sum_{r=1}^{n}\frac1r
 \sum_i \phi(x_i,x_{i+r}),                                   \tag{12.11}
\]
with a bounded nonconstant \(\phi\).  Each \(\Psi_n\) has finite
range \(n\), while every unweighted per-site interaction sum grows
like \(\sum_{r\le n}r^{-1}\), and every positive exponential
weight grows faster.
\end{proposition}

\begin{proof}
At fixed \(n\), only finitely many distances occur.  A fixed site
participates in two interactions at each distance \(r\), with
oscillation proportional to \(1/r\).  The harmonic sum diverges,
and inserting \(e^{\omega r}\) cannot improve convergence.
\end{proof}

\begin{assessment}[Source population after the ledger]
\label{ass:v15v12-source}
The literal source now populates:
\begin{enumerate}
\item a bounded bare vertex--edge--plaquette motif family on its
      regular finite complexes;
\item exact quartet cancellation on the protected chart;
\item finite-cutoff conditional laws and exact projective shell
      transport.
\end{enumerate}
It does not populate, for the unbounded projective sequence:
\begin{enumerate}
\item uniform \(s,m,r_0\) for every mediator, route, geometry, and
      control-modulus term;
\item a uniform fermionic/control-section locality estimate;
\item \(\eta_\omega^{\rm tail}<\infty\) or an equivalent
      target-tail condition for the shell counterterm;
\item local update support for the marked global line and sector
      labels.
\end{enumerate}
The source word \emph{quasilocal} records the intended support
class but does not supply the numerical inequality \((12.6)\).
\end{assessment}

\begin{decision}[V13 fermionic resolvent locality]
\label{dec:v15v12-next}
Attack the first non-bosonic tail quantitatively.  For a
finite-range chiral or Hermitian fermion operator with a uniform
singular or spectral gap, prove a Combes--Thomas inverse decay
bound.  Insert it into the first and mixed second derivatives of
\(\log\det(D^*D)\) and derive an exponential boundary-influence
norm of the form \((12.6)\).  Then audit whether the protected
source floors and range bounds are uniform on the unbounded RPESM
sequence.  If the physical operator has zeros, separate its
zero-mode line before applying the hard-sector estimate.
\end{decision}

\begin{firewall}[Claim boundary after V12]
\label{fire:v15v12}
The bare motif family has a finite-color cover and the exact
weighted tail norm sufficient for V11 is identified.  The source
does not yet provide that norm for the determinant/control modulus
or exact shell tail, so the full cluster gap remains unpopulated.
\end{firewall}

\section*{Version V12 append-only record}
\addcontentsline{toc}{section}{Version V12 append-only record}
The complete V11 source was retained before this continuation.  It
had \(98741\) bytes, \(2674\) lines, and SHA--256
\[
 \texttt{F648C7F4102DC66EA270BBC89EE29CBF532043852A6220366D8DDD031A5E715B}.
\]
V12 classified the literal motif supports, proved the finite
coloring and coverage bound for the bare head, and isolated a
weighted boundary-influence norm whose absence prevents a
quasilocal-tail gap claim.




\section{Fifteenth-cycle V13: singular-gap resolvent decay and determinant boundary influence}
\label{sec:v15v13}

V12 identified the determinant/control modulus as the first
non-bosonic contribution requiring a quantitative tail norm.  We
now prove that a uniform singular gap and weighted locality of the
underlying finite operator give the required decay.  Normality is
not assumed.

Let \((\Lambda,d)\) be a finite graph and
\[
 {\cal H}_\Lambda=\bigoplus_{x\in\Lambda}\mathbb C^q.
\]
For an operator \(D=(D_{xy})\), define the weighted Schur norm
\[
 \|D\|_{\alpha_0,{\rm S}}
 =\max\left\{
 \sup_x\sum_y e^{\alpha_0d(x,y)}\|D_{xy}\|,
 \sup_y\sum_x e^{\alpha_0d(x,y)}\|D_{xy}\|
 \right\}.                                                   \tag{13.1}
\]

\begin{theorem}[Nonnormal Combes--Thomas inverse bound]
\label{thm:v15v13-ct}
Assume
\[
 s_{\min}(D)\ge\tau>0,\qquad
 \|D\|_{\alpha_0,{\rm S}}\le K.                              \tag{13.2}
\]
For \(0<\alpha<\alpha_0\), put
\[
 \varepsilon_{\alpha,\alpha_0}
 =\sup_{r\ge0}(e^{\alpha r}-1)e^{-\alpha_0r}.                 \tag{13.3}
\]
If
\[
 K\varepsilon_{\alpha,\alpha_0}\le\tau/2,                    \tag{13.4}
\]
then for all subsets \(A,B\subset\Lambda\),
\[
 \boxed{
 \|P_AD^{-1}P_B\|
 \le\frac2\tau e^{-\alpha d(A,B)}.}                          \tag{13.5}
\]
For a range-\(R_0\) operator with unweighted Schur norm at most
\(K_0\), condition \((13.4)\) may be replaced by
\[
 K_0(e^{\alpha R_0}-1)\le\tau/2.                             \tag{13.6}
\]
\end{theorem}

\begin{proof}
Fix \(A\) and let
\[
 (M_\alpha f)(x)=e^{-\alpha d(x,A)}f(x),\qquad
 D_\alpha=M_\alpha DM_\alpha^{-1}.
\]
The block multiplier in \(D_\alpha-D\) has absolute value at most
\[
 e^{\alpha|d(x,A)-d(y,A)|}-1
 \le e^{\alpha d(x,y)}-1.
\]
The Schur test and \((13.1)\) therefore give
\[
 \|D_\alpha-D\|
 \le K\varepsilon_{\alpha,\alpha_0}.                         \tag{13.7}
\]
Since \(\|D^{-1}\|\le\tau^{-1}\), the Neumann lemma and
\((13.4)\) give
\[
 \|D_\alpha^{-1}\|\le2\tau^{-1}.                             \tag{13.8}
\]
Because
\[
 D^{-1}=M_\alpha^{-1}D_\alpha^{-1}M_\alpha,
\]
\(M_\alpha^{-1}=I\) on \(A\), while
\(\|M_\alpha P_B\|\le e^{-\alpha d(A,B)}\).  This proves
\((13.5)\).  If \(D_{xy}=0\) for \(d(x,y)>R_0\), the same Schur
calculation gives
\(\|D_\alpha-D\|\le K_0(e^{\alpha R_0}-1)\), proving
\((13.6)\).
\end{proof}

Let \(D(z)\) depend smoothly on local field parameters.  On a
zero-free chart define
\[
 F(z)=\log\det(D(z)^*D(z)).                                  \tag{13.9}
\]
For parameter directions \(u,v\),
\[
 \partial_uF=2\operatorname{Re}\operatorname{Tr}(D^{-1}D_u),
                                                                  \tag{13.10}
\]
\[
 \partial_v\partial_uF
 =2\operatorname{Re}\operatorname{Tr}
 \left(D^{-1}D_{uv}-D^{-1}D_vD^{-1}D_u\right).               \tag{13.11}
\]

\begin{theorem}[Exponential mixed determinant response]
\label{thm:v15v13-det-hessian}
Assume the hypotheses of
Theorem~\ref{thm:v15v13-ct} hold uniformly in \(z\).
Suppose directions \(u,v\) are supported in sets \(A,B\) such
that
\[
 D_u=P_AD_uP_A,\qquad D_v=P_BD_vP_B,\qquad D_{uv}=0,          \tag{13.12}
\]
and
\[
 \|D_u\|\le K_u,\qquad\|D_v\|\le K_v.
\]
Then
\[
 \boxed{
 |\partial_v\partial_uF|
 \le
 \frac{8\,\min\{\operatorname{rank}P_A,
                 \operatorname{rank}P_B\}}
      {\tau^2}
 K_uK_v\,e^{-2\alpha d(A,B)}.}                              \tag{13.13}
\]
If \(D_{uv}\) has bounded support only when
\(d(A,B)\le R_1\), the same estimate holds beyond \(R_1\);
inside that range its direct trace is absorbed into the bounded
head.
\end{theorem}

\begin{proof}
Under \((13.12)\), cyclicity of the finite trace rewrites the
second term of \((13.11)\) as
\[
 \operatorname{Tr}\left(
 P_AD^{-1}P_BD_vP_BD^{-1}P_AD_u
 \right).                                                    \tag{13.14}
\]
Bound the trace by the smaller support rank times the product of
the four operator norms.  Each cross-support inverse is at most
\(2\tau^{-1}e^{-\alpha d(A,B)}\) by \((13.5)\).
The factor two in \((13.11)\) gives \((13.13)\).
The last assertion follows by retaining \(D_{uv}\) among the
finite-range head terms.
\end{proof}

We now convert the differential estimate into the likelihood-ratio
oscillation used in V12.  Let each compact parameter block \(B\)
have intrinsic path diameter at most \(L\).  Suppose its derivative
support has rank at most \(N_0\), its first derivative norm is at
most \(K_1\), and enlarging from parameter support to operator
support costs at most graph radius \(R_D\).

\begin{corollary}[Determinant contribution to the tail norm]
\label{cor:v15v13-tail}
For two parameter blocks \(B,C\) with
\(d(B,C)>2R_D+R_1\), the determinant potential
\(-F\) has boundary oscillation
\[
 \boxed{
 \Delta_{BC}^{F}
 \le
 \frac{8N_0L^2K_1^2}{\tau^2}
 e^{-2\alpha(d(B,C)-2R_D)}.}                                \tag{13.15}
\]
If the block graph has polynomial sphere growth
\[
 \#\{C:r\le d(B,C)<r+1\}\le C_{\rm gr}(1+r)^{d_0-1},          \tag{13.16}
\]
then for every \(0<\omega<2\alpha\),
\[
 \eta_\omega^{F}<\infty                                     \tag{13.17}
\]
with a bound depending only on
\(\tau,\alpha,N_0,L,K_1,R_D,R_1,C_{\rm gr},d_0,\omega\),
not on \(|\Lambda|\).
\end{corollary}

\begin{proof}
Choose paths of length at most \(L\) joining two values in each
parameter block.  The oscillation of
\[
 F(u,v)-F(u,v')
\]
as \(u\) varies is the integral of the mixed derivative over the
resulting parameter rectangle.  Its area is at most \(L^2\).
The operator supports are separated by at least
\(d(B,C)-2R_D\), so \((13.13)\) gives \((13.15)\).
Multiplying by \(e^{\omega d(B,C)}\) leaves exponential decay
\(e^{-(2\alpha-\omega)d(B,C)}\).  Its sum against the polynomial
sphere bound \((13.16)\) converges uniformly, proving
\((13.17)\).
\end{proof}

\begin{corollary}[Insertion into the cluster gap]
\label{cor:v15v13-cluster}
Suppose the bounded head and captured motifs of V11 have influence
constant \(\alpha_{\rm head}<1\), and all remaining positive
determinant-modulus dependence satisfies the hypotheses above.
Then its influence debit is bounded by
\[
 \alpha_F
 \le\frac14\sup_B\sum_C\Delta_{BC}^{F}.                      \tag{13.18}
\]
If
\[
 \alpha_{\rm head}+\alpha_F+\alpha_{\rm other}<1,             \tag{13.19}
\]
where \(\alpha_{\rm other}\) is the separately proved shell and
control-tail debit, the V11 refresh-free cluster gap follows.
\end{corollary}

\begin{proof}
Use \(\tanh x\le x\), sum \((13.15)\), and apply
Theorem~\ref{thm:v15v12-tail-bound} followed by V11.
\end{proof}

Physical zero modes require a separate statement.

\begin{proposition}[Hard determinant and zero-mode line separation]
\label{prop:v15v13-zero}
Let \(Q(z)\) be a constant-rank projection and suppose the hard
operator
\[
 D_Q(z)=Q(z)D(z)Q(z)|_{\operatorname{Ran}Q(z)}
\]
has singular floor \(\tau\).  The preceding theorems apply to
\(\det(D_Q^*D_Q)\) only after the locality of \(Q\), \(Q_u\), and
\(Q_{uv}\) has been included in the derivative support bounds.
The complementary kernel determinant line, its orientation, and
its transport are not part of the positive hard determinant and
cannot be assigned the gap \(\tau\).
\end{proposition}

\begin{proof}
Differentiating \(D_Q=QDQ\) produces terms containing derivatives
of \(Q\).  They enter \(D_u,D_{uv}\) and must satisfy the same
weighted support estimates used in \((13.12)\).  On the
complement, \(D\) has zero singular values, so
Theorem~\ref{thm:v15v13-ct} is inapplicable.  Exterior powers of
that complement form the separately transported zero-mode line.
\end{proof}

\begin{assessment}[Literal source applicability]
\label{ass:v15v13-source}
The protected finite shell states a chiral singular floor
\(\tau_*>0\), a Wilson-kernel window, compact Higgs parameters,
and finite-screen derivative bounds.  These are compatible with
V13, but the overlap operator requires its own weighted locality
estimate, and the theorem quoted in the source does not report the
uniform constants \((13.1)\)--\((13.4)\) for the unbounded
projective sequence.

The RPESM population uses massive positive control sections and
transports physical kernels, divisors, and orientations exactly.
Its population theorem does not give:
\begin{enumerate}
\item a cofinal singular floor for the full positive control
      operator;
\item a uniform weighted Schur norm;
\item local derivative support for the control modulus;
\item weighted locality of the hard projection across physical
      zero modes.
\end{enumerate}
Thus V13 supplies the missing implication but does not assert its
premises for the literal RPESM regulator.
\end{assessment}

\begin{decision}[V14 overlap-sign locality or soft-mode return]
\label{dec:v15v13-next}
Apply the same conjugation method first to the finite-range
Hermitian Wilson kernel \(H_W\) with spectral gap \(g_*\).
Use an integral representation of
\(\operatorname{sgn}H_W\) to derive a weighted Schur bound for the
overlap operator and its first two local field derivatives.
Then combine it with the chiral floor \(\tau_*\) on the protected
chart.  If either floor collapses along the unbounded projective
sequence, extract the normalized soft vector and return it as a
physical zero-mode or concentration witness instead of claiming
determinant locality.
\end{decision}

\begin{firewall}[Claim boundary after V13]
\label{fire:v15v13}
A singular gap plus weighted operator locality now yields
exponential inverse decay and a finite determinant boundary-
influence norm.  The protected source lists related finite-chart
floors, but the required cofinal constants for the unbounded
RPESM control modulus remain unpopulated.
\end{firewall}

\section*{Version V13 append-only record}
\addcontentsline{toc}{section}{Version V13 append-only record}
The complete V12 source was retained before this continuation.  It
had \(109613\) bytes, \(2944\) lines, and SHA--256
\[
 \texttt{6158090580BB67FAB0484A43191FD0BB0C8EC0C706BBF36D815FCC758C988226}.
\]
V13 proved nonnormal Combes--Thomas inverse decay, converted it
into exponential mixed determinant response, and populated the
determinant part of the V12 weighted tail norm under explicit
uniform hard-sector hypotheses.

\section{Fifteenth-cycle V14: Wilson-window locality of the overlap sign}
\label{sec:v15v14}

V13 reduced the determinant-tail problem to weighted locality and a
hard singular-value floor.  The protected source supplies a uniform
spectral window for the Hermitian Wilson kernel but does not state the
weighted constants needed there.  This continuation derives those
constants directly.  A contour proof is preferable to a real-axis
integral here: it is valid for the nonnormal weighted conjugate and
gives the first two field derivatives without a divergent absolute
resolvent integral.

Let \(X=(\mathcal V,\mathcal E)\) be a finite graph and
\(\mathcal H=\bigoplus_{x\in\mathcal V}\mathcal H_x\).  For
\(A\subset\mathcal V\), write \(P_A\) for the corresponding
orthogonal projection.  Let \(H=H^*\) satisfy
\[
 \operatorname{spec}H\subset[-M,-g]\cup[g,M],
 \qquad 0<g\le M,                                           \tag{14.1}
\]
and suppose \(H_{xy}=0\) for \(d(x,y)>R\).  Set
\[
 k_0=\max\left\{
 \sup_x\sum_y\|H_{xy}\|,\,
 \sup_y\sum_x\|H_{xy}\|
 \right\}.                                                   \tag{14.2}
\]
For \(\alpha>0\), define
\[
 \delta_\alpha=k_0(e^{\alpha R}-1).                          \tag{14.3}
\]

\begin{lemma}[Finite-range weighted conjugation]
\label{lem:v15v14-conjugation}
If \(\rho:\mathcal V\to\mathbb R\) is one-Lipschitz and
\(W_\alpha=e^{\alpha\rho}\), then
\[
 \|W_\alpha H W_\alpha^{-1}-H\|
 \le\delta_\alpha.                                          \tag{14.4}
\]
\end{lemma}

\begin{proof}
The \(xy\) block of the difference is
\[
 \left(e^{\alpha(\rho(x)-\rho(y))}-1\right)H_{xy}.
\]
It vanishes unless \(d(x,y)\le R\), and on that support its
scalar factor has absolute value at most \(e^{\alpha R}-1\).
Both the block row sum and block column sum are therefore at
most \(\delta_\alpha\).  The Schur test gives (14.4).
\end{proof}

Let \(\Gamma\) be the positively oriented rectangle with vertices
\[
 \frac g2\pm i\frac g2,\qquad
 M+\frac g2\pm i\frac g2,
\]
and write
\[
 L_\Gamma=2M+2g.                                             \tag{14.5}
\]
This contour encloses the positive spectrum of \(H\), excludes the
negative spectrum, and remains at distance at least \(g/2\) from
\(\operatorname{spec}H\).

\begin{theorem}[Uniform locality of the overlap sign]
\label{thm:v15v14-sign}
Assume
\[
 \delta_\alpha\le\frac g4.                                  \tag{14.6}
\]
Let
\[
 P_+(H)=\frac{1}{2\pi i}\int_\Gamma(z-H)^{-1}\,dz,
 \qquad
 \varepsilon(H)=\operatorname{sgn}H=2P_+(H)-I.               \tag{14.7}
\]
Then for all \(z\in\Gamma\) and all vertex sets \(A,B\),
\[
 \|P_A(z-H)^{-1}P_B\|
 \le\frac4g e^{-\alpha d(A,B)}.                             \tag{14.8}
\]
For disjoint \(A,B\),
\[
 \boxed{
 \|P_A\varepsilon(H)P_B\|
 \le\frac{4L_\Gamma}{\pi g}
 e^{-\alpha d(A,B)}.}                                      \tag{14.9}
\]
The constants are independent of the number of vertices.
\end{theorem}

\begin{proof}
For \(z\in\Gamma\),
\(\|(z-H)^{-1}\|\le2/g\).  Put
\(H_\alpha=W_\alpha H W_\alpha^{-1}\).  Lemma
\ref{lem:v15v14-conjugation} and (14.6) give
\[
 \|(z-H_\alpha)^{-1}\|\le\frac4g                            \tag{14.10}
\]
by the resolvent identity and a Neumann series.

To prove (14.8), choose
\(\rho(x)=d(x,B)\).  Since \(W_\alpha P_B=P_B\) and
\(\|P_AW_\alpha^{-1}\|\le e^{-\alpha d(A,B)}\), similarity
of the resolvents gives the displayed estimate.  Holomorphic
functional calculus also gives
\[
 W_\alpha P_+(H)W_\alpha^{-1}=P_+(H_\alpha).
\]
Thus
\[
 \|P_+(H_\alpha)\|
 \le\frac{L_\Gamma}{2\pi}\frac4g
 =\frac{2L_\Gamma}{\pi g}.                                 \tag{14.11}
\]
The same weight argument bounds the off-diagonal block of \(P_+\).
Since the identity has zero \(AB\) block for disjoint sets,
\(\varepsilon=2P_+-I\) yields (14.9).
\end{proof}

The derivative statement requires an exact meaning of a local
variation.  Let \(s,t\) be parameters with cores \(S,T\) and assume
\[
 H_s=P_SH_sP_S,\quad \|H_s\|\le K_s,\qquad
 H_t=P_TH_tP_T,\quad \|H_t\|\le K_t.                         \tag{14.12}
\]
If a mixed bare derivative is present, suppose
\(H_{st}=P_QH_{st}P_Q\) and \(\|H_{st}\|\le K_{st}\).

\begin{theorem}[Local first and mixed overlap response]
\label{thm:v15v14-derivatives}
Under the hypotheses of Theorem \ref{thm:v15v14-sign},
\[
 \boxed{
 \|P_X\varepsilon_sP_Y\|
 \le\frac{16L_\Gamma K_s}{\pi g^2}
 e^{-\alpha[d(X,S)+d(S,Y)]}.}                              \tag{14.13}
\]
Moreover,
\[
\begin{aligned}
 \|P_X\varepsilon_{st}P_Y\|
 \le{}&
 \frac{64L_\Gamma K_sK_t}{\pi g^3}
 \bigl[
 e^{-\alpha[d(X,S)+d(S,T)+d(T,Y)]}\\
 &\hspace{35mm}+
 e^{-\alpha[d(X,T)+d(T,S)+d(S,Y)]}
 \bigr]\\
 &+\frac{16L_\Gamma K_{st}}{\pi g^2}
 e^{-\alpha[d(X,Q)+d(Q,Y)]}.                               \tag{14.14}
\end{aligned}
\]
The last line is omitted when \(H_{st}=0\).
\end{theorem}

\begin{proof}
With \(R_z=(z-H)^{-1}\), differentiation under the finite contour
integral gives
\[
 \varepsilon_s
 =\frac{1}{\pi i}\int_\Gamma R_zH_sR_z\,dz,                 \tag{14.15}
\]
and
\[
 \varepsilon_{st}
 =\frac{1}{\pi i}\int_\Gamma
 \left(
 R_zH_tR_zH_sR_z+
 R_zH_sR_zH_tR_z+
 R_zH_{st}R_z
 \right)\,dz.                                               \tag{14.16}
\]
Insert \(P_S\), \(P_T\), and \(P_Q\) at the derivative cores.
Applying (14.8) to each resolvent factor gives respectively two,
three, and two factors of \(4/g\).  Multiplication by the contour
length \(L_\Gamma\) proves (14.13)--(14.14).
\end{proof}

\begin{corollary}[Weighted Schur locality]
\label{cor:v15v14-schur}
Suppose the graph has the polynomial sphere bound (13.16), and the
dimensions of the site fibres are uniformly bounded.  For every
\(0<\beta<\alpha\), the operators
\(\varepsilon\), \(\varepsilon_s\), and
\(\varepsilon_{st}\) have finite \(\beta\)-weighted block Schur
norms.  For the derivatives, the bounds are uniform after centering
at their displayed cores.  They depend only on
\[
 g,M,R,k_0,\alpha,\beta,K_s,K_t,K_{st}
\]
and the graph-growth constants.
\end{corollary}

\begin{proof}
Multiply (14.9), (14.13), or (14.14) by the
\(\beta\)-weight and sum over polynomially many shells.  Each sum
is dominated by a constant times
\[
 \sum_{r\ge0}(1+r)^{d_0-1}e^{-(\alpha-\beta)r}<\infty.
\]
For (14.14), use the triangle inequality for graph distance along
each displayed chain.
\end{proof}

This establishes the missing operator-locality implication for the
protected finite shell:
\[
 g_*I\preceq |H_{W,n}|\preceq M_*I
 \quad+\quad
 \sup_n R_n<\infty,\ \sup_n k_{0,n}<\infty
 \quad\Longrightarrow\quad
 \operatorname{sgn}H_{W,n}\ \hbox{is uniformly local}.       \tag{14.17}
\]
The hopping range and Schur bound in (14.17) are automatic for a
Wilson stencil with uniformly bounded compact link matrices, once
all lattice graphs use the same combinatorial distance.  The
moving projector
\[
 \widehat P_{-,n}=\frac12(I+\operatorname{sgn}H_{W,n})        \tag{14.18}
\]
and its first two local derivatives inherit the same estimates.

\begin{proposition}[Cutoff-scaling ledger for the overlap block]
\label{prop:v15v14-scaling}
For
\[
 D_{{\rm ov},n}=a_n^{-1}(I+\gamma_5\varepsilon_n),           \tag{14.19}
\]
raw dimensionless-link derivatives carry the prefactor
\(a_n^{-1}\).  Consequently a determinant-influence estimate based
on V13 requires uniform bounds on ratios such as
\[
 \sup_n\frac{K_{1,n}}{\tau_n}<\infty,\qquad
 \sup_n\frac{K_{2,n}}{\tau_n}<\infty,                        \tag{14.20}
\]
rather than only the fixed lower bound \(\tau_n\ge\tau_*>0\).
For a physical one-form variation whose link increment contains a
factor \(a_n\), that factor can cancel (14.19); the cancellation
does not by itself control arbitrary Gibbs-block oscillations.
\end{proposition}

\begin{proof}
Equations (14.13)--(14.14) bound derivatives of the dimensionless
sign.  Differentiating (14.19) therefore multiplies the raw bounds
by \(a_n^{-1}\).  The determinant Hessian contains the dimensionless
combinations \(D^{-1}D_s\) and \(D^{-1}D_{st}\), so the relevant
uniform quantities are the ratios in (14.20).  If the parameter is
a discretized physical one-form, its insertion has size \(a_n\);
the two factors cancel.  This is the algebraic cancellation
displayed in the source's physical overlap-current formula, but
that formula is restricted to its routed physical source bank.
\end{proof}

\begin{assessment}[What V14 does and does not populate]
\label{ass:v15v14-source}
The literal protected shell assumes (14.1) with
\(g=g_*\), \(M=M_*\), defines (14.18)--(14.19), and separately
assumes \(s_{\min}(D_{\chi,n})\ge\tau_*\).  Therefore V14 supplies
the previously missing overlap-sign locality on every such
protected chart, provided the standard Wilson stencil bounds in
(14.17) are retained uniformly.

Two gaps remain before Corollary \ref{cor:v15v13-tail} can be
applied to the literal chiral determinant.  First, derivatives of
the overlap sign are quasilocal around a core, whereas V13 used an
exact finite derivative support.  Second, the source does not state
the cofinal ratios (14.20) for all local Gibbs-coordinate
variations.  Neither gap is filled merely by compactness at each
fixed cutoff.
\end{assessment}

\begin{decision}[V15 weighted trace extension and compulsory audit]
\label{dec:v15v14-next}
At V15, first inspect the newest Yang--Mills and Navier--Stokes
companion notes.  Then replace the exact-support step in V13 by a
weighted-Schur trace lemma adapted to the core-factorized contour
form (14.15)--(14.16).  This will determine the precise influence
debit of the overlap determinant.  If the required cutoff ratios
fail, normalize a blowing-up derivative or a soft singular vector
and return upstream rather than closing the estimate circularly.
\end{decision}

\begin{firewall}[Claim boundary after V14]
\label{fire:v15v14}
A gapped finite-range Hermitian Wilson kernel now yields uniform
exponential locality of its sign, moving chiral projector, and
first two local derivatives, with explicit constants.  Uniform
locality of the full chiral determinant still requires a weighted
trace argument and the cutoff ratios in (14.20).
\end{firewall}

\section*{Version V14 append-only record}
\addcontentsline{toc}{section}{Version V14 append-only record}
The complete V13 source was retained before this continuation.  It
had \(119854\) bytes, \(3247\) lines, and SHA--256
\[
 \texttt{AF25CC72E64D28DE83BF3DAE390557A351CC3CF0874EB656BBF80CCD558D6AB1}.
\]
V14 derived finite-volume-uniform overlap-sign locality and its
first two local response estimates from the protected Wilson
window, then isolated the remaining weighted-trace and cutoff-
scaling obligations.
\section{Fifteenth-cycle V15: companion audit and weighted trace cycles}
\label{sec:v15v15}

This compulsory checkpoint reads the newest companion notes before
continuing the determinant-tail argument.  It then removes the
exact-derivative-support hypothesis left in V13.  The replacement is
a trace-cycle estimate: two local variation cores force two
exponentially decaying bridges around every determinant Hessian
trace, even when all operators between those cores are quasilocal.

\subsection{Companion checkpoint}

\begin{record}[Files inspected at V15]
\label{rec:v15v15-companions}
The current Yang--Mills companion is
\begin{center}
\path{Renewal_Yang_Mills_Mass_Gap_Research_Notes_V35.tex},
\end{center}
with \(289778\) bytes, \(7915\) lines, and SHA--256
\[
 \texttt{4EA177BF18DD0ED14474AE54E22DE26FF178C904813F84F53D135BE7DA989C04}.
                                                               \tag{15.1}
\]
The current Navier--Stokes companion is
\begin{center}
\path{Renewal_NS_Stability_Research_Notes_V26.tex},
\end{center}
with \(278283\) bytes, \(5319\) lines, and SHA--256
\[
 \texttt{82C887F8C2C69E4F28FAD7C3DC8DE5B9099CB5A4BF431EBA1FFF3E91935978AD}.
                                                               \tag{15.2}
\]
\end{record}

YM V34 proves pointwise coverage of one three-dimensional phase
cube by eight certified charts.  YM V35 then uses exact trace-square
invariants to exclude an assumed full transverse permutation
symmetry.  The transferable rule is that chart adjacency is weaker
than pointwise coverage and apparent symmetry cannot transport a
floor until an intertwiner or invariant test justifies it.  No YM
pencil or floor is a chiral-determinant estimate.

NS V26 computes norms restricted to rank-one variations and finds
that they can be strictly smaller than norms over all admissible
inputs.  The transferable rule is directly relevant to V14:
physical-source derivative bounds cannot be substituted for
arbitrary Gibbs-block derivative bounds.  No NS kernel constant
acts on the Standard Model field space.

\begin{audit}[V15 companion verdict]
\label{aud:v15v15-verdict}
No numerical constant or theorem is imported from either companion.
The SM proof will retain two disciplines: prove actual protected-
chart coverage before taking a uniform infimum, and keep physical
source variations distinct from the full local coordinate class
used in a Gibbs influence coefficient.
\end{audit}

\subsection{The weighted Schur trace mechanism}

For \(\beta>0\), let \(\mathfrak S_\beta\) be the algebra of block
operators with norm
\[
 \|A\|_\beta
 =
 \max\left\{
 \sup_x\sum_y e^{\beta d(x,y)}\|A_{xy}\|,\,
 \sup_y\sum_x e^{\beta d(x,y)}\|A_{xy}\|
 \right\}.                                                   \tag{15.3}
\]

\begin{lemma}[Weighted Schur algebra]
\label{lem:v15v15-algebra}
For \(A,B\in\mathfrak S_\beta\),
\[
 \|AB\|_\beta\le\|A\|_\beta\|B\|_\beta.                      \tag{15.4}
\]
For vertex sets \(S,T\),
\[
 \|P_SAP_T\|
 \le e^{-\beta d(S,T)}\|A\|_\beta.                           \tag{15.5}
\]
\end{lemma}

\begin{proof}
Use \(d(x,z)\le d(x,y)+d(y,z)\), expand the block row sum
of \(AB\), and apply the two row bounds.  The column calculation is
identical, proving (15.4).  Restricting the weighted row and column
sums to \(S\times T\) removes a factor at least
\(e^{\beta d(S,T)}\); the ordinary Schur test gives (15.5).
\end{proof}

\begin{lemma}[Two-core trace cycle]
\label{lem:v15v15-cycle}
Let
\[
 K_S=P_SK_SP_S,\qquad K_T=P_TK_TP_T,
\]
and suppose both cores have rank at most \(r\).  Then
\[
 \boxed{
 |\operatorname{Tr}(A K_S B K_T)|
 \le
 r e^{-2\beta d(S,T)}
 \|A\|_\beta\|B\|_\beta\|K_S\|\|K_T\|.}                    \tag{15.6}
\]
\end{lemma}

\begin{proof}
Insert the support projections and cycle the trace:
\[
 \operatorname{Tr}(A K_S B K_T)
 =
 \operatorname{Tr}\left(
 P_TAP_SK_SP_SBP_TK_T
 \right).
\]
The trace is supported on a space of rank at most \(r\).  Bound it
by rank times operator norm and apply (15.5) once to each of the
two bridges \(P_TAP_S\) and \(P_SBP_T\).
\end{proof}

We encode quasilocal derivatives in a form preserved by contour
calculus.  A first derivative with core \(S\) is core-factorized if
\[
 D_s=\int A_{s,\xi}K_{s,\xi}B_{s,\xi}\,d\nu_s(\xi),
 \qquad
 K_{s,\xi}=P_SK_{s,\xi}P_S,                                 \tag{15.7}
\]
where the integral may be a finite sum, and
\[
 {\cal L}_s(D)
 :=
 \int
 \|A_{s,\xi}\|_\beta\|K_{s,\xi}\|
 \|B_{s,\xi}\|_\beta\,d|\nu_s|(\xi)<\infty.                 \tag{15.8}
\]
For disjoint cores \(S,T\), a mixed derivative is two-core
factorized if it is a sum of terms
\[
 \int A_\xi K_{S,\xi}B_\xi K_{T,\xi}C_\xi\,d\nu(\xi)         \tag{15.9}
\]
and terms with \(S,T\) reversed.  Let
\({\cal M}_{st}(D)\) be the sum of the corresponding integrals of
\[
 \|A_\xi\|_\beta\|K_{S,\xi}\|\|B_\xi\|_\beta
 \|K_{T,\xi}\|\|C_\xi\|_\beta.                              \tag{15.10}
\]

\begin{lemma}[Core-factorization calculus]
\label{lem:v15v15-calculus}
The classes (15.7) and (15.9) are closed under sums, adjoints, and
products with operators in \(\mathfrak S_\beta\).  In particular,
\[
 {\cal L}_s(XY)
 \le {\cal L}_s(X)\|Y\|_\beta+
      \|X\|_\beta{\cal L}_s(Y),                              \tag{15.11}
\]
and, for separated \(S,T\),
\[
\begin{aligned}
 {\cal M}_{st}(XY)
 \le{}&{\cal M}_{st}(X)\|Y\|_\beta+
       \|X\|_\beta{\cal M}_{st}(Y)\\
 &+{\cal L}_s(X){\cal L}_t(Y)
  +{\cal L}_t(X){\cal L}_s(Y).                              \tag{15.12}
\end{aligned}
\]
\end{lemma}

\begin{proof}
Use the Leibniz rules
\[
 (XY)_s=X_sY+XY_s,
\]
and
\[
 (XY)_{st}=X_{st}Y+X_sY_t+X_tY_s+XY_{st}.
\]
Place the undifferentiated factors into the adjacent \(A,B,C\)
terms of (15.7) or (15.9), then use (15.4).  Taking adjoints reverses
the factor order and preserves every norm and core support.
\end{proof}

\begin{theorem}[Determinant Hessian for quasilocal core derivatives]
\label{thm:v15v15-determinant}
Let \(D\) be invertible with \(D,D^{-1}\in\mathfrak S_\beta\), and
put
\[
 G_\beta=\|D^{-1}\|_\beta,\qquad
 F=\log\det(D^*D).                                          \tag{15.13}
\]
Suppose \(s,t\) have separated cores \(S,T\) of rank at most \(r\),
the first derivatives have costs
\({\cal L}_s(D),{\cal L}_t(D)\), and the mixed derivative has cost
\({\cal M}_{st}(D)\).  Then
\[
 \boxed{
 |F_{st}|
 \le
 2r e^{-2\beta d(S,T)}
 \left[
 G_\beta{\cal M}_{st}(D)
 +G_\beta^2{\cal L}_s(D){\cal L}_t(D)
 \right].}                                                  \tag{15.14}
\]
\end{theorem}

\begin{proof}
Jacobi differentiation gives
\[
 F_{st}
 =
 2\operatorname{Re}\operatorname{Tr}
 \left(D^{-1}D_{st}-D^{-1}D_tD^{-1}D_s\right).              \tag{15.15}
\]
For a term in (15.9), cyclicity turns the first trace into
\[
 \operatorname{Tr}
 \left(
 K_S(CD^{-1}A)K_TB
 \right)
\]
or the reversed analogue.  Lemmas
\ref{lem:v15v15-algebra}--\ref{lem:v15v15-cycle} bound its absolute
value by its contribution to
\[
 r e^{-2\beta d(S,T)}G_\beta{\cal M}_{st}(D).
\]
Expanding both first derivatives by (15.7) and cycling the second
trace gives terms
\[
 \operatorname{Tr}\left(
 K_S(B_sD^{-1}A_t)K_T(B_tD^{-1}A_s)
 \right).
\]
The same lemmas give
\[
 r e^{-2\beta d(S,T)}
 G_\beta^2{\cal L}_s(D){\cal L}_t(D).
\]
Integrate against total variation and restore the factor two in
(15.15).
\end{proof}

\begin{corollary}[Boundary oscillation and summability]
\label{cor:v15v15-influence}
If the two compact parameter blocks have path diameters at most
\(L_S,L_T\), then
\[
 \Delta_{ST}^{F}
 \le
 2rL_SL_T e^{-2\beta d(S,T)}
 \left[
 G_\beta{\cal M}_{st}(D)
 +G_\beta^2{\cal L}_s(D){\cal L}_t(D)
 \right].                                                   \tag{15.16}
\]
If the bracketed quantities are uniform and the block graph has
polynomial sphere growth, then the weighted tail norm is finite for
every \(0<\omega<2\beta\).
\end{corollary}

\begin{proof}
Integrate (15.14) over the parameter rectangle, as in Corollary
\ref{cor:v15v13-tail}.  Multiplication by
\(e^{\omega d(S,T)}\) leaves a summable exponential
\(e^{-(2\beta-\omega)d(S,T)}\).
\end{proof}

\subsection{Application to the overlap contour}

Let
\[
 {\cal R}_\beta
 :=
 \sup_{z\in\Gamma}\|(z-H_W)^{-1}\|_\beta.                   \tag{15.17}
\]
Corollary \ref{cor:v15v14-schur} makes this finite, uniformly in
volume, whenever \(0<\beta<\alpha\) and the V14 Wilson constants are
uniform.  Equations (14.15)--(14.16) give the explicit costs
\[
 {\cal L}_s(\varepsilon)
 \le\frac{L_\Gamma}{\pi}{\cal R}_\beta^2K_s,                 \tag{15.18}
\]
and, when the bare mixed kernel derivative vanishes,
\[
 {\cal M}_{st}(\varepsilon)
 \le\frac{2L_\Gamma}{\pi}{\cal R}_\beta^3K_sK_t.             \tag{15.19}
\]
Thus the quasilocal derivative of the sign is exactly of the class
required by Theorem \ref{thm:v15v15-determinant}.

There is also a frame-free way to treat the moving source subspace.
Let \(P=\widehat P_-\), let \(T=TP\) be the full-space extension of
the chiral map from \(\operatorname{Ran}P\), and define
\[
 {\cal A}=T^*T+I-P.                                         \tag{15.20}
\]

\begin{proposition}[Identity-completed moving-domain determinant]
\label{prop:v15v15-completion}
Suppose \(T|_{\operatorname{Ran}P}\) has singular floor \(\tau>0\).
Then
\[
 \det{\cal A}
 =
 \det\nolimits_{\operatorname{Ran}P}(T^*T),                 \tag{15.21}
\]
and
\[
 {\cal A}\succeq \widehat\tau I,\qquad
 \widehat\tau=\min\{1,\tau^2\}.                             \tag{15.22}
\]
If \(P,T\in\mathfrak S_\beta\), their first derivatives are
core-factorized, and their separated mixed derivatives are
two-core factorized, then the same statements hold for
\({\cal A}_s,{\cal A}_{st}\), with costs computed by
(15.11)--(15.12).
\end{proposition}

\begin{proof}
Relative to
\(\mathcal H=\operatorname{Ran}P\oplus\ker P\), the identity
\(T=TP\) gives
\[
 {\cal A}
 =
 \begin{pmatrix}
 T^*T|_{\operatorname{Ran}P}&0\\
 0&I
 \end{pmatrix}.
\]
This proves (15.21)--(15.22).  The derivative assertion follows by
applying Lemma \ref{lem:v15v15-calculus} to
\(T^*T+I-P\).
\end{proof}

For the literal protected chiral block, (15.20) removes any need to
choose a global basis of \(\operatorname{Ran}\widehat P_{-,n}\).
The V14 contour bounds populate the projector part.  Local Yukawa
multiplication is already core-factorized.  Once the cutoff-scaled
costs of \(T\) and the weighted norm of \({\cal A}^{-1}\) are
uniform, (15.14)--(15.16) rigorously populate the chiral
determinant contribution to the V12 tail norm.

\begin{assessment}[Remaining source obligations after V15]
\label{ass:v15v15-source}
The exact-support defect in V13 is closed: overlap derivatives may
be quasilocal, and the two-core trace cycle still supplies twice
the exponential separation.  The moving chiral domain is also
handled intrinsically by (15.20).

The literal source still does not populate three cofinal numerical
inputs:
\begin{enumerate}
\item a pointwise cover of the unbounded regulator population by
      Wilson and chiral protected charts with uniform margins;
\item the Gibbs-coordinate derivative costs in (15.18)--(15.19),
      including the cutoff ratios isolated in (14.20);
\item a uniform weighted norm of \({\cal A}^{-1}\), obtained from
      (15.22), weighted locality of \({\cal A}\), and the V13
      conjugation condition.
\end{enumerate}
Fixed-cutoff compactness does not make these quantities uniform.
\end{assessment}

\begin{decision}[V16 dimensionless cancellation screen]
\label{dec:v15v15-next}
Use the exact identity
\[
 P_+D_{{\rm ov},n}\widehat P_{-,n}
 =\frac{2}{a_n}P_+\widehat P_{-,n}
\]
to compute the local Gibbs-coordinate costs of the completed
operator (15.20), retaining every power of \(a_n\).  Separate link
angle, Higgs, and Yukawa directions.  The outcome must be either a
uniform dimensionless ratio, an explicit additional scaling
hypothesis, or a normalized blowing-up direction; a physical-source
cancellation will not be reused outside its stated variation class.
\end{decision}

\begin{firewall}[Claim boundary after V15]
\label{fire:v15v15}
V15 proves a volume-uniform determinant Hessian bound for
quasilocal core-factorized derivatives and applies it to the overlap
contour.  It does not establish the missing cofinal chart cover or
cutoff-scaled derivative ratios for the source's unbounded RPESM
population.
\end{firewall}

\section*{Version V15 append-only record}
\addcontentsline{toc}{section}{Version V15 append-only record}
The complete V14 source was retained before this continuation.  It
had \(130499\) bytes, \(3552\) lines, and SHA--256
\[
 \texttt{DC79A3DF30322DF9EC34031B52B928753C36766AAB6C18C302768BEEC0C7C41A}.
\]
V15 completed the compulsory companion audit, replaced exact
derivative support by a weighted two-core trace cycle, and gave a
frame-free identity completion of the moving chiral determinant.
\section{Fifteenth-cycle V16: exact cutoff normalization and the chiral soft-mode fork}
\label{sec:v15v16}

V14 retained the inverse lattice spacing in the raw overlap
derivative ledger.  V15 supplied the correct quasilocal trace
estimate.  This continuation now performs the algebra at the
determinant level.  On a fixed-index sector the scalar
\(2/a_n\) can be removed exactly: it changes the determinant by a
field-independent constant and contributes nothing to any boundary
influence.  The surviving question is the singular floor of the
dimensionless normalized chiral map.

Fix one cutoff and suppress \(n\).  Put
\[
 P=\widehat P_-=\frac12(I+\varepsilon),\qquad
 c=\frac2a,\qquad b=c^{-1}=\frac a2.                        \tag{16.1}
\]
Let \(P_+=(I+\gamma_5)/2\).  Since \(\varepsilon P=P\),
\[
 P_+D_{\rm ov}P
 =a^{-1}P_+(I+\gamma_5\varepsilon)P
 =\frac2aP_+P=cP_+P.                                       \tag{16.2}
\]
Regard the Yukawa map \(Y\) as a map on \(\operatorname{Ran}P\);
its full-space extension is \(YP\).  The chiral map and its
dimensionless normalization are therefore
\[
 T=(cP_++Y)P=c\widetilde T,\qquad
 \widetilde T=(P_++bY)P.                                   \tag{16.3}
\]

\begin{theorem}[Exact determinant normalization]
\label{thm:v15v16-normalization}
Suppose \(\operatorname{rank}P=r\) is constant on the protected
sector.  Then
\[
 \boxed{
 \log\det\nolimits_{\operatorname{Ran}P}(T^*T)
 =
 2r\log c+
 \log\det\nolimits_{\operatorname{Ran}P}
       (\widetilde T^*\widetilde T).}                       \tag{16.4}
\]
Consequently every first or mixed field derivative, every
likelihood-ratio oscillation, and every Dobrushin influence
coefficient of the two determinant potentials is identical.
\end{theorem}

\begin{proof}
On the \(r\)-dimensional moving source space,
\(T^*T=c^2\widetilde T^*\widetilde T\).  Taking determinants gives
a factor \(c^{2r}\), proving (16.4).  At fixed cutoff, \(a,c,r\)
are independent of all field coordinates.  Their term
\(2r\log c\) vanishes under every field derivative and cancels from
every conditional likelihood ratio.
\end{proof}

The same conclusion can be expressed without choosing a frame of
the moving source space.  Define
\[
 {\cal A}=T^*T+I-P,\qquad
 \widetilde{\cal A}
 =\widetilde T^*\widetilde T+I-P.                           \tag{16.5}
\]

\begin{corollary}[Normalized identity completion]
\label{cor:v15v16-completion}
One has
\[
 \det{\cal A}=c^{2r}\det\widetilde{\cal A},                 \tag{16.6}
\]
and hence
\[
 \partial_s\partial_t\log\det{\cal A}
 =
 \partial_s\partial_t\log\det\widetilde{\cal A}.            \tag{16.7}
\]
If
\[
 \widetilde\tau=s_{\min}
 \left(\widetilde T|_{\operatorname{Ran}P}\right),
\]
then
\[
 \widetilde{\cal A}\succeq
 \min\{1,\widetilde\tau^2\}I,\qquad
 \widetilde\tau=b\,s_{\min}
 \left(T|_{\operatorname{Ran}P}\right).                    \tag{16.8}
\]
\end{corollary}

\begin{proof}
Use the block decomposition from Proposition
\ref{prop:v15v15-completion}.  On \(\operatorname{Ran}P\), the two
positive blocks differ by \(c^2\); on \(\ker P\), both equal the
identity.  The singular-value identity in (16.8) follows from
\(T=c\widetilde T\).
\end{proof}

\subsection{Local-coordinate derivative ledger}

All derivatives below are field derivatives at fixed cutoff, so
\(a,b,c\) are constants.  Put
\[
 B=P_++bY,\qquad \widetilde T=BP.                           \tag{16.9}
\]
For an arbitrary local coordinate \(s\),
\[
 \widetilde T_s=bY_sP+BP_s,                                \tag{16.10}
\]
and for two coordinates \(s,t\),
\[
 \widetilde T_{st}
 =bY_{st}P+bY_sP_t+bY_tP_s+BP_{st}.                        \tag{16.11}
\]

\begin{proposition}[Core costs after normalization]
\label{prop:v15v16-costs}
In the notation of V15, suppose \(P,Y\in\mathfrak S_\beta\) and
their local derivatives have the required one- and two-core
factorizations.  Then
\[
\begin{aligned}
 {\cal L}_s(\widetilde T)
 \le{}&
 b\,{\cal L}_s(Y)\|P\|_\beta+
 \bigl(\|P_+\|_\beta+b\|Y\|_\beta\bigr){\cal L}_s(P),       \tag{16.12}\\
 {\cal M}_{st}(\widetilde T)
 \le{}&
 b\,{\cal M}_{st}(Y)\|P\|_\beta+
 \bigl(\|P_+\|_\beta+b\|Y\|_\beta\bigr)
       {\cal M}_{st}(P)\\
 &+b\,{\cal L}_s(Y){\cal L}_t(P)
  +b\,{\cal L}_t(Y){\cal L}_s(P).                          \tag{16.13}
\end{aligned}
\]
The completed operator
\(\widetilde{\cal A}=\widetilde T^*\widetilde T+I-P\)
has first- and mixed-derivative costs bounded explicitly from
(16.12)--(16.13) by (15.11)--(15.12).  No negative power of \(a\)
occurs in these bounds.
\end{proposition}

\begin{proof}
Equations (16.10)--(16.11), the triangle inequality, and the
submultiplicativity of Lemma \ref{lem:v15v15-algebra} give
(16.12)--(16.13).  Apply Lemma
\ref{lem:v15v15-calculus} to
\(\widetilde T^*\widetilde T+I-P\).  Every Yukawa derivative is
multiplied by \(b=a/2\), while every projector derivative is
dimensionless.  Thus normalization has removed all \(a^{-1}\)
factors.
\end{proof}

The three relevant variation classes are now distinct.

\begin{enumerate}
\item For a link coordinate, a site-local Yukawa multiplication has
      \(Y_s=0\).  Thus
      \[
       \widetilde T_s=BP_s,\qquad
       \widetilde T_{st}=BP_{st}                            \tag{16.14}
      \]
      for two separated link cores whose bare Wilson mixed
      derivative vanishes.  V14--V15 bound \(P_s,P_{st}\).

\item For a Higgs or Yukawa coordinate, the Wilson projector is
      fixed, so
      \[
       \widetilde T_s=bY_sP.                                \tag{16.15}
      \]
      Separated ultralocal Yukawa coordinates have
      \(Y_{st}=0\), hence no direct mixed derivative.

\item For a link coordinate \(s\) and a separated Higgs coordinate
      \(t\),
      \[
       \widetilde T_{st}=bY_tP_s,                           \tag{16.16}
      \]
      which is a two-core factorization of exactly the V15 type.
\end{enumerate}

\begin{theorem}[Uniform normalized determinant-tail criterion]
\label{thm:v15v16-tail-criterion}
Assume along a cofinal sequence:
\begin{enumerate}
\item the V14 Wilson constants and a choice \(0<\beta<\alpha\) are
      uniform;
\item
      \[
       \inf_n\widetilde\tau_n>0;                            \tag{16.17}
      \]
\item the weighted norms and local derivative costs of
      \(b_nY_n\) are uniform over the field domain used in the
      conditional kernels;
\item the local core ranks, parameter path diameters, and graph
      growth constants are uniform.
\end{enumerate}
Then \(\widetilde{\cal A}_n\) and its first two local derivatives
have uniform weighted Schur and core-factorization bounds.
Moreover, after possibly reducing the weight from \(\beta\) to
\(\beta'>0\), V13 gives
\[
 \sup_n\|\widetilde{\cal A}_n^{-1}\|_{\beta'}<\infty.        \tag{16.18}
\]
The V15 estimate then yields a regulator-uniform exponentially
summable chiral-determinant boundary influence.
\end{theorem}

\begin{proof}
V14 gives uniform weighted bounds for \(P_n\) and its first two
local derivatives.  Proposition \ref{prop:v15v16-costs} and the
assumed bounds on \(b_nY_n\) give the same for
\(\widetilde T_n\) and \(\widetilde{\cal A}_n\).  By (16.8) and
(16.17), \(\widetilde{\cal A}_n\) has a uniform singular floor.
The weighted conjugation perturbation tends to zero with its
exponent, uniformly when the stronger weighted norm is uniform.
Choose \(\beta'>0\) small enough to satisfy the V13 perturbation
condition for every \(n\).  This proves (16.18).  Insert the
resulting constants into Theorem
\ref{thm:v15v15-determinant} and Corollary
\ref{cor:v15v15-influence}.
\end{proof}

\subsection{The unavoidable fork}

The protected source assumes
\[
 s_{\min}(T_n)\ge\tau_*>0.                                 \tag{16.19}
\]
After the exact normalization, this gives only
\[
 \widetilde\tau_n
 =\frac{a_n}{2}s_{\min}(T_n)
 \ge\frac{a_n\tau_*}{2}.                                   \tag{16.20}
\]
The lower bound on the right tends to zero when \(a_n\to0\); it
neither proves nor disproves (16.17).

\begin{proposition}[Normalized hard-gap or soft-vector alternative]
\label{prop:v15v16-fork}
Exactly one of the following holds after passage to a subsequence:
\begin{enumerate}
\item there is \(\widetilde\tau_*>0\) such that
      \(s_{\min}(\widetilde T_n)\ge\widetilde\tau_*\);
\item there are unit vectors
      \(v_n\in\operatorname{Ran}P_n\) such that
      \[
       \|\widetilde T_nv_n\|
       =s_{\min}(\widetilde T_n)\longrightarrow0.           \tag{16.21}
      \]
\end{enumerate}
In the second case, the V13 inverse-locality argument cannot use a
uniform normalized singular floor.  The vectors \(v_n\) are the
required returned soft-mode witnesses.
\end{proposition}

\begin{proof}
Let \(\ell=\liminf_n s_{\min}(\widetilde T_n)\).  If
\(\ell>0\), pass to a tail and take any smaller positive lower
bound.  If \(\ell=0\), choose a subsequence on which the smallest
singular value tends to zero and choose a normalized right singular
vector at each finite cutoff.  Finite-dimensional singular-value
decomposition gives (16.21).
\end{proof}

There is a second, independent obstruction in the unbounded
projective population.  Even if (16.17) holds, a Yukawa map linear
in an unrestricted Higgs coordinate has
\[
 \sup_H\|b_nY_n(H)\|_\beta=\infty                            \tag{16.22}
\]
at every fixed \(a_n>0\).  Thus a worst-case Dobrushin coefficient
cannot be populated from the quartic Higgs tail merely by citing
finite moments.  One must either retain a genuinely compact
conditional field domain, prove a good-set/bad-set mixing theorem,
or change the block metric to one with a proved Lyapunov weight.

\begin{assessment}[Literal source status after normalization]
\label{ass:v15v16-source}
The protected finite shell has a fixed compact Higgs ball, so its
Yukawa weighted norms and local derivatives are finite.  V16 shows
that its overlap prefactor causes no determinant-influence
divergence after exact normalization.  The shell still assumes only
the unnormalized floor (16.19), and the source does not state the
stronger normalized alternative (16.17).

The later RPESM population removes the compact Higgs restriction.
For that population, both (16.17) and a replacement for the
worst-case bound (16.22) remain unproved.  These are separate
obligations: a chiral soft mode is not the same obstruction as a
large Higgs amplitude.
\end{assessment}

\begin{decision}[V17 quartic good-set return]
\label{dec:v15v16-next}
Go upstream to the one-site quartic Higgs confinement proved in
V1--V4.  Replace the impossible global supremum in (16.22) by an
explicit conditional tail estimate outside a radius \(R\), uniform
in the boundary source in the admissible small-coupling regime.
Then formulate a two-level cluster step: use the V15 determinant
influence on the good set and charge the bad set through a
Lyapunov/minorization debit.  If the boundary source destroys a
uniform quartic tail, record the resulting concentration witness
instead of assuming compactness.
\end{decision}

\begin{firewall}[Claim boundary after V16]
\label{fire:v15v16}
The inverse lattice-spacing prefactor has been removed exactly from
all determinant likelihood ratios.  A uniform determinant tail now
follows conditionally from a normalized chiral floor and weighted
bounds on \(a_nY_n\).  Neither the normalized floor nor the required
unbounded-Higgs mixing replacement is supplied by the literal
cofinal source.
\end{firewall}

\section*{Version V16 append-only record}
\addcontentsline{toc}{section}{Version V16 append-only record}
The complete V15 source was retained before this continuation.  It
had \(143354\) bytes, \(3953\) lines, and SHA--256
\[
 \texttt{7AC768EABA53DBE78763E5EF2A32CE8EA1D334DCE9020482928F2DFDC0CEC2FC}.
\]
V16 removed the overlap prefactor exactly at the determinant level,
computed the normalized local derivative costs, and isolated the
normalized hard-gap/soft-vector fork and the independent unbounded-
Higgs obstruction.
\section{Fifteenth-cycle V17: quartic conditional escape and global Lyapunov return}
\label{sec:v15v17}

V16 returned upstream because a Yukawa map on an unbounded Higgs
field cannot have a finite worst-case supremum.  The first possible
repair was a conditional return to a fixed Higgs ball.  That repair
is false: a large boundary tilt moves the entire quartic conditional
law to radius of order \(|t|^{1/3}\).  This continuation proves the
failure and replaces it by a global fourth-moment Lyapunov drift.

The proof is stated for the real \(m\)-component Higgs coordinate.
Let
\[
 \nu_t(\dd h)
 =
 Z_t^{-1}
 \exp\{-\lambda|h|^4+b|h|^2-t\mathbin{\cdot}h\}\,\dd h,
 \qquad h,t\in\mathbb R^m,\quad\lambda>0.                   \tag{17.1}
\]
Write \(b_+=\max\{b,0\}\).

\begin{lemma}[Uniform conditional fourth-moment growth]
\label{lem:v15v17-moment}
There are explicit finite constants \(A_{\lambda,b,m}\) and
\(B_\lambda\) such that
\[
 \boxed{
 \int|h|^4\,\nu_t(\dd h)
 \le A_{\lambda,b,m}
     +B_\lambda |t|^{4/3}}                                 \tag{17.2}
\]
for every \(t\in\mathbb R^m\).  One may take
\[
 A_{\lambda,b,m}
 =\frac{m+b_+^2/\lambda}{2\lambda},
 \qquad
 B_\lambda
 =\frac{3}{2\,4^{4/3}}\lambda^{-4/3}.                       \tag{17.3}
\]
\end{lemma}

\begin{proof}
Integration of the divergence of
\(h\exp\{-\lambda|h|^4+b|h|^2-t\cdot h\}\) gives
\[
 4\lambda\,\mathbb E_t|h|^4
 =
 m+2b\,\mathbb E_t|h|^2-t\cdot\mathbb E_t h.                \tag{17.4}
\]
The boundary term vanishes by quartic decay.  Pointwise,
\[
 2b_+|h|^2
 \le\lambda|h|^4+\frac{b_+^2}{\lambda},                     \tag{17.5}
\]
and maximization over \(r\ge0\) gives
\[
 |t|r
 \le\lambda r^4+
 \frac{3}{4^{4/3}}\lambda^{-1/3}|t|^{4/3}.                  \tag{17.6}
\]
Insert (17.5)--(17.6) in (17.4), move
\(2\lambda\mathbb E_t|h|^4\) to the left, and divide by
\(2\lambda\).
\end{proof}

The exponent \(4/3\) is forced by quartic confinement; it is not an
artifact of the estimate.

\begin{theorem}[No uniform conditional return to an absolute ball]
\label{thm:v15v17-no-ball}
For every fixed \(R<\infty\),
\[
 \boxed{
 \inf_{t\in\mathbb R^m}\nu_t\{|h|\le R\}=0.}                \tag{17.7}
\]
More precisely, if \(t=T e\), \(|e|=1\), and
\(a=(T/(4\lambda))^{1/3}\), then
\[
 \frac{h}{a}\longrightarrow -e
 \quad\hbox{in probability under }\nu_{Te}
 \quad(T\to\infty).                                         \tag{17.8}
\]
\end{theorem}

\begin{proof}
Set \(h=aq\).  Apart from normalization, the rescaled density has
exponent
\[
 -\lambda a^4\bigl(|q|^4+4e\cdot q\bigr)+ba^2|q|^2.         \tag{17.9}
\]
The function
\[
 \Phi_e(q)=|q|^4+4e\cdot q
\]
has the unique minimizer \(-e\): its critical equation is
\(|q|^2q=-e\), and its Hessian at \(-e\) is positive definite.
For every neighborhood \(U\) of \(-e\), coercivity and uniqueness
give
\[
 \inf_{q\notin U}\Phi_e(q)
 >
 \Phi_e(-e).                                                \tag{17.10}
\]
On bounded \(q\)-sets the \(ba^2|q|^2\) term is lower order than
the \(a^4\) gap.  Outside a sufficiently large bounded set, the
quartic term dominates it uniformly.  Comparing the integral on a
small ball about \(-e\) with its complement therefore shows that
the probability of \(q\notin U\) tends to zero.  This proves
(17.8).  Since \(a\to\infty\), any fixed ball
\(\{|h|\le R\}\) corresponds to \(\{|q|\le R/a\}\), which is
eventually disjoint from a fixed neighborhood of \(-e\).  Its
probability tends to zero, proving (17.7).
\end{proof}

Thus neither the tilt-uniform Poincare estimate of V3 nor the
Wasserstein response of V1 implies a uniform absolute-ball
minorization.  Those results control fluctuations and response;
they do not keep the conditional center bounded.

\subsection{A full-system Lyapunov replacement}

Let \(I\) be a finite site set.  At site \(i\), let
\(x_i\in\mathbb R^{m_i}\), \(\lambda_i>0\), and suppose its
heat-bath conditional is (17.1) with
\[
 t_i(x)=\sum_{j\ne i}J_{ij}x_j,                             \tag{17.11}
\]
where \(J_{ij}\) are linear maps.  Let the update rate be \(r_i>0\).
Put
\[
 V(x)=1+\sum_{i\in I}|x_i|^4.                               \tag{17.12}
\]
From Lemma \ref{lem:v15v17-moment}, choose
\[
 A_i=A_{\lambda_i,b_i,m_i},\qquad
 B_i=B_{\lambda_i}.                                         \tag{17.13}
\]
Let \(d_i=\#\{j:J_{ij}\ne0\}\).

\begin{lemma}[Boundary-tilt absorption]
\label{lem:v15v17-absorption}
For every \(\eta>0\),
\[
 |t_i(x)|^{4/3}
 \le
 d_i^{1/3}\sum_{j:J_{ij}\ne0}
 \|J_{ij}\|^{4/3}
 \left(
 \eta |x_j|^4+\frac{2}{3\sqrt{3\eta}}
 \right).                                                   \tag{17.14}
\]
\end{lemma}

\begin{proof}
Convexity gives
\[
 \left(\sum_{j=1}^{d_i}a_j\right)^{4/3}
 \le d_i^{1/3}\sum_{j=1}^{d_i}a_j^{4/3}.
\]
Use \(a_j=\|J_{ij}\||x_j|\).  Finally, for \(u\ge0\),
\[
 u^{4/3}\le\eta u^4+\frac{2}{3\sqrt{3\eta}},
\]
because the maximum of \(z-\eta z^3\) over \(z\ge0\) is
\(2/(3\sqrt{3\eta})\).
\end{proof}

Define
\[
 q_j(\eta)
 =
 r_j-\eta\sum_{i:i\ne j}
 r_iB_i d_i^{1/3}\|J_{ij}\|^{4/3}.                          \tag{17.15}
\]

\begin{theorem}[Quartic heat-bath Foster--Lyapunov drift]
\label{thm:v15v17-drift}
If
\[
 q_*:=\inf_{j\in I}q_j(\eta)>0                              \tag{17.16}
\]
for some \(\eta>0\), then the heat-bath generator of V1 satisfies
\[
 \boxed{
 LV(x)
 \le
 -q_*\sum_{j\in I}|x_j|^4+C_\eta |I|,}                     \tag{17.17}
\]
provided the rates, \(A_i,B_i,d_i\), and the incident
\(\|J_{ij}\|\) sums have uniform bounds.  The constant
\(C_\eta\) depends only on those bounds and \(\eta\), not on
\(|I|\).

For any invariant finite-volume Gibbs law \(\mu_I\),
\[
 \frac1{|I|}\sum_{j\in I}\mathbb E_{\mu_I}|x_j|^4
 \le\frac{C_\eta}{q_*}.                                     \tag{17.18}
\]
In a translation-invariant system this gives
\[
 \mu_I\{|x_j|>R\}
 \le\frac{C_\eta}{q_*R^4}.                                  \tag{17.19}
\]
\end{theorem}

\begin{proof}
A heat-bath update at \(i\) changes \(V\) in conditional
expectation by
\[
 \mathbb E[|x_i'|^4\mid x_{\ne i}]-|x_i|^4
 \le A_i+B_i|t_i(x)|^{4/3}-|x_i|^4.                         \tag{17.20}
\]
Multiply by \(r_i\), sum over \(i\), insert (17.14), and collect the
coefficient of each \(|x_j|^4\).  It is
\(-q_j(\eta)\), proving (17.17); the remaining constants are a
bounded sum per site.

Apply the drift first to the truncation
\(1+\sum_i\min\{|x_i|^4,N\}\), integrate against the invariant law,
and pass to the limit using quartic integrability and monotone
convergence.  Since \(\int LV\,d\mu_I=0\), (17.18) follows.
Equation (17.19) is Markov's inequality.
\end{proof}

The smallness condition (17.16) is mild in a bounded-degree
finite-range system: because the power \(4/3\) is below \(4\),
\(\eta\) may be chosen small enough whenever the displayed incident
sums are finite.  The price is a larger return constant
\(C_\eta\), not loss of the negative quartic drift.

\begin{corollary}[Application to the bosonic Higgs edge term]
\label{cor:v15v17-higgs}
For a compact gauge group acting unitarily on the Higgs fibre, the
edge energy
\[
 \kappa_H\|H_y-\rho_H(U_{xy})H_x\|^2
\]
produces conditional linear maps \(J_{xy}\) whose norms are bounded
uniformly in the link field.  If degrees, quartic coefficients, and
update rates are uniform, the bosonic Higgs heat-bath system
therefore satisfies (17.17) with constants uniform in volume and
in the compact gauge background.
\end{corollary}

\begin{proof}
Unitary gauge transport has operator norm one.  Expanding the
squared edge term gives a quadratic on-site contribution and a
bilinear neighbor term with norm bounded by a constant depending
only on \(\kappa_H\).  The hypotheses of Theorem
\ref{thm:v15v17-drift} then hold uniformly after choosing
\(\eta\) small enough.
\end{proof}

\begin{assessment}[What the Lyapunov return supplies]
\label{ass:v15v17-reading}
V17 proves a volume-uniform fourth-moment return for the bosonic
quartic Higgs heat-bath chain.  It also proves that no argument can
upgrade this to a boundary-uniform conditional return to a fixed
absolute ball.

The estimate is not yet a spectral gap for the determinant-weighted
measure.  The determinant changes a one-site conditional by an
unbounded but finite-rank likelihood ratio.  To preserve
(17.17), its growth must be bounded quantitatively and absorbed by
the quartic term.  A moment estimate under the bosonic measure alone
does not justify that change of measure.
\end{assessment}

\begin{decision}[V18 finite-rank determinant growth]
\label{dec:v15v17-next}
Use the normalized hard operator of V16.  For a local Higgs change,
factor the determinant ratio as a finite-dimensional Fredholm
determinant
\[
 \det(I+\widetilde T^{-1}\Delta\widetilde T)
\]
on the hard source space.  Bound its growth by a fixed power of
\(1+|H_i|\) using the local core rank and normalized inverse bound.
Then test whether the resulting logarithmic conditional correction
preserves the quartic moment estimate and drift.  Keep the
soft-mode branch separate, since the inverse bound fails there.
\end{decision}

\begin{firewall}[Claim boundary after V17]
\label{fire:v15v17}
The bosonic quartic system has a uniform global Lyapunov return, but
its conditional centers escape every fixed ball under large
boundary tilts.  The determinant-weighted drift remains conditional
on a finite-rank likelihood-growth estimate and the normalized hard
gap.
\end{firewall}

\section*{Version V17 append-only record}
\addcontentsline{toc}{section}{Version V17 append-only record}
The complete V16 source was retained before this continuation.  It
had \(155295\) bytes, \(4281\) lines, and SHA--256
\[
 \texttt{AA5C0775777B494EDAB39C602242D55F21F373BF776AAF336738D0556D5F210A}.
\]
V17 disproved uniform conditional return to an absolute Higgs ball,
proved the sharp quartic conditional moment-growth exponent, and
derived a volume-uniform bosonic Foster--Lyapunov drift.
\section{Fifteenth-cycle V18: finite-rank determinant force preserves quartic return}
\label{sec:v15v18}

V17 proved a global quartic Lyapunov drift for the bosonic Higgs
heat-bath chain and isolated the determinant change of measure.
On the normalized hard branch of V16, a local Higgs variation has
uniformly bounded rank.  This gives both a polynomial determinant
ratio and a bounded logarithmic determinant force.  The latter is
exactly what is needed to preserve the fourth-moment drift.

Fix the exterior fields and a local Higgs block
\(h\in\mathbb R^m\).  Let
\[
 \widetilde T(h):\operatorname{Ran}P\longrightarrow{\cal K}
\]
be the normalized chiral map of (16.3).  The Wilson projector \(P\)
is independent of \(h\).  Assume throughout this section that
\[
 s_{\min}(\widetilde T(h))\ge\widetilde\tau>0
 \quad\hbox{for every }h                                   \tag{18.1}
\]
and that the local Yukawa dependence is affine:
\[
 \widetilde T(h)-\widetilde T(k)
 ={\cal Y}(h-k)P.                                           \tag{18.2}
\]
Suppose
\[
 \operatorname{rank}{\cal Y}(u)P\le r_H,\qquad
 \|{\cal Y}(u)P\|\le K_H|u|.                                \tag{18.3}
\]
Set
\[
 w(h)=\det\nolimits_{\operatorname{Ran}P}
       \bigl(\widetilde T(h)^*\widetilde T(h)\bigr)>0.       \tag{18.4}
\]

\begin{lemma}[Finite-rank determinant ratio]
\label{lem:v15v18-ratio}
For all \(h,k\in\mathbb R^m\),
\[
 \boxed{
 \left(1+\frac{K_H|h-k|}{\widetilde\tau}\right)^{-2r_H}
 \le\frac{w(h)}{w(k)}
 \le
 \left(1+\frac{K_H|h-k|}{\widetilde\tau}\right)^{2r_H}.}     \tag{18.5}
\]
\end{lemma}

\begin{proof}
On the moving source space, which is fixed during a Higgs
variation,
\[
 \frac{\det\widetilde T(h)}{\det\widetilde T(k)}
 =
 \det\left(
 I+\widetilde T(k)^{-1}
       [\widetilde T(h)-\widetilde T(k)]
 \right).                                                   \tag{18.6}
\]
For an operator \(K\) of rank at most \(r_H\),
\[
 |\det(I+K)|\le(1+\|K\|)^{r_H}.                              \tag{18.7}
\]
Indeed, factor \(K=AB\) through an \(r_H\)-dimensional space by a
singular-value decomposition, use
\(\det(I+AB)=\det(I+BA)\), and bound the finite determinant by the
product of its singular values.  Equations (18.1)--(18.3) give the
upper bound in (18.5) after taking absolute squares.  Interchange
\(h\) and \(k\) to obtain the lower bound.
\end{proof}

The next estimate is stronger for the moment calculation.  Let
\(\partial_\alpha\) denote the \(m\) coordinate derivatives and put
\[
 K_{\nabla H}
 =
 \left(\sum_{\alpha=1}^m
 \|\partial_\alpha\widetilde T\|^2\right)^{1/2}.             \tag{18.8}
\]

\begin{lemma}[Bounded determinant force]
\label{lem:v15v18-force}
If each \(\partial_\alpha\widetilde T\) has rank at most \(r_H\),
then
\[
 \boxed{
 |\nabla_h\log w(h)|
 \le Q_H:=
 \frac{2r_HK_{\nabla H}}{\widetilde\tau}.}                  \tag{18.9}
\]
The bound is uniform in \(h\) and in the exterior fields whenever
the constants on the right are uniform.
\end{lemma}

\begin{proof}
Jacobi differentiation on the hard source space gives
\[
 \partial_\alpha\log w
 =
 2\operatorname{Re}\operatorname{Tr}
 \left(
 \widetilde T^{-1}\partial_\alpha\widetilde T
 \right).                                                   \tag{18.10}
\]
The trace has rank at most \(r_H\), so
\[
 |\partial_\alpha\log w|
 \le
 2r_H\widetilde\tau^{-1}
 \|\partial_\alpha\widetilde T\|.
\]
Take the Euclidean norm over \(\alpha\).
\end{proof}

Consider the determinant-weighted one-block conditional
\[
 \pi_t(\dd h)
 =
 {\cal Z}_t^{-1}
 e^{-\lambda|h|^4+b|h|^2-t\cdot h}w(h)\,\dd h.              \tag{18.11}
\]
Lemma \ref{lem:v15v18-ratio} ensures only polynomial determinant
growth relative to any reference \(h=k\), so every integration by
parts below is justified by quartic domination.

\begin{theorem}[Determinant-weighted fourth moment]
\label{thm:v15v18-moment}
Under (18.1)--(18.3) and (18.9),
\[
\boxed{
 \mathbb E_{\pi_t}|h|^4
 \le
 \frac{
 m+b_+^2/\lambda+
 \dfrac{3}{4^{4/3}}\lambda^{-1/3}(|t|+Q_H)^{4/3}
 }{2\lambda}.}                                             \tag{18.12}
\]
In particular,
\[
 \mathbb E_{\pi_t}|h|^4
 \le
 \widehat A_{\lambda,b,m,Q_H}
 +\widehat B_\lambda |t|^{4/3},                             \tag{18.13}
\]
where one may take
\[
\begin{aligned}
 \widehat A_{\lambda,b,m,Q_H}
 &=
 \frac{
 m+b_+^2/\lambda+
 \dfrac{3\,2^{1/3}}{4^{4/3}}
 \lambda^{-1/3}Q_H^{4/3}
 }{2\lambda},\\
 \widehat B_\lambda
 &=
 \frac{3\,2^{1/3}}{2\,4^{4/3}}\lambda^{-4/3}.               \tag{18.14}
\end{aligned}
\]
\end{theorem}

\begin{proof}
Integrate the divergence of
\[
 h\,e^{-\lambda|h|^4+b|h|^2-t\cdot h}w(h).
\]
The resulting identity is
\[
 4\lambda\mathbb E_{\pi_t}|h|^4
 =
 m+2b\mathbb E_{\pi_t}|h|^2
 -t\cdot\mathbb E_{\pi_t}h
 +\mathbb E_{\pi_t}
   [h\cdot\nabla\log w(h)].                                 \tag{18.15}
\]
By (18.9), the last two terms are at most
\((|t|+Q_H)\mathbb E_{\pi_t}|h|\).  Apply (17.5) and
(17.6), with \(|t|\) replaced by \(|t|+Q_H\), and absorb
\(2\lambda\mathbb E|h|^4\) into the left side.  This proves
(18.12).  Finally use
\[
 (u+v)^{4/3}\le2^{1/3}(u^{4/3}+v^{4/3})
\]
to obtain (18.13)--(18.14).
\end{proof}

\begin{corollary}[Hard-branch Higgs Lyapunov drift]
\label{cor:v15v18-drift}
Consider the finite Higgs heat-bath system of V17, now with each
conditional multiplied by its normalized chiral determinant.
Assume the local ranks, \(K_{\nabla H}\), and the normalized floor
\(\widetilde\tau\) are uniform, as are the bosonic finite-range
data.  Replace \(A_i,B_i\) in (17.13) by the constants
\(\widehat A_i,\widehat B_i\) of (18.14).  If
\[
 \widehat q_j(\eta)
 =
 r_j-\eta\sum_{i:i\ne j}
 r_i\widehat B_i d_i^{1/3}\|J_{ij}\|^{4/3}
\]
has a positive uniform infimum \(\widehat q_*\), then
\[
 \boxed{
 LV
 \le
 -\widehat q_*\sum_j|x_j|^4+
 \widehat C_\eta|I|.}                                      \tag{18.16}
\]
Consequently the determinant-weighted finite-volume invariant laws
have a uniform fourth moment per site.
\end{corollary}

\begin{proof}
Theorem \ref{thm:v15v18-moment} has exactly the
\(A+B|t|^{4/3}\) form used in (17.20).  Repeat Lemma
\ref{lem:v15v17-absorption} and Theorem
\ref{thm:v15v17-drift} with the hatted constants.
\end{proof}

\begin{proposition}[Local good-set debit]
\label{prop:v15v18-good-set}
Under the hypotheses of Corollary \ref{cor:v15v18-drift}, every
site satisfies
\[
 \mu_I\{|x_i|>R\}
 \le\frac{C_4}{R^4}                                        \tag{18.17}
\]
in a translation-invariant finite volume, with \(C_4\) independent
of \(I\).  On the local good set \(|x_i|\le R\), all factors
\(b_nY_n(x_i)\) in V16 are bounded by a constant growing at most
linearly in \(R\).
\end{proposition}

\begin{proof}
Equation (18.17) is Markov's inequality applied to the per-site
bound from (18.16).  The Yukawa map is affine in the Higgs field,
and its normalized coefficient is \(b_n=a_n/2\), giving the stated
good-set bound.
\end{proof}

The estimate (18.17) is a probability debit, not a worst-case
Dobrushin coefficient.  A further theorem is required to combine it
with the good-set determinant influence of V15.  In particular, a
global Harris minorization on
\(\sum_i|x_i|^4\le C|I|\) would normally deteriorate exponentially
with \(|I|\) and would not yield the required cofinal local gap.

\begin{assessment}[Hard branch versus soft branch]
\label{ass:v15v18-branches}
On the normalized hard branch, the determinant force adds only a
bounded local term to the quartic moment identity, and the global
Higgs Lyapunov return survives.  This repairs the change-of-measure
gap identified in V17.

On the V16 soft branch,
\(\widetilde\tau_n\to0\), both (18.5) and (18.9) lose uniformity.
The normalized singular vectors (16.21) must then be analyzed as
physical soft modes; the hard-branch drift proof cannot be extended
by replacing \(\widetilde\tau_n^{-1}\) with a fixed constant.
\end{assessment}

\begin{assessment}[Literal source applicability]
\label{ass:v15v18-source}
The protected compact shell has finite local ranks and affine local
Yukawa dependence.  If its Wilson and normalized chiral margins are
uniform, V18 rigorously controls the determinant-weighted Higgs
tails without imposing an artificial compact cutoff.

The unbounded RPESM population still lacks the normalized floor
(16.17).  It also lacks a published local weak-Harris or
Lyapunov-weighted influence theorem that combines (18.17) with the
V15 good-set contraction uniformly in volume.  V18 supplies the
inputs for such a theorem but does not assert its conclusion.
\end{assessment}

\begin{decision}[V19 local two-level contraction]
\label{dec:v15v18-next}
Construct a one-sweep coupling for the colored cluster updates.
Inside \(|H|\le R\), use the V11 local gap and the V15 determinant
influence.  Outside, use the fourth-moment return (18.16) with a
concave weighted distance whose derivative decreases at large
amplitude.  Track the bad-set debit locally, avoiding a
volume-sized small set.  If no uniform contraction follows, isolate
the exact inequality between the good-set margin, return rate, and
local tail probability that remains to be populated.
\end{decision}

\begin{firewall}[Claim boundary after V18]
\label{fire:v15v18}
A uniform normalized chiral gap makes the local determinant force
bounded and preserves quartic Higgs Lyapunov return.  The remaining
hard-branch gap is the volume-uniform combination of this return
with good-set cluster contraction; the normalized soft branch
remains separate.
\end{firewall}

\section*{Version V18 append-only record}
\addcontentsline{toc}{section}{Version V18 append-only record}
The complete V17 source was retained before this continuation.  It
had \(165200\) bytes, \(4583\) lines, and SHA--256
\[
 \texttt{7CC1BFDA7CA5EF7475956D773DD391260B8618FA41BAD75283F55DC00C20B1E6}.
\]
V18 proved finite-rank determinant-ratio and determinant-force
bounds and used them to preserve the quartic Higgs Lyapunov drift
on the normalized hard branch.
\section{Fifteenth-cycle V19: determinant-weighted Higgs Poincare and response}
\label{sec:v15v19}

V18 bounded the first Higgs derivative of the hard-branch
determinant.  The same finite-rank calculation bounds its Hessian.
Quartic curvature at infinity then gives a conditional Poincare
constant uniform in every linear boundary tilt.  This restores both
the transport response of V1 and the \(L^2\) block assembly of V6
for the unbounded Higgs sector, conditional on the normalized hard
gap.

Retain the notation of V18 and put
\[
 F(h)=\log w(h).
\]
For \(u\in\mathbb R^m\), write
\[
 \widetilde T_u
 =
 \left.\frac{\dd}{\dd s}\widetilde T(h+su)\right|_{s=0}.
\]
Assume
\[
 \operatorname{rank}\widetilde T_u\le r_H,\qquad
 \|\widetilde T_u\|\le K_H|u|.                              \tag{19.1}
\]

\begin{lemma}[Uniform Hessian of the local log determinant]
\label{lem:v15v19-hessian}
Under the normalized floor (18.1) and affine Higgs dependence,
\[
 \boxed{
 \|\nabla_h^2F(h)\|
 \le Q_{H,2}:=
 \frac{2r_HK_H^2}{\widetilde\tau^2}}                        \tag{19.2}
\]
for every \(h\).
\end{lemma}

\begin{proof}
For directions \(u,v\), affine dependence gives
\(\widetilde T_{uv}=0\).  The second Jacobi identity is therefore
\[
 F_{uv}
 =
 -2\operatorname{Re}\operatorname{Tr}
 \left(
 \widetilde T^{-1}\widetilde T_v
 \widetilde T^{-1}\widetilde T_u
 \right).                                                   \tag{19.3}
\]
The product has rank at most \(r_H\).  Hence
\[
 |F_{uv}|
 \le
 2r_H\widetilde\tau^{-2}K_H^2|u||v|,
\]
which is (19.2).
\end{proof}

We need the finite-dimensional extension of the V3
curvature-at-infinity lemma.

\begin{lemma}[Multidimensional curvature at infinity]
\label{lem:v15v19-curvature}
Let \(U_\theta\in C^2(\mathbb R^m)\) satisfy, uniformly in
\(\theta\),
\[
 \nabla^2U_\theta\ge-KI\quad\hbox{on }\mathbb R^m,\qquad
 \nabla^2U_\theta\ge\mu I\quad\hbox{on }|h|\ge R,            \tag{19.4}
\]
where \(K,\mu,R>0\).  If
\(e^{-U_\theta(h)-t\cdot h}\) is integrable for every \(t\), then
there is a finite \(C_*(K,\mu,R,m)\) such that
\[
 \operatorname{Var}_{\pi_{\theta,t}}f
 \le
 C_*(K,\mu,R,m)
 \int|\nabla f|^2\,\dd\pi_{\theta,t}                        \tag{19.5}
\]
for all \(\theta,t\), where
\(\pi_{\theta,t}\) is the normalized tilted law.
\end{lemma}

\begin{proof}
Choose \(0<\mu_0<\mu/2\).  Construct a smooth radial function
\(\chi(h)=\psi(|h|)\) as follows.  On \(|h|\le R\), take
\(\chi=(K+\mu_0)|h|^2/2\).  In a collar outside \(R\), decrease
\(\psi'\) smoothly to zero over a sufficiently large width, keeping
the negative radial second derivative larger than
\(-(\mu-\mu_0)\); the tangential eigenvalue
\(\psi'(|h|)/|h|\) remains nonnegative.  Keep \(\chi\) constant
outside the collar.  The collar width and
\(\operatorname{osc}\chi\) depend only on
\(K,\mu,R,m\).  Then
\[
 \nabla^2(U_\theta+t\cdot h+\chi)\ge\mu_0I.                 \tag{19.6}
\]
Bakry--Emery gives Poincare constant \(\mu_0^{-1}\) for the law
with the additional factor \(e^{-\chi}\).  The two normalized laws
differ by a density ratio whose oscillation is at most
\(e^{\operatorname{osc}\chi}\).  Holley--Stroock comparison yields
(19.5).
\end{proof}

Consider again
\[
 \pi_t(\dd h)
 \propto
 \exp\{-\lambda|h|^4+b|h|^2-t\cdot h+F(h)\}\,\dd h.          \tag{19.7}
\]
Its untilted potential is
\[
 U(h)=\lambda|h|^4-b|h|^2-F(h).                             \tag{19.8}
\]

\begin{theorem}[Tilt-uniform determinant-weighted Higgs Poincare]
\label{thm:v15v19-poincare}
Under (18.1) and (19.1), there is a finite constant
\[
 C_H=C_H(\lambda,b,m,Q_{H,2})                               \tag{19.9}
\]
such that
\[
 \boxed{
 \sup_{t\in\mathbb R^m}C_P(\pi_t)\le C_H.}                  \tag{19.10}
\]
The constant is also uniform in exterior fields when
\(\widetilde\tau,r_H,K_H\) are uniform.
\end{theorem}

\begin{proof}
The Hessian of the bosonic part is
\[
 \nabla^2(\lambda|h|^4-b|h|^2)
 =
 (4\lambda|h|^2-2b)I+8\lambda hh^*.                        \tag{19.11}
\]
Together with Lemma \ref{lem:v15v19-hessian}, this gives
\[
 \nabla^2U(h)
 \ge
 \bigl(4\lambda|h|^2-2b-Q_{H,2}\bigr)I.                    \tag{19.12}
\]
Thus (19.4) holds with a global lower bound depending only on
\(b_+,Q_{H,2}\) and with any fixed positive exterior curvature
after choosing \(R\) large enough.  Quartic domination and Lemma
\ref{lem:v15v18-ratio} give integrability for every linear tilt.
Apply Lemma \ref{lem:v15v19-curvature}.
\end{proof}

The Poincare estimate also gives a transport response without the
one-dimensional stochastic-order argument used in V1.

\begin{theorem}[Poincare-to-\(W_1\) response]
\label{thm:v15v19-response}
Let \(\{\pi_s:0\le s\le1\}\) have densities proportional to
\(e^{-U_s(h)}\), and suppose \(C_P(\pi_s)\le C_H\).  If
\[
 \sup_h|\nabla_h\dot U_s(h)|\le L_s,                        \tag{19.13}
\]
then
\[
 \boxed{
 W_1(\pi_0,\pi_1)
 \le C_H\int_0^1L_s\,\dd s.}                               \tag{19.14}
\]
\end{theorem}

\begin{proof}
For a smooth one-Lipschitz \(f\),
\[
 \frac{\dd}{\dd s}\int f\,\dd\pi_s
 =
 -\operatorname{Cov}_{\pi_s}(f,\dot U_s).                  \tag{19.15}
\]
Cauchy--Schwarz and the Poincare inequality give
\[
 |\operatorname{Cov}_{\pi_s}(f,\dot U_s)|
 \le
 C_H
 \left(\int|\nabla f|^2\,\dd\pi_s\right)^{1/2}
 \left(\int|\nabla\dot U_s|^2\,\dd\pi_s\right)^{1/2}
 \le C_HL_s.                                                \tag{19.16}
\]
Integrate in \(s\), take the supremum over smooth
one-Lipschitz functions, and use approximation in the
Kantorovich--Rubinstein dual formula.
\end{proof}

\subsection{The hard-branch Higgs influence matrix}

Let \(i,j\) be Higgs blocks.  Write \(J_{ij}\) for the bosonic
bilinear coupling and assume the determinant mixed response obeys
\[
 \|\nabla_{h_i}\nabla_{h_j}F\|
 \le\Xi_{ij}.                                               \tag{19.17}
\]
The V15 two-core theorem gives
\[
 \Xi_{ij}\le C_{\det}e^{-2\beta d(i,j)}                     \tag{19.18}
\]
under its uniform normalized hard-branch inputs.

\begin{corollary}[Uniform Higgs conditional response]
\label{cor:v15v19-influence}
If two exterior configurations differ only at Higgs block \(j\),
then the conditional laws at \(i\) satisfy
\[
 \boxed{
 W_1\bigl(\pi_i(\,\cdot\,|x_j),\pi_i(\,\cdot\,|y_j)\bigr)
 \le
 C_{H,i}
 \bigl(\|J_{ij}\|+\Xi_{ij}\bigr)|x_j-y_j|.}                 \tag{19.19}
\]
\end{corollary}

\begin{proof}
Interpolate linearly from \(x_j\) to \(y_j\).  The derivative of
the conditional potential has \(h_i\)-gradient bounded by
\[
 \bigl(\|J_{ij}\|+\Xi_{ij}\bigr)|x_j-y_j|.
\]
Use Theorem \ref{thm:v15v19-response}.
\end{proof}

Thus the V1 weighted path coupling applies with the determinant-
corrected matrix
\[
 {\cal C}^{\rm hard}_{ij}
 =
 C_{H,i}\bigl(\|J_{ij}\|+\Xi_{ij}\bigr).                   \tag{19.20}
\]
In particular, if
\[
 \kappa_w^{\rm hard}
 =
 \min_j\left[
 r_j-\frac1{w_j}
 \sum_{i\ne j}r_iw_i{\cal C}^{\rm hard}_{ij}
 \right]>0,                                                 \tag{19.21}
\]
the determinant-weighted Higgs heat-bath semigroup contracts
exponentially in the weighted \(W_1\) metric at rate
\(\kappa_w^{\rm hard}\), uniformly in the field amplitudes.

There is also an \(L^2\) conclusion.  Let \(C_i\) be conditional
Poincare matrices for the Higgs blocks and let
\[
 K_{ij}=\|J_{ij}\|+\Xi_{ij}.                                \tag{19.22}
\]
Form the even--odd or colored cross matrix exactly as in V6.

\begin{corollary}[Hard-branch Higgs \(L^2\) gap]
\label{cor:v15v19-gap}
If the V6 normalized cross-Hessian norm satisfies
\[
 q_{\rm hard}<1,                                            \tag{19.23}
\]
then the local Higgs heat-bath generator has
\[
 \boxed{
 \operatorname{gap}A_H
 \ge r_*(1-q_{\rm hard}).}                                 \tag{19.24}
\]
The bound is volume-uniform when the data in
(19.9), (19.18), and the graph geometry are uniform.
\end{corollary}

\begin{proof}
Theorem \ref{thm:v15v19-poincare} supplies the uniform conditional
Poincare matrices.  The total mixed Hessian is bounded by
(19.22).  Apply Theorem \ref{thm:v15v6-local-block-gap}; no new
tensorization argument is needed.
\end{proof}

\begin{assessment}[What has closed on the hard branch]
\label{ass:v15v19-closed}
The unbounded Higgs amplitude is no longer an obstruction to the
Higgs-only hard-branch heat bath.  V18 gives quartic return, while
V19 gives a tilt-uniform conditional Poincare constant, a
Wasserstein response matrix, and a conditional \(L^2\) gap under
the explicit small-influence condition (19.23).

This does not yet give the full gauge--Higgs--fermion cluster gap.
A link derivative contains
\[
 (P_++bY(H))P_s,
\]
whose norm grows with the neighboring Higgs amplitude.  Compact
link conditionals may also have the higher-character multiwell
obstruction of V7--V9.  The Higgs theorem cannot be relabeled as a
uniform gauge-block theorem.
\end{assessment}

\begin{assessment}[Literal source status]
\label{ass:v15v19-source}
For the protected shell, compact Higgs support and assumed Wilson
and chiral floors make all V19 constants finite.  For the unbounded
RPESM sequence, the theorem becomes uniform only if the normalized
floor, determinant tail constant \(C_{\det}\), and the numerical
margin \(1-q_{\rm hard}\) are populated.  The source states none of
these cofinal numerical margins.

The result is nevertheless a complete implication: once those
three inputs are checked, no separate compact truncation or
global minorization is needed for the unbounded Higgs block.
\end{assessment}

\begin{decision}[V20 compulsory audit and gauge--Higgs amplitude metric]
\label{dec:v15v19-next}
At V20, inspect the newest YM and NS companions.  Then address the
remaining mixed link derivative with a state-dependent product
metric or a joint bounded-diameter gauge--Higgs--detail block.
Test whether quartic conditional Poincare control can absorb the
factor \(1+|H|\) in mean square while preserving a volume-uniform
cross norm.  If only a moment-averaged coefficient results, do not
insert it into the worst-case V6 theorem; formulate the exact
weighted alternative required.
\end{decision}

\begin{firewall}[Claim boundary after V19]
\label{fire:v15v19}
On the normalized hard branch, V19 closes the unbounded Higgs
conditional Poincare and response problem under an explicit
determinant influence margin.  Gauge-block amplitudes, retained
compact multiwell details, the normalized soft branch, and the
cofinal numerical constants remain open.
\end{firewall}

\section*{Version V19 append-only record}
\addcontentsline{toc}{section}{Version V19 append-only record}
The complete V18 source was retained before this continuation.  It
had \(175161\) bytes, \(4890\) lines, and SHA--256
\[
 \texttt{CE1DC215150F0B84897EEF06D5A167A14540C3FC42A8CEEF620EB94397F52EAC}.
\]
V19 proved a tilt-uniform determinant-weighted Higgs Poincare
inequality, a general Poincare-to-\(W_1\) response theorem, and
hard-branch Higgs transport and \(L^2\) gap criteria.
\section{Fifteenth-cycle V20: companion audit and the gauge--Higgs source metric}
\label{sec:v15v20}

This is the second compulsory companion checkpoint of the compact
fifteenth cycle.  The companion endpoints have not changed since
V15.  The mathematical continuation identifies the exact metric
behind the remaining gauge--Higgs obstruction.

\subsection{Companion checkpoint}

\begin{record}[Files inspected at V20]
\label{rec:v15v20-companions}
The current Yang--Mills file remains
\begin{center}
\path{Renewal_Yang_Mills_Mass_Gap_Research_Notes_V35.tex},
\end{center}
with \(289778\) bytes, \(7915\) lines, and SHA--256
\[
 \texttt{4EA177BF18DD0ED14474AE54E22DE26FF178C904813F84F53D135BE7DA989C04}.
                                                               \tag{20.1}
\]
The current Navier--Stokes file remains
\begin{center}
\path{Renewal_NS_Stability_Research_Notes_V26.tex},
\end{center}
with \(278283\) bytes, \(5319\) lines, and SHA--256
\[
 \texttt{82C887F8C2C69E4F28FAD7C3DC8DE5B9099CB5A4BF431EBA1FFF3E91935978AD}.
                                                               \tag{20.2}
\]
\end{record}

YM V34--V35 still contributes only the verified distinction between
pointwise chart coverage and an overlap graph, and the requirement
to test a proposed symmetry before transporting a floor.  NS V26
still contributes only the distinction between a restricted
variation norm and the norm over all inputs.  Neither companion has
a new operator acting on the V19 hard-branch Higgs problem.

\begin{audit}[V20 companion verdict]
\label{aud:v15v20-verdict}
No companion theorem or numerical constant is imported.  The V19
Higgs conditional theorem is retained.  The next bottleneck is
analyzed in the actual gauge--Higgs source coordinate instead of
assuming that an ordinary endpoint metric controls it.
\end{audit}

\subsection{The natural source of a compact link}

For \(z\in\mathbb C\), let
\[
 \mu_z(\dd\theta)
 =
 Z_z^{-1}
 \exp\{\operatorname{Re}(ze^{-i\theta})\}\,\dd\theta,
 \qquad\theta\in\mathbb R/(2\pi\mathbb Z).                  \tag{20.3}
\]
This is the translated von Mises family of V7.  Equip the circle
with geodesic distance.

\begin{theorem}[Uniform response in the complex source]
\label{thm:v15v20-source-response}
Let \(\lambda_*>0\) be the uniform Poincare floor of Theorem
\ref{thm:v15v7-von-mises}.  Then
\[
 \boxed{
 W_1(\mu_z,\mu_{z'})
 \le\lambda_*^{-1}|z-z'|}                                  \tag{20.4}
\]
for all \(z,z'\in\mathbb C\).
\end{theorem}

\begin{proof}
Interpolate by \(z_s=(1-s)z+sz'\).  The conditional potential is
\[
 U_s(\theta)=-\operatorname{Re}(z_se^{-i\theta}).
\]
Its parameter derivative obeys
\[
 \left|\partial_\theta\dot U_s(\theta)\right|
 \le|z-z'|.                                                 \tag{20.5}
\]
Every \(\mu_{z_s}\) has Poincare constant at most
\(\lambda_*^{-1}\), including \(z_s=0\).  Apply the
Poincare-to-\(W_1\) response theorem
\ref{thm:v15v19-response} on the circle.
\end{proof}

For scalar endpoints, the gauge--Higgs edge source is the product
\[
 z(x,y)=\gamma x\overline y                                \tag{20.6}
\]
up to a fixed coupling and representation convention.  The same
phenomenon occurs for matrix coefficients in a nonabelian
representation.

\begin{proposition}[No uniform endpoint Lipschitz coefficient]
\label{prop:v15v20-no-endpoint}
Fix any \(\epsilon>0\).  For real \(x,y\) and
\(z(x,y)=xy\), the map
\[
 x\longmapsto\mu_{xy}
\]
has no Lipschitz constant uniform in \(y\) when the input uses
ordinary Euclidean distance.  More precisely, with
\[
 y_R=R,\qquad x_R^\pm=\pm\frac{\epsilon}{R},
\]
one has
\[
 |x_R^+-x_R^-|=\frac{2\epsilon}{R}\longrightarrow0,         \tag{20.7}
\]
while
\[
 W_1(\mu_{x_R^+y_R},\mu_{x_R^-y_R})
 \ge2m(\epsilon)>0,                                        \tag{20.8}
\]
where
\[
 m(\epsilon)=\int\cos\theta\,\mu_\epsilon(\dd\theta)>0.
\]
Hence the quotient of output distance by endpoint distance is at
least \(Rm(\epsilon)/\epsilon\).
\end{proposition}

\begin{proof}
The two link sources are \(\epsilon\) and \(-\epsilon\).
The function \(f(\theta)=\cos\theta\) is one-Lipschitz for geodesic
distance.  Symmetry gives
\[
 \int f\,\dd\mu_{-\epsilon}
 =
 -\int f\,\dd\mu_\epsilon.
\]
Moreover,
\[
 m(\epsilon)
 =
 \frac{\dd}{\dd\epsilon}\log
 \int_0^{2\pi}e^{\epsilon\cos\theta}\,\dd\theta>0
\]
for \(\epsilon>0\).  Kantorovich--Rubinstein duality proves
(20.8), and division by (20.7) gives the final assertion.
\end{proof}

This obstruction is independent of the link Poincare gap.  Both
output laws in (20.8) have the uniform gap of V7.  The failure lies
in the coordinate map from endpoints to the compact-link source.

\begin{corollary}[Correct augmented edge distance]
\label{cor:v15v20-edge-metric}
For endpoint pairs define
\[
 d_\sharp\bigl((x,y),(x',y')\bigr)
 =
 |x-x'|+|y-y'|+
 |x\overline y-x'\overline{y'}|.                            \tag{20.9}
\]
Then
\[
 W_1(\mu_{\gamma x\overline y},
     \mu_{\gamma x'\overline{y'}})
 \le
 \lambda_*^{-1}|\gamma|\,
 d_\sharp\bigl((x,y),(x',y')\bigr).                         \tag{20.10}
\]
The product term in (20.9) is necessary in the sense of
Proposition \ref{prop:v15v20-no-endpoint}.
\end{corollary}

\begin{proof}
Apply Theorem \ref{thm:v15v20-source-response} and use the final
term of (20.9).  The counterexample (20.7)--(20.8) excludes a
uniform replacement by a constant multiple of
\(|x-x'|+|y-y'|\).
\end{proof}

For a unitary representation \(\rho\) and vector Higgs endpoints,
the analogous source is the finite-rank matrix coefficient
\[
 Z(H_x,H_y)=\gamma H_yH_x^*                                \tag{20.11}
\]
restricted to the representation carried by the link.  It obeys
\[
\begin{aligned}
 \|Z(H_x,H_y)-Z(H_x',H_y')\|
 \le|\gamma|\bigl(
 &|H_y|\,|H_x-H_x'|\\
 &+|H_x'|\,|H_y-H_y'|
 \bigr),                                                    \tag{20.12}
\end{aligned}
\]
so the same state-dependent product term is unavoidable.

\begin{proposition}[Moment averaging is not a V6 hypothesis]
\label{prop:v15v20-moment-not-sup}
A fourth-moment bound such as (18.17) does not produce a finite
worst-case endpoint Lipschitz coefficient for (20.6), nor a finite
pointwise cross-Hessian supremum for (20.12).  Therefore it cannot
be inserted directly into the intrinsic cross norm \(q\) of
Theorem \ref{thm:v15v6-local-block-gap}.
\end{proposition}

\begin{proof}
The moment bound is an integral statement under the invariant law;
it does not remove configurations with arbitrarily large \(y\).
Proposition \ref{prop:v15v20-no-endpoint} uses such configurations
and makes the pointwise coefficient arbitrarily large.  The V6
hypothesis is a uniform bound on the conditional cross derivative,
so the quantifiers differ.
\end{proof}

\begin{assessment}[Direction after the V20 audit]
\label{ass:v15v20-direction}
The single-harmonic link sector has a uniform conditional gap and a
uniform response in its natural source.  The ordinary product
metric on unbounded Higgs endpoints is the wrong geometry.  The
minimal repair is to include the edge product source in the
transport distance or to update the responsible endpoints jointly
with the link.

Higher-character and hidden-detail terms remain subject to the
multiwell analysis of V7--V9.  The source metric repairs a response
coefficient; it does not identify or remove distinct wells.
\end{assessment}

\begin{decision}[V21 propagation of the product-source metric]
\label{dec:v15v20-next}
Test the augmented distance (20.9) under a Higgs heat-bath refresh.
Derive a coupling estimate for
\[
 |X'Y-X''Y'|
\]
using the quartic conditional response and the fourth-moment bounds
of V18.  Determine whether a concave truncation of the product term
gives a uniform contraction matrix.  If the estimate necessarily
depends on the untouched endpoint amplitude, enlarge the update to
the full edge block and state the precise bounded-diameter cluster
condition.
\end{decision}

\begin{firewall}[Claim boundary after V20]
\label{fire:v15v20}
V20 proves uniform compact-link response in the natural source and
a sharp obstruction to any uniform ordinary-endpoint coefficient.
It does not yet prove contraction of the augmented metric or the
full gauge--Higgs--detail cluster gap.
\end{firewall}

\section*{Version V20 append-only record}
\addcontentsline{toc}{section}{Version V20 append-only record}
The complete V19 source was retained before this continuation.  It
had \(186132\) bytes, \(5230\) lines, and SHA--256
\[
 \texttt{BBA47883344C3D9B1D8A7C92E853DFDF96B0062F0349970068F6A247CE6317EF}.
\]
V20 completed the compulsory companion audit, proved uniform
single-harmonic link response in its complex source, and proved
that ordinary endpoint distance cannot control that response
uniformly over unbounded Higgs amplitudes.
\section{Fifteenth-cycle V21: failure of the untruncated product metric under a one-site refresh}
\label{sec:v15v21}

V20 found the natural compact-link source metric
\(|xy-x'y'|\).  This continuation tests whether that metric is
stable under a single-endpoint quartic heat-bath update.  It is not.
The obstruction is different from V20: even after the correct source
coordinate is included, refreshing only one endpoint can amplify a
vanishing change in the untouched endpoint into an order-one change
of the refreshed source.

Return to the scalar quartic conditional
\[
 \nu_t(\dd x)
 =
 Z_t^{-1}e^{-\lambda x^4+bx^2-tx}\,\dd x,
 \qquad\lambda>0.                                           \tag{21.1}
\]
Let
\[
 m(t)=\int x\,\nu_t(\dd x),\qquad
 v(t)=\operatorname{Var}_{\nu_t}(x).                        \tag{21.2}
\]

\begin{lemma}[Large-tilt mean and variance]
\label{lem:v15v21-laplace}
As \(t\to+\infty\), with
\[
 a=\left(\frac{t}{4\lambda}\right)^{1/3},
\]
one has
\[
 m(t)=-a+O(a^{-1}),\qquad
 v(t)=\frac{1}{12\lambda a^2}+O(a^{-4}).                    \tag{21.3}
\]
Moreover,
\[
 m'(t)=-v(t).                                               \tag{21.4}
\]
\end{lemma}

\begin{proof}
Set \(x=-a+a^{-1}u\).  Taylor expansion of the exponent about
\(-a\) gives, after removal of its constant value,
\[
 -6\lambda u^2-2bu+
 a^{-2}P_3(u)+a^{-4}P_4(u),                                \tag{21.5}
\]
where the coefficients of the fixed polynomials \(P_3,P_4\) depend
only on \(\lambda,b\) and remain bounded as \(a\to\infty\).
Outside \(|u|\le a^{1/3}\), quartic coercivity gives an integrable
bound \(Ce^{-c u^2}\) after enlarging \(a\); farther out it gives
\(Ce^{-c|u|^4/a^4}\) together with a strictly positive action gap.
Dominated expansion of the first two centered moments therefore
gives
\[
 \mathbb E u=O(1),\qquad
 \operatorname{Var}(u)=\frac{1}{12\lambda}+O(a^{-2}).
\]
Since \(x=-a+a^{-1}u\), this proves (21.3).  Differentiating the
normalized exponential family under the quartically dominated
integral gives
\[
 m'(t)=-\operatorname{Var}_{\nu_t}(x),
\]
which is (21.4).
\end{proof}

Fix \(J>0\).  On \(\mathbb R^2\), let \(K_x\) be the one-site
heat-bath kernel
\[
 K_x((x,y),\dd x'\dd y')
 =
 \nu_{Jy}(\dd x')\,\delta_y(\dd y').                        \tag{21.6}
\]
Thus the second endpoint is held fixed while the first is refreshed
with its bilinear boundary tilt.  Equip the pair with the V20
augmented distance
\[
 d_\sharp((x,y),(\bar x,\bar y))
 =
 |x-\bar x|+|y-\bar y|+|xy-\bar x\bar y|.                   \tag{21.7}
\]

\begin{theorem}[Untruncated source metric is not stable under a
one-endpoint refresh]
\label{thm:v15v21-no-lipschitz}
There is no finite \(L\) such that
\[
 W_{1,d_\sharp}
 \bigl(K_x(z,\cdot),K_x(\bar z,\cdot)\bigr)
 \le Ld_\sharp(z,\bar z)                                    \tag{21.8}
\]
for all \(z,\bar z\in\mathbb R^2\).  More precisely, set
\[
 z_R=(0,R),\qquad
 \bar z_R=(0,R+R^{-1/3}).                                   \tag{21.9}
\]
Then
\[
 d_\sharp(z_R,\bar z_R)=R^{-1/3}\longrightarrow0,           \tag{21.10}
\]
whereas
\[
 \liminf_{R\to\infty}
 W_{1,d_\sharp}
 \bigl(K_x(z_R,\cdot),K_x(\bar z_R,\cdot)\bigr)
 \ge
 \frac43\left(\frac{J}{4\lambda}\right)^{1/3}>0.            \tag{21.11}
\]
\end{theorem}

\begin{proof}
The observable
\[
 f(x,y)=xy
\]
is one-Lipschitz for \(d_\sharp\), because
\(|f(z)-f(\bar z)|\) is the final term of (21.7).  Therefore
Kantorovich--Rubinstein duality gives
\[
 W_{1,d_\sharp}(K_x(z_R,\cdot),K_x(\bar z_R,\cdot))
 \ge|p(R+R^{-1/3})-p(R)|,                                  \tag{21.12}
\]
where
\[
 p(y)=y\,m(Jy).                                             \tag{21.13}
\]
By (21.4),
\[
 p'(y)=m(Jy)-Jy\,v(Jy).                                     \tag{21.14}
\]
If
\[
 a_y=\left(\frac{Jy}{4\lambda}\right)^{1/3},
\]
Lemma \ref{lem:v15v21-laplace} gives
\[
 p'(y)
 =
 -a_y-\frac{Jy}{12\lambda a_y^2}+O(a_y^{-1})
 =
 -\frac43a_y+O(a_y^{-1}).                                  \tag{21.15}
\]
The mean-value theorem on
\([R,R+R^{-1/3}]\) now yields
\[
 |p(R+R^{-1/3})-p(R)|
 \longrightarrow
 \frac43\left(\frac{J}{4\lambda}\right)^{1/3}.              \tag{21.16}
\]
This proves (21.11), and division by (21.10) excludes (21.8).
\end{proof}

The counterexample is already present in the purely bosonic scalar
edge.  A determinant tail or compact-link multiwell is not needed
to produce it.

\begin{corollary}[Any untruncated source-dominating metric fails]
\label{cor:v15v21-source-dominating}
Let \(d\) be a metric satisfying
\[
 d((x,y),(\bar x,\bar y))
 \ge c_0|xy-\bar x\bar y|                                   \tag{21.17}
\]
for some \(c_0>0\), and suppose
\[
 d((0,R),(0,R+\delta))\le C_0|\delta|                       \tag{21.18}
\]
uniformly in \(R\).  Then \(K_x\) has no uniform Lipschitz
coefficient in \(W_{1,d}\).
\end{corollary}

\begin{proof}
Use the test observable \(c_0xy\), which is one-Lipschitz under
(21.17), and repeat (21.12)--(21.16).  The input distance is at
most \(C_0R^{-1/3}\), while the output distance has a positive
lower limit.
\end{proof}

This result specifies what a successful metric must change.  It
must either truncate the edge-source separation at large amplitude,
weight changes of the untouched endpoint by a factor growing at
least like \(R^{1/3}\), or refresh the responsible endpoints
together so that (21.18) is not the relevant comparison.

\begin{proposition}[Minimal update alternatives]
\label{prop:v15v21-alternatives}
An unbounded product-source cost can enter a uniform path-coupling
proof only if at least one of the following is proved:
\begin{enumerate}
\item a state-dependent endpoint cost for which
      \[
       d((0,R),(0,R+\delta))
       \gtrsim R^{1/3}|\delta|                              \tag{21.19}
      \]
      along the V21 sequence, together with contraction of that
      cost under every other update;
\item a block update containing both endpoints of each charged
      product source, with a block Poincare or coupling estimate
      uniform in all exterior tilts;
\item a bounded concave source cost and a weak-Harris estimate that
      combines local contraction with the quartic Lyapunov return.
\end{enumerate}
The V1 ordinary \(W_1\) contraction and V18 fourth moment alone
prove none of these three statements.
\end{proposition}

\begin{proof}
Theorem \ref{thm:v15v21-no-lipschitz} and Corollary
\ref{cor:v15v21-source-dominating} exclude an untruncated
source-dominating metric with ordinary endpoint scaling.  The three
listed changes remove, respectively, the small input in (21.10),
the one-endpoint kernel (21.6), or the unbounded test observable
used in (21.12).  Each introduces a new obligation not present in
V1 or V18.
\end{proof}

\begin{assessment}[Consequence for the SM cluster construction]
\label{ass:v15v21-sm}
The V20 source coordinate is correct for the compact-link
conditional, but its untruncated distance cannot simply be appended
to the V1 Higgs metric.  A single-site colored heat-bath proof in
that metric is impossible even before gauge constraints,
determinants, or retained details are added.

The cluster architecture of V11 remains viable only through the
second or third alternative of Proposition
\ref{prop:v15v21-alternatives}.  For the second, every
gauge--Higgs edge source responsible for an unbounded amplitude must
be internal to the refreshed motif often enough, and the joint
block estimate must be uniform in the remaining boundary tilts.
For the third, a local weak-Harris theorem must retain a
volume-independent contraction rate.
\end{assessment}

\begin{decision}[V22 bounded source cost and local return]
\label{dec:v15v21-next}
Test the bounded concave source cost
\[
 \rho_\ell(r)=1-e^{-r/\ell}.                                \tag{21.20}
\]
First derive a uniform upper response bound for the quartic
one-endpoint kernel in the metric
\[
 |x-\bar x|\wedge1+|y-\bar y|\wedge1+
 \rho_\ell(|xy-\bar x\bar y|).
\]
Use the determinant-weighted fourth moment of V18 to control the
large-amplitude part.  If its contraction coefficient still reaches
one, close the single-endpoint route and formulate the joint
edge-block Poincare problem with all gauge copies quotiented.
\end{decision}

\begin{firewall}[Claim boundary after V21]
\label{fire:v15v21}
The natural untruncated gauge--Higgs source metric fails uniformly
under a one-endpoint quartic refresh.  V21 proves the obstruction
and the minimal alternatives; it does not yet establish a bounded-
cost weak-Harris contraction or a joint edge-block gap.
\end{firewall}

\section*{Version V21 append-only record}
\addcontentsline{toc}{section}{Version V21 append-only record}
The complete V20 source was retained before this continuation.  It
had \(194928\) bytes, \(5490\) lines, and SHA--256
\[
 \texttt{1EB4CA83504374B1B53BB4C13DF3D2925D1C376AE767E5516FE8C192AD010552}.
\]
V21 proved that the untruncated product-source metric is not
uniformly stable under a one-endpoint quartic refresh and isolated
the three remaining admissible repairs.
\section{Fifteenth-cycle V22: bounded source cost restores finite one-endpoint response}
\label{sec:v15v22}

V21 proved that the untruncated product-source distance cannot be
stable under a one-endpoint refresh.  The bounded concave cost
suggested there behaves differently.  At large boundary amplitude,
the refreshed product has a broad distribution; its law changes in
total variation only at rate \(O(|y|^{-1/3})\).  This decay offsets
the escaping conditional center and yields a finite global response
coefficient.

For \(\ell>0\), put
\[
 \rho_\ell(r)=1-e^{-r/\ell},\qquad r\ge0,                   \tag{22.1}
\]
and define
\[
\begin{aligned}
 d_\ell((x,y),(\bar x,\bar y))
 ={}&
 (|x-\bar x|\wedge1)+(|y-\bar y|\wedge1)\\
 &+\rho_\ell(|xy-\bar x\bar y|).                           \tag{22.2}
\end{aligned}
\]
The function \(\rho_\ell\) is increasing, concave,
\(\rho_\ell(0)=0\), and subadditive.  Hence (22.2) is a bounded
metric.

We allow the normalized hard determinant to modify the refreshed
quartic law.  Let \(F\in C^1(\mathbb R)\) satisfy
\[
 |F'(x)|\le Q                                               \tag{22.3}
\]
and
\[
 \pi_t(\dd x)
 =
 {\cal Z}_t^{-1}
 e^{-\lambda x^4+bx^2-tx+F(x)}\,\dd x.                     \tag{22.4}
\]
The V18 finite-rank hard-branch determinant satisfies (22.3) with
\(Q=Q_H\).

For \(J>0\) and \(y\ne0\), let \(\eta_y\) be the law of
\[
 Z_y=yX_y,\qquad X_y\sim\pi_{Jy}.                           \tag{22.5}
\]

\begin{lemma}[Large-source product-law derivative]
\label{lem:v15v22-product-tv}
There are \(C<\infty\) and \(y_0\ge1\), depending only on
\(\lambda,b,J,Q\), such that
\[
 \boxed{
 \|\partial_y\eta_y\|_{\rm TV}
 \le C|y|^{-1/3}}
 \qquad(|y|\ge y_0).                                       \tag{22.6}
\]
Consequently, whenever
\[
 |y-\bar y|<1,\qquad
 \min\{|y|,|\bar y|\}\ge y_0,
\]
one has
\[
 \|\eta_y-\eta_{\bar y}\|_{\rm TV}
 \le
 C|y-\bar y|
 \min\{|y|,|\bar y|\}^{-1/3}.                              \tag{22.7}
\]
\end{lemma}

\begin{proof}
It suffices to take \(y>0\); reflection treats \(y<0\).
The density of \(Z_y\) is
\[
 q_y(z)
 =
 \frac{1}{y{\cal Z}_{Jy}}
 \exp\left\{
 -\lambda\frac{z^4}{y^4}
 +b\frac{z^2}{y^2}
 -Jz+F(z/y)
 \right\}.                                                  \tag{22.8}
\]
Its logarithmic \(y\)-score, written at \(z=yx\), is centered and
equals
\[
 S_y(x)
 =
 \frac1y
 \left[
 4\lambda x^4-2bx^2-xF'(x)
 -\mathbb E_{\pi_{Jy}}
  (4\lambda X^4-2bX^2-XF'(X))
 \right].                                                   \tag{22.9}
\]
The identity follows by differentiating (22.8); centering follows
from differentiation of the normalizing constant.

Put
\[
 a_y=\left(\frac{Jy}{4\lambda}\right)^{1/3},
 \qquad x=-a_y+a_y^{-1}u.                                  \tag{22.10}
\]
The quartic exponent, after subtraction of its value at
\(-a_y\), is
\[
 -6\lambda u^2-2bu+O(a_y^{-1}|u|)
 +O(a_y^{-2}(1+|u|^3))
 +O(a_y^{-4}|u|^4),                                        \tag{22.11}
\]
uniformly for bounded \(u\).  The \(O(a_y^{-1}|u|)\) term uses
(22.3).  Quartic coercivity supplies uniform Gaussian moments for
\(u\) on the central region and an exponentially negligible
complement, exactly as in Lemma
\ref{lem:v15v21-laplace}.

Let
\[
 A(x)=4\lambda x^4-2bx^2-xF'(x).
\]
Although \(F'\) need not be differentiable, (22.3) shows that the
last term in \(A\) has centered fluctuation \(O(a_y)\).
Expansion of the polynomial part under (22.10) gives centered
fluctuation \(O(a_y^2)\).  Hence
\[
 \mathbb E_{\pi_{Jy}}|A(X)-\mathbb EA(X)|
 \le C a_y^2.                                               \tag{22.12}
\]
Equations (22.9)--(22.10) imply
\[
 \mathbb E|S_y|
 \le C\frac{a_y^2}{y}
 \le C y^{-1/3}.                                            \tag{22.13}
\]
Since
\[
 \|\partial_y\eta_y\|_{\rm TV}
 =\frac12\int|S_y|q_y\,\dd z,
\]
this proves (22.6).  Integrating along the interval between
\(y\) and \(\bar y\) proves (22.7).
\end{proof}

\begin{remark}[Regularity of the determinant force]
\label{rem:v15v22-regularity}
For the literal finite-dimensional determinant,
\(F\) is smooth on the hard chart.  The proof used only the bounded
first derivative because the product-law score contains
\(xF'(x)\).  A uniform second derivative, available from V19,
gives the same estimate by direct Taylor expansion.
\end{remark}

We also need a first moment of the product.  The V18 fourth-moment
bound gives
\[
 \mathbb E_{\pi_{Jy}}|X_y|
 \le C(1+|y|^{1/3}),                                       \tag{22.14}
\]
and therefore
\[
 \mathbb E_{\eta_y}|Z_y|
 \le C(|y|+|y|^{4/3}).                                     \tag{22.15}
\]

Define the one-endpoint kernel
\[
 K_x((x,y),\dd x'\dd y')
 =
 \pi_{Jy}(\dd x')\delta_y(\dd y').                          \tag{22.16}
\]

\begin{theorem}[Finite global response in the bounded source metric]
\label{thm:v15v22-finite-response}
There is \(L_\ell<\infty\), depending only on
\(\lambda,b,J,Q,\ell\), such that
\[
 \boxed{
 W_{1,d_\ell}
 \bigl(K_x(z,\cdot),K_x(\bar z,\cdot)\bigr)
 \le L_\ell d_\ell(z,\bar z)}                               \tag{22.17}
\]
for every \(z,\bar z\in\mathbb R^2\).
\end{theorem}

\begin{proof}
The output law depends only on \(y\), so put
\(\delta=|y-\bar y|\).  If \(\delta\ge1\), the right side of
(22.17) is at least \(L_\ell\), while the output distance is at
most \(3\); choose \(L_\ell\ge3\).

Suppose \(\delta<1\).  On a fixed compact \(y\)-interval, the
quartically dominated densities depend smoothly on \(y\).
Couple \(X_y,X_{\bar y}\) by one-dimensional monotone quantiles.
The tilt-uniform response of V19 and uniform compact-tilt moments
give
\[
 \mathbb E(|X_y-X_{\bar y}|\wedge1)
 +
 \mathbb E\rho_\ell(|yX_y-\bar yX_{\bar y}|)
 \le C_\ell\delta.                                         \tag{22.18}
\]
Together with the unchanged-endpoint cost \(\delta\), this proves
the compact-\(y\) case.

It remains to consider \(\min\{|y|,|\bar y|\}\ge y_0\).
The two numbers have the same sign.  Maximally couple the product
laws \(\eta_y,\eta_{\bar y}\).  By (22.7), the mismatch probability
is at most
\[
 C\delta r^{-1/3},\qquad
 r=\min\{|y|,|\bar y|\}.                                    \tag{22.19}
\]
On the matched event set \(Z_y=Z_{\bar y}=Z\).  The source cost is
zero there, and
\[
 \left|\frac Z y-\frac Z{\bar y}\right|
 \le\frac{\delta|Z|}{|y\bar y|}.                            \tag{22.20}
\]
The common part of the two densities has first moment bounded by
either side of (22.15).  Since \(|y|\) and \(|\bar y|\) differ by
less than one, its contribution to the truncated endpoint cost is
at most
\[
 C\delta(r^{-1}+r^{-2/3}).                                  \tag{22.21}
\]
On the mismatch event, both the truncated endpoint cost and the
bounded source cost are at most one.  Adding twice (22.19),
(22.21), and the unchanged-endpoint cost \(\delta\) gives
\[
 W_{1,d_\ell}
 \le C_\ell\delta.                                         \tag{22.22}
\]
Finally, \(d_\ell(z,\bar z)\ge\delta\) when \(\delta<1\).
This proves (22.17).
\end{proof}

The estimate is finite but does not claim contraction.  Indeed, the
second endpoint is unchanged by \(K_x\), so
\[
 W_{1,d_\ell}
 \bigl(K_x((x,y),\cdot),K_x((\bar x,\bar y),\cdot)\bigr)
 \ge |y-\bar y|\wedge1.                                    \tag{22.23}
\]
A strict contraction can arise only after updates that also refresh
the second endpoint.

\begin{corollary}[Determinant hard-branch applicability]
\label{cor:v15v22-hard}
On the normalized hard branch, the scalar gauge--Higgs prototype
satisfies Theorem \ref{thm:v15v22-finite-response} uniformly in
cutoff whenever the quartic coefficients, local determinant rank,
normalized singular floor, and normalized Yukawa derivative are
uniform.  A quasilocal exterior determinant dependence may be
included if its mixed logarithmic response contributes a uniformly
summable addition to the score bound (22.9).
\end{corollary}

\begin{proof}
Lemmas \ref{lem:v15v18-force} and
\ref{lem:v15v19-hessian} populate \(Q\) and its smoothness.
The remaining constants enter only through the V18 moment and V19
tilt-response estimates.  A mixed exterior derivative adds its
centered score to (22.9); V15 controls the sum when its weighted
tail norm is uniform.
\end{proof}

\begin{assessment}[What bounded concavity repairs]
\label{ass:v15v22-reading}
The bounded source cost eliminates the V21 discontinuity of the
untruncated product observable and restores a finite global
one-endpoint response coefficient.  This is enough to make a
path-coupling matrix well-defined.

It is not enough to make that matrix contractive.  The coefficient
\(L_\ell\) is not shown to be below one, and (22.23) gives an exact
unit direction for the untouched endpoint.  The next calculation
must compose the two endpoint updates and the compact-link refresh,
not inspect any one kernel in isolation.
\end{assessment}

\begin{decision}[V23 full edge-sweep coupling matrix]
\label{dec:v15v22-next}
Analyze the ordered edge sweep
\[
 K_{\rm edge}=K_\theta K_yK_x,
\]
where the endpoint kernels use determinant-weighted quartic
conditionals and \(K_\theta\) is the compact-link heat bath.
Compute a three-component response matrix for the two truncated
endpoint costs and the bounded source cost.  A strict spectral
radius below one will give a local Wasserstein contraction; failure
will identify whether the source weight, endpoint coupling, or
multiwell detail forces a joint simultaneous refresh.
\end{decision}

\begin{firewall}[Claim boundary after V22]
\label{fire:v15v22}
Bounded concavity restores a finite uniform response coefficient
for a one-endpoint quartic refresh on the normalized hard branch.
No strict edge-sweep contraction or full gauge--Higgs cluster gap
has yet been proved.
\end{firewall}

\section*{Version V22 append-only record}
\addcontentsline{toc}{section}{Version V22 append-only record}
The complete V21 source was retained before this continuation.  It
had \(203941\) bytes, \(5762\) lines, and SHA--256
\[
 \texttt{7A4873C14ECDDBA469074EA0C3217859D0FFD77E23C2DF667DC425548FCDF5B0}.
\]
V22 proved a large-source total-variation decay for the refreshed
edge product and a finite global response coefficient in the
bounded concave source metric.
\section{Fifteenth-cycle V23: edge-sweep response matrices and the missing mixed sources}
\label{sec:v15v23}

V22 made the one-endpoint response finite after bounding the
gauge--Higgs source cost.  A complete local contraction requires
composition of all updates.  This continuation gives the exact
matrix compiler for that composition.  It also proves that ordinary
link-angle distance leaves an infinite angle-to-endpoint response;
the metric must include every mixed source that appears in an
endpoint conditional.

\subsection{A finite response-matrix compiler}

Let \(d_1,\ldots,d_k\) be bounded component semimetrics on a state
space \(\Omega\), and write
\[
 {\bf d}(z,\bar z)
 =
 (d_1(z,\bar z),\ldots,d_k(z,\bar z))^T.                   \tag{23.1}
\]
For a Markov kernel \(K\), say that a nonnegative matrix \(M_K\)
is a response matrix if, for every \(z,\bar z\), there is a coupling
\((Z',\bar Z')\) of \(K(z,\cdot)\) and \(K(\bar z,\cdot)\) such
that
\[
 \mathbb E{\bf d}(Z',\bar Z')
 \le M_K{\bf d}(z,\bar z)                                  \tag{23.2}
\]
componentwise.

\begin{theorem}[Response matrix for an ordered sweep]
\label{thm:v15v23-matrix}
If \(K_1,\ldots,K_m\) have response matrices
\(M_1,\ldots,M_m\), then the ordered sweep
\[
 K=K_m\cdots K_1
\]
has response matrix
\[
 M=M_m\cdots M_1.                                          \tag{23.3}
\]
There is a positive weight vector \(w\in(0,\infty)^k\) and
\(q<1\) such that the weighted metric
\[
 d_w(z,\bar z)=w^T{\bf d}(z,\bar z)                         \tag{23.4}
\]
obeys
\[
 W_{1,d_w}(K(z,\cdot),K(\bar z,\cdot))
 \le qd_w(z,\bar z)                                        \tag{23.5}
\]
if and only if
\[
 \rho(M)<1.                                                 \tag{23.6}
\]
\end{theorem}

\begin{proof}
Couple the first update using (23.2), condition on its output, and
couple the second update by its response inequality.  Iteration
gives (23.3).

If \(w^TM\le qw^T\) componentwise, multiply (23.2) by \(w^T\)
to obtain (23.5).  Conversely, such an inequality implies
\(\rho(M)\le q\) by the Collatz--Wielandt bound.  If
\(\rho(M)<1\), define
\[
 w^T={\bf 1}^T(I-M)^{-1}
 ={\bf 1}^T\sum_{n\ge0}M^n.
\]
Then \(w>0\) and \(w^TM=w^T-{\bf 1}^T\).  Hence
\[
 q=\max_j\frac{(w^TM)_j}{w_j}<1,
\]
which proves the reverse implication.
\end{proof}

\subsection{The four-channel edge skeleton}

Use four aggregate channels:
\[
 (d_x,d_y,d_s,d_\theta),                                   \tag{23.7}
\]
where \(d_x,d_y\) are truncated endpoint distances,
\(d_s\) is a sum of bounded concave distances of all source
coordinates, and \(d_\theta\) is a residual compact-link/detail
distance.  Suppose the endpoint and link kernels admit the bounds
\[
\begin{array}{c|cccc}
 &d_x&d_y&d_s&d_\theta\\ \hline
 K_x:d_x'&0&a&0&p\\
 K_x:d_y'&0&1&0&0\\
 K_x:d_s'&0&b&0&r\\
 K_x:d_\theta'&0&0&0&1
\end{array}                                                 \tag{23.8}
\]
and, by the corresponding reversed-endpoint calculation,
\[
\begin{array}{c|cccc}
 &d_x&d_y&d_s&d_\theta\\ \hline
 K_y:d_x'&1&0&0&0\\
 K_y:d_y'&a&0&0&p\\
 K_y:d_s'&b&0&0&r\\
 K_y:d_\theta'&0&0&0&1.
\end{array}                                                 \tag{23.9}
\]
Finally suppose the compact-link refresh erases its old channel and
has source response
\[
 d_\theta'\le c\,d_s,                                      \tag{23.10}
\]
while leaving the other three channels unchanged.  The coefficients
are nonnegative.  V22 proves finiteness of \(a,b\) for the scalar
quartic product source; \(c\) is the compact-source response of V20.
The coefficients \(p,r\) require all endpoint--link sources to be
present in \(d_s\), as proved below.

The response matrices are
\[
 M_x=
 \begin{pmatrix}
 0&a&0&p\\
 0&1&0&0\\
 0&b&0&r\\
 0&0&0&1
 \end{pmatrix},
\quad
 M_y=
 \begin{pmatrix}
 1&0&0&0\\
 a&0&0&p\\
 b&0&0&r\\
 0&0&0&1
 \end{pmatrix},
\quad
 M_\theta=
 \begin{pmatrix}
 1&0&0&0\\
 0&1&0&0\\
 0&0&1&0\\
 0&0&c&0
 \end{pmatrix}.                                            \tag{23.11}
\]

\begin{proposition}[Exact edge-sweep feedback block]
\label{prop:v15v23-feedback}
For
\[
 K_{\rm edge}=K_\theta K_yK_x,
\]
the product \(M_{\rm edge}=M_\theta M_yM_x\) has two zero
eigenvalues.  Its remaining eigenvalues are those of
\[
 \boxed{
 B_{\rm edge}
 =
 \begin{pmatrix}
 a^2&p(a+1)\\
 cab&c(bp+r)
 \end{pmatrix}.}                                           \tag{23.12}
\]
Consequently the edge sweep contracts in some positive weighted
sum of the four channels if and only if
\[
 \rho(B_{\rm edge})<1.                                     \tag{23.13}
\]
Equivalently,
\[
\begin{gathered}
 a^2<1,\qquad c(bp+r)<1,                                   \tag{23.14}\\
 (1-a^2)\bigl(1-c(bp+r)\bigr)
 -cpab(a+1)>0.                                             \tag{23.15}
\end{gathered}
\]
\end{proposition}

\begin{proof}
Direct multiplication gives
\[
 M_{\rm edge}
 =
 \begin{pmatrix}
 0&a&0&p\\
 0&a^2&0&p(a+1)\\
 0&ab&0&bp+r\\
 0&cab&0&c(bp+r)
 \end{pmatrix}.                                            \tag{23.16}
\]
The first and third columns vanish.  The nonzero spectrum is the
spectrum of the compression from the input channels
\((d_y,d_\theta)\) to the output channels
\((d_y,d_\theta)\), which is (23.12).  Theorem
\ref{thm:v15v23-matrix} proves (23.13).  For a nonnegative
\(2\times2\) matrix, \(\rho(B)<1\) is equivalent to positivity of
the diagonal entries and determinant of \(I-B\); this is
(23.14)--(23.15).
\end{proof}

The matrix calculation is exact, but its coefficients must be
defined using the complete interaction source list.  A plain angle
distance is insufficient.

\subsection{A second source-coordinate obstruction}

Consider the scalar edge energy
\[
 -Jxy\cos\theta.
\]
At fixed \(y,\theta\), the \(x\)-conditional is
\(\nu_{Jy\cos\theta}\).  Let
\[
 \theta_R=\frac\pi2,\qquad
 \bar\theta_R=\frac\pi2+R^{-1/2},\qquad y_R=R.              \tag{23.17}
\]

\begin{theorem}[No uniform ordinary-angle-to-endpoint response]
\label{thm:v15v23-no-angle}
Let the endpoint output use truncated distance
\(|x-\bar x|\wedge1\), and let the input link use geodesic angle
distance.  Then no finite \(p\) can bound the response of the
\(x\)-refresh uniformly in \(y\).  In fact,
\[
 d_{\mathbb S^1}(\theta_R,\bar\theta_R)=R^{-1/2}\to0,        \tag{23.18}
\]
but
\[
 W_{1,|\cdot|\wedge1}
 \left(
 \nu_{Jy_R\cos\theta_R},
 \nu_{Jy_R\cos\bar\theta_R}
 \right)\longrightarrow1.                                 \tag{23.19}
\]
\end{theorem}

\begin{proof}
The first tilt is zero.  Since
\[
 Jy_R\cos\bar\theta_R
 =
 -J R^{1/2}+O(R^{-1/2}),                                   \tag{23.20}
\]
the second tilt tends to \(-\infty\).  By the large-tilt
concentration theorem \ref{thm:v15v17-no-ball}, the second law
escapes to \(+\infty\), whereas \(\nu_0\) is fixed and tight.
For every \(A>0\), couple the two laws arbitrarily.  With
probability tending to at least \(\nu_0([-A,A])\), the first
variable lies in \([-A,A]\) while the second lies beyond \(A+1\);
the truncated cost is then one.  Let \(R\to\infty\) and then
\(A\to\infty\) to obtain (23.19).
\end{proof}

The missing input coordinate is the endpoint--link source
\[
 s_x(y,\theta)=y\cos\theta.                                \tag{23.21}
\]
Along (23.17),
\[
 |s_x(y_R,\theta_R)-s_x(y_R,\bar\theta_R)|
 \asymp R^{1/2},                                           \tag{23.22}
\]
so its bounded concave cost is order one and correctly records the
large change of conditional law.  Similarly,
\[
 s_y(x,\theta)=x\cos\theta                                 \tag{23.23}
\]
is required for the \(y\)-refresh, while
\[
 s_\theta(x,y)=xy                                          \tag{23.24}
\]
is required for the link refresh.

\begin{corollary}[Complete scalar edge source list]
\label{cor:v15v23-source-list}
For the scalar cosine edge, the aggregate source channel in
(23.7) must contain
\[
\boxed{
 d_s=
 \rho_{\ell_x}(|y\cos\theta-\bar y\cos\bar\theta|)
 +\rho_{\ell_y}(|x\cos\theta-\bar x\cos\bar\theta|)
 +\rho_{\ell_\theta}(|xy-\bar x\bar y|).}                  \tag{23.25}
\]
With only ordinary endpoint and angle distances, at least one of
the feedback coefficients in (23.8)--(23.10) is infinite.
\end{corollary}

\begin{proof}
Theorem \ref{thm:v15v23-no-angle} proves necessity of the first
source, and endpoint reversal proves necessity of the second.
Proposition \ref{prop:v15v20-no-endpoint} proves necessity of the
third.  Each is exactly the coefficient multiplying the variable
updated by the corresponding conditional.
\end{proof}

For a Standard Model representation, (23.21)--(23.24) become the
finite list of matrix coefficients obtained by differentiating each
local edge motif with respect to the variable being updated.  V12's
literal motif ledger is therefore the correct place to enumerate
them.  Higher-character retained details add their own source
coordinates and cannot be hidden in \(d_\theta\).

\begin{assessment}[What V23 decides]
\label{ass:v15v23-reading}
The complete edge sweep has an exact, finite-dimensional contraction
test, (23.13)--(23.15).  The calculation also shows why a matrix
built only from endpoint and link distances would be invalid: its
coefficient \(p\) is infinite already in the scalar cosine model.

No numerical contraction margin is yet claimed.  The complete
bounded source channel (23.25) must first be propagated through
each of the three refreshes.  In the full SM motif, the analogous
matrix coefficients and retained details must be enumerated before
testing (23.15).
\end{assessment}

\begin{decision}[V24 propagate the three scalar sources]
\label{dec:v15v23-next}
For the scalar cosine edge, bound the response of each term in
(23.25) under \(K_x,K_y,K_\theta\).  Use the large-source maximal
coupling of V22 for the endpoint refreshes and the uniform
single-harmonic link response of V20.  Populate finite
\(a,b,p,r,c\), or enlarge the source list if a new conditional
coefficient appears.  Only after the scalar matrix closes should
the nonabelian representation motif be attempted.
\end{decision}

\begin{firewall}[Claim boundary after V23]
\label{fire:v15v23}
V23 proves the response-matrix compiler, its exact edge-sweep
spectral test, and the necessity of three scalar mixed sources.  It
does not yet populate a strict contraction margin or the full
Standard Model source matrix.
\end{firewall}

\section*{Version V23 append-only record}
\addcontentsline{toc}{section}{Version V23 append-only record}
The complete V22 source was retained before this continuation.  It
had \(214154\) bytes, \(6077\) lines, and SHA--256
\[
 \texttt{8755BCCD8B7D9E4AAD8975D340F73D722CF8C02C8C82C743DE037A1626155E4D}.
\]
V23 proved an exact finite response-matrix compiler for an ordered
edge sweep, reduced its contraction test to a \(2\times2\)
feedback block, and proved that all three scalar interaction
sources must enter the bounded metric.
\section{Fifteenth-cycle V24: closure of the scalar bounded-source channel}
\label{sec:v15v24}

V23 identified the three scalar sources
\[
 s_x=y\cos\theta,\qquad
 s_y=x\cos\theta,\qquad
 s_\theta=xy.                                               \tag{24.1}
\]
This continuation proves that their bounded concave costs are closed
under the two endpoint refreshes and the link refresh.  Thus the
edge response matrix has finite entries; no infinite hierarchy of
new scalar source coordinates is generated.

\subsection{Two transformed-law estimates}

Retain the determinant-weighted quartic family
\[
 \pi_t(\dd x)
 \propto e^{-\lambda x^4+bx^2-tx+F(x)}\,\dd x,              \tag{24.2}
\]
where the normalized hard-branch estimates give
\[
 |F'(x)|\le Q,\qquad
 \sup_t C_P(\pi_t)\le C_H.                                  \tag{24.3}
\]
For \(c,t\in\mathbb R\), let
\[
 \zeta_{c,t}=\operatorname{Law}(cX_t),
 \qquad X_t\sim\pi_t.                                       \tag{24.4}
\]

\begin{lemma}[Bounded response of a transformed quartic source]
\label{lem:v15v24-quartic-transform}
Fix \(J_0<\infty\) and \(\ell>0\).  There is
\(C_{\rm q}<\infty\), depending only on
\(\lambda,b,Q,C_H,J_0,\ell\), such that on
\[
 {\cal W}_{J_0}
 =
 \{(c,t):|t|\le J_0(1+|c|)\},                               \tag{24.5}
\]
\[
\boxed{
 W_{1,\rho_\ell(|\cdot-\cdot|)}
 (\zeta_{c,t},\zeta_{\bar c,\bar t})
 \le C_{\rm q}\left[
 (|c-\bar c|\wedge1)+
 \rho_\ell(|t-\bar t|)
 \right].}                                                  \tag{24.6}
\]
The same right side controls the truncated endpoint distance
between \(X_t\) and \(X_{\bar t}\).
\end{lemma}

\begin{proof}
First vary \(t\) at fixed \(c\).  The score of the exponential
family is \(-(X_t-\mathbb EX_t)\), so
\[
 \|\partial_t\pi_t\|_{\rm TV}
 \le\frac12\sqrt{\operatorname{Var}_{\pi_t}X_t}
 \le\frac12\sqrt{C_H}.                                      \tag{24.7}
\]
Maximally couple \(X_t,X_{\bar t}\).  On the matched event both the
endpoint and \(cX\) agree; on the mismatch event their bounded
costs are at most one.  Hence this part is bounded by
\(C\min\{1,|t-\bar t|\}\), which is at most a constant times
\(\rho_\ell(|t-\bar t|)\).

Now vary \(c\) at fixed \(t\).  For \(c\ne0\), the density of
\(Z=cX_t\), evaluated at \(z=cx\), has centered scaling score
\[
 S_{c,t}(x)
 =
 \frac1c\left[
 4\lambda x^4-2bx^2+tx-xF'(x)-1
 \right].                                                   \tag{24.8}
\]
Its expectation is zero by integration of
\(\partial_x(xe^{-\lambda x^4+bx^2-tx+F(x)})\).

On compact \((c,t)\)-sets, quartic domination gives a uniform
total-variation derivative.  If \(|c|\) is large and
\(|t|\le J_0(1+|c|)\), rescale about a global minimizer of the
quartic phase.  Its magnitude is \(O(1+|t|^{1/3})\), its
fluctuation scale is \(O((1+|t|)^{-1/3})\), and the centered
numerator in (24.8) is \(O(1+|t|^{2/3})\) in first absolute
moment.  This is the same Laplace estimate as (22.12), now uniform
for the compact ratio \(t/(1+|c|)\).  Therefore
\[
 \|\partial_c\zeta_{c,t}\|_{\rm TV}
 \le C\left(
 |c|^{-1}+|c|^{-1/3}
 \right)                                                    \tag{24.9}
\]
when \(|c|\) is large.  Integrating (24.9) over a change
\(|c-\bar c|<1\), and using the trivial bound one for larger
changes, gives a constant times \(|c-\bar c|\wedge1\).

Under the maximal coupling of the scaled laws, the source cost
vanishes on the matched part.  The induced endpoint values are
\(z/c,z/\bar c\); their matched truncated distance is bounded by
\[
 \frac{|c-\bar c|}{|c\bar c|}
 \int|z|\,\dd(\zeta_{c,t}\wedge\zeta_{\bar c,t})(z).         \tag{24.10}
\]
The V18 fourth-moment estimate makes (24.10) uniformly bounded by
a constant times \(|c-\bar c|\) on the wedge (24.5).  Compact
values and \(c=0\) are handled directly by coupling \(X_t\) to
itself.  The triangle inequality between the \(t\)- and
\(c\)-variations proves (24.6) and the endpoint assertion.
\end{proof}

For the compact update, let
\[
 \mu_a(\dd\theta)
 =
 Z_a^{-1}e^{a\cos\theta}\,\dd\theta                         \tag{24.11}
\]
for signed \(a\in\mathbb R\), and let
\[
 \xi_{c,a}=\operatorname{Law}(c\cos\Theta_a),
 \qquad\Theta_a\sim\mu_a.                                  \tag{24.12}
\]

\begin{lemma}[Bounded response of a transformed compact source]
\label{lem:v15v24-compact-transform}
For every \(\ell>0\), there is \(C_{\rm c}<\infty\) such that
\[
\boxed{
 W_{1,\rho_\ell(|\cdot-\cdot|)}
 (\xi_{c,a},\xi_{\bar c,\bar a})
 \le C_{\rm c}\left[
 (|c-\bar c|\wedge1)+
 \rho_\ell(|a-\bar a|)
 \right]}                                                   \tag{24.13}
\]
for all \(a,\bar a,c,\bar c\in\mathbb R\).
\end{lemma}

\begin{proof}
At fixed \(c\), the \(a\)-score is
\(\cos\Theta_a-\mathbb E\cos\Theta_a\), whose absolute value is
at most two.  Hence
\[
 \|\mu_a-\mu_{\bar a}\|_{\rm TV}
 \le\min\{1,|a-\bar a|\}.                                   \tag{24.14}
\]
Under a maximal coupling, \(c\cos\Theta\) agrees on the matched
event and has bounded cost on the mismatch event.  Convert the
right side of (24.14) to a constant times
\(\rho_\ell(|a-\bar a|)\).

At fixed \(a\), use the same angle in both variables:
\[
 \rho_\ell(
 |c\cos\Theta_a-\bar c\cos\Theta_a|)
 \le\rho_\ell(|c-\bar c|)
 \le C_\ell(|c-\bar c|\wedge1).                             \tag{24.15}
\]
The triangle inequality proves (24.13).
\end{proof}

\subsection{Closure under the three refreshes}

Choose positive scales \(\ell_x,\ell_y,\ell_\theta\) and define
\[
\begin{aligned}
 d_s(z,\bar z)
 ={}&
 \rho_{\ell_x}(|s_x(z)-s_x(\bar z)|)\\
 &+\rho_{\ell_y}(|s_y(z)-s_y(\bar z)|)\\
 &+\rho_{\ell_\theta}(|s_\theta(z)-s_\theta(\bar z)|).
                                                               \tag{24.16}
\end{aligned}
\]
Let endpoint conditionals be
\[
 X'\sim\pi_{J s_x},\qquad
 Y'\sim\pi_{J s_y},                                        \tag{24.17}
\]
with the other variables held fixed by the corresponding update,
and let the link conditional be
\[
 \Theta'\sim\mu_{J s_\theta}.                              \tag{24.18}
\]

\begin{theorem}[Finite scalar source-channel closure]
\label{thm:v15v24-closure}
For each of \(K_x,K_y,K_\theta\), every output component among
\[
 |x-\bar x|\wedge1,\quad
 |y-\bar y|\wedge1,\quad
 d_s,\quad
 d_{\mathbb S^1}(\theta,\bar\theta)\wedge1                 \tag{24.19}
\]
has a finite response bound by a nonnegative linear combination
of the four input components.  The coefficients depend only on
\[
 \lambda,b,J,Q,C_H,\ell_x,\ell_y,\ell_\theta               \tag{24.20}
\]
and not on \(x,y,\theta\).  In particular, the response matrices
of Theorem \ref{thm:v15v23-matrix} are finite and the source list
(24.1) is closed.
\end{theorem}

\begin{proof}
Consider \(K_x\).  The refreshed endpoint law depends only on
\(s_x=y\cos\theta\).  The truncated endpoint response follows from
(24.7) and the bounded concave cost of the \(s_x\) difference.
The source \(s_x\) itself is unchanged.  The refreshed sources are
\[
 s_y'=X'\cos\theta,\qquad
 s_\theta'=X'y.                                             \tag{24.21}
\]
For the first, apply Lemma
\ref{lem:v15v24-quartic-transform} with
\(c=\cos\theta\), \(t=Js_x\); the coefficient is bounded.
For the second, take \(c=y\), \(t=Js_x\).  Since
\[
 |t|=J|y\cos\theta|\le J|y|,
\]
the wedge condition (24.5) holds.  Differences of \(c\) are
controlled by the truncated \(y\)-component, and differences of
\(t\) by the \(s_x\)-component of (24.16).  The angle itself and
\(y\) are unchanged.  This proves finite response for \(K_x\).
Endpoint reversal proves the same statement for \(K_y\).

For \(K_\theta\), the endpoints and \(s_\theta=xy\) are unchanged.
The link law depends only on \(s_\theta\), and Theorem
\ref{thm:v15v20-source-response}, followed by bounded truncation,
controls its angle response.  The refreshed sources are
\[
 s_x'=y\cos\Theta',\qquad
 s_y'=x\cos\Theta'.                                        \tag{24.22}
\]
Apply Lemma \ref{lem:v15v24-compact-transform} to the first with
\(c=y,a=Js_\theta\), and to the second with
\(c=x,a=Js_\theta\).  Coefficient differences are controlled by
the endpoint components; parameter differences are controlled by
the \(s_\theta\)-component of (24.16).  Thus every output is a
finite linear combination of the four input channels.
\end{proof}

\begin{corollary}[Finite edge-sweep matrix]
\label{cor:v15v24-matrix}
The scalar single-harmonic edge sweep has a finite nonnegative
response matrix
\[
 M_{\rm edge}(\lambda,b,J,Q,C_H,\ell_x,\ell_y,\ell_\theta).
                                                               \tag{24.23}
\]
If a direct evaluation or upper bounding matrix
\(\widehat M_{\rm edge}\ge M_{\rm edge}\) satisfies
\[
 \rho(\widehat M_{\rm edge})<1,                             \tag{24.24}
\]
then the sweep contracts a positive weighted bounded metric
uniformly over all scalar endpoint amplitudes.
\end{corollary}

\begin{proof}
Theorem \ref{thm:v15v24-closure} supplies finite response matrices
for the three kernels.  Multiply them and apply Theorem
\ref{thm:v15v23-matrix}.  Monotonicity of the spectral radius for
nonnegative matrices permits the upper matrix
\(\widehat M_{\rm edge}\).
\end{proof}

The closure theorem concerns one first-harmonic scalar motif.  In a
nonabelian Standard Model edge, the local source list consists of
the finitely many representation matrix coefficients appearing in
the Higgs edge action, the normalized chiral determinant
derivatives, and every retained detail character.  Finite
dimensionality of a representation makes that list finite, but
multiwell characters still require the joint-detail treatment of
V9--V11.

\begin{assessment}[What remains numerical and structural]
\label{ass:v15v24-reading}
V24 closes the scalar source algebra and proves that every entry of
the bounded response matrix is finite.  This removes the two
infinite-coefficient obstructions of V20--V23.

Finiteness does not imply (24.24).  The estimates in the proof are
not optimized, and the literal source supplies no numerical
small-coupling margin for the scalar prototype, still less for the
full representation motif.  A multiwell retained detail can also
create an additional feedback cycle absent from (24.1).
\end{assessment}

\begin{decision}[V25 compulsory audit and contraction margin]
\label{dec:v15v24-next}
At V25, inspect the newest YM and NS companions.  Then choose scales
\(\ell_x,\ell_y,\ell_\theta\), write an explicit upper response
matrix for the scalar first-harmonic edge, and test a transparent
small-coupling regime in which its spectral radius is below one.
If the bounded-source estimates are too coarse to produce a margin,
return the exact failed row and replace it by a simultaneous edge
refresh rather than declaring contraction from finiteness.
\end{decision}

\begin{firewall}[Claim boundary after V24]
\label{fire:v15v24}
The three scalar bounded source coordinates are closed under the
full edge sweep and yield a finite response matrix.  No strict
spectral-radius margin, nonabelian motif matrix, retained-detail
cycle, or full Standard Model cluster gap has yet been proved.
\end{firewall}

\section*{Version V24 append-only record}
\addcontentsline{toc}{section}{Version V24 append-only record}
The complete V23 source was retained before this continuation.  It
had \(225043\) bytes, \(6422\) lines, and SHA--256
\[
 \texttt{13DF8FBF2140D8A6DB25546DEC7A1D8F10018259610ECF1E876C58948E0B4686}.
\]
V24 proved transformed quartic and compact-source response lemmas
and used them to close the three-source scalar edge channel with a
finite response matrix.
\section{Fifteenth-cycle V25: companion audit and saturation of worst-case bounded response}
\label{sec:v15v25}

This compulsory checkpoint first reads the current companion
endpoints.  It then tests the proposed small-coupling contraction of
the V24 bounded response matrix.  The test fails for a precise
reason: bounded source and endpoint costs both saturate along
opposite infinite tilts, for every nonzero coupling.  Small coupling
alone cannot make the global worst-case coefficient strictly less
than one.

\subsection{Companion checkpoint}

\begin{record}[Files inspected at V25]
\label{rec:v15v25-companions}
The current Yang--Mills endpoint is
\begin{center}
\path{Renewal_Yang_Mills_Mass_Gap_Research_Notes_V36.tex},
\end{center}
with \(304825\) bytes, \(8318\) lines, and SHA--256
\[
 \texttt{5FDD7DC94BFA4A7654858355A537EBBBB99C03E4AB633910212B932BFD6BE537}.
                                                               \tag{25.1}
\]
The Navier--Stokes endpoint remains
\begin{center}
\path{Renewal_NS_Stability_Research_Notes_V26.tex},
\end{center}
with \(278283\) bytes, \(5319\) lines, and SHA--256
\[
 \texttt{82C887F8C2C69E4F28FAD7C3DC8DE5B9099CB5A4BF431EBA1FFF3E91935978AD}.
                                                               \tag{25.2}
\]
\end{record}

YM V36 closes a proposed coordinate-symmetry shortcut in two ways:
a cubic trace invariant separates the remaining pencils, and an
exact re-tree calculation shows that a moving frame preserves the
Hessian while its naively differentiated response congruence fails.
The transferable identity is the product rule
\[
 \partial(T^*HT)
 =
 T^*(\partial H)T+
 (\partial T^*)HT+T^*H(\partial T).                         \tag{25.3}
\]
Thus a nonabelian SM source atlas may transport response constants
through a moving gauge frame only after the two connection terms
are included.  V15--V16 avoided this issue for the chiral
determinant by the frame-free identity completion; no YM numerical
floor is imported.

NS V26 is unchanged.  Its restricted rank-one kernel norms still
do not act on the gauge--Higgs response matrix.

\begin{audit}[V25 companion verdict]
\label{aud:v15v25-verdict}
Import only the moving-frame rule (25.3).  Do not import any
companion constant or infer a nonabelian source symmetry from
similar local floors.  The scalar edge test proceeds independently.
\end{audit}

\subsection{The saturation test}

Let \(\pi_t\) be the determinant-weighted quartic conditional
(22.4), with bounded determinant force.  Equip its state line with
\[
 d_{\rm e}(x,\bar x)=|x-\bar x|\wedge1,                     \tag{25.4}
\]
and the tilt-source line with
\[
 d_{\rm s}(s,\bar s)=\rho_\ell(|s-\bar s|).                 \tag{25.5}
\]
For a fixed coupling \(J\ne0\), define the conditional kernel
\[
 {\cal K}_J(s,\dd x)=\pi_{Js}(\dd x).                       \tag{25.6}
\]

\begin{theorem}[Worst-case source response saturates]
\label{thm:v15v25-saturation}
For every \(J\ne0\),
\[
 \boxed{
 \sup_{s\ne\bar s}
 \frac{
 W_{1,d_{\rm e}}(\pi_{Js},\pi_{J\bar s})
 }{
 \rho_\ell(|s-\bar s|)
 }
 \ge1.}                                                     \tag{25.7}
\]
More precisely,
\[
 \lim_{S\to\infty}
 W_{1,d_{\rm e}}(\pi_{JS},\pi_{-JS})=1,\qquad
 \lim_{S\to\infty}\rho_\ell(2S)=1.                         \tag{25.8}
\]
Hence the global coefficient in (25.7) does not tend to zero as
\(J\to0\) through nonzero values.
\end{theorem}

\begin{proof}
The bounded determinant force changes the quartic critical equation
by at most a constant.  The Laplace argument of Theorem
\ref{thm:v15v17-no-ball} therefore still gives
\[
 X_t\longrightarrow-\operatorname{sgn}(t)\infty
 \quad\hbox{in probability as }|t|\to\infty.                \tag{25.9}
\]
Fix \(A>0\).  Under any coupling of
\(X_{JS}\) and \(X_{-JS}\), the first variable is eventually below
\(-A-1\) with probability tending to one and the second is above
\(A+1\) with probability tending to one.  Thus their truncated
distance tends to one in probability and in expectation.  Taking
the infimum over couplings proves the first limit in (25.8).  The
second follows directly from (22.1).  Their quotient tends to one,
proving (25.7).
\end{proof}

\begin{corollary}[No global small-coupling conclusion from V24]
\label{cor:v15v25-no-small-j}
The finite response coefficients proved in V24 cannot all be
replaced by bounds that vanish with \(|J|\) uniformly over the
unbounded source space.  In particular, a proof of
\(\rho(M_{\rm edge})<1\) cannot follow solely by declaring \(J\)
small and estimating every row by a worst-case bounded-source
Lipschitz coefficient.
\end{corollary}

\begin{proof}
The endpoint conditional is one of the rows of any complete edge
response system.  Its source-to-endpoint coefficient is bounded
below by (25.7) for every \(J\ne0\).  Therefore that row has no
uniform \(O(|J|)\) estimate.
\end{proof}

This obstruction does not say that the Gibbs sampler fails to mix.
It says that a globally bounded, worst-case one-step metric discards
the amplitude decrease created by quartic confinement.  An
extreme source is sent to an endpoint of order \(|Js|^{1/3}\), a
large reduction in the Lyapunov scale, while both bounded distances
remain equal to one.

\begin{proposition}[The missing information is Lyapunov decrease]
\label{prop:v15v25-lyapunov}
Let \(V(x)=1+|x|^4\).  Under \({\cal K}_J(s,\cdot)\),
\[
 \mathbb E V(X)
 \le C_0+C_1|Js|^{4/3}.                                    \tag{25.10}
\]
Thus the conditional amplitude grows sublinearly relative to a
quartic source scale, even along the saturation sequence in
Theorem \ref{thm:v15v25-saturation}.  No bounded source metric can
record this decrease without an explicit Lyapunov factor.
\end{proposition}

\begin{proof}
Equation (25.10) is Theorem
\ref{thm:v15v18-moment}.  Along \(s=S\), its right side is
\(O(S^{4/3})\), whereas a quartic amplitude assigned to the source
endpoint is \(O(S^4)\).  The costs (25.4)--(25.5) truncate both
scales at one and erase that distinction.
\end{proof}

\begin{assessment}[Direction after V25]
\label{ass:v15v25-direction}
The V24 source list and finite response matrix remain correct, but a
strict global bounded-metric margin is the wrong target.  The
quartic drift of V18 must enter the distance.  On a bounded
Lyapunov sublevel, small coupling can produce a strict local
response margin because the source range is compact.  Outside that
sublevel, the drift must pay for the missing contraction.

The YM audit adds a separate constraint for the later nonabelian
step: any moving-frame source coordinate must include connection
terms before its local response is entered in the matrix.
\end{assessment}

\begin{decision}[V26 Lyapunov-weighted contraction compiler]
\label{dec:v15v25-next}
Prove a two-region coupling theorem for
\[
 D_\beta(z,\bar z)
 =
 \sqrt{
 d(z,\bar z)
 [1+\beta V(z)+\beta V(\bar z)]
 }.
\]
Combine a finite global response coefficient, strict contraction on
a compact Lyapunov sublevel, and a drift coefficient strong enough
that their product is below one.  Then verify the scalar quartic
edge hypotheses: V24 supplies global finiteness, V18 supplies drift,
and compactness plus small coupling supplies the local margin.
\end{decision}

\begin{firewall}[Claim boundary after V25]
\label{fire:v15v25}
Worst-case bounded source response saturates for every nonzero
coupling, so V24 alone cannot prove a strict global edge contraction.
The programme remains viable through a Lyapunov-weighted distance;
that compiler and its scalar application are not yet proved here.
\end{firewall}

\section*{Version V25 append-only record}
\addcontentsline{toc}{section}{Version V25 append-only record}
The complete V24 source was retained before this continuation.  It
had \(236645\) bytes, \(6742\) lines, and SHA--256
\[
 \texttt{85A1C89550BB48426E8D97E55EEAE84F773A8A70B0715C29A74F16A0D782984F}.
\]
V25 completed the compulsory companion audit and proved that
worst-case bounded source response saturates, forcing the quartic
Lyapunov decrease into the contraction metric.
\section{Fifteenth-cycle V26: a Lyapunov-weighted contraction compiler}
\label{sec:v15v26}

V25 showed that no globally bounded worst-case source metric can
record the amplitude decrease produced by quartic confinement.  The
repair is to multiply the bounded separation by a Lyapunov weight.
This continuation proves the required two-region coupling theorem
and applies it to the scalar determinant-weighted edge sweep.

\subsection{The abstract two-region theorem}

Let \(P\) be a Markov kernel on \(\Omega\), let
\(d:\Omega^2\to[0,1]\) be a bounded separation cost, and let
\(V:\Omega\to[0,\infty)\).  Assume the drift
\[
 PV(z)\le\lambda V(z)+B,\qquad 0<\lambda<1,\quad B<\infty.   \tag{26.1}
\]
Suppose there are \(L<\infty\), \(q<1\), and \(R>0\) such that
for every \(z,\bar z\) one can couple
\((Z',\bar Z')\) with the correct \(P\)-marginals and
\[
 \mathbb E d(Z',\bar Z')
 \le
 \begin{cases}
 q\,d(z,\bar z),
   &V(z)+V(\bar z)\le R,\\
 L\,d(z,\bar z),
   &V(z)+V(\bar z)>R.
 \end{cases}                                                \tag{26.2}
\]
Define
\[
 D_\beta(z,\bar z)
 =
 \sqrt{
 d(z,\bar z)
 [1+\beta V(z)+\beta V(\bar z)]
 }.                                                         \tag{26.3}
\]

\begin{theorem}[Lyapunov-weighted two-region contraction]
\label{thm:v15v26-compiler}
Assume
\[
 L\lambda<1.                                                \tag{26.4}
\]
Choose \(\beta>0\) so that
\[
 q(1+2\beta B)<1.                                          \tag{26.5}
\]
Then choose \(R\) sufficiently large that
\[
 L\,
 \frac{1+2\beta B+\beta\lambda R}
      {1+\beta R}<1.                                       \tag{26.6}
\]
Under (26.1)--(26.2), there is \(\kappa<1\) such that
\[
 \boxed{
 W_{1,D_\beta}(P(z,\cdot),P(\bar z,\cdot))
 \le\kappa D_\beta(z,\bar z)}                              \tag{26.7}
\]
for all \(z,\bar z\).
\end{theorem}

\begin{proof}
For the coupling in (26.2), Cauchy--Schwarz gives
\[
\begin{aligned}
 \mathbb E D_\beta(Z',\bar Z')
 &\le
 \left(\mathbb E d(Z',\bar Z')\right)^{1/2}\\
 &\quad\times
 \left(
 1+\beta\mathbb EV(Z')+\beta\mathbb EV(\bar Z')
 \right)^{1/2}.                                            \tag{26.8}
\end{aligned}
\]
Put \(S=V(z)+V(\bar z)\).  By (26.1), the second squared
factor is at most
\[
 1+2\beta B+\beta\lambda S.                                \tag{26.9}
\]
After division by \(D_\beta(z,\bar z)\), the squared contraction
ratio is bounded by
\[
 q\,g_\beta(S)\quad(S\le R),\qquad
 L\,g_\beta(S)\quad(S>R),                                  \tag{26.10}
\]
where
\[
 g_\beta(S)
 =
 \frac{1+2\beta B+\beta\lambda S}{1+\beta S}.               \tag{26.11}
\]
Because \(\lambda<1\),
\[
 g_\beta'(S)
 =
 \frac{\beta(\lambda-1-2\beta B)}
      {(1+\beta S)^2}<0.                                   \tag{26.12}
\]
Thus on the good region,
\[
 qg_\beta(S)\le qg_\beta(0)=q(1+2\beta B)<1.                \tag{26.13}
\]
On the bad region,
\[
 Lg_\beta(S)\le Lg_\beta(R)<1                              \tag{26.14}
\]
by (26.6).  Such an \(R\) exists because
\(g_\beta(R)\to\lambda\) and (26.4) holds.  Take
\(\kappa^2\) to be the larger right side of
(26.13)--(26.14).
\end{proof}

\begin{remark}[Role of the weighted cost]
\label{rem:v15v26-role}
The cost \(D_\beta\) need only be a lower semicontinuous transport
cost for (26.7); the proof does not require its triangle inequality.
It dominates \(\sqrt d\) and records both bounded separation and
the amount of Lyapunov mass carried by that separation.
\end{remark}

\subsection{Scalar quartic edge drift}

Consider the scalar edge state \(z=(x,y,\theta)\) with
\[
 V(z)=1+|x|^4+|y|^4.                                       \tag{26.15}
\]
Let \(P_J=K_\theta K_yK_x\) be the ordered heat-bath sweep of V23,
with determinant-weighted quartic endpoint conditionals satisfying
the V18 hard-branch bounds and a single-harmonic compact-link
conditional.  Assume all cross-source parameters are multiplied by
a scalar coupling \(J\).

\begin{lemma}[Arbitrarily strong full-sweep quartic drift]
\label{lem:v15v26-edge-drift}
Fix \(J_{\max}<\infty\).  For every \(0<\lambda_0<1\), there is
\(B_{\lambda_0,J_{\max}}<\infty\) such that
\[
 \boxed{
 P_JV(z)
 \le\lambda_0V(z)+B_{\lambda_0,J_{\max}}}                  \tag{26.16}
\]
for all \(|J|\le J_{\max}\) and all \(z\).
\end{lemma}

\begin{proof}
The \(x\)-refresh erases the old \(x\).  The V18 fourth-moment
bound and \(|\cos\theta|\le1\) give
\[
 \mathbb E(|X'|^4\mid y,\theta)
 \le A+C|J|^{4/3}|y|^{4/3}.                                \tag{26.17}
\]
For every \(\epsilon>0\), Young's inequality gives
\[
 |y|^{4/3}\le\epsilon|y|^4+C_\epsilon.                     \tag{26.18}
\]
Hence
\[
 \mathbb E|X'|^4
 \le C_{\epsilon,J_{\max}}+\epsilon C|y|^4.                \tag{26.19}
\]

The \(y\)-refresh erases the old \(y\) and obeys
\[
 \mathbb E(|Y'|^4\mid X',\theta)
 \le A+C|J|^{4/3}|X'|^{4/3}.                               \tag{26.20}
\]
By Jensen and (26.19),
\[
 \mathbb E|X'|^{4/3}
 \le(\mathbb E|X'|^4)^{1/3}
 \le C_{\epsilon,J_{\max}}
      +C\epsilon^{1/3}|y|^{4/3}.                           \tag{26.21}
\]
Apply (26.18) once more, with an independently chosen Young
parameter, to make the coefficient of \(|y|^4\) in
\(\mathbb E|Y'|^4\) arbitrarily small.  The link refresh does not
change \(V\).  Thus
\[
 P_JV(z)\le C+\delta|y|^4
 \le\lambda_0V(z)+B_{\lambda_0,J_{\max}}
\]
after choosing the two absorption parameters so that
\(\delta\le\lambda_0\).  The old \(|x|^4\) has zero coefficient,
which proves (26.16).
\end{proof}

Let \(d\) be the sum of the two truncated endpoint costs, the
truncated link cost, and the three bounded source costs (24.16),
divided by its maximum so that \(0\le d\le1\).

\begin{lemma}[Global finite response and compact local contraction]
\label{lem:v15v26-two-regions}
For every \(J_{\max}<\infty\), there is \(L_*<\infty\) such that
\[
 W_{1,d}(P_J(z,\cdot),P_J(\bar z,\cdot))
 \le L_*d(z,\bar z)                                        \tag{26.22}
\]
for all \(|J|\le J_{\max}\).  Suppose additionally that every
cross dependence, including any determinant mixed response, vanishes
at \(J=0\).  Then for every \(R<\infty\) and every prescribed
\(q_0\in(0,1)\), there is \(J_{R,q_0}>0\) such that
\[
 W_{1,d}(P_J(z,\cdot),P_J(\bar z,\cdot))
 \le q_0d(z,\bar z)                                        \tag{26.23}
\]
whenever
\[
 V(z)+V(\bar z)\le R,\qquad |J|\le J_{R,q_0}.
\]
\end{lemma}

\begin{proof}
Theorem \ref{thm:v15v24-closure} gives finite response matrices
whose constants are uniform on the compact coupling interval
\([-J_{\max},J_{\max}]\).  Their product gives (26.22).

On the sublevel \(V(z)+V(\bar z)\le R\), both endpoints and every
source coefficient range over a compact set.  At \(J=0\), all three
refresh laws are independent of the incoming state; couple the
same draws in the two chains, giving response zero.  Differentiation
of the normalized finite-dimensional conditional densities gives
response coefficients continuous in \(J\), uniformly on the
compact sublevel.  Hence their product matrix has norm tending to
zero as \(J\to0\).  Choose \(J_{R,q_0}\) so that a weighted matrix norm is
less than the prescribed \(q_0\), and apply Theorem
\ref{thm:v15v23-matrix}.
\end{proof}

\begin{theorem}[Scalar hard-branch edge contraction]
\label{thm:v15v26-scalar}
Under the normalized hard-branch hypotheses of V18--V19 and the
cross-scaling hypothesis of Lemma
\ref{lem:v15v26-two-regions}, there is \(J_0>0\) such that for
every \(|J|\le J_0\), one can choose \(\beta>0\), \(R<\infty\),
and \(\kappa<1\) with
\[
 \boxed{
 W_{1,D_\beta}(P_J(z,\cdot),P_J(\bar z,\cdot))
 \le\kappa D_\beta(z,\bar z)}                              \tag{26.24}
\]
uniformly over all scalar endpoint amplitudes.
\end{theorem}

\begin{proof}
Choose a finite \(J_{\max}\) and the global coefficient \(L_*\)
from (26.22).  By Lemma \ref{lem:v15v26-edge-drift}, choose
\(\lambda_0>0\) so small that
\[
 L_*\lambda_0<1.                                           \tag{26.25}
\]
The same lemma supplies the corresponding \(B\).  Fix, for
example, \(q_0=1/2\).  Choose \(\beta>0\) so that
\[
 q_0(1+2\beta B)<1,
\]
and then choose \(R\) large enough to satisfy (26.6).  Lemma
\ref{lem:v15v26-two-regions} supplies a positive
\(J_{R,q_0}\) for this fixed compact sublevel.  Take
\(J_0=\min\{J_{\max},J_{R,q_0}\}\) and apply Theorem
\ref{thm:v15v26-compiler}.
\end{proof}

The selection is noncircular:
\(L_*,\lambda_0,B,q_0,\beta,R\) are fixed in that order, and
compact continuity then supplies the
strictly positive radius \(J_{R,q_0}\).

\begin{assessment}[What V26 proves]
\label{ass:v15v26-reading}
The scalar first-harmonic, normalized-hard-branch edge sweep has a
genuine uniform Wasserstein contraction in a Lyapunov-weighted cost
for a nonempty small-coupling interval.  This resolves the
worst-case saturation of V25 without imposing a compact Higgs
cutoff.

The theorem does not yet give an \(L^2\) spectral gap for the full
lattice.  It also assumes that every cross dependence vanishes with
the scalar coupling.  In the literal Standard Model, gauge,
Yukawa, determinant, plaquette, and retained-detail couplings have
different fixed strengths; their compact-sublevel response matrix
must be populated rather than collapsed into one \(J\).
\end{assessment}

\begin{decision}[V27 lattice assembly of weighted edge costs]
\label{dec:v15v26-next}
Place one copy of the bounded source cost and quartic Lyapunov
weight on every Higgs edge.  Prove a finite-overlap comparison
between the sum of edge costs and a global lattice cost.  Compose
the colored edge-cluster sweeps, tracking the global response and
drift constants without a volume factor.  Then isolate the
finite-dimensional compact-sublevel matrix whose spectral radius
must be checked for the literal SM couplings.
\end{decision}

\begin{firewall}[Claim boundary after V26]
\label{fire:v15v26}
V26 proves a Lyapunov-weighted contraction compiler and a
small-coupling scalar hard-branch edge application.  It does not
yet assemble overlapping lattice edges, prove a volume-uniform
\(L^2\) gap, or populate the fixed Standard Model coupling matrix.
\end{firewall}

\section*{Version V26 append-only record}
\addcontentsline{toc}{section}{Version V26 append-only record}
The complete V25 source was retained before this continuation.  It
had \(244698\) bytes, \(6956\) lines, and SHA--256
\[
 \texttt{6D5B76B3820086DF91BE0D0CE273D18FEFDDA9CEFB42BD714B0F4ED23CE82128}.
\]
V26 proved the Lyapunov-weighted two-region contraction theorem,
an arbitrarily strong scalar edge-sweep quartic drift, and a genuine
small-coupling hard-branch contraction despite global bounded-cost
saturation.
\section{Fifteenth-cycle V27: finite-overlap lattice assembly of local Lyapunov costs}
\label{sec:v15v27}

V26 proved contraction for one scalar edge in a Lyapunov-weighted
cost.  A direct sum with one global Lyapunov function would attach
an additive drift constant proportional to volume.  The correct
lattice object keeps a local Lyapunov weight on each edge
neighborhood.  Finite overlap and a sparse response matrix then
make every assembly constant volume-independent.

Let \(G=(\mathcal V,\mathcal E)\) be a finite graph with
\[
 1\le\deg v\le\Delta.                                      \tag{27.1}
\]
For a field pair \(z,\bar z\), let
\(\delta_v\in[0,1]\) be the truncated Higgs separation at \(v\).
For an edge \(e=\{x,y\}\), let \(\delta_e^U\in[0,1]\) be its
compact-link separation and let \(\delta_{e,a}^s\in[0,1]\) be the
bounded costs of the finitely many local source coordinates.  Put
\[
 d_e
 =
 \delta_x+\delta_y+\delta_e^U+\sum_{a=1}^{N_s}\delta_{e,a}^s.
                                                               \tag{27.2}
\]
A harmless fixed normalization may be used to make \(d_e\le1\).

\begin{proposition}[Finite-overlap comparison]
\label{prop:v15v27-overlap}
Define
\[
\begin{aligned}
 d_{\rm base}
 &=
 \sum_{v\in\mathcal V}\delta_v+
 \sum_{e\in\mathcal E}
 \left(\delta_e^U+\sum_a\delta_{e,a}^s\right),\\
 d_{\rm edge}
 &=\sum_{e\in\mathcal E}d_e.
\end{aligned}                                               \tag{27.3}
\]
Then
\[
 \boxed{
 d_{\rm base}\le d_{\rm edge}\le\Delta d_{\rm base}.}       \tag{27.4}
\]
For
\[
 V_e(z)=1+|H_x|^4+|H_y|^4,                                 \tag{27.5}
\]
one has
\[
 \sum_{e\in\mathcal E}[V_e(z)-1]
 =
 \sum_{v\in\mathcal V}\deg(v)|H_v|^4,                       \tag{27.6}
\]
and hence
\[
 \sum_v|H_v|^4
 \le\sum_e[V_e-1]
 \le\Delta\sum_v|H_v|^4.                                   \tag{27.7}
\]
\end{proposition}

\begin{proof}
Every link and source term occurs once in both costs.  The vertex
term \(\delta_v\) occurs once in \(d_{\rm base}\) and
\(\deg(v)\) times in \(d_{\rm edge}\), proving (27.4).
Equation (27.6) is the same incidence count with \(|H_v|^4\);
(27.1) gives (27.7).
\end{proof}

For \(\beta>0\), define the local weighted edge cost
\[
 D_{e,\beta}(z,\bar z)
 =
 \sqrt{
 d_e(z,\bar z)
 [1+\beta V_e(z)+\beta V_e(\bar z)]
 }.                                                         \tag{27.8}
\]
Given positive weights \(w_e\), the global cost is
\[
 {\cal D}_{\beta,w}(z,\bar z)
 =
 \sum_{e\in\mathcal E}w_eD_{e,\beta}(z,\bar z).             \tag{27.9}
\]
Each additive Lyapunov constant in (27.8) is local.  No factor
\(|\mathcal V|\) appears inside a square root.

\subsection{Coloring the overlapping edge blocks}

Let \(L(G)\) be the line graph of \(G\).  Its maximum degree is at
most \(2\Delta-2\).

\begin{lemma}[Uniform edge coloring]
\label{lem:v15v27-color}
The edge set admits a proper coloring
\[
 \mathcal E=\mathcal E_1\sqcup\cdots\sqcup\mathcal E_m,
 \qquad m\le2\Delta-1,                                     \tag{27.10}
\]
so that the endpoint pairs of distinct edges in one color are
disjoint.
\end{lemma}

\begin{proof}
Greedy vertex coloring of \(L(G)\) uses at most
\(\deg_{\max}L(G)+1\le2\Delta-1\) colors.  Vertices of \(L(G)\)
are edges of \(G\), and adjacency means sharing an endpoint.
\end{proof}

For a color \(c\), let \(P_c\) update in parallel all edge clusters
centered at \(e\in\mathcal E_c\).  Disjointness makes their random
draws conditionally independent given the exterior.  Let
\[
 P_{\rm sweep}=P_m\cdots P_1.                              \tag{27.11}
\]

Index the vector of local costs by edges:
\[
 {\bf D}_\beta(z,\bar z)
 =
 (D_{e,\beta}(z,\bar z))_{e\in\mathcal E}.                 \tag{27.12}
\]
Suppose a coupled color update satisfies
\[
 \mathbb E{\bf D}_\beta(Z',\bar Z')
 \le M_c{\bf D}_\beta(z,\bar z)                            \tag{27.13}
\]
for a nonnegative matrix \(M_c\).

\begin{theorem}[Colored lattice response compiler]
\label{thm:v15v27-compiler}
The full sweep has response matrix
\[
 M_{\rm sweep}=M_m\cdots M_1.                              \tag{27.14}
\]
If positive weights obey
\[
 w^TM_{\rm sweep}\le q_{\rm lat}w^T,\qquad q_{\rm lat}<1,  \tag{27.15}
\]
then
\[
 \boxed{
 W_{1,{\cal D}_{\beta,w}}
 (P_{\rm sweep}(z,\cdot),P_{\rm sweep}(\bar z,\cdot))
 \le q_{\rm lat}{\cal D}_{\beta,w}(z,\bar z).}             \tag{27.16}
\]
Equivalently, such weights exist at fixed volume if and only if
\[
 \rho(M_{\rm sweep})<1.                                    \tag{27.17}
\]
If
\[
 \sup_{f\in\mathcal E}\sum_{e\in\mathcal E}
 (M_{\rm sweep})_{ef}\le q_{\rm lat}<1,                    \tag{27.18}
\]
then one may take \(w_e=1\), and the bound is immediately
volume-uniform.
\end{theorem}

\begin{proof}
Iterated conditional coupling gives (27.14), exactly as in Theorem
\ref{thm:v15v23-matrix}.  Multiply (27.13) by \(w^T\), iterate the
colors, and use (27.15) to obtain (27.16).  The
Perron--Frobenius equivalence follows from the resolvent construction
in the proof of that theorem.  Finally, (27.18) is precisely
\({\bf 1}^TM_{\rm sweep}\le q_{\rm lat}{\bf 1}^T\).
\end{proof}

\subsection{Why locality removes volume}

Let \(d_L\) denote line-graph distance.  Suppose
\[
 (M_c)_{ef}=0\quad\hbox{if }d_L(e,f)>R_0,                  \tag{27.19}
\]
and, for \(0\le r\le R_0\),
\[
 (M_c)_{ef}\le a_r
 \quad\hbox{when }d_L(e,f)=r.                              \tag{27.20}
\]
The number of edges at line distance \(r\) from a fixed edge is at
most
\[
 N_\Delta(r)
 =
 \begin{cases}
 1,&r=0,\\
 (2\Delta-2)(2\Delta-3)^{r-1},&r\ge1,
 \end{cases}                                                \tag{27.21}
\]
after replacing the right side by a larger constant when
\(\Delta=1\).

\begin{proposition}[Sparse response bound]
\label{prop:v15v27-sparse}
Under (27.19)--(27.20),
\[
 \sup_f\sum_e(M_c)_{ef}
 \le
 A_{\Delta,R_0}:=
 \sum_{r=0}^{R_0}N_\Delta(r)a_r.                           \tag{27.22}
\]
For the \(m\)-color sweep,
\[
 \sup_f\sum_e(M_{\rm sweep})_{ef}
 \le\prod_{c=1}^m A_{\Delta,R_0}^{(c)}.                    \tag{27.23}
\]
Both bounds are independent of
\(|\mathcal V|\) and \(|\mathcal E|\).
\end{proposition}

\begin{proof}
For a fixed input edge \(f\), group the nonzero entries of its
matrix column by line distance and use (27.21), proving (27.22).
The maximum column-sum norm is submultiplicative, which gives
(27.23).
\end{proof}

The product bound (27.23) can be very coarse; the exact sparse
product in (27.14) is the correct object for a numerical margin.
The point of Proposition \ref{prop:v15v27-sparse} is that neither
calculation grows merely because the finite lattice grows.

\subsection{Local drift without an extensive additive debit}

A radius-\(R_V\) edge neighborhood contains at most
\(C_{\Delta,R_V}\) vertices.  Define
\[
 V_e^{(R_V)}
 =
 1+\sum_{v:d(v,e)\le R_V}|H_v|^4.                          \tag{27.24}
\]
Suppose an update centered at \(f\) obeys, for every affected
\(e\),
\[
 P_fV_e^{(R_V)}
 \le
 \lambda_{ef}V_f^{(R_V)}+B_{ef},                           \tag{27.25}
\]
where \(\lambda_{ef}=B_{ef}=0\) beyond a fixed line-graph range and
the nonzero constants are uniform.  Applying the V26
Cauchy--Schwarz argument to each \(D_{e,\beta}\) produces a local
response coefficient from \(f\) to \(e\).  At most
\(C_{\Delta,R_V}\) such additive constants meet any fixed edge.

\begin{corollary}[No volume factor in the Lyapunov assembly]
\label{cor:v15v27-no-volume}
Under finite range, bounded degree, and uniform local estimates
(27.25), every column sum of the Lyapunov-weighted response matrix
is bounded by a constant depending only on
\[
 \Delta,R_0,R_V,\beta,\{\lambda_{ef}\},\{B_{ef}\},          \tag{27.26}
\]
not on the lattice volume.
\end{corollary}

\begin{proof}
For each pair \((e,f)\), repeat (26.8)--(26.14) with the local
Lyapunov functions in (27.24).  Finite range makes the coefficient
zero outside a fixed line-graph ball.  Sum the remaining uniformly
bounded entries and use (27.21).
\end{proof}

\begin{theorem}[Conditional volume-uniform scalar lattice
contraction]
\label{thm:v15v27-scalar-lattice}
Consider a bounded-degree scalar first-harmonic lattice whose
edge-cluster conditionals satisfy the normalized hard-branch
hypotheses of V26 uniformly.  Suppose:
\begin{enumerate}
\item all source and determinant tails entering an edge update have
      a common finite response radius or a summable line-graph
      response profile;
\item the local quartic drift estimates (27.25) hold uniformly;
\item the exact colored-sweep response matrix satisfies either
      (27.15) with uniformly comparable weights or the column
      margin (27.18).
\end{enumerate}
Then the colored edge-cluster sweep contracts
\({\cal D}_{\beta,w}\) at a rate independent of finite lattice
volume.
\end{theorem}

\begin{proof}
The local V26 construction gives each response entry.  Corollary
\ref{cor:v15v27-no-volume} and the assumed summable tails make all
columns uniformly finite.  Theorem \ref{thm:v15v27-compiler} and
the stated margin give the contraction.
\end{proof}

\begin{assessment}[The remaining lattice datum]
\label{ass:v15v27-reading}
V27 proves that overlapping edge costs, local quartic returns, and
colored cluster updates can be assembled without a volume factor.
The remaining head-sector datum is finite-dimensional: the exact
compact-sublevel response matrix, including incidence
multiplicities, must have a uniform Perron--Frobenius margin.

The theorem is conditional because V26 populated only a scalar
small-coupling edge.  The literal Standard Model has fixed
representation couplings, plaquette motifs, normalized determinant
tails, and retained details.  Their response entries have not been
enumerated, and a local edge contraction alone does not imply
(27.18).
\end{assessment}

\begin{decision}[V28 head margin under a summable tail]
\label{dec:v15v27-next}
Split the lattice response matrix into a finite-range head and the
quasilocal determinant/shell tail of V12--V15.  Prove a perturbation
theorem: a head column-sum or weighted Perron--Frobenius margin
survives when the weighted tail norm is smaller than that margin.
Express the exact debit in the same edge weights used in
\({\cal D}_{\beta,w}\), so no incompatible norm is silently
substituted.
\end{decision}

\begin{firewall}[Claim boundary after V27]
\label{fire:v15v27}
V27 proves a volume-uniform lattice assembly compiler for local
Lyapunov-weighted edge costs.  It does not populate the literal SM
head response matrix, its strict margin, or the full determinant and
shell-tail debit.
\end{firewall}

\section*{Version V27 append-only record}
\addcontentsline{toc}{section}{Version V27 append-only record}
The complete V26 source was retained before this continuation.  It
had \(255353\) bytes, \(7273\) lines, and SHA--256
\[
 \texttt{ADCCFCDC2064A60D5ADEAE7412E56DD44D6F5F306624A3E60DA253545924A6EF}.
\]
V27 proved finite-overlap comparison, uniform edge coloring, sparse
colored response assembly, and local Lyapunov bookkeeping without a
volume-dependent additive debit.
\section{Fifteenth-cycle V28: head-margin stability under determinant and shell tails}
\label{sec:v15v28}

V27 reduced lattice contraction to a weighted response-matrix
margin.  V12--V15 control quasilocal tails first through likelihood
oscillation and total variation.  Those quantities cannot be
inserted unchanged into the Lyapunov-weighted transport matrix.
A maximal coupling followed by Cauchy--Schwarz takes a square root
of the total-variation coefficient.  The determinant response in
V15 decays with twice the inverse distance, so the resulting
transport tail remains exponentially summable.

\subsection{The common weighted matrix norm}

For positive edge weights \(w=(w_e)\), define the nonnegative-matrix
norm
\[
 \|A\|_w
 =
 \sup_f\frac1{w_f}\sum_e w_eA_{ef}.                         \tag{28.1}
\]

\begin{lemma}[Weighted column norm]
\label{lem:v15v28-norm}
For nonnegative matrices \(A,B\),
\[
 \|AB\|_w\le\|A\|_w\|B\|_w.                                \tag{28.2}
\]
Moreover,
\[
 w^TA\le q\,w^T
 \quad\Longleftrightarrow\quad
 \|A\|_w\le q.                                              \tag{28.3}
\]
\end{lemma}

\begin{proof}
For every input column \(f\),
\[
\begin{aligned}
 \sum_e w_e(AB)_{ef}
 &=
 \sum_g\left(\sum_e w_eA_{eg}\right)B_{gf}\\
 &\le
 \|A\|_w\sum_g w_gB_{gf}
 \le
 \|A\|_w\|B\|_w w_f.
\end{aligned}
\]
This proves (28.2).  Equation (28.3) is the definition written
componentwise.
\end{proof}

Let \(H_c\) be a finite-range head response matrix for color \(c\)
and \(T_c\) a nonnegative tail debit, so that the full response is
bounded by
\[
 M_c\le H_c+T_c.                                           \tag{28.4}
\]
Put
\[
 h_c=\|H_c\|_w,\qquad\eta_c=\|T_c\|_w,                      \tag{28.5}
\]
and
\[
 H_{\rm sweep}=H_m\cdots H_1.                              \tag{28.6}
\]

\begin{theorem}[Tail perturbation of a colored head margin]
\label{thm:v15v28-perturb}
Suppose
\[
 \|H_{\rm sweep}\|_w\le q_{\rm head}<1.                    \tag{28.7}
\]
Define
\[
 {\cal E}_{\rm tail}
 =
 \sum_{j=1}^m
 \eta_j
 \left[\prod_{k=j+1}^m(h_k+\eta_k)\right]
 \left[\prod_{k=1}^{j-1}h_k\right].                        \tag{28.8}
\]
Then
\[
 \boxed{
 \|M_m\cdots M_1\|_w
 \le q_{\rm head}+{\cal E}_{\rm tail}.}                    \tag{28.9}
\]
In particular, the full sweep contracts the V27 global cost if
\[
 q_{\rm head}+{\cal E}_{\rm tail}<1.                       \tag{28.10}
\]
If the tail is added only after a completed head sweep,
\(M_{\rm sweep}\le H_{\rm sweep}+T\), the debit reduces to
\[
 \|M_{\rm sweep}\|_w\le q_{\rm head}+\|T\|_w.               \tag{28.11}
\]
\end{theorem}

\begin{proof}
Use the telescoping identity
\[
\begin{aligned}
 &(H_m+T_m)\cdots(H_1+T_1)-H_m\cdots H_1\\
 &\quad=
 \sum_{j=1}^m
 (H_m+T_m)\cdots(H_{j+1}+T_{j+1})
 T_jH_{j-1}\cdots H_1.                                    \tag{28.12}
\end{aligned}
\]
Apply Lemma \ref{lem:v15v28-norm}, the triangle inequality, and
(28.5), which gives (28.8)--(28.9).  Equations
(28.10)--(28.11) follow from Theorem
\ref{thm:v15v27-compiler}.
\end{proof}

\subsection{From likelihood influence to a local weighted cost}

Let \(K_e^0\) be a head conditional update and \(K_e\) the same
update with a quasilocal tail.  Suppose changing exterior edge block
\(f\) while keeping all other exterior variables fixed gives a
maximal-coupling mismatch probability
\[
 p_{ef}(z,\bar z)
 \le a_{ef}\,d_f(z,\bar z),                                \tag{28.13}
\]
where \(0\le d_f\le1\) is the bounded local separation in (27.8).
Assume the coupled output has the local moment bound
\[
 1+\beta\mathbb E[
 V_e(Z')+V_e(\bar Z')]
 \le
 C_{ef}
 [1+\beta V_f(z)+\beta V_f(\bar z)].                       \tag{28.14}
\]

\begin{lemma}[Square-root transport debit]
\label{lem:v15v28-square-root}
The tail contribution to the expected local cost satisfies
\[
 \boxed{
 \mathbb E[
 D_{e,\beta}(Z',\bar Z')\,{\bf 1}_{\rm mismatch}]
 \le
 \sqrt{a_{ef}C_{ef}}\,
 D_{f,\beta}(z,\bar z).}                                   \tag{28.15}
\]
For several exterior blocks, the corresponding upper debits add:
\[
 (T_{\rm tr})_{ef}
 =
 \sqrt{a_{ef}C_{ef}}.                                      \tag{28.16}
\]
\end{lemma}

\begin{proof}
Since \(d_e\le1\),
\[
 D_{e,\beta}^2
 \le1+\beta V_e(Z')+\beta V_e(\bar Z').
\]
Cauchy--Schwarz, (28.13), and (28.14) give
\[
\begin{aligned}
 \mathbb E[D_{e,\beta}{\bf 1}_{\rm mismatch}]
 &\le
 p_{ef}^{1/2}
 \left(\mathbb ED_{e,\beta}^2\right)^{1/2}\\
 &\le
 \sqrt{a_{ef}d_fC_{ef}
 [1+\beta V_f(z)+\beta V_f(\bar z)]},
\end{aligned}
\]
which is (28.15).  Couple successive exterior perturbations one at
a time and use the triangle inequality for expected costs; this
sums (28.15) and gives (28.16).
\end{proof}

The square root is essential.  A total-variation mismatch may place
the two outputs at large Higgs amplitudes, and the local cost carries
the square root of their Lyapunov mass.  V18 supplies (28.14) on the
normalized hard branch.

\subsection{Exponential determinant response remains summable}

Suppose the V15 determinant Hessian and compact block paths give
\[
 a_{ef}^{\det}
 \le A_{\det}e^{-2\alpha d_L(e,f)}                          \tag{28.17}
\]
and the moment-transfer factors obey
\[
 C_{ef}\le C_{\rm mom}e^{2\gamma d_L(e,f)}                 \tag{28.18}
\]
for some \(0\le\gamma<\alpha\).  The finite-range case has
\(\gamma=0\).  Then
\[
 (T_{\det})_{ef}
 \le
 \sqrt{A_{\det}C_{\rm mom}}\,
 e^{-(\alpha-\gamma)d_L(e,f)}.                             \tag{28.19}
\]

\begin{theorem}[Weighted determinant transport norm]
\label{thm:v15v28-det-tail}
Assume the line graph has polynomial sphere growth
\[
 \#\{e:r\le d_L(e,f)<r+1\}
 \le C_{\rm gr}(1+r)^{d_0-1}.                              \tag{28.20}
\]
If the edge weights satisfy
\[
 \frac{w_e}{w_f}\le C_we^{\omega d_L(e,f)}                 \tag{28.21}
\]
for some
\[
 0\le\omega<\alpha-\gamma,                                 \tag{28.22}
\]
then
\[
\boxed{
 \|T_{\det}\|_w
 \le
 C_w\sqrt{A_{\det}C_{\rm mom}}\,
 C_{\rm gr}
 \sum_{r\ge0}(1+r)^{d_0-1}
 e^{-(\alpha-\gamma-\omega)r}<\infty.}                     \tag{28.23}
\]
The bound is independent of finite lattice volume.
\end{theorem}

\begin{proof}
Insert (28.19) and (28.21) in (28.1), group the sum by
line-graph shells, and apply (28.20).  The remaining
polynomially weighted exponential series converges by (28.22).
\end{proof}

Thus the twice-distance decay obtained from the two-core determinant
trace in V15 is strong enough to survive the square-root conversion
to Lyapunov transport.

For a shell or control tail, define directly
\[
 a_{ef}^{\rm sh}
 \le A_{\rm sh}(d_L(e,f)).                                 \tag{28.24}
\]
The compatible transport debit is
\[
 \eta_{\rm sh}^{\rm tr}
 =
 \sup_f\frac1{w_f}\sum_e
 w_e\sqrt{
 A_{\rm sh}(d_L(e,f))C_{ef}}.                              \tag{28.25}
\]
A likelihood tail summable only before the square root need not
make (28.25) finite.  This is a stronger and now explicit source
obligation.

\begin{corollary}[Head plus determinant plus shell criterion]
\label{cor:v15v28-full}
Suppose the finite-range head sweep has weighted margin
\(1-q_{\rm head}>0\).  Let
\(\eta_{\det}^{\rm tr}\) be bounded by (28.23), and let
\(\eta_{\rm sh}^{\rm tr}\) include the shell, control, and retained-
detail tail debits in the same norm.  If the tails are added after
the head and
\[
 \boxed{
 q_{\rm head}
 +\eta_{\det}^{\rm tr}
 +\eta_{\rm sh}^{\rm tr}<1,}                               \tag{28.26}
\]
then the V27 lattice Lyapunov cost contracts uniformly in volume.
For tails inserted color by color, replace their sum by the exact
telescoping debit (28.8).
\end{corollary}

\begin{proof}
Add the nonnegative tail matrices and apply Theorem
\ref{thm:v15v28-perturb}, followed by Theorem
\ref{thm:v15v27-compiler}.
\end{proof}

\begin{assessment}[Literal source ledger after V28]
\label{ass:v15v28-source}
On the normalized hard branch, V13--V19 conditionally populate
\(A_{\det},\alpha\), and local moment factors.  V28 converts those
quantities into the exact transport norm required by the lattice
assembly.

The literal source still does not give:
\begin{enumerate}
\item the finite-range SM head margin \(1-q_{\rm head}\);
\item a cofinal normalized chiral floor, hence uniform
      \(A_{\det}\) and \(\alpha\);
\item a shell/control likelihood profile strong enough that the
      square-root norm (28.25) is finite;
\item the retained-detail response entries.
\end{enumerate}
The earlier claim of a merely summable likelihood influence is
therefore insufficient for the Lyapunov-weighted programme unless
its decay survives (28.25).
\end{assessment}

\begin{decision}[V29 normalized soft branch and tail dichotomy]
\label{dec:v15v28-next}
Return to the normalized chiral fork of V16.  If the hard floor
holds, record the exact determinant transport debit (28.23).  If it
fails, use the singular vectors (16.21) and the weighted locality
of the normalized chiral operator to decide whether a subsequence
is tight in physical space or escapes to ultraviolet scales.
Formulate the two alternatives without assigning the hard-branch
determinant contraction to the soft branch.
\end{decision}

\begin{firewall}[Claim boundary after V28]
\label{fire:v15v28}
V28 proves stability of a colored head contraction under tails
measured in the same weighted edge norm and derives the necessary
square-root conversion from likelihood influence.  The literal
head margin, cofinal normalized chiral floor, and shell/control
transport tail remain unpopulated.
\end{firewall}

\section*{Version V28 append-only record}
\addcontentsline{toc}{section}{Version V28 append-only record}
The complete V27 source was retained before this continuation.  It
had \(266422\) bytes, \(7610\) lines, and SHA--256
\[
 \texttt{8BBBB19ED1FB13AAA5A2B34CF8FABF621BDE0B3D19947DD3AE1DF2C99C7B2EEB}.
\]
V28 proved the common weighted head-tail perturbation theorem,
the square-root likelihood-to-Lyapunov transport conversion, and
volume-uniform determinant-tail summability under explicit decay.
\section{Fifteenth-cycle V29: physical fermions force the normalized soft branch}
\label{sec:v15v29}

V16 separated a normalized hard branch from a soft-vector branch.
The continuum target decides which branch is physically relevant.
A fermion with bounded physical Dirac energy is automatically soft
for the dimensionless normalized operator
\[
 \widetilde T_n=\frac{a_n}{2}T_n.
\]
Thus a uniform normalized gap for the entire chiral sector would
remove every finite-energy fermion from that sector.  The correct
next object is a low-energy physical bundle plus a separate
ultraviolet determinant, together with blocks measured in physical
rather than lattice distance.

\subsection{The exact scaling obstruction}

Let \(a_n\downarrow0\), and let \(T_n\) denote the physical chiral
operator on the index-zero source space.  Put
\[
 \widetilde T_n=\frac{a_n}{2}T_n.                           \tag{29.1}
\]

\begin{theorem}[Finite physical energy implies normalized softness]
\label{thm:v15v29-physical-soft}
Suppose there are unit vectors \(v_n\) and a constant \(E<\infty\)
such that
\[
 \|T_nv_n\|\le E.                                          \tag{29.2}
\]
Then
\[
 \boxed{
 s_{\min}(\widetilde T_n)
 \le\|\widetilde T_nv_n\|
 \le\frac{a_nE}{2}\longrightarrow0.}                       \tag{29.3}
\]
Conversely, if
\[
 s_{\min}(\widetilde T_n)\ge\widetilde\tau_*>0,             \tag{29.4}
\]
then every unit vector in that sector obeys
\[
 \|T_nv\|\ge\frac{2\widetilde\tau_*}{a_n}.                  \tag{29.5}
\]
Hence the sector in (29.4) contains no uniformly bounded
physical-energy sequence.
\end{theorem}

\begin{proof}
Equation (29.3) follows immediately from (29.1)--(29.2) and the
variational definition of the smallest singular value.  Reversing
the same identity gives
\[
 \|T_nv\|=\frac2{a_n}\|\widetilde T_nv\|
 \ge\frac{2\widetilde\tau_*}{a_n},
\]
which is (29.5).
\end{proof}

\begin{corollary}[Meaning of the V16 fork]
\label{cor:v15v29-fork}
A Standard Model continuum containing a fermion state of finite
physical energy must follow the V16 normalized soft branch unless
that state has first been removed from the determinant hard sector.
The V16 hard-branch determinant estimates can apply uniformly only
to an ultraviolet complement whose physical singular values are
of order \(a_n^{-1}\).
\end{corollary}

\begin{proof}
Apply Theorem \ref{thm:v15v29-physical-soft} to the lattice
approximants of the finite-energy state.  On any complement obeying
(29.4), (29.5) gives the stated ultraviolet scale.
\end{proof}

This is not a failure of the overlap normalization.  It is the
expected distinction between a dimensionless lattice operator and a
physical Dirac operator.

\subsection{Why graph-distance determinant locality degenerates}

Let \(H_n\) be a finite-range normalized operator with range
\(R_0\), uniform Schur bound \(k_0\), and singular floor
\(\tau_n\).  The finite-range V13 conjugation condition is
\[
 k_0(e^{\alpha_nR_0}-1)\le\frac{\tau_n}{2}.                 \tag{29.6}
\]
The largest exponent certified by this condition is
\[
 \alpha_n^{\rm CT}
 =
 \frac1{R_0}\log\left(1+\frac{\tau_n}{2k_0}\right).          \tag{29.7}
\]

\begin{proposition}[Physical mass gives a vanishing lattice
exponent]
\label{prop:v15v29-exponent}
If a physical singular floor is \(m_*>0\), then the normalized floor
has scale
\[
 \tau_n=\frac{a_nm_*}{2}.                                   \tag{29.8}
\]
Consequently
\[
 \alpha_n^{\rm CT}
 =
 \frac{a_nm_*}{4k_0R_0}+O(a_n^2),                           \tag{29.9}
\]
so the V13 graph-distance decay length is of order
\((a_nm_*)^{-1}\).
\end{proposition}

\begin{proof}
Equation (29.8) is (29.1) at the singular-value level.  Substitute
it into (29.7) and use
\(\log(1+x)=x+O(x^2)\).
\end{proof}

On a four-dimensional bounded-geometry lattice,
\[
 \#\{y:r\le d(x,y)<r+1\}\le C(1+r)^3.                       \tag{29.10}
\]
The exponential series in the V28 transport tail then obeys
\[
 \sum_{r\ge0}(1+r)^3e^{-c a_nr}
 \asymp a_n^{-4}.                                          \tag{29.11}
\]

\begin{lemma}[Four-dimensional exponential sum]
\label{lem:v15v29-sum}
For every fixed \(c>0\), there are \(0<C_1<C_2<\infty\) such that
for all sufficiently small \(a>0\),
\[
 C_1a^{-4}
 \le
 \sum_{r=0}^\infty(1+r)^3e^{-car}
 \le
 C_2a^{-4}.                                                 \tag{29.12}
\]
\end{lemma}

\begin{proof}
Compare the sum with
\[
 \int_0^\infty(1+x)^3e^{-cax}\,\dd x.
\]
After \(u=ax\), the leading term is
\(a^{-4}\int_0^\infty u^3e^{-cu}\,\dd u\); the lower powers give
smaller orders.  Standard monotone sum--integral comparison gives
the two-sided bound.
\end{proof}

\begin{assessment}[Limit of the graph-scale hard estimate]
\label{ass:v15v29-graph}
If the normalized floor is only the physical scale (29.8), the
V13--V28 upper bound for a graph-distance tail carries a sum of
order \(a_n^{-4}\).  Therefore that proof does not yield a uniform
single-edge influence norm.  Equation (29.11) is a statement about
the certified bound, not a lower bound on the exact determinant
response; cancellations could improve the latter and would need a
separate proof.

A uniform exponent in graph distance is naturally compatible with
an ultraviolet floor of order one for \(\widetilde T_n\), which by
(29.5) is a physical floor of order \(a_n^{-1}\).  It is not
compatible with retaining the full finite-energy fermion sector in
the same hard determinant estimate.
\end{assessment}

\subsection{Low-energy projection and physical-scale blocks}

Choose a physical threshold \(E_*>0\) and let
\[
 P_n^{\rm low}
 =
 {\bf 1}_{[0,E_*]}(|T_n|),\qquad
 Q_n^{\rm high}=I-P_n^{\rm low}.                           \tag{29.13}
\]
The determinant factorization is
\[
 \det(T_n^*T_n)
 =
 \det\nolimits_{\operatorname{Ran}P_n^{\rm low}}(T_n^*T_n)
 \,
 \det\nolimits_{\operatorname{Ran}Q_n^{\rm high}}(T_n^*T_n),
                                                               \tag{29.14}
\]
whenever neither restricted factor contains an untreated zero mode.
Kernel lines and orientations remain separate as in V13.

\begin{proposition}[A fixed physical cutoff is not a normalized
hard cutoff]
\label{prop:v15v29-fixed-cutoff}
On \(Q_n^{\rm high}\), a fixed physical threshold gives only
\[
 s_{\min}(\widetilde T_n|_{\operatorname{Ran}Q_n^{\rm high}})
 \ge\frac{a_nE_*}{2},                                      \tag{29.15}
\]
which still vanishes.  To obtain a uniform normalized floor, the
high threshold must instead have the form
\[
 E_{*,n}\ge\frac{2\widetilde\tau_*}{a_n}.                  \tag{29.16}
\]
Thus all modes below a fixed fraction of the ultraviolet cutoff
must be treated outside the V16 uniform hard determinant.
\end{proposition}

\begin{proof}
Scale the singular values by \(a_n/2\), as in (29.1).
Equations (29.15)--(29.16) follow.
\end{proof}

The other natural repair is spatial rather than spectral.  Let a
coarse block have fixed physical diameter \(\ell>0\), hence lattice
radius
\[
 R_n^{\rm block}\asymp\frac{\ell}{a_n}.                     \tag{29.17}
\]
If physical inverse decay is \(e^{-m_*a_nd_{\rm lat}}\), then between
coarse blocks at distance \(d_{\rm c}\),
\[
 e^{-m_*a_nd_{\rm lat}}
 \lesssim e^{-m_*\ell d_{\rm c}}.                          \tag{29.18}
\]
The exponent is now uniform in coarse distance.

\begin{proposition}[Necessary prefactor condition after coarse
rescaling]
\label{prop:v15v29-coarse}
Suppose the response between coarse blocks satisfies
\[
 A_{BC}^{(n)}
 \le C_n e^{-m_*\ell d_{\rm c}(B,C)}.                      \tag{29.19}
\]
On a coarse graph of uniform polynomial growth, its weighted row or
column sum is uniform in \(n\) if
\[
 \sup_n C_n<\infty.                                        \tag{29.20}
\]
Physical rescaling alone does not prove (29.20): local core ranks,
block boundary areas, and determinant trace factors grow with
\(R_n^{\rm block}\).
\end{proposition}

\begin{proof}
Under (29.20), sum (29.19) over coarse shells exactly as in Theorem
\ref{thm:v15v28-det-tail}; the fixed exponent
\(m_*\ell\) makes the series uniform.  The final statement follows
because (29.17) makes the number of fine degrees of freedom in a
coarse block diverge.
\end{proof}

\begin{assessment}[Required reorganization of the continuum proof]
\label{ass:v15v29-reorganization}
The full continuum programme cannot place physical finite-energy
fermions inside a determinant sector with a uniform normalized
singular floor.  It must do at least one of the following:
\begin{enumerate}
\item transport a growing low-energy fermion bundle explicitly and
      prove its continuum determinant or correlation limit
      separately from the ultraviolet complement;
\item prove cancellations that replace the absolute trace-rank
      bounds of V15 by uniform determinant densities;
\item use scale-dependent blocks of fixed physical diameter and
      prove the prefactor condition (29.20).
\end{enumerate}
The literal source transports physical kernels, divisors, and
orientation data, but it does not provide this growing low-energy
split or a uniform coarse-block determinant density.
\end{assessment}

\begin{decision}[V30 audit, then physical-block determinant density]
\label{dec:v15v29-next}
At V30, inspect the newest YM and NS companions.  Then replace the
rank factor in the V15 two-core trace estimate by a normalized local
trace density on blocks of fixed physical diameter.  Identify the
boundary-area scaling and the subtraction needed for a uniform
coarse-block prefactor.  Keep the growing low-energy determinant
outside the ultraviolet hard estimate.
\end{decision}

\begin{firewall}[Claim boundary after V29]
\label{fire:v15v29}
V29 proves that finite-energy physical fermions force collapse of
the normalized chiral floor and that the graph-scale determinant
tail bound then loses uniformity in four dimensions.  It identifies
the required low/high split and physical-block prefactor condition;
neither is populated by the source.
\end{firewall}

\section*{Version V29 append-only record}
\addcontentsline{toc}{section}{Version V29 append-only record}
The complete V28 source was retained before this continuation.  It
had \(276459\) bytes, \(7945\) lines, and SHA--256
\[
 \texttt{8B42D833C958335F48355E204344C55F50D94A3978835F86D28E6EA1F0C5617E}.
\]
V29 proved the physical-energy normalized-softness theorem,
quantified the vanishing graph-distance exponent, and isolated the
low-energy bundle and coarse-block determinant-density obligations.
\section{Fifteenth-cycle V30: companion audit and the determinant-density firewall}
\label{sec:v15v30}

This compulsory checkpoint audits the current companion endpoints
before the physical-block determinant calculation.  It also rules
out a tempting but invalid shortcut: multiplying the lattice log
determinant by a cell volume makes a finite trace density, but it
changes the quantum measure and is not a renormalization of the
given determinant.

\subsection{Companion checkpoint}

\begin{record}[Files inspected at V30]
\label{rec:v15v30-companions}
The current Yang--Mills endpoint is
\begin{center}
\path{Renewal_Yang_Mills_Mass_Gap_Research_Notes_V37.tex},
\end{center}
with \(313763\) bytes, \(8605\) lines, and SHA--256
\[
 \texttt{C40198AC7E2569A8DA4756FA98F52BC95AD6F2BC6E31BDB3FBCA3AA1B10477A5}.
                                                               \tag{30.1}
\]
The Navier--Stokes endpoint remains
\begin{center}
\path{Renewal_NS_Stability_Research_Notes_V26.tex},
\end{center}
with \(278283\) bytes, \(5319\) lines, and SHA--256
\[
 \texttt{82C887F8C2C69E4F28FAD7C3DC8DE5B9099CB5A4BF431EBA1FFF3E91935978AD}.
                                                               \tag{30.2}
\]
\end{record}

YM V37 constructs four independent power-three charts and proves
that their tubes cover the residual strip from the last rational
centres to the actual quarter-cell face.  The transferable rule is:
a finite atlas aimed at a physical block must quantify the residual
between its discrete centres and the geometric boundary.  An
overlap graph or last-centre certificate is insufficient.  No YM
pencil floor acts on the chiral determinant.

NS V26 is unchanged and contains no physical-block determinant
estimate.

\begin{audit}[V30 companion verdict]
\label{aud:v15v30-verdict}
Import only the physical-coverage discipline from YM V37.  A future
coarse SM atlas must cover each entire physical block, including
the residual strips between lattice-aligned charts and its geometric
faces.  No companion coefficient or theorem is inserted into the
determinant response.
\end{audit}

\subsection{Why cell-volume multiplication is not renormalization}

Let the positive determinant weight be
\[
 W_n(\varphi)
 =
 \det(D_n(\varphi)^*D_n(\varphi))
 =
 e^{F_n(\varphi)}.                                         \tag{30.3}
\]
A normalized trace density would replace \(F_n\) by
\[
 \overline F_n=a_n^4F_n                                    \tag{30.4}
\]
in four dimensions.

\begin{proposition}[Determinant-density firewall]
\label{prop:v15v30-firewall}
Unless \(F_n\) is field-independent or \(a_n^4=1\), replacing
\(F_n\) by \(\overline F_n\) changes every nontrivial likelihood
ratio:
\[
 \frac{e^{\overline F_n(\varphi)}}{e^{\overline F_n(\psi)}}
 =
 \left[
 \frac{e^{F_n(\varphi)}}{e^{F_n(\psi)}}
 \right]^{a_n^4}.                                         \tag{30.5}
\]
It therefore changes the finite-cutoff Gibbs measure and cannot be
called a counterterm subtraction or an equivalent normalization.
\end{proposition}

\begin{proof}
Equation (30.5) is algebraic.  If
\(F_n(\varphi)-F_n(\psi)\ne0\), exponentiation by \(a_n^4\ne1\)
changes the ratio.  A field-independent additive constant would
cancel from the ratio; multiplication does not.
\end{proof}

The quantity \(a_n^4F_n\) may still be used as a diagnostic trace
density.  To retain the original determinant weight, divergent
local contributions must instead be removed by explicit additive
counterterms:
\[
 F_n^{\rm ren}(\varphi)
 =
 F_n(\varphi)-C_n^{\rm loc}(\varphi),                      \tag{30.6}
\]
with \(C_n^{\rm loc}\) proved to be local, gauge invariant, and
compatible with reflection and projective transport.  Its
coefficients must be fixed by renormalization conditions rather
than chosen to make an estimate finite.

\begin{lemma}[Mixed separated response ignores one-block
counterterms]
\label{lem:v15v30-counterterm}
Let
\[
 C_n^{\rm loc}=\sum_A C_{n,A}(\varphi_A)
\]
be a finite-range local counterterm.  If variations \(s,t\) have
supports separated farther than the counterterm range, then
\[
 \partial_s\partial_t C_n^{\rm loc}=0.                     \tag{30.7}
\]
Hence the separated mixed response of \(F_n^{\rm ren}\) equals that
of \(F_n\).
\end{lemma}

\begin{proof}
No local set \(A\) meets both variation supports.  Each summand is
therefore independent of at least one parameter, so its mixed
derivative vanishes.  Sum over \(A\).
\end{proof}

For adjacent coarse blocks, counterterms supported near their common
boundary can contribute.  Their coefficients and boundary scaling
must be retained explicitly.  The rank estimate in V15 counts an
entire block volume even though a connected two-core trace crosses
the block interface twice.  The next version replaces that rank by
a sum of exponentially weighted pairs of local cores; for a
uniformly gapped ultraviolet operator, this sum scales with the
interface area.

\begin{assessment}[Position after V30]
\label{ass:v15v30-position}
The companion audit changes no SM constant.  V29's low/high
fermion split remains necessary.  V30 rules out cell-volume
multiplication as a measure-preserving cure and identifies the
legitimate route: connected local trace expansion plus explicit
additive counterterms.

The source declares counterterms and projective transport but does
not give a local heat-kernel coefficient ledger or the uniform
boundary response needed here.
\end{assessment}

\begin{decision}[V31 connected boundary-layer trace]
\label{dec:v15v30-next}
Assume each coarse variation is a sum of uniformly bounded
site-local cores.  Expand the V15 two-core trace before applying a
rank bound.  Prove
\[
 |\operatorname{Tr}(AK_BBK_C)|
 \le C\sum_{x\in B,y\in C}e^{-2\alpha d(x,y)}.
\]
For adjacent boxes and fixed ultraviolet exponent, show that the
pair sum grows like their lattice boundary area, not their volume.
Then state exactly which boundary counterterm density would need a
uniform cofinal coefficient.
\end{decision}

\begin{firewall}[Claim boundary after V30]
\label{fire:v15v30}
Multiplying the log determinant by \(a_n^4\) is proved to change the
measure.  Only explicit local additive counterterms and a connected
boundary-layer estimate can address the coarse prefactor; those
estimates are not yet proved here.
\end{firewall}

\section*{Version V30 append-only record}
\addcontentsline{toc}{section}{Version V30 append-only record}
The complete V29 source was retained before this continuation.  It
had \(286944\) bytes, \(8247\) lines, and SHA--256
\[
 \texttt{2C8138D4C880F45CC084173398C12F1B402610428FE6D1F3D22029B471A62E76}.
\]
V30 completed the compulsory companion audit, imported the
physical-face coverage discipline, and proved that cell-volume
multiplication is not a valid determinant renormalization.
\section{Fifteenth-cycle V31: connected determinant traces have boundary scaling}
\label{sec:v15v31}

V29 identified a coarse-block prefactor problem, and V30 ruled out
multiplying the determinant action by a cell volume.  The rank bound
in V15 can nevertheless be sharpened when a coarse variation is a
sum of site-local cores.  Expanding the trace before estimating it
shows that a connected two-core response between adjacent boxes is
supported in an exponential boundary layer and scales with interface
area.

Let
\[
 {\cal H}=\bigoplus_{x\in\Lambda}{\cal H}_x,\qquad
 \sup_x\dim{\cal H}_x\le d_*                               \tag{31.1}
\]
on a finite graph.  Suppose two operators obey pointwise kernel
bounds
\[
 \|P_xAP_y\|\le C_Ae^{-\alpha d(x,y)},\qquad
 \|P_xBP_y\|\le C_Be^{-\alpha d(x,y)}.                     \tag{31.2}
\]
For finite sets \(E,F\subset\Lambda\), let
\[
 K_E=\sum_{x\in E}K_x,\qquad
 L_F=\sum_{y\in F}L_y,                                     \tag{31.3}
\]
where
\[
 K_x=P_xK_xP_x,\quad L_y=P_yL_yP_y,\qquad
 \|K_x\|\le k_E,\quad\|L_y\|\le k_F.                       \tag{31.4}
\]

\begin{theorem}[Connected two-core pair-sum bound]
\label{thm:v15v31-pair}
Under (31.1)--(31.4),
\[
 \boxed{
 |\operatorname{Tr}(AK_EBL_F)|
 \le
 d_*C_AC_Bk_Ek_F
 \sum_{x\in E}\sum_{y\in F}
 e^{-2\alpha d(x,y)}.}                                    \tag{31.5}
\]
The same conclusion holds for cores of fixed radius \(R_c\), after
enlarging \(E,F\) by \(R_c\) and changing the constant by the
bounded neighborhood multiplicity.
\end{theorem}

\begin{proof}
Expand (31.3), insert support projections, and cycle each finite
trace:
\[
 \operatorname{Tr}(AK_xBL_y)
 =
 \operatorname{Tr}\left(
 P_yAP_xK_xP_xBP_yL_y
 \right).                                                   \tag{31.6}
\]
The trace acts on a space of dimension at most \(d_*\).  Hence its
absolute value is at most
\[
 d_*
 \|P_yAP_x\|\|K_x\|
 \|P_xBP_y\|\|L_y\|,
\]
which is the summand in (31.5).  Sum over \(x,y\).  A fixed-radius
core splits into a uniformly bounded number of site blocks on a
bounded-degree graph, giving the last assertion.
\end{proof}

Unlike the rank estimate (15.6), (31.5) retains the location of
every local core and counts only pairs joined by two decaying
bridges.

\subsection{Exact interface-area order on boxes}

Take the nearest-neighbor graph on \(\mathbb Z^d\).  Let
\[
\begin{aligned}
 E_L&=\{-L+1,\ldots,0\}\times\{1,\ldots,L\}^{d-1},\\
 F_L&=\{1,\ldots,L\}\times\{1,\ldots,L\}^{d-1}.             \tag{31.7}
\end{aligned}
\]
These boxes share a lattice interface.  Put \(q=e^{-\kappa}\).

\begin{lemma}[Exponential pair sum is boundary order]
\label{lem:v15v31-area}
For every \(\kappa>0\),
\[
 \boxed{
 e^{-\kappa}L^{d-1}
 \le
 \sum_{x\in E_L}\sum_{y\in F_L}e^{-\kappa|x-y|_1}
 \le
 C_{\kappa,d}L^{d-1},}                                    \tag{31.8}
\]
where
\[
 C_{\kappa,d}
 =
 \frac{q}{(1-q)^2}
 \left(\frac{1+q}{1-q}\right)^{d-1}.                       \tag{31.9}
\]
If a lattice gap of \(g\ge0\) layers separates the boxes, the upper
bound gains the factor \(e^{-\kappa g}\).
\end{lemma}

\begin{proof}
For the lower bound, retain only pairs with equal tangential
coordinates and normal coordinates \(0,1\).  There are
\(L^{d-1}\) such pairs, each at distance one.

Write \(x_1=-i\), \(0\le i<L\), and \(y_1=j\),
\(1\le j\le L\).  The normal sum is bounded by
\[
 \left(\sum_{i\ge0}q^i\right)
 \left(\sum_{j\ge1}q^j\right)
 =
 \frac{q}{(1-q)^2}.                                       \tag{31.10}
\]
For every fixed tangential \(x'\), the infinite tangential sum is
\[
 \sum_{z\in\mathbb Z^{d-1}}q^{|z|_1}
 =
 \left(\frac{1+q}{1-q}\right)^{d-1}.                       \tag{31.11}
\]
There are \(L^{d-1}\) choices of \(x'\).  Multiplication proves the
upper bound.  A gap of \(g\) adds at least \(g\) to every normal
distance.
\end{proof}

\begin{corollary}[Ultraviolet hard determinant has interface
scaling]
\label{cor:v15v31-uv}
Suppose every bridge operator occurring in the V15 determinant
Hessian on an ultraviolet sector satisfies (31.2) with
\(\alpha\ge\alpha_*>0\), and all site-core norms are uniform.  For
adjacent \(d\)-dimensional boxes of lattice side \(L\),
\[
 |\partial_E\partial_FF_{\rm UV}|
 \le C_{\rm UV}L^{d-1}.                                    \tag{31.12}
\]
In four dimensions, a block of fixed physical side \(\ell\) has
\(L_n\asymp\ell/a_n\), so
\[
 |\partial_E\partial_FF_{{\rm UV},n}|
 \le C_{\rm UV}\ell^3a_n^{-3}.                             \tag{31.13}
\]
\end{corollary}

\begin{proof}
Every term in the first and second Jacobi traces factors into two
local cores joined cyclically by exponentially local products of
resolvents, normalized chiral factors, and projectors.  Weighted
Schur submultiplicativity preserves (31.2), with a possibly smaller
uniform exponent.  Apply Theorem \ref{thm:v15v31-pair} and Lemma
\ref{lem:v15v31-area}.  Substitute \(L_n\asymp\ell/a_n\) when
\(d=4\).
\end{proof}

The conclusion applies only to the ultraviolet complement with a
uniform normalized gap.  The physical low-energy sector of V29 has
a decay exponent of order \(a_n\) and is not covered by
Corollary \ref{cor:v15v31-uv}.

\subsection{Boundary-extensive transport removes the artificial
prefactor}

Let \(\partial(E,F)\) be the set of microscopic faces crossing a
coarse interface.  Define its transport cost by summation,
\[
 {\cal D}_{\partial(E,F)}
 =
 \sum_{e\in\partial(E,F)}D_{e,\beta},                      \tag{31.14}
\]
rather than truncating the whole coarse interface to one bounded
cost.

\begin{theorem}[Uniform coarse response in boundary-extensive norm]
\label{thm:v15v31-boundary-metric}
Suppose the microscopic ultraviolet determinant response matrix has
an exponentially summable column profile, uniformly in \(n\), as
in V28.  Group microscopic edges into coarse interfaces and use
(31.14).  Then the response from one coarse interface to an adjacent
one has a coefficient bounded independently of
\(|\partial(E,F)|\), \(L_n\), and \(a_n\).
\end{theorem}

\begin{proof}
Let \(T_{ef}\) be the microscopic response coefficient.  For an
input interface \(I\) and output interface \(J\),
\[
 \sum_{e\in J}\mathbb ED_{e,\beta}'
 \le
 \sum_{f\in I}
 \left(\sum_{e\in J}T_{ef}\right)D_{f,\beta}
 \le
 \eta_{\rm UV}\sum_{f\in I}D_{f,\beta},                    \tag{31.15}
\]
where
\[
 \eta_{\rm UV}=\sup_f\sum_eT_{ef}<\infty
\]
is the V28 column norm.  The number of faces occurs on both sides
through the sums and cancels from the coefficient.  No division or
multiplication of the determinant action is performed.
\end{proof}

This identifies why the prefactor \(C_n\) in (29.19) appeared:
collapsing an interface containing \(O(a_n^{-3})\) microscopic
sources into one bounded coarse coordinate discards its extensive
input cost.  The correct coarse metric retains the interface sum.

\subsection{What counterterm density is actually required}

For continuum convergence of the determinant action itself, define
a local counterterm \(C_n^{\rm loc}\) as in (30.6).  For adjacent
physical blocks with interface area
\[
 |\Sigma|_{\rm phys}=a_n^3|\partial(E,F)|,                  \tag{31.16}
\]
the renormalized connected face density is
\[
 {\mathfrak b}_n(s,t;\Sigma)
 =
 \frac{a_n^3}{|\Sigma|_{\rm phys}}
 \partial_s\partial_t
 \bigl(F_n-C_n^{\rm loc}\bigr).                            \tag{31.17}
\]

\begin{proposition}[Boundary-density obligation]
\label{prop:v15v31-counterterm}
A uniform coarse determinant-density claim requires
\[
 \sup_n\sup_{\Sigma,s,t}
 |{\mathfrak b}_n(s,t;\Sigma)|<\infty                      \tag{31.18}
\]
for the declared physical variation class.  The subtraction in
(31.17) is legitimate only if:
\begin{enumerate}
\item \(C_n^{\rm loc}\) is a sum of gauge-invariant,
      reflection-compatible local motifs already present as
      admissible counterterms;
\item its coefficients are fixed by stated renormalization
      conditions;
\item motifs crossing an artificial coarse interface are obtained
      by restricting the same bulk local functional, rather than
      adding a new physical boundary action.
\end{enumerate}
Corollary \ref{cor:v15v31-uv} proves the unrenormalized
boundary order but does not populate (31.18).
\end{proposition}

\begin{proof}
The factor \(a_n^3/|\Sigma|_{\rm phys}\) is exactly the reciprocal
number of microscopic faces per unit physical interface area.
Thus (31.18) is the statement that the connected mixed response per
physical area remains bounded.  The three listed conditions are
necessary for subtracting \(C_n^{\rm loc}\) without changing gauge
covariance, reflection positivity, or the bulk theory at artificial
block boundaries.  V30 proves that a multiplicative trace-density
replacement would not satisfy these conditions.
\end{proof}

\begin{assessment}[Gain and remaining low-sector problem]
\label{ass:v15v31-reading}
For a uniformly normalized-gapped ultraviolet determinant, V31
replaces the coarse volume rank \(L^4\) by the correct interface
order \(L^3\) and shows that a boundary-extensive transport metric
has a volume-independent response coefficient.  This resolves the
artificial prefactor in the transport assembly.

It does not renormalize the determinant action and does not control
the growing physical low-energy bundle.  The literal source has no
cofinal local counterterm ledger proving (31.18), and its transported
kernel/divisor data do not give a continuum determinant for all
modes below order \(a_n^{-1}\).
\end{assessment}

\begin{decision}[V32 physical low-mode determinant]
\label{dec:v15v31-next}
Treat the low-energy projection \(P_n^{\rm low}\) of V29 as a
finite-rank bundle at each cutoff.  Derive a basis-free relative
determinant under Kato transport and state trace-class convergence
conditions sufficient for its continuum limit.  Separate:
\begin{enumerate}
\item fixed finite physical windows, where ranks may stabilize on a
      compact continuum geometry;
\item growing windows required to reconstruct the full determinant;
\item exact zero modes and their determinant line.
\end{enumerate}
Do not infer the growing-window limit from fixed-window convergence.
\end{decision}

\begin{firewall}[Claim boundary after V31]
\label{fire:v15v31}
Connected ultraviolet determinant response has boundary rather than
volume scaling and is uniform in a boundary-extensive transport
norm.  The renormalized face density and the physical low-energy
determinant remain unproved.
\end{firewall}

\section*{Version V31 append-only record}
\addcontentsline{toc}{section}{Version V31 append-only record}
The complete V30 source was retained before this continuation.  It
had \(293741\) bytes, \(8431\) lines, and SHA--256
\[
 \texttt{18BF386E8C4F1916100B76166348548CF6EAEE2EF1AB499319AB6843DE7F4BB4}.
\]
V31 proved the connected two-core pair-sum estimate, exact
interface-area scaling, and uniform ultraviolet response in a
boundary-extensive coarse metric.
\section{Fifteenth-cycle V32: basis-free physical low-mode determinants}
\label{sec:v15v32}

V29 proved that finite-energy fermions cannot remain in the
uniformly normalized-gapped ultraviolet determinant.  V31 handled
the connected ultraviolet boundary response.  This continuation
constructs the complementary low-energy determinant in a
basis-free form and states the exact convergence needed when the
energy window grows.

Let \(T_n:{\cal H}_n^-\to{\cal H}_n^+\) be the physical chiral
operator after removal of exact zero modes.  For a fixed physical
energy \(E>0\), put
\[
 P_{n,E}={\bf 1}_{(0,E]}(|T_n|),\qquad
 {\cal L}_{n,E}=\operatorname{Ran}P_{n,E}.                 \tag{32.1}
\]
Assume for the moment that
\[
 \dim{\cal L}_{n,E}=r_E                                   \tag{32.2}
\]
is independent of \(n\) on a cofinal tail.

\subsection{Canonical transport of a fixed window}

After embedding the cutoff spaces in a common Hilbert space, let
\(P_E\) be a reference projection and suppose
\[
 \|P_{n,E}-P_E\|<1.                                       \tag{32.3}
\]
Define
\[
 S_{n,E}
 =
 P_EP_{n,E}+(I-P_E)(I-P_{n,E})                             \tag{32.4}
\]
and its polar unitary
\[
 U_{n,E}
 =
 S_{n,E}(S_{n,E}^*S_{n,E})^{-1/2}.                        \tag{32.5}
\]
Then
\[
 U_{n,E}P_{n,E}U_{n,E}^*=P_E.                             \tag{32.6}
\]

\begin{lemma}[Canonical Kato transport]
\label{lem:v15v32-kato}
Under (32.3), \(S_{n,E}\) is invertible, (32.5) is unitary, and
(32.6) holds.  The construction depends only on the two
projections and chooses no basis of either low-energy space.
\end{lemma}

\begin{proof}
A direct calculation gives
\[
 S_{n,E}^*S_{n,E}
 =
 I-(P_{n,E}-P_E)^2.                                       \tag{32.7}
\]
Condition (32.3) makes this operator strictly positive, so the polar
factor is defined and unitary.  The intertwining identity
\[
 P_ES_{n,E}=S_{n,E}P_{n,E}
\]
commutes with the positive polar factor in the required functional
calculus and yields (32.6).
\end{proof}

Let
\[
 A_{n,E}
 =
 P_{n,E}T_n^*T_nP_{n,E}|_{{\cal L}_{n,E}}                 \tag{32.8}
\]
and transport it to the fixed space
\({\cal L}_E=\operatorname{Ran}P_E\):
\[
 \widehat A_{n,E}
 =
 U_{n,E}A_{n,E}U_{n,E}^*|_{{\cal L}_E}.                   \tag{32.9}
\]
Choose a positive reference operator \(A_E^0\) on \({\cal L}_E\)
and define
\[
 R_{n,E}
 =
 (A_E^0)^{-1/2}\widehat A_{n,E}(A_E^0)^{-1/2}.             \tag{32.10}
\]

\begin{definition}[Fixed-window relative determinant]
\label{def:v15v32-relative}
The basis-free fixed-window determinant is
\[
 \boxed{
 \Delta_{n,E}
 =
 \det_{{\cal L}_E}R_{n,E}
 =
 \frac{\det_{{\cal L}_{n,E}}A_{n,E}}
      {\det_{{\cal L}_E}A_E^0}.}                           \tag{32.11}
\]
\end{definition}

The second equality follows from unitary invariance of determinant.
Changing an orthonormal basis in either window does not change
\(\Delta_{n,E}\).

\begin{theorem}[Fixed-window determinant convergence]
\label{thm:v15v32-fixed}
Suppose, for a fixed \(E\),
\[
 \|R_{n,E}-R_E\|_1\longrightarrow0,\qquad
 R_{n,E},R_E\ge\delta_EI>0.                                \tag{32.12}
\]
Then
\[
 \boxed{
 \Delta_{n,E}\longrightarrow
 \det_{{\cal L}_E}R_E>0.}                                  \tag{32.13}
\]
Moreover,
\[
 |\log\Delta_{n,E}-\log\det R_E|
 \le
 \delta_E^{-1}\|R_{n,E}-R_E\|_1.                           \tag{32.14}
\]
\end{theorem}

\begin{proof}
Along
\[
 R_E(s)=(1-s)R_E+sR_{n,E},
\]
the lower bound in (32.12) remains valid.  Jacobi differentiation
gives
\[
 \log\det R_{n,E}-\log\det R_E
 =
 \int_0^1
 \operatorname{Tr}\left(
 R_E(s)^{-1}(R_{n,E}-R_E)
 \right)\,\dd s.
\]
Trace duality bounds the integrand by
\(\delta_E^{-1}\|R_{n,E}-R_E\|_1\), proving
(32.14) and hence (32.13).
\end{proof}

\subsection{When fixed-window ranks stabilize}

Let \(|T|\) be a continuum operator with compact resolvent and
suppose \(E\notin\operatorname{spec}|T|\).  Assume, after the
declared embeddings, norm-resolvent convergence
\[
 (|T_n|-z)^{-1}\longrightarrow(|T|-z)^{-1}                 \tag{32.15}
\]
uniformly for \(z\) on a contour \(\Gamma_E\) enclosing
\((0,E]\) and no other continuum spectrum.

\begin{proposition}[Spectral-window stabilization]
\label{prop:v15v32-window}
Under (32.15),
\[
 P_{n,E}
 =
 \frac{1}{2\pi i}\int_{\Gamma_E}(z-|T_n|)^{-1}\,\dd z
 \longrightarrow
 P_E
 =
 \frac{1}{2\pi i}\int_{\Gamma_E}(z-|T|)^{-1}\,\dd z        \tag{32.16}
\]
in operator norm.  For all sufficiently large \(n\), (32.2)--(32.3)
hold and
\[
 \dim{\cal L}_{n,E}=\dim\operatorname{Ran}P_E<\infty.       \tag{32.17}
\]
\end{proposition}

\begin{proof}
Integrate the uniform resolvent convergence around
\(\Gamma_E\) to obtain norm convergence of the Riesz projections.
Projections at distance less than one have isomorphic ranges by
Lemma \ref{lem:v15v32-kato}, so their ranks agree.  Compact
resolvent makes the continuum window finite-dimensional.
\end{proof}

This proposition supplies projection transport, not (32.12).
Trace-norm convergence of the transported positive operators is a
separate analytic input.

\subsection{Exact zero modes remain determinant-line data}

Let
\[
 Z_n^-=\ker T_n,\qquad Z_n^+=\ker T_n^*.                    \tag{32.18}
\]
The positive pseudo-determinant is
\[
 \det{}'(T_n^*T_n)
 =
 \det_{\,(Z_n^-)^\perp}(T_n^*T_n).                         \tag{32.19}
\]
Its omitted zero-mode data live in the line
\[
 {\cal D}_n
 =
 \det Z_n^+\otimes(\det Z_n^-)^*.                          \tag{32.20}
\]

\begin{proposition}[Zero-mode separation]
\label{prop:v15v32-zero}
If \(\dim Z_n^\pm\) are constant and their projections converge in
norm, Kato transport identifies the lines (32.20) up to their
physical phase or orientation transport.  The positive determinant
(32.19) does not determine that phase or orientation.

If nonzero singular values merely tend to zero, they remain factors
of (32.19) until a declared spectral-flow or divisor rule moves them
into (32.20).
\end{proposition}

\begin{proof}
Apply exterior power to the Kato unitaries on \(Z_n^\pm\).
Exterior power transports a determinant line but does not choose its
phase relative to an independent physical orientation.  A positive
singular value is not in the kernel, so algebraically it remains in
the pseudo-determinant.  Moving it to the kernel line requires a
specified crossing rule.
\end{proof}

\subsection{Growing windows and Fredholm determinants}

Fixed physical windows do not reconstruct the full determinant.
Let \({\cal L}\) be a reference Hilbert space and suppose transported
relative operators have the form
\[
 R_n=I+K_n,\qquad K_n\in{\mathfrak S}_1({\cal L}),          \tag{32.21}
\]
where \({\mathfrak S}_1\) is the trace class.

\begin{theorem}[Trace-class relative determinant limit]
\label{thm:v15v32-fredholm}
Suppose
\[
 \|K_n-K\|_1\longrightarrow0,\qquad
 I+K_n,I+K\ge\delta I>0.                                   \tag{32.22}
\]
Then the Fredholm determinants satisfy
\[
 \boxed{
 \det_{\rm F}(I+K_n)\longrightarrow\det_{\rm F}(I+K)>0,}    \tag{32.23}
\]
and
\[
 \left|
 \log\det_{\rm F}(I+K_n)
 -\log\det_{\rm F}(I+K)
 \right|
 \le\delta^{-1}\|K_n-K\|_1.                                \tag{32.24}
\]
\end{theorem}

\begin{proof}
The convex path between \(I+K\) and \(I+K_n\) has inverse norm at
most \(\delta^{-1}\).  The trace-class Jacobi formula and trace
duality give (32.24), exactly as in Theorem
\ref{thm:v15v32-fixed}.  Continuity of the exponential gives
(32.23).
\end{proof}

Let \(K_n^{(E)}\) be the transported restriction to a fixed energy
window.

\begin{corollary}[Sufficient growing-window condition]
\label{cor:v15v32-growing}
Assume:
\begin{enumerate}
\item for every admissible fixed \(E\),
      \[
       K_n^{(E)}\longrightarrow K^{(E)}
       \quad\hbox{in trace norm};                           \tag{32.25}
      \]
\item the discarded relative tails are uniformly trace-small,
      \[
       \lim_{E\to\infty}\sup_n
       \|K_n-K_n^{(E)}\|_1=0;                               \tag{32.26}
      \]
\item
      \[
       \lim_{E\to\infty}\|K-K^{(E)}\|_1=0,                 \tag{32.27}
      \]
      and all relative operators have a common positive floor.
\end{enumerate}
Then (32.22)--(32.23) hold.
\end{corollary}

\begin{proof}
For every \(E\),
\[
\begin{aligned}
 \|K_n-K\|_1
 \le{}&
 \|K_n-K_n^{(E)}\|_1+
 \|K_n^{(E)}-K^{(E)}\|_1\\
 &+\|K^{(E)}-K\|_1.
\end{aligned}
\]
First choose \(E\) so the first and third terms are uniformly small,
then let \(n\to\infty\) in the middle term.  Apply Theorem
\ref{thm:v15v32-fredholm}.
\end{proof}

\begin{proposition}[Fixed windows do not control an escaping
determinant factor]
\label{prop:v15v32-counterexample}
There are positive trace-class perturbations \(I+K_n\) such that
every fixed coordinate window converges exactly to the identity,
while
\[
 \det_{\rm F}(I+K_n)=2                                     \tag{32.28}
\]
for every \(n\).
\end{proposition}

\begin{proof}
On \(\ell^2(\mathbb N)\), let
\[
 K_n=|e_n\rangle\langle e_n|.
\]
For the projection onto the first \(E\) basis vectors,
\(P_EK_nP_E=0\) whenever \(n>E\), so every fixed window converges
to zero.  But \(I+K_n\) has one eigenvalue two and all remaining
eigenvalues one, proving (32.28).  Here
\(\|K_n\|_1=1\), so the missing uniform-tail condition (32.26)
fails exactly.
\end{proof}

\begin{assessment}[Literal source applicability]
\label{ass:v15v32-source}
The source contains projective transport of physical kernels,
divisors, determinant lines, and orientations.  Those structures
are compatible with Lemma \ref{lem:v15v32-kato} and Proposition
\ref{prop:v15v32-zero}.  It does not supply:
\begin{enumerate}
\item norm-resolvent convergence of the physical chiral operators;
\item trace-norm convergence (32.12) on each fixed window;
\item the uniform growing-window tail (32.26);
\item a common positive relative floor after the exact zero modes
      are removed.
\end{enumerate}
Therefore neither a fixed-window nor a full low-energy continuum
determinant is presently populated by the literal source.
\end{assessment}

\begin{decision}[V33 local variation of the low spectral bundle]
\label{dec:v15v32-next}
Differentiate the Riesz projection (32.16) under a local field
variation.  Prove its quasilocal kernel bound when the energy-window
contour has a uniform spectral separation, and insert the Kato
connection into the relative determinant derivative.  If the
window edge accumulates spectrum and the contour gap closes, return
the resulting crossing modes to the divisor/spectral-flow ledger
rather than asserting smooth bundle transport.
\end{decision}

\begin{firewall}[Claim boundary after V32]
\label{fire:v15v32}
V32 constructs basis-free fixed-window and Fredholm relative
determinants and proves sufficient convergence conditions.  The
source does not populate the required resolvent, trace-norm, or
uniform growing-tail estimates, and fixed-window convergence alone
is proved insufficient.
\end{firewall}

\section*{Version V32 append-only record}
\addcontentsline{toc}{section}{Version V32 append-only record}
The complete V31 source was retained before this continuation.  It
had \(304666\) bytes, \(8744\) lines, and SHA--256
\[
 \texttt{2E837F901E1D1F9C71162CF89A27856ECB7C8C1C28A42909E65F23A0BC33E6EC}.
\]
V32 constructed the Kato-transported low-energy relative
determinant, proved fixed- and growing-window convergence criteria,
and separated exact zero-mode determinant lines from soft nonzero
modes.
\section{Fifteenth-cycle V33: local variation and crossings of the low-energy bundle}
\label{sec:v15v33}

V32 constructed the low-energy determinant after Kato transport but
did not control its variation with the fields.  Work with the
positive operator
\[
 B=T^*T                                                     \tag{33.1}
\]
and a Riesz contour around the squared physical energy window.
A uniform contour separation gives quasilocality of the projection
and its derivative.  The Kato connection is necessary for the
transported operator, although its commutator cancels from the
scalar determinant response.

Let \(B=B^*\ge0\) act on
\({\cal H}=\bigoplus_x{\cal H}_x\).  Suppose
\[
 \|B\|_{\alpha_0}
 =
 \max\left\{
 \sup_x\sum_y e^{\alpha_0d(x,y)}\|B_{xy}\|,
 \sup_y\sum_x e^{\alpha_0d(x,y)}\|B_{xy}\|
 \right\}
 \le K_0.                                                   \tag{33.2}
\]
Let \(\Gamma\) be a positively oriented contour of length
\(L_\Gamma\) such that
\[
 \operatorname{dist}(\Gamma,\operatorname{spec}B)\ge\delta>0.
                                                               \tag{33.3}
\]
It encloses the desired positive low-energy spectrum and excludes
zero and the high spectrum.  Put
\[
 P=\frac{1}{2\pi i}\int_\Gamma(z-B)^{-1}\,\dd z.            \tag{33.4}
\]

For a one-Lipschitz weight \(\rho\), let
\(W_\alpha=e^{\alpha\rho}\).  Define
\[
 \epsilon_\alpha
 =
 \sup_\rho\|W_\alpha BW_\alpha^{-1}-B\|.                   \tag{33.5}
\]
Weighted Schur locality makes \(\epsilon_\alpha\to0\) as
\(\alpha\downarrow0\).

\begin{theorem}[Quasilocal low-energy projection]
\label{thm:v15v33-projection}
If
\[
 \epsilon_\alpha\le\frac\delta2,                            \tag{33.6}
\]
then for every \(z\in\Gamma\) and vertex sets \(X,Y\),
\[
 \|P_X(z-B)^{-1}P_Y\|
 \le\frac2\delta e^{-\alpha d(X,Y)},                       \tag{33.7}
\]
and
\[
 \boxed{
 \|P_XPP_Y\|
 \le\frac{L_\Gamma}{\pi\delta}
 e^{-\alpha d(X,Y)}}                                      \tag{33.8}
\]
for disjoint \(X,Y\).
\end{theorem}

\begin{proof}
Since \(B\) is self-adjoint, (33.3) gives
\(\|(z-B)^{-1}\|\le\delta^{-1}\).  The weighted conjugate differs
from \(B\) by at most \(\delta/2\), so its contour resolvent has norm
at most \(2/\delta\) by a Neumann series.  Anchor the Lipschitz
weight at \(Y\), as in V13--V14, to obtain (33.7).  Integrate
(33.7) over \(\Gamma\) and divide by \(2\pi\), proving (33.8).
\end{proof}

Let \(s\) be a local field parameter with core \(S\), and assume
\[
 B_s=P_SB_sP_S,\qquad\|B_s\|\le K_s.                       \tag{33.9}
\]
The fixed-radius and core-factorized variants follow by the V15
calculus.

\begin{theorem}[Local derivative of the spectral projection]
\label{thm:v15v33-derivative}
Under (33.3), (33.6), and (33.9),
\[
 P_s
 =
 \frac{1}{2\pi i}
 \int_\Gamma
 (z-B)^{-1}B_s(z-B)^{-1}\,\dd z,                           \tag{33.10}
\]
and
\[
 \boxed{
 \|P_XP_sP_Y\|
 \le
 \frac{2L_\Gamma K_s}{\pi\delta^2}
 e^{-\alpha[d(X,S)+d(S,Y)]}.}                              \tag{33.11}
\]
If \(B_s\) is core-factorized rather than exactly supported, the
same conclusion holds with the V15 factorization cost replacing
\(K_s\).
\end{theorem}

\begin{proof}
Differentiate the finite contour integral (33.4).  The resolvent
identity gives (33.10).  Insert \(P_S\) on both sides of \(B_s\),
apply (33.7) to the two resolvents, and multiply by
\(L_\Gamma/(2\pi)\).  This gives (33.11).  Integrating a
core-factorized derivative term by term proves the last statement.
\end{proof}

For \(B=T^*T\),
\[
 B_s=T_s^*T+T^*T_s.                                       \tag{33.12}
\]
Thus the normalized overlap and Yukawa derivative estimates of
V14--V19 populate (33.9) or its core-factorized version on a
uniformly separated window.

\subsection{Quasilocal Kato transport}

Let \(P(s)\) be a differentiable family of projections of constant
rank and define the anti-Hermitian Kato generator
\[
 {\cal K}(s)=[P_s(s),P(s)].                                \tag{33.13}
\]
Let \(U(s)\) solve
\[
 U_s={\cal K}U,\qquad U(0)=I.                              \tag{33.14}
\]
Then
\[
 U(s)P(0)U(s)^*=P(s).                                      \tag{33.15}
\]

\begin{proposition}[Locality of the Kato connection]
\label{prop:v15v33-kato-local}
If \(P,P_s\in{\mathfrak S}_\beta\) uniformly for
\(0<\beta<\alpha\), then
\[
 \|{\cal K}(s)\|_\beta
 \le2\|P_s(s)\|_\beta\|P(s)\|_\beta.                       \tag{33.16}
\]
Along a parameter path of length one,
\[
 \|U(s)\|_\beta
 \le
 \exp\left(\int_0^s\|{\cal K}(r)\|_\beta\,\dd r\right).     \tag{33.17}
\]
Hence Kato transport is quasilocal whenever the right side is
uniform.
\end{proposition}

\begin{proof}
The weighted Schur algebra estimate (15.4) applied to the two terms
of the commutator gives (33.16).  The integral equation for
(33.14), iteration, and submultiplicativity give the exponential
Dyson-series bound (33.17).
\end{proof}

Let
\[
 A(s)=P(s)B(s)P(s)|_{\operatorname{Ran}P(s)}
\]
and transport it to \(\operatorname{Ran}P(0)\):
\[
 \widehat A(s)=U(s)^*A(s)U(s).                             \tag{33.18}
\]

\begin{theorem}[Connection cancellation in the determinant
response]
\label{thm:v15v33-connection}
One has
\[
 \widehat A_s
 =
 U^*\bigl(A_s+[A,{\cal K}]\bigr)U,                         \tag{33.19}
\]
but
\[
 \boxed{
 \frac{\dd}{\dd s}\log\det\widehat A
 =
 \operatorname{Tr}_{\operatorname{Ran}P(s)}
 (A^{-1}A_s).}                                             \tag{33.20}
\]
Thus the connection is required to differentiate the transported
operator, while its commutator contributes zero to the scalar
positive determinant.
\end{theorem}

\begin{proof}
Differentiate (33.18), use \(U_s={\cal K}U\) and
\({\cal K}^*=-{\cal K}\), and obtain (33.19).  Since \(A=PAP\) and \(P{\cal K}P=0\), compression to the moving
low-energy space gives
\[
 P[A,{\cal K}]P=0.                                         \tag{33.21}
\]
Jacobi's formula therefore receives no connection contribution.
Unitary transport identifies the remaining trace with (33.20).
\end{proof}

This cancellation concerns the scalar determinant.  The connection
does not disappear from transported eigenvectors, determinant
lines, mixed source Gramians, or operator trace-norm convergence.

\subsection{Crossings at the window edge}

Suppose a simple eigenvalue \(\lambda(s)>0\) of \(B(s)\) crosses
the window boundary \(E^2\) at \(s=s_0\).  Then the rank of
\[
 P_E(s)={\bf 1}_{(0,E^2]}(B(s))                             \tag{33.22}
\]
changes at \(s_0\), and no contour satisfying (33.3) exists on a
neighborhood of the crossing.

\begin{proposition}[Low/high determinant transfer at a crossing]
\label{prop:v15v33-crossing}
Away from \(s_0\), factor the positive determinant as
\[
 \det{}'B(s)
 =
 \Delta_{\rm low}(s)\Delta_{\rm high}(s).                  \tag{33.23}
\]
When \(\lambda(s)\) enters the low window,
\(\Delta_{\rm low}\) gains the factor \(\lambda(s)\) and
\(\Delta_{\rm high}\) loses the same factor; when it exits, the
transfer is reversed.  The product (33.23) is unchanged, but neither
factor defines a smooth fixed-rank bundle through \(s_0\) without
this transfer record.
\end{proposition}

\begin{proof}
Diagonalize \(B(s)\) at points away from the crossing.  The
pseudo-determinant is the product of all positive eigenvalues.
The two factors in (33.23) partition that product according to
(33.22).  Moving one eigenvalue across the partition gives exactly
the stated reciprocal transfer.  Rank discontinuity precludes Kato
transport of a fixed-rank range through the crossing.
\end{proof}

If \(\lambda(s)\) crosses zero, it is no longer merely a low/high
transfer: Proposition \ref{prop:v15v32-zero} requires the kernel
determinant line, divisor, phase, and orientation ledger.

\begin{corollary}[Uniform bundle or crossing alternative]
\label{cor:v15v33-alternative}
Along any compact field path and fixed window edge \(E\), either:
\begin{enumerate}
\item a contour separation \(\delta>0\) holds, and
      Theorems \ref{thm:v15v33-projection}--\ref{thm:v15v33-connection}
      give a smooth quasilocal low bundle and determinant response;
\item the separation tends to zero, in which case a sequence of
      eigenvalues approaches \(E^2\) and must be entered in the
      low/high crossing ledger.
\end{enumerate}
\end{corollary}

\begin{proof}
The continuous spectral distance from a compact contour is either
bounded below or has a sequence tending to zero.  In the latter
case, compactness of the contour produces spectral points
approaching it; because the contour separates at \(E^2\), these are
window-edge crossings or accumulations.
\end{proof}

\begin{assessment}[Fixed and growing windows]
\label{ass:v15v33-windows}
For a fixed physical window on a compact continuum geometry,
norm-resolvent convergence and a continuum spectral gap at \(E\)
can populate a uniform contour separation.  V33 then supplies
quasilocal projection transport.

For growing windows, spectral spacings may shrink with energy and
volume.  V33 gives no uniform \(\delta\) in that regime.  A sharp
cutoff may need to be replaced by a smooth spectral partition, or
crossing transfers must be controlled together with the uniform
trace-class tail (32.26).  Fixed-window locality cannot be promoted
to the growing window without those estimates.
\end{assessment}

\begin{assessment}[Literal source applicability]
\label{ass:v15v33-source}
The source's kernel/divisor/orientation transport is structurally
compatible with the Kato and crossing ledgers above.  It does not
supply a norm-resolvent limit, a uniform energy-window contour gap,
or the trace-class bounds needed to control the growing collection
of crossings.  The low physical determinant therefore remains a
conditional construction.
\end{assessment}

\begin{decision}[V34 smooth spectral partition]
\label{dec:v15v33-next}
Replace the sharp window projection by a smooth function
\(\chi_E(B)\) with a transition band.  Use a resolvent or
Helffer--Sjoestrand formula to prove quasilocality without a spectral
gap at the artificial window edge.  Define complementary low/high
relative determinant contributions with an overlap correction, and
identify the trace norm required as the transition band moves to
higher energy.  Exact zero modes remain excluded from smoothing.
\end{decision}

\begin{firewall}[Claim boundary after V33]
\label{fire:v15v33}
A uniformly separated fixed low-energy window now has a quasilocal
Riesz projection, local Kato connection, and basis-free determinant
response.  Closing contour gaps require explicit low/high crossing
records; growing-window control remains unproved.
\end{firewall}

\section*{Version V33 append-only record}
\addcontentsline{toc}{section}{Version V33 append-only record}
The complete V32 source was retained before this continuation.  It
had \(316123\) bytes, \(9131\) lines, and SHA--256
\[
 \texttt{5F5DB8A4980342343B1C6E9CA4BEC3F36396DA24637B94F55B2314514B4E34AC}.
\]
V33 proved quasilocality and local variation of a separated
low-energy Riesz bundle, cancellation of the Kato commutator in the
scalar determinant, and the exact low/high transfer rule at window
crossings.
\n+\chapter{Fifteenth-cycle continuation V34: smooth spectral partitions without an artificial edge gap}
\label{sec:v15v34}

\section{Purpose and correction of the sharp-window route}

V33 isolated the exact obstruction created by a sharp spectral
window: when an eigenvalue crosses the artificial edge, the Riesz
projection changes rank and its contour gap closes.  The present
continuation replaces that projection by a smooth spectral
multiplier.  This removes the rank jump at the window edge.  It does
not create a spectral gap, a trace-class ultraviolet tail, or a rule
for exact zero modes.

There is a necessary correction to the tentative wording of
Decision \ref{dec:v15v33-next}.  For a general
\(C^\infty_c\) cutoff, functional calculus gives decay faster than
any prescribed power, but it need not give exponential decay with a
fixed exponent.  Exponential decay would require additional
analyticity.  Rapid polynomial decay is sufficient for the lattice
assembly whenever its order is chosen larger than the metric growth
dimension.

\section{A weighted propagation estimate}

Let \((\mathfrak X,d)\) be a countable metric set and
\[
 {cal H}=\bigoplus_{x\in\mathfrak X}{\cal H}_x,
 \qquad P_x:{\cal H}\longrightarrow{cal H}_x.             \tag{34.1}
\]
For an operator \(A\), set
\[
 \|A\|_\alpha=
 \max\left\{
 \sup_x\sum_y e^{\alpha d(x,y)}\|P_xAP_y\|,
 \sup_y\sum_x e^{\alpha d(x,y)}\|P_xAP_y\|
 \right\}.                                                \tag{34.2}
\]
The Schur test makes this a Banach-algebra norm.  Let
\(B=B^*\) and suppose \(\|B\|_{\alpha_0}<\infty\) for some
\(\alpha_0>0\).  For \(0<\alpha\leq\alpha_0\), define the
conjugation debit
\[
 \varepsilon_\alpha=
 \max\left\{
 \sup_x\sum_y
   (e^{\alpha d(x,y)}-1)\|P_xBP_y\|,
 \sup_y\sum_x
   (e^{\alpha d(x,y)}-1)\|P_xBP_y\|
 \right\}.                                                \tag{34.3}
\]
It is finite and tends to zero with \(\alpha\downarrow0\) by
dominated convergence.

\begin{lemma}[Weighted unitary propagation]
\label{lem:v15v34-propagation}
For subsets \(X,Y\subset\mathfrak X\), with
\(D=d(X,Y)\),
\[
 \|P_Xe^{itB}P_Y\|
 \leq \exp\{-\alpha D+\varepsilon_\alpha|t|\}.             \tag{34.4}
\]
The bound is also valid with the right side replaced by its minimum
with one.
\end{lemma}

\begin{proof}
Choose the bounded, one-Lipschitz function
\[
 \phi(z)=\min\{d(z,Y),D\}
\]
and let \(M\) act on \({\cal H}_z\) as multiplication by
\(e^{\alpha\phi(z)}\).  The Schur test and
\(|\phi(x)-\phi(y)|\leq d(x,y)\) give
\[
 \|MBM^{-1}-B\|\leq\varepsilon_\alpha.                    \tag{34.5}
\]
Writing \(MBM^{-1}=B+C\), the Duhamel formula and Gronwall's
inequality imply
\[
 \|Me^{itB}M^{-1}\|
 =\|e^{it(B+C)}\|
 \leq e^{|t|\|C\|}
 \leq e^{\varepsilon_\alpha|t|}.                          \tag{34.6}
\]
On \(Y\), \(M=I\), while on \(X\),
\(M=e^{\alpha D}I\).  Compressing (34.6) therefore proves
(34.4).  Unitarity gives the additional bound by one.
\end{proof}

\section{Rapid locality of a smooth spectral multiplier}

For \(\chi\in C_c^\infty(\mathbb R)\), use the convention
\[
 \widehat\chi(t)=\int_{\mathbb R}e^{-it\lambda}
                 \chi(\lambda)\,\dd\lambda,
 \qquad
 \chi(B)=\frac1{2\pi}\int_{\mathbb R}
                 \widehat\chi(t)e^{itB}\,\dd t.           \tag{34.7}
\]

\begin{theorem}[Gap-free rapid locality]
\label{thm:v15v34-smooth-locality}
Under the hypotheses of Lemma
\ref{lem:v15v34-propagation}, for every integer \(N\geq0\) there
is a finite constant \(C_{N,\chi,B,\alpha}\) such that
\[
 \boxed{
 \|P_X\chi(B)P_Y\|
 \leq C_{N,\chi,B,\alpha}(1+d(X,Y))^{-N}.}                 \tag{34.8}
\]
No separation of the spectrum from the transition region of
\(\chi\) is assumed.
\end{theorem}

\begin{proof}
Put \(D=d(X,Y)\).  The assertion is immediate for bounded \(D\),
so suppose \(D>0\).  If \(\varepsilon_\alpha>0\), split (34.7) at
\[
 T_D=\frac{\alpha D}{2\varepsilon_\alpha}.                 \tag{34.9}
\]
For \(|t|\leq T_D\), (34.4) is at most
\(e^{-\alpha D/2}\).  For \(|t|>T_D\), use unitarity.  Hence
\[
 \|P_X\chi(B)P_Y\|
 \leq \frac{e^{-\alpha D/2}}{2\pi}\|\widehat\chi\|_1
 +\frac1{2\pi}\int_{|t|>T_D}|\widehat\chi(t)|\,\dd t.    \tag{34.10}
\]
The Fourier transform of a compactly supported smooth function is
Schwartz.  Repeated integration by parts therefore bounds the last
integral by \(C_N(1+T_D)^{-N}\), and the exponential term satisfies
the same type of bound.  If \(\varepsilon_\alpha=0\), (34.4) has no
time debit and the first term alone proves the result.
\end{proof}

\begin{corollary}[Summable smooth window]
\label{cor:v15v34-summable}
Suppose balls in \(\mathfrak X\) obey
\[
 \#\{y:d(x,y)\leq r\}\leq C_{\rm vol}(1+r)^d.              \tag{34.11}
\]
If \(N>d\), then the kernel in (34.8) is uniformly summable in
either endpoint.  Thus \(\chi(B)\) belongs to every polynomial
Schur class of order smaller than the chosen \(N-d\).
\end{corollary}

\begin{proof}
Decompose \(\mathfrak X\) into unit annuli about one endpoint and
use (34.11).  The resulting majorant is a constant times
\(\sum_{n\geq0}(1+n)^{d-1-N}\), which converges for \(N>d\).
The column estimate is identical.
\end{proof}

The theorem distinguishes two notions that must not be conflated.
Smoothness of \(\chi\) removes the artificial spectral-edge gap.
Analyticity of \(\chi\), absent for a nonzero compactly supported
cutoff, would be needed for a fixed exponential spatial rate by this
argument.

\section{Local variation of the smooth window}

Let \(s\mapsto B(s)=B(s)^*\) be differentiable.  Suppose a
variation is supported in a set \(S\):
\[
 B_s=P_SB_sP_S,
 \qquad \|B_s\|\leq K_s.                                  \tag{34.12}
\]
Finite enlargements of \(S\) only alter constants.

\begin{theorem}[Core-factorized derivative locality]
\label{thm:v15v34-derivative}
Assume the constants in (34.3) are uniform along the path.  For
every integer \(N\geq0\),
\[
 \boxed{
 \|P_X\partial_s\chi(B(s))P_Y\|
 \leq C'_{N,\chi,B,\alpha}K_s
 (1+d(X,S)+d(S,Y))^{-N}.}                                  \tag{34.13}
\]
The same conclusion holds for a sum of core-factorized variations,
after summing their right sides.
\end{theorem}

\begin{proof}
For \(t\geq0\), Duhamel's formula gives
\[
 \partial_s e^{itB}
 =i\int_0^t e^{i(t-r)B}B_se^{irB}\,\dd r,                 \tag{34.14}
\]
with the corresponding oriented integral for \(t<0\).  Insert
\(P_S\) on both sides of \(B_s\).  Two applications of (34.4)
give, with
\(D_S=d(X,S)+d(S,Y)\),
\[
 \|P_X\partial_s e^{itB}P_Y\|
 \leq |t|K_s e^{-\alpha D_S+\varepsilon_\alpha|t|}.        \tag{34.15}
\]
The elementary operator-norm bound is \(|t|K_s\).  Differentiate
(34.7), justified by
\(\int |t\widehat\chi(t)|\,\dd t<\infty\), and split the
integral at
\(\alpha D_S/(2\varepsilon_\alpha)\).  The short-time part is an
exponential times a polynomial in \(D_S\); the long-time part is a
Schwartz tail of \(|t\widehat\chi(t)|\).  Both are bounded by the
right side of (34.13), after changing the constant.  The case
\(\varepsilon_\alpha=0\) and the sum statement follow as in
Theorem \ref{thm:v15v34-smooth-locality}.
\end{proof}

For \(B=T^*T\), formula (33.12) and the normalized local variation
bounds from V14--V19 supply the core terms in (34.12).  Unlike the
Riesz derivative (33.10), (34.13) has no inverse power of a contour
gap at the artificial energy edge.

\section{An exact smooth low/high determinant identity}

Work first in finite volume.  Let \(B\geq0\), let
\(P_0={\bf1}_{\{0\}}(B)\), and put \(Q=I-P_0\).  Assume on the
stratum under consideration that
\[
 \sigma(B|_{Q\mathcal H})\subset[m^2,M^2],
 \qquad 0<m<M<\infty.                                     \tag{34.16}
\]
The positive logarithm
\(\log_QB\) is then defined on \(Q\mathcal H\).  Choose smooth
real functions \(\eta_L,\eta_H,\eta_O\) on a neighborhood of
\([m^2,M^2]\) satisfying
\[
 \eta_L+\eta_H-\eta_O=1.                                  \tag{34.17}
\]
Here \(\eta_L\) is a low-energy shoulder, \(\eta_H\) is a
high-energy shoulder, and \(\eta_O\), supported where the shoulders
overlap, prevents double counting.  The complementary choice
\(\eta_H=1-\eta_L\), \(\eta_O=0\) is allowed.

Define
\[
 \begin{split}
 F_L(B)&=\operatorname{Tr}_Q
   [\eta_L(B)\log_QB],\\
 F_H(B)&=\operatorname{Tr}_Q
   [\eta_H(B)\log_QB],\\
 F_O(B)&=\operatorname{Tr}_Q
   [\eta_O(B)\log_QB].                                    \tag{34.18}
 \end{split}
\]

\begin{theorem}[Overlap-corrected smooth factorization]
\label{thm:v15v34-factorization}
The pseudo-determinant has the exact factorization
\[
 \boxed{
 \det{}'B
 =\exp F_L(B)\,\exp F_H(B)\,\exp[-F_O(B)].}               \tag{34.19}
\]
Along a differentiable path of constant kernel dimension,
\[
 \frac{\dd}{\dd s}F_j(B(s))
 =\operatorname{Tr}_Q[f_j'(B)B_s],
 \quad f_j(\lambda)=\eta_j(\lambda)\log\lambda,           \tag{34.20}
\]
and consequently
\[
 \frac{\dd}{\dd s}(F_L+F_H-F_O)
 =\operatorname{Tr}_Q(B^{-1}B_s).                          \tag{34.21}
\]
Thus a spectral crossing inside a smooth transition band produces
no rank jump and no Kato transport term in the scalar determinant.
\end{theorem}

\begin{proof}
By (34.17), functional calculus gives
\[
 [\eta_L(B)+\eta_H(B)-\eta_O(B)]\log_QB=\log_QB.           \tag{34.22}
\]
Taking the trace and exponentiating proves (34.19).  To prove
(34.20), diagonalize at one parameter value and use the standard
divided-difference formula for the Fr\'echet derivative of
\(f_j(B)\).  Only diagonal matrix elements survive the trace, so
the divided difference reduces to \(f_j'\) on each eigenspace.
This proves (34.20), including at eigenvalue multiplicities.  The
derivative of
\(f_L+f_H-f_O=\log\) is \(1/\lambda\), giving (34.21).
Constant kernel dimension permits Kato identification of the
positive subspaces.  As in (33.21), its commutator has zero trace.
\end{proof}

Exact zero modes were removed before (34.18).  A crossing through
zero violates (34.16) and must still be handled by the determinant
line, divisor, phase, and orientation data of V32.  Smoothing an
artificial positive-energy boundary does not smooth that physical
singularity.

\section{Relative determinants and the moving transition band}

Let \(B_a\) and \(B_\star\) be regulator and reference operators,
with their exact kernels removed and their positive spaces
identified by \(J_a\).  For
\(f_{j,E}(\lambda)=\eta_{j,E}(\lambda)\log\lambda\), define the
relative pieces whenever the difference is trace class:
\[
 {cal F}_{j,E}(a)=
 \operatorname{Tr}\!left[
 f_{j,E}(B_a)-J_af_{j,E}(B_\star)J_a^*
 \right].                                                 \tag{34.23}
\]
The pointwise identity (34.17) implies the exact finite-regulator
relation
\[
 {cal F}_{L,E}(a)+{cal F}_{H,E}(a)-{cal F}_{O,E}(a)
 =\operatorname{Tr}\!left[
 \log B_a-J_a\log B_\star J_a^*
 \right].                                                 \tag{34.24}
\]

Let \(\tau_E\) be a smooth function equal to one on the union of
the transition supports and zero away from a slightly larger band.
A sufficient uniform condition for moving the smooth boundary to
higher energy is
\[
 \boxed{
 \lim_{E\to\infty}\sup_{0<a<a_0}
 \left\|
 \tau_E(B_a)\log B_a
 -J_a\tau_E(B_\star)\log B_\star J_a^*
 \right\|_1=0.}                                          \tag{34.25}
\]
Together with fixed-window trace-norm convergence, (34.25) makes the
overlap correction and the dependence on the precise shoulder tend
to zero uniformly.  It is the smooth counterpart of the
growing-window tail (32.26).

\begin{proposition}[Cutoff independence from a transition tail]
\label{prop:v15v34-cutoff-independence}
Let two smooth triples satisfying (34.17) agree outside the band
where \(\tau_E=1\), and suppose all their multipliers are uniformly
bounded.  Under (34.25), the difference between their relative low
pieces, after adding the corresponding overlap allocation, tends to
zero uniformly in \(a\) as \(E\to\infty\).
\end{proposition}

\begin{proof}
The difference multiplier \(h_E\) is supported where
\(\tau_E=1\) and has uniformly bounded supremum norm.  Since all
functions involved are functions of the same operator,
\[
 h_E(B)\log B=h_E(B)\tau_E(B)\log B.                       \tag{34.26}
\]
After subtracting the reference expression, insert and subtract
\(h_E(B_a)J_a\tau_E(B_\star)\log B_\star J_a^*\).  The first
term is bounded in trace norm by the left side of (34.25) times
\(\|h_E\|_\infty\).  The second is precisely the fixed-window
functional-calculus comparison required in (34.23); by hypothesis
it tends to zero.  Taking traces and using
\(|\operatorname{Tr}A|\leq\|A\|_1\) proves the claim.
\end{proof}

The second term in this proof is why a transition-band estimate
alone does not replace fixed-window norm-resolvent or trace-norm
convergence.  Both inputs are required.

\begin{assessment}[What smoothing closes]
\label{ass:v15v34-closes}
The smooth window closes the artificial crossing problem of V33:
there is no changing rank at a positive transition edge, and local
variations have summable rapid decay with no edge-gap denominator.
The exact factorization (34.19) also records overlap without double
counting.

It does not prove (34.25).  Nor does it control a vanishing positive
eigenvalue, identify the Standard Model continuum Dirac operator, or
construct the Einstein-coupled relative determinant.  These remain
upstream analytic inputs.
\end{assessment}

\begin{assessment}[Literal source applicability]
\label{ass:v15v34-source}
The source supplies a finite-cutoff operator and a kernel ledger but
does not state a uniform trace-norm transition-tail estimate of the
form (34.25).  Its qualitative ultraviolet language cannot be used
as that estimate.  V34 therefore upgrades the low/high
decomposition but leaves its growing-window continuum passage
conditional.
\end{assessment}

\begin{decision}[V35 five-version audit]
\label{dec:v15v34-next}
At V35 inspect the newest Yang--Mills and Navier--Stokes companion
notes.  Then audit V31--V34 against the literal Standard
Model--Einstein source and identify the smallest upstream estimate
that can imply both fixed-window trace convergence and the moving
transition tail (34.25).  The audit must be followed by a
substantive V36 continuation.
\end{decision}

\begin{firewall}[Claim boundary after V34]
\label{fire:v15v34}
Smooth positive-energy windows now have gap-free, summable
quasilocality and an exact overlap-corrected determinant identity.
Uniform transition-tail control, exact-zero transport, and the
coupled continuum limit remain unproved.
\end{firewall}

\section*{Version V34 append-only record}
\addcontentsline{toc}{section}{Version V34 append-only record}
The complete V33 source was retained before this continuation.  It
had \(327270\) bytes, \(9457\) lines, and SHA--256
\[
 \texttt{6D1C08248B9E77CB32E2976DE859C49B191DAAAC89EB963158C111F9668B9EBC}.
\]
V34 proved gap-free rapid locality and local variation for smooth
spectral multipliers, gave an exact overlap-corrected low/high
pseudo-determinant factorization, and isolated the uniform
transition-band trace tail required for a growing smooth window.

\chapter{Fifteenth-cycle continuation V35: compulsory companion and source audit}
\label{sec:v15v35}

\section{Checkpoint and audit rule}

This is the fifth version since the fifteenth-cycle restart.  In
accordance with the standing protocol, the newest Yang--Mills and
Navier--Stokes notes in the shared folder were inspected before
choosing the next Standard Model--Einstein step.  The review is
read-only and transfers no theorem across programmes merely by
analogy.

The exact companion checkpoints are
\[
\begin{array}{c|c|c|c}
\text{programme}&\text{file}&\text{lines}&\text{SHA--256}\\ \hline
\mathrm{YM}&
\texttt{Renewal\_Yang\_Mills\_Mass\_Gap\_Research\_Notes\_V38.tex}&
8835&
\texttt{E0D1E18FF192D0F8BE7437B451DDA05ECA569ABDEF8ED085C3E032A8506EF33F}\\
\mathrm{NS}&
\texttt{Renewal\_NS\_Stability\_Research\_Notes\_V26.tex}&
5319&
\texttt{82C887F8C2C69E4F28FAD7C3DC8DE5B9099CB5A4BF431EBA1FFF3E91935978AD}
\end{array}                                                \tag{35.1}
\]

The originally named physical-source V30 file is not present under
that exact name in the current attachment directory.  The grand
tensor manuscript V067 remains present.  For a conservative
same-series check, the available evolved physical-source V35 was
also inspected, without treating a later statement as if it had
occurred in the missing V30.  Their digests are
\[
\begin{split}
 \operatorname{SHA256}(\mathrm{GrandTensor\ V067})
 &=\texttt{FACD058EC8BB210AD8B296C9D7F4D78A229D08D7A2AE378808700655CB8BC0FE},\\
 \operatorname{SHA256}(\mathrm{PhysicalSource\ V35})
 &=\texttt{AE1FBC0A8125C37465E38E7649E2992725F4579A86F0BB17483B9AE829CFB17C}.
                                                               \tag{35.2}
\end{split}
\]

\section{Yang--Mills finding: cover the whole band}

The current YM continuation constructs eight independent rational
charts and proves pointwise positivity on
\[
 [0,\pi/4]\times[0,2t]\times[0,t].                       \tag{35.3}
\]
Its proof explicitly treats one smaller-radius exceptional chart
and the transcendental cap between the last rational centre and
\(\pi/4\).  It does not infer coverage of the full cell from
positivity at a grid of centres.

No YM positivity constant is an SM determinant estimate.  The
transferable proof rule is narrower:

\begin{principle}[Transition-band coverage discipline]
\label{prin:v15v35-band-cover}
A trace-tail certificate for a smooth energy transition must cover
every energy in the transition band.  A finite family of local
estimates is sufficient only after its neighborhoods cover the
band pointwise, and any smaller-radius or near-divisor subband must
be treated explicitly.
\end{principle}

This rule prevents a finite list of sampled energy values from being
mistaken for the uniform supremum in (34.25).

\section{Navier--Stokes finding: nuclear size is separate}

The current NS continuation computes closed-form pointwise rank-one
norms for first and second kernel derivatives and postpones the
certified integral of those norms to its next version.  Its result
does not concern fermion determinants.  It nevertheless identifies
the relevant functional-analytic distinction: pointwise operator
bounds become an integrated nuclear estimate only after the
rank-one magnitudes are summable.

\begin{principle}[Nuclear summation discipline]
\label{prin:v15v35-nuclear}
If an operator is represented as
\[
 K=\sum_{\nu}u_\nu\otimes v_\nu^*,
 \qquad
 \sum_\nu\|u_\nu\|\,\|v_\nu\|<\infty,                    \tag{35.4}
\]
then
\[
 \|K\|_1\leq\sum_\nu\|u_\nu\|\,\|v_\nu\|.                \tag{35.5}
\]
Operator-norm or off-diagonal locality estimates without the
summation in (35.4) do not imply a trace-norm bound.
\end{principle}

\begin{proof}
Each rank-one term has trace norm
\(\|u_\nu\|\,\|v_\nu\|\).  The triangle inequality proves (35.5)
for partial sums, and completeness of the trace ideal proves the
claim.  The projections
\(K_n=|e_n\rangle\langle e_n|\) have norm one and perfect
locality but \(\|\sum_{n=1}^NK_n\|_1=N\), proving the last
statement.
\end{proof}

\section{Literal source audit}

The grand tensor manuscript contains three statements directly
relevant to V31--V34:
\begin{enumerate}
\item ordinary norm-resolvent convergence on a selected finite
      carrier does not imply convergence in a weighted Schur norm;
\item compact resolvent and trace-class heat are named analytic
      rows, not consequences of the renewal grammar;
\item the determinant/Pfaffian line, ultraviolet row, reflection
      row, and noncollapse row remain independent gates in the
      conditional OS handoff.
\end{enumerate}
Its weighted-conjugation and Duhamel results are compatible with
the locality proof in V34, but they do not populate the
trace-norm limit (34.25).

The evolved physical-source V35 likewise treats compact resolvent
and trace-class heat as imported hypotheses for a diagnostic
Dirichlet model.  Its final firewall says that the exact populated
three-level thermal control is not the active
gauge--Higgs--Dirac--Yukawa history, and its acquisition request
still asks for a selected physical RPESM energy fibre.  Consequently
none of its finite five-matrix data can be inserted into the
continuum determinant.

\begin{theorem}[V35 nonimport and direction decision]
\label{thm:v15v35-audit-decision}
The companion and source review supplies no estimate that proves
fixed-window relative determinant convergence or the moving-band
tail (34.25).  The next non-circular target is a common-carrier
trace-ideal criterion with all three of the following features:
\begin{enumerate}
\item a Bochner-integrable trace-norm resolvent difference;
\item convergence under one regulator-independent integrable
      majorant;
\item an energy partition defined on one common comparison
      operator, so trace-class compactness implies a uniform
      high-energy tail.
\end{enumerate}
\end{theorem}

\begin{proof}
The YM note proves a different matrix inequality and contributes
only Principle \ref{prin:v15v35-band-cover}.  The NS note proves
rank-one kernel formulas for a different PDE and contributes only
Principle \ref{prin:v15v35-nuclear}.  The literal source statements
listed above leave the determinant and ultraviolet rows independent.
Thus importing a numerical or positivity conclusion would be
invalid.  The three requested features directly target the two
unproved trace statements and do not assume a contour gap at the
smooth window edge.
\end{proof}

\section{Audit of V31--V34}

\begin{assessment}[No duplication and no hidden promotion]
\label{ass:v15v35-cycle-audit}
V31 proved boundary-area scaling for a connected two-core trace.
V32 formulated basis-free fixed-window relative determinants and
exhibited the escape counterexample.  V33 handled sharp separated
windows and their crossing ledger.  V34 removed the artificial edge
gap with a smooth multiplier and isolated the transition tail.

None of those versions proves trace-class convergence of a
renormalized relative logarithm.  Conversely, the existing source
and companion results do not already contain the common-carrier
trace-ideal compiler selected in Theorem
\ref{thm:v15v35-audit-decision}.  V36 can therefore address a new
upstream bottleneck without repeating the previous notes.
\end{assessment}

\begin{decision}[V36 common-carrier trace-ideal compiler]
\label{dec:v15v35-next}
Let \(B_a,C_a\) be the positive regulator and reference operators
after exact kernels and predeclared local counterterms are removed.
Use the integral representation of the logarithm to give explicit
Schatten-one hypotheses on
\[
 (B_a+t)^{-1}-(C_a+t)^{-1}.                              \tag{35.6}
\]
Prove trace-norm convergence of the renormalized relative
logarithm.  Then prove that a smooth partition made with one common
comparison Hamiltonian has an automatic uniform moving-window tail.
Record separately why heat-trace finiteness at each cutoff is not a
uniform version of those hypotheses.
\end{decision}

\begin{firewall}[Claim boundary after V35]
\label{fire:v15v35}
V35 is an audit and a direction theorem.  It imports no YM or NS
mathematical conclusion into the SM measure and proves no new
continuum determinant.  The trace-ideal hypotheses selected for V36
remain to be populated on the physical regulator.
\end{firewall}

\section*{Version V35 append-only record}
\addcontentsline{toc}{section}{Version V35 append-only record}
The complete V34 source was retained before this continuation.  It
had \(342352\) bytes, \(9870\) lines, and SHA--256
\[
 \texttt{D321D8CEA3627AEB5A0C6AB66AEB484CC2E1BE3C85F74EB422FA2DB93B98D865}.
\]
V35 performed the compulsory companion review, checked the literal
source boundary, and selected a common-carrier resolvent
trace-ideal estimate as the smallest non-circular V36 target.

\chapter{Fifteenth-cycle continuation V36: a common-carrier trace-ideal compiler}
\label{sec:v15v36}

\section{Objective}

V34 formulated separate fixed-window and moving-transition
requirements because its smooth multipliers were attached separately
to the regulator and reference operators.  V35 identified the
upstream trace-ideal input missing from that formulation.  This
continuation places the renormalized relative logarithm on one common
Hilbert carrier and partitions that trace-class operator with one
comparison Hamiltonian.  In this formulation, trace-norm convergence
implies both fixed-window convergence and a uniform moving-window
tail.

The construction concerns the positive determinant modulus.  For a
chiral fermion operator \(T\), it is applied to \(B=T^*T\), with a
factor \(1/2\) in the logarithm.  It does not orient the complex
determinant line or the real Pfaffian line.

\section{Renormalized relative logarithms from resolvents}

Let \({\cal H}\) be a separable Hilbert space.  For every regulator
index \(n\), let \(B_n,C_n\) be positive self-adjoint operators on
\({\cal H}\) satisfying
\[
 B_n\geq m^2I,\qquad C_n\geq m^2I                       \tag{36.1}
\]
for one \(m>0\).  Exact kernels have therefore already been removed
and transported as determinant-line data.  Let
\({\cal E}_n(t)\in{\mathfrak S}_1({\cal H})\) be a predeclared
resolvent counterterm and define
\[
 {\cal R}_n(t)=
 (C_n+t)^{-1}-(B_n+t)^{-1}-{\cal E}_n(t),
 \qquad t>0.                                             \tag{36.2}
\]
The counterterm bank is fixed before the field is integrated; it is
not selected after inspecting the determinant.

\begin{theorem}[Dominated resolvent compiler]
\label{thm:v15v36-resolvent-compiler}
Suppose:
\begin{enumerate}
\item \(t\mapsto{\cal R}_n(t)\) is strongly measurable with values
      in \({\mathfrak S}_1\);
\item there is \(g\in L^1(0,\infty)\), independent of \(n\), such
      that
      \[
       \|{\cal R}_n(t)\|_1\leq g(t)                       \tag{36.3}
      \]
      for almost every \(t\);
\item for almost every \(t>0\), there is
      \({\cal R}(t)\in{\mathfrak S}_1\) with
      \[
       \|{\cal R}_n(t)-{\cal R}(t)\|_1\longrightarrow0.   \tag{36.4}
      \]
\end{enumerate}
Then the Bochner integrals
\[
 K_n=\int_0^\infty{\cal R}_n(t)\,\dd t,\qquad
 K=\int_0^\infty{\cal R}(t)\,\dd t                       \tag{36.5}
\]
belong to \({\mathfrak S}_1\) and
\[
 \boxed{\|K_n-K\|_1\longrightarrow0.}                    \tag{36.6}
\]
In finite volume, when \({\cal E}_n=0\),
\[
 K_n=\log B_n-\log C_n.                                  \tag{36.7}
\]
For nonzero counterterms, (36.5) is the renormalized relative
logarithm defined by the chosen subtraction scheme.
\end{theorem}

\begin{proof}
Assumptions (36.3) and strong measurability give Bochner
integrability.  Fatou's lemma applied to (36.4) gives
\(\|{\cal R}(t)\|_1\leq g(t)\) almost everywhere.  Moreover
\[
 \|{\cal R}_n(t)-{\cal R}(t)\|_1\leq2g(t).
\]
The scalar dominated convergence theorem therefore yields
\[
 \|K_n-K\|_1
 \leq\int_0^\infty
       \|{\cal R}_n(t)-{\cal R}(t)\|_1\,\dd t
 \longrightarrow0.
\]
For positive numbers \(b,c\),
\[
 \log b-\log c
 =\int_0^\infty\left(\frac1{c+t}-\frac1{b+t}\right)\dd t.
                                                               \tag{36.8}
\]
Applying spectral calculus separately to \(B_n\) and \(C_n\),
first on bounded spectral intervals and then by monotone exhaustion,
proves (36.7).
\end{proof}

\begin{corollary}[Relative determinant convergence]
\label{cor:v15v36-det-convergence}
Under Theorem \ref{thm:v15v36-resolvent-compiler},
\[
 {\mathfrak D}_n=\exp(\operatorname{Tr}K_n)
 \longrightarrow
 {\mathfrak D}=\exp(\operatorname{Tr}K).                  \tag{36.9}
\]
For \(B_n=T_n^*T_n\), the positive fermion modulus uses
\(\exp(\tfrac12\operatorname{Tr}K_n)\).
\end{corollary}

\begin{proof}
The trace is continuous on \({\mathfrak S}_1\), so
\[
 |\operatorname{Tr}(K_n-K)|\leq\|K_n-K\|_1\to0.
\]
Continuity of the exponential proves the result.  The square-root
factor follows from
\(\log|\det T|=\tfrac12\log\det(T^*T)\) off the divisor.
\end{proof}

\section{A nuclear sufficient condition}

The resolvent criterion becomes checkable when the perturbation has
a summable factorization.  Suppose the resolvent identity holds in
the form
\[
 (C_n+t)^{-1}-(B_n+t)^{-1}
 =(C_n+t)^{-1}W_n(B_n+t)^{-1},\qquad W_n=B_n-C_n.       \tag{36.10}
\]

\begin{proposition}[Rank-one resolvent bridge]
\label{prop:v15v36-nuclear-bridge}
Assume
\[
 W_n=\sum_{\nu\in{\cal I}_n}u_{\nu,n}\otimes v_{\nu,n}^*
                                                               \tag{36.11}
\]
in a sense for which (36.10) is valid, and put
\[
 q_n(t)=
 \sum_{\nu\in{\cal I}_n}
 \|(C_n+t)^{-1}u_{\nu,n}\|\,
 \|(B_n+t)^{-1}v_{\nu,n}\|.                              \tag{36.12}
\]
If \(q_n(t)<\infty\), then the unsubtracted resolvent difference is
trace class and
\[
 \boxed{
 \|(C_n+t)^{-1}-(B_n+t)^{-1}\|_1\leq q_n(t).}            \tag{36.13}
\]
If \(q_n(t)+\|{\cal E}_n(t)\|_1\leq g(t)\) for one
\(g\in L^1(0,\infty)\), and the resolvent differences and
counterterms converge pointwise in trace norm, then Theorem
\ref{thm:v15v36-resolvent-compiler} applies.
\end{proposition}

\begin{proof}
Substitution of (36.11) into (36.10) gives a sum of rank-one
operators
\[
 \sum_\nu
 \bigl((C_n+t)^{-1}u_{\nu,n}\bigr)
 \otimes
 \bigl((B_n+t)^{-1}v_{\nu,n}\bigr)^*.
                                                               \tag{36.14}
\]
The trace norm of \(x\otimes y^*\) is \(\|x\|\,\|y\|\).
Principle \ref{prin:v15v35-nuclear} proves (36.13), and the final
claim follows from (36.2)--(36.4).
\end{proof}

There is a useful Hilbert--Schmidt variant.  If
\(W_n=X_nY_n^*\) and the two sandwiched factors are Hilbert--Schmidt,
then
\[
 \|(C_n+t)^{-1}X_nY_n^*(B_n+t)^{-1}\|_1
 \leq
 \|(C_n+t)^{-1}X_n\|_2
 \|(B_n+t)^{-1}Y_n\|_2.                                  \tag{36.15}
\]
This is often the natural route for a local perturbation of an
elliptic operator: the individual resolvents need not be trace
class, while their product around a local insertion can be.

\section{One common spectral partition}

Let \(H_{\rm cmp}\geq0\) be a fixed self-adjoint comparison
Hamiltonian on \({\cal H}\).  Choose
\(\chi\in C^\infty([0,\infty);[0,1])\) with
\[
 \chi(r)=1\quad(0\leq r\leq1),\qquad
 \chi(r)=0\quad(r\geq2),                                  \tag{36.16}
\]
and set
\[
 S_E=\chi(H_{\rm cmp}/E),\qquad Q_E=I-S_E.               \tag{36.17}
\]
Then \(S_E\to I\) and \(Q_E\to0\) strongly as \(E\to\infty\).

\begin{lemma}[Strong cutoffs are uniform on trace-compact sets]
\label{lem:v15v36-trace-compact-tail}
Let \(Q_E\) be uniformly bounded and converge strongly to zero.
If \({\cal K}\subset{\mathfrak S}_1({\cal H})\) is compact, then
\[
 \lim_{E\to\infty}\sup_{A\in{\cal K}}
 \left(\|Q_EA\|_1+\|AQ_E^*\|_1\right)=0.                 \tag{36.18}
\]
\end{lemma}

\begin{proof}
First fix \(A\in{\mathfrak S}_1\).  Choose a finite-rank \(F\) with
\(\|A-F\|_1<\epsilon\).  Strong convergence is uniform on the
finite-dimensional range of \(F\), so
\(\|Q_EF\|_1\to0\).  Uniform boundedness of \(Q_E\) gives
\[
 \|Q_E(A-F)\|_1\leq
 \left(\sup_E\|Q_E\|\right)\epsilon.                     \tag{36.19}
\]
Thus \(\|Q_EA\|_1\to0\); applying the same argument to \(A^*\)
proves the right-multiplication statement.  Cover the compact set
\({\cal K}\) by a finite trace-norm \(\epsilon\)-net and use the
same uniform bound.  Letting \(\epsilon\downarrow0\) proves (36.18).
\end{proof}

\begin{theorem}[Trace-norm convergence closes both window limits]
\label{thm:v15v36-window-compiler}
If \(K_n\to K\) in \({\mathfrak S}_1\), define
\[
 F_{L,E}^{(n)}=\operatorname{Tr}(S_EK_n),\qquad
 F_{H,E}^{(n)}=\operatorname{Tr}(Q_EK_n).                \tag{36.20}
\]
Then:
\begin{align}
 \operatorname{Tr}K_n
 &=F_{L,E}^{(n)}+F_{H,E}^{(n)}
 &&\text{for every \(n,E\)},                             \tag{36.21}\\
 F_{L,E}^{(n)}
 &\longrightarrow\operatorname{Tr}(S_EK)
 &&\text{for every fixed \(E\)},                         \tag{36.22}\\
 \lim_{E\to\infty}\sup_{n\geq1}|F_{H,E}^{(n)}|
 &=0.                                                     \tag{36.23}
\end{align}
More generally, if
\(S_{L,E}+S_{H,E}-S_{O,E}=I\) are uniformly bounded smooth
functions of \(H_{\rm cmp}\), then
\[
 \operatorname{Tr}K_n
 =\operatorname{Tr}(S_{L,E}K_n)
 +\operatorname{Tr}(S_{H,E}K_n)
 -\operatorname{Tr}(S_{O,E}K_n),                         \tag{36.24}
\]
and every piece supported above an energy tending to infinity
vanishes uniformly in \(n\).
\end{theorem}

\begin{proof}
Equation (36.21) is linearity of the trace and
\(S_E+Q_E=I\).  Since \(\|S_E\|\leq1\),
\[
 |\operatorname{Tr}(S_E(K_n-K))|
 \leq\|K_n-K\|_1,
\]
which proves (36.22).  The set
\[
 {\cal K}=\{K\}\cup\{K_n:n\geq1\}                        \tag{36.25}
\]
is compact in trace norm because it is a convergent sequence
together with its limit.  Lemma
\ref{lem:v15v36-trace-compact-tail} gives
\[
 \sup_n|F_{H,E}^{(n)}|
 \leq\sup_n\|Q_EK_n\|_1\longrightarrow0.
\]
The overlap identity follows in the same way.  A uniformly bounded
high-energy multiplier is dominated by a cutoff \(Q_E\) after
enlarging the transition band, so (36.18) proves the last claim.
\end{proof}

\begin{corollary}[The separate transition condition can be replaced]
\label{cor:v15v36-replaces-transition}
For a selected cofinal regulator sequence, the trace-norm
convergence (36.6) on one common carrier is sufficient for both the
fixed-window limit and the uniform transition-tail conclusion sought
in (34.25), provided the smooth partition is made with
\(H_{\rm cmp}\).  No spectral gap at the moving edge and no crossing
ledger are needed.
\end{corollary}

\begin{proof}
Apply Theorem \ref{thm:v15v36-window-compiler}.  Smooth functional
calculus for \(H_{\rm cmp}\), together with V34, gives summable
spatial locality when the comparison Hamiltonian has the required
weighted locality.  The trace statements themselves use only
trace-norm convergence and strong convergence of the cutoffs.
\end{proof}

The use of a common comparison operator matters.  If the cutoff on
the first term is \(\chi(B_n/E)\) and that on the second is
\(\chi(C_n/E)\), an additional difference of functional calculi
appears, as in the proof of Proposition
\ref{prop:v15v34-cutoff-independence}.  Equations (36.20)--(36.24)
avoid that extra ultraviolet row while retaining an exact split of
the relative trace.

\section{Heat-trace finiteness is not the uniform estimate}

For positive operators, the parallel representation
\[
 \log B-\log C
 =\int_0^\infty\frac{e^{-sC}-e^{-sB}}{s}\,\dd s           \tag{36.26}
\]
is valid under the corresponding trace-norm integrability at
\(s=0\) and \(s=\infty\).  Merely knowing that each heat operator is
trace class for each fixed regulator does not give such a uniform
integrable majorant.

\begin{countertheorem}[Cutoffwise heat trace does not control the determinant]
\label{cth:v15v36-heat-not-uniform}
On \({\cal H}_n=\mathbb C^n\), let
\[
 B_n=2I_n,\qquad C_n=I_n.                                \tag{36.27}
\]
For every \(s>0\), both heat operators are trace class, and both
operators have a uniform positive floor.  Nevertheless
\[
 \|\log B_n-\log C_n\|_1=n\log2\longrightarrow\infty,     \tag{36.28}
\]
and
\[
 \int_0^\infty
 \frac{\|e^{-sC_n}-e^{-sB_n}\|_1}{s}\,\dd s
 =n\log2.                                                \tag{36.29}
\]
Thus cutoffwise compact resolvent and heat-trace finiteness do not
populate (36.3).
\end{countertheorem}

\begin{proof}
All assertions follow by diagonalization.  The scalar Frullani
integral
\[
 \int_0^\infty\frac{e^{-s}-e^{-2s}}s\,\dd s=\log2
\]
gives (36.29).
\end{proof}

\section{Physical insertion point}

For the lattice Standard Model--Einstein regulator, the intended
common-carrier packet is:
\[
\begin{gathered}
 B_n=J_nT_n^*T_nJ_n^*,\qquad
 C_n=J_nT_{\star,n}^*T_{\star,n}J_n^*,\\
 {\cal R}_n(t)=
 (C_n+t)^{-1}-(B_n+t)^{-1}-{\cal E}_n(t),                \tag{36.30}\\
 \int_0^\infty\|{\cal R}_n(t)\|_1\,\dd t
 \leq\int_0^\infty g(t)\,\dd t<\infty.
\end{gathered}
\]
Here \(J_n\) is the declared identification with one physical
carrier, \(T_{\star,n}\) is the reference Dirac--Yukawa operator,
and \({\cal E}_n\) is built from the fixed local counterterm bank.
Exact kernels are excluded before (36.30).  The normalized
operator convention of V16 is retained, so no multiplication of the
log determinant by a vanishing cell volume is allowed.

\begin{assessment}[What V36 closes and what it moves upstream]
\label{ass:v15v36-closes}
V36 closes a logical gap in the window construction: trace-norm
convergence of one renormalized relative logarithm automatically
controls all fixed common windows and their uniform high-energy
tail.  It also gives two explicit sufficient routes to that
convergence: the dominated resolvent criterion and the nuclear
factorization (36.11)--(36.15).

The physical source does not supply the common carrier \(J_n\), the
counterterm resolvent \({\cal E}_n(t)\), or the uniform integrable
majorant \(g\).  Those data cannot be inferred from cutoffwise heat
traces.  The continuum determinant therefore remains conditional,
but its next bottleneck is now one concrete estimate rather than two
separate window limits.
\end{assessment}

\begin{decision}[V37 short-time heat and subtraction order]
\label{dec:v15v36-next}
Translate (36.3) into a heat-kernel criterion.  Prove an abstract
short-time subtraction theorem: if the relative heat trace has a
finite local asymptotic head and an integrable trace-norm remainder,
then its Mellin integral defines \(K_n\) and satisfies the V36
compiler.  Identify, in four dimensions, which powers and logarithms
must be assigned to predeclared cosmological, Einstein,
gauge--Higgs, and curvature counterterm channels.  Do not claim that
the source has populated their coefficients.
\end{decision}

\begin{firewall}[Claim boundary after V36]
\label{fire:v15v36}
The common-carrier trace-ideal compiler is proved.  No physical
uniform resolvent majorant, renormalized heat remainder, determinant
phase, interacting measure, or Einstein continuum equation has yet
been constructed.
\end{firewall}

\section*{Version V36 append-only record}
\addcontentsline{toc}{section}{Version V36 append-only record}
The complete V35 source was retained before this continuation.  It
had \(351146\) bytes, \(10083\) lines, and SHA--256
\[
 \texttt{4945A56020FA4DC89968889006B3DFC0A50DD446BD7B4CEDDE07FCA4011D2F25}.
\]
V36 proved trace-norm convergence from a dominated relative
resolvent integral, supplied nuclear and Hilbert--Schmidt sufficient
conditions, and proved that one common smooth comparison partition
makes the moving high-energy trace tail automatic.

\chapter{Fifteenth-cycle continuation V37: four-dimensional heat subtraction and the trace-norm remainder}
\label{sec:v15v37}

\section{Proper-time form of the bottleneck}

V36 reduced smooth-window control to convergence of one
renormalized relative logarithm in trace norm.  The present
continuation expresses a sufficient condition in proper time and
separates two logically different claims:
\begin{enumerate}
\item convergence of the scalar logarithmic determinant;
\item convergence of the relative logarithm as a trace-class
      operator on the common carrier.
\end{enumerate}
The first can follow from a bound on a relative heat trace.  The
second requires a trace-norm bound on the heat-operator remainder.
Cancellation inside a scalar trace cannot supply the second.

Let \(B_n,C_n\geq m^2I\) be the positive operators of V36 and write
\[
 {\cal H}_n(s)=e^{-sC_n}-e^{-sB_n},\qquad s>0.            \tag{37.1}
\]
For scalars \(b,c>0\), Frullani's formula is
\[
 \log b-\log c
 =\int_0^\infty\frac{e^{-sc}-e^{-sb}}s\,\dd s.           \tag{37.2}
\]

\section{Scalar subtraction theorem}

Fix \(\vartheta\in C_c^\infty([0,\infty))\) with
\(\vartheta(s)=1\) for \(0\leq s\leq1\).  In four dimensions the
short-time head of a second-order Laplace-type problem has the form
\[
 p_n(s)=c_{0,n}s^{-2}+c_{2,n}s^{-1}+c_{4,n}.             \tag{37.3}
\]
The subscripts name the usual local heat coefficients, rather than
powers of \(s\).

\begin{theorem}[Dominated scalar proper-time subtraction]
\label{thm:v15v37-scalar-heat}
Suppose \(h_n(s)=\operatorname{Tr}{\cal H}_n(s)\) is defined for
every \(s>0\) and
\[
 r_n(s)=h_n(s)-\vartheta(s)p_n(s).                       \tag{37.4}
\]
Assume:
\begin{enumerate}
\item \(r_n(s)\to r(s)\) for almost every \(s>0\);
\item there is \(g\in L^1(0,\infty)\) such that
      \[
       \frac{|r_n(s)|}{s}\leq g(s)                       \tag{37.5}
      \]
      for every \(n\) and almost every \(s\).
\end{enumerate}
Then
\[
 {\cal L}_n^{\rm ren}
 =\int_0^\infty\frac{r_n(s)}s\,\dd s
 \longrightarrow
 {\cal L}^{\rm ren}
 =\int_0^\infty\frac{r(s)}s\,\dd s.                      \tag{37.6}
\]
Consequently
\(\exp({\cal L}_n^{\rm ren})\to
  \exp({\cal L}^{\rm ren})\).
\end{theorem}

\begin{proof}
Equation (37.5) makes every integral absolutely convergent.  The
ordinary dominated convergence theorem proves (37.6), and
continuity of the exponential proves the last assertion.
\end{proof}

The cutoff \(\vartheta\) fixes a finite subtraction scheme.  Replacing
it by another such cutoff changes (37.6) by the finite local
quantity
\[
 -\int_0^\infty
 \frac{\vartheta_1(s)-\vartheta_2(s)}s\,p_n(s)\,\dd s,    \tag{37.7}
\]
so its coefficients must be assigned to the predeclared local
counterterm bank.  The scheme is not allowed to depend on the
observed determinant.

\section{Operator-valued subtraction theorem}

For the common-window theorem of V36, the scalar condition is not
enough.  Let
\({\cal P}_n(s)\in{\mathfrak S}_1({\cal H})\) be a
predeclared operator-valued local subtraction and set
\[
 {\cal A}_n(s)={\cal H}_n(s)-{\cal P}_n(s).              \tag{37.8}
\]

\begin{theorem}[Trace-norm heat compiler]
\label{thm:v15v37-operator-heat}
Suppose:
\begin{enumerate}
\item \(s\mapsto{\cal A}_n(s)\) is strongly measurable in
      \({\mathfrak S}_1\);
\item for almost every \(s>0\),
      \[
       \|{\cal A}_n(s)-{\cal A}(s)\|_1\longrightarrow0;   \tag{37.9}
      \]
\item one \(g\in L^1(0,\infty)\) satisfies
      \[
       \frac{\|{\cal A}_n(s)\|_1}{s}\leq g(s).            \tag{37.10}
      \]
\end{enumerate}
Then
\[
 K_n^{\rm heat}
 =\int_0^\infty\frac{{\cal A}_n(s)}s\,\dd s
 \longrightarrow
 K^{\rm heat}
 =\int_0^\infty\frac{{\cal A}(s)}s\,\dd s               \tag{37.11}
\]
in trace norm.  Hence every conclusion of Theorem
\ref{thm:v15v36-window-compiler} holds with
\(K_n=K_n^{\rm heat}\).
\end{theorem}

\begin{proof}
The proof is the Bochner dominated convergence argument of Theorem
\ref{thm:v15v36-resolvent-compiler}, with measure \(\dd s/s\).
Trace-norm convergence invokes Theorem
\ref{thm:v15v36-window-compiler}.
\end{proof}

\begin{proposition}[Resolvent and heat compilers agree]
\label{prop:v15v37-heat-resolvent}
Assume the heat remainders satisfy (37.10), their Laplace transforms
may be interchanged with the Bochner integral, and the resolvent
counterterm is
\[
 {\cal E}_n(t)=\int_0^\infty e^{-st}{\cal P}_n(s)\,\dd s.
                                                               \tag{37.12}
\]
Then
\[
 \int_0^\infty{\cal R}_n(t)\,\dd t
 =\int_0^\infty\frac{{\cal A}_n(s)}s\,\dd s.             \tag{37.13}
\]
\end{proposition}

\begin{proof}
The Laplace representation
\[
 (B+t)^{-1}=\int_0^\infty e^{-st}e^{-sB}\,\dd s          \tag{37.14}
\]
gives
\[
 {\cal R}_n(t)=
 \int_0^\infty e^{-st}{\cal A}_n(s)\,\dd s.
\]
Bochner--Fubini and
\(\int_0^\infty e^{-st}\,\dd t=s^{-1}\) prove (37.13).
\end{proof}

\section{Why scalar cancellation is insufficient}

\begin{countertheorem}[Zero heat trace can hide a nonzero trace ideal]
\label{cth:v15v37-trace-cancellation}
On \(\mathbb C^2\), set
\[
 B=\begin{pmatrix}1&0\\0&2\end{pmatrix},\qquad
 C=\begin{pmatrix}2&0\\0&1\end{pmatrix}.                 \tag{37.15}
\]
Then
\[
 \operatorname{Tr}(e^{-sC}-e^{-sB})=0
 \quad\text{for every }s>0,                              \tag{37.16}
\]
but
\[
 \|e^{-sC}-e^{-sB}\|_1=2(e^{-s}-e^{-2s})                 \tag{37.17}
\]
and
\[
 \|\log B-\log C\|_1=2\log2.                             \tag{37.18}
\]
Thus a perfect scalar heat-trace estimate does not imply (37.10)
or convergence of local operator responses.
\end{countertheorem}

\begin{proof}
All matrices are diagonal.  Their heat traces agree because their
eigenvalue multisets agree, while the two diagonal differences have
opposite signs and equal magnitudes.  Taking the sum of singular
values gives (37.17)--(37.18).
\end{proof}

\section{The four-dimensional local head}

Let \(P\) be a Laplace-type operator on a closed Riemannian
four-manifold, with convention
\[
 P=-\bigl(g^{\mu\nu}\nabla_\mu\nabla_\nu+E\bigr),\qquad
 \Omega_{\mu\nu}=[\nabla_\mu,\nabla_\nu].                \tag{37.19}
\]
Its heat trace has the formal local expansion
\[
 \operatorname{Tr}(e^{-sP})
 \sim(4\pi s)^{-2}
 \int_M\operatorname{tr}
 \bigl(a_0+s\,a_2+s^2a_4+\cdots\bigr)\,\dd{\rm vol}_g.
                                                               \tag{37.20}
\]
With the convention (37.19),
\[
 a_0=I,\qquad a_2=E+\frac16RI,                            \tag{37.21}
\]
and, up to total divergences on a closed manifold,
\[
 \begin{split}
 a_4=\frac1{360}\bigl(&60RE+180E^2+30\Omega_{\mu\nu}\Omega^{\mu\nu}\\
 &+(5R^2-2R_{\mu\nu}R^{\mu\nu}
       +2R_{\mu\nu\rho\sigma}R^{\mu\nu\rho\sigma})I
 \bigr).                                                 \tag{37.22}
 \end{split}
\]
The omitted local divergences are \(60\nabla^2E+12\nabla^2R\);
they become boundary terms when a boundary is present.

With a proper-time cutoff \(\epsilon=\Lambda^{-2}\), substitution
of (37.3) gives
\[
 \int_\epsilon^1\frac{p_n(s)}s\,\dd s
 =
 \frac{c_{0,n}}2(\Lambda^4-1)
 +c_{2,n}(\Lambda^2-1)
 +2c_{4,n}\log\Lambda.                                  \tag{37.23}
\]
Therefore a second-order four-dimensional determinant has exactly
the following local divergence channels:
\[
\begin{array}{c|c|c}
\text{heat coefficient}&\text{cutoff order}&
\text{local action channel}\\ \hline
a_0&\Lambda^4&\displaystyle\int\sqrt g
\quad\text{(cosmological)}\\
a_2&\Lambda^2&
\displaystyle\int\sqrt g\,R
\ \text{and dimension-two endomorphism terms}\\
a_4&\log\Lambda&
\displaystyle\int\sqrt g\,
(R^2,\ |\mathrm{Ric}|^2,\ |\mathrm{Riem}|^2,\
\operatorname{tr}\Omega^2,\ \operatorname{tr}E^2,\
R\operatorname{tr}E).
\end{array}                                              \tag{37.24}
\]

For a Dirac--Yukawa square, \(\Omega\) contains the spin and gauge
curvatures and \(E\) contains curvature and Higgs--Yukawa
endomorphisms.  Consequently the \(a_4\) row contains gauge kinetic,
Higgs kinetic/quartic, mixed curvature--Higgs, and curvature-squared
terms after the internal trace is evaluated.  The \(a_2\) row
contains the Einstein term and possible Higgs mass-type terms.  The
precise coefficients depend on the actual representation,
normalization, chirality, ghost content, and boundary conditions;
none may be copied from (37.22) without evaluating that physical
bundle.

\begin{theorem}[Four-dimensional subtraction order]
\label{thm:v15v37-subtraction-order}
For a closed four-dimensional Laplace-type relative determinant,
an integrable proper-time remainder at \(s=0\) requires subtraction
through \(a_4\), unless some of \(a_0,a_2,a_4\) cancel by a proved
relative identity.  After those terms are removed, a remainder
\[
 \|{\cal A}_n(s)\|_1\leq Cs^\delta
 \quad(0<s\leq1)                                        \tag{37.25}
\]
with any \(\delta>0\) is integrable against \(\dd s/s\).
For \(s\geq1\), a uniform bound
\[
 \|{\cal A}_n(s)\|_1\leq Ce^{-\gamma s}                 \tag{37.26}
\]
closes the large-time end.
\end{theorem}

\begin{proof}
Before subtraction, the three terms in (37.3), divided by \(s\),
have powers \(s^{-3},s^{-2},s^{-1}\), none integrable at zero.
After subtraction, (37.25) gives
\(\int_0^1Cs^{\delta-1}\dd s=C/\delta\).
Equation (37.26) gives
\(\int_1^\infty Ce^{-\gamma s}\dd s/s<\infty\).
\end{proof}

If a physical boundary is present, half-integer heat coefficients
can occur and (37.24) is incomplete.  Boundary conditions and their
local boundary counterterms must then be included before applying
Theorem \ref{thm:v15v37-subtraction-order}.

\section{Uniformity and local coefficient landing}

It is not enough that every fixed \(n\) admit an asymptotic
expansion.  The V36 compiler requires:
\[
\begin{gathered}
\sup_n\frac{\|{\cal H}_n(s)-{\cal P}_n(s)\|_1}{s}
\leq g(s),\qquad g\in L^1(0,\infty),                    \tag{37.27}\\
\|{\cal H}_n(s)-{\cal P}_n(s)-{\cal A}(s)\|_1\to0
\quad\text{for almost every }s.                         \tag{37.28}
\end{gathered}
\]
The coefficients entering \({\cal P}_n\) must converge in their
declared local action channels.  Principle
\ref{prin:v15v35-band-cover} requires estimates covering the whole
proper-time interval, with the short, intermediate, and long-time
regions joined explicitly.

\begin{assessment}[Literal source applicability]
\label{ass:v15v37-source}
The source's references to compact resolvent and trace-class heat at
each regulator do not state (37.27).  Nor do they populate the
physical \(a_0,a_2,a_4\) coefficients of the full
gauge--Higgs--Dirac--Yukawa bundle on the fluctuating metric.  The
counterterm channels in (37.24) are therefore a required acquisition
list, not a completed renormalization theorem.
\end{assessment}

\begin{decision}[V38 differentiable determinant and stress response]
\label{dec:v15v37-next}
Assume a parameter-dependent trace-norm heat remainder satisfying a
uniform integrable bound for its first variation.  Prove
differentiability of the renormalized relative logarithm under the
integral, derive the resolvent and Duhamel stress formulas, and state
a precise covariance hypothesis under which the resulting metric
variation obeys a Ward identity.  Keep local counterterm variations
and determinant-line phases explicit.
\end{decision}

\begin{firewall}[Claim boundary after V37]
\label{fire:v15v37}
V37 identifies the exact four-dimensional subtraction order and
proves scalar and operator-valued proper-time convergence theorems.
It does not establish the required uniform remainder for the
physical regulator or compute its Standard Model heat
coefficients.
\end{firewall}

\section*{Version V37 append-only record}
\addcontentsline{toc}{section}{Version V37 append-only record}
The complete V36 source was retained before this continuation.  It
had \(365908\) bytes, \(10501\) lines, and SHA--256
\[
 \texttt{0F21D22DC92381E0EC25F59E20C87BE1CF9DB188F7345C3A1FBE73834B5128B1}.
\]
V37 proved dominated scalar and trace-norm heat subtraction
theorems, exhibited the failure of scalar cancellation to control
the trace ideal, and identified the quartic, quadratic, and
logarithmic local counterterm rows of a four-dimensional
Laplace-type determinant.

\chapter{Fifteenth-cycle continuation V38: differentiable relative determinants and the covariant stress row}
\label{sec:v15v38}

\section{Why a first-variation theorem is separate}

The scalar and trace-class limits of V36--V37 control determinant
values.  The coupled Einstein equation requires derivatives with
respect to the metric, vierbein, gauge connection, and Higgs field.
Convergence of functions does not by itself imply convergence of
their derivatives.  This continuation states the additional
trace-norm hypotheses, derives the renormalized response formulas,
and identifies the covariance premise needed for a Ward identity.

Let \(u\) range in an open interval and let
\[
 {\cal A}_n(s,u)=e^{-sC_n(u)}-e^{-sB_n(u)}
                 -{\cal P}_n(s,u)                       \tag{38.1}
\]
be the subtracted heat remainder of V37 on one common carrier.

\section{Differentiation under the proper-time integral}

\begin{theorem}[Trace-norm \(C^1\) heat compiler]
\label{thm:v15v38-c1-heat}
Fix a compact parameter interval \(I\).  Suppose:
\begin{enumerate}
\item for almost every \(s>0\), the map
      \(u\mapsto{\cal A}_n(s,u)\) is \(C^1\) with values in
      \({\mathfrak S}_1\);
\item there are \(g_0,g_1\in L^1(0,\infty)\), independent of
      \(n\) and \(u\in I\), such that
      \[
       \frac{\|{\cal A}_n(s,u)\|_1}{s}\leq g_0(s),
       \qquad
       \frac{\|\partial_u{\cal A}_n(s,u)\|_1}{s}
       \leq g_1(s);                                     \tag{38.2}
      \]
\item for almost every \(s\), both
      \({\cal A}_n(s,u)\) and
      \(\partial_u{\cal A}_n(s,u)\) converge in trace norm,
      uniformly for \(u\in I\), to
      \({\cal A}(s,u)\) and \({\cal A}_u(s,u)\).
\end{enumerate}
Then
\[
 K_n(u)=\int_0^\infty\frac{{\cal A}_n(s,u)}s\,\dd s      \tag{38.3}
\]
belongs to \(C^1(I;{\mathfrak S}_1)\),
\[
 K_n'(u)=
 \int_0^\infty\frac{\partial_u{\cal A}_n(s,u)}s\,\dd s, \tag{38.4}
\]
and
\[
 \sup_{u\in I}\bigl(
 \|K_n(u)-K(u)\|_1+
 \|K_n'(u)-K'(u)\|_1\bigr)\longrightarrow0.             \tag{38.5}
\]
In particular,
\[
 \frac{\dd}{\dd u}\operatorname{Tr}K_n(u)
 =\operatorname{Tr}K_n'(u)
 \longrightarrow\operatorname{Tr}K'(u)                 \tag{38.6}
\]
uniformly on \(I\).
\end{theorem}

\begin{proof}
The Banach-valued fundamental theorem of calculus gives, for
\(u,v\in I\),
\[
 {\cal A}_n(s,u)-{\cal A}_n(s,v)
 =\int_v^u\partial_r{\cal A}_n(s,r)\,\dd r.              \tag{38.7}
\]
Divide by \(s\), integrate, and use (38.2) to justify
Bochner--Fubini.  This proves differentiability and (38.4).
Dominated convergence in the Banach space
\({\mathfrak S}_1\), followed by the assumed uniformity in \(u\),
proves (38.5).  Continuity of the trace proves (38.6).
\end{proof}

\begin{countertheorem}[Determinant convergence does not imply stress convergence]
\label{cth:v15v38-values-not-derivatives}
The functions
\[
 f_n(u)=\frac1n\sin(nu)                                  \tag{38.8}
\]
converge uniformly to zero, while
\[
 f_n'(u)=\cos(nu)                                        \tag{38.9}
\]
does not converge even pointwise for generic \(u\).  Thus convergence
of \(\operatorname{Tr}K_n(u)\), including locally uniform
convergence, cannot replace the derivative majorant in (38.2).
\end{countertheorem}

\section{Duhamel and resolvent response formulas}

Assume \(B_n(u)\) and \(C_n(u)\) have a common form domain and
differentiable perturbations for which the following integrals are
trace class.  Duhamel differentiation gives
\[
 \begin{split}
 \partial_u e^{-sB_n}
 &=-\int_0^s e^{-(s-r)B_n}B_{n,u}e^{-rB_n}\,\dd r,\\
 \partial_u e^{-sC_n}
 &=-\int_0^s e^{-(s-r)C_n}C_{n,u}e^{-rC_n}\,\dd r.
                                                               \tag{38.10}
 \end{split}
\]
Therefore
\[
 \begin{split}
 \partial_u{\cal A}_n(s,u)
={}&-\int_0^s e^{-(s-r)C_n}C_{n,u}e^{-rC_n}\,\dd r\\
 &+\int_0^s e^{-(s-r)B_n}B_{n,u}e^{-rB_n}\,\dd r
 -\partial_u{\cal P}_n(s,u).                            \tag{38.11}
 \end{split}
\]

In the resolvent representation (36.2),
\[
 \begin{split}
 \partial_u{\cal R}_n(t,u)
={}&-(C_n+t)^{-1}C_{n,u}(C_n+t)^{-1}\\
 &+(B_n+t)^{-1}B_{n,u}(B_n+t)^{-1}
 -\partial_u{\cal E}_n(t,u).                            \tag{38.12}
 \end{split}
\]

\begin{proposition}[Equivalent first-variation formula]
\label{prop:v15v38-response}
Suppose (38.11) or (38.12) is dominated in trace norm by an
integrable parameter-independent majorant, and the heat and
resolvent subtractions are related by (37.12).  Then
\[
 \boxed{
 \begin{split}
 \partial_u\operatorname{Tr}K_n
 &=
 \int_0^\infty\frac{\operatorname{Tr}
       [\partial_u{\cal A}_n(s,u)]}{s}\,\dd s\\
 &=
 \int_0^\infty
       \operatorname{Tr}[\partial_u{\cal R}_n(t,u)]\,\dd t.
 \end{split}}                                            \tag{38.13}
\]
When the unsubtracted finite-volume expressions are integrable,
\[
 \partial_u\operatorname{Tr}(\log B_n-\log C_n)
 =
 \operatorname{Tr}(B_n^{-1}B_{n,u}
                   -C_n^{-1}C_{n,u}).                   \tag{38.14}
\]
The renormalized formula is (38.14) minus the derivative of the
predeclared local counterterm.
\end{proposition}

\begin{proof}
Theorem \ref{thm:v15v38-c1-heat} gives the first equality.
Proposition \ref{prop:v15v37-heat-resolvent} and Bochner--Fubini
give the second.  For (38.14), cyclicity of the trace turns the
\(B_n\) term in (38.12) into
\[
 \operatorname{Tr}\bigl(B_{n,u}(B_n+t)^{-2}\bigr).
\]
Spectral calculus gives
\[
 \int_0^\infty(B_n+t)^{-2}\,\dd t=B_n^{-1}.              \tag{38.15}
\]
The \(C_n\) term is identical with the opposite sign.  The stated
integrability justifies all interchanges.
\end{proof}

For \(B_n=T_n^*T_n\),
\[
 B_{n,u}=T_{n,u}^*T_n+T_n^*T_{n,u}.                     \tag{38.16}
\]
The positive fermion modulus response is one half of (38.13).  The
phase and possible spectral flow remain in the determinant-line
packet and do not follow from (38.16).

\section{Metric variation and the stress distribution}

Let \(\Phi\) denote all nonmetric background fields.  Define the
renormalized fermion contribution with one fixed sign convention by
\[
 \Gamma_{\rm f}^{\rm ren}[g,\Phi]
 =-\frac12\operatorname{Tr}K[g,\Phi].                   \tag{38.17}
\]
If its metric derivative extends continuously to compactly
supported smooth symmetric tensors, define the Euclidean stress
distribution by
\[
 D_g\Gamma_{\rm f}^{\rm ren}[h]
 =\frac12\int_M\sqrt g\,T_{\rm f}^{\mu\nu}h_{\mu\nu}.    \tag{38.18}
\]
At the regulator level, (38.13) with \(u\mapsto g+uh\)
is the definition of the corresponding smeared stress.  It avoids
asserting a pointwise tensor before the distributional derivative
has been proved.

\begin{theorem}[Stress convergence from the \(C^1\) trace ideal]
\label{thm:v15v38-stress-convergence}
Let \({\cal T}\) be a locally convex space of test variations and
suppose the hypotheses of Theorem \ref{thm:v15v38-c1-heat} hold
uniformly for \(h\) in every bounded subset of \({\cal T}\), with
\(u\mapsto g+uh\).  Then
\[
 D_g\Gamma_{{\rm f},n}^{\rm ren}[h]
 \longrightarrow
 D_g\Gamma_{\rm f}^{\rm ren}[h]                          \tag{38.19}
\]
uniformly on bounded subsets of \({\cal T}\).  Hence the regulated
stress distributions converge in the strong dual
\({\cal T}'\).
\end{theorem}

\begin{proof}
Equation (38.6), multiplied by \(-1/2\), gives pointwise
convergence in \(h\).  The uniform derivative majorant on bounded
sets gives equicontinuity and uniform convergence there, which is
exactly convergence in the strong dual topology.
\end{proof}

\section{Covariance and the Ward identity}

Trace convergence does not by itself imply conservation.  Let
\(\varphi_\tau\) be the flow of a compactly supported vector field
\(X\).  The required covariance statement is
\[
 \Gamma_{\rm f}^{\rm ren}
 [\varphi_\tau^*g,\varphi_\tau^*\Phi]
 =
 \Gamma_{\rm f}^{\rm ren}[g,\Phi].                       \tag{38.20}
\]
This equality includes the carrier identification, zero-mode line,
regularization, and local counterterms.  Covariance of the bare
differential expression alone is insufficient.

\begin{theorem}[Distributional diffeomorphism Ward identity]
\label{thm:v15v38-ward}
Assume (38.20), differentiability in the sense of (38.18), and
write the nonmetric first variation as
\[
 D_\Phi\Gamma_{\rm f}^{\rm ren}[\Psi]
 =\sum_A\langle{\cal E}_A,\Psi^A\rangle.                 \tag{38.21}
\]
Then for every compactly supported vector field \(X\),
\[
 \boxed{
 \frac12\int_M\sqrt g\,T_{\rm f}^{\mu\nu}
              {\cal L}_Xg_{\mu\nu}
 +\sum_A\langle{\cal E}_A,{\cal L}_X\Phi^A\rangle=0.}    \tag{38.22}
\]
If the nonmetric fields satisfy their renormalized equations
\({\cal E}_A=0\), then
\[
 \nabla_\mu T_{\rm f}^{\mu}{}_\nu=0                     \tag{38.23}
\]
in the distributional sense.
\end{theorem}

\begin{proof}
Differentiate (38.20) at \(\tau=0\).  The derivatives of pullback
are Lie derivatives, giving (38.22).  Since
\[
 {\cal L}_Xg_{\mu\nu}=\nabla_\mu X_\nu+\nabla_\nu X_\mu,
\]
symmetry of the metric derivative and distributional integration by
parts turn the first term into
\[
 -\int_M\sqrt g\,(\nabla_\mu T_{\rm f}^{\mu}{}_\nu)X^\nu.
\]
On the nonmetric equations, this vanishes for every compactly
supported \(X\), proving (38.23).
\end{proof}

\begin{corollary}[Covariance-defect criterion]
\label{cor:v15v38-covariance-defect}
Suppose the regulator is not exactly diffeomorphism covariant, but
for every test vector field
\[
 {\cal W}_n(X)=
 D\Gamma_{{\rm f},n}^{\rm ren}
 [\,{\cal L}_Xg_n,{\cal L}_X\Phi_n\,]
 \longrightarrow0,                                     \tag{38.24}
\]
and the first variations converge as in Theorem
\ref{thm:v15v38-stress-convergence}.  Then the limit obeys
(38.22).
\end{corollary}

\begin{proof}
Pass to the limit in (38.24) using convergence of every first
variation.  The limiting expression is the left side of (38.22).
\end{proof}

Gauge covariance has the parallel form.  Differentiating an exact
or asymptotically vanishing gauge-covariance defect gives the
covariant current Ward identity, with Higgs and fermion equations
included off shell.

\section{Counterterm and anomaly ledger}

The variations of the \(a_0,a_2,a_4\) subtractions in V37 contribute
local tensors to (38.18).  They must be varied together with the
determinant remainder.  A failure of (38.20) after all allowed local
counterterms have been fixed is an anomaly or a regulator covariance
defect; it cannot be removed by declaring (38.23).

For a chiral Standard Model family, local gauge and mixed
gauge--gravitational anomaly cancellation requires the signed
representation traces to vanish.  Even when the familiar algebraic
cancellation holds, one still needs a regulator and determinant-line
orientation implementing it coherently across sectors.  V38 proves
the analytic Ward implication, not that physical implementation.

\begin{assessment}[Einstein handoff remains conditional]
\label{ass:v15v38-einstein}
If the gravitational and bosonic effective actions possess
convergent first variations and the total covariance defect tends to
zero, their stress tensors may be added to
\(T_{\rm f}^{\mu\nu}\).  A stationary point would then satisfy the
renormalized Euclidean Einstein equation in weak form.  V38 does not
prove existence of such a stationary point, tightness of an
interacting metric measure, Lorentzian continuation, or noncollapse.
\end{assessment}

\begin{assessment}[Literal source applicability]
\label{ass:v15v38-source}
The source contains formal current, stress, determinant-line, and
Einstein rows but does not provide the uniform \(C^1\)
trace-ideal majorant (38.2) for the selected full
Dirac--Yukawa operator.  Nor does it prove the covariance-defect
limit (38.24) on the same carrier.  The stress and Ward conclusions
therefore remain conditional on explicit physical estimates.
\end{assessment}

\begin{decision}[V39 from determinant response to measure tightness]
\label{dec:v15v38-next}
Return to the coupled bosonic measure.  Formulate a finite-volume
effective action consisting of the gauge--Higgs--gravity action,
the renormalized relative fermion logarithm, and its fixed
counterterms.  Prove a coercive exponential-tightness criterion on a
gauge-fixed, diffeomorphism-controlled configuration space, and show
exactly where a determinant lower bound, Gribov/metric quotient, and
conformal gravitational direction enter.  Do not infer positivity
of the chiral determinant from its modulus.
\end{decision}

\begin{firewall}[Claim boundary after V38]
\label{fire:v15v38}
V38 proves differentiability, stress convergence, and the Ward
identity under explicit trace-ideal and covariance hypotheses.
Those hypotheses, the determinant-line phase, the interacting
measure, and the Einstein continuum remain unproved for the physical
regulator.
\end{firewall}

\section*{Version V38 append-only record}
\addcontentsline{toc}{section}{Version V38 append-only record}
The complete V37 source was retained before this continuation.  It
had \(378087\) bytes, \(10858\) lines, and SHA--256
\[
 \texttt{AC287AE5ED84C8DA82C60DB6996E818926AA3310C6971D5F9B76FF635773A3B1}.
\]
V38 proved a \(C^1\) trace-ideal determinant theorem, derived its
Duhamel and resolvent responses, obtained distributional stress
convergence, and proved the on-shell Ward identity from an exact or
vanishing covariance defect.

\chapter{Fifteenth-cycle continuation V39: coercive tightness for the modulus-weighted interacting measure}
\label{sec:v15v39}

\section{From a determinant to a measure}

V36--V38 concern the relative fermion determinant and its first
variation.  A full Euclidean Standard Model--Einstein continuum also
requires a tight family of interacting measures.  This continuation
proves an abstract exponential-tightness theorem on a common
gauge-fixed or quotient carrier and records exactly which physical
estimates must populate it.

The positive fermion modulus can define a probability weight.  The
phase of a chiral determinant or the sign of a Pfaffian cannot be
discarded when reconstructing the physical theory.  Thus the
probability measure below is first the modulus envelope; a separate
phase, reflection, and anomaly argument is still required.

\section{Common quotient carrier}

Let \(({\cal X},d_{\cal X})\) be a Polish space representing the
chosen continuum configuration quotient or a proved global atlas
for it.  Each finite regulator has a measurable configuration space
\({\cal X}_n\), a probability reference measure \(\mu_n\), and an
identification
\[
 \iota_n:{\cal X}_n\longrightarrow{\cal X}.              \tag{39.1}
\]
All tightness statements refer to the pushforwards under
\(\iota_n\).  A regulator-dependent coordinate norm with no common
embedding does not define continuum tightness.

Let \(V:{\cal X}\to[0,\infty]\) be lower semicontinuous with compact
sublevel sets
\[
 K_R=\{x\in{\cal X}:V(x)\leq R\}.                        \tag{39.2}
\]
Such a \(V\) is a tightness function.  On \({\cal X}_n\), write
\(V_n=V\circ\iota_n\).

The modulus-weighted effective action is
\[
 {\cal A}_n(x)=
 S_{{\rm bos},n}(x)+S_{{\rm ct},n}(x)
 -\frac12\operatorname{Tr}K_n(x)
 -\log J_n^{\rm quot}(x),                               \tag{39.3}
\]
where \(J_n^{\rm quot}\) is the proved quotient or slice Jacobian.
The probability envelope is
\[
 \dd\nu_n(x)=Z_n^{-1}e^{-{\cal A}_n(x)}\,\dd\mu_n(x),
 \qquad
 Z_n=\int e^{-{\cal A}_n}\,\dd\mu_n.                    \tag{39.4}
\]

\section{Coercivity plus an anchor}

\begin{theorem}[Uniform exponential-tightness compiler]
\label{thm:v15v39-tightness}
Assume there are constants \(c>0\), \(C<\infty\), an anchor level
\(R_0<\infty\), a mass \(p_0>0\), and an anchor ceiling
\(C_0<\infty\), all independent of \(n\), such that
\begin{align}
 {\cal A}_n(x)&\geq cV_n(x)-C
 &&\text{for \(\mu_n\)-almost every \(x\)},              \tag{39.5}\\
 \mu_n(V_n\leq R_0)&\geq p_0,                            \tag{39.6}\\
 {\cal A}_n(x)&\leq C_0
 &&\text{for \(\mu_n\)-almost every \(x\) with
             \(V_n(x)\leq R_0\)}.                       \tag{39.7}
\end{align}
Then
\[
 Z_n\geq p_0e^{-C_0}                                    \tag{39.8}
\]
and, for every \(0<\eta<c\),
\[
 \boxed{
 \int e^{\eta V_n}\,\dd\nu_n
 \leq p_0^{-1}e^{C_0+C}.}                               \tag{39.9}
\]
Consequently
\[
 \boxed{
 \nu_n(V_n>R)
 \leq p_0^{-1}e^{C_0+C-\eta R}}                         \tag{39.10}
\]
and the pushforward family
\(\{(\iota_n)_\#\nu_n\}\) is tight on \({\cal X}\).
\end{theorem}

\begin{proof}
On the anchor set, (39.7) gives
\[
 Z_n\geq\int_{\{V_n\leq R_0\}}e^{-{\cal A}_n}\,\dd\mu_n
 \geq p_0e^{-C_0}.
\]
Since \(\mu_n\) is a probability measure and \(V_n\geq0\),
(39.5) gives, for \(0<\eta<c\),
\[
 \int e^{\eta V_n-{\cal A}_n}\,\dd\mu_n
 \leq e^C\int e^{-(c-\eta)V_n}\,\dd\mu_n
 \leq e^C.
\]
Division by (39.8) proves (39.9).  Markov's inequality proves
(39.10).  The sublevel \(K_R\) is compact and the right side of
(39.10) tends to zero uniformly, which is tightness.
\end{proof}

\begin{corollary}[Subsequential modulus measure]
\label{cor:v15v39-prokhorov}
Under Theorem \ref{thm:v15v39-tightness}, the pushforward measures
have a weakly convergent subsequence.
\end{corollary}

\begin{proof}
Prokhorov's theorem applies on the Polish space \({\cal X}\).
\end{proof}

This corollary gives existence of a subsequential probability
envelope only.  Identification of its correlation functions,
nontriviality, reflection positivity, and independence of the
subsequence require further estimates.

\section{Separate bosonic, fermionic, and quotient debits}

\begin{proposition}[Concrete coercivity ledger]
\label{prop:v15v39-ledger}
Suppose
\begin{align}
 S_{{\rm bos},n}+S_{{\rm ct},n}
 &\geq c_{\rm b}V_n-C_{\rm b},                           \tag{39.11}\\
 \frac12\operatorname{Tr}K_n
 &\leq c_{\rm f}V_n+C_{\rm f},                           \tag{39.12}\\
 \log J_n^{\rm quot}
 &\leq c_{\rm q}V_n+C_{\rm q},                           \tag{39.13}
\end{align}
with
\[
 c_{\rm b}>c_{\rm f}+c_{\rm q}.                         \tag{39.14}
\]
Then (39.5) holds with
\[
 c=c_{\rm b}-c_{\rm f}-c_{\rm q},\qquad
 C=C_{\rm b}+C_{\rm f}+C_{\rm q}.                       \tag{39.15}
\]
\end{proposition}

\begin{proof}
Substitute (39.11)--(39.13) into (39.3) and collect terms.
\end{proof}

The determinant estimate in (39.12) is an upper growth bound.  A
global positive lower bound for the determinant is not required for
tail coercivity.  It enters the anchor:

\begin{proposition}[Where a determinant lower bound is used]
\label{prop:v15v39-anchor-det}
Let \(A_n=\{V_n\leq R_0\}\).  Suppose on \(A_n\)
\[
 S_{{\rm bos},n}+S_{{\rm ct},n}\leq M_{\rm b},\qquad
 J_n^{\rm quot}\geq j_0>0,\qquad
 \exp\!\left(\frac12\operatorname{Tr}K_n\right)
 \geq d_0>0.                                             \tag{39.16}
\]
Then \({\cal A}_n\leq
M_{\rm b}-\log j_0-\log d_0\) on \(A_n\), so (39.7) holds.
\end{proposition}

\begin{proof}
Exponentiating the negative of (39.3) shows
\[
 e^{-{\cal A}_n}
 =e^{-S_{{\rm bos},n}-S_{{\rm ct},n}}
  e^{\frac12\operatorname{Tr}K_n}J_n^{\rm quot}.
\]
The three anchor bounds give the claimed lower bound on this density
and hence the ceiling for \({\cal A}_n\).
\end{proof}

It is enough to choose the anchor away from the determinant divisor.
Zeros elsewhere reduce the modulus density and do not damage the
upper tail estimate.  A moving divisor still matters for the phase,
sector weights, differentiability, and physical nontriviality.

\section{Gauge and diffeomorphism quotients}

For a compact lattice gauge group at fixed graph, the link
configuration space is compact before the continuum limit, but its
dimension grows.  Gauge fixing by a local slice introduces a
Faddeev--Popov or orbit Jacobian.  If a slice meets an orbit more than
once, its multiplicity and determinant cannot be hidden in
\(\mu_n\).  The alternatives needed for (39.1)--(39.3) are:
\begin{enumerate}
\item a directly metrized orbit space with a proved measurable
      quotient measure;
\item a finite or countable atlas with partition-of-unity and
      overlap multiplicities;
\item a global slice with a uniform Jacobian and inverse bound.
\end{enumerate}
The third alternative generally fails at a Gribov horizon.  In that
case the atlas or orbit-space route must record the horizon rather
than impose a fictitious global determinant floor.

The metric quotient is harder.  A family of lattice edge lengths or
coframes requires noncollapse, coordinate control, and an embedding
into a common topology such as measured Gromov--Hausdorff plus field
Sobolev data, or a fixed-manifold gauge with uniform geometric
bounds.  Without one of these constructions, \(V_n\) may be
coercive in coordinates while (39.2) fails in the physical quotient.

\section{The conformal Einstein obstruction}

The Euclidean Einstein--Hilbert action alone cannot supply
(39.11).  With the convention
\[
 S_{\rm EH}^{E}[g]
 =-\frac1{16\pi G}\int_M\sqrt g\,(R-2\Lambda),           \tag{39.17}
\]
write \(g=e^{2\sigma}\bar g\) on a closed four-manifold.  The
conformal transformation law and integration by parts give
\[
 \int_M\sqrt g\,R
 =
 \int_M\sqrt{\bar g}\,e^{2\sigma}
 \left(\bar R+6|\bar\nabla\sigma|^2\right).              \tag{39.18}
\]
Therefore
\[
 S_{\rm EH}^{E}[e^{2\sigma}\bar g]
 =
 -\frac1{16\pi G}
 \int_M\sqrt{\bar g}\,
 \left[
 e^{2\sigma}\bigl(\bar R+6|\bar\nabla\sigma|^2\bigr)
 -2\Lambda e^{4\sigma}
 \right].                                                \tag{39.19}
\]

\begin{countertheorem}[Einstein--Hilbert conformal noncoercivity]
\label{cth:v15v39-conformal}
On any coordinate patch where \(\bar g\) is smooth, there is a
bounded-amplitude sequence \(\sigma_k\) with frequency tending to
infinity such that
\[
 S_{\rm EH}^{E}[e^{2\sigma_k}\bar g]\longrightarrow-\infty.
                                                               \tag{39.20}
\]
Hence the bare Euclidean Einstein--Hilbert term does not imply
(39.11) for a tightness function controlling a first Sobolev norm of
the conformal factor.
\end{countertheorem}

\begin{proof}
Choose a nonzero smooth cutoff \(\rho\) in the patch and put
\(\sigma_k=\varepsilon\rho\sin(kx_1)\) with fixed sufficiently small
\(\varepsilon>0\).  The factors \(e^{2\sigma_k}\) and
\(e^{4\sigma_k}\) remain uniformly bounded above and below, while
\[
 \int e^{2\sigma_k}|\bar\nabla\sigma_k|^2
 =c\varepsilon^2k^2+O(k)
\]
for some \(c>0\).  In (39.19) this term has a negative coefficient.
The curvature and cosmological terms are \(O(1)\), so the action
tends to \(-\infty\).
\end{proof}

A positive higher-derivative term, a justified contour deformation,
or a different constructive gravitational measure could repair this
direction.  Each choice changes the measure and must be checked for
reflection positivity, unitarity, and the intended Einstein limit.
The \(R^2\) counterterm channel identified in V37 is not
automatically a physical coercive completion.

\section{Complex phase and total variation}

Let \(\Theta_n(x)\) be the determinant/Pfaffian phase.  The formal
physical measure is
\[
 \dd\omega_n=e^{i\Theta_n}\,\dd\nu_n.                   \tag{39.21}
\]
As an unnormalized complex measure its total variation equals
\(\nu_n\).  If one normalizes by
\[
 {\cal Z}_n^{\rm ph}=\int e^{i\Theta_n}\,\dd\nu_n,       \tag{39.22}
\]
then the total variation of the normalized complex measure is
\(|{\cal Z}_n^{\rm ph}|^{-1}\).  Thus tightness of \(\nu_n\) does
not prevent a phase-cancellation collapse
\({\cal Z}_n^{\rm ph}\to0\).

\begin{proposition}[Phase-safe subsequential criterion]
\label{prop:v15v39-phase}
If Theorem \ref{thm:v15v39-tightness} holds,
\(|{\cal Z}_n^{\rm ph}|\geq z_0>0\), and
\(e^{i\Theta_n}\) converges in \(L^1(\nu_n)\) after transport along
a tight subsequence, then the normalized complex measures have
uniformly bounded total variation and a weakly convergent
subsequence on bounded continuous observables.
\end{proposition}

\begin{proof}
The phase lower bound gives total variation at most \(z_0^{-1}\).
Tightness of the variation measures and the \(L^1\) convergence
permit approximation by compactly supported bounded observables.
Weak compactness of finite complex measures then gives the claimed
subsequence.
\end{proof}

This proposition does not give Osterwalder--Schrader positivity.
That is a separate constraint on reflected observables.

\begin{assessment}[Physical population status]
\label{ass:v15v39-status}
The earlier scalar and local-block gap estimates can contribute to
(39.11)--(39.12) on selected regions, but they do not provide a
common compactness function for the gauge, Higgs, and metric
quotient.  The source also lacks a global quotient construction, a
coercive conformal gravitational measure, a uniform anchor mass, and
a phase lower bound.  Consequently Theorem
\ref{thm:v15v39-tightness} is a precise compiler, not yet a
constructed SM--Einstein measure.
\end{assessment}

\begin{decision}[V40 compulsory audit and bottleneck choice]
\label{dec:v15v39-next}
At V40 inspect the newest YM and NS companion files.  Audit whether
their current compactness, chart, or kernel estimates can populate
any one of (39.1), (39.11), or (39.16) on the physical shared
carrier.  Then continue at V41 with the smallest uncovered item,
giving priority to a common topology and a noncollapse/tightness
function over another formal effective-action identity.
\end{decision}

\begin{firewall}[Claim boundary after V39]
\label{fire:v15v39}
V39 proves a coercivity-and-anchor theorem for the positive modulus
measure and isolates the conformal, quotient, determinant, and phase
debits.  No physical tight interacting measure, OS-positive chiral
measure, noncollapsed metric limit, or Einstein equation has been
constructed.
\end{firewall}

\section*{Version V39 append-only record}
\addcontentsline{toc}{section}{Version V39 append-only record}
The complete V38 source was retained before this continuation.  It
had \(391388\) bytes, \(11230\) lines, and SHA--256
\[
 \texttt{273CF5028368FE295D511C70F540E227416D73CCA57DCF838DB9EC503DC9CF0E}.
\]
V39 proved uniform exponential tightness from a common coercive
action and anchor, separated fermion and quotient debits, proved the
Einstein--Hilbert conformal noncoercivity, and isolated the
additional phase-cancellation bound needed for complex measures.

\chapter{Fifteenth-cycle continuation V40: compulsory compactness audit}
\label{sec:v15v40}

\section{Current companion checkpoints}

The newest active companion files were inspected read-only:
\[
\begin{array}{c|c|c|c}
\text{programme}&\text{file}&\text{lines}&\text{SHA--256}\\ \hline
\mathrm{YM}&
\texttt{Renewal\_Yang\_Mills\_Mass\_Gap\_Research\_Notes\_V39.tex}&
9076&
\texttt{9B44B0FBCA6037F3319E86036753B0F3B76BDD06F6045F00747E05DC01E23F9D}\\
\mathrm{NS}&
\texttt{Renewal\_NS\_Stability\_Research\_Notes\_V26.tex}&
5319&
\texttt{82C887F8C2C69E4F28FAD7C3DC8DE5B9099CB5A4BF431EBA1FFF3E91935978AD}
\end{array}                                                \tag{40.1}
\]

\section{Audit findings}

YM V39 proves exact pointwise positivity on
\[
 [0,\pi/4]^2\times[0,t]                                  \tag{40.2}
\]
by independently certified rational charts and explicit cap
coverage.  Its firewall states that no interacting continuum
measure, reflection-positive limit, or global mass gap follows.
Thus it reinforces the pointwise coverage discipline of V35 but
does not populate the common carrier (39.1), bosonic coercivity
(39.11), or anchor (39.16).

NS V26 computes pointwise rank-one norms of two Oseen-kernel
derivatives.  It has not yet integrated the refined second-derivative
norm into its announced global constant.  Its fixed Euclidean
kernel carrier and PDE estimates provide neither a quotient topology
for gauge--metric configurations nor a gravitational noncollapse
bound.  Its nuclear-summation lesson was already extracted in V35
and is not duplicated here.

\begin{theorem}[V40 nonimport decision]
\label{thm:v15v40-nonimport}
Neither current companion note proves any of (39.1), (39.11), or
(39.16) for the Standard Model--Einstein regulator.  No companion
constant or positivity floor is imported.  The smallest useful next
step is to construct a common compactness topology under explicit
fixed-manifold Sobolev-slice and metric-ellipticity hypotheses.
\end{theorem}

\begin{proof}
The YM theorem concerns a finite Wilson pencil in momentum space and
its own text excludes an interacting continuum measure.  The NS
theorem concerns contractions of an Oseen kernel with rank-one
tensors.  Neither statement supplies a map from regulated
gauge--Higgs--metric configurations into a common Polish space or
controls metric volume ratios.  The proposed Sobolev construction
addresses exactly those missing definitions without treating the
companion calculations as premises.
\end{proof}

\begin{decision}[V41 fixed-manifold compactness and noncollapse]
\label{dec:v15v40-next}
On a fixed compact four-manifold with background metric, define a
common \(H^{s'}\) configuration carrier for metric, connection, and
Higgs representatives.  For \(s>s'>3\), prove compactness of
\(H^s\)-bounded, uniformly elliptic sublevels by Rellich's theorem.
Prove an explicit lower volume ratio from bilipschitz metric bounds.
State precisely the residual gauge/diffeomorphism slice hypothesis
needed before this construction can populate V39.
\end{decision}

\begin{firewall}[Claim boundary after V40]
\label{fire:v15v40}
V40 is a compulsory audit.  It changes the next construction but
adds no physical continuum estimate and imports no conclusion from
the YM or NS programmes.
\end{firewall}

\section*{Version V40 append-only record}
\addcontentsline{toc}{section}{Version V40 append-only record}
The complete V39 source was retained before this continuation.  It
had \(404359\) bytes, \(11584\) lines, and SHA--256
\[
 \texttt{350D8D159E5E34E1AF6A16DE7515ACF80707011124F3C54B5B402692FB07E55B}.
\]
V40 performed the required companion review and selected a
fixed-manifold Sobolev compactness and noncollapse theorem for V41.

\chapter{Fifteenth-cycle continuation V41: a fixed-manifold Sobolev carrier and quantitative noncollapse}
\label{sec:v15v41}

\section{A concrete common topology}

Let \(M\) be a fixed closed \(d\)-manifold, eventually \(d=4\),
with smooth background metric \(g_0\).  Fix the gauge and Higgs
bundles and smooth background connections so that a configuration
representative can be written
\[
 x=(h,a,\phi),\qquad g=g_0+h.                            \tag{41.1}
\]
Choose
\[
 s>s'>\frac d2+1.                                       \tag{41.2}
\]
Define the common Hilbert carrier
\[
 {\cal X}^{s'}=
 H^{s'}(S^2T^*M)\times
 H^{s'}(T^*M\otimes\operatorname{ad}P)\times
 H^{s'}(E_H).                                            \tag{41.3}
\]
Sobolev embedding makes the metric component \(C^1\), so positivity
and uniform ellipticity are stable under convergence in
\({\cal X}^{s'}\).

Fix \(0<\lambda<1<\Lambda<\infty\) and restrict to representatives
satisfying
\[
 \lambda g_0\leq g_0+h\leq\Lambda g_0.                  \tag{41.4}
\]
Define the extended tightness function
\[
 V_s(h,a,\phi)=
 \begin{cases}
 \|h\|_{H^s}^2+\|a\|_{H^s}^2+\|\phi\|_{H^s}^2,
   &\text{if the fields belong to \(H^s\) and (41.4) holds},\\
 +\infty,&\text{otherwise}.
 \end{cases}                                             \tag{41.5}
\]

\begin{theorem}[Compact Sobolev sublevels]
\label{thm:v15v41-compact-sublevels}
For every \(R<\infty\), the sublevel
\[
 {\cal K}_R=\{x\in{\cal X}^{s'}:V_s(x)\leq R\}           \tag{41.6}
\]
is compact in \({\cal X}^{s'}\).
\end{theorem}

\begin{proof}
Let \(x_n=(h_n,a_n,\phi_n)\) lie in \({\cal K}_R\).  Each component
is bounded in \(H^s\).  Weak compactness gives a subsequence
converging weakly in \(H^s\), and the Rellich compact embedding
\[
 H^s\Subset H^{s'}                                      \tag{41.7}
\]
gives strong convergence in every component of (41.3).  Lower
semicontinuity of the \(H^s\) norm places the limit in the radius
\(R\) ball.  Since \(s'>d/2\), the convergence is uniform.  Passing
to the limit in
\(\lambda g_0(v,v)\leq(g_0+h_n)(v,v)\leq
\Lambda g_0(v,v)\) for every tangent vector \(v\) preserves
(41.4).  Thus every sequence has a convergent subsequence whose
limit remains in \({\cal K}_R\).
\end{proof}

\begin{corollary}[Population of the abstract tightness function]
\label{cor:v15v41-populate-v39}
If regulator interpolants
\[
 \iota_n:{\cal X}_n\longrightarrow{\cal X}^{s'}         \tag{41.8}
\]
land in the elliptic class (41.4), then
\(V_n=V_s\circ\iota_n\) has compact sublevels.  Therefore the
coercivity and anchor hypotheses (39.5)--(39.7) imply tightness of
the interpolated modulus measures in \({\cal X}^{s'}\).
\end{corollary}

\begin{proof}
Theorem \ref{thm:v15v41-compact-sublevels} verifies (39.2).
Theorem \ref{thm:v15v39-tightness} then applies.
\end{proof}

\section{Quantitative metric noncollapse}

Uniform ellipticity alone gives a useful metric-measure bound.  From
(41.4),
\[
 \sqrt\lambda\,d_{g_0}(x,y)
 \leq d_g(x,y)
 \leq\sqrt\Lambda\,d_{g_0}(x,y),                        \tag{41.9}
\]
and
\[
 \lambda^{d/2}\,\dd{\rm vol}_{g_0}
 \leq\dd{\rm vol}_g
 \leq\Lambda^{d/2}\,\dd{\rm vol}_{g_0}.                 \tag{41.10}
\]

\begin{theorem}[Uniform volume noncollapse]
\label{thm:v15v41-noncollapse}
There are \(r_0,c_0>0\), depending only on the fixed smooth metric
\(g_0\), such that every metric satisfying (41.4) obeys
\[
 \boxed{
 {\rm vol}_g(B_g(x,r))
 \geq
 c_0\left(\frac{\lambda}{\Lambda}\right)^{d/2}r^d}
 \qquad(0<r\leq r_0\sqrt\Lambda).                       \tag{41.11}
\]
In dimension four the prefactor is
\(c_0\lambda^2\Lambda^{-2}\).
\end{theorem}

\begin{proof}
Compactness and smoothness of \((M,g_0)\) give
\[
 {\rm vol}_{g_0}(B_{g_0}(x,\rho))\geq c_0\rho^d
 \quad(0<\rho\leq r_0)                                  \tag{41.12}
\]
uniformly in \(x\).  The upper distance inequality in (41.9) gives
\[
 B_{g_0}(x,r/\sqrt\Lambda)\subset B_g(x,r).
\]
Use the lower volume-form inequality in (41.10) and then (41.12):
\[
 \begin{split}
 {\rm vol}_g(B_g(x,r))
 &\geq\lambda^{d/2}
 {\rm vol}_{g_0}(B_{g_0}(x,r/\sqrt\Lambda))\\
 &\geq c_0(\lambda/\Lambda)^{d/2}r^d.
 \end{split}
\]
\end{proof}

\begin{corollary}[Diameter and total-volume control]
\label{cor:v15v41-diameter-volume}
Every metric in the elliptic class obeys
\[
 \operatorname{diam}_gM
 \leq\sqrt\Lambda\,\operatorname{diam}_{g_0}M,\qquad
 \lambda^{d/2}{\rm vol}_{g_0}M
 \leq{\rm vol}_gM
 \leq\Lambda^{d/2}{\rm vol}_{g_0}M.                     \tag{41.13}
\]
\end{corollary}

\begin{proof}
Apply (41.9) to the diameter and integrate (41.10).
\end{proof}

Theorem \ref{thm:v15v41-noncollapse} is stronger than a total-volume
lower bound: it prevents volume from disappearing at every scale
below the fixed background radius.  It does not control curvature
unless the Sobolev bound in (41.5) is also used.

\section{Residual quotient theorem}

A global gauge or diffeomorphism slice is not automatic.  The
following statement identifies the exact harmless residual case.

\begin{proposition}[Compact residual quotient preserves compact sublevels]
\label{prop:v15v41-compact-quotient}
Let a compact topological group \(G_{\rm res}\) act continuously on
\({\cal X}^{s'}\), preserve \(V_s\), and suppose the chosen slice
meets every physical orbit with residual ambiguity exactly
\(G_{\rm res}\).  Then
\[
 {\cal X}_{\rm phys}={\cal X}^{s'}/G_{\rm res}           \tag{41.14}
\]
is metrizable and the image of every \({\cal K}_R\) is compact.
Consequently \(V_s\) descends to a tightness function on the
physical quotient.
\end{proposition}

\begin{proof}
A continuous image of a compact set is compact.  A continuous
action of a compact group on a metrizable space has a metrizable
Hausdorff orbit space; for example, average a bounded compatible
metric over the group and take the induced orbit distance.
Invariance makes the value of \(V_s\) independent of the
representative.
\end{proof}

For the full nonabelian gauge and diffeomorphism groups, the premise
of Proposition \ref{prop:v15v41-compact-quotient} requires proof.
A local Coulomb or harmonic slice must have a uniform inverse for
its Faddeev--Popov or gauge-fixing operator on the stated region.
If Gribov copies or metric stabilizers occur, one needs a compatible
atlas and overlap multiplicities.  Simply declaring a gauge
condition does not define (41.14).

\section{Discrete interpolants}

The maps (41.8) must be stable and consistent.  A sufficient packet
is:
\[
\begin{gathered}
\|\iota_nx\|_{H^s}\leq
C_{\rm int}\|x\|_{{\rm disc},s},                         \tag{41.15}\\
\lambda g_0\leq g(\iota_nx)\leq\Lambda g_0,              \tag{41.16}\\
\|\iota_{n+1}R_{n+1,n}x-\iota_nx\|_{H^{s'}}
\leq\epsilon_n(1+\|x\|_{{\rm disc},s}),
\qquad\epsilon_n\to0.                                   \tag{41.17}
\end{gathered}
\]
Here \(R_{n+1,n}\) is the declared refinement transport.  Equation
(41.15) transfers action coercivity, (41.16) transfers
noncollapse, and (41.17) prevents the common topology from changing
with the cutoff.  None follows solely from combinatorial incidence.

\begin{proposition}[Cofinal consistency]
\label{prop:v15v41-cofinal}
If \(V_n(x_n)\leq R\) and (41.15)--(41.17) hold, then every
subsequence of interpolants has a further subsequence converging in
\({\cal X}^{s'}\), and adjacent refinements have the same set of
limit points.
\end{proposition}

\begin{proof}
Compactness follows from Theorem
\ref{thm:v15v41-compact-sublevels}.  On a bounded discrete
\(s\)-ball, the right side of (41.17) tends to zero.  Therefore an
adjacent refinement and its predecessor have vanishing
\({\cal X}^{s'}\) distance and hence identical limit points.
\end{proof}

\section{The regularity debit}

The simple carrier has a significant cost.  In four dimensions,
(41.2) requires
\[
 s>s'>3.                                                 \tag{41.18}
\]
The ordinary Yang--Mills and Higgs kinetic actions control roughly
one weak derivative, while curvature-squared gravity controls
roughly two derivatives of the metric after gauge fixing.  They do
not directly give an \(H^s\) bound with \(s>3\).  Elliptic
bootstrapping could provide it for classical solutions with
controlled sources, but a path-integral measure is not supported
only on smooth solutions.

\begin{assessment}[Meaning of the V41 branch]
\label{ass:v15v41-meaning}
V41 gives a complete high-regularity route to the common topology
and noncollapse requirements of V39.  It is conditional on uniform
ellipticity, a proved slice or atlas, stable interpolants, and
coercivity above three Sobolev derivatives.  The current source does
not provide those estimates.  Thus this branch is mathematically
closed as a compiler but physically unpopulated.
\end{assessment}

\begin{decision}[V42 weaken the topology]
\label{dec:v15v41-next}
Go upstream from the excessive Sobolev debit.  Formulate a weaker
metric-measure carrier using uniform ellipticity or direct distance
comparison, \(L^p\) curvature control, and weak Sobolev fields.
Prove which part of noncollapse and compactness survives at the
renormalizable derivative count.  If no compact embedding is
available, identify the minimal added moment or fractional
regularity rather than assuming smoothness.
\end{decision}

\begin{firewall}[Claim boundary after V41]
\label{fire:v15v41}
V41 proves compact Sobolev sublevels, quantitative volume
noncollapse, and compact residual quotient descent under explicit
hypotheses.  Those hypotheses are stronger than the presently
proved physical action bounds, so no tight Standard Model--Einstein
measure has yet been obtained.
\end{firewall}

\section*{Version V41 append-only record}
\addcontentsline{toc}{section}{Version V41 append-only record}
The complete V40 source was retained before this continuation.  It
had \(408094\) bytes, \(11673\) lines, and SHA--256
\[
 \texttt{6F3AA866A77F55AA39CCB2DE4FBF4A190CBC8837C29DAB4DB1F08014C90ACF64}.
\]
V41 constructed a fixed-manifold Sobolev carrier with compact
sublevels, proved an explicit four-dimensional volume noncollapse
bound, and isolated the slice, interpolation, ellipticity, and
higher-regularity inputs still missing from the physical regulator.

\chapter{Fifteenth-cycle continuation V42: a supercritical metric carrier at nearly the curvature-squared scale}
\label{sec:v15v42}

\section{Lowering the regularity requirement}

V41 used \(H^s\) with \(s>3\) in four dimensions so that compact
sublevels converged in a \(C^1\) metric topology.  That requirement
is substantially stronger than the two-derivative scale suggested
by curvature-squared control.  The present continuation lowers the
metric hypothesis to
\[
 W^{2,p},\qquad p>\frac d2,                              \tag{42.1}
\]
and uses only \(H^1\) for the gauge and Higgs representatives.
In dimension four, (42.1) is \(p=2+\delta\) for any
\(\delta>0\).  This is still an additional moment beyond the
critical \(L^2\)-curvature scale, but it is the smallest direct
Sobolev route to uniform metric coefficients.

Retain the fixed compact manifold, bundles, background metric
\(g_0\), and ellipticity constants of V41.  Choose
\[
 p>\frac d2,\qquad
 0<\beta<2-\frac dp.                                    \tag{42.2}
\]
Define
\[
 {\cal Y}^{\beta}=
 C^{0,\beta}(S^2T^*M)\times
 L^2(T^*M\otimes\operatorname{ad}P)\times L^2(E_H).
                                                               \tag{42.3}
\]
For representatives \(x=(h,a,\phi)\), put
\[
 V_{p,1}(x)=
 \begin{cases}
 \|h\|_{W^{2,p}}^p+\|a\|_{H^1}^2+\|\phi\|_{H^1}^2,
 &\lambda g_0\leq g_0+h\leq\Lambda g_0,\\
 +\infty,&\text{otherwise}.
 \end{cases}                                             \tag{42.4}
\]

\begin{theorem}[Nearly critical compact carrier]
\label{thm:v15v42-compact}
Every sublevel
\[
 {\cal C}_R=\{x\in{\cal Y}^{\beta}:V_{p,1}(x)\leq R\}    \tag{42.5}
\]
is compact in \({\cal Y}^{\beta}\).
\end{theorem}

\begin{proof}
Morrey--Rellich compactness gives
\[
 W^{2,p}(M)\Subset C^{0,\beta}(M)
 \quad\text{for }\beta<2-d/p.                           \tag{42.6}
\]
The Rellich theorem also gives
\[
 H^1(M)\Subset L^2(M).                                  \tag{42.7}
\]
Thus any sequence in \({\cal C}_R\) has a subsequence converging
strongly in all three factors of (42.3).  Weak lower
semicontinuity retains the bounds in (42.4).  Uniform convergence
of \(h_n\) preserves the two ellipticity inequalities.  The limit
therefore belongs to \({\cal C}_R\).
\end{proof}

\begin{corollary}[Tight modulus measures at reduced regularity]
\label{cor:v15v42-tight}
If regulator interpolants land in (42.3) and their effective actions
satisfy
\[
 {\cal A}_n(x)\geq cV_{p,1}(\iota_nx)-C                 \tag{42.8}
\]
together with the anchor hypotheses (39.6)--(39.7), then their
modulus-weighted measures are tight on \({\cal Y}^{\beta}\).
\end{corollary}

\begin{proof}
Theorem \ref{thm:v15v42-compact} supplies the compact sublevels
required by Theorem \ref{thm:v15v39-tightness}.
\end{proof}

\section{Metric and measure convergence}

The metric component in (42.3) is strong enough to identify the
underlying metric-measure limit.

\begin{proposition}[Uniform convergence of distances and volumes]
\label{prop:v15v42-distance}
Let \(g_n,g\) satisfy (41.4) and
\(\|g_n-g\|_{C^0(g_0)}\to0\).  Then
\[
 \sup_{x,y\in M}|d_{g_n}(x,y)-d_g(x,y)|\longrightarrow0, \tag{42.9}
\]
and
\[
 \left\|\frac{\dd{\rm vol}_{g_n}}{\dd{\rm vol}_{g_0}}
 -\frac{\dd{\rm vol}_{g}}{\dd{\rm vol}_{g_0}}\right\|_{C^0}
 \longrightarrow0.                                     \tag{42.10}
\]
Consequently \((M,d_{g_n},{\rm vol}_{g_n})\) converges to
\((M,d_g,{\rm vol}_g)\) in measured Gromov--Hausdorff topology
using the identity correspondence.
\end{proposition}

\begin{proof}
Uniform ellipticity turns the coefficient convergence into numbers
\(\epsilon_n\downarrow0\) such that
\[
 (1-\epsilon_n)g\leq g_n\leq(1+\epsilon_n)g.             \tag{42.11}
\]
Lengths of every absolutely continuous curve are therefore between
\(\sqrt{1-\epsilon_n}\) and
\(\sqrt{1+\epsilon_n}\) times its \(g\)-length.  Taking infima over
curves proves (42.9), uniformly because \(d_g\) has bounded
diameter.  The map from a uniformly positive symmetric matrix to
the square root of its determinant is uniformly continuous on the
compact ellipticity interval, which proves (42.10).  These two
uniform convergences give measured Gromov--Hausdorff convergence.
\end{proof}

\begin{corollary}[Noncollapse survives]
\label{cor:v15v42-noncollapse}
Every metric in a sublevel (42.5) obeys the volume lower bound
(41.11), independently of its \(W^{2,p}\) norm.  This lower bound
passes to every limit in Theorem \ref{thm:v15v42-compact}.
\end{corollary}

\begin{proof}
The proof of Theorem \ref{thm:v15v41-noncollapse} uses only uniform
ellipticity.  Proposition \ref{prop:v15v42-distance} identifies the
limit metric and measure.
\end{proof}

\section{From curvature to the coordinate bound}

An \(L^p\) curvature estimate is geometric, whereas (42.4) is
coordinate dependent.  Their bridge requires a controlled gauge.
In a harmonic coordinate chart the metric satisfies schematically
\[
 -\frac12g^{ab}\partial_a\partial_b g_{ij}
 ={\rm Ric}_{ij}+Q_{ij}(g^{-1},\partial g,\partial g).   \tag{42.12}
\]

\begin{proposition}[Uniform harmonic-coordinate bridge]
\label{prop:v15v42-harmonic}
Suppose a finite harmonic-coordinate atlas has uniformly controlled
chart geometry, the metrics satisfy (41.4), and
\[
 \|{\rm Ric}(g)\|_{L^p}
 +\|\partial g\|_{L^{2p}}^2
 +\|g\|_{L^p}\leq C                                    \tag{42.13}
\]
in those charts, with \(p>d/2\).  Then
\[
 \|g\|_{W^{2,p}}\leq C'                                 \tag{42.14}
\]
for a constant depending only on the displayed data and the atlas.
\end{proposition}

\begin{proof}
Uniform ellipticity makes the principal part in (42.12) uniformly
elliptic.  The quadratic term is bounded in \(L^p\) by
\(C\|\partial g\|_{L^{2p}}^2\).  Interior
Calderon--Zygmund estimates on smaller charts give
\[
 \|g\|_{W^{2,p}}
 \leq C_{\rm CZ}
 \bigl(\|{\rm Ric}\|_{L^p}
       +\|\partial g\|_{L^{2p}}^2+\|g\|_{L^p}\bigr).
                                                               \tag{42.15}
\]
A finite subordinate cover and the assumed uniform chart geometry
give (42.14).
\end{proof}

The existence of such a uniform harmonic atlas is a separate
geometric theorem, usually obtained from injectivity-radius,
noncollapse, and curvature hypotheses.  It cannot be inferred from
the coordinate formula (42.12), since that would be circular.

\section{The critical four-dimensional failure}

At \(p=d/2\), the embedding (42.6) reaches its critical endpoint and
does not give compactness in \(C^0\).

\begin{countertheorem}[Bounded \(W^{2,2}\) metrics need not be
coefficient compact]
\label{cth:v15v42-critical}
Let \(d=4\), choose a coordinate ball, and take
\(\psi\in C_c^\infty(\mathbb R^4)\) with \(\psi(0)=1\).
For fixed sufficiently small \(\varepsilon>0\), define in the ball
\[
 \sigma_k(x)=\varepsilon\psi(kx),\qquad
 g_k=e^{2\sigma_k}g_0,                                  \tag{42.16}
\]
extended by \(g_0\) outside.  Then the metrics are uniformly
elliptic and
\[
 \sup_k\|\sigma_k\|_{W^{2,2}}<\infty,                   \tag{42.17}
\]
but no subsequence of \(\sigma_k\) converges uniformly.
\end{countertheorem}

\begin{proof}
For a derivative of order \(j\),
\[
 \|D^j\psi(k\cdot)\|_{L^2(\mathbb R^4)}
 =k^{j-2}\|D^j\psi\|_{L^2}.                             \tag{42.18}
\]
Thus the second derivative is uniformly bounded and the lower
derivatives even tend to zero.  Small fixed \(\varepsilon\) gives
uniform ellipticity.  Away from the concentration point,
\(\sigma_k(x)\to0\), while \(\sigma_k(0)=\varepsilon\).
Any uniform limit would be continuous and equal to zero off the
origin, hence zero at the origin, a contradiction.
\end{proof}

This counterexample does not prove that every metric-measure
topology fails at the critical exponent; its shrinking perturbation
can be small in distance.  It proves the narrower statement needed
here: critical \(W^{2,2}\) control cannot populate the
coefficient-level compact carrier (42.3).

\section{Minimal additional regularity}

There are two direct repairs:
\begin{enumerate}
\item any uniform \(W^{2,2+\delta}\) moment with \(\delta>0\),
      giving Theorem \ref{thm:v15v42-compact};
\item a critical \(W^{2,2}\) bound together with an independent
      uniform modulus of continuity for \(g\), giving \(C^0\)
      precompactness by Arzela--Ascoli.
\end{enumerate}
The first is a higher integrability row.  The second is a
nonconcentration row.  Neither follows merely from finiteness of a
curvature-squared action at each cutoff.

\begin{assessment}[Derivative-count improvement]
\label{ass:v15v42-improvement}
Compared with V41, the metric requirement has dropped from more than
three \(L^2\) derivatives to slightly more than two \(L^2\)-scale
derivatives, and the gauge--Higgs carrier uses their natural
\(H^1\)-to-\(L^2\) compact embedding.  Uniform ellipticity continues
to give noncollapse.  The price is one explicit supercritical
curvature moment or an independent equicontinuity estimate.
\end{assessment}

\begin{assessment}[Literal source applicability]
\label{ass:v15v42-source}
The source does not supply a uniform \(L^{2+\delta}\) curvature
moment, a global controlled harmonic atlas, or an equicontinuity
estimate for the interpolated metric.  It also does not construct
the interpolation maps into (42.3).  V42 therefore narrows the
compactness bottleneck but does not populate it.
\end{assessment}

\begin{decision}[V43 probabilistic higher-integrability bridge]
\label{dec:v15v42-next}
Work at the level of measures rather than deterministic action
sublevels.  Prove that a uniform exponential moment of a
\(W^{2,2+\delta}\) metric norm and \(H^1\) field norms yields
exponential tightness in (42.3), even if the bare action is only
known to control the critical norm.  Then isolate a local reverse
Holder or concentration estimate that would generate the extra
\(\delta\), and show why expectation bounds without uniform
integrability are insufficient.
\end{decision}

\begin{firewall}[Claim boundary after V42]
\label{fire:v15v42}
V42 proves a nearly critical compact metric--field carrier,
metric-measure convergence, and noncollapse under
\(W^{2,2+\delta}\), \(H^1\), ellipticity, and gauge-control
hypotheses.  The physical regulator has not yet been shown to obey
the required supercritical moment or harmonic-coordinate estimate.
\end{firewall}

\section*{Version V42 append-only record}
\addcontentsline{toc}{section}{Version V42 append-only record}
The complete V41 source was retained before this continuation.  It
had \(418280\) bytes, \(11959\) lines, and SHA--256
\[
 \texttt{0F1A3CAEF25C994A0C1771C357887786F25070C0079705192555B20DD39978AD}.
\]
V42 replaced the high Sobolev carrier by a
\(W^{2,2+\delta}\)-metric and \(H^1\)-field carrier, proved compact
sublevels and metric-measure convergence, and exhibited the critical
\(W^{2,2}\) concentration that necessitates an additional
higher-integrability or equicontinuity row.

\chapter{Fifteenth-cycle continuation V43: probabilistic higher integrability from a local reverse-Hölder row}
\label{sec:v15v43}

\section{Critical energy versus compactness energy}

V42 showed that in four dimensions the coefficient-level metric
carrier becomes compact once a \(W^{2,2+\delta}\) bound is
available.  A curvature-squared action naturally controls the
critical quantity
\[
 E_2(g)=\int_M|{\rm Rm}(g)|^2\,\dd{\rm vol}_g,           \tag{43.1}
\]
not its supercritical analogue.  The present continuation gives a
probabilistic bridge from \(E_2\) to compactness, conditional on one
local reverse-Hölder estimate.  It also proves that no bridge follows
from critical moments alone.

Let
\[
 q=1+\frac{\delta}{2}>1,\qquad
 E_{2+\delta}(g)=
 \int_M|{\rm Rm}(g)|^{2+\delta}\,\dd{\rm vol}_g.         \tag{43.2}
\]
Use a fixed finite cover by coordinate balls
\(\{B_j\}_{j=1}^J\), with doubled balls \(2B_j\) of overlap at most
\(N_{\rm ov}\).

\section{A finite-cover reverse-Hölder compiler}

\begin{definition}[Physical reverse-Hölder row]
\label{def:v15v43-rh}
A metric \(g\) passes the \((q,C_{\rm RH})\) row with defect
\(D(g)\geq0\) if, for every \(j\),
\[
 \left(\fint_{B_j}|{\rm Rm}|^{2q}\,\dd{\rm vol}_g\right)^{1/q}
 \leq
 C_{\rm RH}\fint_{2B_j}|{\rm Rm}|^2\,\dd{\rm vol}_g
 d_j(g),                                                \tag{43.3}
\]
where
\[
 \sum_{j=1}^J {\rm vol}_g(B_j)\,d_j(g)^q\leq D(g).       \tag{43.4}
\]
\end{definition}

\begin{proposition}[Local-to-global higher integrability]
\label{prop:v15v43-local-global}
Assume uniform ellipticity (41.4), the fixed cover above, and
(43.3)--(43.4).  Then
\[
 \boxed{
 E_{2+\delta}(g)
 \leq C\bigl(1+E_2(g)^q+D(g)\bigr),}                    \tag{43.5}
\]
where \(C\) depends only on the cover, ellipticity constants,
\(q\), and \(C_{\rm RH}\).
\end{proposition}

\begin{proof}
Raise (43.3) to the power \(q\), multiply by
\({\rm vol}_g(B_j)\), and use
\((u+v)^q\leq2^{q-1}(u^q+v^q)\).  Uniform ellipticity makes the
ratios of volumes of \(B_j\) and \(2B_j\) uniformly bounded above
and below.  Therefore
\[
 \int_{B_j}|{\rm Rm}|^{2q}
 \leq C\left[
 \left(\int_{2B_j}|{\rm Rm}|^2\right)^q
 {\rm vol}_g(B_j)d_j^q\right].                          \tag{43.6}
\]
Sum over \(j\).  Each local critical integral is at most
\(E_2(g)\), and the finite overlap and (43.4) give
\[
 \int_M|{\rm Rm}|^{2q}
 \leq C(E_2(g)^q+D(g)).
\]
Since \(2q=2+\delta\), this is (43.5), with the harmless constant
covering the zero-energy case.
\end{proof}

The estimate (43.3) must be proved from the physical regulated
equations, a conditional Gibbs estimate, or another local mechanism.
It is not a generic property of metrics with finite \(E_2\).

\section{Probabilistic tightness theorem}

Let \(\nu_n\) be the modulus measures of V39, transported to a
common fixed-manifold representative space.  Let \(V_\delta\) denote
the compactness function (42.4), with the \(W^{2,2+\delta}\) metric
norm controlled by Proposition \ref{prop:v15v42-harmonic} and the
\(H^1\) gauge--Higgs terms retained.

\begin{theorem}[Critical moments plus reverse Hölder imply tightness]
\label{thm:v15v43-prob-tight}
Assume:
\begin{enumerate}
\item uniform ellipticity and a uniform controlled harmonic atlas;
\item every sampled metric satisfies (43.3);
\item for some \(\eta>0\),
      \[
       \sup_n\int e^{\eta E_2(g)}\,\dd\nu_n<\infty;       \tag{43.7}
      \]
\item
      \[
       \sup_n\int D(g)\,\dd\nu_n<\infty;                 \tag{43.8}
      \]
\item the coordinate and field terms omitted from (43.1) satisfy
      \[
       \sup_n\int
       \bigl(\|\partial g\|_{L^{4+2\delta}}^{4+2\delta}
       +\|a\|_{H^1}^2+\|\phi\|_{H^1}^2\bigr)\,\dd\nu_n
       <\infty.                                          \tag{43.9}
      \]
\end{enumerate}
Then
\[
 \sup_n\int V_\delta\,\dd\nu_n<\infty                  \tag{43.10}
\]
and the family \(\{\nu_n\}\) is tight in the carrier (42.3).
More precisely,
\[
 \nu_n(V_\delta>R)\leq \frac{C}{R}.                     \tag{43.11}
\]
\end{theorem}

\begin{proof}
The exponential moment (43.7) bounds every polynomial moment of
\(E_2\) uniformly.  Indeed, for \(x\geq0\),
\[
 x^q\leq C_{q,\eta}e^{\eta x}.                          \tag{43.12}
\]
Proposition \ref{prop:v15v43-local-global} and (43.8) therefore
give a uniform expectation of \(E_{2+\delta}\).  The harmonic
coordinate estimate (42.15), with (43.9), gives a uniform
\(W^{2,2+\delta}\) metric moment.  Adding the two field moments
proves (43.10).  Markov's inequality gives (43.11), and compactness
of the \(V_\delta\) sublevels from Theorem
\ref{thm:v15v42-compact} proves tightness.
\end{proof}

\begin{corollary}[When exponential tightness is recovered]
\label{cor:v15v43-exp-tight}
If, more strongly,
\[
 \sup_n\int e^{\eta_\delta V_\delta}\,\dd\nu_n<\infty    \tag{43.13}
\]
for some \(\eta_\delta>0\), then
\[
 \nu_n(V_\delta>R)\leq Ce^{-\eta_\delta R}.              \tag{43.14}
\]
\end{corollary}

\begin{proof}
Apply Markov's inequality to \(e^{\eta_\delta V_\delta}\).
\end{proof}

Condition (43.7) alone generally supplies polynomial moments of
\(E_2\), not the superlinear exponential moment
\(\exp(cE_2^q)\) needed to deduce (43.13) from (43.5).  Ordinary
tightness in Theorem \ref{thm:v15v43-prob-tight} is therefore the
correct conclusion.

\section{Critical control alone fails}

\begin{countertheorem}[A critical-energy family can escape the
coefficient topology]
\label{cth:v15v43-critical-random}
Let \(g_k=e^{2\varepsilon\psi(kx)}g_0\) be the metrics of
Countertheorem \ref{cth:v15v42-critical}, and let
\(\nu_k=\delta_{g_k}\).  Then, after choosing the coordinate patch
small and \(\varepsilon\) fixed,
\[
 \sup_k E_2(g_k)<\infty,\qquad
 \sup_k\int e^{\eta E_2(g)}\,\dd\nu_k<\infty             \tag{43.15}
\]
for every fixed \(\eta>0\), but
\[
 E_{2+\delta}(g_k)\asymp k^{2\delta}\longrightarrow\infty
                                                               \tag{43.16}
\]
and \(\{\nu_k\}\) is not tight in a uniform-coefficient topology.
\end{countertheorem}

\begin{proof}
For the conformal perturbation, curvature contains a term linear in
\(D^2\sigma_k\) and lower-order quadratic gradient terms.  The
scaling relation (42.18) makes the \(L^2\) norm of the linear term
bounded, while all lower terms are bounded for small fixed
\(\varepsilon\).  For \(p=2+\delta\),
\[
 \|D^2\psi(k\cdot)\|_{L^p}^p
 =k^{2p-4}\|D^2\psi\|_{L^p}^p
 =k^{2\delta}\|D^2\psi\|_{L^p}^p.                       \tag{43.17}
\]
Choosing \(\psi\) so its linear curvature contribution does not
vanish gives (43.16) after reducing \(\varepsilon\) if necessary.
The lack of a uniformly convergent metric subsequence was proved in
Countertheorem \ref{cth:v15v42-critical}; Dirac masses supported on
a nonprecompact set are not tight.
\end{proof}

Thus even a uniform exponential moment of the critical energy does
not manufacture the reverse-Hölder row.

\section{A conditional Gibbs route}

One possible physical mechanism is a local conditional estimate.
Let \({\cal F}_{j}^{\rm out}\) be the sigma algebra outside \(2B_j\).
If the interacting law obeys, uniformly in \(n,j\),
\[
 \mathbb E_{\nu_n}\!\left[
 \left.
 \left(\fint_{B_j}|{\rm Rm}|^{2q}\right)^{1/q}
 \right|{\cal F}_{j}^{\rm out}\right]
 \leq
 C_{\rm RH}\fint_{2B_j}|{\rm Rm}|^2+d_{j,n},             \tag{43.18}
\]
with the defect moments in (43.8), then taking expectations and
repeating (43.6) gives the moment conclusion of Theorem
\ref{thm:v15v43-prob-tight}.  This route does not require every
sample to satisfy a deterministic elliptic equation.  Proving
(43.18), however, requires an actual conditional law with its
fermion and quotient weights; a classical Euler--Lagrange estimate
cannot be substituted for it.

\begin{assessment}[What is gained]
\label{ass:v15v43-gain}
V43 converts the supercritical deterministic demand of V42 into a
local probabilistic acquisition: a reverse-Hölder row plus a defect
moment.  Once that row lands, the critical exponential energy tail
already yields the polynomial moments needed for ordinary
tightness.  The result is compatible with nonsolution path-integral
samples through the conditional form (43.18).
\end{assessment}

\begin{assessment}[Literal source applicability]
\label{ass:v15v43-source}
The source does not state a curvature reverse-Hölder estimate,
conditional or deterministic, for its interacting metric law.
Nor does it supply the controlled harmonic atlas or the field
moment (43.9).  The critical-to-supercritical bridge therefore
remains a named physical acquisition rather than a completed
continuum measure.
\end{assessment}

\begin{decision}[V44 identify subsequential limits by observables]
\label{dec:v15v43-next}
Assuming tightness on a common carrier, prove a measure
identification theorem from a countable convergence-determining
algebra of bounded local observables.  Include projective
consistency under refinement and show when all subsequences have the
same limit.  Keep OS positivity, Ward identities, and phase
transport as separate closed conditions on that algebra.
\end{decision}

\begin{firewall}[Claim boundary after V43]
\label{fire:v15v43}
V43 proves probabilistic tightness from a local reverse-Hölder row
and uniform moments, and proves that critical curvature energy alone
cannot supply coefficient compactness.  The physical conditional
estimate and hence the interacting continuum measure remain
unproved.
\end{firewall}

\section*{Version V43 append-only record}
\addcontentsline{toc}{section}{Version V43 append-only record}
The complete V42 source was retained before this continuation.  It
had \(429192\) bytes, \(12252\) lines, and SHA--256
\[
 \texttt{385146D01DD806ACB2522D390C43CB2D07943F4A51B4BDFFFDD4C8D902BF6EA1}.
\]
V43 proved a finite-cover reverse-Hölder compiler and a
probabilistic tightness theorem at nearly critical regularity, and
exhibited a critical-energy family whose coefficient topology
escapes through concentration.

\chapter{Fifteenth-cycle continuation V44: identifying the subsequential measure by a local observable algebra}
\label{sec:v15v44}

\section{Tightness is not identification}

V39--V43 give conditional routes to tightness of the positive
modulus measures on a common Polish carrier \({\cal X}\).  Tightness
only produces subsequential limits.  The sequence
\[
 \mu_n=\delta_{(-1)^n}\quad\text{on }\{-1,1\}             \tag{44.1}
\]
is tight and has two different subsequential limits.  To construct a
continuum law one must identify its expectations on a
measure-determining family.

Let
\[
 {\cal A}_0=\{F_1,F_2,\ldots\}\subset C_b({\cal X})       \tag{44.2}
\]
be a countable unital self-adjoint algebra that determines finite
Borel measures: equality of integrals on \({\cal A}_0\) implies
equality of measures.  On a locally compact second-countable carrier,
one may take a countable uniformly dense algebra in
\(C_0({\cal X})\), together with the constant function.  On a
general Polish carrier, a countable bounded-Lipschitz determining
class is sufficient.

\section{Transported observable theorem}

The regulator observable need not equal the restriction of a
continuum function.  Let
\[
 \widehat F_{j,n}:{\cal X}_n\longrightarrow\mathbb C     \tag{44.3}
\]
be its physical discretization and let
\(\iota_n:{\cal X}_n\to{\cal X}\) be the common-carrier map.

\begin{theorem}[Tightness plus determining moments]
\label{thm:v15v44-identification}
Let \(\nu_n\) be probability measures on \({\cal X}_n\).  Suppose:
\begin{enumerate}
\item the pushforwards
      \(\bar\nu_n=(\iota_n)_\#\nu_n\) are tight;
\item for every \(j\),
      \[
       \int_{{\cal X}_n}
       |\widehat F_{j,n}-F_j\circ\iota_n|\,\dd\nu_n
       \longrightarrow0;                                \tag{44.4}
      \]
\item for every \(j\), the limit
      \[
       \ell_j=\lim_{n\to\infty}
       \int_{{\cal X}_n}\widehat F_{j,n}\,\dd\nu_n       \tag{44.5}
      \]
      exists.
\end{enumerate}
Then there is a unique probability measure \(\nu\) on \({\cal X}\)
such that
\[
 \int_{\cal X}F_j\,\dd\nu=\ell_j
 \quad\text{for every }j,                               \tag{44.6}
\]
and the full sequence satisfies
\[
 \boxed{\bar\nu_n\Longrightarrow\nu.}                   \tag{44.7}
\]
\end{theorem}

\begin{proof}
Tightness and Prokhorov's theorem give a weakly convergent
subsequence \(\bar\nu_{n_k}\Rightarrow\nu\).  Since \(F_j\) is
bounded and continuous,
\[
 \int F_j\,\dd\bar\nu_{n_k}\longrightarrow
 \int F_j\,\dd\nu.                                      \tag{44.8}
\]
Equation (44.4) shows that the left side differs by \(o(1)\) from
the expectation in (44.5), proving (44.6).  Any other subsequential
limit has the same integrals \(\ell_j\), so the determining property
of \({\cal A}_0\) makes it equal to \(\nu\).  If the full sequence
did not converge to \(\nu\), some subsequence would remain outside a
weak neighborhood of \(\nu\); tightness would give that subsequence
a further weak limit, necessarily \(\nu\), a contradiction.
\end{proof}

\begin{corollary}[Uniform-on-compacts consistency is enough]
\label{cor:v15v44-compact-consistency}
Assume the tightness function \(V\) obeys
\[
 \sup_n\nu_n(V\circ\iota_n>R)\longrightarrow0
 \quad(R\to\infty),                                     \tag{44.9}
\]
the observables are uniformly bounded, and
\[
 \sup_{\{x:V(\iota_nx)\leq R\}}
 |\widehat F_{j,n}(x)-F_j(\iota_nx)|
 \longrightarrow0                                      \tag{44.10}
\]
for every fixed \(R,j\).  Then (44.4) holds.
\end{corollary}

\begin{proof}
Split the integral in (44.4) over the sublevel and its complement.
The first part tends to zero by (44.10).  The second is bounded by
twice the common supremum norm times the probability in (44.9).
First let \(n\to\infty\), then \(R\to\infty\).
\end{proof}

\section{Approximate refinement consistency}

Let \(R_{m,n}:{\cal X}_m\to{\cal X}_n\) be the declared coarse
restriction for \(m\geq n\).  Exact projective consistency is
\[
 (R_{m,n})_\#\nu_m=\nu_n.                               \tag{44.11}
\]
For an interacting renormalization flow one may only have
approximate consistency on local observables.

\begin{proposition}[Cofinal moment criterion]
\label{prop:v15v44-cofinal}
Suppose for every \(F_j\) there are coarse observables
\(\widehat F_{j,n}\) such that
\[
 \sup_{m\geq n}
 \left|
 \int\widehat F_{j,n}\circ R_{m,n}\,\dd\nu_m
 -\int\widehat F_{j,n}\,\dd\nu_n
 \right|
 \leq\epsilon_{j,n},\qquad\epsilon_{j,n}\to0,            \tag{44.12}
\]
and the interpolation defects (44.4) vanish uniformly for
\(m\geq n\) after coarse restriction.  Then the expectations in
(44.5) are Cauchy and hence converge.
\end{proposition}

\begin{proof}
For \(m\geq n\), add and subtract the two coarse expectations in
(44.12).  The consistency term is at most \(\epsilon_{j,n}\); the
two interpolation terms tend to zero uniformly by hypothesis.
Thus, for every \(\varepsilon>0\), all sufficiently large \(m,n\)
have expectation difference below \(\varepsilon\).
\end{proof}

Equation (44.12) must use the interacting measures, including the
fermion weight, quotient Jacobian, and counterterms.  Projectivity of
the unweighted renewal skeleton does not imply it.

\section{Locality of the determining algebra}

For field theory, \({\cal A}_0\) should be generated by bounded
local cylinder functions.  A typical element is
\[
 F(x)=f\bigl(
 \langle a,\omega_1\rangle,\ldots,
 \langle\phi,\zeta_r\rangle,
 \langle g,\kappa_1\rangle,\ldots\bigr),                 \tag{44.13}
\]
where \(f\) is bounded continuous and the test tensors have compact
support.  A countable dense family of rational test tensors and
bounded functions produces a countable algebra.

\begin{proposition}[Local cylinders determine the weak field law]
\label{prop:v15v44-cylinders}
On a separable Banach or Hilbert distribution carrier, bounded
cylinder functions built from a countable total family of continuous
linear functionals determine Borel probability measures.
\end{proposition}

\begin{proof}
Equality of their expectations gives equality of the joint
characteristic functions of every finite collection of the total
linear functionals.  Hence all finite-dimensional pushforwards
agree.  The sigma algebra generated by a countable total family is
the Borel sigma algebra of a separable Banach space, so the measures
agree on a generating pi-system and therefore everywhere.
\end{proof}

Gauge-invariant observables may fail to separate gauge orbits if the
chosen Wilson and matter grammar is incomplete.  The determining
claim then applies only to the sigma algebra it generates, not to
the full quotient law.

\section{Passing reflection positivity}

Let \(\Theta:{\cal X}\to{\cal X}\) be the continuum reflection and
\({\cal A}_{0,+}\subset{\cal A}_0\) the positive-time local
subalgebra.  Write
\[
 {\cal Q}_\nu(F)=
 \int_{\cal X}\overline{F(\Theta x)}F(x)\,\dd\nu(x).     \tag{44.14}
\]

\begin{theorem}[Closedness of OS positivity on the core]
\label{thm:v15v44-os}
Suppose \(\bar\nu_n\Rightarrow\nu\), the transported reflection
maps converge on every compact tightness sublevel, and for every
\(F\in{\cal A}_{0,+}\) there are bounded regulator
approximants \(\widehat F_n\) satisfying (44.4) for \(F\) and for
the reflected product.  If
\[
 {\cal Q}_{\nu_n}(\widehat F_n)\geq-\epsilon_n,
 \qquad\epsilon_n\to0,                                  \tag{44.15}
\]
then
\[
 {\cal Q}_{\nu}(F)\geq0.                                \tag{44.16}
\]
The result extends by polarization to every finite linear
combination in \({\cal A}_{0,+}\).
\end{theorem}

\begin{proof}
The reflected-product consistency and weak convergence give
\[
 {\cal Q}_{\nu_n}(\widehat F_n)\longrightarrow
 {\cal Q}_{\nu}(F).
\]
Passing to the limit in (44.15) proves (44.16).  Apply the same
argument to a finite linear combination, or equivalently to its
finite Gram matrix.
\end{proof}

Positivity of the modulus measure does not imply (44.15).  The
reflected determinant phase and fermionic boundary convention must
be included in the regulator OS form.

\section{Passing Ward identities}

Let \({\cal D}_X\) be the derivation on local observables generated
by a compactly supported infinitesimal gauge transformation or
diffeomorphism.  Suppose
\({\cal D}_{X,n}\widehat F_{j,n}\) approximates
\({\cal D}_XF_j\) in the sense of (44.4).

\begin{proposition}[Ward residual closure]
\label{prop:v15v44-ward}
If
\[
 \left|\int{\cal D}_{X,n}\widehat F_{j,n}\,\dd\nu_n\right|
 \leq\epsilon_{j,X,n},\qquad
 \epsilon_{j,X,n}\to0,                                  \tag{44.17}
\]
and the differentiated observables are uniformly integrable, then
\[
 \int{\cal D}_XF_j\,\dd\nu=0.                           \tag{44.18}
\]
\end{proposition}

\begin{proof}
Uniform integrability permits the same compact-sublevel truncation
as Corollary \ref{cor:v15v44-compact-consistency}.  On the compact
part, consistency and weak convergence pass the differentiated
observable to the limit.  Equation (44.17) then gives (44.18).
\end{proof}

The stress identity of V38 is the first-variation version of this
closure.  It remains separate from OS positivity.

\section{Complex phase transport}

Let
\[
 \dd\omega_n=
 ({\cal Z}_n^{\rm ph})^{-1}e^{i\Theta_n}\,\dd\nu_n       \tag{44.19}
\]
be normalized complex measures and assume
\(|{\cal Z}_n^{\rm ph}|\geq z_0>0\).

\begin{theorem}[Identification of a complex continuum measure]
\label{thm:v15v44-complex}
Suppose the total-variation measures \(|\omega_n|\) are tight and
uniformly bounded, and the limits
\[
 \lim_{n\to\infty}\int\widehat F_{j,n}\,\dd\omega_n
 =\ell_j^{\rm ph}                                       \tag{44.20}
\]
exist with the consistency (44.4) measured in total variation.
Then \(\omega_n\) converges weakly to a unique finite complex Borel
measure \(\omega\) determined by
\[
 \int F_j\,\dd\omega=\ell_j^{\rm ph}.                   \tag{44.21}
\]
\end{theorem}

\begin{proof}
Uniform boundedness and tightness of total variation give weak
relative compactness of the complex measures by applying
Prokhorov to the positive variation measures and weak compactness
to their bounded Radon--Nikodym densities.  Every subsequential
limit has moments (44.21).  A determining algebra for finite signed
measures determines real and imaginary parts separately, hence the
limit is unique.  The final subsequence argument is the same as in
Theorem \ref{thm:v15v44-identification}.
\end{proof}

This theorem transports a complex measure, not a positive Hilbert
space.  OS positivity still requires Theorem
\ref{thm:v15v44-os} for the complete reflected form.

\begin{assessment}[Continuum-measure compiler]
\label{ass:v15v44-compiler}
V44 closes the abstract step from tightness to a unique continuum
measure, provided local observable expectations, interpolation, and
refinement defects converge.  It also gives closed conditions for
OS positivity and Ward identities and a separate complex-phase
identification theorem.

The source does not populate a countable source-complete local
grammar with the uniform consistency (44.4), the interacting
refinement estimate (44.12), or the reflected positivity defect
(44.15).  No continuum measure is therefore claimed.
\end{assessment}

\begin{decision}[V45 compulsory audit]
\label{dec:v15v44-next}
At V45 inspect the newest YM and NS notes.  Test whether either now
supplies an interacting refinement, local-observable consistency,
or reflection-positive limit usable in (44.12)--(44.15).  Then
continue at V46 with the first missing physical row, prioritizing a
quantitative adjacent-cutoff comparison for bounded local
observables.
\end{decision}

\begin{firewall}[Claim boundary after V44]
\label{fire:v15v44}
V44 proves measure identification and limit-closure theorems under
explicit tightness, consistency, reflection, Ward, and phase
hypotheses.  Those physical hypotheses remain unproved; neither a
unique interacting SM--Einstein law nor an OS Hilbert space has yet
been constructed.
\end{firewall}

\section*{Version V44 append-only record}
\addcontentsline{toc}{section}{Version V44 append-only record}
The complete V43 source was retained before this continuation.  It
had \(439149\) bytes, \(12528\) lines, and SHA--256
\[
 \texttt{A101B4ED4ECE11E433D956CA5AA8AD73B947F73FDAF20B66E098DFFBCA9CACD7}.
\]
V44 proved uniqueness of the continuum measure from a countable
transported local algebra, an adjacent-refinement moment criterion,
and separate limit-closure theorems for OS positivity, Ward
identities, and the determinant phase.

\chapter{Fifteenth-cycle continuation V45: compulsory refinement audit}
\label{sec:v15v45}

\section{Companion checkpoints}

The active companion files at this checkpoint are
\[
\begin{array}{c|c|c|c}
\text{programme}&\text{file}&\text{lines}&\text{SHA--256}\\ \hline
\mathrm{YM}&
\texttt{Renewal\_Yang\_Mills\_Mass\_Gap\_Research\_Notes\_V39.tex}&
9076&
\texttt{9B44B0FBCA6037F3319E86036753B0F3B76BDD06F6045F00747E05DC01E23F9D}\\
\mathrm{NS}&
\texttt{Renewal\_NS\_Stability\_Research\_Notes\_V26.tex}&
5319&
\texttt{82C887F8C2C69E4F28FAD7C3DC8DE5B9099CB5A4BF431EBA1FFF3E91935978AD}.
\end{array}                                                \tag{45.1}
\]

YM V39 proves a finite-dimensional Wilson-pencil floor on a
two-direction quarter-face plate.  Its claim boundary explicitly
excludes an interacting continuum measure and a reflection-positive
limit.  It therefore provides no estimate of the interacting
refinement defect (44.12) or OS defect (44.15).

NS V26 remains a pointwise rank-one Oseen-kernel calculation.  It
contains no gauge--Higgs--metric carrier, refinement kernel, or
reflected field measure.  Its rank-one norm distinction was already
used in the trace-ideal direction of V35--V36.

\begin{theorem}[V45 nonimport and next target]
\label{thm:v15v45-audit}
Neither companion note populates any physical hypothesis of V44.
No positivity floor, kernel constant, or compactness statement is
imported.  The next target remains a direct adjacent-cutoff estimate
for bounded local observables under the full interacting weights.
\end{theorem}

\begin{proof}
The stated companion theorems have different carriers and both
exclude the interacting continuum conclusion needed here.  Hence
they cannot imply (44.12)--(44.15).  An adjacent-cutoff comparison
addresses the first missing condition without assuming the desired
continuum law.
\end{proof}

\begin{decision}[V46 adjacent-cutoff likelihood comparison]
\label{dec:v15v45-next}
Put two successive interacting measures on one fine carrier by a
declared coarse lift.  Bound every bounded-local expectation defect
by a coarse-map observable error plus the total variation distance
between the lifted likelihoods.  Control that distance by relative
entropy or an \(L^2\) likelihood ratio, and decompose the log
likelihood into bosonic, determinant, quotient, and phase rows.
\end{decision}

\begin{firewall}[Claim boundary after V45]
\label{fire:v15v45}
V45 is a compulsory audit and proves no new physical continuum
estimate.  The companion programmes remain logically separate.
\end{firewall}

\section*{Version V45 append-only record}
\addcontentsline{toc}{section}{Version V45 append-only record}
The complete V44 source was retained before this continuation.  It
had \(451686\) bytes, \(12872\) lines, and SHA--256
\[
 \texttt{A764C9B0BAC774DFA022F6D93DB9415C105574DD7462BA88D18605A3E1C1FA9A}.
\]
V45 performed the required companion audit and retained the direct
adjacent-cutoff likelihood comparison as the V46 target.

\chapter{Fifteenth-cycle continuation V46: adjacent-cutoff likelihood and local entropy comparison}
\label{sec:v15v46}

\section{Putting successive cutoffs on one carrier}

Let \({\cal X}_{n+1}\) be the finer configuration space and
\[
 R_n:{\cal X}_{n+1}\longrightarrow{\cal X}_n            \tag{46.1}
\]
the declared coarse restriction.  Let \(L_n(x,\dd y)\) be a
measurable lift kernel supported on
\(\{y:R_ny=x\}\).  The lifted coarse measure is
\[
 \widetilde\nu_n(\dd y)
 =\int_{{\cal X}_n}L_n(x,\dd y)\,\nu_n(\dd x),           \tag{46.2}
\]
so
\[
 (R_n)_\#\widetilde\nu_n=\nu_n.                         \tag{46.3}
\]
The kernel is part of the regulator comparison.  Choosing it after
examining an observable would not establish projective consistency.

Use the total variation convention
\[
 d_{\rm TV}(\mu,\nu)=\sup_A|\mu(A)-\nu(A)|
 =\frac12\left\|\frac{\dd\mu}{\dd\rho}
              -\frac{\dd\nu}{\dd\rho}\right\|_{L^1(\rho)}.
                                                               \tag{46.4}
\]

\begin{theorem}[Adjacent-cutoff observable bound]
\label{thm:v15v46-adjacent}
Let \(F_n:{\cal X}_n\to\mathbb C\) and
\(F_{n+1}:{\cal X}_{n+1}\to\mathbb C\) satisfy
\(\|F_n\|_\infty\leq M\).  Then
\[
 \boxed{
 \begin{split}
 &\left|
 \int F_{n+1}\,\dd\nu_{n+1}
 -\int F_n\,\dd\nu_n\right|\\
 &\quad\leq
 \int|F_{n+1}-F_n\circ R_n|\,\dd\nu_{n+1}
 +2M\,d_{\rm TV}(\nu_{n+1},\widetilde\nu_n).
 \end{split}}                                           \tag{46.5}
\]
If \(\nu_{n+1}\ll\widetilde\nu_n\), then Pinsker's inequality gives
\[
 \left|\mathbb E_{\nu_{n+1}}F_{n+1}
       -\mathbb E_{\nu_n}F_n\right|
 \leq e_{F,n}
 +M\sqrt{2\,D_{\rm KL}(\nu_{n+1}\Vert\widetilde\nu_n)}, \tag{46.6}
\]
where
\[
 e_{F,n}=
 \mathbb E_{\nu_{n+1}}
 |F_{n+1}-F_n\circ R_n|.                                \tag{46.7}
\]
\end{theorem}

\begin{proof}
Insert and subtract
\(\int F_n\circ R_n\,\dd\nu_{n+1}\).  The first difference is
bounded by (46.7).  By (46.3),
\[
 \int F_n\circ R_n\,\dd\widetilde\nu_n
 =\int F_n\,\dd\nu_n.                                   \tag{46.8}
\]
The dual characterization of total variation bounds the remaining
difference by \(2M d_{\rm TV}\), proving (46.5).  Pinsker's
inequality
\[
 d_{\rm TV}(\mu,\nu)
 \leq\sqrt{\frac12D_{\rm KL}(\mu\Vert\nu)}               \tag{46.9}
\]
gives (46.6).
\end{proof}

\section{Local entropy is the relevant quantity}

If \(F_n\) is supported in a region \(O\), let
\(\pi_{O,n+1}\) be the fine local readout on the portion needed to
evaluate \(F_n\circ R_n\).  Put
\[
 \nu_{n+1}^O=(\pi_{O,n+1})_\#\nu_{n+1},\qquad
 \widetilde\nu_n^O=(\pi_{O,n+1})_\#\widetilde\nu_n.      \tag{46.10}
\]

\begin{corollary}[Localized adjacent comparison]
\label{cor:v15v46-local}
In (46.5)--(46.6), the global total variation and entropy may be
replaced by those of the local marginals (46.10).  In particular,
\[
 \left|\mathbb E_{\nu_{n+1}}F_{n+1}
       -\mathbb E_{\nu_n}F_n\right|
 \leq e_{F,n}
 +M\sqrt{2D_{\rm KL}(\nu_{n+1}^O
                     \Vert\widetilde\nu_n^O)}.           \tag{46.11}
\]
\end{corollary}

\begin{proof}
The second term in (46.5) is the difference of expectations of a
function measurable with respect to \(\pi_{O,n+1}\), so it depends
only on the two local marginals.  Apply Pinsker to those marginals.
Equivalently, the data-processing inequality makes their entropy no
larger than the global entropy.
\end{proof}

This localization is essential in increasing volume.  A global
relative entropy may grow like the number of cells even when every
fixed local observable has a continuum limit.

\section{Likelihood-ratio criteria}

Assume
\(\nu_{n+1}\ll\widetilde\nu_n\) and put
\[
 L_n=\frac{\dd\nu_{n+1}}{\dd\widetilde\nu_n},\qquad
 \mathbb E_{\widetilde\nu_n}L_n=1.                      \tag{46.12}
\]
The chi-square estimate gives
\[
 d_{\rm TV}(\nu_{n+1},\widetilde\nu_n)
 \leq\frac12
 \left(\mathbb E_{\widetilde\nu_n}(L_n-1)^2\right)^{1/2}.
                                                               \tag{46.13}
\]

\begin{proposition}[Uniform log-likelihood oscillation]
\label{prop:v15v46-oscillation}
Suppose
\[
 L_n=\frac{e^{X_n}}
           {\mathbb E_{\widetilde\nu_n}e^{X_n}}          \tag{46.14}
\]
and there is a constant \(c_n\) such that
\[
 |X_n-c_n|\leq\epsilon_n
 \quad\widetilde\nu_n\text{-almost surely}.              \tag{46.15}
\]
Then
\[
 e^{-2\epsilon_n}\leq L_n\leq e^{2\epsilon_n},\qquad
 d_{\rm TV}(\nu_{n+1},\widetilde\nu_n)
 \leq\frac12(e^{2\epsilon_n}-1).                        \tag{46.16}
\]
\end{proposition}

\begin{proof}
The denominator in (46.14) lies between
\(e^{c_n-\epsilon_n}\) and \(e^{c_n+\epsilon_n}\), proving the
pointwise likelihood bounds.  Hence
\[
 \mathbb E_{\widetilde\nu_n}|L_n-1|
 \leq e^{2\epsilon_n}-1,
\]
and (46.4) proves the total variation estimate.
\end{proof}

\begin{proposition}[Subgaussian log-likelihood criterion]
\label{prop:v15v46-subgaussian}
Let \(Y_n=X_n-\mathbb E_{\widetilde\nu_n}X_n\).  If
\[
 \log\mathbb E_{\widetilde\nu_n}e^{tY_n}
 \leq\frac12\sigma_n^2t^2
 \quad\text{for }t=1,2,                                 \tag{46.17}
\]
then
\[
 \chi^2(\nu_{n+1}\Vert\widetilde\nu_n)
 =\mathbb E_{\widetilde\nu_n}(L_n-1)^2
 \leq e^{2\sigma_n^2}-1,                                \tag{46.18}
\]
and
\[
 d_{\rm TV}(\nu_{n+1},\widetilde\nu_n)
 \leq\frac12\sqrt{e^{2\sigma_n^2}-1}.                   \tag{46.19}
\]
\end{proposition}

\begin{proof}
Direct calculation gives
\[
 1+\chi^2
 =\frac{\mathbb E e^{2X_n}}
        {(\mathbb E e^{X_n})^2}.                         \tag{46.20}
\]
The numerator is at most
\(e^{2\mathbb EX_n+2\sigma_n^2}\) by (46.17), while Jensen's
inequality makes the denominator at least
\(e^{2\mathbb EX_n}\).  This proves (46.18), and (46.13) proves
(46.19).
\end{proof}

\section{Interacting log-likelihood ledger}

Suppose the lifted and fine modulus measures are written against one
fine reference measure as
\[
 \dd\widetilde\nu_n=
 \widetilde Z_n^{-1}e^{-\widetilde{\cal A}_n}\,\dd\rho_n,
 \qquad
 \dd\nu_{n+1}=
 Z_{n+1}^{-1}e^{-{\cal A}_{n+1}}\,\dd\rho_n.             \tag{46.21}
\]
Then (46.14) holds with
\[
 X_n=-( {\cal A}_{n+1}-\widetilde{\cal A}_n).            \tag{46.22}
\]
Using (39.3), its unnormalized action increment is
\[
 \begin{split}
 {\cal A}_{n+1}-\widetilde{\cal A}_n
={}&\Delta S_{\rm bos}+\Delta S_{\rm ct}
 -\frac12\Delta\operatorname{Tr}K\\
 &-\Delta\log J^{\rm quot}.                             \tag{46.23}
 \end{split}
\]
Every term is evaluated after the same coarse lift.  The normalized
fermion convention and fixed heat/resolvent counterterm scheme of
V36--V38 are required in the third term.

\begin{theorem}[Four-row adjacent compiler]
\label{thm:v15v46-four-row}
If centered versions of the four terms in (46.23) are jointly
subgaussian and
\[
 \log\mathbb E_{\widetilde\nu_n}
 \exp\left[
 t\left(X_n-\mathbb E_{\widetilde\nu_n}X_n\right)\right]
 \leq\frac12\sigma_n^2t^2
 \quad(t=1,2),                                          \tag{46.24}
\]
then every bounded local observable obeys (46.11) with its local
entropy term bounded by
\[
 M\sqrt{e^{2\sigma_{O,n}^2}-1},                         \tag{46.25}
\]
where \(\sigma_{O,n}\) is the corresponding local-marginal
log-likelihood parameter.
\end{theorem}

\begin{proof}
Apply Proposition \ref{prop:v15v46-subgaussian} to the local
likelihood ratio and insert (46.19) in Theorem
\ref{thm:v15v46-adjacent}.  The decomposition (46.23) identifies
the quantities whose joint exponential moment is required.
\end{proof}

Separate variance bounds on the four rows imply a joint bound by
Holder's inequality only after accounting for their dependence.
Adding four individual variances as if the rows were independent is
not valid for the interacting measure.

\section{Summable adjacent errors}

\begin{corollary}[Direct population of the V44 Cauchy row]
\label{cor:v15v46-summable}
For a bounded local observable family \(F_n\), suppose
\[
 \sum_{n=1}^\infty
 \left[
 e_{F,n}
 +M\sqrt{2D_{\rm KL}(\nu_{n+1}^O
                     \Vert\widetilde\nu_n^O)}
 \right]<\infty.                                        \tag{46.26}
\]
Then \(\{\mathbb E_{\nu_n}F_n\}\) is Cauchy.  If this holds for
every member of the countable determining algebra of V44 and the
measures are tight, the full continuum measure is unique.
\end{corollary}

\begin{proof}
Sum (46.11) from \(n=N\) to \(m-1\).  The tail of the convergent
series tends to zero uniformly in \(m>N\), proving the Cauchy
property.  Theorem \ref{thm:v15v44-identification} gives uniqueness.
\end{proof}

Square-summability of entropies alone is generally insufficient
because Pinsker introduces their square roots.  The required
summability rate must be checked rather than inferred from
termwise convergence.

\section{Phase comparison}

Let the fine and lifted complex measures have phases
\(\Theta_{n+1}\) and \(\widetilde\Theta_n\), and suppose their phase
partition functions have modulus at least \(z_0>0\).  Define
\[
 \epsilon_n^{\rm ph}
 =\int|e^{i\Theta_{n+1}}-e^{i\widetilde\Theta_n}|\,
       \dd\widetilde\nu_n.                               \tag{46.27}
\]

\begin{proposition}[Adjacent complex-measure bound]
\label{prop:v15v46-phase}
For every \(|F|\leq M\), with
\(\delta_n=d_{\rm TV}(\nu_{n+1},\widetilde\nu_n)\),
\[
 \left|
 \mathbb E_{\omega_{n+1}}F-
 \mathbb E_{\widetilde\omega_n}F
 \right|
 \leq
 M(2\delta_n+\epsilon_n^{\rm ph})
 \left(z_0^{-1}+z_0^{-2}\right).                        \tag{46.28}
\]
\end{proposition}

\begin{proof}
The difference of the unnormalized phase-weighted numerators is at
most \(M(2\delta_n+\epsilon_n^{\rm ph})\), by first changing the
positive measure and then the phase.  The same estimate with
\(M=1\) bounds the difference of phase partition functions.  For
numerators \(N_1,N_0\) and denominators \(Z_1,Z_0\),
\[
 \left|\frac{N_1}{Z_1}-\frac{N_0}{Z_0}\right|
 \leq\frac{|N_1-N_0|}{z_0}
 +\frac{|N_0|\,|Z_1-Z_0|}{z_0^2},
\]
and \(|N_0|\leq M\), proving (46.28).
\end{proof}

\begin{assessment}[Physical status]
\label{ass:v15v46-status}
V46 reduces interacting projective consistency to three explicit
quantities: observable interpolation error, local likelihood
entropy, and phase transport.  Its four-row decomposition shows
where the bosonic action, renormalized determinant, quotient
Jacobian, and counterterms enter.

The source does not specify a coarse lift \(L_n\), a common fine
reference law, a local entropy rate, or a summable phase defect.
Consequently (46.26) has not been populated for the physical
regulator.
\end{assessment}

\begin{decision}[V47 interpolate the adjacent action]
\label{dec:v15v46-next}
Introduce a path
\({\cal A}_{n,\theta}=\widetilde{\cal A}_n+
\theta({\cal A}_{n+1}-\widetilde{\cal A}_n)\).
Use differentiation of log partition functions to express the
relative entropy through the variance of the action increment.
Prove a local entropy bound from a uniform Poincare or logarithmic
Sobolev estimate for the interpolating laws, and identify the
determinant-response norm needed in that variance.
\end{decision}

\begin{firewall}[Claim boundary after V46]
\label{fire:v15v46}
V46 proves local adjacent-cutoff comparison and a summable
consistency compiler.  No physical lift, local entropy rate,
interpolating functional inequality, or phase comparison is yet
available for the complete SM--Einstein law.
\end{firewall}

\section*{Version V46 append-only record}
\addcontentsline{toc}{section}{Version V46 append-only record}
The complete V45 source was retained before this continuation.  It
had \(454679\) bytes, \(12945\) lines, and SHA--256
\[
 \texttt{2EFB59B399099769FFB7DDD20880439083D4538C604649F2CDE83B1DC917ACBC}.
\]
V46 proved adjacent-cutoff bounds by total variation, local relative
entropy, likelihood oscillation, and subgaussian increments; it
decomposed the full interacting likelihood and gave a summable
criterion that directly yields convergence of every determining
local moment.

\chapter{Fifteenth-cycle continuation V47: entropy as integrated action-increment variance}
\label{sec:v15v47}

\section{The adjacent Gibbs path}

Retain the fine carrier and lifted coarse measure of V46.  Write
\[
 \dd\nu_0=\dd\widetilde\nu_n,\qquad
 \Delta_n={\cal A}_{n+1}-\widetilde{\cal A}_n.           \tag{47.1}
\]
For \(0\leq\theta\leq1\), define
\[
 \dd\nu_\theta=
 \exp[-\theta\Delta_n-\psi_n(\theta)]\,\dd\nu_0,
 \qquad
 \psi_n(\theta)=
 \log\int e^{-\theta\Delta_n}\,\dd\nu_0.                \tag{47.2}
\]
Thus \(\nu_1=\nu_{n+1}\), provided both endpoint actions are written
against the same fine reference measure.

\begin{lemma}[Log-partition derivatives]
\label{lem:v15v47-logpartition}
If \(\Delta_n^2e^{-\theta\Delta_n}\) is integrable uniformly on
compact subintervals of \([0,1]\), then
\[
 \psi_n'(\theta)=-\mathbb E_{\nu_\theta}\Delta_n,\qquad
 \psi_n''(\theta)=
 \operatorname{Var}_{\nu_\theta}(\Delta_n).             \tag{47.3}
\]
\end{lemma}

\begin{proof}
Differentiate the integral in (47.2) twice under the sign of
integration.  The first derivative is the normalized first moment
of \(-\Delta_n\).  Differentiating that quotient gives the centered
second moment.
\end{proof}

\section{Exact entropy identity}

\begin{theorem}[Thermodynamic interpolation identity]
\label{thm:v15v47-entropy-identity}
Under Lemma \ref{lem:v15v47-logpartition},
\[
 \boxed{
 D_{\rm KL}(\nu_1\Vert\nu_0)
 =
 \int_0^1\theta\,
 \operatorname{Var}_{\nu_\theta}(\Delta_n)\,\dd\theta,} \tag{47.4}
\]
and
\[
 \boxed{
 D_{\rm KL}(\nu_0\Vert\nu_1)
 =
 \int_0^1(1-\theta)\,
 \operatorname{Var}_{\nu_\theta}(\Delta_n)\,\dd\theta.} \tag{47.5}
\]
In particular, if
\[
 \sup_{0\leq\theta\leq1}
 \operatorname{Var}_{\nu_\theta}(\Delta_n)\leq v_n^2,   \tag{47.6}
\]
then both relative entropies are at most \(v_n^2/2\).
\end{theorem}

\begin{proof}
From (47.2),
\[
 \log\frac{\dd\nu_1}{\dd\nu_0}
 =-\Delta_n-\psi_n(1).
\]
Therefore
\[
 D_{\rm KL}(\nu_1\Vert\nu_0)
 =-\mathbb E_{\nu_1}\Delta_n-\psi_n(1)
 =\psi_n'(1)-\psi_n(1),                                 \tag{47.7}
\]
because \(\psi_n(0)=0\).  Integration by parts gives
\[
 \int_0^1\theta\psi_n''(\theta)\,\dd\theta
 =\psi_n'(1)-\psi_n(1)+\psi_n(0),
\]
which proves (47.4).  Similarly,
\[
 D_{\rm KL}(\nu_0\Vert\nu_1)
 =\mathbb E_{\nu_0}\Delta_n+\psi_n(1)
 =-\psi_n'(0)+\psi_n(1)
 =\int_0^1(1-\theta)\psi_n''(\theta)\,\dd\theta.
\]
The last claim follows by integrating the weights \(\theta\) and
\(1-\theta\).
\end{proof}

\begin{corollary}[Variance rate for local moments]
\label{cor:v15v47-moment-rate}
For a bounded local observable as in V46,
\[
 \left|\mathbb E_{\nu_{n+1}}F_{n+1}
 -\mathbb E_{\nu_n}F_n\right|
 \leq e_{F,n}+Mv_{O,n},                                 \tag{47.8}
\]
provided (47.6) holds with \(v_{O,n}\) for a comparison whose
outside marginal agrees and whose action increment is the local
conditional increment relevant to \(O\).
\end{corollary}

\begin{proof}
Theorem \ref{thm:v15v47-entropy-identity} gives local relative
entropy at most \(v_{O,n}^2/2\).  Substitute it in (46.11).
\end{proof}

The outside-marginal qualification is essential.  A global
interpolation variance can be volume extensive and need not control
a summable local refinement rate.

\section{Poincare control of the variance}

Let the fine configuration coordinates carry a declared carré du
champ \(\Gamma_n\).  A probability law \(\nu_\theta\) has Poincare
constant \(C_{P,n}(\theta)\) if
\[
 \operatorname{Var}_{\nu_\theta}(f)
 \leq C_{P,n}(\theta)
 \int\Gamma_n(f)\,\dd\nu_\theta                         \tag{47.9}
\]
on a common core.

\begin{theorem}[Poincare-to-entropy compiler]
\label{thm:v15v47-poincare}
If
\[
 \sup_{\theta\in[0,1]}C_{P,n}(\theta)\leq C_{P,n},
 \qquad
 \sup_{\theta\in[0,1]}
 \int\Gamma_n(\Delta_n)\,\dd\nu_\theta\leq G_n^2,        \tag{47.10}
\]
then
\[
 D_{\rm KL}(\nu_1\Vert\nu_0)
 \leq\frac12C_{P,n}G_n^2                                \tag{47.11}
\]
and
\[
 d_{\rm TV}(\nu_1,\nu_0)
 \leq\frac12\sqrt{C_{P,n}}\,G_n.                        \tag{47.12}
\]
\end{theorem}

\begin{proof}
Apply (47.9) to \(f=\Delta_n\) and use (47.10) in (47.4).
Pinsker's inequality gives (47.12).
\end{proof}

\begin{corollary}[Summable adjacent Poincare criterion]
\label{cor:v15v47-summable}
If for every determining local observable
\[
 \sum_n\left(e_{F,n}
 +M\sqrt{C_{P,O,n}}\,G_{O,n}\right)<\infty,             \tag{47.13}
\]
then its expectations converge.  With tightness, the continuum
measure is unique by V44.
\end{corollary}

\begin{proof}
Combine (47.8), (47.10), and Corollary
\ref{cor:v15v46-summable}.
\end{proof}

\section{Logarithmic Sobolev strengthening}

Use the convention that an LSI with constant \(C_{\rm LS}\) is
\[
 \operatorname{Ent}_\nu(f^2)
 \leq2C_{\rm LS}\int\Gamma(f)\,\dd\nu.                  \tag{47.14}
\]

\begin{proposition}[Herbst likelihood bound]
\label{prop:v15v47-herbst}
If every \(\nu_\theta\) obeys (47.14) with a common constant
\(C_{{\rm LS},n}\), and
\[
 \Gamma_n(\Delta_n)\leq L_n^2
 \quad\text{almost surely},                             \tag{47.15}
\]
then \(\Delta_n-\mathbb E_{\nu_\theta}\Delta_n\) is
subgaussian:
\[
 \log\mathbb E_{\nu_\theta}
 e^{t(\Delta_n-\mathbb E_{\nu_\theta}\Delta_n)}
 \leq\frac12C_{{\rm LS},n}L_n^2t^2.                    \tag{47.16}
\]
Consequently Proposition \ref{prop:v15v46-subgaussian} applies with
\(\sigma_n^2=C_{{\rm LS},n}L_n^2\).
\end{proposition}

\begin{proof}
Apply (47.14) to
\(f=e^{t\Delta_n/2}\).  Dividing the resulting differential
inequality by the exponential moment and integrating in \(t\) is
the Herbst argument, which gives (47.16).  The final statement is
the definition of the subgaussian parameter used in V46.
\end{proof}

\section{The four-row gradient ledger}

From (46.23),
\[
 \nabla\Delta_n=
 \nabla\Delta S_{\rm bos}
 +\nabla\Delta S_{\rm ct}
 -\frac12\nabla\Delta\operatorname{Tr}K
 -\nabla\Delta\log J^{\rm quot}.                        \tag{47.17}
\]
For a quadratic carré du champ,
\[
 \Gamma_n(\Delta_n)
 \leq4\left[
 \Gamma_n(\Delta S_{\rm bos})
 +\Gamma_n(\Delta S_{\rm ct})
 +\frac14\Gamma_n(\Delta\operatorname{Tr}K)
 +\Gamma_n(\Delta\log J^{\rm quot})
 \right].                                               \tag{47.18}
\]

\begin{proposition}[Determinant-gradient insertion]
\label{prop:v15v47-det-gradient}
Let \(z\) be a fine configuration coordinate.  Under the \(C^1\)
trace-ideal hypotheses of V38,
\[
 \begin{split}
 \partial_z\operatorname{Tr}K_n
 =\int_0^\infty\operatorname{Tr}\bigl[
 &(B_n+t)^{-1}B_{n,z}(B_n+t)^{-1}\\
 &-(C_n+t)^{-1}C_{n,z}(C_n+t)^{-1}
 -{\cal E}_{n,z}(t)\bigr]\,\dd t.                       \tag{47.19}
 \end{split}
\]
If the coordinate carré du champ is
\(\Gamma_n(f)=\sum_z w_{z,n}|\partial_zf|^2\), then the fermion row
needed in (47.18) is exactly
\[
 \sum_z w_{z,n}
 \left|\partial_z\Delta\operatorname{Tr}K\right|^2.     \tag{47.20}
\]
\end{proposition}

\begin{proof}
Equation (47.19) is (38.12)--(38.13) in the coordinate direction
\(z\).  Substitution into the definition of the carré du champ
gives (47.20).
\end{proof}

Operator-norm locality of the determinant response does not by
itself bound (47.20): the coordinate sum is a Hilbert-space or
nuclear summation problem.  V31's boundary-area scaling and
V35's nuclear discipline are relevant to estimating it, but no
uniform physical bound was proved there.

\section{Conditional entropy localization}

Disintegrate the fine and lifted laws into an outside configuration
\(\xi\) and an inside configuration \(\zeta\) on a collar \(O^+\):
\[
 \nu_i(\dd\xi,\dd\zeta)
 =\nu_i^{\rm out}(\dd\xi)\nu_i(\dd\zeta\mid\xi),
 \qquad i=0,1.                                           \tag{47.21}
\]
The entropy chain rule is
\[
 \begin{split}
 D_{\rm KL}(\nu_1\Vert\nu_0)
 ={}&D_{\rm KL}(\nu_1^{\rm out}\Vert\nu_0^{\rm out})\\
 &+\int D_{\rm KL}\bigl(
 \nu_1(\cdot\mid\xi)\Vert\nu_0(\cdot\mid\xi)\bigr)
 \,\nu_1^{\rm out}(\dd\xi).                             \tag{47.22}
 \end{split}
\]
If the lift makes the outside marginals equal, the first term
vanishes.  Applying Theorem \ref{thm:v15v47-entropy-identity} to
each conditional Gibbs path gives a genuinely local entropy bound.
This is the appropriate place to insert the finite-block gap
compilers of V11, V19, and V27, provided their hypotheses hold
uniformly for every outside boundary condition.

\begin{assessment}[Gap inventory]
\label{ass:v15v47-gap}
Earlier versions proved conditional Poincare estimates for selected
Higgs and bounded-remainder blocks and exposed their failure in
unbounded gauge--Higgs or multiwell sectors.  They do not supply a
uniform Poincare or LSI constant for the complete interpolating law,
especially in the conformal metric direction.  The entropy identity
therefore identifies a use for those local gaps but does not promote
them to a full-measure gap.
\end{assessment}

\begin{assessment}[Literal source applicability]
\label{ass:v15v47-source}
The source does not define the common fine Gibbs path (47.2), a
conditional local lift with equal outside marginal, or uniform
Poincare constants along that path.  It also lacks the square-summed
determinant response (47.20).  Thus no numerical adjacent entropy
rate is currently available.
\end{assessment}

\begin{decision}[V48 shell-supported interpolation]
\label{dec:v15v47-next}
Exploit locality rather than a global gap.  Decompose the action
increment into a collar-supported head and a distant tail.  Prove a
conditional entropy bound for the head using a uniform block
Poincare constant, and debit the tail by the quasilocal
determinant/action estimates of V28, V31, and V34.  Derive the
collar growth required for summability across cutoffs.
\end{decision}

\begin{firewall}[Claim boundary after V47]
\label{fire:v15v47}
V47 proves the exact entropy--variance identity and Poincare/LSI
adjacent-cutoff compilers.  The physical local lift, full
interpolating gap, determinant-gradient sum, and summable rate
remain unproved.
\end{firewall}

\section*{Version V47 append-only record}
\addcontentsline{toc}{section}{Version V47 append-only record}
The complete V46 source was retained before this continuation.  It
had \(466578\) bytes, \(13318\) lines, and SHA--256
\[
 \texttt{94D305883B5B73D7D9DC8164E57628762853B71C242ECE9E6072B8C90B0C1CF9}.
\]
V47 expressed adjacent relative entropy exactly as integrated
action-increment variance, bounded it by local Poincare and LSI
estimates, and identified the square-summed renormalized determinant
response required in the complete interacting action gradient.

\chapter{Fifteenth-cycle continuation V48: shell-supported entropy and optimized collar growth}
\label{sec:v15v48}

\section{Localizing the adjacent action increment}

Let \(O\) support a bounded local observable and let \(O^R\) be its
metric \(R\)-collar on the fine regulator.  Use the conditional lift
of V47 so that the fine and lifted measures have the same outside
marginal.  For each outside configuration \(\xi\), decompose the
conditional action increment as
\[
 \Delta_{n,\xi}=H_{n,R,\xi}+T_{n,R,\xi},                 \tag{48.1}
\]
where the head \(H_{n,R,\xi}\) depends only on variables in
\(O^R\), while \(T_{n,R,\xi}\) is the cumulative contribution of
interactions or determinant responses reaching from \(O\) beyond
that collar.

Let \(\nu_{\theta,\xi}\) be the conditional interpolating Gibbs law
on \(O^R\).  All bounds below are uniform in the outside boundary
condition \(\xi\).  This uniformity is necessary before the
conditional estimates can be integrated over the outside law.

\section{Head--tail variance bound}

\begin{theorem}[Conditional shell entropy]
\label{thm:v15v48-shell-entropy}
Suppose, for every \(\theta\in[0,1]\) and almost every \(\xi\),
\[
 \operatorname{Var}_{\nu_{\theta,\xi}}(f)
 \leq C_{P,n}(R)
 \int\Gamma_{n,R}(f)\,\dd\nu_{\theta,\xi},              \tag{48.2}
\]
and
\begin{align}
 \int\Gamma_{n,R}(H_{n,R,\xi})\,\dd\nu_{\theta,\xi}
 &\leq G_n(R)^2,                                        \tag{48.3}\\
 \int|T_{n,R,\xi}|^2\,\dd\nu_{\theta,\xi}
 &\leq\tau_n(R)^2.                                      \tag{48.4}
\end{align}
Then the local-marginal relative entropy satisfies
\[
 \boxed{
 D_{\rm KL}(\nu_{n+1}^O\Vert\widetilde\nu_n^O)
 \leq C_{P,n}(R)G_n(R)^2+\tau_n(R)^2.}                  \tag{48.5}
\]
Consequently a bounded local observable obeys
\[
 \left|\mathbb E_{\nu_{n+1}}F_{n+1}
       -\mathbb E_{\nu_n}F_n\right|
 \leq e_{F,n}
 +M\sqrt{2C_{P,n}(R)G_n(R)^2+2\tau_n(R)^2}.             \tag{48.6}
\]
\end{theorem}

\begin{proof}
For arbitrary square-integrable random variables,
\[
 \operatorname{Var}(H+T)
 \leq2\operatorname{Var}(H)+2\operatorname{Var}(T)
 \leq2\operatorname{Var}(H)+2\mathbb E|T|^2.            \tag{48.7}
\]
Equations (48.2)--(48.4) therefore give
\[
 \operatorname{Var}_{\nu_{\theta,\xi}}(\Delta_{n,\xi})
 \leq2C_{P,n}(R)G_n(R)^2+2\tau_n(R)^2.                  \tag{48.8}
\]
Apply the entropy identity (47.4); integration of \(\theta\) over
\([0,1]\) contributes the factor \(1/2\), giving (48.5) for each
conditional law.  The entropy chain rule (47.22), with identical
outside marginals, integrates the bound over \(\xi\).  Data
processing gives the \(O\)-marginal inequality.  Finally insert
(48.5) in (46.11).
\end{proof}

\section{Tail from a summable quasilocal kernel}

Assume the metric graph has polynomial growth dimension \(d\):
\[
 \#\{z:r\leq d(z,O)<r+1\}\leq C_O(1+r)^{d-1}.           \tag{48.9}
\]
Suppose the distant action increment is a sum
\[
 T_{n,R}=\sum_{z:d(z,O)>R}t_{z,n}                       \tag{48.10}
\]
and, uniformly along the interpolating laws,
\[
 \|t_{z,n}\|_{L^2(\nu_{\theta,\xi})}
 \leq b_n(1+d(z,O))^{-N}.                               \tag{48.11}
\]

\begin{proposition}[Polynomial collar tail]
\label{prop:v15v48-tail}
If \(N>d\), then
\[
 \boxed{
 \tau_n(R)\leq
 \frac{C_Ob_n}{N-d}(1+R)^{d-N}.}                        \tag{48.12}
\]
\end{proposition}

\begin{proof}
Minkowski's inequality, (48.9), and (48.11) give
\[
 \|T_{n,R}\|_2
 \leq C_Ob_n
 \sum_{r>R}(1+r)^{d-1-N}.                               \tag{48.13}
\]
The integral test bounds the sum by
\((N-d)^{-1}(1+R)^{d-N}\).
\end{proof}

The smooth spectral multiplier theorem of V34 supplies arbitrarily
large fixed \(N\) at a fixed smooth energy window.  Its constant can
depend on the cutoff function and comparison operator.  That
dependence is part of \(b_n\) and must be controlled when the energy
window moves.

If an exponential kernel
\(\|t_{z,n}\|_2\leq b_ne^{-\alpha d(z,O)}\) is available instead,
the same proof gives
\[
 \tau_n(R)\leq C_{\alpha,d}b_n(1+R)^{d-1}e^{-\alpha R}. \tag{48.14}
\]

\section{Head growth}

Suppose \(O^R\) has at most \(C(1+R)^d\) fine blocks and the
per-block carré-du-champ contribution of the adjacent head is at
most \(a_n^2\).  Then
\[
 G_n(R)^2\leq C a_n^2(1+R)^d.                           \tag{48.15}
\]
If the conditional block Poincare constant is uniform,
\[
 C_{P,n}(R)\leq C_P,                                    \tag{48.16}
\]
Theorem \ref{thm:v15v48-shell-entropy} and Proposition
\ref{prop:v15v48-tail} give
\[
 \sqrt{2D_{\rm KL}(\nu_{n+1}^O\Vert\widetilde\nu_n^O)}
 \leq
 A\,a_n(1+R)^{d/2}
 +B\,b_n(1+R)^{-(N-d)}.                                 \tag{48.17}
\]

\begin{proposition}[Optimized polynomial collar]
\label{prop:v15v48-optimize}
Let \(N>d\), \(a_n>0\), and \(b_n>0\).  Up to constants independent
of \(n\), the two terms in (48.17) balance at
\[
 1+R_n\asymp
 \left(\frac{b_n}{a_n}\right)^{1/(N-d/2)},              \tag{48.18}
\]
provided \(b_n\geq a_n\).  At that radius,
\[
 \boxed{
 \sqrt{2D_{\rm KL}(\nu_{n+1}^O\Vert\widetilde\nu_n^O)}
 \leq
 C\,
 a_n^{\frac{N-d}{N-d/2}}
 b_n^{\frac{d/2}{N-d/2}}.}                              \tag{48.19}
\]
\end{proposition}

\begin{proof}
Set
\(p=d/2\) and \(q=N-d\).  The two powers in (48.17) are
\(a_nR^p\) and \(b_nR^{-q}\).  Equating them gives
\(R^{p+q}=b_n/a_n\), which is (48.18) because
\(p+q=N-d/2\).  Substitution gives
\[
 a_nR^p
 =a_n^{q/(p+q)}b_n^{p/(p+q)},
\]
and the tail has the same value.
\end{proof}

\begin{corollary}[Summability requirement]
\label{cor:v15v48-summability}
If \(b_n\leq b\), then the local entropy contribution is summable
whenever
\[
 \sum_n
 a_n^{\frac{N-d}{N-d/2}}<\infty.                        \tag{48.20}
\]
For \(a_n\leq Ch_n^\kappa\) and geometrically decreasing lattice
spacing \(h_n\leq C\rho^n\), \(0<\rho<1\), condition (48.20) holds
for every \(\kappa>0\).  For \(h_n\asymp n^{-1}\), it requires
\[
 \kappa\,\frac{N-d}{N-d/2}>1.                           \tag{48.21}
\]
\end{corollary}

\begin{proof}
Equation (48.19) and boundedness of \(b_n\) give (48.20).
The geometric series converges for every positive exponent.  The
polynomial case is the \(p\)-series test.
\end{proof}

\section{Boundary-extensive determinant response}

V31 showed that a connected two-core determinant trace can scale
with the boundary area of a block rather than its volume.  If the
head gradient is boundary supported, (48.15) improves to
\[
 G_n(R)^2\leq Ca_n^2(1+R)^{d-1}.                        \tag{48.22}
\]
Repeating Proposition \ref{prop:v15v48-optimize} replaces
\(d/2\) in (48.18)--(48.21) by \((d-1)/2\).  This improves the
summability exponent, but only after a connected-cluster
cancellation has been proved for the complete determinant,
counterterm, and quotient increment.  The two-core calculation
alone does not justify (48.22).

\section{Four physical tail rows}

The tail in (48.1) decomposes according to (46.23):
\[
 T_{n,R}=T_{n,R}^{\rm bos}
 +T_{n,R}^{\rm ct}
 -\frac12T_{n,R}^{\rm det}
 -T_{n,R}^{\rm quot}.                                   \tag{48.23}
\]
Minkowski gives
\[
 \tau_n(R)\leq
 \tau_n^{\rm bos}(R)+\tau_n^{\rm ct}(R)
 +\frac12\tau_n^{\rm det}(R)+\tau_n^{\rm quot}(R).      \tag{48.24}
\]
The bosonic finite-range action has zero tail once its interaction
range lies inside the collar.  A predeclared local counterterm bank
has the same property.  The determinant tail uses V34's smooth
functional calculus plus the trace/nuclear summation of V31,
V35, and V36.  The quotient tail requires a local gauge-fixing or
orbit-Jacobian theorem; it is not controlled by determinant
locality.

\begin{countertheorem}[Pointwise decay without uniform constants is
insufficient]
\label{cth:v15v48-nonuniform}
Let \(t_{z,n}=b_n(1+|z|)^{-N}\) on \(\mathbb Z^d\), with
\(b_n\to\infty\).  Every fixed \(n\) has a summable tail for
\(N>d\), but for any preassigned collar sequence \(R_n\),
\[
 \sum_{|z|>R_n}|t_{z,n}|
 \asymp b_nR_n^{d-N}.                                   \tag{48.25}
\]
Choosing \(b_n\geq R_n^{N-d}\) prevents the tail from tending to
zero.  Thus fixed-cutoff quasilocality alone does not populate a
cofinal entropy rate.
\end{countertheorem}

\begin{proof}
Comparison with the radial integral gives (48.25), and the stated
choice of \(b_n\) makes its right side bounded below.
\end{proof}

\begin{assessment}[Use of earlier gap results]
\label{ass:v15v48-gap-use}
The one-site and finite-color gap compilers of V9, V11, V19, and
V27 can supply (48.2) for selected conditional heads only where
their one-well, bounded-remainder, small-coupling, and boundary
uniformity assumptions hold.  They do not cover the metric
conformal direction or a determinant phase.  V48 therefore shows
how local gaps would enter without claiming a complete conditional
gap.
\end{assessment}

\begin{assessment}[Physical status]
\label{ass:v15v48-status}
V48 replaces a global interpolating gap by a conditional collar gap
and an explicit quasilocal tail.  It derives the exact tradeoff
between collar growth, adjacent head size, locality order, and
summability.  The source supplies none of the uniform constants
\(C_P,a_n,b_n,N\) for the complete lifted law, and no physical
coarse lift has yet landed.
\end{assessment}

\begin{decision}[V49 reflected collar consistency]
\label{dec:v15v48-next}
Apply the same head--tail decomposition to the finite OS Gram
matrices generated by positive-time local observables.  Prove that a
strict reflected Gram floor survives adjacent refinement when the
collar interpolation and phase defects are smaller than that floor.
Distinguish preservation on a finite determining core from
reflection positivity on its completed algebra.
\end{decision}

\begin{firewall}[Claim boundary after V48]
\label{fire:v15v48}
V48 proves shell-supported entropy bounds and optimized collar
rates.  The uniform physical conditional gap, determinant and
quotient tail constants, phase transport, and summable adjacent
rate remain unproved.
\end{firewall}

\section*{Version V48 append-only record}
\addcontentsline{toc}{section}{Version V48 append-only record}
The complete V47 source was retained before this continuation.  It
had \(477101\) bytes, \(13655\) lines, and SHA--256
\[
 \texttt{E9471DB7FEE59773F3C633B949166FDDFA38F6F9E49665B9B587853AEBF7852D}.
\]
V48 localized the adjacent action increment to a conditional
collar, bounded its entropy by a block gap plus a quasilocal tail,
optimized the collar radius, and derived the exact summability
condition along geometric and polynomial refinement schedules.

\chapter{Fifteenth-cycle continuation V49: reflected collar consistency and survival of finite OS floors}
\label{sec:v15v49}

\section{Finite reflected Gram matrices}

Let \(\Theta_n\) be the regulator reflection and let
\[
 {\bf F}_n=(F_{1,n},\ldots,F_{m,n})                     \tag{49.1}
\]
be bounded observables supported in the positive-time half.  For
the complete normalized measure \(\omega_n\), including its
determinant phase and fermionic boundary convention, define
\[
 (G_n)_{ij}=
 \int\overline{F_{i,n}(\Theta_nx)}F_{j,n}(x)\,
       \dd\omega_n(x).                                  \tag{49.2}
\]
Reflection positivity on this span is \(G_n\succeq0\), after
hermiticity has been proved.

Let \(\widetilde G_n\) be the matrix obtained after lifting the
cutoff-\(n\) law and observables to the cutoff-\(n+1\) carrier.
Write
\[
 E_n=G_{n+1}-\widetilde G_n.                            \tag{49.3}
\]

\begin{theorem}[Finite-core floor stability]
\label{thm:v15v49-floor}
Suppose \(\widetilde G_n=\widetilde G_n^*\) and
\[
 \widetilde G_n\succeq\gamma_nI_m,\qquad
 \|E_n\|_{\rm op}\leq\varepsilon_n.                     \tag{49.4}
\]
If \(G_{n+1}\) is Hermitian, then
\[
 \boxed{G_{n+1}\succeq(\gamma_n-\varepsilon_n)I_m.}     \tag{49.5}
\]
In particular, a positive floor survives whenever
\(\varepsilon_n<\gamma_n\).
\end{theorem}

\begin{proof}
For every \(c\in\mathbb C^m\),
\[
 c^*G_{n+1}c
 =c^*\widetilde G_nc+c^*E_nc
 \geq(\gamma_n-\|E_n\|_{\rm op})\|c\|_2^2.
\]
Use (49.4).
\end{proof}

\begin{corollary}[Accumulated floor]
\label{cor:v15v49-accumulated}
Suppose all lifted spaces are coherently identified,
\(G_{n_0}\succeq\gamma_0I_m\), and
\[
 \sum_{n=n_0}^\infty\|E_n\|_{\rm op}<\gamma_0.           \tag{49.6}
\]
Then every later Gram matrix and its limit satisfy
\[
 G_n\succeq
 \left(\gamma_0-\sum_{k=n_0}^\infty\|E_k\|_{\rm op}\right)I_m
 >0.                                                     \tag{49.7}
\]
\end{corollary}

\begin{proof}
Iterate Theorem \ref{thm:v15v49-floor} and pass to the matrix
limit, using norm convergence implied by (49.6).
\end{proof}

\section{From entrywise collar estimates to an operator bound}

Assume \(\|F_{i,n}\|_\infty\leq M_i\), and let
\(\eta_{ij,n}\) bound the adjacent expectation error for the
reflected product in (49.2), including observable, reflection,
modulus-likelihood, and phase defects.

\begin{proposition}[Gram defect from reflected-product defects]
\label{prop:v15v49-entry}
If
\[
 |(E_n)_{ij}|\leq\eta_{ij,n},                            \tag{49.8}
\]
then
\[
 \|E_n\|_{\rm op}
 \leq\left(\sum_{i,j=1}^m\eta_{ij,n}^2\right)^{1/2}
 \leq m\max_{i,j}\eta_{ij,n}.                           \tag{49.9}
\]
\end{proposition}

\begin{proof}
The operator norm is bounded by the Frobenius norm, which is the
first expression on the right.  The second inequality has \(m^2\)
terms.
\end{proof}

For a common bound \(M_i\leq M\), Proposition
\ref{prop:v15v46-phase} and Theorem
\ref{thm:v15v48-shell-entropy} give the schematic explicit row
\[
 \begin{split}
 \eta_{ij,n}\leq{}&
 e_{ij,n}^{\rm obs}+e_{ij,n}^{\rm refl}\\
 &+M^2(2\delta_{O^\Theta,n}+\epsilon_{O^\Theta,n}^{\rm ph})
 (z_0^{-1}+z_0^{-2}),                                   \tag{49.10}
 \end{split}
\]
where \(O^\Theta\) contains the support of \(F_j\), the reflected
support of \(F_i\), and the collar joining them, while
\[
 \delta_{O^\Theta,n}\leq
 \sqrt{\frac12\left[
 C_{P,n}(R)G_n(R)^2+\tau_n(R)^2\right]}.                \tag{49.11}
\]
Equations (49.9)--(49.11) turn the V48 collar estimates into the
quantity used in (49.5).

\section{Hermiticity is an independent row}

For a complex measure, (49.2) need not be Hermitian.  Define
\[
 H_n=\frac12(G_n+G_n^*),\qquad
 A_n=\frac1{2i}(G_n-G_n^*).                             \tag{49.12}
\]

\begin{proposition}[Hermiticity-defect closure]
\label{prop:v15v49-hermiticity}
If
\[
 \|A_n\|_{\rm op}\longrightarrow0,\qquad
 H_n\longrightarrow H
 \quad\text{in operator norm},                          \tag{49.13}
\]
and \(H_n\succeq-\epsilon_nI\) with \(\epsilon_n\to0\),
then \(G_n\to H=H^*\) and \(H\succeq0\).
\end{proposition}

\begin{proof}
Equation (49.12) gives
\(\|G_n-H\|\leq\|H_n-H\|+\|A_n\|\), proving convergence.
For every \(c\),
\[
 c^*Hc=\lim_n c^*H_nc\geq
 -\lim_n\epsilon_n\|c\|^2=0.
\]
\end{proof}

Thus a positive modulus and a small likelihood entropy do not
replace the reflected phase relation required to make
\(A_n\to0\).

\section{Null spaces and strict floors}

If \(G_n\succeq0\) but has null vectors, small perturbations can
produce negative eigenvalues.  A strict floor is therefore useful
on a fixed finite core.  It must not be extrapolated to the full
algebra.

\begin{countertheorem}[Positive semidefinite matrices have no
uniform perturbative reserve]
\label{cth:v15v49-null}
Let
\[
 G=\begin{pmatrix}1&0\\0&0\end{pmatrix},\qquad
 E_\epsilon=\begin{pmatrix}0&0\\0&-\epsilon\end{pmatrix}. \tag{49.14}
\]
Then \(G\succeq0\), \(\|E_\epsilon\|=\epsilon\to0\), but
\(G+E_\epsilon\) is not positive for any \(\epsilon>0\).
\end{countertheorem}

\begin{proof}
The second eigenvalue is \(-\epsilon\).
\end{proof}

For a semidefinite continuum target, the correct route is not a
uniform finite-cutoff strict floor.  It is the closedness theorem
(44.15)--(44.16): prove a negative defect tending to zero for each
finite polynomial, then quotient the limiting null space.

\section{Countable core and completion}

Let
\({\cal A}_{+,0}=\bigcup_{m\geq1}{\cal V}_m\) be an increasing
countable union of finite-dimensional positive-time observable
spaces.  Suppose the limiting OS form is
\[
 \langle F,G\rangle_{\rm OS}
 =\int\overline{F(\Theta x)}G(x)\,\dd\omega(x).          \tag{49.15}
\]

\begin{theorem}[From finite Grams to the OS pre-Hilbert space]
\label{thm:v15v49-completion}
Assume:
\begin{enumerate}
\item every finite Gram matrix on \({\cal V}_m\) is Hermitian and
      positive semidefinite;
\item the form (49.15) is consistent on intersections;
\item there is a normed positive-time algebra \({\cal A}_+\) in
      which \({\cal A}_{+,0}\) is dense and
      \[
       |\langle F,G\rangle_{\rm OS}|
       \leq C\|F\|_+\|G\|_+.                             \tag{49.16}
      \]
\end{enumerate}
Then (49.15) extends to a positive semidefinite continuous form on
\({\cal A}_+\).  Its null space
\[
 {\cal N}=\{F:\langle F,F\rangle_{\rm OS}=0\}            \tag{49.17}
\]
is a linear subspace, and the completion of
\({\cal A}_+/{\cal N}\) is a Hilbert space.
\end{theorem}

\begin{proof}
Finite Gram positivity gives
\(\langle F,F\rangle_{\rm OS}\geq0\) on the algebraic union.
Continuity (49.16) and density extend the form and its positivity to
\({\cal A}_+\).  Cauchy--Schwarz for a positive semidefinite form
follows by applying positivity to \(F+zG\).  Hence every
\(F\in{\cal N}\) is orthogonal to all \(G\), so \({\cal N}\) is
linear and the induced form on the quotient is positive definite.
Metric completion gives a Hilbert space.
\end{proof}

Positivity on one fixed finite determining family does not prove
item 1 for the union.  Every finite polynomial in the countable
positive-time core must be reached.

\section{Reflection-plane collar}

If a positive-time observable approaches the reflection plane, its
reflected product can have supports separated by only one lattice
spacing.  The collar \(O^\Theta\) in (49.10) must therefore include:
\begin{enumerate}
\item the positive support and its reflected copy;
\item all interaction terms crossing the reflection plane;
\item the fermionic hopping/boundary factorization;
\item quotient and counterterm terms meeting either support.
\end{enumerate}
The tail estimate begins outside this doubled collar.  A locality
bound centered only on the positive support misses the reflected
hopping terms and cannot establish (49.10).

\begin{assessment}[Physical status]
\label{ass:v15v49-status}
V49 proves quantitative survival of strict finite-core reflected
floors and the semidefinite limit/completion route.  The source does
not supply the complete normalized reflected complex measure, the
phase partition lower bound, the hermiticity defect, or the
summable doubled-collar estimates.  The current YM finite Wilson
floor is not an OS Gram floor and cannot be inserted into (49.4).
\end{assessment}

\begin{decision}[V50 compulsory audit]
\label{dec:v15v49-next}
At V50 inspect the newest YM and NS companions for a reflected
boundary factorization, conditional shell estimate, or completed
continuum measure.  Then continue at V51 by deriving a reflected
transfer-operator criterion that implies hermiticity and OS
positivity on the local core, with the chiral determinant phase and
zero-mode line kept explicit.
\end{decision}

\begin{firewall}[Claim boundary after V49]
\label{fire:v15v49}
V49 proves finite reflected-Gram stability and abstract OS
completion under explicit adjacent, phase, and hermiticity bounds.
Those physical bounds and a source-complete positive-time algebra
remain unproved.
\end{firewall}

\section*{Version V49 append-only record}
\addcontentsline{toc}{section}{Version V49 append-only record}
The complete V48 source was retained before this continuation.  It
had \(487668\) bytes, \(13964\) lines, and SHA--256
\[
 \texttt{7BC6419F79FBB94D969B4289EDFF6315D7C75A3E25EE9D3C5CA9EE3B44C175ED}.
\]
V49 converted reflected-product collar errors into Gram operator
bounds, proved survival of finite OS floors, separated hermiticity
and null-space issues, and constructed the OS pre-Hilbert completion
from a positive countable core.

\chapter{Fifteenth-cycle continuation V50: compulsory reflected-boundary audit}
\label{sec:v15v50}

\section{Checkpoints and decision}

The active companions remain
\[
\begin{array}{c|c|c|c}
\text{programme}&\text{file}&\text{lines}&\text{SHA--256}\\ \hline
\mathrm{YM}&
\texttt{Renewal\_Yang\_Mills\_Mass\_Gap\_Research\_Notes\_V39.tex}&
9076&
\texttt{9B44B0FBCA6037F3319E86036753B0F3B76BDD06F6045F00747E05DC01E23F9D}\\
\mathrm{NS}&
\texttt{Renewal\_NS\_Stability\_Research\_Notes\_V26.tex}&
5319&
\texttt{82C887F8C2C69E4F28FAD7C3DC8DE5B9099CB5A4BF431EBA1FFF3E91935978AD}.
\end{array}                                                \tag{50.1}
\]
YM V39 proves pointwise positivity for a finite momentum-space
Wilson pencil and explicitly excludes reflection-positive and
interacting continuum limits.  NS V26 proves pointwise rank-one
Oseen-kernel norms and contains no reflected field boundary.  Neither
supplies a fermionic hopping factorization, determinant-line
orientation, conditional shell estimate, or continuum OS measure.

\begin{theorem}[V50 nonimport]
\label{thm:v15v50-audit}
No companion theorem populates the hypotheses of V49.  The
reflected transfer factorization must be derived on the actual
SM--Einstein boundary carrier, with its phase and zero-mode data.
\end{theorem}

\begin{proof}
The companion carriers and conclusions described above do not
contain a reflected positive-time algebra or its Gram form.  Their
positive matrices therefore cannot be identified with (49.2).
\end{proof}

\begin{decision}[V51 reflected transfer criterion]
\label{dec:v15v50-next}
Factor the reflected boundary kernel through a Hilbert boundary
space.  Prove that a representation
\(K(\Theta x_+,y_+)=V(x_+)^*V(y_+)\) gives hermiticity and OS
positivity.  Extend the statement to fermionic exterior powers and
show precisely why a nonfactorable chiral phase or a changing
zero-mode determinant line obstructs the argument.
\end{decision}

\begin{firewall}[Claim boundary after V50]
\label{fire:v15v50}
V50 is a compulsory audit.  It imports no companion result and
proves no reflected SM--Einstein continuum statement.
\end{firewall}

\section*{Version V50 append-only record}
\addcontentsline{toc}{section}{Version V50 append-only record}
The complete V49 source was retained before this continuation.  It
had \(497199\) bytes, \(14250\) lines, and SHA--256
\[
 \texttt{B1DD8EF1057D7DE14C879A5B6C544E9E684DFBC3690B0A48055CC6DA94CC3730}.
\]
V50 performed the required companion audit and retained the
reflected transfer-factorization problem for V51.

\chapter{Fifteenth-cycle continuation V51: reflected boundary-kernel factorization and determinant-line obstruction}
\label{sec:v15v51}

\section{Boundary reduction of the OS form}

Fix a reflection hypersurface and write a regulated configuration as
\[
 x=(x_-,b,x_+),\qquad
 \Theta(x_-,b,x_+)=(\theta x_+,b,\theta x_-).            \tag{51.1}
\]
After integrating all variables strictly inside the two half-spaces,
suppose the reflected form of positive-time observables has the
boundary representation
\[
 {\cal Q}(F,G)=
 \int_{\Omega_+}\int_{\Omega_+}
 \overline{F(u)}\,K(u,v)\,G(v)\,
 \lambda(\dd u)\lambda(\dd v).                          \tag{51.2}
\]
Here \(u,v\) include the boundary data and the remaining
positive-half variables needed by the observables.  The complete
bosonic action, quotient Jacobian, fermion boundary amplitude,
counterterms, and determinant-line pairing are part of \(K\).

\section{The factorization criterion}

\begin{theorem}[Reflected boundary factorization]
\label{thm:v15v51-factorization}
Let \({\cal H}_\partial\) be a Hilbert space and suppose there is a
weakly measurable map
\[
 V:\Omega_+\longrightarrow{\cal H}_\partial             \tag{51.3}
\]
such that
\[
 K(u,v)=\langle V(u),V(v)\rangle_{{\cal H}_\partial}    \tag{51.4}
\]
and all integrals below exist.  Then \(K(v,u)=\overline{K(u,v)}\)
and
\[
 \boxed{
 {\cal Q}(F,F)=
 \left\|\int_{\Omega_+}F(v)V(v)\,\lambda(\dd v)
 \right\|_{{\cal H}_\partial}^2\geq0.}                  \tag{51.5}
\]
Thus the regulated law is reflection positive on every
positive-time observable for which the Bochner integral exists.
\end{theorem}

\begin{proof}
Conjugate symmetry of the Hilbert inner product gives hermiticity.
Substitute (51.4) into (51.2), use Bochner--Fubini, and recognize the
inner product of the vector integral with itself.
\end{proof}

\begin{theorem}[Converse Kolmogorov factorization]
\label{thm:v15v51-converse}
Suppose \(K\) is a Hermitian positive-definite kernel:
\[
 \sum_{i,j=1}^m\overline{c_i}c_jK(u_i,u_j)\geq0         \tag{51.6}
\]
for every finite family.  Then there is a Hilbert space
\({\cal H}_K\) and vectors \(V(u)\in{\cal H}_K\) satisfying
(51.4).
\end{theorem}

\begin{proof}
Take the complex vector space generated by formal symbols
\(\{e_u\}\) and define
\[
 \left\langle\sum_ic_ie_{u_i},\sum_jd_je_{v_j}\right\rangle
 =\sum_{i,j}\overline{c_i}d_jK(u_i,v_j).                \tag{51.7}
\]
Hermiticity makes this sesquilinear and (51.6) makes it positive
semidefinite.  Quotient its null space and complete.  The class
\(V(u)=[e_u]\) then satisfies (51.4).
\end{proof}

Hence boundary factorization is equivalent to positivity of every
finite reflected boundary Gram.  Positivity of a single matrix or a
finite test list is insufficient.

\section{Operator version}

If \(K\) defines a bounded integral operator
\({\mathsf K}\) on \(L^2(\lambda)\), then
\[
 {\cal Q}(F,F)=\langle F,{\mathsf K}F\rangle_{L^2}.      \tag{51.8}
\]

\begin{corollary}[Positive boundary transfer]
\label{cor:v15v51-operator}
The form (51.8) is reflection positive exactly when
\({\mathsf K}={\mathsf K}^*\succeq0\).  In that case
\[
 {\mathsf K}=({\mathsf K}^{1/2})^*{\mathsf K}^{1/2}     \tag{51.9}
\]
is the operator factorization.
\end{corollary}

\begin{proof}
Positivity of the quadratic form and the polarization identity imply
self-adjointness and positivity.  The positive square root follows
from the spectral theorem.  The converse is immediate.
\end{proof}

\section{Bosonic and quotient factors}

Suppose the boundary kernel is a pointwise or Schur product
\[
 K=K_{\rm bos}K_{\rm quot}K_{\rm f}.                    \tag{51.10}
\]
If each factor is positive definite, their product is positive
definite by the Schur product theorem on every finite boundary
matrix.  This provides a modular proof only when the three factors
occur on the same boundary carrier and with the same reflection.
A positive bulk density does not imply that its reduced boundary
kernel is positive definite.

For a quotient factor, positivity requires an actual orbit or slice
kernel.  A Faddeev--Popov determinant that changes sign cannot be
absorbed into a positive \(K_{\rm quot}\).  An absolute value changes
the theory unless it is part of the declared quotient measure.

\section{Fermionic exterior powers}

Let \(T:{\cal H}_1\to{\cal H}_1\) be a bounded positive one-particle
boundary transfer.  Its fermionic second quantization is
\[
 \Gamma_-(T)=
 \bigoplus_{k=0}^{\dim{\cal H}_1}\bigwedge\nolimits^kT   \tag{51.11}
\]
in finite dimension, with the evident Fock-space version when the
sum is controlled.

\begin{theorem}[Positive one-particle transfer gives a positive
fermion boundary kernel]
\label{thm:v15v51-fermion}
If \(T=T^*\succeq0\), then
\[
 \bigwedge\nolimits^kT\succeq0
 \quad\text{for every }k,\qquad
 \Gamma_-(T)\succeq0.                                   \tag{51.12}
\]
Moreover,
\[
 \bigwedge\nolimits^kT
 =
 \left(\bigwedge\nolimits^kT^{1/2}\right)^*
 \left(\bigwedge\nolimits^kT^{1/2}\right),              \tag{51.13}
\]
and on the top exterior power,
\[
 \bigwedge\nolimits^{\dim{\cal H}_1}T=\det(T)\geq0.      \tag{51.14}
\]
\end{theorem}

\begin{proof}
The exterior-power functor respects products and adjoints.  Since
\(T=(T^{1/2})^*T^{1/2}\),
\[
 \bigwedge\nolimits^kT
 =
 \bigwedge\nolimits^k(T^{1/2})^*
 \bigwedge\nolimits^kT^{1/2},
\]
which is (51.13) and is positive.  Direct sums of positive operators
are positive.  The top exterior action is multiplication by the
determinant.
\end{proof}

The theorem applies to a proved positive reflected hopping transfer.
Replacing a chiral operator \(D\) by \(D^*D\) proves positivity of
its modulus sector, not of the chiral determinant phase.

\section{Phase kernels}

Let \(K_0\) be a positive-definite boundary kernel and let
\[
 K(u,v)=q(u,v)K_0(u,v),\qquad |q(u,v)|=1.               \tag{51.15}
\]

\begin{proposition}[Positive phase criterion]
\label{prop:v15v51-phase}
If \(q\) is itself a positive-definite Hermitian kernel, then \(K\)
is positive definite.  In particular, if
\[
 q(u,v)=\overline{s(u)}s(v),\qquad |s(u)|=1,             \tag{51.16}
\]
then the phase is a boundary coboundary and
\[
 K(u,v)=
 \langle s(u)V_0(u),s(v)V_0(v)\rangle                  \tag{51.17}
\]
is a factorization.
\end{proposition}

\begin{proof}
On every finite set, the matrices of \(q\) and \(K_0\) are positive
semidefinite.  Their entrywise product is positive by the Schur
product theorem.  Equation (51.17) proves the coboundary case
directly.
\end{proof}

\begin{countertheorem}[Unit modulus does not preserve reflection
positivity]
\label{cth:v15v51-phase}
On three boundary points, let \(K_0(u_i,u_j)=1\) and
\[
 Q=(q(u_i,u_j))_{i,j=1}^3=
 \begin{pmatrix}
 1&-1&-1\\
 -1&1&-1\\
 -1&-1&1
 \end{pmatrix}.                                         \tag{51.18}
\]
Every entry has unit modulus and \(Q=Q^*\), but its eigenvalues are
\(2,2,-1\).  Hence the phase-twisted kernel \(K=Q\) is not
reflection positive.
\end{countertheorem}

\begin{proof}
The vector \((1,1,1)\) has eigenvalue \(-1\); the orthogonal
two-plane has eigenvalue \(2\).
\end{proof}

Thus hermiticity and absolute determinant positivity do not suffice.
The determinant phase must be shown to define a positive boundary
cocycle, a coboundary, or a larger fermionic factorization before
the Grassmann variables are integrated out.

\section{Determinant lines and reflection}

Suppose the chiral determinant at boundary datum \(u\) lies in a
complex line \({\cal L}_u\).  A scalar boundary kernel requires:
\begin{enumerate}
\item an anti-linear reflection isometry
      \[
       J_u:{\cal L}_u\longrightarrow{\cal L}_{\theta u}; \tag{51.19}
      \]
\item a parallel pairing of the reflected negative line with the
      positive line;
\item trivial holonomy, or a proved positive phase kernel, on every
      boundary loop used by the local algebra.
\end{enumerate}
Without these data, choosing a local scalar representative can
change by phases on overlaps and (51.4) is not globally defined.

\begin{proposition}[Line-valued factorization]
\label{prop:v15v51-line}
Let \({\cal L}\to\Omega_+\) be a Hermitian line bundle and suppose
there is a Hilbert space \({\cal H}_\partial\) and fibrewise linear
maps
\[
 W(u):{\cal L}_u\longrightarrow{\cal H}_\partial        \tag{51.20}
\]
such that the reflected boundary amplitude is
\[
 K(u,v)=W(u)^*W(v):
 {\cal L}_v\longrightarrow{\cal L}_u.                  \tag{51.21}
\]
Then the OS form on compactly supported line-valued observables is
positive.
\end{proposition}

\begin{proof}
For a section \(F(v)\in{\cal L}_v\), substitution of (51.21) gives
\[
 {\cal Q}(F,F)=
 \left\|\int W(v)F(v)\,\lambda(\dd v)\right\|^2\geq0.
\]
\end{proof}

This is the invariant version of (51.4).  A scalar determinant is
not needed if the line-valued boundary pairing is retained.

\section{Zero modes and changing exterior degree}

If \(T(u)\) has a zero mode, (51.12) remains positive semidefinite,
but its top exterior determinant vanishes.  Across a change of kernel
dimension, the reduced determinant belongs to
\[
 \det(\ker T(u))\otimes\det(\operatorname{coker}T(u))^*
                                                               \tag{51.22}
\]
together with the nonzero exterior product.  These line fibres can
change degree.

\begin{theorem}[Zero-mode alternative]
\label{thm:v15v51-zero}
Along a connected boundary stratum, exactly one of the following is
available:
\begin{enumerate}
\item the kernel dimension is constant and the kernel/cokernel
      determinant lines admit the reflected isometries and
      factorization (51.20)--(51.21);
\item the kernel dimension changes, in which case no fixed-degree
      scalar reduced determinant can define one continuous boundary
      kernel without a divisor transition map.
\end{enumerate}
A pseudo-determinant that simply drops each zero eigenvalue does not
supply the missing transition map.
\end{theorem}

\begin{proof}
Constant rank gives vector bundles of kernels and cokernels by the
standard spectral-projection construction, so their determinant
lines have fixed degree.  At a rank jump, the dimension of the top
exterior fibre changes.  Vector spaces of different exterior degree
have no canonical nonzero identification.  The vanishing eigenvalue
defines the divisor factor needed to relate the neighboring
strata; deleting it removes precisely that factor.
\end{proof}

\section{From one slab to a Hamiltonian}

Reflection positivity at one slab is not yet a continuum dynamics.
Let \({\mathsf K}_t\) be the positive boundary transfer across a
slab of thickness \(t\).

\begin{theorem}[Positive transfer semigroup]
\label{thm:v15v51-semigroup}
Suppose on the OS Hilbert space
\[
 {\mathsf K}_0=I,\qquad
 {\mathsf K}_{s+t}={\mathsf K}_s{\mathsf K}_t,\qquad
 {\mathsf K}_t={\mathsf K}_t^*,\quad
 0\preceq{\mathsf K}_t\preceq I,                        \tag{51.23}
\]
and \(t\mapsto{\mathsf K}_t\) is strongly continuous.  Then there is
a unique self-adjoint \(H\geq0\) such that
\[
 {\mathsf K}_t=e^{-tH}.                                 \tag{51.24}
\]
\end{theorem}

\begin{proof}
The strongly continuous contraction semigroup has a closed
dissipative generator \(-H\).  Self-adjointness of every
\({\mathsf K}_t\) makes the generator self-adjoint, and positivity
of the contractions gives spectrum in \([0,1]\), hence
\(H\geq0\).  The spectral theorem gives (51.24) and uniqueness.
\end{proof}

The mass spectrum, Lorentzian interpretation, and nontrivial
interacting fields require further spectral and reconstruction
steps after (51.24).

\begin{assessment}[Physical status]
\label{ass:v15v51-status}
V51 gives necessary and sufficient boundary-kernel positivity,
fermionic exterior-power positivity, a phase criterion with a
counterexample, the invariant determinant-line formulation, and
the positive transfer-semigroup handoff.

The source does not provide one complete boundary kernel containing
the gauge--Higgs--gravity action, quotient factor, chiral line,
zero-mode transition maps, and reflected hopping factorization on
the same carrier.  Its determinant modulus cannot be substituted
for Proposition \ref{prop:v15v51-phase}.  OS positivity and the
Hamiltonian therefore remain conditional.
\end{assessment}

\begin{decision}[V52 cofinal transfer-semigroup convergence]
\label{dec:v15v51-next}
Assume finite-cutoff reflected transfer operators exist.  Prove a
common-core theorem that strong convergence at a dense set of slab
times, uniform contraction, and approximate semigroup consistency
produce a limiting positive strongly continuous semigroup.  State
the additional compact-resolvent or spectral-projection estimates
needed to transport a mass gap without confusing the regulator
parameter with physical mass.
\end{decision}

\begin{firewall}[Claim boundary after V51]
\label{fire:v15v51}
V51 proves the abstract reflected boundary and transfer theorems.
No physical chiral phase factorization, zero-mode line transport,
cofinal OS semigroup, or Standard Model--Einstein Hamiltonian has
yet been constructed.
\end{firewall}

\section*{Version V51 append-only record}
\addcontentsline{toc}{section}{Version V51 append-only record}
The complete V50 source was retained before this continuation.  It
had \(499756\) bytes, \(14315\) lines, and SHA--256
\[
 \texttt{FEB708E5FF240E2BC89D3EC867A5070F12A00D26D0C616FED6E3EA48C77123C0}.
\]
V51 proved the reflected boundary factorization criterion, its
fermionic exterior-power and line-valued forms, the positive phase
criterion and obstruction, the zero-mode alternative, and the
positive transfer-semigroup reconstruction theorem.

\chapter{Fifteenth-cycle continuation V52: cofinal positive semigroups and the spectral-gap handoff}
\label{sec:v15v52}

\section{A common OS Hilbert carrier}

Assume the finite-cutoff reflected kernels of V51 have already been
transported to one Hilbert space \({\cal H}_{\rm OS}\).  Let
\[
 K_n(t):{\cal H}_{\rm OS}\to{\cal H}_{\rm OS},
 \qquad t\geq0,                                          \tag{52.1}
\]
be self-adjoint positive contractions.  They may have an
approximate, rather than exact, semigroup law because restriction,
interpolation, and counterterm matching occur between cutoffs.

\section{Strong semigroup convergence}

\begin{theorem}[Dense-time cofinal semigroup theorem]
\label{thm:v15v52-semigroup}
Let \(D\subset[0,\infty)\) be a countable dense additive set
containing zero.  Suppose:
\begin{enumerate}
\item \(K_n(0)=I\) and \(0\preceq K_n(t)\preceq I\);
\item for every \(t\in D\), \(K_n(t)\) converges strongly to an
      operator \(K(t)\);
\item the family is uniformly strongly equicontinuous: for every
      \(f\in{\cal H}_{\rm OS}\) and \(T<\infty\),
      \[
       \lim_{\delta\downarrow0}
       \sup_n\sup_{\substack{s,t\in[0,T]\\|s-t|\leq\delta}}
       \|(K_n(t)-K_n(s))f\|=0;                           \tag{52.2}
      \]
\item for every \(s,t\in D\) and \(f\in{\cal H}_{\rm OS}\),
      \[
       \|[K_n(s+t)-K_n(s)K_n(t)]f\|\longrightarrow0.     \tag{52.3}
      \]
\end{enumerate}
Then there is a unique strongly continuous semigroup
\(\{K(t)\}_{t\geq0}\) such that
\[
 K_n(t)\longrightarrow K(t)
 \quad\text{strongly, locally uniformly in \(t\)},       \tag{52.4}
\]
and
\[
 K(0)=I,\quad K(s+t)=K(s)K(t),\quad
 0\preceq K(t)\preceq I.                                \tag{52.5}
\]
Consequently there is a unique self-adjoint \(H\geq0\) with
\[
 K(t)=e^{-tH}.                                          \tag{52.6}
\]
\end{theorem}

\begin{proof}
Fix \(f\) and \(t\geq0\).  Choose \(t_j\in D\) with \(t_j\to t\).
Equation (52.2) makes \(K_n(t_j)f\) uniformly Cauchy in \(j\),
independently of \(n\).  Strong convergence on \(D\) therefore
defines \(K(t)f=\lim_jK(t_j)f\), independently of the approximating
sequence.  A three-epsilon argument using (52.2) gives (52.4) and
strong continuity.

Strong limits preserve the contraction bound.  For \(f,g\),
\[
 \langle f,K(t)g\rangle
 =\lim_n\langle f,K_n(t)g\rangle
 =\lim_n\langle K_n(t)f,g\rangle
 =\langle K(t)f,g\rangle,
\]
so \(K(t)\) is self-adjoint; the same limit of quadratic forms
preserves positivity.  Pass to the limit in (52.3) on \(D\), using
uniform boundedness to handle the product.  Strong continuity
extends the semigroup identity from \(D\) to all \(s,t\geq0\).
Theorem \ref{thm:v15v51-semigroup} gives (52.6).
\end{proof}

\section{Identified changing Hilbert spaces}

If the cutoff OS spaces are \({\cal H}_n\), the hypotheses above
apply after isometries
\[
 J_n:{\cal H}_n\longrightarrow{\cal H}_{\rm OS},\qquad
 P_n=J_nJ_n^*\longrightarrow I\quad\text{strongly},     \tag{52.7}
\]
by replacing \(K_n(t)\) with
\[
 \widehat K_n(t)=J_nK_n(t)J_n^*+(I-P_n).                \tag{52.8}
\]
The identity on the unused complement is only an embedding device;
physical observables must lie asymptotically in \(P_n{\cal H}_{\rm
OS}\).  A choice of \(J_n\) is part of the continuum construction,
not a consequence of equal finite dimensions.

\section{Transporting the vacuum gap}

Let \(P_n^{\rm vac}\) be the projection onto the regulator vacuum
sector, transported to the common carrier, and suppose
\[
 P_n^{\rm vac}\longrightarrow P^{\rm vac}
 \quad\text{strongly}.                                  \tag{52.9}
\]

\begin{theorem}[Uniform semigroup decay transports a gap]
\label{thm:v15v52-gap}
Assume Theorem \ref{thm:v15v52-semigroup}, the vacuum projections
commute with \(K_n(t)\), and for some \(m>0\),
\[
 \boxed{
 \|K_n(t)-P_n^{\rm vac}\|_{\rm op}\leq e^{-mt}
 \quad\text{for every \(n,t\geq0\)}.}                   \tag{52.10}
\]
Then
\[
 K(t)P^{\rm vac}=P^{\rm vac},\qquad
 \|K(t)-P^{\rm vac}\|_{\rm op}\leq e^{-mt},             \tag{52.11}
\]
and the generator obeys
\[
 \sigma(H|_{(I-P^{\rm vac}){\cal H}_{\rm OS}})
 \subset[m,\infty).                                     \tag{52.12}
\]
\end{theorem}

\begin{proof}
Strong convergence of \(K_n(t)\) and \(P_n^{\rm vac}\) passes the
vacuum identity to the limit.  For every unit vector \(f\),
\[
 \|(K(t)-P^{\rm vac})f\|
 =\lim_n\|(K_n(t)-P_n^{\rm vac})f\|
 \leq e^{-mt}.
\]
Taking the supremum gives (52.11).  On the vacuum orthogonal
subspace, spectral calculus yields
\[
 \sup_{\lambda\in\sigma(H)\setminus\{0\}}e^{-t\lambda}
 \leq e^{-mt}
\]
for every \(t>0\), which is equivalent to (52.12).
\end{proof}

\begin{corollary}[Finite-volume version]
\label{cor:v15v52-finite-gap}
If every \(H_n\) has vacuum projection \(P_n^{\rm vac}\) and a
uniform spectral gap at least \(m\), then (52.10) holds.  Hence the
gap passes to the limit provided (52.9) and the common-semigroup
hypotheses also hold.
\end{corollary}

\begin{proof}
On the vacuum complement, the spectral theorem gives
\(\|e^{-tH_n}\|\leq e^{-mt}\).
\end{proof}

\section{Why weaker spectral evidence is insufficient}

\begin{countertheorem}[Soft modes can escape every fixed vector]
\label{cth:v15v52-escape}
On \(\ell^2(\mathbb N)\), define
\[
 H_ne_n=n^{-1}e_n,\qquad H_ne_k=e_k\quad(k\neq n).       \tag{52.13}
\]
Then
\[
 e^{-tH_n}\longrightarrow e^{-t}I
 \quad\text{strongly for every \(t>0\)},                 \tag{52.14}
\]
although
\[
 \inf\sigma(H_n)=n^{-1}\longrightarrow0.                \tag{52.15}
\]
Thus regulator soft eigenvalues can escape through moving
eigenvectors and are invisible to fixed-vector strong convergence.
\end{countertheorem}

\begin{proof}
For \(f=(f_k)\),
\[
 \|(e^{-tH_n}-e^{-t}I)f\|
 =|e^{-t/n}-e^{-t}|\,|f_n|\longrightarrow0,
\]
because \(f_n\to0\).  Equation (52.15) is immediate.
\end{proof}

This example also shows why a list of finite-cutoff eigenvalues is
not a continuum spectral theorem.  Uniform operator decay such as
(52.10), or an equivalent resolvent estimate with compactness of the
relevant spectral screens, is needed.

\section{Resolvent formulation}

For \(\lambda>0\),
\[
 (H+\lambda)^{-1}
 =\int_0^\infty e^{-\lambda t}K(t)\,\dd t.              \tag{52.16}
\]
Local uniform strong convergence of the contractions and dominated
convergence give
\[
 (H_n+\lambda)^{-1}\longrightarrow(H+\lambda)^{-1}
 \quad\text{strongly}.                                  \tag{52.17}
\]
Norm-resolvent convergence requires operator-norm control and does
not follow from Theorem \ref{thm:v15v52-semigroup}.  Compact
resolvent likewise requires collective compactness of the images of
the unit ball.  These are independent of the OS positivity proved
by the boundary factorization.

\section{Physical mass versus regulator parameters}

The number \(m\) in (52.10) is a Hamiltonian energy gap after OS
reconstruction.  A Wilson mass parameter, a resolvent shift
\(\lambda\), an inverse correlation length in lattice units, and an
eigenvalue of a finite boundary pencil are different quantities.
Identifying a physical particle mass additionally requires:
\begin{enumerate}
\item spatial translations and a joint energy--momentum
      representation;
\item the forward spectrum condition after Lorentzian
      reconstruction;
\item a nonzero local-field matrix element in the relevant spectral
      sector;
\item scaling of lattice units to the common physical metric.
\end{enumerate}
Theorem \ref{thm:v15v52-gap} proves only the energy-gap handoff under
(52.10).

\begin{assessment}[Physical status]
\label{ass:v15v52-status}
V52 proves that approximate finite-cutoff transfers yield a positive
continuum semigroup under dense-time convergence, equicontinuity,
and vanishing semigroup defect.  It also gives a sufficient uniform
operator estimate for transporting a vacuum gap and exhibits
spectral escape under weaker convergence.

The source does not construct the common OS Hilbert identifications,
the full reflected transfers, the strong equicontinuity (52.2), or
the uniform vacuum decay (52.10).  No Standard Model mass spectrum
or Einstein-coupled Hamiltonian has therefore been obtained.
\end{assessment}

\begin{decision}[V53 common-core generator convergence]
\label{dec:v15v52-next}
Derive practical sufficient conditions for (52.2)--(52.3) from
finite-cutoff quadratic forms.  Use Mosco convergence on the OS
carrier to obtain strong resolvent and semigroup convergence, then
add a collectively compact low-energy screen to prevent spectral
escape.  Keep the interacting locality and field-visibility rows
separate from spectral convergence.
\end{decision}

\begin{firewall}[Claim boundary after V52]
\label{fire:v15v52}
V52 proves the cofinal semigroup and uniform-gap handoff theorems.
Their reflected-transfer, common-carrier, equicontinuity, vacuum,
and uniform-decay hypotheses remain unproved for the physical
SM--Einstein regulator.
\end{firewall}

\section*{Version V52 append-only record}
\addcontentsline{toc}{section}{Version V52 append-only record}
The complete V51 source was retained before this continuation.  It
had \(513390\) bytes, \(14709\) lines, and SHA--256
\[
 \texttt{89786BBBC8C1233EB297FB1BB1930F39090C603C90A92C09715706E3BFD3AEE3}.
\]
V52 proved cofinal convergence of positive approximate transfer
semigroups, transported a uniform vacuum gap by operator semigroup
decay, and exhibited escaping soft modes that defeat weaker
fixed-vector spectral evidence.

\chapter{Fifteenth-cycle continuation V53: Mosco forms, collective compactness, and stable low-energy screens}
\label{sec:v15v53}

\section{Closed forms on the common OS carrier}

Let \({\cal H}\) be the common OS Hilbert space of V52.  For each
cutoff let \(q_n\) be a densely defined, closed, nonnegative
quadratic form with associated self-adjoint operator \(H_n\geq0\).
Let \(q\) be another such form with operator \(H\).

\begin{definition}[Mosco convergence]
\label{def:v15v53-mosco}
The forms \(q_n\) converge to \(q\) in the Mosco sense if:
\begin{enumerate}
\item whenever \(u_n\rightharpoonup u\) weakly in \({\cal H}\),
      \[
       q(u)\leq\liminf_{n\to\infty}q_n(u_n);             \tag{53.1}
      \]
\item for every \(u\in{\cal H}\), there are \(u_n\to u\) strongly
      such that
      \[
       q(u)\geq\limsup_{n\to\infty}q_n(u_n).             \tag{53.2}
      \]
\end{enumerate}
Forms take the value \(+\infty\) outside their domains.
\end{definition}

\section{Resolvent and semigroup convergence}

\begin{theorem}[Mosco resolvent theorem]
\label{thm:v15v53-mosco-resolvent}
If \(q_n\to q\) in the Mosco sense, then for every \(\lambda>0\),
\[
 \boxed{
 (H_n+\lambda)^{-1}f\longrightarrow(H+\lambda)^{-1}f
 \quad\text{for every }f\in{\cal H}.}                   \tag{53.3}
\]
Consequently, for every \(t\geq0\),
\[
 e^{-tH_n}f\longrightarrow e^{-tH}f                    \tag{53.4}
\]
strongly, locally uniformly for \(t\) in compact subsets of
\([0,\infty)\).
\end{theorem}

\begin{proof}
For fixed \(f\) and \(\lambda>0\), the vector
\[
 u_n=(H_n+\lambda)^{-1}f
\]
is the unique minimizer of
\[
 {\cal J}_n(v)=q_n(v)+\lambda\|v\|^2
 -2\operatorname{Re}\langle f,v\rangle.                 \tag{53.5}
\]
Comparison with \(v=0\) gives
\(\lambda\|u_n\|\leq2\|f\|\), so the minimizers are weakly
precompact.  If \(u_{n_k}\rightharpoonup u_*\), (53.1) gives the
liminf inequality for (53.5).  For any \(v\), a recovery sequence
from (53.2) gives the matching limsup inequality.  Hence \(u_*\)
minimizes
\[
 {\cal J}(v)=q(v)+\lambda\|v\|^2
-2\operatorname{Re}\langle f,v\rangle,
\]
so \(u_*=(H+\lambda)^{-1}f\).  Uniqueness removes the need to pass
to a subsequence.  Applying the liminf and recovery inequalities to
the minimum values shows
\(\|u_n\|\to\|u_*\|\); weak convergence plus convergence of norms is
strong convergence, proving (53.3).

Strong resolvent convergence of nonnegative self-adjoint operators
implies strong convergence of every bounded continuous function of
the operators that has a limit at infinity.  Apply this to
\(x\mapsto e^{-tx}\).  Uniformity away from \(t=0\) follows from
uniform approximation by resolvent polynomials; uniformity at zero
follows first on a dense set of form recovery vectors by (53.7)
below, and then on all of \({\cal H}\) by contraction.
\end{proof}

\section{A practical equicontinuity estimate}

\begin{lemma}[Form-energy control at time zero]
\label{lem:v15v53-equicontinuity}
For \(u\in D(q_n)\),
\[
 \boxed{
 \|(e^{-tH_n}-I)u\|^2\leq t\,q_n(u).}                   \tag{53.6}
\]
Thus if a dense core has recovery vectors \(u_n\to u\) with
\(\sup_nq_n(u_n)<\infty\), then the semigroups are uniformly
strongly continuous at zero on that core.
\end{lemma}

\begin{proof}
Let \(\mu_u\) be the spectral measure of \(H_n\) at \(u\).  Since
\((1-e^{-x})^2\leq x\) for \(x\geq0\),
\[
 \|(e^{-tH_n}-I)u\|^2
 =\int(1-e^{-t\xi})^2\,\dd\mu_u(\xi)
 \leq t\int\xi\,\dd\mu_u(\xi)
 =tq_n(u).
                                                               \tag{53.7}
\]
\end{proof}

This supplies the equicontinuity input of V52 once a common dense
form core and bounded-energy recovery sequences have been built.

\section{Collective compactness}

Strong resolvent convergence alone permits the moving-mode behavior
of Countertheorem \ref{cth:v15v52-escape}.  The missing condition is
compactness uniform in \(n\).

\begin{definition}[Collectively compact resolvents]
\label{def:v15v53-collective}
For one \(\lambda>0\), the resolvents
\[
 R_n=(H_n+\lambda)^{-1}                                 \tag{53.8}
\]
are collectively compact if
\[
 \bigcup_nR_n\{f:\|f\|\leq1\}                           \tag{53.9}
\]
has compact closure in \({\cal H}\).
\end{definition}

\begin{theorem}[Strong plus collective compactness gives norm
resolvent convergence]
\label{thm:v15v53-norm-resolvent}
If \(R_n\to R=(H+\lambda)^{-1}\) strongly, the \(R_n\) are
self-adjoint, and (53.9) is relatively compact, then \(R\) is
compact and
\[
 \boxed{\|R_n-R\|_{\rm op}\longrightarrow0.}            \tag{53.10}
\]
\end{theorem}

\begin{proof}
The strong image under \(R\) of each unit vector is a limit of
points in the compact closure of (53.9), so \(R\) maps the unit ball
into that closure and is compact.

If (53.10) failed, choose unit vectors \(f_n\) and \(\epsilon>0\)
such that
\(\|(R_n-R)f_n\|\geq\epsilon\).  Pass to a subsequence with
\(f_n\rightharpoonup f\) and \(R_nf_n\to y\) strongly, using
collective compactness.  For every \(g\),
\[
 \langle g,y\rangle
 =\lim_n\langle R_ng,f_n\rangle
 =\langle Rg,f\rangle
 =\langle g,Rf\rangle,                                  \tag{53.11}
\]
where self-adjointness and strong convergence on the fixed vector
\(g\) were used.  Thus \(y=Rf\).  Compactness of \(R\) gives
\(Rf_n\to Rf\) strongly, contradicting the lower bound.
\end{proof}

\begin{proposition}[Uniform compact form embedding]
\label{prop:v15v53-compact-embedding}
Let \({\cal V}\) be a Hilbert space compactly embedded in
\({\cal H}\).  If
\[
 q_n(u)+\|u\|^2\geq c_{\cal V}\|u\|_{\cal V}^2          \tag{53.12}
\]
for every \(n\) and \(u\in D(q_n)\), with one
\(c_{\cal V}>0\), then the resolvents (53.8) are collectively
compact.
\end{proposition}

\begin{proof}
For \(\|f\|\leq1\), let \(u=R_nf\).  The resolvent equation tested
against \(u\) gives
\[
 q_n(u)+\lambda\|u\|^2
 =\operatorname{Re}\langle f,u\rangle
 \leq\|u\|.                                             \tag{53.13}
\]
Also \(\|u\|\leq\lambda^{-1}\).  Combining (53.12)--(53.13)
bounds \(\|u\|_{\cal V}\) uniformly in \(n\) and \(f\).
The compact embedding makes the union of these images relatively
compact in \({\cal H}\).
\end{proof}

\section{Stable low-energy spectral projections}

\begin{theorem}[Isolated screen stability]
\label{thm:v15v53-screen}
Assume norm-resolvent convergence (53.10).  Let
\(I=[a,b]\) be a bounded interval whose endpoints lie in the
resolvent set of \(H\), and suppose a contour \(\Gamma\) surrounds
\(\sigma(H)\cap I\) while avoiding \(\sigma(H)\).  Then for all
large \(n\), \(\Gamma\) avoids \(\sigma(H_n)\), and
\[
 P_n(I)=\frac1{2\pi i}\int_\Gamma(z-H_n)^{-1}\,\dd z
 \longrightarrow
 P(I)=\frac1{2\pi i}\int_\Gamma(z-H)^{-1}\,\dd z         \tag{53.14}
\]
in operator norm.  If \(P(I)\) has finite rank, then
\[
 \operatorname{rank}P_n(I)=\operatorname{rank}P(I)      \tag{53.15}
\]
for all large \(n\).
\end{theorem}

\begin{proof}
Norm-resolvent convergence at one point and the resolvent identity
give uniform norm convergence on the compact contour \(\Gamma\).
Integrating proves (53.14).  Projections at operator-norm distance
less than one have isomorphic ranges, giving (53.15).
\end{proof}

\begin{corollary}[No low-energy escape under collective compactness]
\label{cor:v15v53-no-escape}
Under Mosco convergence and Proposition
\ref{prop:v15v53-compact-embedding}, every isolated finite-rank
low-energy screen of \(H\) is approximated with the correct
multiplicity by the corresponding screens of \(H_n\).  Moving
eigenvectors cannot disappear as in (52.13).
\end{corollary}

\begin{proof}
Theorem \ref{thm:v15v53-mosco-resolvent},
Theorem \ref{thm:v15v53-norm-resolvent}, and Theorem
\ref{thm:v15v53-screen} apply successively.
\end{proof}

\section{Field visibility remains separate}

Let \(\Omega_n\) be vacuum vectors and \(F_n\) transported local
fields.  Spectral existence in a screen does not show that a local
field creates that state.

\begin{proposition}[Visibility under norm projection convergence]
\label{prop:v15v53-visibility}
If
\[
 P_n(I)\to P(I)\ \text{in norm},\qquad
 F_n\Omega_n\to\psi\ \text{strongly},                  \tag{53.16}
\]
then
\[
 P_n(I)F_n\Omega_n\longrightarrow P(I)\psi.             \tag{53.17}
\]
In particular, if
\(\|P_n(I)F_n\Omega_n\|\geq c>0\), then
\(\|P(I)\psi\|\geq c\).
\end{proposition}

\begin{proof}
Add and subtract \(P_n(I)\psi\) and use the projection norm bound:
\[
 \|P_nF_n\Omega_n-P\psi\|
 \leq\|F_n\Omega_n-\psi\|+\|P_n-P\|\,\|\psi\|.
\]
Take the limit.
\end{proof}

The local-field convergence in (53.16) belongs to the observable
consistency programme of V44--V49, not to Mosco convergence alone.

\section{Vacuum and gap screens}

If zero is an isolated eigenvalue of \(H\), Theorem
\ref{thm:v15v53-screen} stabilizes its finite-dimensional vacuum
projection.  To prove a positive gap \(m\), one still needs a
uniform spectral exclusion or semigroup decay such as (52.10).
Collective compactness prevents eigenvectors from escaping but does
not forbid a genuine sequence of eigenvalues converging to zero
inside a compact region.

\begin{assessment}[Physical form packet]
\label{ass:v15v53-packet}
The practical packet needed from the regulator is now:
\begin{enumerate}
\item one common OS Hilbert carrier and closed forms \(q_n\);
\item the Mosco liminf and recovery sequence;
\item a compact form-domain estimate (53.12);
\item a vacuum projection and uniform gap or decay estimate;
\item convergence of local-field vectors for visibility.
\end{enumerate}
Together these yield a positive continuum Hamiltonian, stable
low-energy screens, and a visible uniform energy gap.
\end{assessment}

\begin{assessment}[Literal source applicability]
\label{ass:v15v53-source}
The source discusses form/resolvent transport abstractly but does
not construct the complete interacting OS forms after quotient,
phase, and zero-mode reduction.  It supplies neither (53.1)--(53.2)
nor the uniform compact embedding (53.12) for that physical
carrier.  V53 therefore closes the operator theorem but not its
physical hypotheses.
\end{assessment}

\begin{decision}[V54 form convergence from local cylinders]
\label{dec:v15v53-next}
Connect the measure construction to the form packet.  Define
finite-cutoff transfer forms on the countable local OS core and
prove a core criterion for Mosco convergence using convergence of
Gram matrices, a uniform graph-norm coercivity, and core density.
Identify why pointwise convergence of form values without a
liminf estimate is insufficient.
\end{decision}

\begin{firewall}[Claim boundary after V53]
\label{fire:v15v53}
V53 proves Mosco-to-semigroup convergence, the collective
compactness upgrade to norm resolvents, stable isolated spectral
screens, and a field-visibility handoff.  The physical OS forms,
compact embedding, gap, and local-field convergence remain
unproved.
\end{firewall}

\section*{Version V53 append-only record}
\addcontentsline{toc}{section}{Version V53 append-only record}
The complete V52 source was retained before this continuation.  It
had \(522860\) bytes, \(14979\) lines, and SHA--256
\[
 \texttt{8E4DAEB5F6CF90CD44BF4078EE181F53EEF6BE36AA49481AF4818613B3AF467B}.
\]
V53 proved Mosco convergence of OS forms implies strong resolvent
and semigroup convergence, derived a uniform compact-embedding
criterion for norm resolvents, and stabilized isolated low-energy
screens and their local-field visibility.

\chapter{Fifteenth-cycle continuation V54: local-core form convergence requires a uniform tail}
\label{sec:v15v54}

\section{From transfer Grams to generator forms}

Let \(K_n(t)=e^{-tH_n}\) be positive contraction transfers on the
common OS Hilbert space \({\cal H}\).  Their bounded time-\(t\)
forms are
\[
 q_{n,t}(u,v)=
 \frac1t\langle u,(I-K_n(t))v\rangle.                   \tag{54.1}
\]
For \(u\in D(H_n^{1/2})\), spectral calculus gives
\[
 q_n(u)=\|H_n^{1/2}u\|^2
 =\sup_{t>0}q_{n,t}(u,u).                               \tag{54.2}
\]
Thus convergence of reflected Gram matrices at fixed slab times
controls \(q_{n,t}\), but passage to the generator form additionally
requires uniform control as \(t\downarrow0\).

\begin{proof}
For a spectral value \(\lambda\geq0\),
\[
 \frac{1-e^{-t\lambda}}t\uparrow\lambda
 \quad\text{as }t\downarrow0.
\]
Integrate this monotone scalar convergence against the spectral
measure of \(H_n\).
\end{proof}

\section{A common graph space}

Let \({\cal V}\) be a Hilbert space densely and continuously
embedded in \({\cal H}\).  Suppose the closed nonnegative forms
\(q_n\) and \(q\) have common domain \({\cal V}\).  Write
\[
 a_n(u,v)=q_n(u,v)+\langle u,v\rangle_{\cal H},\qquad
 a(u,v)=q(u,v)+\langle u,v\rangle_{\cal H}.             \tag{54.3}
\]
Assume the uniform graph-norm bounds
\[
 c\|u\|_{\cal V}^2
 \leq a_n(u,u)
 \leq C\|u\|_{\cal V}^2                                \tag{54.4}
\]
and the same inequalities for \(a\), with \(0<c\leq C<\infty\).

\begin{theorem}[Uniform form convergence implies Mosco convergence]
\label{thm:v15v54-uniform-form}
If
\[
 |a_n(u,v)-a(u,v)|
 \leq\epsilon_n\|u\|_{\cal V}\|v\|_{\cal V},
 \qquad\epsilon_n\to0,                                  \tag{54.5}
\]
then \(q_n\to q\) in the Mosco sense.  Moreover, for every
\(\lambda>0\),
\[
 \|(H_n+\lambda)^{-1}-(H+\lambda)^{-1}\|_{
 {\cal H}\to{\cal V}}\longrightarrow0.                  \tag{54.6}
\]
\end{theorem}

\begin{proof}
For the Mosco liminf, let \(u_n\rightharpoonup u\) in
\({\cal H}\) and suppose \(\liminf q_n(u_n)<\infty\).
After taking a subsequence realizing the liminf, (54.4) makes
\(u_n\) bounded in \({\cal V}\).  Pass to a weak
\({\cal V}\) limit.  Continuity of the embedding identifies that
limit with \(u\).  Weak lower semicontinuity of \(a\) and (54.5)
give
\[
 a(u,u)\leq\liminf a(u_n,u_n)
 \leq\liminf[a_n(u_n,u_n)+
 \epsilon_n\|u_n\|_{\cal V}^2].
\]
The last error tends to zero, proving (53.1).  For recovery take
\(u_n=u\); (54.5) gives \(q_n(u)\to q(u)\).

For (54.6), let \(u_n=(H_n+\lambda)^{-1}f\) and
\(u=(H+\lambda)^{-1}f\).  The variational equations differ by the
form perturbation in (54.5).  Test their difference with
\(u_n-u\), use coercivity of \(q_n+\lambda\|\cdot\|^2\), and the
uniform bound \(\|u\|_{\cal V}\leq C_\lambda\|f\|\).  This yields
\[
 \|u_n-u\|_{\cal V}
 \leq C_\lambda\epsilon_n\|f\|,
\]
which proves (54.6).
\end{proof}

\section{Finite local cores plus a uniform tail}

Let \({\cal D}\subset{\cal V}\) be the countable local OS core of
V49.  Choose finite-rank \({\cal V}\)-orthogonal projections
\[
 P_m{\cal V}\subset{\cal D},\qquad P_m\to I
 \quad\text{strongly on }{\cal V}.                      \tag{54.7}
\]
Let \(A_n,A:{\cal V}\to{\cal V}^*\) represent \(a_n,a\).

\begin{theorem}[Core Gram and tail criterion]
\label{thm:v15v54-core-tail}
Assume (54.4) and:
\begin{align}
 \|P_m^*(A_n-A)P_m\|_{{\cal V}\to{\cal V}^*}
 &\longrightarrow0
 &&\text{for each fixed \(m\)},                         \tag{54.8}\\
 \sup_n\|A_n-P_m^*A_nP_m\|_{{\cal V}\to{\cal V}^*}
 +\|A-P_m^*AP_m\|_{{\cal V}\to{\cal V}^*}
 &\leq\eta_m,\qquad \eta_m\to0.                         \tag{54.9}
\end{align}
Then (54.5) holds and hence \(q_n\to q\) in the Mosco sense.
\end{theorem}

\begin{proof}
Insert the two compressed forms:
\[
 \begin{split}
 \|A_n-A\|
 \leq{}&\|A_n-P_m^*A_nP_m\|\\
 &+\|P_m^*(A_n-A)P_m\|
 +\|P_m^*AP_m-A\|.                                     \tag{54.10}
 \end{split}
\]
For fixed \(m\), the middle term tends to zero by (54.8); the two
outer terms are at most \(\eta_m\).  First take \(n\to\infty\), then
\(m\to\infty\), obtaining operator-norm convergence
\(A_n\to A\).  This is (54.5), and Theorem
\ref{thm:v15v54-uniform-form} applies.
\end{proof}

The finite matrices in (54.8) are precisely generator-form Grams on
a finite local observable span.  Equation (54.9) is the missing
source-complete tail: no unit graph-norm vector may hide most of its
form interaction outside every finite local core.

\section{Why finite Gram convergence is insufficient}

\begin{countertheorem}[A moving low-cost direction defeats the
Mosco liminf]
\label{cth:v15v54-moving}
Let \({\cal H}={\cal V}=\ell^2(\mathbb N)\), choose
\(0<c<1\), and define
\[
 q_n(u)=c\|u\|^2+|u_1-u_n|^2,\qquad
 q(u)=c\|u\|^2+|u_1|^2.                                \tag{54.11}
\]
On every fixed finite-coordinate core,
\[
 q_n(u,v)\longrightarrow q(u,v),                       \tag{54.12}
\]
and all forms are uniformly coercive and bounded.  Nevertheless,
with \(v_n=e_1+e_n\rightharpoonup e_1\),
\[
 \liminf_nq_n(v_n)=2c<c+1=q(e_1).                       \tag{54.13}
\]
Thus the Mosco liminf fails.
\end{countertheorem}

\begin{proof}
For fixed finite-support \(u\), its \(n\)-th coordinate vanishes for
large \(n\), proving (54.12), and polarization gives the bilinear
version.  Uniform equivalence to the \(\ell^2\) norm is immediate.
For \(v_n\), the difference \(u_1-u_n\) vanishes, so
\(q_n(v_n)=2c\).  The weak limit is \(e_1\), whose limiting form
value is \(c+1\).
\end{proof}

The example is the form analogue of spectral escape in V52.  It
proves that core density and pointwise Gram convergence do not
replace (54.9).

\section{Obtaining the generator tail from finite times}

Suppose finite-time transfer forms satisfy a tail bound
\[
 \sup_n\sup_{0<t\leq t_0}
 \|Q_m^*q_{n,t}Q_m\|_{{\cal V}\to{\cal V}^*}
 \leq\eta_m,\qquad Q_m=I-P_m,                           \tag{54.14}
\]
and mixed head--tail blocks have the same vanishing bound.  If
(54.14) is uniform as \(t\downarrow0\), then (54.2) and monotone
convergence transfer it to the generator form.  A bound only at one
fixed slab time does not: high generator energies are compressed by
\((1-e^{-t\lambda})/t\leq1/t\).

\begin{proposition}[Small-time graph recovery]
\label{prop:v15v54-small-time}
Suppose for \(u\in{\cal D}\)
\[
 \sup_n|q_n(u)-q_{n,t}(u)|\leq\omega_u(t),
 \qquad\omega_u(t)\to0,                                 \tag{54.15}
\]
and \(q_{n,t}(u,v)\to q_t(u,v)\) for each fixed \(t>0\) and
\(u,v\in{\cal D}\).  Then \(q_n(u,v)\) converges on
\({\cal D}\) to
\[
 q(u,v)=\lim_{t\downarrow0}q_t(u,v).                    \tag{54.16}
\]
\end{proposition}

\begin{proof}
For diagonal values, add and subtract \(q_{n,t}(u)\).  At fixed
\(t\), pass \(n\to\infty\); then use (54.15) and let
\(t\downarrow0\).  Polarization gives the bilinear conclusion.
\end{proof}

\section{Compactness and core tails}

If \({\cal V}\Subset{\cal H}\), Theorem
\ref{thm:v15v53-norm-resolvent} follows after Mosco convergence.
Compact embedding does not itself imply (54.9), since (54.9)
controls the forms in graph norm rather than the embedding in
\({\cal H}\).  Both conditions are needed for the strongest route:
\[
 \boxed{
 \text{core Gram convergence}
 +\text{uniform form tail}
 +\text{compact graph embedding}.}                      \tag{54.17}
\]
They yield, respectively, identification, Mosco convergence, and
stable isolated spectral screens.

\begin{assessment}[Physical population]
\label{ass:v15v54-status}
The local observable convergence of V44 and reflected Gram estimates
of V49 can address (54.8) at fixed slab time.  The shell estimates
of V48 suggest a route to (54.9), but only if their constants are
uniform in the graph-norm unit ball and as slab time tends to zero.
No previous version proves that stronger uniformity.

The source does not supply the common generator graph space
\({\cal V}\), the small-time estimate (54.15), or the uniform
form-tail (54.9).  Mosco convergence of the physical OS generator
remains conditional.
\end{assessment}

\begin{decision}[V55 compulsory audit]
\label{dec:v15v54-next}
At V55 inspect the newest YM and NS companion notes for a
source-complete tail, compact graph embedding, or small-time
transfer estimate.  Then continue at V56 with a quantitative
finite-time-to-generator theorem based on a uniform second spectral
moment on the local core.
\end{decision}

\begin{firewall}[Claim boundary after V54]
\label{fire:v15v54}
V54 proves a sufficient core-plus-tail criterion for Mosco and norm
resolvent convergence and exhibits failure of the liminf under
fixed-core Gram convergence alone.  The physical uniform form tail
and small-time generator control remain unproved.
\end{firewall}

\section*{Version V54 append-only record}
\addcontentsline{toc}{section}{Version V54 append-only record}
The complete V53 source was retained before this continuation.  It
had \(534252\) bytes, \(15307\) lines, and SHA--256
\[
 \texttt{069C494CE25127B54AE6C9609A6B6F88DDEF458A1FD8C4DDD8EAC9B708B5943F}.
\]
V54 proved uniform form and finite-core-plus-tail Mosco criteria,
gave a small-time recovery condition for generator Grams, and
constructed a uniformly coercive moving-direction counterexample
showing why fixed-core convergence alone fails.

\chapter{Fifteenth-cycle continuation V55: compulsory small-time form audit}
\label{sec:v15v55}

\section{Companion review}

The inspected active companion checkpoints are
\[
\begin{array}{c|c|c|c}
\text{programme}&\text{file}&\text{lines}&\text{SHA--256}\\ \hline
\mathrm{YM}&
\texttt{Renewal\_Yang\_Mills\_Mass\_Gap\_Research\_Notes\_V39.tex}&
9076&
\texttt{9B44B0FBCA6037F3319E86036753B0F3B76BDD06F6045F00747E05DC01E23F9D}\\
\mathrm{NS}&
\texttt{Renewal\_NS\_Stability\_Research\_Notes\_V26.tex}&
5319&
\texttt{82C887F8C2C69E4F28FAD7C3DC8DE5B9099CB5A4BF431EBA1FFF3E91935978AD}.
\end{array}                                                \tag{55.1}
\]
YM V39 remains an exact finite Wilson-pencil atlas with no
interacting OS transfer or continuum measure.  NS V26 remains a
pointwise Oseen-kernel norm calculation.  Neither provides a common
OS graph space, a uniform small-time transfer expansion, a compact
graph embedding, or the source-complete tail (54.9).

\begin{theorem}[V55 nonimport]
\label{thm:v15v55-audit}
No companion result populates the V54 Mosco packet.  The next
estimate must be proved within the SM--Einstein transfer programme.
\end{theorem}

\begin{proof}
The companion statements have neither the carrier nor any of the
form hypotheses listed in Assessment \ref{ass:v15v53-packet}.
\end{proof}

\begin{decision}[V56 second-moment generator recovery]
\label{dec:v15v55-next}
For \(q_{n,t}=t^{-1}\langle\cdot,(I-e^{-tH_n})\cdot\rangle\),
bound \(q_n-q_{n,t}\) by a uniform second spectral moment on the
local core.  Extend the estimate to mixed Grams by polarization and
derive a joint choice of slab time and cutoff that populates
(54.15)--(54.16).
\end{decision}

\begin{firewall}[Claim boundary after V55]
\label{fire:v15v55}
V55 is a compulsory audit and imports no companion theorem into the
physical OS form.
\end{firewall}

\section*{Version V55 append-only record}
\addcontentsline{toc}{section}{Version V55 append-only record}
The complete V54 source was retained before this continuation.  It
had \(543569\) bytes, \(15575\) lines, and SHA--256
\[
 \texttt{FC4925AFD4CD72C0D5BF013C78DC67958E8F8D4FC28C74C4B6CA181BDE2A9BBC}.
\]
V55 performed the required companion audit and retained the
second-spectral-moment generator estimate as the V56 target.

\chapter{Fifteenth-cycle continuation V56: second spectral moments recover the generator from finite slabs}
\label{sec:v15v56}

\section{The scalar remainder}

For \(\lambda,t\geq0\), define
\[
 r_t(\lambda)=
 \lambda-\frac{1-e^{-t\lambda}}t.                       \tag{56.1}
\]
The elementary inequalities
\[
 1-e^{-x}\leq x,\qquad
 1-e^{-x}\geq x-\frac{x^2}{2}                           \tag{56.2}
\]
give
\[
 \boxed{0\leq r_t(\lambda)\leq\frac t2\lambda^2.}       \tag{56.3}
\]

\section{Diagonal and mixed form estimates}

Let \(H\geq0\) be self-adjoint,
\[
 q(u,v)=\langle H^{1/2}u,H^{1/2}v\rangle,\qquad
 q_t(u,v)=\frac1t\langle u,(I-e^{-tH})v\rangle.          \tag{56.4}
\]

\begin{theorem}[Second-moment finite-slab error]
\label{thm:v15v56-second-moment}
For \(u\in D(H)\),
\[
 \boxed{
 0\leq q(u,u)-q_t(u,u)
 \leq\frac t2\|Hu\|^2.}                                 \tag{56.5}
\]
For \(u,v\in D(H)\),
\[
 \boxed{
 |q(u,v)-q_t(u,v)|
 \leq\frac t2\|Hu\|\,\|Hv\|.}                           \tag{56.6}
\]
\end{theorem}

\begin{proof}
Let \(E_H\) be the spectral resolution of \(H\).  Then
\[
 q(u,u)-q_t(u,u)
 =\int_0^\infty r_t(\lambda)\,
       \dd\langle u,E_H(\lambda)u\rangle.
\]
Equation (56.3) gives (56.5).  The operator
\(r_t(H)\) is positive, so Cauchy--Schwarz for its quadratic form
and (56.5) give
\[
 |\langle u,r_t(H)v\rangle|
 \leq
 \langle u,r_t(H)u\rangle^{1/2}
 \langle v,r_t(H)v\rangle^{1/2}
 \leq\frac t2\|Hu\|\,\|Hv\|.
\]
\end{proof}

\begin{corollary}[Uniform recovery on a local core]
\label{cor:v15v56-uniform-core}
Let \(H_n\geq0\), and let \(u_n,v_n\) be transported core vectors
such that
\[
 \sup_n\|H_nu_n\|\leq M_u,\qquad
 \sup_n\|H_nv_n\|\leq M_v.                              \tag{56.7}
\]
Then
\[
 \sup_n|q_n(u_n,v_n)-q_{n,t}(u_n,v_n)|
 \leq\frac t2M_uM_v.                                    \tag{56.8}
\]
Thus (54.15) holds uniformly on every finite core whose second
spectral moments are uniformly bounded.
\end{corollary}

\section{Balancing transfer and time errors}

Suppose limit vectors \(u,v\in D(H)\) satisfy the analogous bounds
and the finite-slab comparison obeys
\[
 |q_{n,t}(u_n,v_n)-q_t(u,v)|
 \leq\frac{\epsilon_n}{t}
 \quad(0<t\leq t_0),                                    \tag{56.9}
\]
where \(\epsilon_n\to0\).  This occurs, for example, when the
underlying transfer matrix element has error at most
\(\epsilon_n\), since (56.4) divides that error by \(t\).

\begin{theorem}[Optimized finite-time-to-generator rate]
\label{thm:v15v56-balance}
Put
\[
 M_{uv}=M_uM_v+\|Hu\|\,\|Hv\|.                          \tag{56.10}
\]
Then
\[
 |q_n(u_n,v_n)-q(u,v)|
 \leq\frac t2M_{uv}+\frac{\epsilon_n}{t}.               \tag{56.11}
\]
If \(M_{uv}>0\) and
\[
 t_n=\sqrt{\frac{2\epsilon_n}{M_{uv}}}\leq t_0,         \tag{56.12}
\]
then
\[
 \boxed{
 |q_n(u_n,v_n)-q(u,v)|
 \leq\sqrt{2M_{uv}\epsilon_n}.}                         \tag{56.13}
\]
\end{theorem}

\begin{proof}
Insert \(q_{n,t}\) and \(q_t\) between the two generator forms.
The two outer differences are bounded by (56.6) and (56.8), and the
middle one by (56.9), proving (56.11).  The two terms on its right
are minimized at (56.12), where their sum is (56.13).
\end{proof}

\begin{corollary}[Diagonal choice without a rate]
\label{cor:v15v56-diagonal}
If for every fixed \(t>0\) the finite-slab forms converge and the
uniform second-moment bounds (56.7) hold, then there is a sequence
\(t_n\downarrow0\) for which the generator Grams converge on every
member of a countable local core.
\end{corollary}

\begin{proof}
Enumerate the core pairs.  For the first \(k\) pairs, choose
\(t^{(k)}>0\) so that their uniform errors in (56.8) are below
\(1/k\).  At that fixed time choose \(N_k\) so the finite-slab
errors for those pairs are below \(1/k\) for \(n\geq N_k\).  Take
\(t_n=t^{(k)}\) for \(N_k\leq n<N_{k+1}\), after replacing by a
decreasing subsequence.  Every fixed pair then converges.
\end{proof}

\section{First moments do not give uniform small-time recovery}

\begin{countertheorem}[Uniform energy without a second moment]
\label{cth:v15v56-first-not-second}
On the one-dimensional Hilbert space for each \(n\), let
\[
 H_n=n,\qquad u_n=n^{-1/2},\qquad t_n=n^{-1}.            \tag{56.14}
\]
Then
\[
 q_n(u_n,u_n)=1                                         \tag{56.15}
\]
is uniformly bounded, but
\[
 q_{n,t_n}(u_n,u_n)=1-e^{-1},\qquad
 q_n(u_n,u_n)-q_{n,t_n}(u_n,u_n)=e^{-1}.                \tag{56.16}
\]
Thus \(t_n\to0\) and a uniform first spectral moment do not imply a
uniform version of (54.15).
\end{countertheorem}

\begin{proof}
Substitute (56.14) into (56.4).
\end{proof}

The missing quantity is
\[
 \|H_nu_n\|^2
 =\langle u_n,H_n^2u_n\rangle=n,                        \tag{56.17}
\]
which diverges.

\section{A low-energy-screen alternative}

Let \(P_n(E)={\bf1}_{[0,E]}(H_n)\).

\begin{proposition}[Screened second moment]
\label{prop:v15v56-screen}
For \(u\in D(q_n)\),
\[
 \|H_nP_n(E)u\|^2
 \leq E\,q_n(u,u).                                      \tag{56.18}
\]
Consequently
\[
 0\leq q_n(P_n(E)u)-q_{n,t}(P_n(E)u)
 \leq\frac{tE}{2}q_n(u,u).                              \tag{56.19}
\]
\end{proposition}

\begin{proof}
On \([0,E]\), \(\lambda^2\leq E\lambda\).  Integrate against the
spectral measure and apply (56.5).
\end{proof}

This alternative needs a uniform high-energy form tail
\[
 q_n((I-P_n(E))u)\longrightarrow0
 \quad(E\to\infty)                                      \tag{56.20}
\]
on the local core.  It is the generator analogue of the smooth
determinant-window tail in V34--V36.

\section{Reading the second moment from time correlations}

For \(u\in D(H)\), define
\[
 C_u(t)=\langle u,e^{-tH}u\rangle.                      \tag{56.21}
\]
Then
\[
 C_u'(0+)=-q(u,u),\qquad
 C_u''(0+)=\|Hu\|^2                                    \tag{56.22}
\]
whenever the corresponding moments are finite.  Therefore the
uniform input (56.7) is a uniform curvature bound at the origin of
the reflected two-point function.  Values at finitely many positive
times do not control that curvature without an additional convexity
or tail estimate.

\begin{assessment}[Physical acquisition]
\label{ass:v15v56-status}
V56 turns the small-time issue into a precise local observable
condition:
\[
 \sup_n\|H_nF_n\Omega_n\|<\infty.                       \tag{56.23}
\]
Equivalently, one needs a uniform second spectral moment of each
core field vector, or a screened estimate plus (56.20).  The source
does not provide these moments for the complete interacting
reflected transfer.  Its finite-lag and finite-matrix data cannot
be differentiated twice at zero without this uniform row.
\end{assessment}

\begin{decision}[V57 transfer curvature from local commutators]
\label{dec:v15v56-next}
Relate (56.23) to the generator commutator of a local field with the
vacuum:
\[
 H_nF_n\Omega_n=[H_n,F_n]\Omega_n
\]
when \(H_n\Omega_n=0\).  Prove a local interaction-norm criterion
for a uniform commutator bound and expose the boundary/counterterm
terms that prevent a formal locality claim from populating it.
\end{decision}

\begin{firewall}[Claim boundary after V56]
\label{fire:v15v56}
V56 proves a sharp second-moment finite-slab error and an optimized
joint cutoff/time rate.  The physical second spectral moments,
screen tails, and twice-differentiable reflected correlations remain
unproved.
\end{firewall}

\section*{Version V56 append-only record}
\addcontentsline{toc}{section}{Version V56 append-only record}
The complete V55 source was retained before this continuation.  It
had \(545852\) bytes, \(15636\) lines, and SHA--256
\[
 \texttt{E8B3B52013326864D3DA07EEA10BA5F954B7085D76451E5EC116CC74F1DA7E09}.
\]
V56 proved the \(t\|Hu\|^2/2\) finite-slab error, optimized it
against transfer convergence, gave a diagonal countable-core
construction, and proved that a uniform first energy moment cannot
replace the second spectral moment.

\chapter{Fifteenth-cycle continuation V57: local commutators as second spectral moments}
\label{sec:v15v57}

\section{Vacuum commutator identity}

Let \(H_n\geq0\) be a finite-cutoff OS Hamiltonian with normalized
vacuum \(\Omega_n\) satisfying
\[
 H_n\Omega_n=0.                                         \tag{57.1}
\]
For every bounded observable \(F_n\) preserving the operator domain,
\[
 \boxed{
 H_nF_n\Omega_n=[H_n,F_n]\Omega_n.}                     \tag{57.2}
\]
Consequently
\[
 \|H_nF_n\Omega_n\|
 \leq\|[H_n,F_n]\|.                                     \tag{57.3}
\]
This identifies the second spectral moment (56.23) with a local
commutator bound.

\begin{proof}
Expand the commutator and use (57.1).
\end{proof}

If the transported vacuum is only approximate,
\(\|H_n\Omega_n\|\leq\epsilon_n\), then
\[
 \|H_nF_n\Omega_n\|
 \leq\|[H_n,F_n]\Omega_n\|+\|F_n\|\epsilon_n.            \tag{57.4}
\]

\section{Uniform interaction norm}

Let \(\Lambda_n\) be a metric lattice and write a bounded
finite-volume Hamiltonian as
\[
 H_n=\sum_{X\Subset\Lambda_n}\Phi_n(X),\qquad
 \Phi_n(X)=\Phi_n(X)^*.                                 \tag{57.5}
\]
For \(\mu\geq0\), define
\[
 \|\Phi_n\|_\mu=
 \sup_{x\in\Lambda_n}
 \sum_{X\ni x}|X|e^{\mu\operatorname{diam}X}
 \|\Phi_n(X)\|.                                         \tag{57.6}
\]

\begin{theorem}[Local-field second-moment bound]
\label{thm:v15v57-local}
Let \(F_n\) be supported in \(O_n\subset\Lambda_n\).  If
\[
 \sup_n\|\Phi_n\|_0\leq J,\qquad
 \sup_n|O_n|\leq N_O,\qquad
 \sup_n\|F_n\|\leq M_F,                                 \tag{57.7}
\]
then
\[
 \boxed{
 \|H_nF_n\Omega_n\|
 \leq2M_FN_OJ.}                                         \tag{57.8}
\]
Hence V56's uniform second-moment hypothesis holds for the
transported local-field vectors.
\end{theorem}

\begin{proof}
Terms with \(X\cap O_n=\varnothing\) commute with \(F_n\).
Therefore
\[
 \|[H_n,F_n]\|
 \leq2\|F_n\|
 \sum_{X:X\cap O_n\neq\varnothing}\|\Phi_n(X)\|.         \tag{57.9}
\]
Sum first over a point \(x\in O_n\).  Since \(|X|\geq1\),
\[
 \sum_{X:X\cap O_n\neq\varnothing}\|\Phi_n(X)\|
 \leq\sum_{x\in O_n}\sum_{X\ni x}|X|\|\Phi_n(X)\|
 \leq |O_n|\|\Phi_n\|_0.                                \tag{57.10}
\]
Equations (57.2), (57.7), and (57.9)--(57.10) prove (57.8).
\end{proof}

The support count in (57.7) is in physical local degrees of
freedom.  If one fixed physical region contains order \(a_n^{-d}\)
lattice sites, the raw cardinality bound fails and must be replaced
by a scaled interaction norm and a smeared observable estimate.

\section{Smeared lattice fields}

Let
\[
 F_n(f)=\sum_{x\in\Lambda_n}w_{x,n}f(x)F_{x,n}          \tag{57.11}
\]
be a bounded smeared observable, with
\(\|F_{x,n}\|\leq1\).  Directly applying Theorem
\ref{thm:v15v57-local} to the union of supports loses the cell
weights.

\begin{proposition}[Weighted commutator bound]
\label{prop:v15v57-smeared}
Assume each \(F_{x,n}\) is supported in a block of uniformly bounded
cardinality and
\[
 \sup_{x,n}\sum_{X:X\cap\operatorname{supp}F_{x,n}\neq\varnothing}
 \|[\Phi_n(X),F_{x,n}]\|\leq J_F.                       \tag{57.12}
\]
Then
\[
 \boxed{
 \|[H_n,F_n(f)]\|
 \leq J_F\sum_x|w_{x,n}f(x)|.}                          \tag{57.13}
\]
If the quadrature weights give a uniform \(L^1\) bound for the
fixed test function \(f\), the second spectral moment is uniform.
\end{proposition}

\begin{proof}
Expand the commutator using (57.5) and (57.11), apply the triangle
inequality, and use (57.12) for each \(x\).
\end{proof}

This proposition applies to bounded cylinder transforms of fields.
Raw gauge, Higgs, or metric coordinate fields can be unbounded and
require domain and moment estimates before the commutator sum is
defined.

\section{Quasilocal interactions}

The exponential norm in (57.6) also controls a field whose support
has a collar tail.  If
\[
 F_n=\sum_YF_n(Y),\qquad
 \sum_Ye^{\mu d(Y,O)}\|F_n(Y)\|\leq M_{\mu,F},           \tag{57.14}
\]
then the same double-sum argument gives
\[
 \|[H_n,F_n]\|
 \leq C_\mu\|\Phi_n\|_\mu M_{\mu,F}.                    \tag{57.15}
\]
The constant depends on the lattice growth and overlap convention.
For polynomial locality, summability requires the decay order to
exceed the growth dimension, as in V34 and V48.

\section{Scaling obstruction}

\begin{countertheorem}[Ultralocality without a uniform coefficient
does not control the moment]
\label{cth:v15v57-scaling}
On one spin, let
\[
 H_n=n(I-\sigma_z),\qquad
 \Omega_n=|\!\uparrow\rangle,\qquad F_n=\sigma_x.       \tag{57.16}
\]
Then \(H_n\Omega_n=0\), every Hamiltonian is supported on one site,
and \(\|F_n\|=1\), but
\[
 \|H_nF_n\Omega_n\|=2n,\qquad
 \|[H_n,F_n]\|=2n.                                      \tag{57.17}
\]
Thus cutoffwise ultralocality gives no uniform second spectral
moment.
\end{countertheorem}

\begin{proof}
\(F_n\Omega_n=|\!\downarrow\rangle\), on which
\(I-\sigma_z\) has eigenvalue two.
\end{proof}

In lattice field theory, coefficients often contain powers of the
inverse lattice spacing.  The physical smearing and time
renormalization must cancel those powers in (57.13); calling the
term local does not do so.

\section{Complete-generator requirement}

Suppose the Hamiltonian is formally split into
\[
 H_n=H_n^{\rm bos}+H_n^{\rm f}
 +H_n^{\rm ct}+H_n^{\rm quot}.                          \tag{57.18}
\]
Then
\[
 [H_n,F_n]=
 [H_n^{\rm bos},F_n]+[H_n^{\rm f},F_n]
 +[H_n^{\rm ct},F_n]+[H_n^{\rm quot},F_n].              \tag{57.19}
\]
Bounding the four terms separately is sufficient but can be
wasteful if divergent local pieces cancel.  If cancellation is
required, it must be proved for the combined commutator on a common
domain before norms are taken.  A cancellation visible only after
taking expectations does not bound (57.3).

The counterterm and quotient pieces in (57.18) must arise from the
same reflected transfer factorization as the other terms.  Spatial
locality of a Euclidean action does not itself construct a
self-adjoint OS Hamiltonian decomposition.

\section{Higher commutators and a generator core}

For a local observable,
\[
 H_n^2F_n\Omega_n
 =H_n[H_n,F_n]\Omega_n
 =[H_n,[H_n,F_n]]\Omega_n,                              \tag{57.20}
\]
provided all domains are valid.  Iterating gives
\[
 H_n^kF_n\Omega_n=\operatorname{ad}_{H_n}^k(F_n)\Omega_n.
                                                               \tag{57.21}
\]
Uniform factorial bounds on these nested commutators would make
local-field vectors analytic for the generator and give a common
operator core.  V57 only establishes the \(k=1\) route required by
V56.

\begin{assessment}[Relation to previous locality]
\label{ass:v15v57-relation}
V34 proves rapid locality of smooth functions of a spatial
Dirac-square under a weighted operator hypothesis.  V57 concerns
commutators with the reconstructed time-translation Hamiltonian.
The two operators and metrics are not automatically the same.
One locality theorem cannot be substituted for the other without
the OS transfer construction of V51--V53.
\end{assessment}

\begin{assessment}[Physical status]
\label{ass:v15v57-status}
V57 supplies a concrete sufficient estimate for the second spectral
moment on bounded local observables.  The source does not construct
the complete Hamiltonian interaction \(\Phi_n\), its uniform
physical scaling norm (57.6), or the smeared commutator bound
(57.12).  It also does not prove that counterterm and quotient
pieces share the required self-adjoint domain.
\end{assessment}

\begin{decision}[V58 analytic local core]
\label{dec:v15v57-next}
Assume an exponentially local interaction with a uniform physical
norm and derive factorial bounds on nested commutators of a bounded
local observable.  Prove that the resulting local-field vectors are
uniform analytic vectors for \(H_n\), giving all finite spectral
moments and a practical common operator core.  Record the volume
and lattice-spacing dependence explicitly.
\end{decision}

\begin{firewall}[Claim boundary after V57]
\label{fire:v15v57}
V57 reduces V56's second moment to a uniform local commutator and
proves sufficient strict-local and smeared bounds.  The physical OS
Hamiltonian interaction and its cutoff-uniform scaled norm remain
unproved.
\end{firewall}

\section*{Version V57 append-only record}
\addcontentsline{toc}{section}{Version V57 append-only record}
The complete V56 source was retained before this continuation.  It
had \(553716\) bytes, \(15902\) lines, and SHA--256
\[
 \texttt{732A135A752399EDA6D20B7F60A00A0D4B08941906B82C8BB1892264879425DC}.
\]
V57 proved local-interaction and smeared-field commutator bounds
that imply uniform second spectral moments, and exhibited the
cutoff-scaling obstruction hidden by cutoffwise ultralocality.

\chapter{Fifteenth-cycle continuation V58: factorial commutator bounds and a uniform analytic local core}
\label{sec:v15v58}

\section{A scale of local observable algebras}

Let \(\{{\mathfrak A}_\mu:0<\mu\leq\mu_0\}\) be a decreasing scale
of Banach spaces of quasilocal observables, with
\[
 \|A\|_{\mu'}\leq\|A\|_\mu
 \quad(0<\mu'<\mu).                                     \tag{58.1}
\]
For a lattice decomposition \(A=\sum_XA(X)\), a model norm is
\[
 \|A\|_\mu=
 \sup_x\sum_{X\ni x}|X|e^{\mu\operatorname{diam}X}
 \|A(X)\|,                                               \tag{58.2}
\]
supplemented by a finite-support anchor norm when \(A\) is a single
local observable.  Let the Hamiltonian interaction \(\Phi_n\) have
a uniform exponential norm
\[
 \sup_n\|\Phi_n\|_{\mu_0}\leq J.                        \tag{58.3}
\]

\section{One-commutator loss}

\begin{lemma}[Weighted derivation estimate]
\label{lem:v15v58-derivation}
Assume the lattice family has uniformly bounded local incidence and
polynomial volume growth.  There is a geometric constant
\(C_{\rm geo}\) such that, for
\(0<\mu'<\mu\leq\mu_0\),
\[
 \boxed{
 \|[H_n,A]\|_{\mu'}
 \leq
 \frac{C_{\rm geo}\|\Phi_n\|_\mu}{\mu-\mu'}
 \|A\|_\mu.}                                            \tag{58.4}
\]
\end{lemma}

\begin{proof}
Expand
\[
 [H_n,A]=\sum_{X,Y:X\cap Y\neq\varnothing}
 [\Phi_n(X),A(Y)].                                      \tag{58.5}
\]
Each term is bounded by
\(2\|\Phi_n(X)\|\|A(Y)\|\).  If \(X\cap Y\neq\varnothing\),
\[
 \operatorname{diam}(X\cup Y)
 \leq\operatorname{diam}X+\operatorname{diam}Y.         \tag{58.6}
\]
The weighted sum at exponent \(\mu'\) is therefore a convolution of
the two sums at exponent \(\mu\), with a loss
\(e^{-(\mu-\mu')r}\) against the number of possible overlap
anchors at separation or diameter \(r\).  The incidence and volume
growth bound this number by \(C_{\rm geo}(1+r)^d\).  Absorb the
polynomial using
\[
 (1+r)^de^{-(\mu-\mu')r}
 \leq\frac{C_d}{(\mu-\mu')^{d+1}}.
\]
If the factor \(|X|\) in (58.2) is chosen with the corresponding
volume weight, all but one power are absorbed in the norm, leaving
the displayed loss.  Enlarging \(C_{\rm geo}\) covers bounded
diameters.  The bound is uniform in \(n\) because the incidence and
growth constants are uniform.
\end{proof}

The precise power of \((\mu-\mu')^{-1}\) depends on the chosen
weighted algebra.  V58 uses the scale normalized so that (58.4)
holds.  If a physical construction yields a higher power, the
iteration below gives a Gevrey rather than analytic core and must be
recorded accordingly.

\section{Factorial nested commutators}

Put
\[
 \delta_n(A)=[H_n,A].
\]

\begin{theorem}[Uniform analytic commutator bound]
\label{thm:v15v58-factorial}
Fix \(0<\mu_*<\mu_0\) and let
\(\Delta_\mu=\mu_0-\mu_*\).  Under (58.3)--(58.4), every
\(F_n\in{\mathfrak A}_{\mu_0}\) satisfies
\[
 \boxed{
 \|\delta_n^k(F_n)\|_{\mu_*}
 \leq k!
 \left(\frac{eC_{\rm geo}J}{\Delta_\mu}\right)^k
 \|F_n\|_{\mu_0}.}                                      \tag{58.7}
\]
\end{theorem}

\begin{proof}
For fixed \(k\), choose the equally spaced exponents
\[
 \mu_j=\mu_0-\frac{j}{k}\Delta_\mu,
 \qquad 0\leq j\leq k.                                  \tag{58.8}
\]
Apply (58.4) successively from \(\mu_{j-1}\) to \(\mu_j\).
Every loss is \(k/\Delta_\mu\), so
\[
 \|\delta_n^k(F_n)\|_{\mu_*}
 \leq
 \left(\frac{C_{\rm geo}Jk}{\Delta_\mu}\right)^k
 \|F_n\|_{\mu_0}.                                      \tag{58.9}
\]
The elementary Stirling bound \(k^k\leq e^kk!\) gives (58.7).
\end{proof}

\section{Analytic local-field vectors}

Assume \(H_n\Omega_n=0\).  Domain invariance gives
\[
 H_n^kF_n\Omega_n=\delta_n^k(F_n)\Omega_n.              \tag{58.10}
\]
Let
\[
 A_*=\frac{eC_{\rm geo}J}{\Delta_\mu}.                  \tag{58.11}
\]

\begin{corollary}[Uniform analytic radius]
\label{cor:v15v58-analytic}
If
\(\sup_n\|F_n\|_{\mu_0}\leq M_F\), then
\[
 \|H_n^kF_n\Omega_n\|
 \leq M_Fk!A_*^k                                       \tag{58.12}
\]
and
\[
 \sum_{k=0}^\infty\frac{r^k}{k!}
 \|H_n^kF_n\Omega_n\|
 \leq\frac{M_F}{1-rA_*}
 \quad(0\leq r<A_*^{-1})                               \tag{58.13}
\]
uniformly in \(n\).  Thus all these local-field vectors are analytic
for \(H_n\) with one common radius.
\end{corollary}

\begin{proof}
Use (58.10), the operator norm dominated by the local Banach norm,
and (58.7).  Summing the geometric series gives (58.13).
\end{proof}

\begin{corollary}[All local spectral moments]
\label{cor:v15v58-moments}
For every integer \(k\geq1\),
\[
 \int_0^\infty\lambda^{2k}\,
 \dd\langle F_n\Omega_n,E_{H_n}(\lambda)F_n\Omega_n\rangle
 \leq M_F^2(k!)^2A_*^{2k}.                              \tag{58.14}
\]
In particular, V56's second spectral moment is uniform.
\end{corollary}

\begin{proof}
The left side is
\(\|H_n^kF_n\Omega_n\|^2\).  Apply (58.12).
\end{proof}

\section{A common operator core}

Let \({\cal D}_n\) be the span of
\[
 F_{1,n}\cdots F_{r,n}\Omega_n                           \tag{58.15}
\]
for bounded observables in a source-complete local algebra.  Suppose
these vectors are dense in the OS Hilbert space and the products
have a uniform \({\mathfrak A}_{\mu_0}\) bound on every fixed
finite word.

\begin{theorem}[Analytic-vector core]
\label{thm:v15v58-core}
Every vector in \({\cal D}_n\) is analytic for \(H_n\).
Consequently the restriction of \(H_n\) to \({\cal D}_n\) is
essentially self-adjoint and \({\cal D}_n\) is an operator core.
\end{theorem}

\begin{proof}
Finite products of quasilocal observables remain in the weighted
Banach algebra, so Corollary \ref{cor:v15v58-analytic} applies.
Nelson's analytic-vector theorem states that a symmetric operator
with a dense invariant set of analytic vectors is essentially
self-adjoint on that set.  Since \(H_n\) is already self-adjoint,
the closure of its restriction is \(H_n\).
\end{proof}

To obtain one continuum core, the transported words in (58.15) must
converge and remain dense after the OS null quotient.  This is the
observable and Gram consistency requirement of V44 and V49.

\section{Physical scaling of the interaction norm}

The exponent \(\mu_0\) and interaction norm \(J\) must be measured in
physical units.  If \(a_n\) is the lattice spacing and a lattice-step
exponent is \(\widehat\mu_n\), the physical decay rate is
\[
 \mu_{\rm phys,n}=\frac{\widehat\mu_n}{a_n}.             \tag{58.16}
\]
A fixed \(\widehat\mu_n\) becomes stronger in physical distance,
whereas \(\widehat\mu_n\to0\) can collapse the available
\(\Delta_\mu\).  Likewise, if the transfer time step is
\(\tau_n\), the physical generator is scaled by \(\tau_n^{-1}\);
the interaction norm entering (58.3) is the norm after that
rescaling.

For a smeared physical field, the anchor norm
\(\|F_n\|_{\mu_0}\) must include the cell quadrature weights.  A
support containing \(a_n^{-d}\) sites can still have a uniform norm
when those weights and commutator strengths combine as in
Proposition \ref{prop:v15v57-smeared}; raw site counting cannot be
used.

\section{Gevrey fallback}

If the best derivation estimate is
\[
 \|\delta_nA\|_{\mu'}
 \leq
 \frac{CJ}{(\mu-\mu')^\rho}\|A\|_\mu,\qquad\rho>1,       \tag{58.17}
\]
the same subdivision gives
\[
 \|\delta_n^kF\|_{\mu_*}
 \leq C^k(k!)^\rho\|F\|_{\mu_0}.                        \tag{58.18}
\]
This supplies every finite spectral moment but the series (58.13)
can have zero analytic radius.  Such a Gevrey core may still be a
form core, but Nelson's analytic-vector conclusion cannot be invoked
from (58.18).

\begin{assessment}[Physical status]
\label{ass:v15v58-status}
V58 derives a uniform analytic local core from an exponentially
local physical Hamiltonian interaction with the scale estimate
(58.4).  The source does not construct that Hamiltonian, establish
its cutoff-uniform physical interaction norm, or prove that the
complete gauge--Higgs--gravity--fermion local grammar is dense after
the OS quotient.  The theorem is therefore a populated route only
after V51's reflected transfer has landed.
\end{assessment}

\begin{decision}[V59 spacetime translation and spectral visibility]
\label{dec:v15v58-next}
Assume the continuum measure and OS core are translation covariant.
Construct spatial translation unitaries and combine them with the
positive time semigroup.  State and prove the joint spectral
measure identities needed to distinguish a Hamiltonian gap from a
visible relativistic mass shell.  Record the additional Euclidean
rotation and analytic-continuation hypotheses needed for the
forward cone.
\end{decision}

\begin{firewall}[Claim boundary after V58]
\label{fire:v15v58}
V58 proves factorial nested-commutator bounds and an analytic
local-field core under a uniform physical interaction norm.  That
interaction, its scaling, and the source-complete OS density remain
unproved for the actual SM--Einstein regulator.
\end{firewall}

\section*{Version V58 append-only record}
\addcontentsline{toc}{section}{Version V58 append-only record}
The complete V57 source was retained before this continuation.  It
had \(562455\) bytes, \(16169\) lines, and SHA--256
\[
 \texttt{B0FC59962A41C5F4F98D27B341A32ADEF3125FC45F6ED819F77DDFB9C4E01860}.
\]
V58 derived factorial nested-commutator and spectral-moment bounds
from a weighted local-algebra derivation estimate and constructed a
uniform analytic local operator core, while recording the Gevrey
fallback and physical scaling debits.

\chapter{Fifteenth-cycle continuation V59: joint energy--momentum spectra and visible mass shells}
\label{sec:v15v59}

\section{Spatial translations on the OS space}

This section concerns a translation-invariant flat or asymptotically
translation-invariant sector.  A general curved Einstein background
has no global translation group and requires the local-covariant
alternative stated below.

Let \(\tau_a\), \(a\in\mathbb R^d\), act on continuum
configurations by spatial translation, commute with time reflection,
and preserve the complete continuum measure and positive-time
algebra.  On the OS core define
\[
 U(a)[F]=[F\circ\tau_{-a}].                              \tag{59.1}
\]

\begin{theorem}[Spatial unitary representation]
\label{thm:v15v59-spatial}
If the OS form is translation invariant and
\[
 \|[F\circ\tau_{-a}]-[F]\|_{\rm OS}\longrightarrow0
 \quad(a\to0)                                           \tag{59.2}
\]
on a dense local core, then \(U(a)\) extends to a strongly
continuous unitary representation of \(\mathbb R^d\).  Hence there
are strongly commuting self-adjoint momenta
\({\bf P}=(P_1,\ldots,P_d)\) such that
\[
 U(a)=e^{i a\cdot{\bf P}}.                              \tag{59.3}
\]
\end{theorem}

\begin{proof}
Translation invariance gives
\[
 \langle U(a)F,U(a)G\rangle_{\rm OS}
 =\langle F,G\rangle_{\rm OS},
\]
so \(U(a)\) preserves the null space and extends isometrically.
\(U(-a)\) is its inverse, hence it is unitary.  Strong continuity on
the core follows from (59.2), and the unitary bound extends it to
the completion.  The spectral theorem for strongly continuous
unitary representations of the abelian group \(\mathbb R^d\) gives
(59.3) with strongly commuting generators.
\end{proof}

If the positive time transfer commutes with spatial translations,
\[
 U(a)e^{-tH}=e^{-tH}U(a),                               \tag{59.4}
\]
then \(H,P_1,\ldots,P_d\) strongly commute and admit a joint
projection-valued measure
\[
 E(\dd E,\dd p)
 \quad\text{on }[0,\infty)\times\mathbb R^d.            \tag{59.5}
\]

\section{Euclidean two-point transform}

Let \(\Omega\) be the translation-invariant vacuum and
\(\psi=F\Omega\) for a local field \(F\).  Define the positive
spectral measure
\[
 \mu_\psi(A)=\langle\psi,E(A)\psi\rangle.               \tag{59.6}
\]

\begin{theorem}[Laplace--Fourier representation]
\label{thm:v15v59-two-point}
For \(t\geq0\) and \(a\in\mathbb R^d\),
\[
 \boxed{
 \langle\psi,e^{-tH}U(a)\psi\rangle
 =
 \int e^{-tE}e^{ia\cdot p}\,
       \mu_\psi(\dd E,\dd p).}                           \tag{59.7}
\]
The field sees a Borel spectral set \(A\) exactly when
\[
 \mu_\psi(A)=\|E(A)\psi\|^2>0.                          \tag{59.8}
\]
\end{theorem}

\begin{proof}
Apply the joint spectral theorem to the bounded function
\((E,p)\mapsto e^{-tE}e^{ia\cdot p}\).  Equation (59.8) is the
definition of the scalar spectral measure.
\end{proof}

Thus existence of a spectral subspace and visibility by a chosen
local field are separate.  Proposition
\ref{prop:v15v53-visibility} transports (59.8) from stable
finite-cutoff spectral projections.

\section{Energy gap versus invariant mass}

A Hamiltonian gap is the spectral exclusion
\[
 \sigma(H)\subset\{0\}\cup[m_E,\infty).                 \tag{59.9}
\]
It contains no statement about momentum.  A relativistic
energy--momentum representation additionally obeys the forward-cone
condition
\[
 \operatorname{supp}E\subset
 \{(E,p):E\geq|p|\}.                                    \tag{59.10}
\]
Only then is
\[
 M^2=H^2-|{\bf P}|^2                                   \tag{59.11}
\]
a nonnegative self-adjoint operator by joint functional calculus.

\begin{definition}[Visible isolated mass shell]
\label{def:v15v59-shell}
A local field vector \(\psi\) sees an isolated mass \(m>0\) if
there is \(\epsilon>0\) such that the joint spectrum in the relevant
sector has an isolated component on
\[
 E^2-|p|^2=m^2,\qquad E>0,                              \tag{59.12}
\]
and
\[
 E_{M^2}((m^2-\epsilon,m^2+\epsilon))\psi\neq0.          \tag{59.13}
\]
\end{definition}

\begin{proposition}[Mass-shell contribution to the two-point
function]
\label{prop:v15v59-shell}
If the restriction of \(\mu_\psi\) to the isolated shell has
nonzero measure \(\rho_\psi(\dd p)\), then (59.7) contains the
nonzero term
\[
 \int e^{-t\sqrt{|p|^2+m^2}}e^{ia\cdot p}\,
       \rho_\psi(\dd p).                                 \tag{59.14}
\]
\end{proposition}

\begin{proof}
Restrict the joint spectral integral (59.7) to the graph
\(E=\sqrt{|p|^2+m^2}\).  Nonzero visibility is precisely
\(\rho_\psi\neq0\).
\end{proof}

\section{A gap need not be a particle mass}

\begin{countertheorem}[Gapped nonrelativistic dispersion without an
isolated mass shell]
\label{cth:v15v59-gap-not-mass}
On
\[
 {\cal H}=\mathbb C\Omega\oplus L^2(\mathbb R^d),
\]
let the vacuum have \(H\Omega={\bf P}\Omega=0\), and on the second
summand define
\[
 ({\bf P}f)(p)=p f(p),\qquad
 (Hf)(p)=(m+|p|)f(p),\qquad m>0.                        \tag{59.15}
\]
Then \(H\) has energy gap \(m\) and commutes with translations, but
\[
 H^2-|{\bf P}|^2=m^2+2m|p|                              \tag{59.16}
\]
has continuous spectrum \([m^2,\infty)\) with no isolated positive
mass shell.
\end{countertheorem}

\begin{proof}
The spectrum of multiplication by \(m+|p|\) is
\([m,\infty)\).  Equation (59.16) is multiplication by a continuous
function whose essential range is \([m^2,\infty)\); \(m^2\) is a
threshold, not an isolated spectral component.
\end{proof}

This proves that the uniform semigroup gap in V52, even if
populated, is not yet a Standard Model particle spectrum.

\section{How the forward cone enters}

Spatial translation covariance and \(H\geq0\) give only
\(E\geq0\).  The sharper inequality \(E\geq|p|\) requires the
Euclidean rotation and reconstruction part of the OS theorem:
\begin{enumerate}
\item invariance under the relevant Euclidean group on a
      source-complete local algebra;
\item reflection positivity for every rotated time half-space, or
      the equivalent analytic continuation packet;
\item regularity allowing the local Euclidean representation to
      continue to a unitary Poincare representation;
\item compatibility of spin, chirality, gauge quotient, and
      determinant-line reflection with those rotations.
\end{enumerate}
Under these hypotheses, the reconstructed translations obey the
spectral condition and (59.10).  Positivity for one fixed reflection
plane alone does not prove it.

\section{Cofinal joint spectral screens}

Suppose finite-cutoff commuting pairs
\((H_n,{\bf P}_n)\) have joint spectral projections \(E_n(A)\) on a
common carrier.  If a compact isolated joint spectral set \(A\) has
a separating contour through a bounded continuous joint functional
calculus and
\[
 \|E_n(A)-E(A)\|_{\rm op}\to0,\qquad
 \psi_n\to\psi,                                         \tag{59.17}
\]
then
\[
 \|E_n(A)\psi_n\|\to\|E(A)\psi\|.                       \tag{59.18}
\]
This is the multi-operator version of Proposition
\ref{prop:v15v53-visibility}.  Establishing (59.17) requires joint
resolvent or functional-calculus convergence, not separate lists of
energy and momentum eigenvalues.

\section{General Einstein backgrounds}

On a generic dynamical curved spacetime, global spatial translations
and a global particle mass operator need not exist.  The appropriate
continuum target is then a locally covariant net:
\[
 (M,g)\longmapsto{\cal A}(M,g),                         \tag{59.19}
\]
with causal embeddings mapped to injective algebra morphisms,
locally covariant fields, a microlocal spectrum condition, and
stress-energy satisfying the Ward identity of V38.  Particle masses
can be recovered only in stationary, asymptotically flat, or other
sectors with an adequate asymptotic symmetry group.

Therefore the full SM--Einstein continuum cannot require one global
mass-shell theorem on every metric.  It must construct the local
covariant theory first and state particle interpretations on the
geometric sectors where they are defined.

\begin{assessment}[Physical status]
\label{ass:v15v59-status}
V59 constructs the joint spectral measure and proves the exact
Euclidean two-point representation, field-visibility criterion, and
mass-shell contribution.  It also proves that an energy gap need not
be an isolated invariant mass.

The source does not provide a unique translation-invariant continuum
measure, the complete Euclidean covariance and rotated-reflection
packet, joint spectral convergence, or a stationary/asymptotic
Einstein sector.  No physical Standard Model mass shell is claimed.
\end{assessment}

\begin{decision}[V60 compulsory audit]
\label{dec:v15v59-next}
At V60 inspect the newest YM and NS notes for a completed global
translation, reflected continuum, or local-covariant construction.
Then continue at V61 with the curved-spacetime alternative: define a
locally covariant net on a category of controlled globally
hyperbolic backgrounds and prove descent of the Ward-compatible
local observable algebra under isometric causal embeddings.
\end{decision}

\begin{firewall}[Claim boundary after V59]
\label{fire:v15v59}
V59 proves joint energy--momentum and visible-shell criteria under
translation and OS reconstruction hypotheses.  Those hypotheses,
the forward cone, an isolated shell, and the locally covariant
Einstein theory remain unproved.
\end{firewall}

\section*{Version V59 append-only record}
\addcontentsline{toc}{section}{Version V59 append-only record}
The complete V58 source was retained before this continuation.  It
had \(571857\) bytes, \(16447\) lines, and SHA--256
\[
 \texttt{FABD9A4018557661D630811622F3472114FC9A02053B04B7E07B32742D1072AC}.
\]
V59 constructed spatial translations and the joint
energy--momentum spectral measure, proved the Euclidean two-point
transform and visible mass-shell criterion, and separated the
translation-invariant particle branch from the local-covariant
curved-spacetime target.

\chapter{Fifteenth-cycle continuation V60: compulsory covariance audit}
\label{sec:v15v60}

\section{Companion result}

The active checkpoints remain
\[
\begin{array}{c|c|c|c}
\text{programme}&\text{file}&\text{lines}&\text{SHA--256}\\ \hline
\mathrm{YM}&
\texttt{Renewal\_Yang\_Mills\_Mass\_Gap\_Research\_Notes\_V39.tex}&
9076&
\texttt{9B44B0FBCA6037F3319E86036753B0F3B76BDD06F6045F00747E05DC01E23F9D}\\
\mathrm{NS}&
\texttt{Renewal\_NS\_Stability\_Research\_Notes\_V26.tex}&
5319&
\texttt{82C887F8C2C69E4F28FAD7C3DC8DE5B9099CB5A4BF431EBA1FFF3E91935978AD}.
\end{array}                                                \tag{60.1}
\]
YM V39 stops at a finite Wilson-pencil atlas and explicitly excludes
an interacting reflected continuum.  NS V26 remains a Euclidean
kernel-norm calculation.  Neither constructs a global translation
representation, forward-cone spectrum, locally covariant net, or
causal embedding functor.

\begin{theorem}[V60 nonimport]
\label{thm:v15v60-audit}
No companion theorem populates the translation or local-covariance
hypotheses of V59.  The curved-background net must be constructed
independently from the SM--Einstein local observable grammar.
\end{theorem}

\begin{proof}
Neither companion carrier includes globally hyperbolic background
embeddings or a functorial local observable algebra.
\end{proof}

\begin{decision}[V61 local-covariant descent]
\label{dec:v15v60-next}
Define a controlled category of globally hyperbolic backgrounds and
a local-symbol algebra on each object.  Prove that natural relation
ideals give a covariant algebra functor, that exact pullback of the
ideals gives injectivity, and that causal commutation and the
time-slice property descend.  Keep state selection and a dynamical
metric separate from the background-covariant algebra.
\end{decision}

\begin{firewall}[Claim boundary after V60]
\label{fire:v15v60}
V60 is a compulsory audit and imports no companion result into the
SM--Einstein continuum construction.
\end{firewall}

\section*{Version V60 append-only record}
\addcontentsline{toc}{section}{Version V60 append-only record}
The complete V59 source was retained before this continuation.  It
had \(581822\) bytes, \(16728\) lines, and SHA--256
\[
 \texttt{410C40DEBAA35334DA2A09EC2FB087DBCF5F7ABBE881169985285BCA60ADEB0A}.
\]
V60 performed the required companion audit and retained the
locally covariant curved-background algebra as the V61 target.

\chapter{Fifteenth-cycle continuation V61: locally covariant descent on controlled curved backgrounds}
\label{sec:v15v61}

\section{The controlled background category}

Let \({\bf Loc}_{\rm ESM}\) have as objects
\[
 {\mathsf M}=(M,g,\mathfrak o,\mathfrak t,P,E_H,\mathfrak s)   \tag{61.1}
\]
where \((M,g)\) is an oriented, time-oriented globally hyperbolic
four-manifold, \(P\) is the Standard Model gauge bundle, \(E_H\)
is the Higgs bundle, and \(\mathfrak s\) is the spin/chiral
structure and determinant-line datum.  The controlled designation
means that the geometric and bundle bounds required by the chosen
local construction are included in the object class.

A morphism
\[
 \psi:{\mathsf M}\longrightarrow{\mathsf N}             \tag{61.2}
\]
is an orientation- and time-orientation-preserving isometric
embedding with causally convex image, together with compatible
bundle and spin maps.  Composition is ordinary composition.

\section{Free local symbols and relation ideals}

Let \({\mathfrak F}({\mathsf M})\) be the unital star algebra
generated by bounded local gauge-invariant cylinder symbols
\[
 \Phi_{\mathsf M}(f),\qquad
 f\in{\cal T}_{\rm c}({\mathsf M}),                     \tag{61.3}
\]
including the declared gauge, Higgs, fermion-line, stress, and
curvature test types.  Pushforward of test data gives a free
algebra homomorphism
\[
 {\mathfrak F}(\psi):
 {\mathfrak F}({\mathsf M})\longrightarrow
 {\mathfrak F}({\mathsf N}),\qquad
 \Phi_{\mathsf M}(f)\mapsto\Phi_{\mathsf N}(\psi_*f).   \tag{61.4}
\]
Let \({\mathfrak I}({\mathsf M})\) be the two-sided star ideal
generated by the proved local field equations, gauge relations,
Ward identities, determinant-line patch relations, and causal
relations on \({\mathsf M}\).  Define
\[
 {\mathfrak A}({\mathsf M})
 ={\mathfrak F}({\mathsf M})/{\mathfrak I}({\mathsf M}). \tag{61.5}
\]

\section{Functorial quotient theorem}

\begin{theorem}[Natural ideals give a covariant algebra]
\label{thm:v15v61-functor}
Suppose for every morphism \(\psi:{\mathsf M}\to{\mathsf N}\),
\[
 {\mathfrak F}(\psi){\mathfrak I}({\mathsf M})
 \subset{\mathfrak I}({\mathsf N}).                     \tag{61.6}
\]
Then
\[
 {\mathfrak A}(\psi)[A]=[{\mathfrak F}(\psi)A]          \tag{61.7}
\]
is a well-defined unital star homomorphism, and
\({\mathfrak A}\) is a covariant functor.  It is injective exactly
when
\[
 \boxed{
 {\mathfrak F}(\psi)^{-1}{\mathfrak I}({\mathsf N})
 ={\mathfrak I}({\mathsf M}).}                          \tag{61.8}
\]
\end{theorem}

\begin{proof}
If \(A-B\in{\mathfrak I}({\mathsf M})\), (61.6) implies
\({\mathfrak F}(\psi)(A-B)\in{\mathfrak I}({\mathsf N})\), so
(61.7) is well defined.  Preservation of products, adjoints, units,
identities, and composition descends from the free maps.

The kernel of (61.7) consists of classes \([A]\) such that
\({\mathfrak F}(\psi)A\in{\mathfrak I}({\mathsf N})\).
It is zero precisely when every such \(A\) already belongs to
\({\mathfrak I}({\mathsf M})\), which is (61.8).
\end{proof}

Condition (61.6) expresses covariance of known relations.  The
stronger condition (61.8) says that embedding into a larger
spacetime creates no new relation among observables already
supported in the smaller one.  This exact pullback row is required
for isotony.

\section{Einstein causality}

\begin{theorem}[Causal commutation descends]
\label{thm:v15v61-causality}
Suppose \(\psi_i:{\mathsf M}_i\to{\mathsf N}\),
\(i=1,2\), have causally disjoint images and
\[
 [{\mathfrak F}(\psi_1)A_1,
   {\mathfrak F}(\psi_2)A_2]
 \in{\mathfrak I}({\mathsf N})                          \tag{61.9}
\]
for every \(A_i\in{\mathfrak F}({\mathsf M}_i)\), with the graded
commutator for fermionic symbols.  Then
\[
 [{\mathfrak A}(\psi_1)[A_1],
   {\mathfrak A}(\psi_2)[A_2]]=0                        \tag{61.10}
\]
in \({\mathfrak A}({\mathsf N})\).
\end{theorem}

\begin{proof}
The commutator in the quotient is the class of the free commutator
in (61.9), which vanishes because it belongs to the ideal.
\end{proof}

Locality of Euclidean kernels is not itself (61.9).  Causal
commutation is a Lorentzian or reconstructed algebraic relation and
must be transported through the OS continuation.

\section{Time-slice property}

Let \(\psi:{\mathsf M}\to{\mathsf N}\) have image containing a
Cauchy surface of \({\mathsf N}\).

\begin{theorem}[Algebraic time slice]
\label{thm:v15v61-timeslice}
Assume (61.8) and
\[
 {\mathfrak F}({\mathsf N})
 ={\mathfrak F}(\psi){\mathfrak F}({\mathsf M})
  +{\mathfrak I}({\mathsf N}).                          \tag{61.11}
\]
Then \({\mathfrak A}(\psi)\) is a star-algebra isomorphism.
\end{theorem}

\begin{proof}
Equation (61.8) gives injectivity.  For every
\(B\in{\mathfrak F}({\mathsf N})\), (61.11) writes
\(B={\mathfrak F}(\psi)A+I\) with
\(I\in{\mathfrak I}({\mathsf N})\).  Hence
\([B]={\mathfrak A}(\psi)[A]\), proving surjectivity.
\end{proof}

Equation (61.11) is normally proved by hyperbolic propagation of
the interacting field equations.  A Euclidean finite-volume
projective limit does not imply it.

\section{Naturality of fields and stress}

A locally covariant field is a family of maps
\[
 \Phi_{\mathsf M}:{\cal T}_{\rm c}({\mathsf M})
 \longrightarrow{\mathfrak A}({\mathsf M})              \tag{61.12}
\]
satisfying
\[
 {\mathfrak A}(\psi)\Phi_{\mathsf M}(f)
 =\Phi_{\mathsf N}(\psi_*f).                            \tag{61.13}
\]

\begin{proposition}[Ward-compatible stress descent]
\label{prop:v15v61-stress}
Suppose the renormalized stress symbols satisfy (61.13) and
\[
 T_{\mathsf M}(\nabla_{(\mu}X_{\nu)})=0                 \tag{61.14}
\]
in the quotient for every compactly supported vector field \(X\),
after the nonmetric equations are imposed.  Then the stress is
distributionally conserved on every object, and this conservation
is preserved by every morphism.
\end{proposition}

\begin{proof}
Equation (61.14) is the weak statement
\(\nabla_\mu T^{\mu\nu}=0\), by integration by parts.  Isometric
embeddings commute with covariant differentiation and pushforward
of compactly supported tensors.  Apply (61.13).
\end{proof}

The analytic trace-ideal and covariance-defect hypotheses of V38
are one route to (61.14).  They must hold naturally for the whole
background class, not only at one metric.

\section{Compatible states are additional data}

A locally covariant algebra does not select a state.  A compatible
state family would satisfy
\[
 \omega_{\mathsf M}
 =\omega_{\mathsf N}\circ{\mathfrak A}(\psi).            \tag{61.15}
\]

\begin{proposition}[State consistency from determining moments]
\label{prop:v15v61-state}
Let a countable natural local algebra determine the states on every
object.  If for every morphism and every determining symbol
\[
 \omega_{\mathsf N}
 ({\mathfrak A}(\psi)\Phi_{\mathsf M}(f))
 =\omega_{\mathsf M}(\Phi_{\mathsf M}(f)),              \tag{61.16}
\]
then (61.15) holds on the generated algebra.
\end{proposition}

\begin{proof}
Linearity and multiplicative generation extend (61.16) to the
countable algebra.  Continuity in the chosen algebra norm extends it
to the completion.
\end{proof}

If regulator measures converge separately on each background,
(61.15) still requires uniform embedding consistency of their local
moments.  A vacuum chosen by global time translations cannot be
natural on every curved spacetime.

\section{Background covariance is not quantum gravity}

The functor \({\mathfrak A}\) assigns an algebra to each classical
background metric.  The interacting SM--Einstein path integral of
V39 instead integrates over metric configurations.  To turn (61.5)
into a dynamical-gravity theory one additionally needs:
\begin{enumerate}
\item a groupoid or quotient of metric configurations with the
      noncollapse and tightness structure of V39--V43;
\item compatible local algebras over that groupoid;
\item an Einstein Schwinger--Dyson or relative-Cauchy-evolution
      identity for metric variations;
\item anomaly-free descent of the determinant line and measure.
\end{enumerate}
The background functor is the local-covariance target against which
those requirements can be stated; it does not solve them.

\section{Regulator descent criterion}

Let \({\mathfrak A}_n({\mathsf M})\) be regulator algebras with
coarse maps and local continuum symbols as in V44.  A sufficient
cofinal packet for (61.6)--(61.8) is:
\[
\begin{gathered}
\text{natural interpolation of every generator},\\
\text{vanishing adjacent moment and product defects},\\
\text{uniform convergence of every defining relation},\\
\text{a source-complete tail excluding new limiting relations}.
                                                               \tag{61.17}
\end{gathered}
\]
The last row is the algebraic counterpart of the form tail (54.9).
Without it, relations may appear only in the large-background
completion and violate injectivity.

\begin{assessment}[Physical status]
\label{ass:v15v61-status}
V61 proves the functorial quotient, injectivity, causal descent, and
time-slice criteria and connects the Ward-compatible stress to the
curved-background net.  It also separates compatible states and a
dynamical metric measure from background covariance.

The source does not provide a source-complete local ideal with exact
pullback (61.8), a hyperbolic propagation proof of (61.11), or the
cofinal regulator packet (61.17).  No locally covariant interacting
SM--Einstein net is therefore claimed.
\end{assessment}

\begin{decision}[V62 relative Cauchy evolution and the Einstein row]
\label{dec:v15v61-next}
Use the time-slice property to construct relative Cauchy evolution
for compactly supported metric perturbations.  Prove that its
derivative is generated by the renormalized stress tensor under an
explicit differentiability hypothesis, and formulate the
background-independence/Einstein condition as a vanishing total
metric response after including gravity, gauge--Higgs, fermion,
ghost, quotient, and counterterm sectors.
\end{decision}

\begin{firewall}[Claim boundary after V61]
\label{fire:v15v61}
V61 proves an abstract locally covariant descent theorem on
controlled backgrounds.  The physical relation ideals, time-slice
propagation, compatible interacting states, and dynamical Einstein
metric measure remain unproved.
\end{firewall}

\section*{Version V61 append-only record}
\addcontentsline{toc}{section}{Version V61 append-only record}
The complete V60 source was retained before this continuation.  It
had \(584244\) bytes, \(16791\) lines, and SHA--256
\[
 \texttt{07B021D3AD701FCE871E73C226E5A2B1113D6C6D5267E6E829EEF9612BAAF266}.
\]
V61 constructed the controlled curved-background category and
proved natural-ideal descent, exact-pullback injectivity, causal
commutation, the time-slice isomorphism, natural stress
conservation, and the separation between a background-covariant net
and a dynamical metric theory.

\chapter{Fifteenth-cycle continuation V62: relative Cauchy evolution and the quantum Einstein response}
\label{sec:v15v62}

\section{Metric perturbations between two Cauchy surfaces}

Let \({\mathsf M}=(M,g,\ldots)\) be an object of
\({\bf Loc}_{\rm ESM}\).  Let \(h\) be a compactly supported
symmetric tensor such that \(g+h\) remains globally hyperbolic and
belongs to the controlled object class.  Denote the perturbed object
by \({\mathsf M}[h]\).

Choose causally convex subspacetimes \({\mathsf M}^\pm\), lying to
the future and past of \(\operatorname{supp}h\), each containing a
Cauchy surface, and let
\[
 i^\pm:{\mathsf M}^\pm\to{\mathsf M},\qquad
 j^\pm:{\mathsf M}^\pm\to{\mathsf M}[h]                 \tag{62.1}
\]
be the canonical embeddings.  All four maps are time-slice
morphisms.

\section{Relative Cauchy evolution}

\begin{theorem}[Metric relative Cauchy evolution]
\label{thm:v15v62-rce}
If the functor of V61 has the time-slice property, then
\[
 \boxed{
 {\rm rce}_{\mathsf M}[h]
 =
 {\mathfrak A}(i^-)\,
 {\mathfrak A}(j^-)^{-1}\,
 {\mathfrak A}(j^+)\,
 {\mathfrak A}(i^+)^{-1}}                               \tag{62.2}
\]
is a star-automorphism of
\({\mathfrak A}({\mathsf M})\).  It is independent of shrinking
\({\mathsf M}^\pm\) within the causal past and future regions,
provided the corresponding diagrams commute.
\end{theorem}

\begin{proof}
The time-slice property makes each algebra morphism in (62.2) an
isomorphism, so their composition is an automorphism.  For two
admissible choices of past or future subspacetime, take a smaller
common Cauchy neighborhood.  Functoriality makes both compositions
factor through the same algebra maps; the extra maps cancel with
their inverses.  Hence the result is independent of the choice.
\end{proof}

This construction compares the same local physics under a compact
change of background metric.  It requires the exact time-slice
property (61.11), not only causal isotony.

\section{Infinitesimal stress generator}

Let \(k\) be a compactly supported symmetric test tensor and suppose
\(s\mapsto{\rm rce}_{\mathsf M}[sk](A)\) is differentiable at zero
on a dense star subalgebra.  Define
\[
 \delta_k(A)=
 \left.\frac{\dd}{\dd s}\right|_{s=0}
 {\rm rce}_{\mathsf M}[sk](A).                          \tag{62.3}
\]

\begin{proposition}[The metric response is a star derivation]
\label{prop:v15v62-derivation}
The map \(\delta_k\) satisfies
\[
 \delta_k(AB)=\delta_k(A)B+A\delta_k(B),\qquad
 \delta_k(A^*)=\delta_k(A)^*.                           \tag{62.4}
\]
\end{proposition}

\begin{proof}
Differentiate preservation of products and adjoints by the
automorphisms in (62.2).
\end{proof}

\begin{theorem}[Stress implementation criterion]
\label{thm:v15v62-stress}
Suppose the relative Cauchy evolution is implemented in a
representation by differentiable unitaries
\[
 {\rm rce}_{\mathsf M}[sk](A)
 =U_k(s)AU_k(s)^*,\qquad
 U_k(0)=I,\qquad
 U_k'(0)=\frac{i}{2}T_{\rm ren}(k),                     \tag{62.5}
\]
on a common invariant core, where \(T_{\rm ren}(k)\) is symmetric.
Then
\[
 \boxed{
 \delta_k(A)=\frac{i}{2}[T_{\rm ren}(k),A].}            \tag{62.6}
\]
\end{theorem}

\begin{proof}
Differentiate (62.5), use
\((U_k^*)'(0)=-\frac{i}{2}T_{\rm ren}(k)\), and collect the two
terms.
\end{proof}

The factor and sign in (62.5) fix the convention.  A different
definition of stress changes both consistently.  Existence of the
unitary implementation and common domain is an additional theorem;
the algebraic derivation (62.3) need not be inner.

\section{Complete renormalized metric-response ledger}

Let the Euclidean effective action be
\[
 \Gamma_{\rm tot}
 =S_{\rm grav}^{\rm ren}
 +S_{\rm gauge,H}^{\rm ren}
 +\Gamma_{\rm f}^{\rm ren}
 +\Gamma_{\rm ghost}^{\rm ren}
 +\Gamma_{\rm quot}^{\rm ren}
 +S_{\rm ct}^{\rm fin}.                                 \tag{62.7}
\]
All terms use the same metric, bundle, regulator identification,
and subtraction scheme.  Define their metric responses by
\[
 D_g\Gamma_j[k]
 =\frac12\int_M\sqrt g\,
 {\cal T}_j^{\mu\nu}k_{\mu\nu}.                         \tag{62.8}
\]
The total Einstein residual is
\[
 {\cal R}_{\rm Ein}^{\mu\nu}
 =\sum_j{\cal T}_j^{\mu\nu}.                            \tag{62.9}
\]
The convention includes the variation of the gravitational action
inside \({\cal T}_{\rm grav}\); after fixing Lorentzian/Euclidean
signs and Newton's constant, (62.9) is the usual Einstein tensor
plus cosmological, higher-curvature, and matter terms.

\begin{theorem}[Classical effective Einstein equation]
\label{thm:v15v62-classical-einstein}
If \(\Gamma_{\rm tot}\) is differentiable in every compactly
supported metric direction and \(g\) is a stationary point, then
\[
 \boxed{{\cal R}_{\rm Ein}^{\mu\nu}=0}                  \tag{62.10}
\]
as a distribution.
\end{theorem}

\begin{proof}
Stationarity says
\[
 0=D_g\Gamma_{\rm tot}[k]
 =\frac12\int\sqrt g\,{\cal R}_{\rm Ein}^{\mu\nu}
 k_{\mu\nu}
\]
for every compactly supported smooth \(k\).  This is precisely the
distributional equality (62.10).
\end{proof}

The fermion contribution in (62.8) is supplied conditionally by
the \(C^1\) trace-ideal theorem of V38.  The \(a_0,a_2,a_4\)
variations from V37 belong to the gravitational, gauge--Higgs, and
counterterm rows and must not be counted twice.

\section{Quantum metric integration by parts}

Let a finite regulator have metric coordinates \(g^A\), a positive
reference density \(\rho_n(g)\), and effective action
\({\cal A}_n(g,\Phi)\).  Include the derivative of the reference
density by defining
\[
 {\cal S}_{n,A}
 =\partial_{g^A}{\cal A}_n
 -\partial_{g^A}\log\rho_n.                             \tag{62.11}
\]

\begin{theorem}[Finite-cutoff metric Schwinger--Dyson identity]
\label{thm:v15v62-sd}
Suppose \(F\rho_ne^{-{\cal A}_n}\) has zero boundary flux in the
\(g^A\) direction and both terms below are integrable.  Then
\[
 \boxed{
 \mathbb E_{\nu_n}[\partial_{g^A}F]
 =
 \mathbb E_{\nu_n}[F\,{\cal S}_{n,A}].}                 \tag{62.12}
\]
\end{theorem}

\begin{proof}
Integrate
\[
 \partial_{g^A}
 \bigl(F\rho_ne^{-{\cal A}_n}\bigr)
\]
over the metric coordinate.  The boundary term vanishes by
hypothesis.  Expanding the derivative and dividing by the partition
function gives (62.12).
\end{proof}

For \(F=1\), (62.12) gives
\[
 \mathbb E_{\nu_n}{\cal S}_{n,A}=0.                     \tag{62.13}
\]
This is an averaged quantum equation.  It does not state that
\({\cal S}_{n,A}=0\) on every sampled metric.

\section{Continuum quantum Einstein identity}

\begin{theorem}[Cofinal Schwinger--Dyson closure]
\label{thm:v15v62-sd-limit}
Let \(F_n\to F\), let the directional derivatives and scores in
(62.12) transport to continuum quantities
\(D_kF\) and \({\cal R}_{\rm Ein}(k)\), and assume the two product
families are uniformly integrable.  If the measures converge as in
V44, then
\[
 \boxed{
 \mathbb E_\nu[D_kF]
 =
 \mathbb E_\nu[F\,{\cal R}_{\rm Ein}(k)]}               \tag{62.14}
\]
for every local cylinder \(F\) and compactly supported metric
variation \(k\).
\end{theorem}

\begin{proof}
Use compact tightness sublevels.  On each sublevel, convergence of
the transported observables and scores passes both expectations in
(62.12) to the limit.  Uniform integrability makes the complement
uniformly negligible.  Let the sublevel grow to obtain (62.14).
\end{proof}

Equation (62.14), together with gauge and diffeomorphism Ward
identities, is the quantum Einstein response.  Equation (62.10) is
the stationary effective-action branch.  They coincide only under
an additional concentration or semiclassical limit.

\section{Background independence}

In a split-background construction, a compact change of the
background accompanied by the inverse change of the dynamical
metric should leave physical observables unchanged.  Infinitesimally
this requires cancellation between the relative Cauchy derivation
and the dynamical metric response:
\[
 \delta_k^{\rm bg}(A)+\delta_{-k}^{\rm dyn}(A)=0.        \tag{62.15}
\]
If both derivations are inner on a common core, (62.15) says the
complete renormalized stress generator, including ghosts, quotient,
and counterterms, acts trivially on physical observables.  Omitting
one row can produce a false background-independence statement.

\begin{proposition}[Pure-gauge perturbations give the Ward row]
\label{prop:v15v62-pure-gauge}
If \(k={\cal L}_Xg\) for a compactly supported vector field \(X\)
and the functor is diffeomorphism covariant without anomaly, then
\[
 \delta_{{\cal L}_Xg}(A)=0                              \tag{62.16}
\]
on diffeomorphism-invariant observables.  Under (62.6), this implies
\[
 [T_{\rm ren}({\cal L}_Xg),A]=0.                        \tag{62.17}
\]
\end{proposition}

\begin{proof}
The perturbed background is related to the original by the
compactly supported diffeomorphism generated by \(X\).  Naturality
identifies the two embeddings, so their relative Cauchy evolution
is the identity on invariant observables.  Differentiate and use
(62.6).
\end{proof}

\section{Anomaly and boundary ledger}

Each of the following can obstruct (62.14)--(62.17):
\begin{enumerate}
\item a determinant/Pfaffian line without covariant reflected
      orientation;
\item a gauge or gravitational anomaly in the chiral measure;
\item a Gribov or metric-quotient boundary flux in (62.12);
\item a nonuniform heat-counterterm derivative;
\item failure of uniform integrability of the stress insertion;
\item a conformal direction in which the metric measure is not
      normalizable.
\end{enumerate}
These are independent rows.  Formal cancellation of representation
anomaly coefficients does not prove absence of quotient boundary
flux or trace-ideal failure.

\begin{assessment}[Physical status]
\label{ass:v15v62-status}
V62 constructs relative Cauchy evolution, proves its stress
implementation, and separates the classical effective Einstein
equation from the quantum Schwinger--Dyson identity.  It also gives
the exact cofinal convergence and background-split conditions.

The source does not prove the physical time-slice property, unitary
implementation (62.5), metric integration-by-parts boundary
condition, uniform stress integrability, or full anomaly-free
measure.  Neither (62.10) nor (62.14) has therefore been populated
for the complete continuum.
\end{assessment}

\begin{decision}[V63 stress-insertion uniform integrability]
\label{dec:v15v62-next}
Prove a sufficient Orlicz or entropy criterion for uniform
integrability of the metric score and its products with bounded
local observables.  Use the determinant \(C^1\) trace-ideal response
of V38 and the bosonic coercivity of V39 to isolate a concrete
exponential stress moment.  Show why convergence in probability of
the score is insufficient for passing (62.12).
\end{decision}

\begin{firewall}[Claim boundary after V62]
\label{fire:v15v62}
V62 proves the relative-Cauchy and Einstein-response compilers.
Their physical time-slice, differentiability, boundary, anomaly,
and uniform-integrability hypotheses remain unproved.
\end{firewall}

\section*{Version V62 append-only record}
\addcontentsline{toc}{section}{Version V62 append-only record}
The complete V61 source was retained before this continuation.  It
had \(595097\) bytes, \(17087\) lines, and SHA--256
\[
 \texttt{CD34919A1695F4E2D4EDEE26E0EEECE27A52357415C89E3511964BBF15575CA0}.
\]
V62 constructed relative Cauchy evolution, derived its stress
generator, formulated the complete classical metric residual, and
proved finite-cutoff and cofinal quantum Einstein
Schwinger--Dyson identities under explicit boundary and uniform
integrability hypotheses.

\chapter{Fifteenth-cycle continuation V63: uniform integrability of the metric score}
\label{sec:v15v63}

\section{Purpose and probability spaces}

V62 reduced passage of the finite-cutoff metric
Schwinger--Dyson identity to one measure-theoretic obstruction:
uniform integrability of the differentiated observable and of the
metric-score insertion.  This chapter supplies a quantitative
compiler for that obstruction.  It does not infer stress
integrability merely from weak convergence of the measures.

For each cutoff \(n\), let \((\Omega_n,\nu_n)\) be the normalized
regulated configuration space and let \(k\) be a fixed compactly
supported smooth metric variation.  Write
\[
 {\cal S}_n(k)
 =
 D_k{\cal A}_n-D_k\log\rho_n                         \tag{63.1}
\]
for the complete score, including gravitational, matter,
determinant, ghost, quotient, reference-density, and finite
counterterm rows.  Let \(V_n\geq 0\) be the coercive control
functional used in the tightness argument of V39--V43.

\section{Three uniform-integrability criteria}

\begin{definition}[Uniform integrability]
\label{def:v15v63-ui}
A family \(X_n\in L^1(\nu_n)\) is uniformly integrable when
\[
 \lim_{R\to\infty}\sup_n
 \int_{\{|X_n|>R\}}|X_n|\,d\nu_n=0.                  \tag{63.2}
\]
The definition allows the probability spaces to vary with \(n\).
\end{definition}

\begin{theorem}[Superlinear moment criterion]
\label{thm:v15v63-dvp}
Let \(\Theta:[0,\infty)\to[0,\infty)\) be increasing and satisfy
\[
 \lim_{t\to\infty}\frac{\Theta(t)}{t}=\infty.         \tag{63.3}
\]
If
\[
 \sup_n\mathbb E_{\nu_n}\Theta(|X_n|)<\infty,         \tag{63.4}
\]
then \((X_n)\) is uniformly integrable.
\end{theorem}

\begin{proof}
Put
\[
 m(R)=\inf_{t\geq R}\frac{\Theta(t)}{t}.
\]
Condition (63.3) gives \(m(R)\to\infty\).  On
\(\{|X_n|>R\}\),
\[
 |X_n|\leq m(R)^{-1}\Theta(|X_n|).
\]
Consequently the left side of (63.2) is at most
\[
 m(R)^{-1}\sup_n\mathbb E_{\nu_n}\Theta(|X_n|),
\]
which tends to zero.
\end{proof}

\begin{corollary}[Exponential stress moment]
\label{cor:v15v63-exp-ui}
If some \(\lambda>0\) obeys
\[
 \sup_n\mathbb E_{\nu_n}e^{\lambda|{\cal S}_n(k)|}
 <\infty,                                             \tag{63.5}
\]
then the score family is uniformly integrable.  If in addition
\(\sup_n\|F_n\|_\infty\leq M\), then
\((F_n{\cal S}_n(k))_n\) is uniformly integrable.
\end{corollary}

\begin{proof}
Use Theorem~\ref{thm:v15v63-dvp} with
\(\Theta(t)=e^{\lambda t}-1\).  For the product, if
\(|F_n{\cal S}_n(k)|>R\), then
\(|{\cal S}_n(k)|>R/M\), and
\[
 |F_n{\cal S}_n(k)|
 \leq M|{\cal S}_n(k)|.
\]
The score tail at \(R/M\) tends to zero uniformly.
\end{proof}

Bounded cylinder observables therefore require only a score
estimate.  Unbounded field insertions require a paired moment
condition.  The following form records the precise pairing.

\begin{proposition}[Orlicz product compiler]
\label{prop:v15v63-orlicz-product}
Let \(\Phi,\Psi\) be complementary Young functions, so
\[
 ab\leq\Phi(a)+\Psi(b),\qquad a,b\geq0.               \tag{63.6}
\]
Suppose there are constants \(a,b>0\) and a Young function
\(\Xi\) with \(\Xi(t)/t\to\infty\) such that
\[
 \sup_n\mathbb E_{\nu_n}\Phi(a|F_n|)<\infty,
 \qquad
 \sup_n\mathbb E_{\nu_n}\Psi(b|{\cal S}_n(k)|)<\infty, \tag{63.7}
\]
and
\[
 \sup_n\mathbb E_{\nu_n}
 \Xi\!\left(
 \Phi(a|F_n|)+\Psi(b|{\cal S}_n(k)|)
 \right)<\infty.                                     \tag{63.8}
\]
Then \((F_n{\cal S}_n(k))_n\) is uniformly integrable.
\end{proposition}

\begin{proof}
Apply (63.6) to \(a|F_n|\) and \(b|{\cal S}_n(k)|\):
\[
 ab|F_n{\cal S}_n(k)|
 \leq
 Y_n:=
 \Phi(a|F_n|)+\Psi(b|{\cal S}_n(k)|).                 \tag{63.9}
\]
Theorem~\ref{thm:v15v63-dvp} and (63.8) make
\((Y_n)\) uniformly integrable.  Domination by \(Y_n/(ab)\)
then makes the product uniformly integrable.
\end{proof}

The extra hypothesis (63.8) is deliberate.  Separate bounded
Orlicz norms alone control \(L^1\) norms of products but need not
control their tails uniformly.  A convenient stronger row is a
uniform \(L^p\)-\(L^q\) estimate with slack: if
\[
 \sup_n\mathbb E|F_n|^{p(1+\epsilon)}<\infty,\qquad
 \sup_n\mathbb E|{\cal S}_n(k)|^{q(1+\epsilon)}<\infty,
 \quad p^{-1}+q^{-1}=1,                               \tag{63.10}
\]
then H\"older gives a uniform \(L^{1+\epsilon}\) bound on the
product, hence uniform integrability.

\section{Coercivity-to-stress transfer}

\begin{theorem}[Coercive domination criterion]
\label{thm:v15v63-coercive-score}
Assume there are constants \(\eta>0\) and \(C_V<\infty\) such that
\[
 \sup_n\mathbb E_{\nu_n}e^{\eta V_n}\leq C_V.          \tag{63.11}
\]
For the fixed variation \(k\), assume the complete score obeys
\[
 |{\cal S}_n(k)|\leq a_kV_n+b_k                       \tag{63.12}
\]
outside a \(\nu_n\)-null set, with \(a_k,b_k\) independent of
\(n\).  If \(a_k>0\), every
\(0<\lambda\leq\eta/a_k\) satisfies
\[
 \sup_n\mathbb E_{\nu_n}e^{\lambda|{\cal S}_n(k)|}
 \leq e^{\lambda b_k}C_V.                             \tag{63.13}
\]
If \(a_k=0\), (63.13) holds for every \(\lambda>0\), with right
side \(e^{\lambda b_k}\).  Thus bounded local observables satisfy
the uniform-integrability row in V62.
\end{theorem}

\begin{proof}
Equation (63.12) gives
\[
 e^{\lambda|{\cal S}_n(k)|}
 \leq e^{\lambda b_k}e^{\lambda a_kV_n}.
\]
For \(\lambda a_k\leq\eta\), the last exponential is bounded by
\(e^{\eta V_n}\).  Take expectations and use (63.11), then apply
Corollary~\ref{cor:v15v63-exp-ui}.
\end{proof}

\begin{remark}[Local variations]
\label{rem:v15v63-local}
The constants may depend on a fixed seminorm of \(k\), for example
\[
 a_k+b_k
 \leq C_K\sum_{j=0}^{r}\|\nabla^jk\|_{L^\infty(K)}
                                                               \tag{63.14}
\]
when \(\operatorname{supp}k\subset K\).  Uniformity is required in
the cutoff \(n\), not over all test tensors at once.
\end{remark}

\section{The determinant-stress row}

Let the renormalized relative fermion action on the chosen branch
have the form
\[
 \Gamma^{\rm ren}_{f,n}(g)
 =
 -\log\det_{\!F}(I+K_n(g))+C_{f,n}(g),                 \tag{63.15}
\]
where \(K_n(g)\) is trace class on the common carrier and
\(C_{f,n}\) contains the prescribed local subtraction.  The sign
can be changed for a different determinant convention without
affecting the estimates below.

\begin{proposition}[Trace-ideal stress estimate]
\label{prop:v15v63-det-stress}
Suppose \(g\mapsto K_n(g)\) is differentiable in trace norm in the
direction \(k\), \(I+K_n(g)\) is invertible, and
\[
 \|(I+K_n)^{-1}\|\leq c_0,                             \tag{63.16}
\]
\[
 \|D_kK_n\|_1+|D_kC_{f,n}|
 \leq a_{f,k}V_n+b_{f,k}                              \tag{63.17}
\]
with cutoff-independent constants.  Then
\[
 |D_k\Gamma^{\rm ren}_{f,n}|
 \leq
 (c_0+1)(a_{f,k}V_n+b_{f,k}).                         \tag{63.18}
\]
\end{proposition}

\begin{proof}
The \(C^1\) Fredholm determinant formula from V38 gives
\[
 D_k\log\det_{\!F}(I+K_n)
 =
 \operatorname{Tr}\bigl((I+K_n)^{-1}D_kK_n\bigr).      \tag{63.19}
\]
The trace ideal inequality yields
\[
 \left|
 \operatorname{Tr}\bigl((I+K_n)^{-1}D_kK_n\bigr)
 \right|
 \leq c_0\|D_kK_n\|_1.                                \tag{63.20}
\]
Add the counterterm derivative and use (63.17).  The displayed
constant \(c_0+1\) is a harmless common upper bound for the two
coefficients.
\end{proof}

Thus the determinant row enters
Theorem~\ref{thm:v15v63-coercive-score} once (63.16)--(63.17)
are proved.  A trace-norm derivative alone is insufficient near a
zero of the determinant: the inverse in (63.16) can diverge.  The
zero-mode sector must instead be removed, saturated by insertions,
or controlled through the determinant-line construction of V51.

\begin{proposition}[Rowwise assembly]
\label{prop:v15v63-rowwise}
Suppose
\[
 {\cal S}_n(k)=\sum_{\alpha=1}^{N}{\cal S}^{(\alpha)}_n(k)
                                                               \tag{63.21}
\]
is the complete finite list of score rows and for each \(\alpha\)
\[
 |{\cal S}^{(\alpha)}_n(k)|
 \leq a_{\alpha,k}V_n+b_{\alpha,k}.                   \tag{63.22}
\]
Then (63.12) holds with
\[
 a_k=\sum_\alpha a_{\alpha,k},\qquad
 b_k=\sum_\alpha b_{\alpha,k}.                        \tag{63.23}
\]
\end{proposition}

\begin{proof}
Use the triangle inequality in (63.21).
\end{proof}

The assembly must include the derivative of the reference density
in (63.1).  Absorbing that derivative into an unnamed measure can
hide a cutoff-growing stress contribution.

\section{Entropy transport of exponential moments}

The coercive exponential moment may first be available under a
comparison measure \(\mu_n\), rather than \(\nu_n\).  The next
inequality imports it quantitatively.

\begin{lemma}[Entropy inequality]
\label{lem:v15v63-entropy}
Let \(\nu\ll\mu\) be probability measures and let
\[
 {\rm Ent}(\nu\mid\mu)
 =
 \int\log\!\left(\frac{d\nu}{d\mu}\right)d\nu<\infty.
\]
For every measurable \(X\) with
\(\mathbb E_\mu e^{\lambda X}<\infty\) and \(\lambda>0\),
\[
 \mathbb E_\nu X
 \leq
 \frac{
 {\rm Ent}(\nu\mid\mu)+
 \log\mathbb E_\mu e^{\lambda X}}
 {\lambda}.                                           \tag{63.24}
\]
\end{lemma}

\begin{proof}
Write \(f=d\nu/d\mu\).  Jensen's inequality under \(\nu\) gives
\[
 \lambda\mathbb E_\nu X-{\rm Ent}(\nu\mid\mu)
 =
 \mathbb E_\nu\log\frac{e^{\lambda X}}{f}
 \leq
 \log\mathbb E_\nu\frac{e^{\lambda X}}{f}
 =
 \log\mathbb E_\mu e^{\lambda X}.
\]
Rearrange.
\end{proof}

\begin{corollary}[Moment transfer under bounded likelihood]
\label{cor:v15v63-entropy-transfer}
Assume for some \(r>1\)
\[
 \sup_n\mathbb E_{\mu_n}e^{r\lambda|X_n|}<\infty,
 \qquad
 \sup_n
 \mathbb E_{\mu_n}
 \left(\frac{d\nu_n}{d\mu_n}\right)^{r/(r-1)}
 <\infty.                                             \tag{63.25}
\]
Then
\[
 \sup_n\mathbb E_{\nu_n}e^{\lambda|X_n|}<\infty.       \tag{63.26}
\]
\end{corollary}

\begin{proof}
Let \(f_n=d\nu_n/d\mu_n\).  H\"older's inequality with conjugate
exponents \(r\) and \(r/(r-1)\) gives
\[
 \mathbb E_{\nu_n}e^{\lambda|X_n|}
 =
 \mathbb E_{\mu_n}\bigl[f_ne^{\lambda|X_n|}\bigr]
 \leq
 \|f_n\|_{r/(r-1)}
 \left(\mathbb E_{\mu_n}e^{r\lambda|X_n|}\right)^{1/r}.
\]
Both factors are uniformly bounded by (63.25).
\end{proof}

Lemma~\ref{lem:v15v63-entropy} controls first moments from relative
entropy.  Corollary~\ref{cor:v15v63-entropy-transfer} records the
stronger likelihood integrability actually needed to transport an
exponential moment.  A mere uniform entropy bound does not by
itself imply (63.26).

\section{Vitali passage to the continuum}

\begin{theorem}[Score expectation closure]
\label{thm:v15v63-vitali}
Place transported variables
\[
 X_n=F_n{\cal S}_n(k)
\]
and \(X=F{\cal R}_{\rm Ein}(k)\) on a common probability space by
a coupling of \(\nu_n\) and \(\nu\).  Suppose
\[
 X_n\longrightarrow X\quad\hbox{in probability}       \tag{63.27}
\]
and \((X_n)\) is uniformly integrable.  Then \(X\in L^1\) and
\[
 \lim_{n\to\infty}\mathbb E_{\nu_n}
 [F_n{\cal S}_n(k)]
 =
 \mathbb E_{\nu}[F{\cal R}_{\rm Ein}(k)].              \tag{63.28}
\]
The same statement applies to the transported derivative
\(D_kF_n\).
\end{theorem}

\begin{proof}
Fix \(\epsilon>0\).  Uniform integrability gives \(R\) such that
the \(X_n\)-tails beyond \(R\) have \(L^1\) mass below
\(\epsilon\).  Fatou's lemma applied along an almost surely
convergent subsequence of (63.27) gives the corresponding
integrability and tail bound for \(X\).  On the region where both
variables have absolute value at most \(R\), convergence in
probability implies \(L^1\) convergence because the variables are
bounded; this follows by taking almost surely convergent
subsequences and applying dominated convergence.  The two tails
contribute at most a constant multiple of \(\epsilon\).  Every
subsequence has a further subsequence with convergent expectations,
so the full sequence satisfies (63.28).
\end{proof}

Together with the exact finite-cutoff identity (62.12),
Theorem~\ref{thm:v15v63-vitali} proves (62.14) whenever the score
domination, observable product, derivative, and boundary rows all
hold.

\begin{counterexample}[Probability convergence is insufficient]
\label{cnt:v15v63-probability}
On \(([0,1],dx)\), set
\[
 X_n(x)=n\,{\bf1}_{[0,1/n]}(x).                       \tag{63.29}
\]
Then \(X_n\to0\) in probability, but
\[
 \mathbb E X_n=1\quad\hbox{for every }n.               \tag{63.30}
\]
Indeed,
\(\mathbb E[X_n{\bf1}_{\{X_n>R\}}]=1\) whenever
\(n>R\), so the family is not uniformly integrable.  This is the
precise concentration defect excluded by (63.5) or (63.10).
\end{counterexample}

\section{Verification ledger}

For each compactly supported \(k\), the continuum quantum Einstein
identity now reduces to the following finite list:
\begin{enumerate}
\item a uniform coercive moment (63.11);
\item rowwise domination (63.22), including reference density;
\item a determinant inverse bound and trace derivative
      (63.16)--(63.17), or a zero-mode replacement;
\item an observable product estimate, bounded or of the form
      (63.10);
\item transported convergence in probability for the score product
      and differentiated observable;
\item vanishing finite-cutoff boundary flux.
\end{enumerate}
Rows one through five close the uniform-integrability part of V62.
The sixth is logically independent and will be attacked next.

\begin{assessment}[Physical status]
\label{ass:v15v63-status}
V63 proves that a cutoff-uniform exponential coercive moment and
linear domination of every metric-score row imply the uniform
integrability required by the continuum Schwinger--Dyson theorem.
It also isolates the determinant inverse bound at zero modes and
proves the relevant Vitali passage.

The programme has not established (63.11)--(63.17) for the complete
interacting Standard Model--Einstein regulator.  In particular,
conformal instability, quotient singularities, and determinant
zeros can violate the proposed domination.  The result is a
rigorous sufficient compiler, not a proof that the physical
continuum measure satisfies it.
\end{assessment}

\begin{decision}[V64 metric-boundary flux]
\label{dec:v15v63-next}
Construct a cutoff-and-exhaustion integration-by-parts theorem for
noncompact metric coordinates.  Derive vanishing outer flux from a
coercive exponential tail and a controlled logarithmic derivative,
while treating quotient or Gribov boundaries as separate interior
faces.  Exhibit a one-dimensional boundary counterexample showing
why tightness alone does not eliminate flux.
\end{decision}

\begin{firewall}[Claim boundary after V63]
\label{fire:v15v63}
V63 closes the measure-theoretic implication from explicit
coercive stress bounds to uniform integrability.  It does not prove
those analytic bounds, finite-cutoff boundary cancellation, anomaly
cancellation, or existence of the full interacting continuum.
\end{firewall}

\section*{Version V63 append-only record}
\addcontentsline{toc}{section}{Version V63 append-only record}
The complete V62 source was retained before this continuation.  It
had \(606706\) bytes, \(17420\) lines, and SHA--256
\[
 \texttt{2E824C93B743A7B4E341565AE0F41C230171274BB409A8E466B078CBC3D2F850}.
\]
V63 proved superlinear and exponential uniform-integrability
criteria, a rowwise coercivity-to-stress theorem, determinant
trace-ideal control, entropy and likelihood moment transport, and
the Vitali passage required for the continuum quantum Einstein
identity.

\chapter{Fifteenth-cycle continuation V64: metric exhaustion and quotient-boundary flux}
\label{sec:v15v64}

\section{Geometric score at finite cutoff}

The coordinate calculation in V62 must be refined when a metric
variation is represented by a nonconstant vector field on the
finite-dimensional regulator.  Let
\(\Omega_n\subset{\mathbb R}^{d_n}\) have piecewise smooth boundary,
let
\[
 d\nu_n=Z_n^{-1}w_n(x)\,dx,\qquad
 w_n=\rho_n e^{-{\cal A}_n},                           \tag{64.1}
\]
and let \(B_n\) be the regulator vector field induced by a fixed
compactly supported continuum variation \(k\).  Define
\[
 {\cal S}_{n,B}
 =
 D_{B_n}{\cal A}_n-D_{B_n}\log\rho_n
 -\operatorname{div}B_n.                              \tag{64.2}
\]
For a constant coordinate vector field the last term vanishes and
(64.2) reduces to (63.1).

\begin{proposition}[Cutoff divergence identity]
\label{prop:v15v64-cutoff-div}
Let \(\chi\) be compactly supported in the closure of \(\Omega_n\).
Assume all displayed terms are integrable and traces on the smooth
boundary faces exist.  Then
\[
 \begin{split}
 \mathbb E_{\nu_n}[\chi D_{B_n}F]
 &=
 \mathbb E_{\nu_n}[\chi F{\cal S}_{n,B}]
 -\mathbb E_{\nu_n}[F D_{B_n}\chi]\\
 &\quad+
 Z_n^{-1}\int_{\partial\Omega_n}
 \chi Fw_n\,B_n\mathbin{\cdot}{\bf n}\,d\sigma .
                                                               \tag{64.3}
 \end{split}
\]
\end{proposition}

\begin{proof}
The divergence theorem applied to
\(\chi Fw_nB_n\) gives
\[
 \int_{\partial\Omega_n}\chi Fw_nB_n\mathbin{\cdot}{\bf n}\,d\sigma
 =
 \int_{\Omega_n}\operatorname{div}(\chi Fw_nB_n)\,dx.  \tag{64.4}
\]
Since
\[
 D_{B_n}\log w_n
 =
 D_{B_n}\log\rho_n-D_{B_n}{\cal A}_n,
\]
expanding (64.4), moving terms, and dividing by \(Z_n\)
gives (64.3).
\end{proof}

The divergence row in (64.2) is part of the metric score.  It
cannot be omitted when a background variation is transported
through a nonlinear chart or a gauge slice.

\section{Uniform removal of the boundary at infinity}

Let \(V_n:\Omega_n\to[0,\infty)\) be proper at fixed cutoff and
choose a smooth function \(\vartheta:[0,\infty)\to[0,1]\) with
\[
 \vartheta=1\ \hbox{on }[0,1],\qquad
 \vartheta=0\ \hbox{on }[2,\infty),\qquad
 \|\vartheta'\|_\infty\leq C_\vartheta.                \tag{64.5}
\]
Set
\[
 \chi_{n,R}(x)=\vartheta(V_n(x)/R).                    \tag{64.6}
\]

\begin{lemma}[Exponential tails absorb polynomial collars]
\label{lem:v15v64-tail}
If
\[
 \sup_n\mathbb E_{\nu_n}e^{\eta V_n}\leq C_V           \tag{64.7}
\]
for some \(\eta>0\), then for every \(q\geq0\) there is
\(C_{q,\eta}<\infty\) such that
\[
 \sup_n\mathbb E_{\nu_n}
 \left[(1+V_n)^q{\bf1}_{\{V_n\geq R\}}\right]
 \leq C_{q,\eta}C_Ve^{-\eta R/2}.                     \tag{64.8}
\]
\end{lemma}

\begin{proof}
The continuous function
\((1+t)^qe^{-\eta t/2}\) is bounded on \([0,\infty)\);
call its supremum \(C_{q,\eta}\).  On \(t\geq R\),
\[
 (1+t)^q
 \leq C_{q,\eta}e^{\eta t/2}
 \leq C_{q,\eta}e^{-\eta R/2}e^{\eta t}.
\]
Substitute \(t=V_n\), integrate, and use (64.7).
\end{proof}

\begin{theorem}[Cofinal exhaustion theorem]
\label{thm:v15v64-exhaustion}
Assume (64.7).  Suppose for fixed \(k\) and some \(p,m\geq0\),
\[
 |F_n|\leq C_F(1+V_n)^p,\qquad
 |D_{B_n}V_n|\leq C_B(1+V_n)^m                       \tag{64.9}
\]
with constants independent of \(n\).  Then
\[
 \lim_{R\to\infty}\sup_n
 \left|\mathbb E_{\nu_n}
 [F_nD_{B_n}\chi_{n,R}]\right|=0.                     \tag{64.10}
\]
If the genuine boundary term in (64.3) vanishes for every \(R\),
and \(D_{B_n}F_n\) and \(F_n{\cal S}_{n,B}\) are uniformly
integrable, then
\[
 \mathbb E_{\nu_n}[D_{B_n}F_n]
 =
 \mathbb E_{\nu_n}[F_n{\cal S}_{n,B}]                 \tag{64.11}
\]
for every \(n\), with the exhaustion limit uniform along the
cofinal sequence.
\end{theorem}

\begin{proof}
By (64.5)--(64.6),
\[
 |D_{B_n}\chi_{n,R}|
 \leq
 \frac{C_\vartheta C_B}{R}(1+V_n)^m
 {\bf1}_{\{R\leq V_n\leq2R\}}.                        \tag{64.12}
\]
Multiplication by the first bound in (64.9), followed by
Lemma~\ref{lem:v15v64-tail} with \(q=p+m\), bounds the
absolute value in (64.10) by
\[
 \frac{C_\vartheta C_BC_FC_{p+m,\eta}C_V}{R}
 e^{-\eta R/2}.                                       \tag{64.13}
\]
This tends to zero uniformly in \(n\).

Use \(\chi=\chi_{n,R}\) in (64.3).  The true boundary term is zero
by hypothesis, and (64.10) removes the artificial collar.
Uniform integrability permits \(\chi_{n,R}\uparrow1\) in the other
two expectations, proving (64.11).  The estimates are uniform in
\(n\).
\end{proof}

\begin{remark}[What coercivity proves]
\label{rem:v15v64-coercivity}
Theorem~\ref{thm:v15v64-exhaustion} removes only the artificial
outer boundary introduced by truncating a noncompact coordinate.
It says nothing about a finite endpoint, a gauge-fixing horizon, a
change-of-chart seam, or a singular quotient stratum.
\end{remark}

\section{Paired quotient faces}

Suppose the regulator is represented by a fundamental region
\(\Omega_n\) whose regular boundary contains face pairs
\((\Sigma_a^+,\Sigma_a^-)\).  Let
\(\tau_a:\Sigma_a^+\to\Sigma_a^-\) be their identification.
Write
\[
 \omega_{n,B}=w_n(B_n\mathbin{\cdot}{\bf n})\,d\sigma \tag{64.14}
\]
for the oriented flux density on a face.

\begin{theorem}[Cancellation of identified faces]
\label{thm:v15v64-face-pair}
Assume, for every paired face, that
\[
 \tau_a^*(F_n|_{\Sigma_a^-})=F_n|_{\Sigma_a^+},        \tag{64.15}
\]
\[
 \tau_a^*(\omega_{n,B}|_{\Sigma_a^-})
 =
 -\omega_{n,B}|_{\Sigma_a^+}.                         \tag{64.16}
\]
Then the two faces make zero total contribution to (64.3).
\end{theorem}

\begin{proof}
Changing variables by \(\tau_a\) on \(\Sigma_a^-\) and using
(64.15)--(64.16) gives
\[
 \int_{\Sigma_a^-}\chi F_n\,\omega_{n,B}
 =
 -\int_{\Sigma_a^+}(\chi\circ\tau_a^{-1})F_n\,
 \omega_{n,B}.                                        \tag{64.17}
\]
Choose the exhaustion invariant under the face identification, so
\(\chi\circ\tau_a^{-1}=\chi\).  The right side cancels the
\(\Sigma_a^+\) integral.
\end{proof}

\begin{corollary}[Reflecting and vanishing-density faces]
\label{cor:v15v64-reflecting}
An unpaired regular face gives zero flux if either
\[
 B_n\mathbin{\cdot}{\bf n}=0                          \tag{64.18}
\]
there, or the trace \(w_nB_n\mathbin{\cdot}{\bf n}\)
vanishes in \(L^1_{\rm loc}(d\sigma)\).  The second condition may
follow from a vanishing Faddeev--Popov density only after the normal
vector field and the remaining factors are shown not to diverge
too rapidly.
\end{corollary}

\begin{proof}
In either case the oriented density (64.14) vanishes almost
everywhere, hence so does its integral against bounded compactly
supported \(\chi F_n\).
\end{proof}

\section{Singular strata and tubular flux}

A quotient singularity need not possess an ordinary boundary
trace.  Let \(S_n\) be a closed singular stratum and
\(N_\epsilon(S_n)\) its \(\epsilon\)-tube inside \(\Omega_n\).

\begin{proposition}[Removable singular stratum]
\label{prop:v15v64-tube}
Suppose the divergence identity holds on
\(\Omega_n\setminus N_\epsilon(S_n)\), and for every fixed
exhaustion radius \(R\),
\[
 \lim_{\epsilon\downarrow0}
 \int_{\partial N_\epsilon(S_n)}
 |\chi_{n,R}F_n|\,w_n
 |B_n\mathbin{\cdot}{\bf n}_\epsilon|\,d\sigma=0.       \tag{64.19}
\]
Then \(S_n\) contributes no defect to (64.3).
\end{proposition}

\begin{proof}
Apply Proposition~\ref{prop:v15v64-cutoff-div} on the punctured
domain.  Its new boundary is
\(\partial N_\epsilon(S_n)\), with the appropriate orientation.
The absolute value of that contribution is bounded by (64.19).
Let \(\epsilon\downarrow0\).
\end{proof}

Codimension at least two is not by itself a proof of (64.19):
the density or the normal component of \(B_n\) may have a
nonintegrable inverse-power singularity.

\section{Boundary defect and the quantum Einstein identity}

Define the regulated defect functional
\[
 {\cal B}_{n,B}(F)
 =
 Z_n^{-1}\int_{\partial\Omega_n}
 Fw_nB_n\mathbin{\cdot}{\bf n}\,d\sigma,              \tag{64.20}
\]
after paired faces and removable strata have been combined.
The correct finite-cutoff equation is
\[
 \mathbb E_{\nu_n}[D_{B_n}F]
 =
 \mathbb E_{\nu_n}[F{\cal S}_{n,B}]
 +{\cal B}_{n,B}(F).                                  \tag{64.21}
\]

\begin{theorem}[Boundary-aware continuum closure]
\label{thm:v15v64-boundary-limit}
Assume the convergence and uniform-integrability hypotheses of
Theorem~\ref{thm:v15v63-vitali}, the exhaustion hypotheses of
Theorem~\ref{thm:v15v64-exhaustion}, and
\[
 {\cal B}_{n,B}(F_n)\longrightarrow{\cal B}_B(F).      \tag{64.22}
\]
Then the continuum identity is
\[
 \mathbb E_\nu[D_kF]
 =
 \mathbb E_\nu[F{\cal R}_{\rm Ein}(k)]
 +{\cal B}_B(F).                                      \tag{64.23}
\]
Consequently the defect-free quantum Einstein identity (62.14)
holds precisely on the tested class when
\({\cal B}_B(F)=0\) there.
\end{theorem}

\begin{proof}
Use (64.21), pass its two expectations to the limit by V63, and
use (64.22) for the remaining scalar term.
\end{proof}

\begin{counterexample}[Tightness does not remove a finite face]
\label{cnt:v15v64-halfline}
Let \(\Omega=[0,\infty)\), \(w(x)=e^{-x}\), \(B=1\), and
\(F=1\).  This is a normalized probability measure with
\[
 \mathbb E e^{\eta x}=\frac1{1-\eta}<\infty
 \quad(0<\eta<1).                                     \tag{64.24}
\]
Here \({\cal A}(x)=x\), \(\rho=1\), and
\({\cal S}_B=1\).  Nevertheless,
\[
 \mathbb E[D_BF]=0,\qquad
 \mathbb E[F{\cal S}_B]=1,                            \tag{64.25}
\]
because the outward normal at \(x=0\) is \(-1\) and
\[
 {\cal B}_B(F)=-1.                                    \tag{64.26}
\]
Thus even an exponential coercive moment does not erase a genuine
finite boundary.
\end{counterexample}

\section{Application ledger for metric and gauge quotients}

At each cutoff the boundary row is discharged only after the
following checks:
\begin{enumerate}
\item choose a proper \(V_n\) and prove the uniform directional
      derivative bound in (64.9);
\item pair chart or gauge-copy faces and verify both observable
      matching (64.15) and oriented-density matching (64.16);
\item impose or derive tangency at a reflecting face;
\item estimate every Gribov-horizon trace, including possible
      inverse determinants from the score;
\item prove tubular flux decay at each singular orbit-type stratum;
\item show the residual defect tends to zero rather than merely
      remaining bounded.
\end{enumerate}

The companion Yang--Mills programme may eventually provide data
for items two through five.  A finite Wilson-pencil atlas alone
does not establish these measure-theoretic quotient identities.

\begin{assessment}[Physical status]
\label{ass:v15v64-status}
V64 proves a uniform exhaustion theorem that removes artificial
metric boundaries at infinity under the coercive moment of V63.
It also gives exact sufficient conditions for paired-face,
reflecting-face, and singular-stratum cancellation, and carries a
nonzero defect into the continuum identity when cancellation
fails.

The complete Standard Model--Einstein regulator has not yet been
shown to possess a global fundamental domain satisfying the face
matching conditions, nor have its Gribov and metric-quotient
tubular fluxes been bounded.  These remain physical construction
obligations.
\end{assessment}

\begin{decision}[V65 cross-program audit]
\label{dec:v15v64-next}
Perform the scheduled five-version audit of the newest
Yang--Mills and Navier--Stokes notes.  Test whether either companion
now supplies quotient-face, local-kernel, or compactness estimates
that strengthen (64.15)--(64.19), update the dependency ledger,
and then continue with a substantive V66.
\end{decision}

\begin{firewall}[Claim boundary after V64]
\label{fire:v15v64}
V64 proves conditional removal and propagation of boundary flux.
It does not prove that the physical gauge and metric quotients have
zero flux, that the conformal measure is coercive, or that the
continuum quantum Einstein identity is populated.
\end{firewall}

\section*{Version V64 append-only record}
\addcontentsline{toc}{section}{Version V64 append-only record}
The complete V63 source was retained before this continuation.  It
had \(622111\) bytes, \(17898\) lines, and SHA--256
\[
 \texttt{EEC16062CBE1A1276D9BFAE071647B07FC6B7E6F200490FF87B18CEE736AE487}.
\]
V64 derived the geometric metric score, proved uniform removal of
the exhaustion collar, characterized cancellation of quotient
faces and singular strata, and retained any residual boundary flux
as an explicit continuum quantum-Einstein defect.

\chapter{Fifteenth-cycle continuation V65: scheduled YM--NS dependency audit}
\label{sec:v15v65}

\section{Audit protocol}

This is the compulsory five-version comparison following V60.
The active companion files in the shared notes directory were read
at their claim boundaries and hashed before choosing the next
Standard Model--Einstein step.  The audited snapshots are
\[
\begin{array}{c|c|c|l}
\hbox{programme}&\hbox{version}&\hbox{lines}&\hbox{SHA--256}\\ \hline
\hbox{Yang--Mills}&V39&9076&
\texttt{9B44B0FBCA6037F3319E86036753B0F3B76BDD06F6045F00747E05DC01E23F9D}\\
\hbox{Navier--Stokes}&V26&5319&
\texttt{82C887F8C2C69E4F28FAD7C3DC8DE5B9099CB5A4BF431EBA1FFF3E91935978AD}.
\end{array}                                             \tag{65.1}
\]
These hashes are identical to the V60 audit.  No companion theorem
has therefore appeared since that checkpoint.

\section{Yang--Mills comparison}

The terminal Yang--Mills theorem gives an exact positive lower
bound for a \(45\times45\) Wilson pencil on
\[
 [0,\pi/4]^2\times[0,t]                               \tag{65.2}
\]
in the specified transverse variables, including the two-face
corner.  Its proof uses finitely many rational chart centres and
exact Gram-residual certificates.  The third transverse direction
still has width \(t\), and the note expressly leaves global torus
propagation, the interacting measure, reflection positivity,
clustering, and the Hamiltonian mass gap open.

\begin{proposition}[What the YM plate can and cannot import]
\label{prop:v15v65-ym-import}
The finite Wilson-pencil floor can serve as a local symbol
invertibility input on the covered momentum plate.  By itself it
implies none of the following SM--Einstein rows:
\[
\begin{gathered}
 \sup_n\|(I+K_n)^{-1}\|<\infty,\\
 \sup_n\mathbb E_{\nu_n}e^{\eta V_n}<\infty,\\
 {\cal B}_{n,B}(F)\to0,\qquad
 \hbox{or a reflection-positive continuum measure}.   \tag{65.3}
\end{gathered}
\]
\end{proposition}

\begin{proof}
The first line in (65.3) is a global inverse estimate on the full
regulated carrier, whereas (65.2) covers only a proper momentum
plate and only the stated finite pencil.  The second and third
lines concern probability tails and boundary traces, data absent
from a pointwise matrix inequality.  Reflection positivity and
continuum convergence require compatible measures across cutoffs,
also absent from the finite plate theorem.
\end{proof}

The exact chart method remains useful as a future certification
mechanism: once the missing global analytic reduction is proved,
finite rational Gram certificates may verify its compact symbol
remainder.  Using those certificates before that reduction would
confuse a local matrix floor with a continuum operator bound.

\section{Navier--Stokes comparison}

The active Navier--Stokes note computes pointwise rank-one norms for
the first and second derivatives of the Oseen kernel.  The first
rank-one norm equals the full symmetric-traceless norm, while the
second is modestly smaller on an intermediate radial range.  Its
own firewall leaves full stability and global regularity open.

\begin{proposition}[No metric-tail or quotient transfer]
\label{prop:v15v65-ns-import}
The Oseen rank-one formulas do not supply (64.7), (64.16), or
(64.19) for the SM--Einstein regulator.
\end{proposition}

\begin{proof}
Those formulas are deterministic convolution-kernel estimates on
\(\mathbb R^3\).  Equation (64.7) is an exponential probability
moment on a metric-field configuration space; (64.16) is an
oriented identification of quotient-face densities; and (64.19)
is a tubular trace estimate near a regulator singular stratum.
No map identifying the corresponding measures, faces, or normal
fields is present in the NS source.
\end{proof}

The reusable lesson is methodological: exploit the tensor rank of
the physical input before taking a full operator norm, but certify
the angular maximization and radial tails independently.  This may
sharpen a future local stress estimate; it does not create the
missing metric probability law.

\section{Dependency decision}

\begin{theorem}[Audit conclusion]
\label{thm:v15v65-conclusion}
Relative to V60, the cross-program dependency graph is unchanged.
Neither active companion discharges a hypothesis of the V63
uniform-integrability theorem or the V64 quotient-boundary theorem.
The nearest upstream SM--Einstein obstruction is therefore the
absence of a coercive Euclidean metric weight, in particular the
conformal direction already isolated in V39.
\end{theorem}

\begin{proof}
The hashes in (65.1) prove that the audited sources are unchanged.
Propositions~\ref{prop:v15v65-ym-import} and
\ref{prop:v15v65-ns-import} compare their strongest terminal
results with the exact missing hypotheses in V63--V64.  None
matches.  Both V63 and V64 depend on the exponential metric moment,
whose known upstream obstruction is conformal noncoercivity.
\end{proof}

\begin{assessment}[Physical status]
\label{ass:v15v65-status}
The audit prevents a false import from finite symbol positivity or
deterministic kernel bounds.  It adds no physical continuum
existence theorem.  The YM and NS programmes remain useful future
sources only after a precise interface theorem is supplied.
\end{assessment}

\begin{decision}[V66 conformal coercivity]
\label{dec:v15v65-next}
Go upstream to the Euclidean conformal factor.  Derive its exact
wrong-sign gradient contribution under \(g=e^{2\phi}\bar g\);
construct an explicit high-frequency sequence driving the
Einstein--Hilbert action downward; then test a positive
\(R^2\)-type stabilizer and fixed-volume or mean-zero slice for a
genuine Sobolev coercivity estimate.  Keep contour rotation,
higher-derivative modification, and Lorentzian reconstruction as
distinct physical choices.
\end{decision}

\begin{firewall}[Claim boundary after V65]
\label{fire:v15v65}
V65 is a dependency audit.  It proves no new continuum measure,
quotient cancellation, conformal stabilization, mass gap, or
global regularity result.
\end{firewall}

\section*{Version V65 append-only record}
\addcontentsline{toc}{section}{Version V65 append-only record}
The complete V64 source was retained before this continuation.  It
had \(634680\) bytes, \(18284\) lines, and SHA--256
\[
 \texttt{CF3D3ED26C054161D5D41744FEBCDBB2262E627956F69ACCFE3F15563953223C}.
\]
V65 performed the scheduled companion audit, recorded exact source
identities and claim boundaries, rejected unsupported transfers,
and selected the Euclidean conformal mode as the next upstream
obstruction.

\chapter{Fifteenth-cycle continuation V66: the conformal factor and coercive completions}
\label{sec:v15v66}

\section{Exact conformal reduction in four dimensions}

Let \((M,\bar g)\) be a closed smooth Riemannian four-manifold and
write
\[
 g=e^{2\phi}\bar g=u^2\bar g,\qquad u=e^\phi>0.         \tag{66.1}
\]
We use \(\bar\Delta=\operatorname{div}_{\bar g}\bar\nabla\), so
\(-\bar\Delta\) is nonnegative.  The conformal transformation laws
are
\[
 dV_g=u^4dV_{\bar g},\qquad
 R_g=u^{-3}L_{\bar g}u,\qquad
 L_{\bar g}:=-6\bar\Delta+\bar R.                      \tag{66.2}
\]

\begin{proposition}[Einstein--Hilbert conformal action]
\label{prop:v15v66-eh}
With \(\kappa=(16\pi G)^{-1}>0\) and Euclidean convention
\[
 S_{\rm EH}^{E}(g)
 =
 -\kappa\int_M(R_g-2\Lambda)\,dV_g,                   \tag{66.3}
\]
one has
\[
 \boxed{
 S_{\rm EH}^{E}(u^2\bar g)
 =
 -\kappa\left[
 \int_M\bigl(\bar Ru^2+6|\bar\nabla u|^2\bigr)
 dV_{\bar g}
 -2\Lambda\int_Mu^4dV_{\bar g}
 \right].}                                            \tag{66.4}
\]
\end{proposition}

\begin{proof}
Equations (66.2) give
\[
 \int_MR_g\,dV_g
 =
 \int_MuL_{\bar g}u\,dV_{\bar g}.
\]
Because \(M\) has no boundary,
\[
 -\int_Mu\bar\Delta u\,dV_{\bar g}
 =
 \int_M|\bar\nabla u|^2\,dV_{\bar g}.
\]
Substitution into (66.3) proves (66.4).
\end{proof}

Equivalently, in the logarithmic variable,
\[
 \int_MR_g\,dV_g
 =
 \int_Me^{2\phi}
 \bigl(\bar R+6|\bar\nabla\phi|^2\bigr)dV_{\bar g}.
                                                               \tag{66.5}
\]
The minus sign in (66.3) therefore makes the conformal gradient
direction unstable.

\section{A fixed-volume instability sequence}

\begin{theorem}[Conformal action is unbounded below at fixed volume]
\label{thm:v15v66-instability}
On a flat four-torus, fix \(V_0>0\) and \(0<a<1\).  There is a
sequence of smooth positive functions \(u_N\) satisfying
\[
 \int_Mu_N^4\,dV_{\bar g}=V_0                           \tag{66.6}
\]
and
\[
 S_{\rm EH}^{E}(u_N^2\bar g)\longrightarrow-\infty.   \tag{66.7}
\]
\end{theorem}

\begin{proof}
Use a periodic coordinate \(x_1\) of period \(2\pi\) and set
\[
 u_N=c_a(1+a\sin Nx_1),                               \tag{66.8}
\]
where \(c_a>0\) is chosen so that (66.6) holds.  The integral of
\((1+a\sin Nx_1)^4\) is independent of the positive integer \(N\),
so the same \(c_a\) works for every \(N\).  Positivity follows from
\(a<1\).  Moreover,
\[
 \int_M|\bar\nabla u_N|^2dV_{\bar g}
 =
 C_aN^2                                                   \tag{66.9}
\]
for a constant \(C_a>0\).  Since \(\bar R=0\) and the cosmological
term is constant on (66.6), (66.4) becomes
\[
 S_{\rm EH}^{E}(u_N^2\bar g)
 =
 -6\kappa C_aN^2+2\kappa\Lambda V_0,
\]
which proves (66.7).
\end{proof}

Thus fixing volume removes the constant conformal direction but
does not repair the oscillatory one.  No probability-tail estimate
of the form (63.11) can follow from the bare Euclidean
Einstein--Hilbert action.

\section{What a curvature-square term controls}

In four dimensions,
\[
 {\cal I}_2(u)
 :=
 \int_MR_g^2\,dV_g
 =
 \int_Mu^{-2}|L_{\bar g}u|^2\,dV_{\bar g}.             \tag{66.10}
\]

\begin{theorem}[Global critical estimate from \(R^2\)]
\label{thm:v15v66-r2-critical}
On the fixed-volume slice (66.6),
\[
 \|L_{\bar g}u\|_{L^{4/3}(\bar g)}
 \leq
 V_0^{1/4}{\cal I}_2(u)^{1/2}.                        \tag{66.11}
\]
Consequently there is a background constant \(C_{\bar g}\) such
that
\[
 \|u\|_{W^{2,4/3}(\bar g)}
 \leq
 C_{\bar g}\left[
 V_0^{1/4}{\cal I}_2(u)^{1/2}
 +V_0^{1/4}
 \right].                                             \tag{66.12}
\]
\end{theorem}

\begin{proof}
Put \(q=u^{-1}L_{\bar g}u\).  Equation (66.10) says
\(\|q\|_2={\cal I}_2(u)^{1/2}\), while
\(L_{\bar g}u=uq\).  H\"older's inequality with
\(1/(4/3)=1/4+1/2\) gives
\[
 \|L_{\bar g}u\|_{4/3}
 \leq\|u\|_4\|q\|_2,
\]
which is (66.11).  The elliptic estimate
\[
 \|u\|_{W^{2,4/3}}
 \leq C_{\bar g}
 \bigl(\|L_{\bar g}u\|_{4/3}+\|u\|_{4/3}\bigr)
\]
and the finite-volume embedding \(L^4\subset L^{4/3}\) give
(66.12), after enlarging \(C_{\bar g}\).
\end{proof}

The exponent is critical:
\[
 W^{2,4/3}(M)\longrightarrow L^4(M).                  \tag{66.13}
\]
The estimate does not give the \(W^{2,2+\delta}\) carrier required
in V42.

\begin{counterexample}[Curvature-square concentration]
\label{cnt:v15v66-mobius}
Let \((S^4,g_0)\) be the round sphere.  In a stereographic chart,
the dilation \(D_\lambda(x)=\lambda x\) satisfies
\[
 D_\lambda^*g_0=u_\lambda^2g_0,\qquad
 u_\lambda(x)=
 \frac{\lambda(1+|x|^2)}{1+\lambda^2|x|^2}.            \tag{66.14}
\]
Because \(D_\lambda\) is a diffeomorphism,
\[
 \int_{S^4}u_\lambda^4\,dV_{g_0}
 =\operatorname{Vol}(S^4,g_0),\qquad
 {\cal I}_2(u_\lambda)
 =\int_{S^4}R_{g_0}^2\,dV_{g_0}.                      \tag{66.15}
\]
For every fixed \(x\ne0\), \(u_\lambda(x)\to0\), while its
\(L^4\) norm is constant.  Hence no subsequence converges strongly
in \(L^4\).  Bounded volume and \(R^2\) action therefore do not give
pre-quotient compactness.
\end{counterexample}

\begin{proof}
The round metric in stereographic coordinates is
\[
 g_0=\frac{4}{(1+|x|^2)^2}\,dx^2.
\]
Pulling it back by \(D_\lambda\) gives
\[
 D_\lambda^*g_0
 =
 \frac{4\lambda^2}{(1+\lambda^2|x|^2)^2}\,dx^2,
\]
which proves (66.14).  Diffeomorphism invariance proves (66.15).
The stated pointwise convergence together with a constant nonzero
\(L^4\) norm rules out strong \(L^4\) convergence.
\end{proof}

This particular concentration lies on a diffeomorphism orbit, so a
correct quotient should identify it.  The example proves that the
gauge slice and its Jacobian are indispensable; it does not prove
physical inequivalence of the concentrating metrics.

\section{Two rigorous coercive completions}

\subsection{A background biharmonic completion}

Define on the fixed-volume slice
\[
 S_{\rm bi}(u)
 =
 S_{\rm EH}^{E}(u^2\bar g)
 +\alpha\|\bar\Delta u\|_2^2,\qquad \alpha>0.          \tag{66.16}
\]

\begin{theorem}[Fixed-volume biharmonic coercivity]
\label{thm:v15v66-biharmonic}
There is a constant
\(C=C(\bar g,V_0,\kappa,\Lambda,\alpha)\) such that
\[
 S_{\rm bi}(u)
 \geq
 \frac{\alpha}{2}\|\bar\Delta u\|_2^2-C               \tag{66.17}
\]
for every smooth positive \(u\) satisfying (66.6).
\end{theorem}

\begin{proof}
The spectral inequality
\[
 \|\bar\nabla u\|_2^2
 \leq
 \epsilon\|\bar\Delta u\|_2^2
 \frac1{4\epsilon}\|u\|_2^2                           \tag{66.18}
\]
holds for every \(\epsilon>0\), because
\(\lambda\leq\epsilon\lambda^2+(4\epsilon)^{-1}\)
on each eigenmode of \(-\bar\Delta\).  Choose
\(\epsilon=\alpha/(12\kappa)\).  Then
\[
 6\kappa\|\bar\nabla u\|_2^2
 \leq
 \frac{\alpha}{2}\|\bar\Delta u\|_2^2
 \frac{18\kappa^2}{\alpha}\|u\|_2^2.                  \tag{66.19}
\]
Also
\[
 \|u\|_2^2
 \leq\operatorname{Vol}_{\bar g}(M)^{1/2}V_0^{1/2}.   \tag{66.20}
\]
Bound the scalar-curvature term in (66.4) by
\(\kappa\|\bar R\|_\infty\|u\|_2^2\), and note that the
cosmological term is fixed.  Combining these estimates with
(66.16) yields (66.17).
\end{proof}

The completion (66.16) is background dependent.  It may be used as
a regulator or gauge-slice penalty only if split independence and
the associated ghost/Jacobian rows are later proved.

\subsection{A curvature-square completion on a bounded slice}

\begin{theorem}[\(R^2\) coercivity under conformal upper control]
\label{thm:v15v66-r2-bounded}
Fix \(M_0<\infty\) and restrict to positive \(u\) satisfying
\[
 \int_Mu^4dV_{\bar g}=V_0,\qquad u\leq M_0
 \quad\hbox{almost everywhere}.                       \tag{66.21}
\]
For every \(\alpha>0\), there are \(c_\alpha>0\) and
\(C_\alpha<\infty\) such that
\[
 S_{\rm EH}^{E}(u^2\bar g)+\alpha{\cal I}_2(u)
 \geq
 c_\alpha\|u\|_{H^2(\bar g)}^2-C_\alpha.              \tag{66.22}
\]
\end{theorem}

\begin{proof}
From (66.10) and \(u\leq M_0\),
\[
 {\cal I}_2(u)
 \geq M_0^{-2}\|L_{\bar g}u\|_2^2.                    \tag{66.23}
\]
Elliptic interpolation gives, for every \(\epsilon>0\),
\[
 \|\bar\nabla u\|_2^2
 \leq
 \epsilon\|L_{\bar g}u\|_2^2+C_{\epsilon,\bar g}\|u\|_2^2.
                                                               \tag{66.24}
\]
Choose \(\epsilon\) so that
\(6\kappa\epsilon\leq\alpha/(2M_0^2)\).
Equations (66.4), (66.20), (66.23), and (66.24) give
\[
 S_{\rm EH}^{E}+\alpha{\cal I}_2
 \geq
 \frac{\alpha}{2M_0^2}\|L_{\bar g}u\|_2^2-C.           \tag{66.25}
\]
Finally the elliptic estimate
\[
 \|u\|_{H^2}^2
 \leq C_{\bar g}\bigl(\|L_{\bar g}u\|_2^2+\|u\|_2^2\bigr)
\]
and (66.20) yield (66.22).
\end{proof}

The upper bound in (66.21) is a real hypothesis.  Fixed volume and
bounded \(R^2\) action do not supply it, as
Counterexample~\ref{cnt:v15v66-mobius} demonstrates before
quotient normalization.

\section{Connection to the exponential metric moment}

\begin{proposition}[Finite-cutoff Gibbs moment compiler]
\label{prop:v15v66-gibbs}
Let \(\mu_n\) be finite-cutoff reference measures and
\[
 d\nu_n=Z_n^{-1}e^{-S_n}\,d\mu_n.
\]
Suppose for some \(c>\eta>0\),
\[
 S_n\geq cV_n-C,                                      \tag{66.26}
\]
and
\[
 \sup_n
 \frac{e^C}{Z_n}
 \int e^{-(c-\eta)V_n}\,d\mu_n<\infty.                \tag{66.27}
\]
Then
\[
 \sup_n\mathbb E_{\nu_n}e^{\eta V_n}<\infty.           \tag{66.28}
\]
\end{proposition}

\begin{proof}
Directly,
\[
 \mathbb E_{\nu_n}e^{\eta V_n}
 =
 Z_n^{-1}\int e^{\eta V_n-S_n}\,d\mu_n
 \leq
 \frac{e^C}{Z_n}
 \int e^{-(c-\eta)V_n}\,d\mu_n.
\]
Use (66.27).
\end{proof}

The lower action bound alone is not a cofinal probability theorem:
the normalized Gaussian-volume ratio (66.27) must also be uniform
as \(d_n\to\infty\).

\section{Physical choice ledger}

There are three distinct responses to the conformal instability:
\begin{enumerate}
\item add a background or covariant higher-derivative term and
      prove its quotient, reflection, and renormalization
      properties;
\item rotate the conformal contour, producing a complex metric
      integration cycle whose phase and reflection properties must
      be reconstructed;
\item use a Lorentzian oscillatory measure, abandoning the positive
      Euclidean probability argument until a separate
      reconstruction theorem is supplied.
\end{enumerate}
Theorem~\ref{thm:v15v66-biharmonic} proves coercivity for the first
route with an explicit background penalty.
Theorem~\ref{thm:v15v66-r2-bounded} proves it for a covariant
\(R^2\) term only on a conformally bounded slice.  Neither theorem
selects the physical route.

\begin{assessment}[Physical status]
\label{ass:v15v66-status}
V66 proves the fixed-volume conformal instability of the bare
Euclidean Einstein action.  It gives the exact global critical
control supplied by \(R^2\), exhibits diffeomorphism-orbit
concentration, and proves two conditional coercive completions.

The full programme still lacks a natural gauge slice that removes
conformal-orbit drift, a covariant global replacement for the
upper bound (66.21), the uniform normalization estimate (66.27),
and compatibility with reflection positivity and the physical
Lorentzian theory.
\end{assessment}

\begin{decision}[V67 quotient-normalized conformal compactness]
\label{dec:v15v66-next}
Attack the gap between the critical estimate (66.12) and the
carrier required in V42.  Formulate a quotient normalization that
removes conformal automorphism drift, prove compactness under an
explicit no-concentration condition, and identify the concentration
measure that remains when the condition fails.  Test whether a
uniform \(L^\infty\) bound can be replaced by equiintegrability of
\(u^4\) plus an elliptic right-hand-side estimate.
\end{decision}

\begin{firewall}[Claim boundary after V66]
\label{fire:v15v66}
V66 proves conformal formulas, instability, critical estimates,
and conditional stabilizations.  It does not construct the
physical Euclidean gravity measure, choose a conformal contour,
prove cofinal normalization, or establish the full
Standard Model--Einstein continuum.
\end{firewall}

\section*{Version V66 append-only record}
\addcontentsline{toc}{section}{Version V66 append-only record}
The complete V65 source was retained before this continuation.  It
had \(641201\) bytes, \(18440\) lines, and SHA--256
\[
 \texttt{F887AE5BCB92B640676D492B23F67E4375FA60989E600611234080BA2DADA6DC}.
\]
V66 derived the four-dimensional conformal action, proved its
fixed-volume high-frequency instability, characterized the
critical control from \(R^2\), exhibited conformal-orbit
concentration, and proved background-biharmonic and bounded-slice
curvature-square coercivity theorems.

\chapter{Fifteenth-cycle continuation V67: quotient-normalized conformal compactness}
\label{sec:v15v67}

\section{Critical conformal data}

Retain the notation of V66 and consider a sequence \(u_j>0\) on
the fixed background \((M,\bar g)\) satisfying
\[
 \int_Mu_j^4\,dV_{\bar g}=V_0.                         \tag{67.1}
\]
Define
\[
 q_j=u_j^{-1}L_{\bar g}u_j.                           \tag{67.2}
\]
Then
\[
 \|q_j\|_2^2
 =
 \int_MR_{g_j}^2\,dV_{g_j},\qquad g_j=u_j^2\bar g.     \tag{67.3}
\]
The V66 estimate gives boundedness of \(u_j\) in
\(W^{2,4/3}\) whenever \(q_j\) is bounded in \(L^2\).

\begin{definition}[No concentration of conformal volume]
\label{def:v15v67-equi}
The conformal volume densities have no concentration when
\((u_j^4)_j\) is uniformly integrable with respect to
\(dV_{\bar g}\); equivalently, for every \(\epsilon>0\) there is
\(\delta>0\) such that
\[
 \operatorname{Vol}_{\bar g}(E)<\delta
 \quad\Longrightarrow\quad
 \sup_j\int_Eu_j^4\,dV_{\bar g}<\epsilon.              \tag{67.4}
\]
\end{definition}

\section{Critical compactness with an explicit defect}

\begin{theorem}[No-concentration compactness]
\label{thm:v15v67-critical-compactness}
Assume (67.1),
\[
 \sup_j\|q_j\|_2<\infty,                              \tag{67.5}
\]
and (67.4).  Then a subsequence and functions
\[
 u\in W^{2,4/3}(M),\qquad q\in L^2(M)                 \tag{67.6}
\]
exist such that
\[
\begin{aligned}
 u_j&\rightharpoonup u &&\hbox{in }W^{2,4/3},\\
 u_j&\to u &&\hbox{strongly in }L^4,\\
 q_j&\rightharpoonup q &&\hbox{in }L^2,               \tag{67.7}
\end{aligned}
\]
and
\[
 L_{\bar g}u=uq                                       \tag{67.8}
\]
in distributions.  In particular,
\[
 \int_Mu^4\,dV_{\bar g}=V_0.                          \tag{67.9}
\]
If \(u>0\) almost everywhere, then
\[
 \int_M|u^{-1}L_{\bar g}u|^2dV_{\bar g}
 \leq\liminf_{j\to\infty}
 \int_M|q_j|^2dV_{\bar g}.                            \tag{67.10}
\]
\end{theorem}

\begin{proof}
Theorem~\ref{thm:v15v66-r2-critical} bounds \(u_j\) in the
reflexive space \(W^{2,4/3}\).  After taking a subsequence,
\(u_j\rightharpoonup u\) there.  Rellich compactness gives
\[
 u_j\to u\quad\hbox{in }L^r\ (1\leq r<4)
 \quad\hbox{and almost everywhere}.                   \tag{67.11}
\]
Uniform integrability of \(u_j^4\), almost-everywhere convergence,
and Vitali's theorem imply
\[
 \|u_j^4-u^4\|_1\to0.
\]
The elementary inequality
\[
 |a-b|^4\leq 8(a^4+b^4),\qquad a,b\geq0,
\]
together with uniform integrability and convergence in measure,
also gives \(u_j\to u\) in \(L^4\).  This proves (67.9).

By (67.5), take a further subsequence with
\(q_j\rightharpoonup q\) in \(L^2\).  For every smooth test
function \(\zeta\), (67.11) gives
\(\zeta u_j\to\zeta u\) strongly in \(L^2\).  Hence
\[
 \int_M\zeta u_jq_j\,dV_{\bar g}
 \longrightarrow
 \int_M\zeta uq\,dV_{\bar g}.                         \tag{67.12}
\]
Since \(L_{\bar g}u_j=u_jq_j\), weak convergence on the left
gives (67.8).  If \(u>0\) almost everywhere, then
\(q=u^{-1}L_{\bar g}u\), and weak lower semicontinuity of the
\(L^2\) norm proves (67.10).
\end{proof}

Positivity of every \(u_j\) does not alone ensure positivity of
the limit.  A nondegenerate continuum metric still requires an
inverse-volume or noncollapse estimate.

\begin{proposition}[Concentration-defect measure]
\label{prop:v15v67-defect}
Assume only (67.1) and (67.5).  After a subsequence,
\[
 u_j^4dV_{\bar g}\stackrel{*}{\rightharpoonup}\mu     \tag{67.13}
\]
as finite measures, and
\[
 \mu=u^4dV_{\bar g}+\omega,\qquad \omega\geq0.         \tag{67.14}
\]
The defect has mass
\[
 \omega(M)=V_0-\int_Mu^4dV_{\bar g}.                  \tag{67.15}
\]
Condition (67.4) forces \(\omega=0\).  Without further
\(\epsilon\)-regularity or classification, \(\omega\) is not
proved to be atomic.
\end{proposition}

\begin{proof}
Fixed total mass gives weak-star precompactness of the measures in
(67.13).  The \(W^{2,4/3}\) bound again gives almost-everywhere
convergence of a subsequence.  For every nonnegative continuous
\(\zeta\), Fatou's lemma gives
\[
 \int_M\zeta u^4dV_{\bar g}
 \leq\liminf_j\int_M\zeta u_j^4dV_{\bar g}
 =
 \int_M\zeta\,d\mu.
\]
Thus the difference in (67.14) is a positive measure, and total
mass gives (67.15).  Under (67.4), Vitali convergence in the proof
of Theorem~\ref{thm:v15v67-critical-compactness} gives
\(\mu=u^4dV_{\bar g}\).
\end{proof}

\section{What quotient normalization removes}

On \(S^4\subset{\mathbb R}^5\), associate to \(u\) the probability
measure
\[
 d\widehat\mu_u
 =
 V_0^{-1}u^4dV_{g_0}.                                 \tag{67.16}
\]
A standard conformal-orbit normalization is the barycenter
condition
\[
 \int_{S^4}x\,d\widehat\mu_u(x)=0\in{\mathbb R}^5.     \tag{67.17}
\]

\begin{proposition}[Barycenter normalization excludes one bubble]
\label{prop:v15v67-barycenter}
If every \(u_j\) satisfies (67.17), then
\(\widehat\mu_{u_j}\) cannot converge weakly to a single point mass
\(\delta_p\).
\end{proposition}

\begin{proof}
Weak convergence to \(\delta_p\) would imply
\[
 0
 =
 \lim_j\int_{S^4}x\,d\widehat\mu_{u_j}(x)
 =
 \int_{S^4}x\,d\delta_p(x)=p,
\]
which is impossible because \(|p|=1\).
\end{proof}

The condition does not exclude balanced multi-point
concentration: for instance
\(\tfrac12(\delta_p+\delta_{-p})\) has zero barycenter.  Thus a
quotient normalization removes the one-parameter M\"obius drift
exhibited in V66 but does not replace the no-concentration estimate
(67.4).

\section{Equiintegrability is below the \(H^2\) threshold}

\begin{counterexample}[No \(H^2\) upgrade from critical compactness]
\label{cnt:v15v67-h2-gap}
In a coordinate ball in a four-manifold, fix
\(0<\alpha<1\), a nonnegative
\(\chi\in C_c^\infty(B_1)\), and \(\epsilon_j\downarrow0\).  Put
\[
 v_j(x)=\epsilon_j^{-\alpha}
 \chi(x/\epsilon_j)                                   \tag{67.18}
\]
in that chart and extend by zero.  After adding \(1\) and
rescaling by constants tending to a nonzero limit, obtain positive
\(u_j\) satisfying (67.1).  Then
\[
 \sup_j\|u_j\|_{W^{2,4/3}}<\infty,                    \tag{67.19}
\]
\((u_j^4)\) is uniformly integrable, and \(u_j\) converges strongly
in \(L^4\), but
\[
 \|\nabla^2u_j\|_2\longrightarrow\infty.              \tag{67.20}
\]
\end{counterexample}

\begin{proof}
Four-dimensional scaling gives
\[
\begin{aligned}
 \|v_j\|_4&=O(\epsilon_j^{1-\alpha}),\\
 \|\nabla v_j\|_{4/3}&=O(\epsilon_j^{2-\alpha}),\\
 \|\nabla^2v_j\|_{4/3}&=O(\epsilon_j^{1-\alpha}),\\
 \|\nabla^2v_j\|_2&=
 \epsilon_j^{-\alpha}\|\nabla^2\chi\|_2.              \tag{67.21}
\end{aligned}
\]
The first three lines prove (67.19) and strong \(L^4\) convergence
after the harmless volume normalization.  Moreover, for a level
\(R\), only indices with
\(\epsilon_j^{-\alpha}\gtrsim R\) contribute to the high-amplitude
tail, whose fourth-power mass is
\(O(\epsilon_j^{4(1-\alpha)})\).  Its supremum tends to zero as
\(R\to\infty\), proving uniform integrability.  The last line in
(67.21) proves (67.20).
\end{proof}

This example is a function-space separation.  It does not assert
that the displayed \(u_j\) have bounded curvature-square energy.
It proves that the conclusions of the critical compactness theorem
alone cannot populate the V42 metric carrier.

\section{A higher-integrability bootstrap}

The missing upgrade has a clean elliptic form: improve \(q_j\) by
any fixed amount beyond \(L^2\).

\begin{theorem}[Strict curvature-factor gain gives the V42 carrier]
\label{thm:v15v67-bootstrap}
Let \(p>2\).  Suppose (67.1) holds and
\[
 \sup_j\|q_j\|_{L^p(\bar g)}\leq Q_p.                 \tag{67.22}
\]
Then
\[
 \sup_j\left(
 \|u_j\|_{L^\infty}
 +\|u_j\|_{W^{2,p}}
 \right)
 \leq C(\bar g,V_0,p,Q_p).                            \tag{67.23}
\]
In particular \(p=2+\delta\) supplies the
\(W^{2,2+\delta}\) conformal carrier required in V42.
\end{theorem}

\begin{proof}
Start with the fixed bound \(\|u_j\|_4=V_0^{1/4}\).
If \(u_j\) is bounded in \(L^{r_k}\), define \(s_k\) by
\[
 \frac1{s_k}=\frac1p+\frac1{r_k}.                     \tag{67.24}
\]
H\"older and the elliptic estimate give
\[
 \|u_j\|_{W^{2,s_k}}
 \leq C\bigl(Q_p\|u_j\|_{r_k}+\|u_j\|_{s_k}\bigr).
                                                               \tag{67.25}
\]
While \(s_k<2\), Sobolev embedding gives a new exponent \(r_{k+1}\)
with
\[
 \frac1{r_{k+1}}
 =
 \frac1{s_k}-\frac12
 =
 \frac1{r_k}-\left(\frac12-\frac1p\right).             \tag{67.26}
\]
Because \(p>2\), the parenthesized quantity is positive.  Starting
from \(r_0=4\), finitely many iterations make
\(1/p+1/r_k<1/2\).  If equality occurs at one step, the critical
Sobolev embedding gives every finite exponent, so choose a slightly
larger \(r_k\) and continue.  Thus eventually \(s_k>2\), and
\[
 W^{2,s_k}(M)\hookrightarrow L^\infty(M).
\]
This proves the first bound in (67.23).  Now
\(\|u_jq_j\|_p\leq\|u_j\|_\infty Q_p\), and the elliptic estimate
applied to
\[
 L_{\bar g}u_j=u_jq_j
\]
gives the uniform \(W^{2,p}\) bound.
\end{proof}

\begin{remark}[Covariant meaning of the stronger row]
\label{rem:v15v67-weight}
Since
\[
 q_j=R_{g_j}u_j^2,\qquad dV_{\bar g}=u_j^{-4}dV_{g_j},
\]
condition (67.22) is
\[
 \int_M|R_{g_j}|^p u_j^{\,2p-4}\,dV_{g_j}\leq Q_p^p.  \tag{67.27}
\]
For \(p=2\) the conformal weight disappears and one recovers the
covariant \(R^2\) action.  For \(p>2\) the weight is real; an
unweighted \(L^p(g_j)\) curvature bound does not automatically
imply (67.22) without additional control of \(u_j\).
\end{remark}

\section{Updated conformal closure ledger}

The conformal sector now has two rigorous routes:
\begin{enumerate}
\item critical route: \(R^2\) bound plus uniform integrability of
      \(u_j^4\) gives strong volume convergence and the weak
      equation (67.8), but only at \(W^{2,4/3}\);
\item supercritical route: the weighted gain (67.22) gives
      \(W^{2,2+\delta}\), \(L^\infty\), and compactness below the
      corresponding H\"older exponent.
\end{enumerate}
Both routes still require noncollapse of \(u_j^{-1}\) to produce a
nondegenerate limiting metric.

\begin{assessment}[Physical status]
\label{ass:v15v67-status}
V67 proves an exact no-concentration compactness theorem, identifies
the conformal-volume defect measure, shows the limited power of
barycenter quotient normalization, and gives a finite elliptic
bootstrap from \(L^{2+\delta}\) control of the conformal curvature
factor to the V42 carrier.

The programme does not yet derive uniform integrability of \(u_j^4\),
the stronger weighted estimate (67.22), or inverse conformal
noncollapse from the physical action.  The critical \(R^2\) term
alone remains insufficient.
\end{assessment}

\begin{decision}[V68 local epsilon regularity]
\label{dec:v15v67-next}
Seek a local, scale-invariant condition that implies both
equiintegrability and the higher-integrability row without assuming
a global \(L^\infty\) bound.  Formulate a Morrey-smallness
\(\epsilon\)-regularity theorem for
\(L_{\bar g}u=qu\), prove the local reverse-H\"older step under an
explicit small \(L^2\) norm of \(q\), and retain concentration
points when that threshold fails.
\end{decision}

\begin{firewall}[Claim boundary after V67]
\label{fire:v15v67}
V67 closes conditional conformal compactness and elliptic
bootstrap implications.  It does not prove the needed
no-concentration, higher-integrability, noncollapse, quotient
Jacobian, or continuum probability estimates for the physical
theory.
\end{firewall}

\section*{Version V67 append-only record}
\addcontentsline{toc}{section}{Version V67 append-only record}
The complete V66 source was retained before this continuation.  It
had \(653736\) bytes, \(18868\) lines, and SHA--256
\[
 \texttt{5075C2C3398526FBF9D1835A2157AE60247EA8288BC101BDB3DCE55E07141F2F}.
\]
V67 proved critical conformal compactness under volume
equiintegrability, identified the defect measure, separated
quotient drift from multi-point concentration, and proved that any
fixed \(L^{2+\delta}\) gain for the conformal curvature factor
bootstraps to a uniform \(W^{2,2+\delta}\) carrier.

\chapter{Fifteenth-cycle continuation V68: critical epsilon regularity and conformal concentration}
\label{sec:v15v68}

\section{Local equation and scale}

The conformal equation from V67 is
\[
 -6\bar\Delta u=(q-\bar R)u.                           \tag{68.1}
\]
In four dimensions the \(L^2\) norm of \(q\) is invariant under
the natural rescaling of (68.1).  This is the critical local
quantity.

\begin{lemma}[Critical local \(L^4\)-to-\(L^6\) gain]
\label{lem:v15v68-local-gain}
There are constants \(\epsilon_*>0\), \(r_*>0\), and \(C<\infty\),
depending only on a fixed finite background atlas and
\(\|\bar R\|_\infty\), with the following property.  Let
\(B_{2r}\) be a coordinate ball with \(0<r\leq r_*\).  If
\(u\geq0\), \(u\in H^1(B_{2r})\cap L^4(B_{2r})\), is a weak
solution of (68.1), \(q\in L^2(B_{2r})\), and
\[
 \int_{B_{2r}}|q|^2dV_{\bar g}\leq\epsilon_*,          \tag{68.2}
\]
then
\[
 \left(\fint_{B_r}u^6dV_{\bar g}\right)^{1/6}
 \leq
 C\left(\fint_{B_{2r}}u^4dV_{\bar g}\right)^{1/4}.
                                                               \tag{68.3}
\]
\end{lemma}

\begin{proof}
Choose a cutoff \(\eta\) supported in \(B_{2r}\), equal to one on
\(B_r\), and satisfying \(|\bar\nabla\eta|\leq C/r\).
Test (68.1), first with bounded truncations and then by monotone
convergence, against \(\eta^2u^2\).  Since
\[
 |\bar\nabla(u^{3/2})|^2=\frac94u|\bar\nabla u|^2,
\]
Cauchy's inequality on the cutoff cross term gives
\[
 \int|\bar\nabla(\eta u^{3/2})|^2
 \leq
 C\int
 \bigl(r^{-2}+|\bar R|\bigr)u^3
 C\int |q|\,\eta^2u^3.                                \tag{68.4}
\]
The four-dimensional local Sobolev inequality gives
\[
 \|\eta u^{3/2}\|_4^2
 \leq C_S\left[
 \int|\bar\nabla(\eta u^{3/2})|^2
 r^{-2}\int\eta^2u^3
 \right].                                             \tag{68.5}
\]
By H\"older,
\[
 \int |q|\,\eta^2u^3
 \leq
 \|q\|_{L^2(B_{2r})}
 \|\eta u^{3/2}\|_4^2.                                \tag{68.6}
\]
Choose \(\epsilon_*\) so that the last term is absorbed into the
left side of (68.5).  Shrink \(r_*\), if necessary, so the
\(\bar R\) term is controlled by the scale term.  We obtain
\[
 \|u\|_{L^6(B_r)}^3
 \leq Cr^{-2}\|u\|_{L^3(B_{2r})}^3.                   \tag{68.7}
\]
H\"older bounds the \(L^3\) norm by the \(L^4\) norm and the volume
of \(B_{2r}\).  Uniform comparability of coordinate-ball volume
with \(r^4\) converts (68.7) to the averaged estimate (68.3).
\end{proof}

The smallness threshold concerns the local norm of \(q\), not its
pointwise size.  The lower-order background curvature becomes
small after restricting to a sufficiently small four-dimensional
ball.

\section{Uniform curvature equiintegrability}

\begin{definition}[No concentration of curvature-square energy]
\label{def:v15v68-q-equi}
A family \(q_j\in L^2(M)\) has equiintegrable critical energy when
\[
 \lim_{\delta\downarrow0}
 \sup_j\sup_{\operatorname{Vol}_{\bar g}(E)\leq\delta}
 \int_E|q_j|^2dV_{\bar g}=0.                           \tag{68.8}
\]
\end{definition}

\begin{theorem}[Curvature equiintegrability implies conformal compactness]
\label{thm:v15v68-global}
Let \(u_j>0\) satisfy
\[
 L_{\bar g}u_j=q_ju_j,\qquad
 \int_Mu_j^4dV_{\bar g}=V_0,                           \tag{68.9}
\]
and suppose
\[
 \sup_j\|q_j\|_2<\infty                                \tag{68.10}
\]
together with (68.8).  Then
\[
 \sup_j\|u_j\|_{L^6}<\infty,\qquad
 \sup_j\|u_j\|_{W^{2,3/2}}<\infty.                    \tag{68.11}
\]
After a subsequence,
\[
 u_j\longrightarrow u\quad\hbox{strongly in }L^4.     \tag{68.12}
\]
Thus the no-concentration hypothesis (67.4) follows from
curvature-energy equiintegrability.
\end{theorem}

\begin{proof}
By (68.8), choose a common small radius so every doubled ball in a
finite background cover satisfies (68.2).  Lemma
\ref{lem:v15v68-local-gain}, (68.9), and the finite cover give the
uniform \(L^6\) bound.

H\"older now yields
\[
 \|q_ju_j\|_{3/2}
 \leq\|q_j\|_2\|u_j\|_6.                              \tag{68.13}
\]
The global elliptic estimate for
\(L_{\bar g}u_j=q_ju_j\), together with the fixed \(L^4\) norm,
gives the \(W^{2,3/2}\) bound in (68.11).  In four dimensions,
\[
 W^{2,3/2}(M)\Subset L^4(M),                           \tag{68.14}
\]
because the critical target of \(W^{2,3/2}\) is \(L^6\), and the
embedding is compact below exponent six.  This proves (68.12).
Strong \(L^4\) convergence makes \(u_j^4\) uniformly integrable.
\end{proof}

This result gains compactness but still stops below \(W^{2,2}\).
The \(L^{2+\delta}\) row of V67 remains necessary for the stronger
metric carrier.

\section{Finite concentration set}

Assume only (68.9)--(68.10).  Pass to a subsequence such that
\[
 |q_j|^2dV_{\bar g}\stackrel{*}{\rightharpoonup}\lambda
                                                               \tag{68.15}
\]
as finite measures and define
\[
 \Sigma
 =
 \{x\in M:\lambda(\{x\})\geq\epsilon_*/4\}.            \tag{68.16}
\]

\begin{proposition}[Energy quantization at the compactness scale]
\label{prop:v15v68-finite-set}
The set \(\Sigma\) is finite and
\[
 \#\Sigma
 \leq
 \frac{4\lambda(M)}{\epsilon_*}
 \leq
 \frac{4\liminf_j\|q_j\|_2^2}{\epsilon_*}.             \tag{68.17}
\]
For every compact \(K\Subset M\setminus\Sigma\), the sequence is
bounded in \(W^{2,3/2}(K)\) and is precompact in \(L^4(K)\).
\end{proposition}

\begin{proof}
Distinct atoms have disjoint mass, so summing
\(\lambda(\{x\})\geq\epsilon_*/4\) proves (68.17).
If \(x\notin\Sigma\), continuity from above gives a ball
\(B_{2r_x}(x)\) whose closure has
\(\lambda\)-mass less than \(\epsilon_*/2\); choose its radius so
the boundary has zero \(\lambda\)-mass.  Weak convergence of the
energy measures then gives, for all sufficiently large \(j\),
\[
 \int_{B_{2r_x}(x)}|q_j|^2dV_{\bar g}<\epsilon_*.
\]
A finite subcover of \(K\), Lemma
\ref{lem:v15v68-local-gain}, and the local elliptic estimate give
the asserted bounds.  Rellich compactness gives local strong
\(L^4\) convergence.
\end{proof}

\begin{theorem}[The conformal-volume defect is atomic]
\label{thm:v15v68-atomic-defect}
Let the conformal-volume defect \(\omega\) be the measure in
(67.14) for the subsequence used in (68.15).  Then
\[
 \operatorname{supp}\omega\subseteq\Sigma,\qquad
 \omega=\sum_{x\in\Sigma}m_x\delta_x,\quad m_x\geq0.   \tag{68.18}
\]
\end{theorem}

\begin{proof}
Proposition~\ref{prop:v15v68-finite-set} gives strong \(L^4\)
convergence on every compact subset of \(M\setminus\Sigma\).
Hence \(u_j^4dV_{\bar g}\to u^4dV_{\bar g}\) in \(L^1\) there, so
the defect measure vanishes off \(\Sigma\).  A positive measure
supported on a finite set has the form displayed in (68.18).
\end{proof}

This is a genuine improvement over the general defect statement
of V67: the conformal-volume loss is now localized at finitely many
points selected by atoms of the curvature-square energy.

\section{Small-energy regions and quotient normalization}

\begin{corollary}[No atoms, no conformal-volume defect]
\label{cor:v15v68-no-atoms}
If the weak limit \(\lambda\) in (68.15) has no atom of mass
\(\epsilon_*/4\) or larger, then
\[
 u_j\to u\quad\hbox{strongly in }L^4(M)                \tag{68.19}
\]
and \(\omega=0\).
\end{corollary}

\begin{proof}
Then \(\Sigma=\varnothing\), so Proposition
\ref{prop:v15v68-finite-set} and Theorem
\ref{thm:v15v68-atomic-defect} apply globally.
\end{proof}

Barycenter normalization from V67 can exclude a single volume
atom carrying all mass, but it does not exclude several atoms.
The curvature-energy measure (68.15) supplies the correct local
bookkeeping for those remaining configurations.

\section{Why bounded total \(R^2\) energy is insufficient}

\begin{counterexample}[Moving critical energy]
\label{cnt:v15v68-moving-energy}
Let \(x_j\) range over distinct points and let \(q_j\) be supported
in a ball of radius \(r_j\downarrow0\), scaled so that
\[
 \|q_j\|_2=1.                                         \tag{68.20}
\]
Then the total critical energy is uniformly bounded, but (68.8)
fails because
\[
 \int_{\operatorname{supp}q_j}|q_j|^2=1
\]
on sets whose volumes tend to zero.  Thus a global \(R^2\) bound
cannot replace local energy control.
\end{counterexample}

The counterexample concerns the potential data.  Solving (68.1)
may impose further structure, but that structure has to be proved;
it cannot be inferred from the norm bound alone.

\section{Interface with the metric carrier}

Away from \(\Sigma\), the conformal factors have
\[
 u_j\ \hbox{bounded in }W^{2,3/2}_{\rm loc}
 \quad\hbox{and strongly compact in }L^4_{\rm loc}.    \tag{68.21}
\]
If, on the same region, \(q_j\) is bounded in
\(L^{2+\delta}\), Theorem~\ref{thm:v15v67-bootstrap}
upgrades (68.21) to
\[
 u_j\ \hbox{bounded in }W^{2,2+\delta}_{\rm loc}
 \cap L^\infty_{\rm loc}.                             \tag{68.22}
\]
The remaining tasks are therefore separated into:
\begin{enumerate}
\item exclude or resolve the finitely many atoms in \(\Sigma\);
\item obtain a strict higher-integrability gain away from them;
\item prove a lower bound for \(u_j\) so the limiting metric does
      not degenerate;
\item transport all estimates through the metric quotient and its
      Jacobian.
\end{enumerate}

\begin{assessment}[Physical status]
\label{ass:v15v68-status}
V68 proves a critical epsilon-regularity lemma, derives global
conformal compactness from equiintegrability of curvature-square
energy, and shows that under a mere total bound any conformal-volume
defect is supported on finitely many curvature-energy atoms.

The physical action has not yet been shown to prevent those atoms,
to improve \(q\) beyond \(L^2\), or to keep the conformal factor
away from zero.  Hence the full metric carrier and continuum
measure remain open.
\end{assessment}

\begin{decision}[V69 conformal noncollapse]
\label{dec:v15v68-next}
Attack the inverse-conformal-factor row.  Prove a quantitative
Harnack and lower-bound compiler for positive solutions of
\(L_{\bar g}u=qu\) when \(q\in L^{2+\delta}\), using the fixed
\(L^4\) normalization to turn a local supremum--infimum estimate
into a global positive lower bound.  Exhibit a critical
\(L^2\) sequence showing why the strict exponent gain is needed.
\end{decision}

\begin{firewall}[Claim boundary after V68]
\label{fire:v15v68}
V68 proves conditional epsilon regularity and concentration
localization.  It does not exclude curvature atoms, establish
noncollapse, derive higher integrability from the Standard Model
action, or construct the physical continuum measure.
\end{firewall}

\section*{Version V68 append-only record}
\addcontentsline{toc}{section}{Version V68 append-only record}
The complete V67 source was retained before this continuation.  It
had \(665768\) bytes, \(19244\) lines, and SHA--256
\[
 \texttt{2C6FF08A1FE060F9D59DD82DBC4482A0CDE3561FE306560545972709021E9DC1}.
\]
V68 proved critical local \(L^4\)-to-\(L^6\) regularity, obtained
global \(W^{2,3/2}\) compactness from curvature-energy
equiintegrability, and confined every remaining conformal-volume
defect to finitely many atoms of the curvature-square energy.

\chapter{Fifteenth-cycle continuation V69: Harnack control and conformal noncollapse}
\label{sec:v15v69}

\section{The supercritical potential equation}

Let \(M\) be closed and connected and let
\[
 L_{\bar g}u=qu,\qquad u>0.                            \tag{69.1}
\]
Equivalently,
\[
 -\bar\Delta u=Vu,\qquad
 V=\frac{q-\bar R}{6}.                                \tag{69.2}
\]
In four dimensions the strict condition
\[
 V\in L^p,\qquad p>2,                                 \tag{69.3}
\]
lies above the scale-invariant Schr\"odinger threshold.

\section{A local Harnack compiler}

\begin{lemma}[Local Harnack inequality]
\label{lem:v15v69-local-harnack}
Fix \(p>2\), \(K<\infty\), and a background coordinate radius
\(r_0>0\) below the injectivity radius.  There is
\[
 C_H=C_H(\bar g,p,K,r_0)<\infty                       \tag{69.4}
\]
such that every positive weak solution of (69.2) on \(B_{2r_0}\)
with
\[
 \|V\|_{L^p(B_{2r_0})}\leq K                          \tag{69.5}
\]
satisfies
\[
 \sup_{B_{r_0/2}}u
 \leq C_H\inf_{B_{r_0/2}}u.                           \tag{69.6}
\]
\end{lemma}

\begin{proof}
We record the three estimates in the Moser argument.

First, if \(w\geq0\) obeys
\[
 -\bar\Delta w\leq |V|w,                              \tag{69.7}
\]
testing with \(\eta^2w^{m-1}\), \(m\geq2\), gives
\[
 \|\eta w^{m/2}\|_4^2
 \leq
 C m^2\int
 \bigl(|\bar\nabla\eta|^2+r_0^{-2}\eta^2\bigr)w^m
 C m\int |V|\eta^2w^m.                                \tag{69.8}
\]
Because \(p>2\),
\[
 \frac{2p}{p-1}<4.
\]
Interpolate the \(L^{2p/(p-1)}\) norm of
\(\eta w^{m/2}\) between \(L^2\) and \(L^4\), use
H\"older in the last term of (69.8), and then Young's inequality.
The \(L^4\) term is absorbed, with a remaining coefficient
depending only on \(\bar g,p,K,r_0,m\).  Iteration on concentric
balls and exponents \(m,2m,4m,\ldots\) yields, for every fixed
\(s>0\),
\[
 \sup_{B_{r_0}}w
 \leq
 C_s\left(\fint_{B_{3r_0/2}}w^s\right)^{1/s}.          \tag{69.9}
\]
The usual truncation \(w+\epsilon\) justifies negative or small
initial exponents before \(\epsilon\downarrow0\).

Second, test (69.2) with \(\eta^2/u\).  Cauchy's inequality gives
\[
 \int\eta^2|\bar\nabla\log u|^2
 \leq
 C\int|\bar\nabla\eta|^2+
 C\int\eta^2|V|.                                      \tag{69.10}
\]
H\"older and (69.5) control the final term.  Poincar\'e's
inequality on all subballs therefore gives a uniform BMO bound for
\(\log u\).  The John--Nirenberg argument, obtained by applying
Chebyshev's inequality to the dyadic Calder\'on--Zygmund
decomposition of its oscillation, supplies an \(s_0>0\) and
\(C_0<\infty\) such that
\[
 \left(\fint_{B_{r_0}}u^{s_0}\right)
 \left(\fint_{B_{r_0}}u^{-s_0}\right)\leq C_0.         \tag{69.11}
\]
All constants depend only on the data in (69.4).

Third, \(v=u^{-1}\) satisfies
\[
 -\bar\Delta v
 =
 -Vv-2u^{-3}|\bar\nabla u|^2
 \leq |V|v.                                           \tag{69.12}
\]
Apply (69.9) to \(u\) and \(u^{-1}\), with exponent \(s_0\), on
nested balls.  Multiplying the two estimates and using (69.11)
gives
\[
 \bigl(\sup_{B_{r_0/2}}u\bigr)
 \bigl(\sup_{B_{r_0/2}}u^{-1}\bigr)\leq C_H,
\]
which is (69.6).
\end{proof}

\section{Global fixed-volume noncollapse}

\begin{theorem}[Quantitative conformal noncollapse]
\label{thm:v15v69-global-harnack}
Let \(p>2\).  Suppose \(u>0\) solves (69.1) and
\[
 \|q\|_{L^p(\bar g)}\leq Q,\qquad
 \int_Mu^4dV_{\bar g}=V_0.                            \tag{69.13}
\]
There is a constant
\[
 H=H(M,\bar g,p,Q)<\infty                             \tag{69.14}
\]
such that
\[
 \sup_Mu\leq H\inf_Mu.                                \tag{69.15}
\]
Consequently
\[
 \frac{V_0^{1/4}}
 {H\,\operatorname{Vol}_{\bar g}(M)^{1/4}}
 \leq u(x)\leq
 \frac{H\,V_0^{1/4}}
 {\operatorname{Vol}_{\bar g}(M)^{1/4}}
 \quad\hbox{for every }x\in M.                        \tag{69.16}
\]
\end{theorem}

\begin{proof}
The \(L^p\) norm of \(V=(q-\bar R)/6\) is bounded in terms of the
data.  Cover the connected compact manifold by finitely many balls
to which Lemma~\ref{lem:v15v69-local-harnack} applies, with
successive balls overlapping.  Iterating the local inequality
along a chain between a maximum and a minimum proves (69.15).

Let \(m=\inf_Mu\).  Since \(u\leq Hm\),
\[
 V_0=\int_Mu^4dV_{\bar g}
 \leq H^4m^4\operatorname{Vol}_{\bar g}(M),
\]
which gives the lower bound in (69.16).  Conversely
\[
 m^4\operatorname{Vol}_{\bar g}(M)\leq V_0,
\]
and \(\sup u\leq Hm\) gives the upper bound.
\end{proof}

\begin{corollary}[Uniform metric equivalence]
\label{cor:v15v69-metric}
Under the hypotheses of
Theorem~\ref{thm:v15v69-global-harnack}, the metric
\(g=u^2\bar g\) satisfies
\[
 c^2\bar g\leq g\leq C^2\bar g                       \tag{69.17}
\]
with positive constants depending only on the data in (69.13).
The volume densities and inverse metrics are uniformly comparable.
\end{corollary}

\begin{proof}
Square the two pointwise bounds in (69.16).  The inverse comparison
follows by reversing the inequalities.
\end{proof}

\begin{corollary}[Compact supercritical conformal carrier]
\label{cor:v15v69-carrier}
For a sequence with common bounds (69.13),
\[
 \sup_j\|u_j\|_{W^{2,p}}<\infty                       \tag{69.18}
\]
and, after a subsequence,
\[
 u_j\to u\quad\hbox{in }C^{0,\beta}(M)                \tag{69.19}
\]
for every
\[
 0<\beta<\min\{1,\,2-4/p\}.
\]
The limit is bounded away from zero and defines a nondegenerate
conformal metric.
\end{corollary}

\begin{proof}
Theorem~\ref{thm:v15v67-bootstrap} proves (69.18).
Morrey--Rellich compactness gives (69.19), and (69.16) passes its
positive lower bound to the uniform limit.
\end{proof}

\section{Sharp failure at the critical exponent}

\begin{counterexample}[\(L^2\) potential does not give Harnack control]
\label{cnt:v15v69-critical}
For the round sphere and the conformal factors \(u_\lambda\) in
(66.14), set
\[
 g_\lambda=u_\lambda^2g_0=D_\lambda^*g_0.
\]
Then
\[
 L_{g_0}u_\lambda=q_\lambda u_\lambda,\qquad
 q_\lambda=R_{g_0}u_\lambda^2,                        \tag{69.20}
\]
and
\[
 \|q_\lambda\|_2^2
 =
 R_{g_0}^2\int_{S^4}u_\lambda^4dV_{g_0}
\]
is independent of \(\lambda\).  Nevertheless,
\[
 \frac{\sup u_\lambda}{\inf u_\lambda}
 =\lambda^2\longrightarrow\infty.                    \tag{69.21}
\]
\end{counterexample}

\begin{proof}
The conformal curvature equation
\(L_{g_0}u_\lambda=R_{g_\lambda}u_\lambda^3\)
and diffeomorphism invariance give
\(R_{g_\lambda}=R_{g_0}\), proving (69.20).
The \(L^2\) identity follows from the fixed volume in (66.15).
Formula (66.14) has value \(\lambda\) at the stereographic origin
and extends to value \(\lambda^{-1}\) at the opposite pole, proving
(69.21).
\end{proof}

The example is a conformal-diffeomorphism orbit, but that is enough
to disprove any pre-quotient Harnack constant depending only on the
critical \(L^2\) norm and volume.  A quotient normalization or a
strict exponent gain is necessary.

\section{Probabilistic noncollapse row}

Let \(u_n,q_n\) be random finite-cutoff conformal data on the
fixed-volume slice.  For \(Q>0\), define the good event
\[
 {\cal G}_{n,Q}=\{\|q_n\|_p\leq Q\}.                  \tag{69.22}
\]

\begin{proposition}[High-probability metric equivalence]
\label{prop:v15v69-probability}
On \({\cal G}_{n,Q}\), the deterministic bounds (69.16)--(69.17)
hold with constants independent of \(n\).  If, for some \(r>0\),
\[
 \sup_n\mathbb E_{\nu_n}\|q_n\|_p^r\leq C_r,          \tag{69.23}
\]
then
\[
 \sup_n\nu_n({\cal G}_{n,Q}^{\,c})
 \leq C_rQ^{-r}.                                      \tag{69.24}
\]
\end{proposition}

\begin{proof}
The first assertion is
Theorem~\ref{thm:v15v69-global-harnack}.  The second is Markov's
inequality.
\end{proof}

To obtain exponential inverse-volume moments, one must combine a
quantitative growth bound for \(H(Q)\) with a correspondingly
strong tail for \(\|q_n\|_p\).  A polynomial moment such as
(69.23) gives high-probability noncollapse but does not by itself
give the exponential stress moment required in V63.

\section{Weight comparison after noncollapse}

\begin{proposition}[Covariant and background curvature norms]
\label{prop:v15v69-weight}
Under (69.16), for every \(s\geq1\),
\[
 \|q\|_{L^s(\bar g)}^s
 =
 \int_M|R_g|^su^{2s-4}\,dV_g                         \tag{69.25}
\]
is comparable, with constants determined by (69.16), to
\[
 \int_M|R_g|^s\,dV_g.                                 \tag{69.26}
\]
\end{proposition}

\begin{proof}
The identity is (67.27).  The upper and lower pointwise bounds for
\(u^{2s-4}\), with the inequalities reversed when \(2s-4<0\),
give two-sided comparison.
\end{proof}

Thus once noncollapse has been established, an unweighted
supercritical scalar-curvature estimate can replace the weighted
condition in V67.  It cannot be used to prove the same noncollapse
without an independent noncircular argument.

\begin{assessment}[Physical status]
\label{ass:v15v69-status}
V69 proves that a strict \(L^p\), \(p>2\), conformal-potential bound
and fixed volume yield global Harnack control, metric
nondegeneracy, and compactness of the supercritical conformal
carrier.  It also gives an exact critical counterexample and a
probabilistic good-event formulation.

The physical Standard Model--Einstein measure has not yet been
shown to provide the supercritical \(q\) moment or the tail needed
for exponential inverse-volume control.  The result remains a
conditional noncollapse compiler.
\end{assessment}

\begin{decision}[V70 cross-program audit]
\label{dec:v15v69-next}
At the scheduled multiple of five, inspect the newest active
Yang--Mills and Navier--Stokes notes.  Test specifically for a new
reverse-H\"older, local-potential, quotient-normalization, or
cofinal-tail estimate that could populate V68--V69.  Record the
comparison, then continue in V71 with the trace Einstein equation
as a possible source of the required \(L^{2+\delta}\) curvature
gain.
\end{decision}

\begin{firewall}[Claim boundary after V69]
\label{fire:v15v69}
V69 proves deterministic and good-event noncollapse under a strict
supercritical potential bound.  It does not derive that bound from
the physical action, remove quotient drift by itself, or construct
the full interacting continuum measure.
\end{firewall}

\section*{Version V69 append-only record}
\addcontentsline{toc}{section}{Version V69 append-only record}
The complete V68 source was retained before this continuation.  It
had \(676823\) bytes, \(19566\) lines, and SHA--256
\[
 \texttt{5021D6AA1F2AD5D8EE905407A3E552E4C2C543587B86B679C803DA389502135E}.
\]
V69 proved local and global supercritical Harnack estimates,
converted fixed volume into uniform upper and lower conformal
bounds, established the compact nondegenerate carrier, and showed
by the round-sphere M\"obius family that critical \(L^2\) control
cannot suffice.

\chapter{Fifteenth-cycle continuation V70: scheduled reverse-H\"older dependency audit}
\label{sec:v15v70}

\section{Audited sources}

The newest active companion notes were hashed and their compactness,
Morrey, quotient, and cofinal claim boundaries were inspected.  The
result is
\[
\begin{array}{c|c|c|c|l}
\hbox{programme}&\hbox{version}&\hbox{bytes}&\hbox{lines}&\hbox{SHA--256}\\ \hline
\hbox{Yang--Mills}&V39&329898&9076&
\texttt{9B44B0FBCA6037F3319E86036753B0F3B76BDD06F6045F00747E05DC01E23F9D}\\
\hbox{Navier--Stokes}&V26&278283&5319&
\texttt{82C887F8C2C69E4F28FAD7C3DC8DE5B9099CB5A4BF431EBA1FFF3E91935978AD}.
\end{array}                                             \tag{70.1}
\]
Both hashes equal those recorded in V65.  There is no new companion
appendix to import.

\section{Targeted Yang--Mills check}

The Yang--Mills terminal theorem remains the finite
Wilson-pencil lower bound on
\([0,\pi/4]^2\times[0,t]\).  Its own firewall excludes a full
transverse-torus cover, a cofinal response sign, an interacting
measure, and reflection-positive regulator removal.

\begin{proposition}[No supercritical conformal input from YM V39]
\label{prop:v15v70-ym}
YM V39 supplies none of
\[
 q\in L^{2+\delta},\qquad
 \sup u/\inf u<\infty,\qquad
 \sup_n\mathbb E e^{\eta V_n}<\infty,                 \tag{70.2}
\]
for the SM--Einstein conformal equation.
\end{proposition}

\begin{proof}
Its certified object is a finite matrix pencil on a proper
momentum plate.  No conformal factor \(u\), elliptic potential \(q\),
metric probability measure, or cofinal normalization occurs in
the theorem.  Therefore none of the quantities in (70.2) is
defined by that result.
\end{proof}

\section{Targeted Navier--Stokes check}

The NS source contains Morrey lower bounds, concentration-point
counts, local propagation, and compact rescaled families.  These
statements are conditional on the Navier--Stokes evolution,
parabolic scaling, vorticity/enstrophy ledgers, and stated
unresolved gates.  Its active V26 addition concerns pointwise
rank-one Oseen-kernel derivative norms.

\begin{proposition}[Why the NS Morrey mechanism does not transfer]
\label{prop:v15v70-ns}
The NS concentration theorems do not prove the small-potential
estimate (68.2), equiintegrability (68.8), or the Harnack hypothesis
(69.13).
\end{proposition}

\begin{proof}
The NS critical quantities are three-dimensional parabolic energy,
enstrophy, and Gaussian or cell-localized flux.  Equations
(68.1) and (69.1) are four-dimensional elliptic equations for a
positive conformal factor, with critical potential norm
\(\|q\|_{L^2}\).  No map in the NS note sends its solution,
measure, scaling, or kernel estimates to \(u\) and \(q\) while
preserving either equation.  Its own summary also states that
rate-free epsilon regularity and the closing Liouville gates remain
unproved.
\end{proof}

The shared structural lesson is valid but already implemented in
V68: bounded critical energy yields a finite concentration ledger,
and regularity holds away from the concentration set.  Importing
the NS conclusion itself would add an unsupported cross-equation
identification.

\section{Audit conclusion}

\begin{theorem}[Dependency state at V70]
\label{thm:v15v70-conclusion}
No companion result added since V65 populates a missing hypothesis
of V68 or V69.  The next noncircular source to test is the traced
Einstein equation, which directly relates scalar curvature to the
renormalized matter trace when the higher-derivative regulator has
been removed.
\end{theorem}

\begin{proof}
The source identities follow from (70.1).
Propositions~\ref{prop:v15v70-ym} and
\ref{prop:v15v70-ns} rule out the only targeted transfers.
The classical Einstein trace contains \(R_g\), and hence the
conformal potential \(q=R_gu^2\), so it is the next direct
interface rather than an analogy with another equation.
\end{proof}

\begin{assessment}[Physical status]
\label{ass:v15v70-status}
V70 prevents a false transfer of parabolic Morrey structure or
finite symbol positivity into the conformal metric problem.  It is
an audit and proves no new continuum existence statement.
\end{assessment}

\begin{decision}[V71 traced Einstein source]
\label{dec:v15v70-next}
Derive the weak trace equation with all stress, anomaly, and
counterterm rows visible.  Determine an explicit matter-trace
integrability condition that bootstraps the V68 \(L^6\) conformal
factor to \(L^\infty\) and \(W^{2,p}\).  Treat an \(R^2\) regulator
as a fourth-order trace term and prove that its vanishing
coefficient does not imply decoupling without a derivative bound.
\end{decision}

\begin{firewall}[Claim boundary after V70]
\label{fire:v15v70}
V70 records a dependency comparison only.  It does not prove the
Einstein trace equation, matter higher integrability, anomaly
cancellation, conformal noncollapse, or a continuum measure.
\end{firewall}

\section*{Version V70 append-only record}
\addcontentsline{toc}{section}{Version V70 append-only record}
The complete V69 source was retained before this continuation.  It
had \(687601\) bytes, \(19915\) lines, and SHA--256
\[
 \texttt{48A49108008ADC82C7EB1378AFD9CD5D6DC7CF2BF41161133322EF4E487CFF81}.
\]
V70 performed the scheduled YM--NS audit, distinguished the NS
Morrey analogy from a valid elliptic transfer, and retained the
traced Einstein equation as the next direct source of conformal
regularity.

\chapter{Fifteenth-cycle continuation V71: the traced Einstein source and its off-shell limit}
\label{sec:v15v71}

\section{Complete traced residual}

On a classical or stationary effective-action branch, write the
renormalized metric equation in the convention
\[
 G_{\mu\nu}+\Lambda g_{\mu\nu}
 =
 8\pi G\,T^{\rm tot}_{\mu\nu}.                         \tag{71.1}
\]
The total tensor includes every differentiated row,
\[
 T^{\rm tot}
 =
 T_{\rm gauge}+T_{\rm Higgs}+T_{\rm fermion}
 +T_{\rm ghost}+T_{\rm quotient}
 +T_{\rm ct}+T_{\rm anomaly}.                         \tag{71.2}
\]
Terms may be regrouped between the two sides of (71.1), but none
may be counted twice.

\begin{proposition}[Four-dimensional trace equation]
\label{prop:v15v71-trace}
For a four-dimensional solution of (71.1),
\[
 R_g=4\Lambda-8\pi G\,T^{\rm tot},\qquad
 T^{\rm tot}=g^{\mu\nu}T^{\rm tot}_{\mu\nu}.           \tag{71.3}
\]
If \(g=u^2\bar g\), its conformal factor satisfies
\[
 \boxed{
 L_{\bar g}u=A\,u^3,\qquad
 A:=4\Lambda-8\pi G\,T^{\rm tot}.}                    \tag{71.4}
\]
\end{proposition}

\begin{proof}
In four dimensions,
\[
 g^{\mu\nu}G_{\mu\nu}
 =
 R_g-\frac42R_g=-R_g.
\]
Taking the trace of (71.1) gives
\(-R_g+4\Lambda=8\pi G\,T^{\rm tot}\), which is
(71.3).  The conformal identity \(R_g=u^{-3}L_{\bar g}u\)
from (66.2) gives (71.4).
\end{proof}

Equation (71.4) is critical: under the four-dimensional scaling,
the cubic power has the same \(L^4\) conformal-volume exponent as
the left side.

\section{A matter-trace bootstrap after epsilon regularity}

\begin{theorem}[Trace-source supercritical bootstrap]
\label{thm:v15v71-trace-bootstrap}
Let \(r>6\).  Suppose positive solutions of (71.4) satisfy
\[
 \sup_j\|u_j\|_{L^6(\bar g)}\leq C_6,\qquad
 \sup_j\|A_j\|_{L^r(\bar g)}\leq C_A.                 \tag{71.5}
\]
Then
\[
 \sup_j\left(
 \|u_j\|_{L^\infty}
 +\|u_j\|_{W^{2,r}}
 \right)<\infty.                                      \tag{71.6}
\]
The same conclusion holds for \(r=\infty\).
\end{theorem}

\begin{proof}
Assume first \(6<r<\infty\).  If \(u_j\) is bounded in
\(L^{s_k}\), H\"older gives
\[
 A_ju_j^3\in L^{t_k},\qquad
 \frac1{t_k}=\frac1r+\frac3{s_k}.                     \tag{71.7}
\]
As long as \(t_k<2\), elliptic regularity and Sobolev embedding
give a new exponent satisfying
\[
 \frac1{s_{k+1}}
 =
 \frac1r+\frac3{s_k}-\frac12.                         \tag{71.8}
\]
Start with \(s_0=6\).  Equation (71.8) gives
\[
 \frac1{s_1}=\frac1r.
\]
Put \(a_k=1/s_k\) and
\[
 a_*=\frac14-\frac1{2r}.
\]
For \(r>6\), \(a_0=1/6<a_*\), and
\[
 a_{k+1}-a_*=3(a_k-a_*).                              \tag{71.9}
\]
Thus after finitely many steps the right side of (71.8) is
nonpositive, equivalently some \(t_k\geq2\).  At equality,
\(W^{2,2}\) embeds into every finite \(L^s\), so one more step makes
\(t_k>2\).  For \(t_k>2\),
\[
 W^{2,t_k}(M)\hookrightarrow L^\infty(M),
\]
with uniform constants throughout the finite iteration.  Then
\[
 \|A_ju_j^3\|_r
 \leq C_A\|u_j\|_\infty^3,
\]
and the \(L^r\) elliptic estimate proves the \(W^{2,r}\) bound.

If \(r=\infty\), (71.7) with \(s_0=6\) gives \(t_0=2\).
Use the critical embedding into an arbitrarily large finite
\(L^s\), repeat once to obtain \(t>2\), and conclude as above.
\end{proof}

\begin{corollary}[Stationary conformal carrier]
\label{cor:v15v71-stationary-carrier}
Assume the no-atom or equiintegrable-curvature hypothesis of V68,
so that the \(L^6\) bound in (71.5) holds.  If the complete traced
source in (71.3) is uniformly bounded in
\(L^r(\bar g)\) for some \(r>6\), then the stationary conformal
metrics form a bounded \(W^{2,r}\) family.  After a subsequence they
are compact in \(C^{1,\beta}\) for every
\[
 0<\beta<1-\frac4r.                                   \tag{71.10}
\]
\end{corollary}

\begin{proof}
Apply Theorem~\ref{thm:v15v71-trace-bootstrap}.  Since \(r>6>4\),
Morrey--Rellich gives compactness in \(C^{1,\beta}\) below the
endpoint exponent.
\end{proof}

The threshold \(r>6\) belongs to this particular bootstrap from the
V68 \(L^6\) input.  The natural energy bounds on Standard Model
fields do not automatically put the complete trace in such a high
space.

\section{Trace-anomaly and regularity ledger}

At the quantum level the renormalized trace may contain
\[
 T_{\rm anomaly}
 =
 c\,W_{\mu\nu\rho\sigma}W^{\mu\nu\rho\sigma}
 -a\,E_4+b\,\Box_gR_g
 +\sum_i\beta_i{\cal O}_i
 +\hbox{mass terms}.                                  \tag{71.11}
\]
The coefficients and operator mixings depend on the renormalized
theory.  Equation (71.11) is a ledger, not an assertion that the
terms cancel in the Standard Model.

\begin{proposition}[Energy-level mismatch]
\label{prop:v15v71-energy-mismatch}
An \(L^2\) curvature bound controls the curvature-quadratic terms in
(71.11) only in \(L^1\).  It does not imply the
\(L^r\), \(r>6\), hypothesis in (71.5).
\end{proposition}

\begin{proof}
If \({\rm Rm}\in L^2\), then
\[
 \||{\rm Rm}|^2\|_1=\|{\rm Rm}\|_2^2.
\]
There is no embedding \(L^1\subset L^r\) for \(r>1\) on a
non-atomic finite-volume space.  Concentrating functions with fixed
\(L^1\) norm give an explicit counterexample.
\end{proof}

The derivative term \(\Box R\) is distributional unless additional
curvature regularity is proved.  It cannot be inserted into
Theorem~\ref{thm:v15v71-trace-bootstrap} as an \(L^r\) function
without that proof.

\section{An \(R^2\) regulator changes the trace equation}

Let
\[
 H^{(R^2)}_{\mu\nu}
 =
 2R R_{\mu\nu}-\frac12R^2g_{\mu\nu}
 +2(g_{\mu\nu}\Box_g-\nabla_\mu\nabla_\nu)R.           \tag{71.12}
\]
In four dimensions,
\[
 g^{\mu\nu}H^{(R^2)}_{\mu\nu}=6\Box_gR.               \tag{71.13}
\]

\begin{proposition}[Regulated trace equation]
\label{prop:v15v71-r2-trace}
If the regulated field equation is
\[
 G_{\mu\nu}+\Lambda g_{\mu\nu}
 +\alpha H^{(R^2)}_{\mu\nu}
 =
 8\pi G\,T_{\mu\nu},                                  \tag{71.14}
\]
then
\[
 -R+4\Lambda+6\alpha\Box_gR=8\pi G\,T.                \tag{71.15}
\]
\end{proposition}

\begin{proof}
Trace (71.14), use
\(g^{\mu\nu}G_{\mu\nu}=-R\), and compute from (71.12)
\[
 2R^2-2R^2+2(4\Box R-\Box R)=6\Box R.
\]
\end{proof}

\begin{theorem}[Distributional decoupling criterion]
\label{thm:v15v71-decoupling}
Let \(\alpha_n\downarrow0\) and let \(g_n\) be smooth metrics on a
closed manifold, uniformly equivalent to a fixed background with
uniformly controlled volume.  If
\[
 \sup_n\alpha_n\int_MR_{g_n}^2\,dV_{g_n}<\infty,       \tag{71.16}
\]
then
\[
 \alpha_n\Box_{g_n}R_{g_n}\longrightarrow0            \tag{71.17}
\]
distributionally against every fixed smooth test function whose
\(\Box_{g_n}\)-norm is uniformly bounded.
\end{theorem}

\begin{proof}
For a test function \(\zeta\), integration by parts twice in the
weak definition gives
\[
 \left|
 \alpha_n\int_M\zeta\Box_{g_n}R_{g_n}\,dV_{g_n}
 \right|
 =
 \left|
 \alpha_n\int_MR_{g_n}\Box_{g_n}\zeta\,dV_{g_n}
 \right|.
\]
By Cauchy--Schwarz, this is at most
\[
 \alpha_n
 \|R_{g_n}\|_2
 \|\Box_{g_n}\zeta\|_2
 \leq
 C_\zeta
 \alpha_n^{1/2}
 \left(\alpha_n\|R_{g_n}\|_2^2\right)^{1/2},
\]
which tends to zero by (71.16).
\end{proof}

\begin{counterexample}[A vanishing coefficient alone is insufficient]
\label{cnt:v15v71-decoupling}
On a fixed flat torus, choose a nonharmonic smooth function \(f\)
and set
\[
 R_n=\alpha_n^{-1}f.                                  \tag{71.18}
\]
Then
\[
 \alpha_n\bar\Delta R_n=\bar\Delta f                  \tag{71.19}
\]
does not tend to zero.  The energy row (71.16) fails, since
\(\alpha_n\|R_n\|_2^2=\alpha_n^{-1}\|f\|_2^2\).
\end{counterexample}

Thus an \(R^2\) regulator can decouple weakly under its natural
bounded-action estimate; the statement does not follow from
\(\alpha_n\to0\) alone.

\section{Off-shell quantum firewall}

The continuum Schwinger--Dyson equation of V62 is
\[
 \mathbb E_\nu[D_kF]
 =
 \mathbb E_\nu[F{\cal R}_{\rm Ein}(k)].                \tag{71.20}
\]
It is not a statement that
\({\cal R}_{\rm Ein}(k)=0\) on almost every configuration.

\begin{proposition}[An averaged residual need not vanish samplewise]
\label{prop:v15v71-off-shell}
Even if (71.20) with \(F=1\) gives
\[
 \mathbb E_\nu{\cal R}_{\rm Ein}(k)=0,                \tag{71.21}
\]
one cannot infer the samplewise trace equation (71.3).
\end{proposition}

\begin{proof}
On a two-point probability space, let a residual \(X\) take values
\(+1\) and \(-1\) with equal probability.  Then
\(\mathbb EX=0\), but \(X\ne0\) everywhere.  Equation (71.21) has
exactly this logical form.  More generally, (71.20) determines
integration-by-parts relations among correlations; it does not
concentrate the measure on the zero set of the score.
\end{proof}

\begin{corollary}[Domain of the trace bootstrap]
\label{cor:v15v71-domain}
Theorem~\ref{thm:v15v71-trace-bootstrap} applies directly to
classical solutions, stationary effective-action configurations,
or a quantum measure already proved to be supported on such
solutions.  It cannot be applied configuration by configuration to
a generic off-shell Euclidean path-integral measure using only
(71.20).
\end{corollary}

\begin{proof}
The bootstrap starts from the samplewise PDE (71.4).
Proposition~\ref{prop:v15v71-off-shell} shows that (71.20) does not
supply that PDE.
\end{proof}

\section{Revised source ledger}

For the stationary branch, the route is
\[
\begin{gathered}
\hbox{no curvature atoms}
\Longrightarrow u\in L^6,\\
T^{\rm tot}\in L^r,\ r>6
\Longrightarrow A\in L^r
\Longrightarrow u\in W^{2,r}\cap L^\infty,            \tag{71.22}\\
\hbox{Harnack for }q=Au^2
\Longrightarrow \inf u>0.
\end{gathered}
\]
For the quantum continuum measure, the route must instead derive
moments from the full Schwinger--Dyson family, convexity, or direct
coercivity of the action.  The stationary PDE cannot be substituted
for that measure-level argument.

\begin{assessment}[Physical status]
\label{ass:v15v71-status}
V71 derives the complete traced stationary equation, proves an
\(L^6\)-plus-\(L^r\) matter-source bootstrap, isolates the
energy-level anomaly mismatch, and proves a sharp weak-decoupling
criterion for an \(R^2\) regulator.  It also closes a possible
circularity by proving that the quantum Schwinger--Dyson identity
does not put individual sampled metrics on shell.

The programme still lacks the \(L^r\) complete matter trace on the
stationary branch and a measure-level mechanism producing
supercritical conformal moments off shell.
\end{assessment}

\begin{decision}[V72 score-moment Schwinger--Dyson compiler]
\label{dec:v15v71-next}
Remain on the quantum branch and test nonlinear observables in the
metric integration-by-parts identity.  Prove that a monotone or
uniformly convex conformal score yields \(L^2\) and exponential
score moments, identify the Hessian and boundary terms required,
and give a counterexample showing that the Ward identity with
\(F=1\) alone gives no variance control.
\end{decision}

\begin{firewall}[Claim boundary after V71]
\label{fire:v15v71}
V71 proves stationary trace and regulator-decoupling compilers and
an off-shell non-implication.  It does not derive physical
matter-trace regularity, score convexity, anomaly cancellation, or
the full quantum continuum measure.
\end{firewall}

\section*{Version V71 append-only record}
\addcontentsline{toc}{section}{Version V71 append-only record}
The complete V70 source was retained before this continuation.  It
had \(692983\) bytes, \(20052\) lines, and SHA--256
\[
 \texttt{DF74414E1C4914AA440FB2408B42FC053B92D852FA97A57F268526405CE17C8E}.
\]
V71 derived the traced stationary Einstein equation, proved the
matter-source conformal bootstrap, established distributional
decoupling of a bounded-action \(R^2\) regulator, and proved that
the averaged quantum Einstein identity cannot be used as a
samplewise elliptic equation.

\chapter{Fifteenth-cycle continuation V72: score moments from metric Schwinger--Dyson identities}
\label{sec:v15v72}

\section{Positive modulus measure and score}

At a finite cutoff, work first with the positive modulus measure
\[
 d\nu_n(x)=Z_n^{-1}e^{-{\cal A}_n(x)}\,dx             \tag{72.1}
\]
on a Euclidean regulator space.  Reference-density and coordinate
Jacobian terms may be absorbed into \({\cal A}_n\) only when their
derivatives are retained.  For a constant normalized direction
\(v\), define
\[
 {\cal S}_{n,v}=D_v{\cal A}_n.                        \tag{72.2}
\]
The complex determinant phase is treated as an observable relative
to (72.1); it is not part of the positive density used for
logarithmic Sobolev estimates.

\section{The exact Fisher--Hessian identity}

\begin{theorem}[Score-square identity]
\label{thm:v15v72-fisher}
Assume the boundary term vanishes and
\({\cal S}_{n,v}\), \({\cal S}_{n,v}^2\), and
\(D_v{\cal S}_{n,v}\) are integrable.  Then
\[
 \mathbb E_{\nu_n}{\cal S}_{n,v}=0,                   \tag{72.3}
\]
and
\[
 \boxed{
 \mathbb E_{\nu_n}{\cal S}_{n,v}^2
 =
 \mathbb E_{\nu_n}D_v^2{\cal A}_n.}                   \tag{72.4}
\]
\end{theorem}

\begin{proof}
The finite-cutoff identity (62.12) with \(F=1\) gives (72.3).
Use it again with \(F={\cal S}_{n,v}\):
\[
 \mathbb E D_v{\cal S}_{n,v}
 =
 \mathbb E{\cal S}_{n,v}^2.
\]
Since \(D_v{\cal S}_{n,v}=D_v^2{\cal A}_n\), this is
(72.4).  Unbounded scores are justified by applying the identity
to a smooth truncation and passing to the limit under the stated
integrability assumptions.
\end{proof}

\begin{corollary}[Directional variance from an upper Hessian bound]
\label{cor:v15v72-variance}
If
\[
 D_v^2{\cal A}_n\leq H_v                              \tag{72.5}
\]
almost everywhere, with \(H_v\) independent of \(n\), then
\[
 \sup_n\mathbb E_{\nu_n}{\cal S}_{n,v}^2\leq H_v.      \tag{72.6}
\]
\end{corollary}

\begin{proof}
Take expectations in (72.5) and use (72.4).
\end{proof}

An upper directional Hessian controls the score variance.  A lower
Hessian bound instead supplies concentration of the measure, but
the two roles must not be interchanged.

\section{Boundary-corrected score identity}

For a nonconstant regulator vector field \(B_n\), use the geometric
score (64.2).  The boundary-aware identity (64.21) gives
\[
 \mathbb E D_{B_n}F
 =
 \mathbb E[F{\cal S}_{n,B}]
 +{\cal B}_{n,B}(F).                                  \tag{72.7}
\]

\begin{proposition}[Boundary defects in the first two moments]
\label{prop:v15v72-boundary}
Whenever the terms exist,
\[
 \mathbb E{\cal S}_{n,B}
 =
 -{\cal B}_{n,B}(1),                                  \tag{72.8}
\]
\[
 \mathbb E{\cal S}_{n,B}^2
 =
 \mathbb E D_{B_n}{\cal S}_{n,B}
 -{\cal B}_{n,B}({\cal S}_{n,B}).                     \tag{72.9}
\]
\end{proposition}

\begin{proof}
Put \(F=1\) and \(F={\cal S}_{n,B}\), respectively, in (72.7), and
rearrange.
\end{proof}

Thus quotient-boundary cancellation is needed even to center the
score.  A small defect for bounded observables does not
automatically control
\({\cal B}_{n,B}({\cal S}_{n,B})\), whose insertion is unbounded.

\section{Uniform logarithmic Sobolev concentration}

\begin{definition}[Logarithmic Sobolev inequality]
\label{def:v15v72-lsi}
A probability measure \(\nu\) satisfies
\({\rm LSI}(C_{\rm LS})\) when every smooth \(f\) obeys
\[
 \operatorname{Ent}_\nu(f^2)
 :=
 \mathbb E_\nu\!\left[
 f^2\log\frac{f^2}{\mathbb E_\nu f^2}
 \right]
 \leq
 2C_{\rm LS}\mathbb E_\nu|\nabla f|^2.                \tag{72.10}
\]
\end{definition}

\begin{theorem}[Herbst score bound]
\label{thm:v15v72-herbst}
Suppose the measures \(\nu_n\) satisfy (72.10) with one constant
\(C_{\rm LS}\), the boundary defect vanishes, and
\[
 |\nabla{\cal S}_{n,v}|\leq L_v                       \tag{72.11}
\]
almost everywhere with \(L_v\) independent of \(n\).  Then for
every \(\lambda\in{\mathbb R}\),
\[
 \mathbb E_{\nu_n}e^{\lambda{\cal S}_{n,v}}
 \leq
 \exp\!\left(\frac12C_{\rm LS}L_v^2\lambda^2\right),   \tag{72.12}
\]
and for every \(\tau>0\),
\[
 \sup_n\mathbb E_{\nu_n}e^{\tau|{\cal S}_{n,v}|}
 \leq
 2\exp\!\left(\frac12C_{\rm LS}L_v^2\tau^2\right).
                                                               \tag{72.13}
\]
\end{theorem}

\begin{proof}
By (72.3), the score is centered.  Put
\[
 \psi_n(\lambda)
 =
 \log\mathbb E_{\nu_n}e^{\lambda{\cal S}_{n,v}}.
\]
Apply (72.10) to
\(f=e^{\lambda{\cal S}_{n,v}/2}\).  Using (72.11) and dividing by
\(\mathbb E e^{\lambda{\cal S}_{n,v}}\) gives
\[
 \lambda\psi_n'(\lambda)-\psi_n(\lambda)
 \leq
 \frac12C_{\rm LS}L_v^2\lambda^2.                    \tag{72.14}
\]
For \(\lambda>0\), divide by \(\lambda^2\) and integrate the
derivative of \(\psi_n(\lambda)/\lambda\) from zero.  Since
\(\psi_n'(0)=0\), this proves (72.12).  Apply the same argument to
\(-{\cal S}_{n,v}\) for negative \(\lambda\).  Finally,
\[
 e^{\tau|s|}\leq e^{\tau s}+e^{-\tau s},
\]
which gives (72.13).
\end{proof}

\begin{corollary}[Uniform-integrability row for Einstein response]
\label{cor:v15v72-ui}
Under Theorem~\ref{thm:v15v72-herbst}, the score and its product
with every cutoff-uniformly bounded local observable are uniformly
integrable.  Therefore the score-product hypothesis of
Theorem~\ref{thm:v15v63-vitali} is satisfied.
\end{corollary}

\begin{proof}
Equation (72.13) is exactly the exponential moment in
Corollary~\ref{cor:v15v63-exp-ui}.
\end{proof}

\section{Strong convexity as an LSI source}

\begin{theorem}[Finite-dimensional Bakry--Emery compiler]
\label{thm:v15v72-bakry}
Suppose
\[
 \nabla^2{\cal A}_n(x)\succeq mI,\qquad m>0,           \tag{72.15}
\]
on the full regulator space, uniformly in \(n\), with no boundary
or with a convex reflecting boundary.  Then \(\nu_n\) satisfies
\[
 {\rm LSI}(1/m).                                      \tag{72.16}
\]
\end{theorem}

\begin{proof}
Let
\[
 {\mathscr L}_n=\Delta-\nabla{\cal A}_n\mathbin{\cdot}\nabla
\]
and let \(P_t\) be its Markov semigroup.  Integration by parts shows
that \(\nu_n\) is invariant and symmetric.  The carré-du-champ
calculation gives
\[
 \Gamma_2(f)
 =
 \|\nabla^2f\|_{\rm HS}^2+
 \langle\nabla^2{\cal A}_n\nabla f,\nabla f\rangle
 \geq m|\nabla f|^2.                                  \tag{72.17}
\]
Consequently the semigroup gradient estimate
\[
 |\nabla P_tf|^2\leq e^{-2mt}P_t|\nabla f|^2           \tag{72.18}
\]
follows by differentiating
\(P_s|\nabla P_{t-s}f|^2\) and using (72.17).

For positive \(h\), entropy dissipation and invariance give
\[
 \operatorname{Ent}_{\nu_n}(h)
 =
 \int_0^\infty
 \mathbb E_{\nu_n}
 \frac{|\nabla P_th|^2}{P_th}\,dt.                    \tag{72.19}
\]
Apply (72.18) and Cauchy--Schwarz for the Markov kernel to bound the
integrand by
\[
 e^{-2mt}
 \mathbb E_{\nu_n}\frac{|\nabla h|^2}{h}.
\]
Integration in \(t\) yields
\[
 \operatorname{Ent}_{\nu_n}(h)
 \leq
 \frac1{2m}
 \mathbb E_{\nu_n}\frac{|\nabla h|^2}{h}.             \tag{72.20}
\]
Taking \(h=f^2\) gives
\[
 \operatorname{Ent}_{\nu_n}(f^2)
 \leq\frac2m\mathbb E_{\nu_n}|\nabla f|^2,
\]
which is (72.16).  Convex reflection adds no adverse boundary term.
\end{proof}

\begin{corollary}[Uniform Hessian corridor]
\label{cor:v15v72-corridor}
If, in addition to (72.15),
\[
 \|\nabla^2{\cal A}_n(x)v\|\leq L_v                   \tag{72.21}
\]
uniformly in \(x,n\), then
\[
 \sup_n\mathbb E_{\nu_n}e^{\tau|{\cal S}_{n,v}|}
 \leq
 2\exp\!\left(\frac{L_v^2\tau^2}{2m}\right).           \tag{72.22}
\]
\end{corollary}

\begin{proof}
The score gradient is
\(\nabla{\cal S}_{n,v}=\nabla^2{\cal A}_n v\).
Use Theorems~\ref{thm:v15v72-bakry} and
\ref{thm:v15v72-herbst}.
\end{proof}

\section{Rowwise Hessian assembly}

\begin{proposition}[Positive core with form-small remainder]
\label{prop:v15v72-form-small}
Write
\[
 \nabla^2{\cal A}_n=H_{0,n}+R_n.                      \tag{72.23}
\]
If
\[
 H_{0,n}\succeq m_0I,\qquad
 \langle R_nh,h\rangle\geq-\theta m_0\|h\|^2,
 \quad0\leq\theta<1,                                  \tag{72.24}
\]
uniformly in \(n,x\), then
\[
 \nabla^2{\cal A}_n\succeq(1-\theta)m_0I.             \tag{72.25}
\]
\end{proposition}

\begin{proof}
Add the two quadratic-form inequalities in (72.24).
\end{proof}

The positive core may come from a gauge-fixed biharmonic
stabilizer.  Every gravitational, Higgs, gauge, determinant, ghost,
quotient, and finite-counterterm Hessian must enter \(R_n\).
The bare Euclidean Einstein conformal direction violates (72.15),
as the V66 high-frequency sequence shows.

\section{Determinant Hessian row}

For the relative determinant in (63.15), differentiating (63.19)
once more gives
\[
\begin{split}
 D_v^2\log\det_{\!F}(I+K_n)
 &=
 \operatorname{Tr}\bigl((I+K_n)^{-1}D_v^2K_n\bigr)\\
 &\quad-
 \operatorname{Tr}\bigl(
 (I+K_n)^{-1}D_vK_n
 (I+K_n)^{-1}D_vK_n
 \bigr).                                              \tag{72.26}
\end{split}
\]

\begin{proposition}[Trace-ideal Hessian bound]
\label{prop:v15v72-det-hessian}
If
\[
 \|(I+K_n)^{-1}\|\leq c_0,                            \tag{72.27}
\]
then
\[
 \left|D_v^2\log\det_{\!F}(I+K_n)\right|
 \leq
 c_0\|D_v^2K_n\|_1+
 c_0^2\|D_vK_n\|_2^2.                                \tag{72.28}
\]
\end{proposition}

\begin{proof}
Use
\(|\operatorname{Tr}(AB)|\leq\|A\|\|B\|_1\)
for the first term.  For the second, the product of two
Hilbert--Schmidt operators is trace class and
\(\|XY\|_1\leq\|X\|_2\|Y\|_2\).  Insert the two bounded inverse
factors to obtain (72.28).
\end{proof}

The inverse bound fails near a determinant zero.  Consequently
neither the Hessian corridor nor the LSI conclusion can cross that
divisor without the zero-mode construction of V51.

\section{The first Ward identity has no variance content}

\begin{counterexample}[Centered scores with unbounded variance]
\label{cnt:v15v72-centered}
Let \(\nu_\sigma\) be the centered Gaussian on \(\mathbb R\) with
variance \(\sigma^2\):
\[
 {\cal A}_\sigma(x)=\frac{x^2}{2\sigma^2},\qquad
 {\cal S}_\sigma(x)=\frac{x}{\sigma^2}.                \tag{72.29}
\]
Then
\[
 \mathbb E_{\nu_\sigma}{\cal S}_\sigma=0
\]
for every \(\sigma>0\), while
\[
 \mathbb E_{\nu_\sigma}{\cal S}_\sigma^2
 =\sigma^{-2}\longrightarrow\infty
 \quad\hbox{as }\sigma\downarrow0.                    \tag{72.30}
\]
\end{counterexample}

\begin{proof}
Both identities follow from
\(\mathbb E x=0\) and \(\mathbb E x^2=\sigma^2\).
\end{proof}

Thus the \(F=1\) Schwinger--Dyson equation centers the score but
does not control its variance.  One needs the \(F={\cal S}\)
identity and an upper Hessian bound, or the LSI--Lipschitz route for
exponential moments.

\begin{assessment}[Physical status]
\label{ass:v15v72-status}
V72 proves the exact score-square/Hessian identity, incorporates
boundary defects, derives cutoff-uniform exponential score moments
from LSI and a Lipschitz score, and supplies strong-convexity,
form-smallness, and determinant trace-ideal compilers.

The complete SM--Einstein modulus action is not known to be
uniformly convex.  Its bare conformal direction is not convex, its
determinant inverse may fail at zero modes, and quotient boundaries
remain.  Hence the exponential stress moment is conditionally
compiled but not physically populated.
\end{assessment}

\begin{decision}[V73 nonconvex perturbation stability]
\label{dec:v15v72-next}
Relax global convexity without losing the LSI.  Prove a bounded
oscillation perturbation theorem for a strongly log-concave
reference action, then formulate a local convexity plus Lyapunov
tail criterion for unbounded interactions.  Test determinant and
Higgs rows against those hypotheses and retain conformal escape as
an explicit failure mode.
\end{decision}

\begin{firewall}[Claim boundary after V72]
\label{fire:v15v72}
V72 proves measure-level score-moment implications under explicit
functional inequalities.  It does not prove those inequalities for
the physical regulator, remove determinant zeros, cancel quotient
flux, or construct the full continuum.
\end{firewall}

\section*{Version V72 append-only record}
\addcontentsline{toc}{section}{Version V72 append-only record}
The complete V71 source was retained before this continuation.  It
had \(704747\) bytes, \(20438\) lines, and SHA--256
\[
 \texttt{2F476441FEB1056311ABF0552CFD1A86DBF310B721D524E4BB3BC0236ED0B3BB}.
\]
V72 tested the score inside the quantum metric
Schwinger--Dyson identity, proved exact second-moment and
boundary-defect formulas, and derived exponential stress moments
from a cutoff-uniform logarithmic Sobolev inequality and a
trace-ideal Hessian corridor.

\chapter{Fifteenth-cycle continuation V73: stability under nonconvex perturbations}
\label{sec:v15v73}

\section{Bounded perturbations of a positive measure}

Let \(\mu\) be a probability measure satisfying
\({\rm LSI}(C_\mu)\), and define
\[
 d\nu=Z_W^{-1}e^{-W}\,d\mu.                           \tag{73.1}
\]
Write
\[
 \operatorname{osc}W=\operatorname*{ess\,sup}W
 -\operatorname*{ess\,inf}W.                         \tag{73.2}
\]

\begin{theorem}[Bounded-oscillation perturbation]
\label{thm:v15v73-holley}
If
\[
 \operatorname{osc}W\leq D<\infty,                    \tag{73.3}
\]
then \(\nu\) satisfies
\[
 {\rm LSI}(e^DC_\mu).                                 \tag{73.4}
\]
\end{theorem}

\begin{proof}
Let \(h=d\nu/d\mu\).  Normalization cancels in ratios, so
\[
 \frac{\operatorname*{ess\,sup}h}
 {\operatorname*{ess\,inf}h}\leq e^D.                 \tag{73.5}
\]
For \(s,a\geq0\), put
\[
 \Phi_a(s)=s\log(s/a)-s+a\geq0.
\]
The variational entropy formula is
\[
 \operatorname{Ent}_\nu(f^2)
 =
 \inf_{a\geq0}\int\Phi_a(f^2)\,d\nu.                  \tag{73.6}
\]
Therefore
\[
\begin{split}
 \operatorname{Ent}_\nu(f^2)
 &\leq
 (\operatorname*{ess\,sup}h)
 \operatorname{Ent}_\mu(f^2)\\
 &\leq
 2C_\mu(\operatorname*{ess\,sup}h)
 \int|\nabla f|^2d\mu\\
 &\leq
 2C_\mu
 \frac{\operatorname*{ess\,sup}h}
 {\operatorname*{ess\,inf}h}
 \int|\nabla f|^2d\nu.
\end{split}
\]
Use (73.5).
\end{proof}

\begin{corollary}[Unbounded convex part plus bounded wells]
\label{cor:v15v73-convex-bounded}
Suppose
\[
 {\cal A}={\cal A}_0+W_c+W_b,                         \tag{73.7}
\]
where
\[
 \nabla^2({\cal A}_0+W_c)\succeq mI,\qquad
 \operatorname{osc}W_b\leq D.                         \tag{73.8}
\]
The term \(W_c\) may be unbounded.  Then the probability measure
with action \({\cal A}\) satisfies
\[
 {\rm LSI}(e^D/m).                                    \tag{73.9}
\]
\end{corollary}

\begin{proof}
Theorem~\ref{thm:v15v72-bakry} applies to
\({\cal A}_0+W_c\).  Apply
Theorem~\ref{thm:v15v73-holley} to \(W_b\).
\end{proof}

Uniformity along the regulator sequence requires \(m>0\) and
\(D<\infty\) independent of the growing cutoff dimension.

\section{Measure concentration does not control the score derivative}

\begin{proposition}[Perturbed-score compiler]
\label{prop:v15v73-score}
Under Corollary~\ref{cor:v15v73-convex-bounded}, fix a direction
\(v\).  If
\[
 \left|
 \nabla D_v({\cal A}_0+W_c+W_b)
 \right|\leq L_v                                      \tag{73.10}
\]
uniformly, then
\[
 \mathbb E_\nu e^{\tau|D_v{\cal A}|}
 \leq
 2\exp\!\left(
 \frac{e^DL_v^2\tau^2}{2m}
 \right)                                              \tag{73.11}
\]
provided the boundary defect vanishes.
\end{proposition}

\begin{proof}
The score is centered by the first Schwinger--Dyson identity.
Combine (73.9), (73.10), and
Theorem~\ref{thm:v15v72-herbst}.
\end{proof}

\begin{counterexample}[Bounded action oscillation with divergent score]
\label{cnt:v15v73-oscillation}
On \(\mathbb R\), let
\[
 {\cal A}_N(x)=\frac{x^2}{2}+\epsilon\sin(Nx),
 \qquad0<\epsilon<1.                                  \tag{73.12}
\]
The measures \(Z_N^{-1}e^{-{\cal A}_N}dx\) satisfy a uniform LSI,
because the perturbation has oscillation \(2\epsilon\).  Their
scores are
\[
 {\cal S}_N(x)=x+\epsilon N\cos(Nx).                  \tag{73.13}
\]
There is \(c_\epsilon>0\) such that
\[
 \mathbb E_{\nu_N}{\cal S}_N^2\geq c_\epsilon N^2
 \quad\hbox{for all sufficiently large }N.            \tag{73.14}
\]
\end{counterexample}

\begin{proof}
The density ratio between \(\nu_N\) and the standard Gaussian
\(\gamma\) is bounded above and below by constants depending only
on \(\epsilon\).  Hence
\[
 \mathbb E_{\nu_N}\cos^2(Nx)
 \geq c_\epsilon\mathbb E_\gamma\cos^2(Nx)
 =
 \frac{c_\epsilon}{2}(1+e^{-2N^2}).                   \tag{73.15}
\]
Also \(\mathbb E_{\nu_N}x^2\) is uniformly bounded.  From
\((a+b)^2\geq\tfrac12b^2-a^2\), with
\(a=x\) and \(b=\epsilon N\cos(Nx)\), one obtains (73.14).
The uniform LSI follows from
Theorem~\ref{thm:v15v73-holley}.
\end{proof}

Thus bounded determinant or counterterm oscillation can preserve
measure concentration while its metric derivative becomes
cutoff-large.  The derivative row (73.10) is independent.

\section{A Lyapunov route to a spectral gap}

For unbounded nonconvex interactions, a bounded-oscillation
decomposition may fail.  Let
\[
 {\mathscr L}=\Delta-\nabla{\cal A}\mathbin{\cdot}\nabla
\]
be symmetric in \(L^2(\nu)\).

\begin{lemma}[Lyapunov integration estimate]
\label{lem:v15v73-lyapunov}
For every smooth \(W\geq1\) and test function \(f\),
\[
 \int\left(-\frac{{\mathscr L}W}{W}\right)f^2d\nu
 \leq\int|\nabla f|^2d\nu.                            \tag{73.16}
\]
\end{lemma}

\begin{proof}
Symmetry and the chain rule give
\[
\begin{split}
 \int\left(-\frac{{\mathscr L}W}{W}\right)f^2d\nu
 &=
 \int\nabla W\mathbin{\cdot}
 \nabla\left(\frac{f^2}{W}\right)d\nu\\
 &=
 \int\left[
 2f\nabla f\mathbin{\cdot}\nabla\log W
 -f^2|\nabla\log W|^2
 \right]d\nu\\
 &\leq\int|\nabla f|^2d\nu,
\end{split}
\]
using \(2a\cdot b-|b|^2\leq|a|^2\).
\end{proof}

\begin{proposition}[Lyapunov plus local Poincar\'e]
\label{prop:v15v73-poincare}
Suppose a set \(K\) has \(a=\nu(K)>0\), satisfies the local
Poincar\'e inequality
\[
 \int_K|f-f_K|^2d\nu
 \leq C_K\int_K|\nabla f|^2d\nu,                      \tag{73.17}
\]
where \(f_K=a^{-1}\int_Kf\,d\nu\), and there is \(W\geq1\) with
\[
 -\frac{{\mathscr L}W}{W}\geq c-b{\bf1}_K.            \tag{73.18}
\]
If
\[
 \gamma:=c-b\frac{1-a}{a}>0,                          \tag{73.19}
\]
then \(\nu\) satisfies
\[
 \operatorname{Var}_\nu(f)
 \leq
 \frac{1+bC_K}{\gamma}
 \int|\nabla f|^2d\nu.                                \tag{73.20}
\]
\end{proposition}

\begin{proof}
Take \(\int f\,d\nu=0\).  Lemma
\ref{lem:v15v73-lyapunov} and (73.18) give
\[
 c\int f^2d\nu
 \leq
 \int|\nabla f|^2d\nu+b\int_Kf^2d\nu.                 \tag{73.21}
\]
The local inequality and the identity
\(\int_Kf=-\int_{K^c}f\) imply
\[
 \int_Kf^2d\nu
 \leq
 C_K\int_K|\nabla f|^2d\nu
 \frac{1-a}{a}\int_{K^c}f^2d\nu.                      \tag{73.22}
\]
Insert (73.22) into (73.21), bound the last integral by
\(\int f^2\), and use (73.19).  This proves (73.20).
\end{proof}

Condition (73.19) can be arranged by enlarging a coercive core
when the tail mass decreases and the drift constant remains
positive.  The constants must be controlled uniformly in the
cutoff dimension.

\section{Tightening a defective logarithmic Sobolev estimate}

\begin{lemma}[Defective LSI plus spectral gap]
\label{lem:v15v73-tighten}
Assume
\[
 \operatorname{Ent}_\nu(f^2)
 \leq
 C_D\int|\nabla f|^2d\nu
 D_D\int f^2d\nu                                      \tag{73.23}
\]
and the Poincar\'e inequality
\[
 \operatorname{Var}_\nu(f)
 \leq C_P\int|\nabla f|^2d\nu.                        \tag{73.24}
\]
Then
\[
 \operatorname{Ent}_\nu(f^2)
 \leq
 \bigl[C_D+(D_D+2)C_P\bigr]
 \int|\nabla f|^2d\nu.                                \tag{73.25}
\]
\end{lemma}

\begin{proof}
Let \(g=f-\mathbb E_\nu f\).  The entropy shift inequality
\[
 \operatorname{Ent}_\nu(f^2)
 \leq
 \operatorname{Ent}_\nu(g^2)+2\mathbb E_\nu g^2       \tag{73.26}
\]
follows from convexity of \(s\log s\) and optimization over the
constant shift.  Apply (73.23) to \(g\), note that
\(\nabla g=\nabla f\), and use (73.24) on both
\(\int g^2\) terms.  This gives (73.25).
\end{proof}

\begin{theorem}[Local-convexity/Lyapunov package]
\label{thm:v15v73-package}
Suppose the regulated action has:
\begin{enumerate}
\item a cutoff-uniform defective estimate (73.23), obtained for
      example by local convexity and a controlled tail patch;
\item a cutoff-uniform local Poincar\'e inequality (73.17);
\item a cutoff-uniform Lyapunov function satisfying
      (73.18)--(73.19).
\end{enumerate}
Then the measures satisfy a cutoff-uniform LSI, with a constant
given explicitly by (73.20) and (73.25).
\end{theorem}

\begin{proof}
Proposition~\ref{prop:v15v73-poincare} supplies \(C_P\) uniformly.
Lemma~\ref{lem:v15v73-tighten} converts the defective inequality to
a tight one.
\end{proof}

This theorem deliberately lists the defective estimate as a
separate input.  Tail drift and local Poincar\'e control alone do
not prove logarithmic Sobolev concentration without a local
entropy estimate.

\section{Physical row tests}

\subsection{Higgs radial well}

For a real finite-dimensional Higgs coordinate \(H\),
\[
 W_H(H)=\lambda(|H|^2-v^2)^2,\qquad\lambda>0,
\]
has Hessian
\[
 \nabla^2W_H
 =
 8\lambda H\otimes H+
 4\lambda(|H|^2-v^2)I
 \succeq-4\lambda v^2I.                               \tag{73.27}
\]
Thus a positive quadratic/kinetic core with Hessian floor
\(m_0>4\lambda v^2\) absorbs the radial nonconvexity by the
form-small theorem of V72.  Gauge zero modes must first be removed
or fixed; otherwise no positive full-space floor exists.

\subsection{Determinant modulus}

The determinant contribution can enter the bounded perturbation
theorem only if its oscillation is uniform in the cutoff.  A
determinant zero sends
\(-\log|\det|\) to \(+\infty\), and growing rank can make even a
pointwise spectral corridor have cutoff-growing total oscillation.
The trace-ideal derivative and inverse estimates (72.27)--(72.28)
remain separate requirements.

\subsection{Conformal metric direction}

The bare Euclidean Einstein term is unbounded below along V66's
fixed-volume sequence, so it supplies neither (73.3) nor a
Lyapunov drift.  A stabilizer must first produce the coercive tail.
If the remaining low-mode wells are confined to a core, the
package in Theorem~\ref{thm:v15v73-package} becomes the appropriate
nonconvex route.

\begin{assessment}[Physical status]
\label{ass:v15v73-status}
V73 proves stability of LSI under bounded action oscillation,
shows that score derivatives still require independent control,
and gives a rigorous Lyapunov/local-Poincar\'e route from a
defective entropy estimate to a full LSI.  It also tests the Higgs,
determinant, and conformal rows against these hypotheses.

No cutoff-uniform decomposition, local entropy estimate, Lyapunov
function, or determinant corridor has yet been constructed for the
complete SM--Einstein modulus action.
\end{assessment}

\begin{decision}[V74 gauge--Higgs quotient Hessian]
\label{dec:v15v73-next}
Analyze the gauge--Higgs sector on a genuine quotient slice.
Separate gauge-orbit zero modes, derive the Hessian of the
gauge-covariant kinetic and Higgs potential terms, and formulate a
uniform transverse lower bound with the Faddeev--Popov determinant
and boundary faces visible.  Use a finite-dimensional counterexample
to show why positivity before quotienting is impossible.
\end{decision}

\begin{firewall}[Claim boundary after V73]
\label{fire:v15v73}
V73 proves perturbative and Lyapunov functional-inequality
compilers.  It does not establish their uniform hypotheses for the
physical regulator, control determinant zeros, or construct the
continuum theory.
\end{firewall}

\section*{Version V73 append-only record}
\addcontentsline{toc}{section}{Version V73 append-only record}
The complete V72 source was retained before this continuation.  It
had \(717226\) bytes, \(20865\) lines, and SHA--256
\[
 \texttt{7CFA46640D85FBB0C666A2DB411385EF4BB84F92B86BC86FBB01F0179A80456E}.
\]
V73 proved bounded-oscillation stability of logarithmic Sobolev
inequalities, separated action and score perturbations, developed
a Lyapunov/local-entropy tightening package, and tested the Higgs,
determinant, and conformal rows against it.

\chapter{Fifteenth-cycle continuation V74: the gauge--Higgs quotient Hessian}
\label{sec:v15v74}

\section{Gauge invariance forces unreduced zero modes}

Let a compact finite-cutoff gauge group \({\cal G}_n\) act smoothly
and isometrically on a Euclidean configuration space \(X_n\).
For \(\xi\) in its Lie algebra, write
\(\xi_X(x)\) for the infinitesimal orbit vector.  Let
\({\cal A}_n:X_n\to{\mathbb R}\) be gauge invariant.

\begin{theorem}[Orbit kernel at a critical configuration]
\label{thm:v15v74-orbit-kernel}
If \(x_*\) is a critical point of \({\cal A}_n\), then
\[
 \nabla^2{\cal A}_n(x_*)[h,\xi_X(x_*)]=0              \tag{74.1}
\]
for every \(h\in T_{x_*}X_n\) and every \(\xi\).  In particular, a
nontrivial gauge orbit prevents
\(\nabla^2{\cal A}_n(x_*)\succeq mI\) with \(m>0\) on the
unreduced space.
\end{theorem}

\begin{proof}
Gauge invariance gives
\[
 D{\cal A}_n(x)[\xi_X(x)]=0                            \tag{74.2}
\]
for every \(x\).  Differentiate (74.2) in the direction \(h\):
\[
 \nabla^2{\cal A}_n(x)[h,\xi_X(x)]
 +D{\cal A}_n(x)[D_h\xi_X]=0.                         \tag{74.3}
\]
At \(x=x_*\), the second term vanishes because
\(D{\cal A}_n(x_*)=0\), proving (74.1).
\end{proof}

\begin{counterexample}[Mexican-hat orbit]
\label{cnt:v15v74-circle}
On \({\mathbb R}^2\), let \(SO(2)\) act by rotations and set
\[
 W(x)=\lambda(|x|^2-v^2)^2.
\]
Every point with \(|x|=v\) is critical, and
\[
 \nabla^2W(x)=8\lambda x\otimes x.                    \tag{74.4}
\]
The tangential orbit direction is in the kernel.  Quotienting the
circle of minima leaves the radial Hessian, but no positive
full-space Hessian exists before reduction.
\end{counterexample}

\section{Local slice and Faddeev--Popov map}

Let a gauge condition
\[
 {\cal F}_n(x)=0                                      \tag{74.5}
\]
define a local slice \({\cal S}_n\).  Its Faddeev--Popov map is
\[
 {\cal M}_{\rm FP}(x)\xi
 =
 D{\cal F}_n(x)[\xi_X(x)].                             \tag{74.6}
\]

\begin{proposition}[Uniform transverse chart]
\label{prop:v15v74-slice}
Suppose on a slice region \(U_n\subset{\cal S}_n\),
\[
 \|{\cal M}_{\rm FP}(x)\xi\|
 \geq\sigma\|\xi\|,\qquad\sigma>0,                    \tag{74.7}
\]
with \(\sigma\) independent of \(n,x\), and the infinitesimal action
and slice projection have uniform operator bounds.  Then the map
\[
 {\cal G}_n\times U_n\longrightarrow X_n,\qquad
 (\gamma,s)\longmapsto\gamma s                        \tag{74.8}
\]
is a uniformly nondegenerate local chart modulo the stabilizer.
Gauge-invariant observables and actions descend to \(U_n\), and
their physical Hessian is the Hessian in transverse directions.
\end{proposition}

\begin{proof}
The derivative of (74.8) in the orbit direction is
\(\xi_X(s)\), while applying \(D{\cal F}_n\) gives
\({\cal M}_{\rm FP}(s)\xi\).  Bound (74.7) provides a uniformly
bounded inverse on the orbit component.  The inverse-function
theorem, together with the assumed projection bounds, gives the
local product chart.  Gauge invariance makes the descended
functions constant in the first factor, so only the slice Hessian
acts on quotient directions.
\end{proof}

At a stabilizer jump, the quotient is an orbifold or a more
singular stratified space.  Proposition~\ref{prop:v15v74-slice}
applies only on one regular orbit-type stratum.

\section{Exact gauge--Higgs second variation}

At one finite regulator, use the schematic but convention-fixed
action
\[
 {\cal A}_{GH}(A,H)
 =
 \frac12\|F_A\|^2+
 \frac12\|D_AH\|^2+
 \lambda\int(|H|^2-v^2)^2,\qquad\lambda>0.             \tag{74.9}
\]
Let \((a,h)\) be a linear variation.  Define
\[
 B_A(a,a)=
 \left.\frac{d^2}{dt^2}\right|_{t=0}F_{A+ta}.          \tag{74.10}
\]
For the usual connection convention,
\(B_A(a,a)=2a\wedge a\).

\begin{theorem}[Gauge--Higgs Hessian formula]
\label{thm:v15v74-gh-hessian}
The second variation of (74.9) is
\[
\begin{split}
 D^2{\cal A}_{GH}(A,H)[(a,h)]^2
 &=
 \|d_Aa\|^2+\langle F_A,B_A(a,a)\rangle\\
 &\quad+
 \|D_Ah+\rho(a)H\|^2
 +2\langle D_AH,\rho(a)h\rangle\\
 &\quad+
 8\lambda\int\langle H,h\rangle^2
 +4\lambda\int(|H|^2-v^2)|h|^2.                      \tag{74.11}
\end{split}
\]
\end{theorem}

\begin{proof}
For a differentiable vector-valued map \(Y(t)\),
\[
 \left.\frac{d^2}{dt^2}\right|_0\frac12\|Y(t)\|^2
 =
 \|Y'(0)\|^2+\langle Y(0),Y''(0)\rangle.               \tag{74.12}
\]
Apply this first to
\[
 F_{A+ta},
\]
whose first and second derivatives are \(d_Aa\) and
\(B_A(a,a)\).  Next,
\[
 D_{A+ta}(H+th)
 =
 D_AH+t(D_Ah+\rho(a)H)+t^2\rho(a)h,                  \tag{74.13}
\]
so its second derivative is \(2\rho(a)h\).  Formula (74.12) gives
the second line of (74.11).  Finally, differentiating
\(\lambda(|H|^2-v^2)^2\) twice gives the last line.
\end{proof}

\section{A transverse lower-bound compiler}

Let \(T_{(A,H)}{\cal S}_n\) be the chosen transverse slice.  Suppose
the positive part in (74.11) has a slice floor
\[
\begin{split}
 &\|d_Aa\|^2+\|D_Ah+\rho(a)H\|^2
 +8\lambda\int\langle H,h\rangle^2\\
 &\hspace{30mm}\geq
 \gamma\bigl(\|a\|^2+\|h\|^2\bigr)                   \tag{74.14}
\end{split}
\]
for \((a,h)\in T_{(A,H)}{\cal S}_n\).

\begin{proposition}[Background-error estimate]
\label{prop:v15v74-error}
Assume regulator-uniform multiplication bounds give
\[
 |\langle F_A,B_A(a,a)\rangle|
 \leq C_F\|F_A\|_\infty\|a\|^2,                       \tag{74.15}
\]
\[
 2|\langle D_AH,\rho(a)h\rangle|
 \leq
 C_H\|D_AH\|_\infty
 \bigl(\|a\|^2+\|h\|^2\bigr).                         \tag{74.16}
\]
Then on the slice
\[
\begin{split}
 D^2{\cal A}_{GH}[(a,h)]^2
 \geq
 \bigl[
 \gamma-C_F\|F_A\|_\infty-C_H\|D_AH\|_\infty
 -4\lambda v^2
 \bigr]
 (\|a\|^2+\|h\|^2).                                   \tag{74.17}
\end{split}
\]
\end{proposition}

\begin{proof}
The final term in (74.11) obeys
\[
 4\lambda(|H|^2-v^2)|h|^2
 \geq-4\lambda v^2|h|^2.                              \tag{74.18}
\]
Combine (74.14)--(74.16) and (74.18), enlarging the negative
coefficient to multiply the full norm.
\end{proof}

\begin{corollary}[Uniform quotient Hessian corridor]
\label{cor:v15v74-corridor}
If the bracket in (74.17) is bounded below by \(m_{GH}>0\)
throughout the slice region, then the gauge--Higgs action is
uniformly strongly convex on physical transverse directions there.
\end{corollary}

\begin{proof}
Equation (74.17) is exactly the quadratic-form definition of the
claim.
\end{proof}

The floor (74.14) is a covariant Poincar\'e/spectral statement.  A
finite Wilson-pencil certificate may verify a compact portion of
it only after a global reduction covers all momenta and all cutoff
scales.  The active YM V39 plate does not yet do so.

\section{The quotient density changes the Hessian}

On a regular gauge slice, the positive modulus density contains
\[
 J_{\rm FP}(x)=|\det{\cal M}_{\rm FP}(x)|.             \tag{74.19}
\]
Thus the reduced action is
\[
 {\cal A}_{\rm red}
 =
 {\cal A}_{GH}
 -\log J_{\rm FP}
 +{\cal A}_{\rm other}.                               \tag{74.20}
\]

\begin{proposition}[Faddeev--Popov Hessian]
\label{prop:v15v74-fp-hessian}
For a slice variation \(z\), if
\({\cal M}_{\rm FP}\) is invertible and differentiable, then
\[
\begin{split}
 D_z^2\log\det{\cal M}_{\rm FP}
 &=
 \operatorname{Tr}\bigl(
 {\cal M}_{\rm FP}^{-1}D_z^2{\cal M}_{\rm FP}
 \bigr)\\
 &\quad-
 \operatorname{Tr}\bigl(
 {\cal M}_{\rm FP}^{-1}D_z{\cal M}_{\rm FP}
 {\cal M}_{\rm FP}^{-1}D_z{\cal M}_{\rm FP}
 \bigr).                                              \tag{74.21}
\end{split}
\]
If
\[
 \|{\cal M}_{\rm FP}^{-1}\|\leq\sigma^{-1},           \tag{74.22}
\]
then
\[
 |D_z^2\log\det{\cal M}_{\rm FP}|
 \leq
 \sigma^{-1}\|D_z^2{\cal M}_{\rm FP}\|_1+
 \sigma^{-2}\|D_z{\cal M}_{\rm FP}\|_2^2.             \tag{74.23}
\]
\end{proposition}

\begin{proof}
Differentiate the logarithmic determinant twice, exactly as in
(72.26).  The trace-ideal inequalities used in (72.28) give
(74.23).
\end{proof}

The quotient contribution must be subtracted according to
(74.20).  Therefore a positive gauge--Higgs Hessian before the
Jacobian does not prove a positive reduced Hessian.

\section{Gribov faces and paired flux}

At a Gribov horizon, the smallest singular value in (74.7) tends
to zero.  Both the chart and (74.23) lose uniformity.  If a
fundamental modular region has gauge-equivalent boundary faces,
Theorem~\ref{thm:v15v64-face-pair} cancels their flux only when:
\begin{enumerate}
\item the gauge-invariant observable has matching traces;
\item the complete density, including \(J_{\rm FP}\), has matching
      oriented traces;
\item the metric variation transports equivariantly across the
      face;
\item the unbounded score insertion has an integrable face trace.
\end{enumerate}
Pointwise vanishing of \(J_{\rm FP}\) is insufficient if
\(D\log J_{\rm FP}\) diverges at the reciprocal rate.

\begin{proposition}[Simple horizon cancellation test]
\label{prop:v15v74-horizon}
Let \(s\geq0\) be a normal coordinate at a regular horizon.  If
\[
 J_{\rm FP}(s)=s^\alpha j(s),\quad j(0)>0,             \tag{74.25}
\]
and the remaining density and normal field are bounded, then the
undifferentiated flux vanishes for \(\alpha>0\).  The score-weighted
flux containing
\[
 D_s\log J_{\rm FP}=\frac{\alpha}{s}+O(1)             \tag{74.26}
\]
vanishes only when \(\alpha>1\), absent additional cancellation.
\end{proposition}

\begin{proof}
The ordinary flux density is \(O(s^\alpha)\).  Multiplication by
the logarithmic derivative in (74.26) makes the score flux
\(O(s^{\alpha-1})\).  Take the trace \(s\downarrow0\).
\end{proof}

\section{Complete transverse LSI ledger}

A gauge--Higgs LSI through the V72--V73 compilers requires:
\begin{enumerate}
\item a regular slice with the cofinal singular-value bound (74.7);
\item a physical spectral floor (74.14);
\item a positive margin in (74.17);
\item trace-ideal control (74.23) for the quotient density;
\item either a bounded-oscillation decomposition or a
      local-entropy/Lyapunov package for remaining interactions;
\item paired or vanishing score flux at all slice boundaries.
\end{enumerate}
These conditions are regulator-uniform statements.  Verification
at one finite lattice does not establish the cofinal row.

\begin{assessment}[Physical status]
\label{ass:v15v74-status}
V74 proves the gauge-orbit kernel theorem, derives the exact
gauge--Higgs Hessian, supplies a transverse lower-bound compiler,
and incorporates the Faddeev--Popov determinant and Gribov
score-flux singularity.

The Standard Model regulator has not yet been shown to admit a
global cofinal slice with (74.7), a uniform spectral floor
(74.14), a positive error margin, or integrable horizon score
flux.  The quotient LSI is therefore not yet populated.
\end{assessment}

\begin{decision}[V75 cross-program audit]
\label{dec:v15v74-next}
Perform the scheduled companion audit.  Inspect the newest
Yang--Mills note specifically for additional momentum coverage or
a cofinal transverse floor relevant to (74.14), and the newest
Navier--Stokes note for a quotient or boundary estimate relevant
to (74.24)--(74.26).  Record the result, then continue in V76 with
a cofinal slice-atlas gluing theorem.
\end{decision}

\begin{firewall}[Claim boundary after V74]
\label{fire:v15v74}
V74 proves finite-cutoff quotient-Hessian implications.  It does
not construct a global gauge slice, prove cofinal Faddeev--Popov
invertibility, cancel Gribov flux, or establish the physical
continuum.
\end{firewall}

\section*{Version V74 append-only record}
\addcontentsline{toc}{section}{Version V74 append-only record}
The complete V73 source was retained before this continuation.  It
had \(728709\) bytes, \(21258\) lines, and SHA--256
\[
 \texttt{805A100F219D75330BBD7A26C128074A37D87E9C5C6DFE77F978510AF2C8D778}.
\]
V74 proved the unreduced gauge-orbit Hessian kernel, derived the
gauge--Higgs transverse second variation, and added the
Faddeev--Popov Hessian, Gribov boundary, and score-flux rows needed
for a quotient logarithmic Sobolev estimate.


\chapter{Fifteenth-cycle continuation V75: scheduled quotient-floor audit}
\label{sec:v15v75}

\section{Companion identities}

The active companion files at this checkpoint are
\[
\begin{array}{c|c|c|c|l}
\hbox{programme}&\hbox{version}&\hbox{bytes}&\hbox{lines}&\hbox{SHA--256}\\ \hline
\hbox{Yang--Mills}&V39&329898&9076&
\texttt{9B44B0FBCA6037F3319E86036753B0F3B76BDD06F6045F00747E05DC01E23F9D}\\
\hbox{Navier--Stokes}&V26&278283&5319&
\texttt{82C887F8C2C69E4F28FAD7C3DC8DE5B9099CB5A4BF431EBA1FFF3E91935978AD}.
\end{array}                                             \tag{75.1}
\]
The hashes agree with V70, so neither companion acquired a new
result during V71--V74.

\section{Transverse-floor comparison}

YM V39 proves exact Wilson-pencil positivity on
\[
 [0,\pi/4]^2\times[0,t]                               \tag{75.2}
\]
and announces a proposed third-layer extension.  The proposal is
not yet part of the active theorem.  Its firewall still excludes a
full transverse cell, global torus propagation, cofinal response,
and an interacting reflection-positive measure.

\begin{proposition}[Status of the V74 spectral input]
\label{prop:v15v75-ym}
The active YM theorem does not prove the transverse slice floor
(74.14) for the SM--Einstein cutoff family.
\end{proposition}

\begin{proof}
Equation (74.14) concerns the complete gauge--Higgs covariant
Hessian on every physical slice direction and every cofinal
cutoff.  The YM result concerns one finite Wilson matrix pencil on
the proper plate (75.2).  The missing momentum region, Higgs
mixing, Faddeev--Popov row, and cutoff-uniform identification are
all outside its theorem.
\end{proof}

\section{Boundary comparison}

NS V26 adds pointwise rank-one norms of Oseen-kernel derivatives.
Its earlier Morrey ledgers concern a parabolic velocity equation,
not gauge slices.

\begin{proposition}[Status of the V74 boundary input]
\label{prop:v15v75-ns}
The active NS theorem supplies no face pairing, Faddeev--Popov
singular-value estimate, or score-weighted Gribov flux bound for
(74.24)--(74.26).
\end{proposition}

\begin{proof}
No gauge group, slice map, quotient density, Gribov horizon, or
metric score occurs in the terminal NS theorem.  Its tensor-kernel
norms therefore have no defined map to the boundary quantities in
V74.
\end{proof}

\section{Direction after the audit}

\begin{theorem}[Audit conclusion]
\label{thm:v15v75-conclusion}
No cross-program result presently populates the local transverse
Hessian or quotient-boundary hypotheses of V74.  The next valid
step is an abstract cofinal gluing theorem that states exactly what
additional overlap and between-chart estimates would convert local
slice results into a global quotient functional inequality.
\end{theorem}

\begin{proof}
The source identities are fixed by (75.1).
Propositions~\ref{prop:v15v75-ym} and
\ref{prop:v15v75-ns} eliminate the two targeted transfers.
\end{proof}

\begin{assessment}[Physical status]
\label{ass:v15v75-status}
V75 is an audit.  It prevents the local Wilson-pencil plate from
being promoted to a global gauge--Higgs quotient floor and adds no
continuum existence claim.
\end{assessment}

\begin{decision}[V76 cofinal slice-atlas gluing]
\label{dec:v15v75-next}
Construct a quotient atlas with exact density cocycles, decompose
entropy into within-chart and between-chart contributions, and
prove a global LSI from local slice LSIs, a discrete chart-graph
LSI, and quantitative overlap bridges.  Track all constants
cofinally and exhibit a separated-mixture counterexample showing
that local chart inequalities alone do not glue.
\end{decision}

\begin{firewall}[Claim boundary after V75]
\label{fire:v15v75}
V75 proves no new spectral floor, quotient atlas, functional
inequality, or continuum measure.
\end{firewall}

\section*{Version V75 append-only record}
\addcontentsline{toc}{section}{Version V75 append-only record}
The complete V74 source was retained before this continuation.  It
had \(740676\) bytes, \(21637\) lines, and SHA--256
\[
 \texttt{5AC3D6B40026BD61B19AF6E08862F2FA2D74B022868430F3FAF26E8B938E580C}.
\]
V75 performed the scheduled companion audit, recorded the
unchanged source hashes, and retained quotient-atlas gluing as the
next noncircular task.

\chapter{Fifteenth-cycle continuation V76: cofinal gluing of a gauge-slice atlas}
\label{sec:v15v76}

\section{Density cocycles on the quotient}

Let \(Q_n=X_n/{\cal G}_n\) be the regular part of a finite-cutoff
quotient.  Suppose it is covered by slice charts
\[
 \kappa_i:U_i\longrightarrow V_i,\qquad i\in I_n.     \tag{76.1}
\]
In chart \(i\), write the positive modulus density as
\[
 d\nu_{n,i}=Z_n^{-1}w_{n,i}(x)\,dx,                   \tag{76.2}
\]
where \(w_{n,i}\) includes the action, Faddeev--Popov factor, metric
Jacobian, and every positive reference-density row.

\begin{definition}[Exact density cocycle]
\label{def:v15v76-cocycle}
On an overlap, let
\[
 \tau_{ij}=\kappa_j\circ\kappa_i^{-1}.
\]
The local densities form an exact cocycle when
\[
 w_{n,i}(x)
 =
 w_{n,j}(\tau_{ij}x)
 |\det D\tau_{ij}(x)|                                 \tag{76.3}
\]
almost everywhere and the transition maps satisfy the usual
triple-overlap identities.
\end{definition}

\begin{proposition}[Global quotient measure]
\label{prop:v15v76-global-measure}
An exact density cocycle defines a unique measure \(\nu_n\) on the
covered regular quotient whose coordinate expression is (76.2).
\end{proposition}

\begin{proof}
Equation (76.3) is precisely the change-of-variables formula on
overlaps.  Hence the integrals of a partition-of-unity
decomposition are independent of the selected chart.  The
triple-overlap identity gives consistency on threefold
intersections, and uniqueness follows because the charts cover.
\end{proof}

\begin{proposition}[Score cocycle]
\label{prop:v15v76-score-cocycle}
Let \(B_{n,i}\) be local vector fields representing one quotient
variation and assume
\[
 D\tau_{ij}B_{n,i}=B_{n,j}\circ\tau_{ij}.              \tag{76.4}
\]
Define the local geometric score as minus the logarithmic Lie
derivative of the density:
\[
 {\cal S}_{n,i}(B)
 =
 -w_{n,i}^{-1}\operatorname{div}(w_{n,i}B_{n,i}).      \tag{76.5}
\]
Then
\[
 {\cal S}_{n,i}(B)
 =
 {\cal S}_{n,j}(B)\circ\tau_{ij}.                     \tag{76.6}
\]
\end{proposition}

\begin{proof}
The density \(w_i\,dx\) is the pullback of \(w_j\,dy\) by (76.3),
and the vector fields are related by (76.4).  Lie differentiation
commutes with pullback.  Dividing the pulled-back Lie derivative by
the pulled-back density gives (76.6).
\end{proof}

Formula (76.5) agrees with (64.2):
\[
 -w^{-1}\operatorname{div}(wB)
 =
 D_B{\cal A}-D_B\log\rho-\operatorname{div}B.
\]

\section{Exact entropy decomposition}

Choose a measurable ownership partition
\[
 Q_n=\bigsqcup_{i\in I_n}C_{n,i},\qquad C_{n,i}\subset U_i,      \tag{76.7}
\]
up to a null set.  Put
\[
 p_{n,i}=\nu_n(C_{n,i})>0,\qquad
 \nu_{n,i}=\nu_n(\,\cdot\,|C_{n,i}).                  \tag{76.8}
\]
For a test function \(f\), define
\[
 m_i=\mathbb E_{\nu_{n,i}}f^2.                        \tag{76.9}
\]

\begin{lemma}[Within-chart/between-chart decomposition]
\label{lem:v15v76-entropy-split}
One has the exact identity
\[
\begin{split}
 \operatorname{Ent}_{\nu_n}(f^2)
 &=
 \sum_i p_{n,i}
 \operatorname{Ent}_{\nu_{n,i}}(f^2)
 +\operatorname{Ent}_{p_n}(m),                        \tag{76.10}\\
 \operatorname{Ent}_{p_n}(m)
 &:=
 \sum_i p_{n,i}m_i
 \log\frac{m_i}{\sum_jp_{n,j}m_j}.
\end{split}
\]
\end{lemma}

\begin{proof}
On cell \(i\), insert
\[
 \log\frac{f^2}{\mathbb E_{\nu_n}f^2}
 =
 \log\frac{f^2}{m_i}
 +\log\frac{m_i}{\sum_jp_{n,j}m_j}.
\]
Integrate with \(\nu_{n,i}\), multiply by \(p_{n,i}\), and sum.
\end{proof}

The second term in (76.10) is invisible to every purely local
chart estimate.

\section{Discrete graph and bridge energies}

Let \(G_n=(I_n,E_n,c^{(n)})\) be a connected weighted graph with
symmetric conductances \(c_{ij}^{(n)}\geq0\).  Define
\[
 {\cal E}_{G_n}(a)
 =
 \frac12\sum_{i,j}c_{ij}^{(n)}(a_i-a_j)^2.             \tag{76.11}
\]

\begin{definition}[Discrete chart LSI]
\label{def:v15v76-graph-lsi}
The chart graph satisfies \({\rm LSI}(C_G)\) when
\[
 \operatorname{Ent}_{p_n}(a^2)
 \leq2C_G{\cal E}_{G_n}(a)                            \tag{76.12}
\]
for every real vector \(a=(a_i)\).
\end{definition}

\begin{definition}[Overlap bridge]
\label{def:v15v76-bridge}
The ownership cells have bridge constant \(C_B\) when
\[
 {\cal E}_{G_n}
 \left(
 \bigl(\sqrt{\mathbb E_{\nu_{n,i}}f^2}\bigr)_i
 \right)
 \leq
 C_B\int_{Q_n}|\nabla_Qf|^2d\nu_n                    \tag{76.13}
\]
for every smooth quotient function \(f\).
\end{definition}

\begin{proposition}[A coupling sufficient condition for bridges]
\label{prop:v15v76-coupling}
For each edge \(ij\), suppose there is a coupling
\(\pi_{ij}\) of \(\nu_{n,i}\) and \(\nu_{n,j}\) such that
\[
 \int|f(x)-f(y)|^2d\pi_{ij}(x,y)
 \leq b_{ij}{\cal D}_{ij}(f),                         \tag{76.14}
\]
where the local Dirichlet forms obey the congestion bound
\[
 \frac12\sum_{i,j}c_{ij}^{(n)}b_{ij}{\cal D}_{ij}(f)
 \leq C_B\int|\nabla_Qf|^2d\nu_n.                    \tag{76.15}
\]
Then (76.13) holds.
\end{proposition}

\begin{proof}
The reverse triangle inequality in
\(L^2(\pi_{ij})\) gives
\[
\begin{split}
 \left|
 \sqrt{\mathbb E_{\nu_{n,i}}f^2}
 -\sqrt{\mathbb E_{\nu_{n,j}}f^2}
 \right|
 &=
 \left|\|f(x)\|_{L^2(\pi_{ij})}
 -\|f(y)\|_{L^2(\pi_{ij})}\right|\\
 &\leq\|f(x)-f(y)\|_{L^2(\pi_{ij})}.
\end{split}
\]
Square, use (76.14), multiply by \(c_{ij}^{(n)}/2\), sum, and
apply (76.15).
\end{proof}

Overlap corridors and gauge-flow paths are intended to construct
the couplings in (76.14).  Their length, density ratio, and
Dirichlet congestion must all be bounded cofinally.

\section{Global quotient logarithmic Sobolev theorem}

\begin{theorem}[Slice-atlas gluing]
\label{thm:v15v76-gluing}
Assume:
\[
 \operatorname{Ent}_{\nu_{n,i}}(f^2)
 \leq2C_0\mathbb E_{\nu_{n,i}}|\nabla_Qf|^2           \tag{76.16}
\]
on every ownership cell, the graph inequality (76.12) holds with
\(C_G\), and the bridge inequality (76.13) holds with \(C_B\).
Then
\[
 \operatorname{Ent}_{\nu_n}(f^2)
 \leq
 2(C_0+C_GC_B)
 \int|\nabla_Qf|^2d\nu_n.                             \tag{76.17}
\]
\end{theorem}

\begin{proof}
The first term on the right of (76.10) is at most
\[
 2C_0\sum_i p_{n,i}
 \mathbb E_{\nu_{n,i}}|\nabla_Qf|^2
 =
 2C_0\int|\nabla_Qf|^2d\nu_n.                         \tag{76.18}
\]
Apply (76.12) to \(a_i=\sqrt{m_i}\) and then (76.13):
\[
 \operatorname{Ent}_{p_n}(m)
 \leq
 2C_GC_B\int|\nabla_Qf|^2d\nu_n.                     \tag{76.19}
\]
Adding (76.18)--(76.19) proves (76.17).
\end{proof}

\begin{corollary}[Cofinal quotient score moment]
\label{cor:v15v76-cofinal}
If \(C_0,C_G,C_B\) are independent of \(n\), the quotient measures
satisfy a uniform LSI.  If a compact physical metric variation
\(k\) has quotient score satisfying
\[
 |\nabla_Q{\cal S}_{n}(k)|\leq L_k                    \tag{76.20}
\]
uniformly and all true boundary defects vanish, then
\[
 \sup_n\mathbb E_{\nu_n}e^{\tau|{\cal S}_n(k)|}
 \leq
 2\exp\!\left[
 \frac12(C_0+C_GC_B)L_k^2\tau^2
 \right].                                             \tag{76.21}
\]
\end{corollary}

\begin{proof}
Use Theorem~\ref{thm:v15v76-gluing}, the score cocycle
(76.6), and Theorem~\ref{thm:v15v72-herbst}.
\end{proof}

\section{Coordinate distortion and local constants}

\begin{proposition}[Chart distortion]
\label{prop:v15v76-distortion}
Suppose a local chart measure satisfies an LSI with coordinate
gradient constant \(C_{\rm coord}\), and
\[
 |\nabla_{\rm coord}f|^2
 \leq\kappa^2|\nabla_Qf|^2.                           \tag{76.22}
\]
Then its quotient-gradient LSI constant is at most
\(\kappa^2C_{\rm coord}\).
\end{proposition}

\begin{proof}
Insert (76.22) into the coordinate LSI.
\end{proof}

Thus the inverse Faddeev--Popov and slice-projection bounds in V74
enter \(C_0\) through \(\kappa\), even before the determinant enters
the density.

\section{Artificial and genuine boundaries}

The ownership boundaries in (76.7) are bookkeeping devices.
Because the density and score cocycles agree, their two oriented
contributions cancel when the global divergence identity is
reassembled.  A Gribov horizon, singular orbit-type stratum, or
unpaired edge of the covered quotient is a genuine boundary and
still requires V64's trace conditions.

\begin{proposition}[Ownership-face cancellation]
\label{prop:v15v76-ownership}
Let a regular interface between cells \(C_i,C_j\) lie inside a
chart overlap.  If the density cocycle (76.3), vector transport
(76.4), and observable trace matching hold, then its two boundary
fluxes cancel.
\end{proposition}

\begin{proof}
The outward normals of the two ownership cells have opposite
orientations.  Equation (76.3) identifies the induced density,
(76.4) identifies the normal vector flux, and the observable is one
global quotient function.  The two surface integrals are therefore
equal with opposite signs, as in
Theorem~\ref{thm:v15v64-face-pair}.
\end{proof}

\section{Why local slice LSIs are insufficient}

\begin{counterexample}[Separated chart mixture]
\label{cnt:v15v76-mixture}
Let
\[
 \nu_R=\frac12{\cal N}(-R,1)+\frac12{\cal N}(R,1)     \tag{76.23}
\]
on \(\mathbb R\).  Each conditional Gaussian component satisfies
\({\rm LSI}(1)\), uniformly in \(R\).  Nevertheless the global
Poincar\'e constant, and hence the global LSI constant, is at least
\[
 R^2+1.                                               \tag{76.24}
\]
\end{counterexample}

\begin{proof}
For \(f(x)=x\),
\[
 \operatorname{Var}_{\nu_R}(f)=R^2+1,\qquad
 \int|f'|^2d\nu_R=1.
\]
Therefore any Poincar\'e constant is at least \(R^2+1\).
Linearizing an LSI at \(f=1+\epsilon x\) gives the Poincar\'e
inequality with the same constant, proving the claim.
\end{proof}

The diverging constant lies in the between-chart term of (76.10).
A growing number of individually good charts can likewise fail if
the discrete graph LSI or bridge congestion degenerates.

\section{Cofinal atlas ledger}

The quotient atlas closes the score-moment row only after proving:
\begin{enumerate}
\item exact density and score cocycles;
\item local quotient LSIs with uniform coordinate distortion;
\item a chart ownership graph with uniform discrete LSI;
\item overlap or gauge-flow couplings with uniform congestion;
\item cancellation of artificial interfaces and zero flux at
      every genuine quotient boundary;
\item a uniform quotient score-gradient bound.
\end{enumerate}

\begin{assessment}[Physical status]
\label{ass:v15v76-status}
V76 proves an exact local/global entropy decomposition and a
cofinal quotient LSI gluing theorem.  It identifies chart-graph
mixing and overlap transport as independent of local slice
convexity, and propagates the result to exponential metric-score
moments.

No physical gauge-slice atlas has yet been shown to satisfy the
uniform density cocycle, graph LSI, coupling congestion, or genuine
boundary conditions.  The theorem is a complete gluing compiler,
not a construction of those inputs.
\end{assessment}

\begin{decision}[V77 overlap bridges from gauge flows]
\label{dec:v15v76-next}
Construct the bridge estimate (76.14) from explicit paths through
slice overlaps.  Prove an \(L^2\) path-coupling bound using
uniformly bounded gauge-flow length, Jacobian density ratios, and
path congestion.  Show how a shrinking Faddeev--Popov singular
value makes the bridge constant diverge.
\end{decision}

\begin{firewall}[Claim boundary after V76]
\label{fire:v15v76}
V76 proves conditional atlas gluing.  It does not supply a cofinal
atlas, local transverse floors, graph mixing, overlap couplings,
Gribov cancellation, or the full continuum measure.
\end{firewall}

\section*{Version V76 append-only record}
\addcontentsline{toc}{section}{Version V76 append-only record}
The complete V75 source was retained before this continuation.  It
had \(744910\) bytes, \(21752\) lines, and SHA--256
\[
 \texttt{DFF3B738E25A1A011D1603AC0E39B3034A7CFD16531C94F2F07E0D3EB42437E3}.
\]
V76 constructed exact density and score cocycles, split entropy
into local and chart-graph components, and proved a cofinal
slice-atlas LSI from local inequalities, a discrete graph LSI, and
quantitative overlap bridges.

\chapter{Fifteenth-cycle continuation V77: overlap bridges from controlled slice paths}
\label{sec:v15v77}

\section{Path couplings}

Let \((Q,d,\nu)\) be a finite-dimensional regular quotient with
Dirichlet form
\[
 {\cal E}(f)=\int_Q|\nabla_Qf|^2d\nu.                 \tag{77.1}
\]
Let \(\nu_i,\nu_j\) be two ownership-cell conditional measures.
A path coupling consists of a coupling \(\pi_{ij}\) and, for
\(\pi_{ij}\)-almost every \((x,y)\), an absolutely continuous path
\[
 \gamma_{xy}:[0,1]\to Q,\qquad
 \gamma_{xy}(0)=x,\quad\gamma_{xy}(1)=y.              \tag{77.2}
\]
Write
\[
 \ell(\gamma_{xy})=\int_0^1|\dot\gamma_{xy}(t)|dt.    \tag{77.3}
\]

\begin{definition}[Path occupation measure]
\label{def:v15v77-occupation}
The occupation measure of the coupling is
\[
 \Omega_{ij}(E)
 =
 \int\!\!\int_0^1
 {\bf1}_E(\gamma_{xy}(t))
 |\dot\gamma_{xy}(t)|\,dt\,d\pi_{ij}(x,y).            \tag{77.4}
\]
\end{definition}

\begin{theorem}[Path-coupling bridge]
\label{thm:v15v77-path}
Assume
\[
 \ell(\gamma_{xy})\leq L_{ij}                         \tag{77.5}
\]
almost surely and
\[
 \Omega_{ij}\leq K_{ij}\nu                            \tag{77.6}
\]
as measures.  Then every smooth \(f\) satisfies
\[
 \int|f(x)-f(y)|^2d\pi_{ij}(x,y)
 \leq
 L_{ij}K_{ij}{\cal E}(f).                             \tag{77.7}
\]
\end{theorem}

\begin{proof}
The fundamental theorem along each path and Cauchy's inequality
with arc-length measure give
\[
\begin{split}
 |f(x)-f(y)|^2
 &\leq
 \left(\int_0^1
 |\nabla_Qf(\gamma)|\,|\dot\gamma|dt\right)^2\\
 &\leq
 \ell(\gamma)
 \int_0^1|\nabla_Qf(\gamma)|^2|\dot\gamma|dt.
\end{split}
\]
Use (77.5), integrate over \(\pi_{ij}\), recognize (77.4), and
apply (77.6).
\end{proof}

\begin{corollary}[Bridge constant from aggregate congestion]
\label{cor:v15v77-aggregate}
For graph edges \(ij\), suppose path couplings satisfy (77.5) and
the weighted occupation estimate
\[
 \frac12\sum_{i,j}c_{ij}^{(n)}L_{ij}\Omega_{ij}
 \leq C_B\nu_n.                                      \tag{77.8}
\]
Then the atlas bridge inequality (76.13) holds with constant
\(C_B\).
\end{corollary}

\begin{proof}
For each coupling, the reverse \(L^2\) triangle inequality gives
\[
 \left|
 \sqrt{\mathbb E_{\nu_i}f^2}
 -\sqrt{\mathbb E_{\nu_j}f^2}
 \right|^2
 \leq
 \int|f(x)-f(y)|^2d\pi_{ij}.                          \tag{77.9}
\]
Use the pointwise path estimate from the proof of
Theorem~\ref{thm:v15v77-path}, multiply by
\(c_{ij}^{(n)}/2\), sum, and apply (77.8).
\end{proof}

The aggregate form (77.8) is sharper than summing independent
global bounds: it permits many paths provided they do not overload
the same quotient region.

\section{Flow transports and density distortion}

Suppose a transport is generated by maps
\[
 T_t:C_i\to Q,\qquad T_0=\operatorname{id},\quad
 T_1{}_\#\nu_i=\nu_j,                                 \tag{77.10}
\]
with velocity field \(Z_t\).  The coupling is
\((\operatorname{id},T_1)_\#\nu_i\), and its paths are
\(\gamma_x(t)=T_t(x)\).

\begin{proposition}[Flow occupation bound]
\label{prop:v15v77-flow}
If
\[
 |Z_t(T_t x)|\leq V_{ij},                             \tag{77.11}
\]
\[
 \frac{d(T_t{}_\#\nu_i)}{d\nu}\leq R_{ij}             \tag{77.12}
\]
for all \(t\), then
\[
 L_{ij}\leq V_{ij},\qquad
 \Omega_{ij}\leq V_{ij}R_{ij}\nu,                    \tag{77.13}
\]
and hence
\[
 \int|f-f\circ T_1|^2d\nu_i
 \leq V_{ij}^2R_{ij}{\cal E}(f).                     \tag{77.14}
\]
\end{proposition}

\begin{proof}
Equation (77.11) integrated over \(t\in[0,1]\) gives the length
bound.  For a nonnegative test function \(h\),
\[
\begin{split}
 \int h\,d\Omega_{ij}
 &=
 \int_0^1\int
 h(T_tx)|Z_t(T_tx)|\,d\nu_i(x)\,dt\\
 &\leq
 V_{ij}R_{ij}\int h\,d\nu,
\end{split}
\]
using (77.12).  Apply Theorem~\ref{thm:v15v77-path}.
\end{proof}

Let \(d\nu=w\,dx\).  Along a smooth ambient flow,
\[
 \frac{d}{dt}
 \log\!\left[
 w(T_tx)|\det DT_t(x)|
 \right]
 =
 \operatorname{div}Z_t+
 D_{Z_t}\log w
 =
 -{\cal S}(Z_t).                                      \tag{77.15}
\]
Thus an integrated bound
\[
 \int_0^1\|{\cal S}(Z_t)\|_\infty dt\leq D_{ij}        \tag{77.16}
\]
controls density distortion by \(e^{D_{ij}}\), apart from the
normalization of the conditional starting cell.  In particular,
\[
 R_{ij}\leq p_i^{-1}e^{D_{ij}}                        \tag{77.17}
\]
whenever the transported restriction has no additional
compression.  Small ownership mass and large score distortion are
therefore genuine bridge costs.

\section{Gauge-corrected slice paths}

Let \({\cal F}(x)=0\) define a slice and let a raw physical path
have velocity \(h_t\).  To keep its representative on the slice,
add an orbit correction \(\xi_t{}_X\):
\[
 \dot x_t=h_t+\xi_t{}_X(x_t).                         \tag{77.18}
\]
Differentiating the gauge condition gives
\[
 0=D{\cal F}(x_t)[h_t]
 +{\cal M}_{\rm FP}(x_t)\xi_t.                        \tag{77.19}
\]

\begin{proposition}[Faddeev--Popov path bound]
\label{prop:v15v77-fp-path}
Assume along the path
\[
 \|{\cal M}_{\rm FP}^{-1}\|\leq\sigma^{-1},\qquad
 \|D{\cal F}\|\leq C_F,\qquad
 \|\xi_X(x)\|\leq C_X\|\xi\|.                         \tag{77.20}
\]
Then the unique gauge correction is
\[
 \xi_t=-{\cal M}_{\rm FP}^{-1}D{\cal F}[h_t],          \tag{77.21}
\]
and
\[
 \|\dot x_t\|
 \leq
 \left(1+\frac{C_XC_F}{\sigma}\right)\|h_t\|.          \tag{77.22}
\]
\end{proposition}

\begin{proof}
Invert (77.19) using (77.20), which gives (77.21) and
\[
 \|\xi_t\|\leq\sigma^{-1}C_F\|h_t\|.
\]
Use the triangle inequality and the orbit-action bound in (77.20)
to obtain (77.22).
\end{proof}

If the raw quotient paths have length at most \(L_0\), their slice
lifts have the estimate
\[
 L_{\rm slice}
 \leq
 \left(1+\frac{C_XC_F}{\sigma}\right)L_0.              \tag{77.23}
\]
Combined with (77.14), this creates at least a quadratic
\(\sigma^{-1}\) loss unless occupation improves simultaneously.

\section{Density derivative along the corrected path}

The slice density contains
\[
 J_{\rm FP}=|\det{\cal M}_{\rm FP}|.
\]
Along (77.18),
\[
 D_{\dot x}\log J_{\rm FP}
 =
 \operatorname{Tr}\bigl(
 {\cal M}_{\rm FP}^{-1}D_{\dot x}{\cal M}_{\rm FP}
 \bigr).                                              \tag{77.24}
\]

\begin{proposition}[Slice-density distortion]
\label{prop:v15v77-density}
If
\[
 \|D_z{\cal M}_{\rm FP}\|_1\leq C_M\|z\|              \tag{77.25}
\]
and (77.20) holds, then
\[
 |D_{\dot x}\log J_{\rm FP}|
 \leq
 \frac{C_M}{\sigma}
 \left(1+\frac{C_XC_F}{\sigma}\right)\|h_t\|.          \tag{77.26}
\]
\end{proposition}

\begin{proof}
The trace inequality in (77.24) gives
\[
 |D_{\dot x}\log J_{\rm FP}|
 \leq\|{\cal M}_{\rm FP}^{-1}\|
 \|D_{\dot x}{\cal M}_{\rm FP}\|_1.
\]
Use (77.20), (77.25), and (77.22).
\end{proof}

Thus the logarithmic density distortion may carry a
\(\sigma^{-2}\) loss even before the action and metric Jacobian are
differentiated.  This is the bridge counterpart of the Hessian
loss in (74.23).

\section{A local ill-conditioned-slice example}

\begin{counterexample}[Shrinking Faddeev--Popov derivative]
\label{cnt:v15v77-slice}
In local coordinates \((\theta,y)\), let the gauge orbit direction
be \(\partial_\theta\) with unit norm and impose
\[
 {\cal F}_\epsilon(\theta,y)=\epsilon\theta+y=0.       \tag{77.27}
\]
Then
\[
 {\cal M}_{\rm FP}=
 D{\cal F}_\epsilon[\partial_\theta]=\epsilon.         \tag{77.28}
\]
For the raw velocity \(h=\partial_y\), maintaining the slice
requires
\[
 \xi=-\epsilon^{-1},\qquad
 \dot x=\partial_y-\epsilon^{-1}\partial_\theta,      \tag{77.29}
\]
whose norm is \((1+\epsilon^{-2})^{1/2}\).
\end{counterexample}

\begin{proof}
Equation (77.19) is
\[
 1+\epsilon\xi=0.
\]
Solve it and compute the Euclidean norm.
\end{proof}

The physical quotient need not be singular in this example; the
chosen slice is ill conditioned.  It proves that a cofinal atlas
must control slice conditioning rather than merely cover every
orbit.

\section{Cofinal overlap criterion}

\begin{theorem}[Uniform bridge from slice flows]
\label{thm:v15v77-uniform}
Suppose every chart-graph edge admits a slice-corrected flow such
that:
\begin{enumerate}
\item raw quotient path lengths and action/reference-density scores
      have uniform bounds;
\item \(\sigma,C_F,C_X,C_M\) in (77.20), (77.25) have uniform
      cofinal bounds with \(\sigma>0\);
\item conditional cell masses and transported density ratios are
      compatible with the graph conductances;
\item the aggregate occupation inequality (77.8) has a uniform
      constant.
\end{enumerate}
Then the bridge constants in V76 are cutoff uniform.
\end{theorem}

\begin{proof}
Proposition~\ref{prop:v15v77-fp-path} gives uniform path lengths.
Proposition~\ref{prop:v15v77-density}, together with the other
score rows and (77.15), gives uniform density distortion.
Proposition~\ref{prop:v15v77-flow} constructs each path-coupling
estimate, and the aggregate hypothesis plus
Corollary~\ref{cor:v15v77-aggregate} gives the uniform bridge.
\end{proof}

\begin{assessment}[Physical status]
\label{ass:v15v77-status}
V77 proves a path-occupation bridge theorem, derives it from
transport flows, and tracks the inverse Faddeev--Popov losses in
both slice length and density distortion.  It provides an explicit
route to the bridge hypothesis of V76.

No cofinal SM gauge-slice flows, uniform singular-value floor,
density-distortion estimate, or aggregate congestion bound has
been constructed.  The Gribov obstruction remains explicit.
\end{assessment}

\begin{decision}[V78 discrete chart-graph mixing]
\label{dec:v15v77-next}
Attack the remaining between-chart constant \(C_G\).  Prove
discrete Poincar\'e and logarithmic Sobolev estimates from
conductance/minorization data, track degradation on a long path of
charts, and identify product or multiscale atlas structures that
can keep the cofinal constant uniform as the chart count grows.
\end{decision}

\begin{firewall}[Claim boundary after V77]
\label{fire:v15v77}
V77 proves conditional overlap-transport estimates.  It does not
establish a physical slice atlas, uniform Faddeev--Popov
invertibility, graph mixing, quotient boundary cancellation, or
the full continuum.
\end{firewall}

\section*{Version V77 append-only record}
\addcontentsline{toc}{section}{Version V77 append-only record}
The complete V76 source was retained before this continuation.  It
had \(756926\) bytes, \(22151\) lines, and SHA--256
\[
 \texttt{8C60BCC1BB678050EEDCFFD10C19E2F05B70BC75F4F0C3325804D3BA395DB1D3}.
\]
V77 constructed path-coupling and flow-occupation bridges, derived
the gauge-corrected slice path, and quantified the
Faddeev--Popov singular-value losses in path length and density
distortion.

\chapter{Fifteenth-cycle continuation V78: cofinal mixing of the chart graph}
\label{sec:v15v78}

\section{Discrete Dirichlet form}

Let \(I\) be a finite chart-label set with probability vector
\(p=(p_i)\) and symmetric conductances \(c_{ij}=c_{ji}\geq0\).
Write
\[
 {\cal E}_G(a)
 =
 \frac12\sum_{i,j}c_{ij}(a_i-a_j)^2.                  \tag{78.1}
\]
The spectral-gap and logarithmic Sobolev constants are defined by
\[
 \operatorname{Var}_p(a)
 \leq C_P{\cal E}_G(a),                               \tag{78.2}
\]
\[
 \operatorname{Ent}_p(a^2)
 \leq2C_G{\cal E}_G(a).                               \tag{78.3}
\]

\begin{lemma}[A spectral-gap-to-LSI bound]
\label{lem:v15v78-gap-lsi}
Let
\[
 p_*=\min_i p_i>0.
\]
If (78.2) holds, then
\[
 C_G
 \leq
 \frac{C_P}{2}\left(\log\frac1{p_*}+2\right).          \tag{78.4}
\]
\end{lemma}

\begin{proof}
Let \(\bar a=\sum_ip_ia_i\) and \(g=a-\bar a\).  The entropy shift
inequality used in V73 gives
\[
 \operatorname{Ent}_p(a^2)
 \leq\operatorname{Ent}_p(g^2)+2\operatorname{Var}_p(a).
                                                               \tag{78.5}
\]
Put \(V=\sum_ip_ig_i^2\).  Since
\[
 g_i^2\leq V/p_i\leq V/p_*,
\]
one has
\[
 \operatorname{Ent}_p(g^2)
 =
 \sum_ip_ig_i^2\log\frac{g_i^2}{V}
 \leq V\log\frac1{p_*}.                               \tag{78.6}
\]
Use (78.2) in (78.5)--(78.6), and compare with (78.3).
\end{proof}

The bound is deliberately crude but cofinally informative:
controlling only the spectral gap can still leave a
\(\log(1/p_*)\) loss.

\section{Minorization and bottlenecks}

Let \(P(i,j)\) be a reversible Markov kernel with stationary law
\(p\), and set
\[
 c_{ij}=p_iP(i,j).                                    \tag{78.7}
\]

\begin{proposition}[Doeblin chart refresh]
\label{prop:v15v78-doeblin}
If
\[
 P(i,j)\geq\epsilon p_j                               \tag{78.8}
\]
for every \(i,j\), then
\[
 C_P\leq\epsilon^{-1},\qquad
 C_G\leq
 \frac{\log(1/p_*)+2}{2\epsilon}.                     \tag{78.9}
\]
\end{proposition}

\begin{proof}
From (78.7)--(78.8),
\[
\begin{split}
 {\cal E}_G(a)
 &\geq
 \frac{\epsilon}{2}
 \sum_{i,j}p_ip_j(a_i-a_j)^2\\
 &=\epsilon\operatorname{Var}_p(a).
\end{split}
\]
This proves the Poincar\'e bound.  Apply
Lemma~\ref{lem:v15v78-gap-lsi}.
\end{proof}

For a subset \(S\subset I\), define
\[
 p(S)=\sum_{i\in S}p_i,\qquad
 Q(S,S^c)=\sum_{i\in S,j\notin S}c_{ij}.              \tag{78.10}
\]

\begin{proposition}[Exact bottleneck lower bounds]
\label{prop:v15v78-bottleneck}
For \(0<p(S)\leq1/2\),
\[
 C_P\geq
 \frac{p(S)(1-p(S))}{Q(S,S^c)},                       \tag{78.11}
\]
\[
 C_G\geq
 \frac{p(S)\log(1/p(S))}
 {2Q(S,S^c)}.                                         \tag{78.12}
\]
\end{proposition}

\begin{proof}
Use \(a={\bf1}_S\) in (78.2).  Its variance is
\(p(S)(1-p(S))\), and its Dirichlet energy is
\(Q(S,S^c)\), proving (78.11).

For (78.12), use
\[
 a=p(S)^{-1/2}{\bf1}_S.
\]
Then \(\mathbb E_pa^2=1\),
\(\operatorname{Ent}_p(a^2)=\log(1/p(S))\), and
\({\cal E}_G(a)=Q(S,S^c)/p(S)\).  Insert these values in
(78.3).
\end{proof}

A cluster of slice charts connected to the rest by small overlap
conductance therefore forces a large between-chart constant.

\section{Long-chain failure}

\begin{counterexample}[A path of charts]
\label{cnt:v15v78-path}
Let \(I_N=\{1,\ldots,N\}\), \(p_i=1/N\), and let the nearest-neighbor
Markov chain jump left or right at rates of order one.  Equivalently,
take
\[
 c_{i,i+1}=1/N,\qquad1\leq i<N.                       \tag{78.13}
\]
Then
\[
 C_P\geq cN^2,\qquad C_G\geq cN^2                    \tag{78.14}
\]
for an absolute \(c>0\).
\end{counterexample}

\begin{proof}
Take \(a_i=i\).  Its variance under the uniform law is
\[
 \operatorname{Var}_p(a)=\frac{N^2-1}{12},
\]
while
\[
 {\cal E}_G(a)
 =
 \sum_{i=1}^{N-1}\frac1N(a_{i+1}-a_i)^2
 =
 \frac{N-1}{N}\leq1.
\]
This proves the Poincar\'e lower bound.  Linearizing (78.3) around
the constant vector shows that an LSI with constant \(C_G\)
implies (78.2) with constant \(C_G\), proving the second bound.
\end{proof}

Thus a cofinal atlas organized as an increasingly long chain cannot
give the uniform constant required in V76, even if every chart and
every adjacent overlap is individually uniform.

\section{Tensor-product chart labels}

Let
\[
 (I,p)=\prod_{\ell=1}^{N}(I_\ell,p_\ell).             \tag{78.15}
\]
For a function \(a\) on \(I\), let
\({\cal E}_\ell(a)\) be the Dirichlet form that updates only label
\(\ell\), with the other labels frozen, averaged over those other
labels.  Put
\[
 {\cal E}_{\rm prod}(a)=\sum_{\ell=1}^N{\cal E}_\ell(a).
                                                               \tag{78.16}
\]

\begin{theorem}[Discrete LSI tensorization]
\label{thm:v15v78-tensor}
If every coordinate measure and coordinate form satisfies
\[
 \operatorname{Ent}_{p_\ell}(b^2)
 \leq2C_*\mathcal E_\ell(b)                           \tag{78.17}
\]
with one constant \(C_*\), then
\[
 \operatorname{Ent}_{p}(a^2)
 \leq2C_*{\cal E}_{\rm prod}(a),                      \tag{78.18}
\]
independently of \(N\).
\end{theorem}

\begin{proof}
Entropy tensorization for a product probability law gives
\[
 \operatorname{Ent}_p(a^2)
 \leq
 \sum_{\ell=1}^N
 \mathbb E_{p_{\widehat\ell}}
 \operatorname{Ent}_{p_\ell}
 \bigl(a^2\,|\,i_{\widehat\ell}\bigr).                \tag{78.19}
\]
Apply (78.17) to each conditional function and sum.  The right side
is \(2C_*{\cal E}_{\rm prod}(a)\).
\end{proof}

The normalization of (78.16) matters.  If only one uniformly
chosen coordinate is refreshed per unit time, the form is
\(N^{-1}{\cal E}_{\rm prod}\) and its LSI constant grows by \(N\).
Uniform tensorization corresponds to a fixed update rate per
coordinate.

\section{Conditional tensorization criterion}

Exact product structure is usually broken by gauge and matter
interactions.  The following statement isolates the additional
input rather than assuming it implicitly.

\begin{theorem}[Approximate chart-label tensorization compiler]
\label{thm:v15v78-approx}
Suppose a nonproduct chart-label law \(p\) satisfies the entropy
factorization
\[
 \operatorname{Ent}_{p}(a^2)
 \leq
 C_{\rm at}
 \sum_{\ell=1}^N
 \mathbb E_p
 \operatorname{Ent}_{p_\ell(\cdot|i_{\widehat\ell})}
 (a^2),                                               \tag{78.20}
\]
and every conditional one-coordinate chain satisfies an LSI with
constant \(C_*\).  Then the full chart graph satisfies
\[
 C_G\leq C_{\rm at}C_*                                \tag{78.21}
\]
for the sum of the conditional coordinate Dirichlet forms.
\end{theorem}

\begin{proof}
Apply each conditional LSI on the right side of (78.20).  Their
averaged sum is the full coordinate-update Dirichlet form.
\end{proof}

Proving a cutoff-uniform \(C_{\rm at}\) is the discrete
weak-dependence problem.  It is not a consequence of bounded local
degree.

\section{Graph scaling cannot improve the glued constant}

\begin{proposition}[Conductance rescaling invariance]
\label{prop:v15v78-rescale}
Replace every graph conductance by
\[
 \widetilde c_{ij}=\alpha c_{ij},\qquad\alpha>0.
\]
Then the optimal graph LSI constant changes as
\[
 \widetilde C_G=C_G/\alpha,                           \tag{78.22}
\]
while the optimal bridge constant in (76.13) changes as
\[
 \widetilde C_B=\alpha C_B.                           \tag{78.23}
\]
Hence
\[
 \widetilde C_G\widetilde C_B=C_GC_B.                 \tag{78.24}
\]
\end{proposition}

\begin{proof}
Both graph energies in (76.12) and (76.13) are multiplied by
\(\alpha\).  The stated transformations follow from their
definitions.
\end{proof}

One cannot manufacture a uniform quotient LSI by multiplying all
chart conductances by a cutoff-dependent factor.

\section{Cofinal graph alternatives}

The between-chart row can now be populated by any one of:
\begin{enumerate}
\item a direct uniform graph LSI;
\item a Doeblin refresh with both \(\epsilon\) and the rare-cell
      entropy cost controlled;
\item exact tensorization with fixed per-coordinate update rate;
\item approximate entropy tensorization (78.20) with a uniform
      weak-dependence constant.
\end{enumerate}
Every alternative must be combined with the bridge congestion of
V77.  A graph with excellent abstract mixing can still fail if its
edges correspond to paths crossing an ill-conditioned Gribov
region.

\begin{assessment}[Physical status]
\label{ass:v15v78-status}
V78 proves discrete spectral and entropy bounds, exact bottleneck
lower bounds, long-chain degeneration, tensorized uniform mixing,
and conductance-rescaling invariance.  It supplies a precise menu
for the graph hypothesis in V76.

The physical slice atlas has not been shown to possess product or
weak-dependence structure, uniform minorization, or bridge-compatible
conductances.  The cofinal quotient LSI remains conditional.
\end{assessment}

\begin{decision}[V79 complex phase over the modulus measure]
\label{dec:v15v78-next}
Return to the complete chiral theory.  Assuming the modulus
measure has the V76 LSI, prove the exact phase-reweighting formula
for observables and Schwinger--Dyson identities, quantify stability
under a uniform lower bound on the average phase, and show that
modulus tightness alone gives no control when the average phase
vanishes cofinally.
\end{decision}

\begin{firewall}[Claim boundary after V78]
\label{fire:v15v78}
V78 proves discrete graph-mixing compilers.  It does not construct
the physical chart graph, prove approximate tensorization, control
Gribov bridges, or establish the continuum theory.
\end{firewall}

\section*{Version V78 append-only record}
\addcontentsline{toc}{section}{Version V78 append-only record}
The complete V77 source was retained before this continuation.  It
had \(767544\) bytes, \(22528\) lines, and SHA--256
\[
 \texttt{CAF3B384CEC5820434602B807B01CC0392E3E6BDBC30AD417298E821CE33EB51}.
\]
V78 proved chart-graph bottleneck and tensorization theorems,
quantified long-chain failure, and showed that conductance
rescaling leaves the glued graph--bridge product invariant.

\chapter{Fifteenth-cycle continuation V79: phase reweighting over the quotient modulus}
\label{sec:v15v79}

\section{Normalized complex functional}

Let \(\mu_n\) be the positive quotient modulus measure produced
conditionally by V76, and let
\[
 z_n=e^{i\theta_n},\qquad |z_n|=1,                    \tag{79.1}
\]
be the determinant/Pfaffian phase.  Define the average phase
\[
 a_n=\mathbb E_{\mu_n}z_n.                            \tag{79.2}
\]
When \(a_n\ne0\), the normalized complex expectation is
\[
 \langle F\rangle_n
 =
 \frac{\mathbb E_{\mu_n}[Fz_n]}{a_n}.                 \tag{79.3}
\]

\begin{proposition}[Total variation of the reweighted measure]
\label{prop:v15v79-variation}
The complex measure
\[
 d\rho_n=\frac{z_n}{a_n}\,d\mu_n                     \tag{79.4}
\]
has mass one and total variation
\[
 \|\rho_n\|_{\rm TV}=\frac1{|a_n|}.                   \tag{79.5}
\]
\end{proposition}

\begin{proof}
Equation (79.2) gives
\(\rho_n(\Omega_n)=1\).  Since \(|z_n|=1\),
\[
 |\rho_n|(\Omega_n)
 =
 \int\left|\frac{z_n}{a_n}\right|d\mu_n
 =
 |a_n|^{-1}.
\]
\end{proof}

Thus the average phase is not merely a diagnostic.  Its reciprocal
is exactly the operator norm of reweighting on bounded
observables.

\section{Cofinal stability under a phase floor}

\begin{theorem}[Stable phase reweighting]
\label{thm:v15v79-stable}
Suppose
\[
 \inf_n|a_n|\geq a_*>0.                               \tag{79.6}
\]
Then for every \(F_n\in L^1(\mu_n)\),
\[
 |\langle F_n\rangle_n|
 \leq a_*^{-1}\mathbb E_{\mu_n}|F_n|.                 \tag{79.7}
\]
If \((F_n)\) is uniformly integrable under \(\mu_n\), then its
reweighted tails are uniformly negligible:
\[
 \sup_n
 \left|
 \int_{\{|F_n|>R\}}F_n\,d\rho_n
 \right|
 \longrightarrow0.                                   \tag{79.8}
\]
If the modulus measures are tight, then the complex measures are
uniformly tight in total variation.
\end{theorem}

\begin{proof}
Use \(|z_n|=1\) and (79.6) in (79.3), which gives (79.7).
The same estimate restricted to \(\{|F_n|>R\}\) proves (79.8).
For a compact set \(K\),
\[
 |\rho_n|(K^c)=|a_n|^{-1}\mu_n(K^c)
 \leq a_*^{-1}\mu_n(K^c),
\]
which transfers tightness.
\end{proof}

\begin{theorem}[Limit of complex expectations]
\label{thm:v15v79-limit}
On transported configuration spaces, suppose
\[
 \mu_n\Longrightarrow\mu,\qquad z_n\to z,\qquad F_n\to F
                                                               \tag{79.9}
\]
in the convergence sense of V44, with
\((F_n)\) uniformly integrable and \(|z_n|=|z|=1\).
If
\[
 a=\mathbb E_\mu z\ne0,                               \tag{79.10}
\]
then
\[
 a_n\to a,\qquad
 \langle F_n\rangle_n
 \longrightarrow
 \frac{\mathbb E_\mu[Fz]}{\mathbb E_\mu z}.            \tag{79.11}
\]
\end{theorem}

\begin{proof}
The phases are bounded, so transported convergence and modulus
tightness give
\(\mathbb E_{\mu_n}z_n\to\mathbb E_\mu z\).
Uniform integrability of \(F_n\) also applies to \(F_nz_n\), giving
convergence of the numerator.  The denominator tends to the
nonzero number in (79.10), so the quotient converges.
\end{proof}

\begin{proposition}[Quantitative quotient error]
\label{prop:v15v79-error}
For complex numbers \(N_n,N,a_n,a\) with
\(|a_n|,|a|\geq a_*>0\),
\[
 \left|\frac{N_n}{a_n}-\frac{N}{a}\right|
 \leq
 \frac{|N_n-N|}{a_*}
 +\frac{|N|\,|a_n-a|}{a_*^2}.                         \tag{79.12}
\]
\end{proposition}


\begin{proof}
Write
\[
 \frac{N_n}{a_n}-\frac Na
 =
 \frac{N_n-N}{a_n}
 +N\frac{a-a_n}{a_na}
\]
and take absolute values.
\end{proof}

\section{Full complex Schwinger--Dyson score}

Let the full density be
\[
 e^{-{\cal A}_{R,n}+i\theta_n}.
\]
For a constant metric direction \(v\), define
\[
 {\cal S}^{\rm full}_{n,v}
 =
 D_v{\cal A}_{R,n}-iD_v\theta_n.                      \tag{79.13}
\]

\begin{theorem}[Phase-corrected metric identity]
\label{thm:v15v79-sd}
Assume \(a_n\ne0\), vanishing boundary flux, and integrability of
the displayed terms.  Then
\[
 \boxed{
 \langle D_vF\rangle_n
 =
 \left\langle
 F\bigl(D_v{\cal A}_{R,n}-iD_v\theta_n\bigr)
 \right\rangle_n.}                                    \tag{79.14}
\]
\end{theorem}

\begin{proof}
Use the modulus Schwinger--Dyson identity with the test function
\(Fz_n\):
\[
 \mathbb E_{\mu_n}D_v(Fz_n)
 =
 \mathbb E_{\mu_n}
 [Fz_nD_v{\cal A}_{R,n}].
\]
Since
\[
 D_v(Fz_n)=z_nD_vF+iFz_nD_v\theta_n,
\]
rearrange and divide by \(a_n\).
\end{proof}

\begin{corollary}[Cofinal full-score integrability]
\label{cor:v15v79-full-ui}
Assume (79.6), a uniform exponential modulus-score moment from
V72, and
\[
 \sup_n\mathbb E_{\mu_n}e^{\tau|D_v\theta_n|}<\infty   \tag{79.15}
\]
for some \(\tau>0\).  Then
\[
 F_n{\cal S}^{\rm full}_{n,v}z_n
\]
is uniformly integrable under \(\mu_n\) for every uniformly bounded
\(F_n\), and the reweighted score insertions are uniformly
integrable in total variation.
\end{corollary}

\begin{proof}
The triangle inequality and
\[
 e^{(\tau/2)(|S_R|+|D_v\theta|)}
 \leq
 \frac12e^{\tau|S_R|}+\frac12e^{\tau|D_v\theta|}
\]
give a uniform exponential moment for the full score.  Bounded
\(F_n\) and \(|z_n|=1\) preserve it.  Apply V63 and then the
factor \(a_*^{-1}\) from Theorem~\ref{thm:v15v79-stable}.
\end{proof}

The derivative of the phase is an independent stress row.  A lower
bound on \(|a_n|\) does not control \(D_v\theta_n\).

\section{Two sufficient phase-floor criteria}

\begin{proposition}[Phase confined to an open semicircle]
\label{prop:v15v79-semicircle}
If there are real \(c_n\) and \(0\leq\alpha<\pi/2\) such that a
real lift obeys
\[
 |\theta_n-c_n|\leq\alpha\quad\mu_n\hbox{-almost surely},        \tag{79.16}
\]
then
\[
 |a_n|\geq\cos\alpha.                                 \tag{79.17}
\]
\end{proposition}

\begin{proof}
Rotate by \(e^{-ic_n}\).  The real part of every rotated phase is
at least \(\cos\alpha\), so
\[
 |a_n|\geq
 \operatorname{Re}(e^{-ic_n}a_n)\geq\cos\alpha.
\]
\end{proof}

\begin{proposition}[Small phase variance]
\label{prop:v15v79-variance}
Suppose \(\theta_n\) has a global real lift and
\[
 \operatorname{Var}_{\mu_n}(\theta_n)\leq\sigma^2<2.  \tag{79.18}
\]
Then
\[
 |a_n|\geq1-\frac{\sigma^2}{2}>0.                     \tag{79.19}
\]
If \(\mu_n\) satisfies \({\rm LSI}(C_{\rm LS})\) and
\[
 |\nabla\theta_n|\leq L_\theta,\qquad
 C_{\rm LS}L_\theta^2<2,                              \tag{79.20}
\]
then (79.19) holds with
\(\sigma^2=C_{\rm LS}L_\theta^2\).
\end{proposition}

\begin{proof}
Let \(c_n=\mathbb E\theta_n\).  Since
\(\cos x\geq1-x^2/2\),
\[
\begin{split}
 |a_n|
 &\geq
 \operatorname{Re}(e^{-ic_n}a_n)\\
 &=
 \mathbb E\cos(\theta_n-c_n)
 \geq1-\frac12\operatorname{Var}(\theta_n).
\end{split}
\]
An LSI implies the Poincar\'e inequality with the same constant,
so (79.20) gives
\(\operatorname{Var}\theta_n\leq C_{\rm LS}L_\theta^2\).
\end{proof}

A global real lift is itself a determinant-line condition.  A
nontrivial phase winding or anomaly can prevent it even when the
local derivative is bounded.

\section{Vanishing average phase}

\begin{counterexample}[Tight modulus with divergent reweighting]
\label{cnt:v15v79-sign}
Let \(\Omega=\{+,-\}\), set
\[
 \mu_n(+)=\frac{1+\epsilon_n}{2},\qquad
 \mu_n(-)=\frac{1-\epsilon_n}{2},\qquad
 \epsilon_n\downarrow0,
\]
and let
\[
 z_n(+)=1,\qquad z_n(-)=-1.                           \tag{79.21}
\]
Then the modulus measures live on one fixed finite space and are
tight, but
\[
 a_n=\epsilon_n,\qquad
 \|\rho_n\|_{\rm TV}=\epsilon_n^{-1}\to\infty.         \tag{79.22}
\]
For the bounded observable \(F_n=z_n\),
\[
 \langle F_n\rangle_n=\epsilon_n^{-1}\to\infty.       \tag{79.23}
\]
\end{counterexample}

\begin{proof}
The first two identities follow by direct summation and
Proposition~\ref{prop:v15v79-variation}.  Since \(z_n^2=1\),
the numerator in (79.3) for \(F_n=z_n\) equals one, proving
(79.23).
\end{proof}

Thus modulus tightness, an LSI, and bounded observables do not
control the normalized complex theory when the average phase
collapses.

\section{Reflection positivity remains separate}

A phase floor controls normalization and total variation.  It does
not imply
\[
 \langle(\Theta F)F\rangle\geq0                       \tag{79.24}
\]
for reflected observables.  The reflected phase must factor or
cancel as in V51.  Likewise, local anomaly-coefficient
cancellation must be supplemented by a global determinant-line
orientation and compatible quotient cocycle.

\begin{assessment}[Physical status]
\label{ass:v15v79-status}
V79 constructs the normalized complex continuum compiler, proves
that the reciprocal average phase is exactly its total variation,
derives the phase-corrected metric Schwinger--Dyson identity, and
gives two sufficient phase-floor criteria plus a sharp
vanishing-phase counterexample.

The chiral Standard Model determinant has not yet been shown to
admit a global real phase lift, a uniform average-phase floor, an
exponential phase-derivative moment, or the reflected factorization
needed for Osterwalder--Schrader positivity.
\end{assessment}

\begin{decision}[V80 cross-program audit]
\label{dec:v15v79-next}
Perform the scheduled companion audit, looking specifically for a
new determinant-orientation, reflection-positive, or cofinal
normalization input.  Record the result, then continue in V81 with
the global anomaly and determinant-line obstruction to a natural
phase lift.
\end{decision}

\begin{firewall}[Claim boundary after V79]
\label{fire:v15v79}
V79 proves phase-reweighting implications under a nonvanishing
average phase.  It does not prove that floor, control the physical
phase derivative, cancel global anomalies, or establish reflection
positivity and the full continuum.
\end{firewall}

\section*{Version V79 append-only record}
\addcontentsline{toc}{section}{Version V79 append-only record}
The complete V78 source was retained before this continuation.  It
had \(777544\) bytes, \(22878\) lines, and SHA--256
\[
 \texttt{43187E1731DFD9EE4DF953F7602B4F57B0F85EEF385F1B7C8899F56C2D6A26D5}.
\]
V79 proved stable complex reweighting under an average-phase floor,
derived the full phase-corrected metric score, and isolated phase
lifting, derivative, normalization, and reflection as distinct
closure rows.


\chapter{Fifteenth-cycle continuation V80: scheduled phase and reflection audit}
\label{sec:v15v80}

\section{Updated companion snapshots}

The active companion identities are now
\[
\begin{array}{c|c|c|c|l}
\hbox{programme}&\hbox{version}&\hbox{bytes}&\hbox{lines}&\hbox{SHA--256}\\ \hline
\hbox{Yang--Mills}&V40&344764&9458&
\texttt{579A346B3193F33B2E082C6682E83271843B665AC950E70D7F7707FE19EAF7E2}\\
\hbox{Navier--Stokes}&V26&278283&5319&
\texttt{82C887F8C2C69E4F28FAD7C3DC8DE5B9099CB5A4BF431EBA1FFF3E91935978AD}.
\end{array}                                             \tag{80.1}
\]
YM advanced from V39 to V40; NS is unchanged.

\section{The new Yang--Mills result}

YM V40 first rejects radius-\(1/8\) isotropic boxes at two exact
centres by negative rational quadratic-form witnesses.  It then
uses two anisotropic repairs and proves pointwise positivity of the
finite \(45\times45\) Wilson pencil on
\[
 [0,\pi/4]^2\times[0,2t].                             \tag{80.2}
\]
The third coordinate still has not reached \(\pi/4\).  Its firewall
continues to exclude a full transverse quarter cell, the transverse
torus, a cofinal operator inequality, an interacting continuum
measure, and reflection positivity.

\begin{proposition}[Valid V40 import]
\label{prop:v15v80-valid-import}
YM V40 strengthens the available local symbol-atlas input from
(75.2) to (80.2), and its anisotropic repair is compatible with the
nonuniform chart geometry allowed in V76--V77.
\end{proposition}

\begin{proof}
The terminal YM theorem explicitly proves the larger prism and
lists positive exact floors for all sixteen new charts.  V76--V77
permit chart-dependent local geometry, provided the resulting
distortion, graph, and bridge constants are tracked.  No isotropic
radius was assumed in those gluing theorems.
\end{proof}

\begin{proposition}[No determinant-phase or reflection import]
\label{prop:v15v80-no-phase-import}
The phase variables in YM V40 do not supply the phase lift,
average-phase floor, phase derivative, determinant-line
orientation, or reflected factorization required in V79.
\end{proposition}

\begin{proof}
The YM phases parametrize momentum entries of a finite Wilson
matrix pencil and are used to certify Hermitian positivity.
The V79 phase is the normalized complex phase of the chiral
fermion determinant as a function on the gauge--metric quotient.
YM V40 constructs no determinant line, complex reweighted measure,
average phase, or Osterwalder--Schrader form.  Its own firewall
leaves the reflection-positive continuum limit open.
\end{proof}

\section{Navier--Stokes comparison}

NS V26 remains the pointwise rank-one Oseen-kernel norm note.
Neither its terminal theorem nor its earlier concentration ledgers
contain a determinant/Pfaffian line, chiral phase, or reflection
operation on a Euclidean field measure.  No phase or reflection
row transfers.

\section{Direction after the audit}

\begin{theorem}[V80 dependency decision]
\label{thm:v15v80-decision}
The dependency graph changes only in the local finite-dimensional
symbol row: the certified prism is now (80.2).  The global complex
phase and reflection rows remain unchanged.  The next direct task
is therefore the topology and connection of the chiral determinant
line over the quotient.
\end{theorem}

\begin{proof}
Equation (80.1) identifies the sources.
Propositions~\ref{prop:v15v80-valid-import} and
\ref{prop:v15v80-no-phase-import} give the positive and negative
parts of the comparison, while the NS source is unchanged.
\end{proof}

\begin{assessment}[Physical status]
\label{ass:v15v80-status}
V80 records a genuine YM improvement without promoting a momentum
phase to a determinant phase.  It adds no global anomaly,
reflection-positive, or continuum-measure theorem.
\end{assessment}

\begin{decision}[V81 determinant-line obstruction]
\label{dec:v15v80-next}
Construct the determinant line over the regular background
quotient, distinguish its first Chern class, connection curvature,
and flat holonomy, and prove the lifting criterion for a global
real phase.  Audit the local Standard Model gauge-anomaly
coefficients and the \(SU(2)\) mod-two count, while keeping global
gauge, gravitational, reflection, and average-phase conditions
separate.
\end{decision}

\begin{firewall}[Claim boundary after V80]
\label{fire:v15v80}
V80 is a cross-program audit.  The larger finite YM symbol prism
does not prove a global determinant orientation, average-phase
floor, reflection-positive measure, or continuum theory.
\end{firewall}

\section*{Version V80 append-only record}
\addcontentsline{toc}{section}{Version V80 append-only record}
The complete V79 source was retained before this continuation.  It
had \(788048\) bytes, \(23251\) lines, and SHA--256
\[
 \texttt{AB7BFE454819298B5D87C5188CEB6DFCDD0FB76DF3C64A454BC77920DA4F38C5}.
\]
V80 performed the scheduled audit, imported YM V40 only as a
larger local Wilson-symbol prism, and separated its momentum phases
from the chiral determinant phase and reflection requirements.

\chapter{Fifteenth-cycle continuation V81: determinant lines, local curvature, and global anomaly}
\label{sec:v15v81}

\section{The determinant line over backgrounds}

Let \({\cal B}\) denote a regular component of the gauge--metric
background quotient and let
\[
 D_b:{\cal H}_b^+\longrightarrow{\cal H}_b^-,
 \qquad b\in{\cal B},                                  \tag{81.1}
\]
be a continuous Fredholm family of chiral Dirac operators.  Its
determinant line is
\[
 {\rm Det}\,D_b
 =
 \Lambda^{\rm top}\ker D_b
 \otimes
 \left(\Lambda^{\rm top}{\rm coker}\,D_b\right)^*.     \tag{81.2}
\]
These fibres form a complex line bundle
\({\cal L}_{\rm det}\to{\cal B}\).

\begin{proposition}[Nonvanishing determinant and trivialization]
\label{prop:v15v81-trivial}
A continuous nowhere-zero determinant section on \({\cal B}\)
trivializes \({\cal L}_{\rm det}\).  Conversely, a trivialization
provides a nowhere-zero section.  On a paracompact base of CW type,
\[
 {\cal L}_{\rm det}\hbox{ is topologically trivial}
 \quad\Longleftrightarrow\quad
 c_1({\cal L}_{\rm det})=0.                            \tag{81.3}
\]
\end{proposition}

\begin{proof}
A nowhere-zero section \(s\) defines the fibrewise isomorphism
\[
 {\cal B}\times{\mathbb C}\longrightarrow{\cal L}_{\rm det},
 \qquad(b,z)\longmapsto zs(b).
\]
The converse uses the image of the constant section \(1\).
Complex line bundles over a paracompact CW-type base are classified
by homotopy classes of maps to
\({\mathbb CP}^{\infty}=K({\mathbb Z},2)\), equivalently by their
first Chern class.  This gives (81.3).
\end{proof}

The regularized determinant section may vanish when \(D_b\) has
zero modes.  At such a point its phase is undefined even though the
line bundle remains meaningful.

\section{Lifting a circle-valued phase}

Suppose a trivialization has been chosen on an invertible locus
and the normalized determinant gives
\[
 z:{\cal B}\longrightarrow U(1).                      \tag{81.4}
\]

\begin{theorem}[Global real-phase lifting criterion]
\label{thm:v15v81-lift}
Assume \({\cal B}\) is connected, locally path connected, and
semilocally simply connected.  There is a continuous
\(\theta:{\cal B}\to{\mathbb R}\) with
\[
 z=e^{i\theta}                                        \tag{81.5}
\]
if and only if
\[
 z_*\pi_1({\cal B})=0\subset\pi_1(U(1))\simeq{\mathbb Z}.        \tag{81.6}
\]
Equivalently, every loop in \({\cal B}\) has zero phase winding.
\end{theorem}

\begin{proof}
The exponential map
\[
 {\mathbb R}\longrightarrow U(1),\qquad t\longmapsto e^{it}
\]
is the universal covering.  The covering-space lifting theorem
says that \(z\) lifts after fixing one basepoint value precisely
when
\[
 z_*\pi_1({\cal B})
 \subset
 \exp_*\pi_1({\mathbb R})=\{0\}.
\]
This is (81.6).  A lift is exactly (81.5).
\end{proof}

Topological triviality of the line and liftability of one chosen
circle-valued section are related but distinct statements.  The
first makes a scalar phase possible; the second removes its winding
in that trivialization.

\section{Connection, local anomaly, and flat holonomy}

Choose local unit frames \(e_i\) of the determinant line.  On an
overlap write
\[
 e_j=e^{i\chi_{ij}}e_i.                               \tag{81.7}
\]
Let a unit determinant section be
\[
 s=e^{i\theta_i}e_i.
\]
Then
\[
 \theta_j=\theta_i-\chi_{ij}\pmod{2\pi}.              \tag{81.8}
\]
For a unitary determinant connection, write
\[
 \nabla e_i=i{\cal A}_i e_i.
\]
The connection forms obey
\[
 {\cal A}_j={\cal A}_i+d\chi_{ij}.                    \tag{81.9}
\]

\begin{proposition}[Invariant phase derivative]
\label{prop:v15v81-covariant-phase}
The local one-forms
\[
 \eta_i=d\theta_i+{\cal A}_i                          \tag{81.10}
\]
agree on overlaps and define a global real one-form \(\eta\) on
the invertible locus.
\end{proposition}

\begin{proof}
Use (81.8)--(81.9):
\[
 d\theta_j+{\cal A}_j
 =
 d\theta_i-d\chi_{ij}+{\cal A}_i+d\chi_{ij}
 =
 d\theta_i+{\cal A}_i.
\]
\end{proof}

Thus the phase-stress insertion in V79 must be interpreted
covariantly as \(\eta(k)\).  The raw derivative \(D_k\theta_i\)
depends on the local determinant frame.

The curvature is
\[
 {\cal F}_{\rm det}=d{\cal A}_i,                      \tag{81.11}
\]
independent of \(i\).  Its gauge and diffeomorphism components are
the local anomaly.

\begin{theorem}[Parallel phase criterion]
\label{thm:v15v81-parallel}
On a connected component, a global unit parallel section exists if
and only if
\[
 {\cal F}_{\rm det}=0                                 \tag{81.12}
\]
and every loop has trivial connection holonomy:
\[
 {\rm Hol}_{\rm det}(\gamma)=1
 \quad\hbox{for all }\gamma\in\pi_1({\cal B}).         \tag{81.13}
\]
\end{theorem}

\begin{proof}
A parallel section forces zero curvature by
\(\nabla^2s=i{\cal F}_{\rm det}s\), and parallel transport around a
loop must return the section, giving (81.13).

Conversely, fix a unit vector in one fibre and define its value at
another point by parallel transport along any path.  Zero curvature
makes transport invariant under contractible path deformations,
and trivial holonomy makes it invariant under changing the path by
an arbitrary loop.  The resulting section is well defined, unit,
and parallel.
\end{proof}

Vanishing curvature does not imply (81.13).  A flat line on a
circle may have holonomy \(e^{i\alpha}\ne1\).  This is the basic
separation between perturbative/local and global anomaly
cancellation.

\section{One-generation Standard Model local audit}

Write all fermions as left-handed Weyl fields with conventional
hypercharge:
\[
\begin{array}{c|c}
\hbox{field}&SU(3)\times SU(2)\times U(1)_Y\\ \hline
Q&(3,2)_{1/6}\\
u^c&(\bar3,1)_{-2/3}\\
d^c&(\bar3,1)_{1/3}\\
L&(1,2)_{-1/2}\\
e^c&(1,1)_{1}.
\end{array}                                             \tag{81.14}
\]
A sterile \(\nu^c\) has \(Y=0\) and contributes zero to every row
below.

\begin{theorem}[Perturbative gauge-anomaly sums]
\label{thm:v15v81-sm-anomaly}
For the representation list (81.14), the
\(SU(3)^3\), \(SU(3)^2U(1)\), \(SU(2)^2U(1)\),
\(U(1)^3\), and mixed gravitational--\(U(1)\) coefficients all
vanish generation by generation.
\end{theorem}

\begin{proof}
Normalize the quadratic index of a fundamental to \(1/2\) and the
cubic \(SU(3)\) fundamental coefficient to \(+1\).  The
contributions in field order
\((Q,u^c,d^c,L,e^c)\) are
\[
\begin{array}{c|rrrrr|r}
&Q&u^c&d^c&L&e^c&\hbox{sum}\\ \hline
SU(3)^3&2&-1&-1&0&0&0\\
SU(3)^2Y&1/6&-1/3&1/6&0&0&0\\
SU(2)^2Y&1/4&0&0&-1/4&0&0\\
Y^3&1/36&-8/9&1/9&-1/4&1&0\\
{\rm grav}^2Y&1&-2&1&-1&1&0.
\end{array}                                             \tag{81.15}
\]
For example, the cubic hypercharge numerator with denominator
\(36\) is
\[
 1-32+4-9+36=0.
\]
The other sums are immediate from the table.  There is no local
\(SU(2)^3\) coefficient because the doublet is pseudoreal.
\end{proof}

\begin{proposition}[\(SU(2)\) mod-two count]
\label{prop:v15v81-witten}
The conventional representation content has an even number of
left-handed \(SU(2)\) doublets: four per generation and twelve for
three generations.  It therefore passes the standard mod-two
doublet-count test.
\end{proposition}

\begin{proof}
The quark doublet occurs in three colours and the lepton doublet
once, giving \(3+1=4\) per generation.
\end{proof}

The arithmetic in (81.15) cancels the perturbative gauge anomaly
polynomial.  Proposition~\ref{prop:v15v81-witten} checks the
standard \(SU(2)\) global test.  These calculations do not
classify every global anomaly for every global form of the gauge
group, spin structure, boundary condition, or dynamical metric
family.

\section{What the local audit does not cancel}

\begin{proposition}[Separation of anomaly rows]
\label{prop:v15v81-separation}
The cancellations in (81.15) do not imply any of:
\[
\begin{gathered}
 c_1({\cal L}_{\rm det})=0\hbox{ on the full quotient},\\
 {\rm Hol}_{\rm det}=1\hbox{ on every mapping-torus loop},\\
 T^\mu{}_\mu=0,\qquad
 \inf_n|\mathbb E_{\mu_n}z_n|>0,\\
 \hbox{or Osterwalder--Schrader positivity}.           \tag{81.16}
\end{gathered}
\]
\end{proposition}

\begin{proof}
The table computes local gauge-anomaly coefficients from
representations.  The first two lines of (81.16) are global
topological and holonomy statements on the background quotient.
The trace anomaly depends on beta functions and curvature terms.
The average phase depends on the probability distribution of the
phase.  Reflection positivity is an inequality for a reflected
functional.  None is algebraically equal to a row of (81.15).
\end{proof}

\section{Trivial topology still permits phase cancellation}

\begin{counterexample}[Global lift with zero average phase]
\label{cnt:v15v81-average}
Let the base be the interval \([0,1]\), with uniform modulus
measure, trivial determinant line, and
\[
 \theta(x)=2\pi x,\qquad z(x)=e^{2\pi ix}.             \tag{81.17}
\]
Then \(\theta\) is a global real lift with bounded derivative, but
\[
 \int_0^1z(x)\,dx=0.                                  \tag{81.18}
\]
\end{counterexample}

\begin{proof}
The interval is contractible, and (81.17) is already a real lift.
Direct integration gives (81.18).
\end{proof}

Therefore determinant-line triviality and anomaly cancellation do
not imply the average-phase floor in V79.

\section{Reflection-real structure}

Let \(\Theta:{\cal B}\to{\cal B}\) be Euclidean time reflection.  A
reflection lift is an antilinear fibre map
\[
 J_b:{\cal L}_{{\rm det},b}
 \longrightarrow{\cal L}_{{\rm det},\Theta b},\qquad
 J_{\Theta b}J_b=1.                                   \tag{81.19}
\]

\begin{proposition}[Reflection-real phase]
\label{prop:v15v81-real}
If a unit determinant section satisfies
\[
 J_bs(b)=s(\Theta b),                                 \tag{81.20}
\]
then in compatible scalar frames its phase satisfies
\[
 z(\Theta b)=\overline{z(b)}.                         \tag{81.21}
\]
\end{proposition}

\begin{proof}
An antilinear map conjugates the scalar multiplying a unit frame.
Apply \(J_b\) to the local expression \(s(b)=z(b)e_b\) and use
(81.20).
\end{proof}

Equation (81.21) gives reflection hermiticity.  It is still weaker
than positivity of the reflected Gram form; the factorization or
exterior-power criterion of V51 remains necessary.

\section{Naturality across backgrounds}

For the locally covariant programme of V61, the determinant lines
must themselves form a functor over admissible background
embeddings, and the trivializing or reflection-real sections must
commute with those maps.  Choosing an unrelated trivialization on
each individual background can leave a categorical \(U(1)\)
cocycle and does not prove background independence.

\begin{assessment}[Physical status]
\label{ass:v15v81-status}
V81 constructs the determinant-line obstruction ledger, proves the
global phase-lifting and parallel-section criteria, identifies the
covariant phase derivative, and explicitly verifies the
conventional perturbative Standard Model anomaly sums and even
\(SU(2)\) doublet count.

The programme has not classified the determinant-line holonomy for
the chosen global gauge group and dynamical background category,
handled all zero-mode strata, constructed a natural
reflection-real trivialization, or proved an average-phase floor.
\end{assessment}

\begin{decision}[V82 mapping-torus holonomy]
\label{dec:v15v81-next}
Turn the global holonomy row into a computable topological
criterion.  Associate a mapping torus to every gauge--diffeomorphism
loop, formulate the exponentiated eta-invariant character, and
prove that triviality on the relevant five-dimensional spin
bordism group removes flat global holonomy.  Keep boundary
zero modes and reflection-real trivialization as separate rows.
\end{decision}

\begin{firewall}[Claim boundary after V81]
\label{fire:v15v81}
V81 proves local anomaly arithmetic and global determinant-line
criteria.  It does not prove the required bordism character is
trivial, construct a natural phase, establish its average floor,
or complete reflection-positive continuum reconstruction.
\end{firewall}

\section*{Version V81 append-only record}
\addcontentsline{toc}{section}{Version V81 append-only record}
The complete V80 source was retained before this continuation.  It
had \(793068\) bytes, \(23376\) lines, and SHA--256
\[
 \texttt{67F701891BD2A0E7E38DF96FAF5294CDB8C2407A9B35AFC1F4D42CCA7782CC87}.
\]
V81 constructed the determinant-line, connection-curvature,
holonomy, phase-lift, and reflection-real ledgers and audited the
local Standard Model gauge anomalies without promoting those
cancellations to global phase or continuum claims.

\chapter{Fifteenth-cycle continuation V82: mapping tori and the global anomaly character}
\label{sec:v15v82}

\section{From a quotient loop to a five-manifold}

Let \(\gamma\) be a loop in the regular gauge--metric quotient.
Choose representatives \(b_t\), \(0\leq t\leq1\), whose endpoints
are related by a gauge--diffeomorphism \(\varphi\).  The associated
mapping torus is
\[
 Y_\gamma
 =
 (M\times[0,1])/
 \bigl((x,1)\sim(\varphi(x),0)\bigr),                 \tag{82.1}
\]
with the induced spin structure and gauge bundle.  It is a closed
five-dimensional background for the self-adjoint Dirac operator
\({\mathscr D}_{Y_\gamma,R}\) in the fermion representation \(R\).

\begin{definition}[Exponentiated reduced eta invariant]
\label{def:v15v82-eta}
Put
\[
 \xi({\mathscr D})
 =
 \frac{\eta({\mathscr D})+\dim\ker{\mathscr D}}{2}
 \pmod{\mathbb Z},                                    \tag{82.2}
\]
and define
\[
 \Xi_R(Y)
 =
 \exp\bigl(2\pi i\,\xi({\mathscr D}_{Y,R})\bigr).      \tag{82.3}
\]
Changing the overall sign convention complex conjugates (82.3)
and does not change the condition \(\Xi_R=1\).
\end{definition}

\begin{theorem}[Mapping-torus holonomy formula]
\label{thm:v15v82-holonomy}
After the local anomaly curvature has been cancelled by the
complete fermion representation and prescribed local
counterterms, the flat determinant-line holonomy around \(\gamma\)
is
\[
 {\rm Hol}_{\rm det}(\gamma)=\Xi_R(Y_\gamma)           \tag{82.4}
\]
in the convention (82.3).
\end{theorem}

\begin{proof}
Cut the mapping torus at \(t=0\).  Parallel transport of the
determinant line along \(b_t\) identifies the endpoint fibre with
the initial fibre by the gauge--diffeomorphism \(\varphi\).
The spectral-flow correction through zero modes is the kernel term
in (82.2), while the regularized spectral asymmetry is
\(\eta/2\).  The product of these two contributions is (82.3).
The local transgression term that normally accompanies this
formula is exactly the integrated determinant-connection curvature;
under the stated local-anomaly cancellation it is removed, leaving
(82.4).
\end{proof}

The theorem remains meaningful when the four-dimensional loop
crosses a zero-mode locus.  A scalar phase chart fails there, but
the determinant line and reduced eta invariant retain spectral
flow.

\section{Bordism descent}

Let \(G\) denote the chosen global form of the Standard Model gauge
group and write
\[
 \Omega^{\rm Spin}_5(BG)                              \tag{82.5}
\]
for the five-dimensional spin bordism group of \(G\)-bundles.

\begin{theorem}[APS bordism relation]
\label{thm:v15v82-aps}
Suppose a closed five-dimensional background \(Y\) bounds a spin
six-manifold \(Z\) to which the gauge bundle extends.  Then
\[
 \Xi_R(Y)
 =
 \exp\left(2\pi i\int_Z I_6(R)\right),                 \tag{82.6}
\]
where \(I_6(R)\) is the degree-six local anomaly polynomial, with
the same convention as (82.3).
\end{theorem}

\begin{proof}
The Atiyah--Patodi--Singer index formula for the chiral Dirac
operator on \(Z\) is
\[
 \operatorname{index}{\mathscr D}_{Z,R}
 =
 \int_ZI_6(R)-\xi({\mathscr D}_{Y,R}).                 \tag{82.7}
\]
Exponentiating \(2\pi i\) times (82.7) removes the integer index
and gives (82.6).
\end{proof}

\begin{corollary}[Global anomaly character]
\label{cor:v15v82-character}
If the total local anomaly polynomial vanishes in the required
integral sense, then \(\Xi_R(Y)\) depends only on the bordism class
\([Y]\) and defines a homomorphism
\[
 \Xi_R:\Omega^{\rm Spin}_5(BG)\longrightarrow U(1).   \tag{82.8}
\]
\end{corollary}

\begin{proof}
If \(Y_0\) and \(Y_1\) are bordant, apply
Theorem~\ref{thm:v15v82-aps} to the bordism with boundary
\(Y_1\sqcup(-Y_0)\).  Vanishing of the total \(I_6\) gives
\(\Xi_R(Y_1)=\Xi_R(Y_0)\).  Disjoint union adds eta invariants and
therefore multiplies their exponentials, proving the homomorphism
property.
\end{proof}

\section{Finite generator test}

\begin{theorem}[Bordism generator criterion]
\label{thm:v15v82-generators}
Suppose classes \(y_\alpha\) generate
\(\Omega^{\rm Spin}_5(BG)\).  If
\[
 \Xi_R(y_\alpha)=1\quad\hbox{for every generator},    \tag{82.9}
\]
then every mapping-torus loop has trivial determinant holonomy.
\end{theorem}

\begin{proof}
By Corollary~\ref{cor:v15v82-character}, \(\Xi_R\) is a group
homomorphism.  A homomorphism equal to one on a generating set is
the trivial homomorphism.  Equation (82.4) then gives trivial
holonomy on each loop.
\end{proof}

This reduces the global anomaly row to two finite or algebraic
tasks: compute the bordism group for the chosen global gauge group
and evaluate the eta character on its generators.  The group
depends on whether the gauge group is the direct product or a
central quotient such as a \(\mathbb Z_6\) form; these choices
cannot be interchanged silently.

\section{The \(SU(2)\) mod-two generator}

\begin{proposition}[Doublet parity character]
\label{prop:v15v82-su2}
On the standard \(SU(2)\) mod-two mapping-torus class, one
left-handed doublet contributes
\[
 \Xi_{\bf2}=-1.                                       \tag{82.10}
\]
For \(N_D\) doublets the combined character is
\[
 \Xi_{\bf2}^{\,N_D}=(-1)^{N_D}.                       \tag{82.11}
\]
The conventional Standard Model count \(N_D=12\) therefore gives
\(+1\) on this class.
\end{proposition}

\begin{proof}
The mod-two index of the five-dimensional pseudoreal Dirac
operator is one on the generating class, so its exponentiated eta
phase is \(-1\).  Direct sums add reduced eta invariants modulo
integers and hence multiply their phases.  Equation (82.11)
follows, and \((-1)^{12}=1\).
\end{proof}

\begin{counterexample}[Local cancellation without global cancellation]
\label{cnt:v15v82-witten}
A theory with one left-handed \(SU(2)\) doublet has no perturbative
\(SU(2)^3\) anomaly polynomial because the representation is
pseudoreal, but (82.10) gives nontrivial global holonomy.
\end{counterexample}

\begin{proof}
The local cubic invariant vanishes, while
Proposition~\ref{prop:v15v82-su2} gives \(\Xi_{\bf2}=-1\) on the
mod-two loop.
\end{proof}

This is the concrete reason the arithmetic in V81 cannot replace
the bordism test.

\section{Local counterterms and global holonomy}

\begin{proposition}[Single-valued counterterms do not change flat holonomy]
\label{prop:v15v82-counterterm}
Let a local counterterm multiply a determinant section by
\[
 e^{iC(b)}
\]
for a globally single-valued real functional \(C\).  It shifts the
local phase connection by \(dC\) but leaves every loop holonomy
unchanged.
\end{proposition}

\begin{proof}
Along a closed loop \(\gamma\), the additional connection integral
is
\[
 \oint_\gamma dC=0.
\]
Therefore the exponentiated holonomy is unchanged.
\end{proof}

A counterterm capable of changing the global character is not a
single-valued function on the quotient; its own transformation law
is part of an anomaly trivialization and must be checked globally.

\section{Naturality of the trivialization}

Let admissible background embeddings act on determinant lines as
in V61.  A vanishing character on each individual quotient
component gives parallel sections after a basepoint choice.
Background covariance further requires those basepoint choices and
sections to commute with embeddings and disjoint unions.

\begin{proposition}[Residual constant cocycle]
\label{prop:v15v82-natural}
Suppose every component has trivial curvature and holonomy.  Any
two unit parallel trivializations differ on that component by a
constant element of \(U(1)\).  Failure to choose these constants
compatibly across background embeddings is a residual categorical
anomaly.
\end{proposition}

\begin{proof}
If \(s_1,s_2\) are unit parallel sections, write
\(s_2=fs_1\) with \(f:{\cal B}\to U(1)\).  Parallelness gives
\(df=0\), so \(f\) is constant on a connected component.
Compatibility under composition of embeddings is precisely the
cocycle condition on these constants.
\end{proof}

\section{Reflection and the average phase remain independent}

The spin bordism character (82.8) concerns oriented gauge--metric
loops.  A reflection-real determinant line involves an antilinear
involution and may require an equivariant or Pin-type bordism
problem.  Triviality of (82.8) does not construct the lift (81.19).

Likewise, a trivial parallel anomaly line allows a globally
consistent scalar determinant, but its dynamical phase may still
have zero modulus expectation, as Counterexample
\ref{cnt:v15v81-average} shows.  The V79 average-phase floor remains
a separate analytic probability estimate.

\section{Global-anomaly verification ledger}

For the complete Standard Model--Einstein continuum one must:
\begin{enumerate}
\item specify the global gauge group \(G\), spin structures, and
      allowed gauge--diffeomorphism backgrounds;
\item prove the total integral anomaly polynomial cancels;
\item compute \(\Omega^{\rm Spin}_5(BG)\);
\item evaluate the combined fermion eta character on generators;
\item extend the result over zero-mode strata in the determinant
      line;
\item choose the resulting trivialization naturally across
      backgrounds;
\item solve the separate reflection-real and average-phase rows.
\end{enumerate}

\begin{assessment}[Physical status]
\label{ass:v15v82-status}
V82 turns flat determinant holonomy into a mapping-torus eta
character, proves its bordism descent under local cancellation, and
gives a finite generator criterion.  It verifies cancellation on
the standard \(SU(2)\) mod-two class using the twelve-doublet
Standard Model count.

The full five-dimensional bordism group and eta character for the
chosen global Standard Model gauge group coupled to dynamical
spin metrics have not been computed here.  Reflection-real
trivialization, zero-mode extension, and the average-phase floor
also remain open.
\end{assessment}

\begin{decision}[V83 reflection-real anomaly line]
\label{dec:v15v82-next}
Construct the reflection-equivariant determinant-line problem.
Double a positive-time background across its boundary, formulate
the antilinear sewing map and its possible sign character, and
prove sufficient conditions under which the phase cancels from the
reflected Gram form.  Keep reflection hermiticity distinct from
strict positivity and zero-mode saturation.
\end{decision}

\begin{firewall}[Claim boundary after V82]
\label{fire:v15v82}
V82 proves the mapping-torus and bordism compilers and one
\(SU(2)\) generator cancellation.  It does not complete the full
bordism calculation, natural anomaly trivialization, reflection
sewing, average-phase estimate, or continuum construction.
\end{firewall}

\section*{Version V82 append-only record}
\addcontentsline{toc}{section}{Version V82 append-only record}
The complete V81 source was retained before this continuation.  It
had \(805649\) bytes, \(23763\) lines, and SHA--256
\[
 \texttt{AC7C5843FBD67EBE38A71578797D9080C83C5CB9D63BE93929F91613E257F045}.
\]
V82 expressed global determinant holonomy by mapping-torus eta
invariants, descended it to a five-dimensional spin-bordism
character, and isolated the remaining generator, naturality,
reflection, zero-mode, and average-phase checks.

\chapter{Fifteenth-cycle continuation V83: reflection-real sewing of the determinant line}
\label{sec:v15v83}

\section{Doubled backgrounds}

Let \(X_+\) be a compact Euclidean positive-time region with
boundary \(\Sigma\).  Its reflected copy is \(X_-\), and the doubled
background is
\[
 X=X_-\cup_\Sigma X_+.                               \tag{83.1}
\]
Let \(\Theta\) exchange the two halves.  Denote by
\({\cal L}_+\) the fermion determinant/Pfaffian line associated
with \(X_+\) and its boundary data, and by \({\cal L}_X\) the line
on the double.

\begin{definition}[Reflection-real sewing data]
\label{def:v15v83-sewing}
Reflection-real sewing consists of:
\begin{enumerate}
\item an antilinear isometry
\[
 J:{\cal L}_+\longrightarrow{\cal L}_-;
                                                               \tag{83.2}
\]
\item a bilinear gluing isomorphism
\[
 {\rm Sew}:{\cal L}_-\otimes{\cal L}_+
 \longrightarrow{\cal L}_X;                           \tag{83.3}
\]
\item a boundary contraction for which
\[
 {\rm Sew}(Jv\otimes v)=\|v\|_\Sigma^2\,{\bf1}_X      \tag{83.4}
\]
in a chosen positive real ray of \({\cal L}_X\).
\end{enumerate}
All maps must be natural under boundary-preserving gauge and
metric embeddings.
\end{definition}

Equation (83.4) is stronger than the phase-conjugation relation
(81.21): it fixes the sewing sign and the positive ray.

\section{Factorization theorem}

Let \(y\) denote boundary fields.  Suppose the positive modulus
measure on the double disintegrates as
\[
 d\mu_X(x_-,x_+,y)
 =
 d\mu_\Sigma(y)\,
 d\mu_-(x_-|y)\,
 d\mu_+(x_+|y),                                      \tag{83.5}
\]
with reflection carrying the plus conditional measure to the
minus conditional measure.  Let
\[
 A_+(x_+,y)
\]
be the full plus-side amplitude, valued in the boundary fermion
line or its exterior-algebra completion.

\begin{theorem}[Reflected norm-square criterion]
\label{thm:v15v83-norm-square}
Assume reflection-real sewing and the factorization
\[
 W_X(x_-,x_+,y)
 =
 \overline{A_+(\Theta x_-,y)}\,A_+(x_+,y)             \tag{83.6}
\]
after the boundary contraction in (83.4).  Then every bounded
observable \(F\) supported in \(X_+\) satisfies
\[
 \boxed{
 \langle(\Theta F)F\rangle_X
 =
 \int
 \left|
 \int F(x_+)A_+(x_+,y)d\mu_+(x_+|y)
 \right|^2d\mu_\Sigma(y)
 \geq0.}                                              \tag{83.7}
\]
\end{theorem}

\begin{proof}
Insert (83.5)--(83.6) into the reflected expectation.  Reflection
turns the minus-side integral into the complex conjugate of the
plus-side integral:
\[
\begin{split}
 \langle(\Theta F)F\rangle_X
 &=
 \int d\mu_\Sigma(y)
 \left[
 \int\overline{F(x_-)A_+(\Theta x_-,y)}
 d\mu_-(x_-|y)
 \right]\\
 &\qquad\qquad\times
 \left[
 \int F(x_+)A_+(x_+,y)d\mu_+(x_+|y)
 \right].
\end{split}
\]
The first bracket is the conjugate of the second.  This is (83.7).
\end{proof}

\begin{corollary}[Positive finite-cutoff OS form]
\label{cor:v15v83-os}
The null space
\[
 {\cal N}_+=
 \{F:\langle(\Theta F)F\rangle_X=0\}                  \tag{83.8}
\]
is a linear subspace, and quotienting the positive-time observable
algebra by \({\cal N}_+\), followed by completion, gives a Hilbert
space.
\end{corollary}

\begin{proof}
Equation (83.7) realizes the sesquilinear form as an \(L^2\) inner
product of boundary amplitudes.  Its null space is the kernel of a
linear map.  The standard quotient and completion of a positive
semidefinite form is a Hilbert space.
\end{proof}

This theorem refines the abstract exterior-power positivity
criterion of V51 by locating the necessary phase cancellation in
the determinant-line sewing map.

\section{Hermiticity is weaker than positivity}

\begin{counterexample}[Reflection-conjugate kernel with a negative direction]
\label{cnt:v15v83-hermitian}
On a two-point positive-time configuration space, let the reflected
kernel be
\[
 K=
 \begin{pmatrix}
 1&2\\
 2&1
 \end{pmatrix}.                                      \tag{83.9}
\]
It is real symmetric and therefore satisfies reflection
hermiticity.  For \(f=(1,-1)\),
\[
 f^*Kf=-2<0.                                          \tag{83.10}
\]
\end{counterexample}

\begin{proof}
The eigenvalues of \(K\) are \(3\) and \(-1\); the displayed vector
lies in the negative eigenspace.
\end{proof}

Thus \(z(\Theta b)=\overline{z(b)}\) can make the reflected form
Hermitian while leaving a negative eigenvalue.  The norm-square
factorization (83.6) is the relevant positivity statement.

\section{Sewing sign and reflection anomaly}

Suppose the gluing law instead gives
\[
 {\rm Sew}(Jv\otimes v)
 =
 \chi(b)\|v\|_\Sigma^2{\bf1}_X,\qquad
 \chi(b)\in U(1).                                     \tag{83.11}
\]

\begin{proposition}[Positive-ray obstruction]
\label{prop:v15v83-sign}
If (83.11) is to be nonnegative for every nonzero \(v\), then
\[
 \chi(b)=1.                                           \tag{83.12}
\]
A locally constant value \(\chi=-1\) is a reflection anomaly that
cannot be removed by rescaling \(v\) by a phase.
\end{proposition}

\begin{proof}
The value in (83.11) must lie in the positive real ray for every
\(\|v\|^2>0\), forcing \(\chi=1\).  Replacing \(v\) by
\(e^{i\alpha}v\) changes \(Jv\) by \(e^{-i\alpha}\) and the second
factor by \(e^{i\alpha}\); the phases cancel, so \(\chi\) is
unchanged.
\end{proof}

The sign character is equivariant data and is not detected by the
oriented mapping-torus character alone.  Its classification may
require reflection or Pin bordism, depending on the chosen
Euclidean reflection structure.

\section{Zero modes and exterior-algebra amplitudes}

If the half-space Dirac operator has zero modes, the scalar
determinant amplitude vanishes.  Let
\[
 {\cal F}_\Sigma=\Lambda^\bullet{\cal H}_\Sigma       \tag{83.13}
\]
be the boundary fermion exterior algebra and interpret \(A_+\) as
a vector in \({\cal F}_\Sigma\).  Fermion insertions saturate the
appropriate zero-mode degree.

\begin{proposition}[Zero-mode-compatible positivity]
\label{prop:v15v83-zero}
If reflection acts antiunitarily on \({\cal F}_\Sigma\) and sewing
uses its Hilbert inner product, then
\[
 {\rm Sew}(JA_+\otimes A_+)=\|A_+\|_{{\cal F}_\Sigma}^2\geq0.    \tag{83.14}
\]
Consequently Theorem~\ref{thm:v15v83-norm-square} remains valid for
zero-mode-saturating observables.
\end{proposition}

\begin{proof}
Expand \(A_+\) in an orthonormal exterior-power basis.  Antiunitary
reflection conjugates its coefficients, and the contraction sums
their squared absolute values.  This is (83.14), degree by degree.
\end{proof}

The unsaturated vacuum amplitude may be zero.  Positivity then
produces a null vector, not a normalized vacuum state; a nonzero
normalization sector must still be identified.

\section{Gauge quotient and boundary faces}

For (83.5) to hold on a gauge quotient:
\begin{enumerate}
\item the gauge condition must be reflection covariant;
\item the Faddeev--Popov determinant and slice Jacobian must split
      or sew into conjugate half factors;
\item gauge-copy ownership faces must cancel as in V64 and V76;
\item residual boundary gauge transformations on \(\Sigma\) must
      be integrated in the boundary measure;
\item Gribov score singularities must have integrable reflected
      traces.
\end{enumerate}
The determinant-line sewing theorem does not remove these bosonic
and quotient obligations.

\section{Cofinal preservation}

\begin{theorem}[Limit of reflected positive forms]
\label{thm:v15v83-limit}
Let \({\cal A}_+\) be a countable positive-time cylinder algebra.
Suppose for every \(F\in{\cal A}_+\),
\[
 \langle(\Theta F)F\rangle_n\geq0                     \tag{83.15}
\]
and the expectations converge:
\[
 \langle(\Theta F)F\rangle_n
 \longrightarrow
 \langle(\Theta F)F\rangle.                           \tag{83.16}
\]
Then the limiting form is positive semidefinite on
\({\cal A}_+\).
\end{theorem}

\begin{proof}
The limit of the nonnegative real numbers in (83.15) is
nonnegative.  Polarization gives the limiting sesquilinear form on
the algebra.
\end{proof}

Convergence in (83.16) requires the phase-reweighting and
uniform-integrability conditions of V79 unless the reflected
amplitude construction defines the normalized measure directly.

\section{Reflection closure ledger}

The fermion reflection row now consists of:
\begin{enumerate}
\item an antilinear lift \(J\) with \(J^2=1\);
\item a sewing map with positive-ray sign \(\chi=1\);
\item norm-square factorization of the full half amplitudes;
\item exterior-algebra treatment of zero modes;
\item nonzero normalization in the chosen sector;
\item quotient-density and boundary compatibility;
\item cofinal convergence on a determining positive-time algebra.
\end{enumerate}

\begin{assessment}[Physical status]
\label{ass:v15v83-status}
V83 proves a determinant-line sewing criterion that turns the
reflected functional into an explicit boundary-amplitude norm
square.  It separates reflection hermiticity from positivity,
isolates a sewing-sign anomaly, and extends the construction to
zero-mode-saturating exterior-power sectors.

The physical Standard Model regulator has not yet been shown to
possess the required natural reflection lift, positive sewing
sign, quotient factorization, nonzero vacuum sector, or cofinal
reflected convergence.
\end{assessment}

\begin{decision}[V84 reflected normalization and vacuum sector]
\label{dec:v15v83-next}
Analyze the nonzero-normalization row.  Express the finite-slab
partition function as the squared norm of a boundary vacuum
amplitude, derive a uniform lower bound from a visible vacuum
component, and show how zero-mode parity or a vanishing overlap can
destroy normalization even though the reflected form remains
positive semidefinite.
\end{decision}

\begin{firewall}[Claim boundary after V83]
\label{fire:v15v83}
V83 proves conditional reflection-positive sewing.  It does not
construct the reflection lift or quotient factorization, prove a
nonzero normalized vacuum, control the average phase, or complete
the continuum.
\end{firewall}

\section*{Version V83 append-only record}
\addcontentsline{toc}{section}{Version V83 append-only record}
The complete V82 source was retained before this continuation.  It
had \(816790\) bytes, \(24076\) lines, and SHA--256
\[
 \texttt{5F46D651B25B298576C5E07A1FBEAAC3C545509EBCFAAF06EC421D9305D8CB7E}.
\]
V83 constructed reflection-real determinant-line sewing, proved
the reflected norm-square formula, separated hermiticity from
positivity, and incorporated sewing signs, zero modes, quotient
faces, and cofinal limits.

\chapter{Fifteenth-cycle continuation V84: reflected normalization and the visible vacuum sector}
\label{sec:v15v84}

\section{Boundary amplitude and unnormalized partition function}

Under the sewing hypotheses of V83, let
\[
 {\cal T}_n:{\cal A}_{+,n}\longrightarrow{\cal H}_{\Sigma,n}    \tag{84.1}
\]
be the boundary-amplitude map, so that
\[
 {\cal B}_n(F,G)
 =
 \langle{\cal T}_nF,{\cal T}_nG\rangle_{\Sigma,n}.     \tag{84.2}
\]
Put
\[
 \Psi_n={\cal T}_n1.                                  \tag{84.3}
\]

\begin{proposition}[Reflected partition norm]
\label{prop:v15v84-partition}
The unnormalized doubled partition function is
\[
 Z_n^{\rm ref}={\cal B}_n(1,1)=\|\Psi_n\|^2.           \tag{84.4}
\]
It is nonnegative, and it is strictly positive exactly when
\(\Psi_n\ne0\).
\end{proposition}

\begin{proof}
Set \(F=G=1\) in (84.2).  A Hilbert norm is nonnegative and
vanishes precisely on the zero vector.
\end{proof}

When \(Z_n^{\rm ref}>0\), the normalized reflected state is
\[
 \omega_n((\Theta F)G)
 =
 \frac{{\cal B}_n(F,G)}{Z_n^{\rm ref}}.                \tag{84.5}
\]
Positivity of \({\cal B}_n\) does not by itself ensure that the
denominator stays away from zero cofinally.

\section{Visible-component lower bound}

\begin{theorem}[Uniform visible vacuum]
\label{thm:v15v84-visible}
Let \(P_{0,n}\) be an orthogonal projection onto the physical
vacuum normalization sector of
\({\cal H}_{\Sigma,n}\).  If
\[
 \|P_{0,n}\Psi_n\|\geq c_*>0                          \tag{84.6}
\]
uniformly in \(n\), then
\[
 Z_n^{\rm ref}\geq c_*^2.                             \tag{84.7}
\]
It is sufficient that there be unit vectors
\(\Omega_n\in\operatorname{Ran}P_{0,n}\) satisfying
\[
 |\langle\Omega_n,\Psi_n\rangle|\geq c_*.             \tag{84.8}
\]
\end{theorem}

\begin{proof}
Orthogonal projection decreases norm, so
\[
 \|\Psi_n\|\geq\|P_{0,n}\Psi_n\|\geq c_*.
\]
This proves (84.7).  For (84.8), Cauchy's inequality gives
\[
 |\langle\Omega_n,\Psi_n\rangle|
 =
 |\langle\Omega_n,P_{0,n}\Psi_n\rangle|
 \leq\|P_{0,n}\Psi_n\|,
\]
and the first part applies.
\end{proof}

The condition is a lower bound on one transported physical
component, not merely an upper bound on the full boundary norm.

\section{Transfer semigroup interpretation}

Suppose a reflected slab of half-width \(T\) has positive transfer
Hamiltonian \(H_n\geq0\), vacuum projection \(P_{0,n}\), and
boundary seed \(b_n\).  Then
\[
 \Psi_{n,T}=e^{-TH_n}b_n.                             \tag{84.9}
\]

\begin{proposition}[Vacuum overlap persists along the slab]
\label{prop:v15v84-semigroup}
If
\[
 H_nP_{0,n}=0,\qquad
 |\langle\Omega_n,b_n\rangle|\geq c_*                 \tag{84.10}
\]
for a unit vacuum vector \(\Omega_n\), then
\[
 \|\Psi_{n,T}\|^2\geq c_*^2                           \tag{84.11}
\]
for every \(T\geq0\).  If additionally
\[
 H_n\geq m(I-P_{0,n}),\qquad m>0,                     \tag{84.12}
\]
then
\[
 \left|
 \|\Psi_{n,T}\|^2-\|P_{0,n}b_n\|^2
 \right|
 \leq e^{-2mT}\|(I-P_{0,n})b_n\|^2.                  \tag{84.13}
\]
\end{proposition}

\begin{proof}
The vacuum component is unchanged:
\[
 P_{0,n}e^{-TH_n}b_n=P_{0,n}b_n.
\]
Apply Theorem~\ref{thm:v15v84-visible}.  Under (84.12), the spectral
theorem gives
\[
 \|e^{-TH_n}(I-P_{0,n})b_n\|
 \leq e^{-mT}\|(I-P_{0,n})b_n\|.
\]
Vacuum and orthogonal components are perpendicular, so squaring
gives (84.13).
\end{proof}

This is the normalization counterpart of the semigroup gap
transport in V52--V54.

\section{Vanishing vacuum overlap}

\begin{counterexample}[Positive slabs with collapsing normalization]
\label{cnt:v15v84-orthogonal}
Let
\[
 H=
 \begin{pmatrix}0&0\\0&m\end{pmatrix},\qquad
 P_0=
 \begin{pmatrix}1&0\\0&0\end{pmatrix},\qquad
 b=
 \begin{pmatrix}0\\1\end{pmatrix}.                   \tag{84.14}
\]
Then every reflected norm is positive,
\[
 \|e^{-TH}b\|^2=e^{-2mT}>0,                           \tag{84.15}
\]
but it tends to zero as \(T\to\infty\).
\end{counterexample}

\begin{proof}
Direct matrix exponentiation gives (84.15), while \(P_0b=0\).
\end{proof}

A spectral gap suppresses excited components; it cannot create a
vacuum component that is absent from the boundary seed.

\section{Fermion parity and zero-mode sectors}

Let the boundary Fock space split into parity sectors,
\[
 {\cal H}_{\Sigma,n}
 =
 {\cal H}^{\rm even}_{\Sigma,n}
 \oplus{\cal H}^{\rm odd}_{\Sigma,n}.                 \tag{84.16}
\]
Suppose the scalar vacuum normalization uses the even projection
\(P_{\rm even}\).

\begin{proposition}[Parity obstruction]
\label{prop:v15v84-parity}
If zero-mode selection rules force
\[
 \Psi_n\in{\cal H}^{\rm odd}_{\Sigma,n},               \tag{84.17}
\]
then
\[
 P_{\rm even}\Psi_n=0.                                \tag{84.18}
\]
The full reflected form on the exterior algebra may remain positive
semidefinite, while the selected scalar vacuum amplitude vanishes.
\end{proposition}

\begin{proof}
The two parity subspaces in (84.16) are orthogonal.  Therefore the
orthogonal projection of an odd vector onto the even subspace is
zero.  Positivity of the full form follows independently from the
norm-square representation in V83.
\end{proof}

Insertions that saturate the zero modes may select a nonzero odd or
higher-degree sector.  Such a sector defines a charged or inserted
state, not automatically the vacuum state requested by the
uninserted continuum theory.

\section{Relation to the average phase}

Let \(Z_n^{\rm mod}>0\) be the modulus partition function and
\(Z_n^{\rm full}\) the full complex one in the same normalization.
Then
\[
 a_n=\mathbb E_{\mu_n}z_n
 =
 \frac{Z_n^{\rm full}}{Z_n^{\rm mod}}.                \tag{84.19}
\]
Under reflection-real sewing with positive sign,
\[
 Z_n^{\rm full}=Z_n^{\rm ref}=\|\Psi_n\|^2.           \tag{84.20}
\]

\begin{theorem}[Visible vacuum implies a phase floor with modulus control]
\label{thm:v15v84-phase-floor}
If (84.6) holds and
\[
 Z_n^{\rm mod}\leq C_{\rm mod}<\infty                 \tag{84.21}
\]
uniformly in the common normalization, then
\[
 a_n\geq\frac{c_*^2}{C_{\rm mod}}>0.                  \tag{84.22}
\]
\end{theorem}

\begin{proof}
Use (84.7), (84.19), and (84.20).  The reflected full partition
function is nonnegative, so no absolute-value ambiguity remains.
\end{proof}

A lower reflected norm without the modulus upper bound does not
give a lower bound on their ratio.  Conversely, a phase floor can
hold without identifying a particular vacuum component.

\section{Cofinal normalized limit}

\begin{theorem}[Normalized reflected convergence]
\label{thm:v15v84-limit}
Suppose on a common countable positive-time algebra
\[
 {\cal B}_n(F,G)\longrightarrow{\cal B}(F,G)           \tag{84.23}
\]
for every \(F,G\), and
\[
 Z_n^{\rm ref}={\cal B}_n(1,1)\geq c_*^2.             \tag{84.24}
\]
Then
\[
 Z_n^{\rm ref}\to Z^{\rm ref}={\cal B}(1,1)
 \geq c_*^2,                                          \tag{84.25}
\]
and
\[
 \frac{{\cal B}_n(F,G)}{Z_n^{\rm ref}}
 \longrightarrow
 \frac{{\cal B}(F,G)}{Z^{\rm ref}}.                   \tag{84.26}
\]
The limiting normalized form is positive semidefinite.
\end{theorem}

\begin{proof}
Use (84.23) with \(F=G=1\) to get convergence of the denominators
and pass (84.24) to the limit.  Quotient convergence gives
(84.26).  Positivity follows from V83's limiting-form theorem.
\end{proof}

\section{Normalization ledger}

The reflected vacuum normalization now requires:
\begin{enumerate}
\item a specified physical boundary sector \(P_{0,n}\);
\item a transported visible component satisfying (84.6);
\item compatibility with zero-mode parity and saturation;
\item cofinal convergence of the unnormalized reflected forms;
\item for the V79 reweighting route, a common normalization and
      the modulus upper bound (84.21).
\end{enumerate}

\begin{assessment}[Physical status]
\label{ass:v15v84-status}
V84 expresses the reflected partition function as a boundary norm,
proves a uniform lower bound from a visible vacuum component,
relates that component to transfer-semigroup decay, and separates
full positivity from parity-sector normalization.  It also derives
an average-phase floor when the modulus partition function is
uniformly bounded.

The physical regulator has not yet supplied the visible component,
sector choice through zero modes, modulus upper bound, or cofinal
boundary-amplitude convergence.
\end{assessment}

\begin{decision}[V85 cross-program audit]
\label{dec:v15v84-next}
Perform the scheduled companion audit.  Check the newest
Yang--Mills note for a new transfer/vacuum or reflection-positive
result and the newest Navier--Stokes note for a normalization or
cofinal compactness input.  Record the result, then continue in V86
with a cofinal boundary-amplitude compactness theorem.
\end{decision}

\begin{firewall}[Claim boundary after V84]
\label{fire:v15v84}
V84 proves conditional nonzero-normalization implications.  It
does not construct the physical vacuum component, settle zero-mode
selection, bound the modulus partition, or complete reflected
continuum reconstruction.
\end{firewall}

\section*{Version V84 append-only record}
\addcontentsline{toc}{section}{Version V84 append-only record}
The complete V83 source was retained before this continuation.  It
had \(827312\) bytes, \(24401\) lines, and SHA--256
\[
 \texttt{70DD89D69A5CDCBC1599EB60E42A0B9EFBFC699AFD894667AA7B356416821CC6}.
\]
V84 proved reflected boundary-norm normalization, a uniform visible
vacuum lower bound, transfer-semigroup persistence, parity-sector
obstructions, and a conditional bridge to the V79 average-phase
floor.

\chapter{Fifteenth-cycle continuation V85: scheduled vacuum and normalization audit}
\label{sec:v15v85}

\section{Companion checkpoints}

The active companion sources are
\[
\begin{array}{c|c|c|c|l}
\hbox{programme}&\hbox{version}&\hbox{bytes}&\hbox{lines}&\hbox{SHA--256}\\ \hline
\hbox{Yang--Mills}&V40&344764&9458&
\texttt{579A346B3193F33B2E082C6682E83271843B665AC950E70D7F7707FE19EAF7E2}\\
\hbox{Navier--Stokes}&V26&278283&5319&
\texttt{82C887F8C2C69E4F28FAD7C3DC8DE5B9099CB5A4BF431EBA1FFF3E91935978AD}.
\end{array}                                             \tag{85.1}
\]
Both hashes equal those in V80.

\section{Targeted normalization comparison}

YM V40 proves a larger finite Wilson-pencil prism and retains exact
negative witnesses for two false isotropic charts.  It constructs
no Euclidean transfer semigroup, reflected boundary Hilbert space,
vacuum projection, boundary seed, or interacting partition
function.

\begin{proposition}[No V84 vacuum import from YM V40]
\label{prop:v15v85-ym}
YM V40 supplies neither the visible component (84.6) nor the
modulus partition bound (84.21).
\end{proposition}

\begin{proof}
Both quantities are norms or normalizations of an interacting
boundary amplitude.  The terminal YM theorem is a pointwise
finite-matrix inequality in momentum variables and explicitly
leaves the interacting reflection-positive measure open.
\end{proof}

NS V26 proves pointwise Oseen-kernel tensor norms and states that
global regularity remains open.  It has no reflected Euclidean
boundary amplitude, fermion parity sector, or complex phase.

\begin{proposition}[No V84 normalization import from NS V26]
\label{prop:v15v85-ns}
NS V26 supplies none of (84.6), (84.21), or the cofinal reflected
form convergence (84.23).
\end{proposition}

\begin{proof}
The three objects belong to a Euclidean quantum-field
reconstruction.  They are absent from the deterministic parabolic
kernel theorem in NS V26.
\end{proof}

\section{Audit conclusion}

\begin{theorem}[V85 dependency state]
\label{thm:v15v85-conclusion}
No companion theorem populates the reflected normalization row.
The next noncircular task remains compactness and convergence of
the boundary-amplitude vectors themselves.
\end{theorem}

\begin{proof}
Equation (85.1) shows that no new companion source exists, and
Propositions~\ref{prop:v15v85-ym} and
\ref{prop:v15v85-ns} rule out the targeted imports.
\end{proof}

\begin{assessment}[Physical status]
\label{ass:v15v85-status}
V85 is a dependency audit.  It adds no vacuum, normalization,
reflection-positive, or continuum existence theorem.
\end{assessment}

\begin{decision}[V86 boundary-amplitude compactness]
\label{dec:v15v85-next}
Place the changing boundary Hilbert spaces in a common carrier.
Prove weak compactness from a uniform norm bound, strong compactness
from a finite-rank tail criterion, convergence of reflected Gram
forms, and persistence of a visible vacuum component.  Exhibit an
escaping orthonormal sequence showing why bounded norms alone are
insufficient.
\end{decision}

\begin{firewall}[Claim boundary after V85]
\label{fire:v15v85}
V85 proves no new physical construction; all reflected
normalization hypotheses of V84 remain to be populated.
\end{firewall}

\section*{Version V85 append-only record}
\addcontentsline{toc}{section}{Version V85 append-only record}
The complete V84 source was retained before this continuation.  It
had \(836853\) bytes, \(24723\) lines, and SHA--256
\[
 \texttt{7F4F75544EA059881826B3FF35C84C0EA11470C7945611C18D1C2A119BAEBEB7}.
\]
V85 performed the scheduled companion audit and retained
boundary-amplitude compactness as the next substantive step.

\chapter{Fifteenth-cycle continuation V86: compactness of boundary amplitudes}
\label{sec:v15v86}

\section{Common boundary carrier}

Let \({\cal H}_{\Sigma,n}\) be the finite-cutoff boundary Hilbert
spaces.  Assume isometric identifications
\[
 J_n:{\cal H}_{\Sigma,n}\longrightarrow{\cal H}_\Sigma
                                                               \tag{86.1}
\]
into one separable Hilbert space.  We suppress \(J_n\) and regard
the reflected boundary amplitudes
\[
 \Psi_n={\cal T}_n1                                  \tag{86.2}
\]
as vectors in \({\cal H}_\Sigma\).

\begin{proposition}[Weak compactness]
\label{prop:v15v86-weak}
If
\[
 \sup_n\|\Psi_n\|\leq C,                              \tag{86.3}
\]
then a subsequence converges weakly in
\({\cal H}_\Sigma\).
\end{proposition}

\begin{proof}
A Hilbert space is reflexive, so its closed bounded balls are weakly
sequentially compact.
\end{proof}

Weak convergence does not imply convergence of the partition norms
\(\|\Psi_n\|^2\).

\section{Finite-rank tail criterion}

Let \(P_R\) be increasing finite-rank orthogonal projections on
\({\cal H}_\Sigma\) satisfying
\[
 P_R\longrightarrow I\quad\hbox{strongly}.             \tag{86.4}
\]

\begin{theorem}[Boundary-amplitude precompactness]
\label{thm:v15v86-tail}
Suppose (86.3) and
\[
 \lim_{R\to\infty}
 \sup_n\|(I-P_R)\Psi_n\|=0.                           \tag{86.5}
\]
Then \((\Psi_n)\) is relatively compact in
\({\cal H}_\Sigma\).  Every weakly convergent subsequence therefore
converges strongly.
\end{theorem}

\begin{proof}
Fix \(\epsilon>0\) and choose \(R\) so the left side of (86.5) is
below \(\epsilon/3\).  The bounded set
\(\{P_R\Psi_n\}\) lies in a finite-dimensional space and has a
finite \(\epsilon/3\)-net.  Adding the two tails shows that the same
centres form an \(\epsilon\)-net for \(\{\Psi_n\}\).  Thus the
family is totally bounded, hence relatively compact.  Strong and
weak limits agree, so a weakly convergent subsequence is strong.
\end{proof}

\begin{proposition}[Weak convergence plus norm convergence]
\label{prop:v15v86-norm}
If
\[
 \Psi_n\rightharpoonup\Psi,\qquad
 \|\Psi_n\|\to\|\Psi\|,                               \tag{86.6}
\]
then \(\Psi_n\to\Psi\) strongly.
\end{proposition}

\begin{proof}
Expand
\[
 \|\Psi_n-\Psi\|^2
 =
 \|\Psi_n\|^2+\|\Psi\|^2
 -2\operatorname{Re}\langle\Psi_n,\Psi\rangle.
\]
The first term converges by hypothesis and the inner product by
weak convergence.
\end{proof}

\section{Spectral sufficient conditions}

Let \(K\geq0\) be a self-adjoint operator on the common carrier
with compact resolvent, and let \(P_R\) be its spectral projection
on \([0,R]\).

\begin{proposition}[Energy tail]
\label{prop:v15v86-energy}
If some \(s>0\) satisfies
\[
 \sup_n\|K^{s/2}\Psi_n\|\leq C_s,                     \tag{86.7}
\]
then
\[
 \|(I-P_R)\Psi_n\|^2
 \leq R^{-s}C_s^2.                                    \tag{86.8}
\]
Hence Theorem~\ref{thm:v15v86-tail} applies.
\end{proposition}

\begin{proof}
On the spectral subspace \((I-P_R){\cal H}_\Sigma\),
\(K^s\geq R^sI\).  Therefore
\[
 R^s\|(I-P_R)\Psi_n\|^2
 \leq
 \|K^{s/2}\Psi_n\|^2.
\]
Use (86.7).
\end{proof}

\begin{proposition}[Uniform semigroup smoothing]
\label{prop:v15v86-semigroup}
Suppose
\[
 \Psi_{n,T}=e^{-TH_n}b_n,\qquad
 \sup_n\|b_n\|\leq C_b,                               \tag{86.9}
\]
and common finite-rank projections \(P_R\) reduce \(H_n\) with
\[
 H_n\geq E_R(I-P_R),\qquad E_R\to\infty               \tag{86.10}
\]
uniformly in \(n\).  Then for every fixed \(T>0\),
\[
 \sup_n\|(I-P_R)\Psi_{n,T}\|
 \leq C_be^{-TE_R}\longrightarrow0.                  \tag{86.11}
\]
\end{proposition}

\begin{proof}
The spectral theorem and (86.10) give
\[
 \|(I-P_R)e^{-TH_n}\|\leq e^{-TE_R}.
\]
Apply this to \(b_n\).
\end{proof}

Condition (86.10) is a collective high-energy compactness
statement.  Strong convergence of each fixed low-energy spectral
projection does not imply it.

\section{Visible vacuum under changing carriers}

Let \(\Omega_n\) be unit vectors in the transported vacuum sectors.

\begin{theorem}[Persistence of visibility]
\label{thm:v15v86-visible}
Suppose
\[
 \Psi_n\rightharpoonup\Psi,\qquad
 \sup_n\|\Psi_n\|<\infty,\qquad
 \Omega_n\to\Omega\quad\hbox{strongly},                \tag{86.12}
\]
and
\[
 |\langle\Omega_n,\Psi_n\rangle|\geq c_*>0.           \tag{86.13}
\]
After a subsequence on which the complex overlaps converge,
\[
 |\langle\Omega,\Psi\rangle|\geq c_*,
 \qquad
 \|\Psi\|\geq c_*.                                    \tag{86.14}
\]
\end{theorem}

\begin{proof}
The difference
\[
 \langle\Omega_n,\Psi_n\rangle-\langle\Omega,\Psi\rangle
 =
 \langle\Omega_n-\Omega,\Psi_n\rangle
 +\langle\Omega,\Psi_n-\Psi\rangle
\]
tends to zero: the first term by strong convergence and the uniform
norm bound, the second by weak convergence.  Pass (86.13) to the
limit and use Cauchy's inequality.
\end{proof}

Visibility prevents the weak limit from vanishing, but does not
prevent additional norm from escaping orthogonally.

\section{Escaping boundary modes}

\begin{counterexample}[Bounded norms with zero weak limit]
\label{cnt:v15v86-escape}
For an orthonormal basis \((e_n)\) of an infinite-dimensional
Hilbert space, set
\[
 \Psi_n=e_n.                                          \tag{86.15}
\]
Then \(\|\Psi_n\|=1\), but
\[
 \Psi_n\rightharpoonup0
\]
and (86.5) fails for every fixed finite-rank exhaustion.
\end{counterexample}

\begin{proof}
Every fixed vector has Fourier coefficients tending to zero, which
gives weak convergence.  For a finite-rank coordinate projection
\(P_R\), all sufficiently large \(e_n\) lie in its orthogonal
complement, so the supremum tail equals one.
\end{proof}

\begin{counterexample}[Visible vacuum with hidden escaping norm]
\label{cnt:v15v86-visible-escape}
Fix a unit vector \(\Omega\) and an orthonormal sequence
\((e_n)\subset\Omega^\perp\).  Let
\[
 \Psi_n=\Omega+e_n.                                   \tag{86.16}
\]
Then
\[
 \langle\Omega,\Psi_n\rangle=1,\qquad
 \Psi_n\rightharpoonup\Omega,
\]
but
\[
 \|\Psi_n\|^2=2\not\longrightarrow\|\Omega\|^2=1.      \tag{86.17}
\]
\end{counterexample}

\begin{proof}
Orthogonality gives the overlap and norm.  Weak convergence follows
from \(e_n\rightharpoonup0\).
\end{proof}

Thus V84's visible component secures nonzero normalization but a
tail estimate is still needed to identify its limiting numerical
value.

\section{Convergence of reflected Gram forms}

Let \({\cal A}_+\) be a countable determining positive-time
algebra, and transport
\[
 {\cal T}_nF\in{\cal H}_\Sigma.                        \tag{86.18}
\]

\begin{theorem}[Strong-amplitude reconstruction]
\label{thm:v15v86-gram}
Suppose for every \(F\in{\cal A}_+\),
\[
 {\cal T}_nF\longrightarrow{\cal T}F
 \quad\hbox{strongly in }{\cal H}_\Sigma.              \tag{86.19}
\]
Then
\[
 {\cal B}_n(F,G)
 =
 \langle{\cal T}_nF,{\cal T}_nG\rangle
 \longrightarrow
 \langle{\cal T}F,{\cal T}G\rangle
 =:{\cal B}(F,G).                                     \tag{86.20}
\]
If
\[
 \|{\cal T}1\|\geq c_*>0,                             \tag{86.21}
\]
the normalized reflected forms converge and the limiting form is
positive semidefinite.
\end{theorem}

\begin{proof}
Strong convergence of both vectors gives convergence of their
inner products.  Take \(F=G=1\) to obtain convergence to the
positive denominator in (86.21), then apply
Theorem~\ref{thm:v15v84-limit}.  Positivity is also immediate from
the final inner-product representation.
\end{proof}

\begin{corollary}[Spectral-tail compiler for the cylinder algebra]
\label{cor:v15v86-algebra-tail}
Assume for every \(F\in{\cal A}_+\):
\[
 \sup_n\|{\cal T}_nF\|<\infty,\qquad
 \lim_{R\to\infty}\sup_n
 \|(I-P_R){\cal T}_nF\|=0,                            \tag{86.22}
\]
and every finite set of coordinates
\(P_R{\cal T}_nF\) converges.  Then (86.19) holds.
\end{corollary}

\begin{proof}
For fixed \(F\), the uniform tail reduces the Cauchy condition to
the convergent finite-dimensional projections.  Hence
\({\cal T}_nF\) is strongly Cauchy.  Apply this separately on the
countable algebra.
\end{proof}

\section{Cofinal boundary-amplitude ledger}

The reflected construction now requires:
\begin{enumerate}
\item isometric or uniformly controlled common-carrier embeddings;
\item finite-coordinate convergence for every determining
      boundary amplitude;
\item a uniform finite-rank, energy, or semigroup tail;
\item convergence of transported vacuum vectors;
\item a visible overlap that survives the limit;
\item compatibility with parity, determinant-line sewing, and
      quotient boundaries.
\end{enumerate}

\begin{assessment}[Physical status]
\label{ass:v15v86-status}
V86 proves weak and strong boundary-amplitude compactness criteria,
spectral and semigroup tail compilers, persistence of a visible
vacuum, and convergence of the entire reflected Gram form.  It
also gives exact escaping-mode counterexamples.

The physical SM--Einstein regulator has not yet supplied a common
boundary carrier, uniform high-energy tail, finite-coordinate
convergence, or visible transported vacuum satisfying these
hypotheses.
\end{assessment}

\begin{decision}[V87 nuclear boundary-map compiler]
\label{dec:v15v86-next}
Derive the tail criterion from operator ideals.  Show that a
cutoff-uniform factorization of the half-amplitude map through a
common Hilbert--Schmidt or trace-class smoothing operator gives
collective compactness, track the singular-value tail, and exhibit
a growing-rank family with bounded operator norm but no collective
compactness.
\end{decision}

\begin{firewall}[Claim boundary after V86]
\label{fire:v15v86}
V86 proves conditional compactness and reflected convergence.  It
does not establish the required common carrier, nuclear smoothing,
vacuum visibility, phase normalization, or full continuum
construction.
\end{firewall}

\section*{Version V86 append-only record}
\addcontentsline{toc}{section}{Version V86 append-only record}
The complete V85 source was retained before this continuation.  It
had \(840507\) bytes, \(24824\) lines, and SHA--256
\[
 \texttt{91D9F6E05A5C4381E61322DE90C0BEF24302351B82EEABADB5C150F44FC54509}.
\]
V86 placed boundary amplitudes in a common carrier, proved
finite-rank and spectral compactness criteria, preserved a visible
vacuum component, and closed convergence of normalized reflected
Gram forms under a uniform high-energy tail.

\chapter{Fifteenth-cycle continuation V87: nuclear smoothing of the boundary map}
\label{sec:v15v87}

\section{Factorized half-amplitude maps}

Let \({\cal E}\) be a common normed carrier for positive-time
observables or boundary seeds, let \({\cal K}\) be an auxiliary
Hilbert space, and let \({\cal H}_\Sigma\) be the common boundary
carrier of V86.  Suppose
\[
 {\cal T}_n=C_nB_n,\qquad
 B_n:{\cal E}\to{\cal K},\quad
 C_n:{\cal K}\to{\cal H}_\Sigma.                      \tag{87.1}
\]

\begin{theorem}[Common compact-factor criterion]
\label{thm:v15v87-common}
Assume \(C:{\cal K}\to{\cal H}_\Sigma\) is compact,
\[
 \|C_n-C\|_{\rm op}\longrightarrow0,                  \tag{87.2}
\]
and for a fixed \(F\in{\cal E}\),
\[
 B_nF\rightharpoonup b_F\quad\hbox{weakly in }{\cal K}.          \tag{87.3}
\]
Then
\[
 {\cal T}_nF=C_nB_nF\longrightarrow Cb_F             \tag{87.4}
\]
strongly in \({\cal H}_\Sigma\).
\end{theorem}

\begin{proof}
Weak convergence implies
\(\sup_n\|B_nF\|<\infty\).  Hence
\[
\begin{split}
 \|C_nB_nF-Cb_F\|
 &\leq
 \|(C_n-C)B_nF\|+\|C(B_nF-b_F)\|.
\end{split}
\]
The first term tends to zero by (87.2).  A compact operator maps
weakly convergent sequences to strongly convergent sequences, so
the second term tends to zero.
\end{proof}

\begin{corollary}[Trace-ideal convergence]
\label{cor:v15v87-ideal-convergence}
The conclusion holds if \(C_n\to C\) in Hilbert--Schmidt or trace
norm.
\end{corollary}

\begin{proof}
Both ideal norms dominate the operator norm, and their limits are
compact.
\end{proof}

The theorem supplies both finite-coordinate convergence and the
tail criterion of V86 for each \(F\), rather than assuming them
separately.

\section{Singular-value tail of one common smoother}

Let \(C\in{\mathfrak S}_2({\cal K},{\cal H}_\Sigma)\) have singular
value decomposition
\[
 C=\sum_{j\geq1}s_j\,|u_j\rangle\langle v_j|,
 \qquad s_1\geq s_2\geq\cdots\geq0.                  \tag{87.5}
\]
Let \(P_R\) project onto
\(\operatorname{span}\{u_1,\ldots,u_R\}\).

\begin{proposition}[Hilbert--Schmidt boundary tail]
\label{prop:v15v87-hs-tail}
For every \(b\in{\cal K}\),
\[
 \|(I-P_R)Cb\|
 \leq
 \left(\sum_{j>R}s_j^2\right)^{1/2}\|b\|.             \tag{87.6}
\]
If \(C\in{\mathfrak S}_1\), also
\[
 \|(I-P_R)Cb\|
 \leq
 \left(\sum_{j>R}s_j\right)\|b\|.                     \tag{87.7}
\]
\end{proposition}

\begin{proof}
The operator norm is bounded by the Hilbert--Schmidt norm:
\[
 \|(I-P_R)C\|_{\rm op}
 \leq\|(I-P_R)C\|_2
 =
 \left(\sum_{j>R}s_j^2\right)^{1/2}.
\]
This proves (87.6).  The trace norm also dominates the operator
norm and equals the singular-value sum, giving (87.7).
\end{proof}

\begin{corollary}[Uniform amplitude tail]
\label{cor:v15v87-amplitude-tail}
If
\[
 \sup_n\|B_nF\|\leq M_F,\qquad C_n=C,                 \tag{87.8}
\]
then
\[
 \sup_n\|(I-P_R){\cal T}_nF\|
 \leq
 M_F\left(\sum_{j>R}s_j^2\right)^{1/2}\to0.           \tag{87.9}
\]
\end{corollary}

\begin{proof}
Apply Proposition~\ref{prop:v15v87-hs-tail} to
\(b=B_nF\) and take the supremum.
\end{proof}

\section{Collectively compact varying smoothers}

\begin{definition}[Common projection tail]
\label{def:v15v87-common-tail}
A family \(C_n:{\cal K}\to{\cal H}_\Sigma\) has a common compact
tail if finite-rank projections \(P_R\uparrow I\) satisfy
\[
 \lim_{R\to\infty}\sup_n
 \|(I-P_R)C_n\|_{\rm op}=0.                           \tag{87.10}
\]
\end{definition}

\begin{proposition}[Collective compactness]
\label{prop:v15v87-collective}
If (87.10) holds and
\[
 \sup_n\|C_n\|_{\rm op}<\infty,                       \tag{87.11}
\]
then
\[
 \bigcup_nC_n\{b:\|b\|\leq1\}
\]
is relatively compact provided each finite-dimensional family
\(\{P_RC_n\}_n\) has relatively compact coefficient matrices.
\end{proposition}

\begin{proof}
Given \(\epsilon>0\), choose \(R\) so the common tail is below
\(\epsilon/3\).  Relative compactness of the finite-dimensional
coefficient family gives a finite \(\epsilon/3\)-net for all
\(P_RC_nb\) with \(\|b\|\leq1\).  Adding the uniform tails gives a
finite \(\epsilon\)-net for the full union.
\end{proof}

\begin{proposition}[Ideal convergence produces a common tail]
\label{prop:v15v87-convergence-tail}
If \(C_n\to C\) in Hilbert--Schmidt norm, then the family has
(87.10).
\end{proposition}

\begin{proof}
Choose a finite-rank projection \(P_R\) with
\(\|(I-P_R)C\|_2<\epsilon/2\).  For all sufficiently large \(n\),
\[
 \|(I-P_R)C_n\|_{\rm op}
 \leq\|C_n-C\|_2+\|(I-P_R)C\|_2<\epsilon.
\]
Enlarge \(P_R\) to control the finitely many remaining \(C_n\),
which are individually compact because they are Hilbert--Schmidt.
\end{proof}

\section{Uniform ideal norm is not enough}

\begin{counterexample}[Drifting rank-one smoother]
\label{cnt:v15v87-rank-one}
Let \({\cal K}={\cal H}_\Sigma=\ell^2\) and set
\[
 C_n=|e_n\rangle\langle e_1|.                         \tag{87.12}
\]
Then
\[
 \|C_n\|_1=\|C_n\|_2=\|C_n\|_{\rm op}=1              \tag{87.13}
\]
for every \(n\), but
\(\{C_ne_1\}=\{e_n\}\) is not relatively compact.
\end{counterexample}

\begin{proof}
Each \(C_n\) has one singular value equal to one, giving (87.13).
The output vectors are pairwise a distance \(\sqrt2\) apart.
\end{proof}

Even a uniform trace-norm bound does not prevent the leading output
singular vector from escaping.

\begin{counterexample}[Growing-rank bounded operators]
\label{cnt:v15v87-growing-rank}
Let
\[
 C_n=P_n
\]
be projection onto
\(\operatorname{span}\{e_1,\ldots,e_n\}\).  Then
\[
 \|C_n\|_{\rm op}=1,                                  \tag{87.14}
\]
but the union of the images of the unit ball contains every
\(e_j\) and is not relatively compact.
\end{counterexample}

\begin{proof}
For \(n\geq j\), \(P_ne_j=e_j\).  The orthonormal sequence has no
strongly convergent subsequence.
\end{proof}

The common-output tail, not an unaligned list of singular values,
is the necessary carrier statement.

\section{Heat-smoothing specialization}

Let \(K\geq0\) be a common boundary operator with compact resolvent
and eigenvalues
\[
 0\leq\lambda_1\leq\lambda_2\leq\cdots\to\infty.
\]
For \(t>0\), put
\[
 C=e^{-tK}.                                           \tag{87.15}
\]

\begin{proposition}[Heat singular values]
\label{prop:v15v87-heat}
If
\[
 \operatorname{Tr}e^{-2tK}
 =
 \sum_je^{-2t\lambda_j}<\infty,                       \tag{87.16}
\]
then \(C\) is Hilbert--Schmidt and
\[
 \|(I-P_R)Cb\|^2
 \leq
 \left(\sum_{j>R}e^{-2t\lambda_j}\right)\|b\|^2.       \tag{87.17}
\]
\end{proposition}

\begin{proof}
The singular values of the positive operator \(e^{-tK}\) are
\(e^{-t\lambda_j}\).  Apply
Proposition~\ref{prop:v15v87-hs-tail}.
\end{proof}

For varying \(K_n\), uniform heat-trace bounds do not align the
eigenvectors.  One needs trace-ideal convergence of the heat
operators or a common spectral projection tail.

\section{Boundary Gram convergence}

\begin{theorem}[Nuclear half-amplitude reconstruction]
\label{thm:v15v87-reconstruction}
Let \({\cal A}_+\) be the countable algebra of V86.  Suppose for
each \(F\):
\[
 B_nF\rightharpoonup b_F,\qquad
 C_n\to C\ \hbox{in }{\mathfrak S}_2,                 \tag{87.18}
\]
with the same compact \(C\), and suppose the visible vacuum
condition of V86 holds.  Then
\[
 {\cal T}_nF\to Cb_F
\]
strongly for every \(F\), the reflected Gram forms converge, and
their normalized limit is positive and nonzero.
\end{theorem}

\begin{proof}
Theorem~\ref{thm:v15v87-common} gives strong convergence.
Theorem~\ref{thm:v15v86-visible} keeps the vacuum limit nonzero,
and Theorem~\ref{thm:v15v86-gram} gives Gram and normalized
reflected convergence.
\end{proof}

\begin{assessment}[Physical status]
\label{ass:v15v87-status}
V87 proves that trace-ideal convergence to one common boundary
smoother turns weak seed convergence into strong amplitude
convergence.  It tracks singular-value tails and proves that
uniform ideal norms, even trace norms, do not suffice when singular
directions drift.

The physical regulator has not yet supplied the common smoother,
trace-ideal convergence, weak seed limits, or aligned spectral tail.
\end{assessment}

\begin{decision}[V88 elliptic heat-tail source]
\label{dec:v15v87-next}
Derive a common Hilbert--Schmidt boundary smoother from uniform
ellipticity and geometry.  Prove a heat-trace tail from a
cutoff-uniform eigenvalue counting bound, transport it through
norm-resolvent convergence, and distinguish a fixed boundary
manifold from growing boundary volume or topology.
\end{decision}

\begin{firewall}[Claim boundary after V87]
\label{fire:v15v87}
V87 proves operator-ideal compactness compilers.  It does not prove
the required geometric spectral counting, trace-ideal convergence,
vacuum visibility, or full continuum reconstruction.
\end{firewall}

\section*{Version V87 append-only record}
\addcontentsline{toc}{section}{Version V87 append-only record}
The complete V86 source was retained before this continuation.  It
had \(850872\) bytes, \(25179\) lines, and SHA--256
\[
 \texttt{676A792D288C956634C677706FFE08731991C74BE33FAF705D0D50F145293965}.
\]
V87 factorized the half-amplitude map through a common nuclear
smoother, proved singular-value and collective compactness tails,
and isolated drifting singular directions and growing rank as
independent cofinal obstructions.

\chapter{Fifteenth-cycle continuation V88: elliptic counting and the common heat smoother}
\label{sec:v15v88}

\section{Uniform boundary spectral data}

Let \(\Sigma\) be a fixed compact \(d\)-dimensional boundary
manifold and let \(K_n\geq0\) be self-adjoint elliptic boundary
operators of order \(m>0\), transported to one Hilbert carrier.
Write their eigenvalues with multiplicity as
\[
 0\leq\lambda_{1,n}\leq\lambda_{2,n}\leq\cdots,       \tag{88.1}
\]
and define
\[
 N_n(\Lambda)=\#\{j:\lambda_{j,n}\leq\Lambda\}.        \tag{88.2}
\]

\begin{definition}[Cofinal counting bound]
\label{def:v15v88-count}
The family has a uniform elliptic counting bound when
\[
 N_n(\Lambda)\leq C_N(1+\Lambda)^\alpha,\qquad
 \alpha=\frac dm,                                     \tag{88.3}
\]
for every \(n\) and \(\Lambda\geq0\).
\end{definition}

\section{Heat trace and its tail}

\begin{theorem}[Uniform heat trace]
\label{thm:v15v88-heat-trace}
Under (88.3), for every \(t>0\),
\[
 \sup_n\operatorname{Tr}e^{-tK_n}<\infty.             \tag{88.4}
\]
More precisely, for \(R\geq1\),
\[
 \sup_n
 \sum_{\lambda_{j,n}>R}e^{-t\lambda_{j,n}}
 \leq
 C_{t,\alpha,C_N}e^{-tR/2}.                           \tag{88.5}
\]
\end{theorem}

\begin{proof}
As a Stieltjes integral,
\[
 \sum_{\lambda_{j,n}>R}e^{-t\lambda_{j,n}}
 =
 \int_{(R,\infty)}e^{-t\lambda}\,dN_n(\lambda).
                                                               \tag{88.6}
\]
Integration by parts and omission of the nonpositive lower
endpoint term give
\[
 \int_{(R,\infty)}e^{-t\lambda}\,dN_n(\lambda)
 \leq
 tC_N\int_R^\infty e^{-t\lambda}(1+\lambda)^\alpha
 d\lambda.                                             \tag{88.7}
\]
The function
\((1+\lambda)^\alpha e^{-t\lambda/2}\) is bounded on
\([0,\infty)\).  Pulling out \(e^{-tR/2}\) in (88.7) proves
(88.5).  Taking \(R=0\), with the finitely many zero modes included
through (88.3), proves (88.4).
\end{proof}

\begin{corollary}[Uniform Hilbert--Schmidt smoother]
\label{cor:v15v88-hs}
For every \(t>0\),
\[
 C_n(t)=e^{-tK_n}
\]
is Hilbert--Schmidt and
\[
 \sup_n\|C_n(t)\|_2^2
 =
 \sup_n\operatorname{Tr}e^{-2tK_n}<\infty.            \tag{88.8}
\]
Its intrinsic squared singular-value tail is bounded by
\[
 C_{2t,\alpha,C_N}e^{-tR}.                            \tag{88.9}
\]
\end{corollary}

\begin{proof}
The singular values are \(e^{-t\lambda_{j,n}}\).  Apply
Theorem~\ref{thm:v15v88-heat-trace} with \(2t\).
\end{proof}

The tail in (88.9) is ordered in the spectral basis of \(K_n\).
Alignment with a common carrier still requires convergence of the
low spectral subspaces.

\section{A geometric source of the counting law}

\begin{proposition}[Uniform ellipticity and bounded geometry]
\label{prop:v15v88-geometry}
Suppose the transported quadratic forms have a common Sobolev
domain \(H^{m/2}(\Sigma)\) and constants \(c_0,C_0>0\) such that
\[
 \langle K_nu,u\rangle+C_0\|u\|_2^2
 \geq c_0\|u\|_{H^{m/2}}^2                            \tag{88.10}
\]
uniformly in \(n\).  Assume the background boundary metrics range
in a bounded-geometry class with uniform volume bound.  Then
(88.3) holds with a cutoff-independent \(C_N\).
\end{proposition}

\begin{proof}
Let \(\mu_j\) be the eigenvalues of a fixed positive reference
operator \((I-\Delta_\Sigma)^{m/2}\).  The min--max principle and
(88.10) give
\[
 \lambda_{j,n}+C_0\geq c_1\mu_j                       \tag{88.11}
\]
with uniform \(c_1>0\).  The reference counting estimate on a
bounded-geometry compact \(d\)-manifold is
\[
 \#\{j:\mu_j\leq L\}\leq C(1+L)^{d/m}.                \tag{88.12}
\]
If \(\lambda_{j,n}\leq\Lambda\), then
\(\mu_j\leq c_1^{-1}(\Lambda+C_0)\).  Insert this in (88.12) and
absorb constants to obtain (88.3).
\end{proof}

Uniform ellipticity without volume and topology control is
insufficient because the number of independent low modes may grow.

\section{Norm-resolvent alignment}

Assume now that
\[
 \|(K_n+1)^{-1}-(K+1)^{-1}\|_{\rm op}\longrightarrow0 \tag{88.13}
\]
for a nonnegative compact-resolvent operator \(K\).

\begin{proposition}[Low spectral clusters converge]
\label{prop:v15v88-clusters}
Let \(\Lambda\) lie outside the spectrum of \(K\).  Then for all
large \(n\), \(\Lambda\) lies outside the spectrum of \(K_n\), the
spectral projections
\[
 P_n(\Lambda)={\bf1}_{[0,\Lambda]}(K_n),\qquad
 P(\Lambda)={\bf1}_{[0,\Lambda]}(K)                   \tag{88.14}
\]
have equal finite rank, and
\[
 \|P_n(\Lambda)-P(\Lambda)\|_{\rm op}\to0.             \tag{88.15}
\]
\end{proposition}

\begin{proof}
Apply the continuous spectral map
\(\lambda\mapsto(\lambda+1)^{-1}\).  The threshold
\((\Lambda+1)^{-1}\) lies in a resolvent gap of
\((K+1)^{-1}\).  The corresponding Riesz projection is a contour
integral of the resolvent.  Operator-norm convergence in (88.13)
gives uniform resolvent convergence on the contour, hence (88.15).
Ranks of projections within operator distance less than one are
equal.
\end{proof}

\begin{theorem}[Heat trace-norm convergence]
\label{thm:v15v88-trace-convergence}
Assume (88.3) and (88.13).  Then for every \(t>0\),
\[
 \|e^{-tK_n}-e^{-tK}\|_1\longrightarrow0.              \tag{88.16}
\]
\end{theorem}

\begin{proof}
Norm-resolvent convergence implies
\[
 \|e^{-tK_n}-e^{-tK}\|_{\rm op}\to0                  \tag{88.17}
\]
by continuous functional calculus for the function
\(e^{-t\lambda}\) vanishing at infinity.

Fix \(\epsilon>0\).  Choose a spectral threshold \(\Lambda\)
outside the spectrum of \(K\) so large that the heat-trace tails
of \(K\) and, by (88.5), all \(K_n\), are below \(\epsilon\).
Proposition~\ref{prop:v15v88-clusters} identifies finite-rank low
spectral spaces of one common rank \(r\).  On their sum, whose rank
is at most \(2r\), (88.17) implies trace-norm convergence because
\[
 \|A\|_1\leq2r\|A\|_{\rm op}.                         \tag{88.18}
\]
The two positive high-energy pieces have trace norms equal to their
traces, each below \(\epsilon\).  Hence the full trace-norm
limsup is at most \(2\epsilon\).  Let \(\epsilon\downarrow0\).
\end{proof}

\begin{corollary}[Common nuclear boundary smoother]
\label{cor:v15v88-common}
Under Theorem~\ref{thm:v15v88-trace-convergence}, the heat factors
\[
 C_n(t)=e^{-tK_n}
\]
converge in trace norm to the common compact operator
\(C(t)=e^{-tK}\).  They therefore satisfy every smoothing and
alignment hypothesis in V87.
\end{corollary}

\begin{proof}
Trace norm convergence implies Hilbert--Schmidt and operator-norm
convergence.  Apply Corollary~\ref{cor:v15v87-ideal-convergence}
and Proposition~\ref{prop:v15v87-convergence-tail}.
\end{proof}

\section{Form convergence sufficient for resolvents}

\begin{proposition}[Uniform form-norm criterion]
\label{prop:v15v88-forms}
Let \(q_n,q\) be closed nonnegative forms on a common domain
\({\cal V}\), and suppose
\[
 |q_n(u,v)-q(u,v)|
 \leq\epsilon_n
 \|u\|_{\cal V}\|v\|_{\cal V},\qquad\epsilon_n\to0,    \tag{88.19}
\]
where the form norms are uniformly equivalent.  Then (88.13)
holds.
\end{proposition}

\begin{proof}
For \(f\) in the carrier, let
\[
 u_n=(K_n+1)^{-1}f,\qquad u=(K+1)^{-1}f.
\]
Subtract the two variational resolvent equations and test with
\(u_n-u\).  Uniform coercivity bounds
\(\|u_n-u\|_{\cal V}\) by
\(C\epsilon_n\|f\|\).  The continuous embedding
\({\cal V}\hookrightarrow{\cal H}_\Sigma\) gives
\[
 \|u_n-u\|\leq C'\epsilon_n\|f\|.
\]
Taking the supremum over unit \(f\) proves (88.13).
\end{proof}

\section{Growing-boundary obstruction}

\begin{counterexample}[Disjoint boundary replication]
\label{cnt:v15v88-volume}
Let \(\Sigma_n\) be the disjoint union of \(n\) identical compact
components and let \(K_n\) be the same elliptic operator on each
component.  Then
\[
 N_n(\Lambda)=nN_1(\Lambda),\qquad
 \operatorname{Tr}e^{-tK_n}
 =
 n\operatorname{Tr}e^{-tK_1}.                        \tag{88.20}
\]
No uniform counting or heat-trace bound holds.
\end{counterexample}

\begin{proof}
The spectrum of a direct sum is the union of the component spectra
with multiplicities added, giving both identities.
\end{proof}

The same obstruction appears if topology creates a growing
zero-mode multiplicity.  A fixed local ellipticity constant does
not control global state count.

\section{Boundary geometric ledger}

The nuclear smoother now follows from:
\begin{enumerate}
\item a fixed or coherently identified compact boundary manifold;
\item uniform ellipticity and bounded geometry;
\item a uniform bound on volume, topology, and zero-mode
      multiplicity;
\item form-norm or norm-resolvent convergence;
\item a positive smoothing time before the reflection interface.
\end{enumerate}

\begin{assessment}[Physical status]
\label{ass:v15v88-status}
V88 proves a uniform heat-trace tail from elliptic eigenvalue
counting and upgrades norm-resolvent convergence to trace-norm
convergence of boundary heat smoothers.  It also isolates growing
boundary volume and topology as independent failures.

The physical fluctuating-boundary geometry has not yet been shown
to stay in the required bounded-geometry class or to yield the
uniform form convergence and topology bounds.
\end{assessment}

\begin{decision}[V89 Fock-space lift]
\label{dec:v15v88-next}
Lift the one-particle heat control to fermionic and bosonic boundary
Fock sectors.  Derive determinant/product formulas for the Fock
partition norms, prove the fermionic trace bound from one-particle
trace class, identify the bosonic need for a positive spectral gap,
and exhibit zero-mode divergence in the bosonic sector.
\end{decision}

\begin{firewall}[Claim boundary after V88]
\label{fire:v15v88}
V88 proves conditional elliptic heat smoothing.  It does not prove
bounded geometry for dynamical boundaries, Fock-space
normalization, vacuum visibility, or the full continuum.
\end{firewall}

\section*{Version V88 append-only record}
\addcontentsline{toc}{section}{Version V88 append-only record}
The complete V87 source was retained before this continuation.  It
had \(860429\) bytes, \(25501\) lines, and SHA--256
\[
 \texttt{AA510097AFE2710178CB3213A7474EA7E9C681579624F53DBDE742C163D4BF50}.
\]
V88 derived uniform elliptic heat tails, proved trace-norm
convergence from counting plus norm-resolvent alignment, and
identified geometric volume, topology, and zero-mode multiplicity
as the remaining boundary spectral inputs.
\chapter{Fifteenth-cycle continuation V89: Fock lifts and the infrared gate}
\label{sec:v15v89}

V88 produced a common trace-class one-particle heat smoother.  The
next question is whether second quantization preserves the nuclear
control needed by a boundary transfer construction.  The answer is
unconditional for fermions and conditional on a strict one-particle
gap for bosons.  This distinction is an infrared fact, not a defect
of the determinant notation.

\section{Second-quantized heat factors}

Let \({\cal h}\) be a separable one-particle Hilbert space and let
\(Q\) be a positive trace-class contraction on \({\cal h}\).  Write
its eigenvalues, with multiplicity, as
\[
 1\geq q_1\geq q_2\geq\cdots\geq0,
 \qquad \sum_jq_j<\infty.                              \tag{89.1}
\]
The fermionic and bosonic Fock spaces are
\[
 {\cal F}_{-}({\cal h})=\bigoplus_{k\geq0}\bigwedge^k{\cal h},
 \qquad
 {\cal F}_{+}({\cal h})=\bigoplus_{k\geq0}{\rm Sym}^k{\cal h},       \tag{89.2}
\]
and the second quantization of \(Q\) is
\[
 \Gamma_\pm(Q)=\bigoplus_{k\geq0}Q^{\otimes k}\big|_{{\cal F}^{(k)}_\pm},
 \qquad Q^{\otimes0}=1.                              \tag{89.3}
\]
The vacuum contribution is therefore one.

\begin{theorem}[Fermionic nuclear lift]
\label{thm:v15v89-fermion}
For every positive trace-class contraction \(Q\),
\(\Gamma_-(Q)\) is positive and trace class, and
\[
 \operatorname{Tr}_{{\cal F}_{-}}\Gamma_-(Q)
 =\prod_{j\geq1}(1+q_j)
 =\det(1+Q)
 \leq \exp(\operatorname{Tr}Q).                       \tag{89.4}
\]
No strict inequality \(\|Q\|_{\rm op}<1\) is required.
\end{theorem}

\begin{proof}
In an eigenbasis of \(Q\), an occupation vector is indexed by a
finite subset \(S\) of the one-particle modes and has eigenvalue
\(\prod_{j\in S}q_j\).  Positivity and monotone convergence give
\[
 \operatorname{Tr}\Gamma_-(Q)
 =\sum_{S\subset\mathbb N\ {
m finite}}\prod_{j\in S}q_j
 =\lim_{m\to\infty}\prod_{j=1}^m(1+q_j).             \tag{89.5}
\]
Since \(\log(1+x)\leq x\) for \(x\geq0\), the product is finite
and obeys (89.4).  The product is the defining positive
trace-class Fredholm determinant.
\end{proof}

\begin{theorem}[Bosonic nuclear lift]
\label{thm:v15v89-boson}
For a positive trace-class contraction \(Q\), the following are
equivalent:
\begin{enumerate}
\item \(\Gamma_+(Q)\) is trace class;
\item \(1\) is not an eigenvalue of \(Q\);
\item \(\|Q\|_{\rm op}<1\).
\end{enumerate}
When these conditions hold,
\[
 \operatorname{Tr}_{{\cal F}_{+}}\Gamma_+(Q)
 =\prod_{j\geq1}(1-q_j)^{-1}
 =\det(1-Q)^{-1},                                     \tag{89.6}
\]
and, if \(\|Q\|_{\rm op}\leq q_*<1\),
\[
 \log\operatorname{Tr}\Gamma_+(Q)
 \leq {\operatorname{Tr}Q\over1-q_*}.                \tag{89.7}
\]
\end{theorem}

\begin{proof}
An occupation vector is indexed by a finitely supported sequence
\((n_j)_{j\geq1}\) of nonnegative integers and has eigenvalue
\(\prod_jq_j^{n_j}\).  Tonelli's theorem gives
\[
 \operatorname{Tr}\Gamma_+(Q)
 =\prod_j\sum_{n_j\geq0}q_j^{n_j}
 =\prod_j(1-q_j)^{-1},                                \tag{89.8}
\]
with the value \(+\infty\) allowed.  A positive compact operator
attains its norm.  Thus norm one is equivalent to an eigenvalue
one, which makes one geometric series in (89.8) divergent.

If \(q_*<1\), then
\[
 -\log(1-q_j)=\sum_{r\geq1}{q_j^r\over r}
 \leq {q_j\over1-q_*}.                               \tag{89.9}
\]
Summing proves convergence and (89.7).  Conversely, finiteness of
(89.8) excludes an eigenvalue one.  The Fredholm determinant
identity follows from the same eigenvalue product.
\end{proof}

\section{Heat interpretation and the exact infrared condition}

Let \(K\geq0\) have compact resolvent and put \(Q(t)=e^{-tK}\).
V88 supplies \(Q(t)\in{\cal L}^1\).  Its norm is
\[
 \|Q(t)\|_{\rm op}=e^{-t\inf\sigma(K)}.              \tag{89.10}
\]

\begin{corollary}[Fock heat criterion]
\label{cor:v15v89-heat}
For every \(t>0\), the fermionic heat factor
\(\Gamma_-(e^{-tK})\) is trace class.  The bosonic heat factor
\(\Gamma_+(e^{-tK})\) is trace class exactly when
\[
 K\geq\mu I\quad\hbox{for some }\mu>0.                \tag{89.11}
\]
Under a uniform gap \(K_n\geq\mu I\), its trace is bounded by
\[
 \operatorname{Tr}\Gamma_+(e^{-tK_n})
 \leq
 \exp\left\{{\operatorname{Tr}e^{-tK_n}\over1-e^{-t\mu}}\right\}. \tag{89.12}
\]
\end{corollary}

\begin{proof}
Apply Theorems~\ref{thm:v15v89-fermion} and
\ref{thm:v15v89-boson}.  Compact resolvent makes the bottom of the
spectrum an eigenvalue, so (89.10) is strictly below one exactly
under (89.11).  Formula (89.7) gives (89.12).
\end{proof}

If a finite-dimensional gauge or constraint kernel is removed by
an orthogonal projection \(P^\perp\), condition (89.11) must hold
for \(K|_{P^\perp{\cal h}}\).  The projection must be part of the
common-carrier data and must converge as in V86--V88.  Merely
deleting a zero eigenvalue separately at every cutoff does not
identify the physical quotient.

\section{Continuity on the common Fock carrier}

The scalar determinant is continuous in the topology already
obtained in V88.

\begin{proposition}[Positive determinant continuity]
\label{prop:v15v89-determinants}
Let \(A,B\geq0\) be trace class.  Then
\[
 |\det(1+A)-\det(1+B)|
 \leq
 \|A-B\|_1\exp\!\bigl(\max\{\operatorname{Tr}A,
                                  \operatorname{Tr}B\}\bigr).       \tag{89.13}
\]
If in addition \(\|A\|_{\rm op},\|B\|_{\rm op}\leq q_*<1\), then
\[
 |\det(1-A)^{-1}-\det(1-B)^{-1}|
 \leq {\|A-B\|_1\over1-q_*}
 \exp\!\left({\max\{\operatorname{Tr}A,
                       \operatorname{Tr}B\}\over1-q_*}\right).    \tag{89.14}
\]
\end{proposition}

\begin{proof}
Put \(C_s=(1-s)A+sB\).  Finite-rank approximation, followed by
trace-norm passage to the limit, justifies the Fredholm derivative
formula
\[
 {d\over ds}\det(1+C_s)
 =\det(1+C_s)\operatorname{Tr}((1+C_s)^{-1}(B-A)).     \tag{89.15}
\]
Here \(\|(1+C_s)^{-1}\|\leq1\) and
\(\det(1+C_s)\leq e^{\operatorname{Tr}C_s}\).  Integration proves
(89.13).  For the inverse determinant, differentiate
\(\det(1-C_s)^{-1}\).  Convexity preserves
\(\|C_s\|\leq q_*\), while
\[
 \|(1-C_s)^{-1}\|\leq(1-q_*)^{-1},\qquad
 \det(1-C_s)^{-1}
 \leq e^{\operatorname{Tr}C_s/(1-q_*)}.              \tag{89.16}
\]
Integration gives (89.14).
\end{proof}

Trace-norm convergence also holds for the full second-quantized
operators, not only for their traces.

\begin{theorem}[Fock trace-norm continuity]
\label{thm:v15v89-fock-continuity}
Suppose \(Q_n,Q\geq0\) and \(\|Q_n-Q\|_1\to0\).
Then
\[
 \|\Gamma_-(Q_n)-\Gamma_-(Q)\|_1\to0.               \tag{89.17}
\]
If also
\[
 \sup_n\|Q_n\|_{\rm op}\leq q_*<1,
 \qquad \|Q\|_{\rm op}\leq q_*,                     \tag{89.18}
\]
then
\[
 \|\Gamma_+(Q_n)-\Gamma_+(Q)\|_1\to0.               \tag{89.19}
\]
\end{theorem}

\begin{proof}
Trace-norm convergence bounds all \(\operatorname{Tr}Q_n\) by one
constant \(M\).  On every fixed particle sector, tensor
telescoping gives
\[
 \|Q_n^{\otimes k}-Q^{\otimes k}\|_1
 \leq k\max(\|Q_n\|_1,\|Q\|_1)^{k-1}\|Q_n-Q\|_1,    \tag{89.20}
\]
so compression to the symmetric or antisymmetric subspace
converges in trace norm.

For fermions, if \(e_k(Q_n)=\operatorname{Tr}\bigwedge^kQ_n\),
then the elementary-symmetric estimate gives
\[
 e_k(Q_n)\leq {M^k\over k!}.                          \tag{89.21}
\]
Consequently the sum of sector trace norms above particle number
\(L\) tends to zero uniformly in \(n\).

For bosons choose \(r\) with \(1<r<q_*^{-1}\).  Writing
\(h_k(Q_n)=\operatorname{Tr}{\rm Sym}^kQ_n\), the product formula
at \(rQ_n\) yields
\[
 \sum_{k\geq0}r^kh_k(Q_n)
 =\det(1-rQ_n)^{-1}
 \leq \exp\left({rM\over1-rq_*}\right).              \tag{89.22}
\]
Thus
\[
 \sum_{k>L}h_k(Q_n)
 \leq r^{-(L+1)}
      \exp\left({rM\over1-rq_*}\right),              \tag{89.23}
\]
uniformly in \(n\), and the same holds for \(Q\).  First truncate
the Fock direct sum at \(L\), use (89.20), and then let
\(L\to\infty\).  This proves (89.17) and (89.19).
\end{proof}

\begin{corollary}[Boundary Fock smoother]
\label{cor:v15v89-boundary}
Under the hypotheses of Theorem~\ref{thm:v15v88-trace-convergence},
\[
 \Gamma_-(e^{-tK_n})\longrightarrow\Gamma_-(e^{-tK})
 \quad\hbox{in trace norm}.                            \tag{89.24}
\]
The same conclusion holds in the bosonic sector if
\(K_n\geq\mu I\) and \(K\geq\mu I\) for one \(\mu>0\).
\end{corollary}

\begin{proof}
V88 gives \(e^{-tK_n}\to e^{-tK}\) in trace norm.  For bosons the
uniform gap gives (89.18) with \(q_*=e^{-t\mu}<1\).  Apply
Theorem~\ref{thm:v15v89-fock-continuity}.
\end{proof}

\section{Sharp infrared failures}

\begin{counterexample}[One bosonic zero mode]
\label{cnt:v15v89-zero}
If \(Ke_0=0\), the vectors with \(m\) particles in \(e_0\) all
have heat eigenvalue one.  Hence
\[
 \operatorname{Tr}\Gamma_+(e^{-tK})\geq\sum_{m\geq0}1=\infty.      \tag{89.25}
\]
The one-particle heat operator may nevertheless be trace class.
\end{counterexample}

\begin{counterexample}[Closing bosonic gap]
\label{cnt:v15v89-closing}
On a fixed one-dimensional carrier take \(K_n=\mu_n I\) with
\(\mu_n\downarrow0\).  Then
\[
 \operatorname{Tr}e^{-tK_n}=e^{-t\mu_n}\leq1,
 \qquad
 \operatorname{Tr}\Gamma_+(e^{-tK_n})
 ={1\over1-e^{-t\mu_n}}\longrightarrow\infty.         \tag{89.26}
\]
Thus a uniform one-particle heat trace does not replace a uniform
bosonic gap.
\end{counterexample}

\begin{counterexample}[Growing fermionic kernel]
\label{cnt:v15v89-fermion-kernel}
If \(\dim\ker K_n=m_n\), the zero modes contribute the factor
\[
 2^{m_n}                                                     \tag{89.27}
\]
to the fermionic Fock trace.  Every fixed finite kernel is harmless,
but unbounded kernel multiplicity destroys uniform normalization.
\end{counterexample}

The three examples separate the needed inputs: one-particle
nuclearity, a bosonic infrared gap after the physical quotient,
and a uniform fermionic zero-mode count.

\section{Physical ledger and next bottleneck}

For a Standard Model boundary sector, V89 permits the following
conditional assembly:
\begin{enumerate}
\item use V86--V88 to put all one-particle boundary operators on a
      common carrier and prove trace-norm heat convergence;
\item apply the fermionic lift directly, retaining the determinant
      line and anomaly constraints of V81--V84;
\item first quotient gauge degeneracies in the bosonic sector, then
      prove a positive gap or introduce an explicitly controlled
      infrared regulator;
\item keep the Higgs, gauge, ghost, and gravitational normalizations
      separate until their quotient and phase factors have been
      shown to combine.
\end{enumerate}

\begin{assessment}[Physical status]
\label{ass:v15v89-status}
V89 proves the exact free-Fock nuclearity and continuity theorem
needed after V88.  It shows that fermionic trace class follows from
one-particle trace class, whereas bosonic trace class additionally
requires a strict spectral gap after quotienting zero modes.

No uniform gap for the fluctuating Standard Model--Einstein
boundary Hamiltonian has been proved.  Interactions have not yet
been inserted into the Fock heat factor.
\end{assessment}

\begin{decision}[V90 companion audit]
\label{dec:v15v89-next}
Perform the compulsory YM/NS companion audit.  Test whether their
newest results supply a quotient gap, a common heat carrier, or an
interaction estimate usable in the Fock construction.  Then choose
the next upstream bottleneck and continue in V91.
\end{decision}

\begin{firewall}[Claim boundary after V89]
\label{fire:v15v89}
Fock nuclearity is a normalization theorem for a conditional free
boundary heat operator.  It is not a construction of the interacting
Standard Model Fock measure, a Yang--Mills mass gap, a dynamical
Einstein boundary gap, or the full continuum.
\end{firewall}

\section*{Version V89 append-only record}
\addcontentsline{toc}{section}{Version V89 append-only record}
The complete V88 source was retained before this continuation.  It
had \(870709\) bytes, \(25819\) lines, and SHA--256
\[
 \texttt{CEEEE14DB8AE166C211B1D705FEE9E6A56085B0A22E6CB924A3FB96318FC99A7}.
\]
V89 lifted the one-particle heat smoother to fermionic Fock space,
proved the sharp gap criterion for bosonic Fock nuclearity, proved
trace-norm continuity on the common Fock carrier, and isolated zero
modes and closing gaps as genuine infrared obstructions.
\chapter{Fifteenth-cycle continuation V90: companion audit and the interaction pivot}
\label{sec:v15v90}

This is the compulsory five-version audit following V85.  The
newest active Yang--Mills and Navier--Stokes notes in the shared
folder were read before choosing the V91 direction.  The audit asks
one narrow question: do the companion results close either the
uniform bosonic gap in V89 or the still-missing interacting Fock
heat estimate?

\section{Frozen companion checkpoints}

The inspected files are
\[
\begin{array}{c|r|r}
\hbox{file}&\hbox{bytes}&\hbox{lines}\\ \hline
\texttt{Renewal\_Yang\_Mills\_Mass\_Gap\_Research\_Notes\_V41.tex}
 &356260&9752\\
\texttt{Renewal\_NS\_Stability\_Research\_Notes\_V26.tex}
 &278283&5319.
\end{array}                                             \tag{90.1}
\]
Their SHA--256 digests are, respectively,
\[
\begin{split}
 &\texttt{0C111A921ABD8DC8970DB8E19412172D72469898CC94670F92D12CC9BB8DD498},\\
 &\texttt{82C887F8C2C69E4F28FAD7C3DC8DE5B9099CB5A4BF431EBA1FFF3E91935978AD}.
\end{split}                                             \tag{90.2}
\]

\section{Yang--Mills V41: a local symbol floor}

YM V41 proves, by exact rational product-hull certificates, a
positive floor for a \(45\)-dimensional Wilson pencil on one
selected transverse quarter cell and obtains the paired
minus-oriented statement by momentum inversion.  It explicitly
does not prove full-torus coverage, a covariant operator
inequality, a cofinal continuum measure, reflection positivity, or
a Hamiltonian mass gap.

The distinction between a cellwise symbol floor and an operator
gap is exact.

\begin{proposition}[Fourier symbol import criterion]
\label{prop:v15v90-symbol}
Let \(T\) be a translation-invariant self-adjoint operator whose
Fourier representation is multiplication by a measurable Hermitian
matrix \(p(k)\) on a compact momentum torus \(\mathbb T^d\).  Then
\[
 T\geq\mu I
 \quad\Longleftrightarrow\quad
 p(k)\geq\mu I\quad\hbox{for almost every }k\in\mathbb T^d.          \tag{90.3}
\]
Consequently, a positive floor on a proper measurable subset does
not by itself give any lower bound for \(T\).
\end{proposition}

\begin{proof}
For \(f\) in Fourier space,
\[
 \langle f,Tf\rangle
 =\int_{\mathbb T^d}\langle f(k),p(k)f(k)\rangle\,dk.   \tag{90.4}
\]
The pointwise floor implies the integrated inequality.  Conversely,
if the pointwise floor fails on a set of positive measure, measurable
spectral projection and approximation by simple vector fields give
an \(f\) supported on that set with
\(\langle f,Tf\rangle<\mu\|f\|^2\).  For the final assertion, take
the scalar symbol equal to \(1\) on the certified subset and
\(-1\) on a positive-measure part of its complement.
\end{proof}

\begin{assessment}[YM import decision]
\label{ass:v15v90-ym}
The V41 rational floor is relevant evidence for one local
one-particle symbol sector, but Proposition~\ref{prop:v15v90-symbol}
prevents its use as the uniform gap (89.11).  Import is deferred
until full-torus coverage, gauge quotient control, and passage from
the finite pencil to one common self-adjoint boundary Hamiltonian
are independently proved.
\end{assessment}

\section{Navier--Stokes V26: a restricted-input norm}

NS V26 derives closed formulas for Oseen-kernel derivative norms
when the input tensor is the traceless part of \(a\otimes a\).
For the first derivative the restricted norm equals the full
symmetric-traceless norm.  For the second derivative it is smaller
on an intermediate radial range.  These are rigorous estimates for
the special nonlinear input occurring in that programme.

\begin{lemma}[Restricted cones require an invariant input theorem]
\label{lem:v15v90-cone}
Let \(X,Y\) be normed spaces, \(S\subset X\), and \(L:X\to Y\)
linear.  A bound
\[
 \|Ls\|_Y\leq C\|s\|_X\qquad(s\in S)                  \tag{90.5}
\]
controls \(Lx\) for a general representation
\(x=\sum_{\ell=1}^m a_\ell s_\ell\) only by
\[
 \|Lx\|_Y
 \leq C\sum_{\ell=1}^m|a_\ell|\,\|s_\ell\|_X.         \tag{90.6}
\]
It yields the full operator bound with constant \(C\) only if the
decomposition gauge on the right is bounded by \(\|x\|_X\) with
constant one.
\end{lemma}

\begin{proof}
The triangle inequality proves (90.6).  The last assertion follows
by taking the infimum over all such representations.  Without the
stated comparison, that infimum defines a generally larger atomic
norm, so (90.5) does not imply a \(C\)-bound in the original norm.
\end{proof}

\begin{assessment}[NS import decision]
\label{ass:v15v90-ns}
No NS estimate is imported into the Fock construction.  The
rank-one restriction is legitimate for the specific classical
Oseen source, but no theorem confines arbitrary quantum interaction
forms to the same input cone.  Lemma~\ref{lem:v15v90-cone} records
the missing bridge and prevents a silent substitution of a
restricted norm for a Fock-space form norm.
\end{assessment}

\section{Audit conclusion and upstream pivot}

\begin{theorem}[V90 companion-audit conclusion]
\label{thm:v15v90-audit}
Neither inspected companion note proves the uniform post-quotient
bosonic gap required by V89, nor does either provide a uniform
relative-form estimate for an interacting Standard Model--Einstein
boundary Hamiltonian.
\end{theorem}

\begin{proof}
YM V41 is excluded as a gap input by
Proposition~\ref{prop:v15v90-symbol} and its own full-torus and
continuum firewalls.  NS V26 concerns a classical Oseen kernel on a
restricted input family; Lemma~\ref{lem:v15v90-cone} identifies the
additional invariant-input theorem that would be needed for a
quantum import.  Neither note constructs the required common
interacting Fock operator.
\end{proof}

The audit therefore changes the next attack from trying to assert a
physical mass gap to proving the strongest useful heat theorem that
does not assume one: uniform nuclearity and convergence under
relative form domination.  This goes upstream of the unresolved
spectral mechanism and isolates exactly what an interaction must
satisfy.

\begin{decision}[V91 interacting heat compiler]
\label{dec:v15v90-next}
Let \(H_{0,n}\) be the free Fock Hamiltonians controlled in V89 and
\(V_n\) the cutoff interactions.  Prove that the uniform form lower
bound
\[
 H_n=H_{0,n}+V_n\geq aH_{0,n}-bI,\qquad a>0,           \tag{90.7}
\]
transfers free heat-trace tails to the interacting operators.  Add
norm-resolvent convergence to prove trace-norm convergence of
\(e^{-tH_n}\), and give counterexamples showing why a mere lower
bound \(H_n\geq-bI\) or pointwise convergence of forms is
insufficient.
\end{decision}

\begin{firewall}[Claim boundary after V90]
\label{fire:v15v90}
The audit imports no mass gap or Navier--Stokes estimate.  It proves
why the available local and restricted-input companion results do
not close the Fock infrared or interaction gates.  The full
Standard Model--Einstein continuum remains open.
\end{firewall}

\section*{Version V90 append-only record}
\addcontentsline{toc}{section}{Version V90 append-only record}
The complete V89 source was retained before this continuation.  It
had \(883040\) bytes, \(26178\) lines, and SHA--256
\[
 \texttt{F37E7EB4D7AB570481CE3F8DA232458FB1BF4785685D0DEE662AF30459D77F99}.
\]
V90 inspected YM V41 and NS V26, froze their checkpoint digests,
proved the precise local-symbol and restricted-input non-import
criteria, and selected a uniform interacting-form heat theorem for
V91.
\chapter{Fifteenth-cycle continuation V91: interacting heat control by form domination}
\label{sec:v15v91}

V89 constructed convergent free Fock heat factors.  V90 showed that
the companion programmes do not yet provide a physical bosonic mass
gap or an interaction bound.  This note goes upstream: it proves an
abstract interacting heat theorem whose hypotheses can be checked
before the full spectral problem is solved.

\section{Uniform relative forms over varying free Hamiltonians}

Let \({\cal F}\) be one common Fock carrier.  For every \(n\), let
\(H_{0,n}\geq0\) be self-adjoint with compact resolvent and closed
quadratic form \(q_{0,n}\).  Let \(v_n\) be a symmetric form on
\({\cal D}(q_{0,n})\) satisfying
\[
 |v_n[u]|
 \leq \alpha q_{0,n}[u]+\beta\|u\|^2,
 \qquad 0\leq\alpha<1,\quad \beta<\infty,              \tag{91.1}
\]
with \(\alpha,\beta\) independent of \(n\).  Denote by \(H_n\) the
form sum \(H_{0,n}\dotplus V_n\).

\begin{lemma}[Compact interacting form sum]
\label{lem:v15v91-klmn}
Under (91.1), \(H_n\) is self-adjoint, bounded below, and has compact
resolvent.  Its eigenvalues in nondecreasing order obey
\[
 \lambda_j(H_n)
 \geq(1-\alpha)\lambda_j(H_{0,n})-\beta
 \qquad(j\geq1).                                      \tag{91.2}
\]
\end{lemma}

\begin{proof}
The KLMN theorem gives a closed lower-bounded form on
\({\cal D}(q_{0,n})\), and
\[
 q_n[u]\geq(1-\alpha)q_{0,n}[u]-\beta\|u\|^2.          \tag{91.3}
\]
The form norms of \(q_n\) and \(q_{0,n}\), after adding a common
constant, are equivalent.  Hence compactness of the embedding of
\({\cal D}(q_{0,n})\) into \({\cal F}\) is preserved, so \(H_n\)
has compact resolvent.  The min--max principle applied to (91.3)
gives (91.2).
\end{proof}

\begin{theorem}[Interacting Gibbs compiler]
\label{thm:v15v91-gibbs}
Assume (91.1) and suppose that for every \(s>0\) there is a
nonnegative self-adjoint \(H_0\) such that
\[
 \|e^{-sH_{0,n}}-e^{-sH_0}\|_1\longrightarrow0.        \tag{91.4}
\]
Then, for every \(t>0\),
\[
 \operatorname{Tr}e^{-tH_n}
 \leq e^{t\beta}
       \operatorname{Tr}e^{-t(1-\alpha)H_{0,n}},       \tag{91.5}
\]
and the right-hand side is uniformly bounded in \(n\).

If, in addition, \(H_n\to H\) in norm resolvent, then
\[
 \|e^{-tH_n}-e^{-tH}\|_1\longrightarrow0
 \qquad(t>0).                                         \tag{91.6}
\]
\end{theorem}

\begin{proof}
Apply the decreasing function \(x\mapsto e^{-tx}\) to the
eigenvalue inequality (91.2) and sum:
\[
 \sum_{j\geq1}e^{-t\lambda_j(H_n)}
 \leq e^{t\beta}\sum_{j\geq1}
       e^{-t(1-\alpha)\lambda_j(H_{0,n})}.             \tag{91.7}
\]
This proves (91.5).  Trace-norm convergence in (91.4), at
\(s=t(1-\alpha)\), gives the uniform bound.

For the convergence assertion, (91.4) also gives a uniform
free-sector eigenvalue tail:
\[
 \lim_{J\to\infty}\ \sup_n
 \sum_{j>J}e^{-t(1-\alpha)\lambda_j(H_{0,n})}=0.        \tag{91.8}
\]
Indeed, trace-norm convergence implies convergence of singular
values in \(\ell^1\), and finitely many exceptional initial indices
can be absorbed by increasing \(J\).  Equations (91.2) and (91.8)
therefore give
\[
 \lim_{J\to\infty}\ \sup_n
 \sum_{j>J}e^{-t\lambda_j(H_n)}=0.                     \tag{91.9}
\]

The lower bound \(H_n\geq-\beta I\) and norm-resolvent convergence
imply
\[
 \|e^{-tH_n}-e^{-tH}\|_{\rm op}\to0                   \tag{91.10}
\]
by the continuous functional calculus on \([-\beta,\infty)\).
The limit resolvent is compact because it is the operator-norm
limit of compact resolvents.  Choose a spectral threshold outside
\(\sigma(H)\).  Norm-resolvent convergence aligns the corresponding
finite-rank Riesz projections.  On the aligned low-energy sum,
(91.10) implies trace-norm convergence because the rank is fixed.
The two high-energy trace norms are uniformly small by (91.9);
the same tail bound passes to \(H\) by convergence of each fixed
eigenvalue.  Letting the threshold tend to infinity proves (91.6).
\end{proof}

\begin{remark}[Only one time is needed for one conclusion]
\label{rem:v15v91-one-time}
To prove (91.5)--(91.6) at a fixed \(t\), it is enough to assume
(91.4) at \(s=t(1-\alpha)\).  The all-\(s\) formulation is used
because a boundary transfer semigroup must be controlled at every
positive collar width.
\end{remark}

\section{Application to the Fock lift}

For a one-particle operator \(K_n\), let
\(d\Gamma_\pm(K_n)\) denote its additive second quantization.  The
functional-calculus identity is
\[
 e^{-s\,d\Gamma_\pm(K_n)}
 =\Gamma_\pm(e^{-sK_n}).                               \tag{91.11}
\]

\begin{corollary}[Interacting boundary Fock heat]
\label{cor:v15v91-fock}
Assume the one-particle hypotheses of V88.  In the fermionic sector
put \(H_{0,n}=d\Gamma_-(K_n)\).  In the bosonic sector put
\(H_{0,n}=d\Gamma_+(K_n)\) and additionally assume one uniform
post-quotient gap \(K_n\geq\mu I\), \(K\geq\mu I\).

If \(v_n\) satisfies (91.1), then the interacting heat factors obey
the uniform trace bound (91.5).  If the form sums \(H_n\) also
converge in norm resolvent, then they obey (91.6).
\end{corollary}

\begin{proof}
By (91.11), Theorem~\ref{thm:v15v89-fock-continuity} is precisely
(91.4) for the free Fock Hamiltonians.  In the bosonic case the
uniform gap is the hypothesis used in V89 to obtain that conclusion.
Apply Theorem~\ref{thm:v15v91-gibbs}.
\end{proof}

\begin{proposition}[A sufficient convergence ledger]
\label{prop:v15v91-form-convergence}
Suppose, after transport to a common form domain \({\cal V}\), that
the full forms \(q_n=q_{0,n}+v_n\) and \(q=q_0+v\) have uniformly
equivalent form norms and
\[
 |q_n(u,w)-q(u,w)|
 \leq\epsilon_n\|u\|_{\cal V}\|w\|_{\cal V},
 \qquad \epsilon_n\to0.                               \tag{91.12}
\]
Then the norm-resolvent hypothesis in
Theorem~\ref{thm:v15v91-gibbs} holds.
\end{proposition}

\begin{proof}
This is Proposition~\ref{prop:v15v88-forms} applied on the common
Fock carrier.  Subtracting the two variational resolvent equations
and using uniform coercivity gives operator-norm convergence of
the resolvents.
\end{proof}

\section{Normalized states and boundary observables}

\begin{corollary}[Gibbs state convergence]
\label{cor:v15v91-states}
Under Theorem~\ref{thm:v15v91-gibbs}, set
\[
 Z_n(t)=\operatorname{Tr}e^{-tH_n},\qquad
 \rho_n(t)=Z_n(t)^{-1}e^{-tH_n}.                       \tag{91.13}
\]
Then \(Z_n(t)\to Z(t)>0\) and
\[
 \|\rho_n(t)-\rho(t)\|_1\to0.                         \tag{91.14}
\]
Consequently, for every bounded boundary observable \(A\),
\[
 \operatorname{Tr}(\rho_n(t)A)
 \longrightarrow\operatorname{Tr}(\rho(t)A).          \tag{91.15}
\]
\end{corollary}

\begin{proof}
The trace is continuous in trace norm, so (91.6) gives convergence
of \(Z_n\).  The positive operator \(e^{-tH}\) is nonzero and trace
class, hence \(Z>0\).  Add and subtract
\(Z_n^{-1}e^{-tH}\) to prove (91.14).  Trace duality gives
\[
 |\operatorname{Tr}((\rho_n-\rho)A)|
 \leq\|\rho_n-\rho\|_1\|A\|_{\rm op},
\]
which proves (91.15).
\end{proof}

\section{Why the hypotheses cannot be weakened formally}

\begin{counterexample}[A lower bound does not control multiplicity]
\label{cnt:v15v91-multiplicity}
On \(\ell^2(\mathbb N)\), define
\[
 H_ne_j=
 \begin{cases}
 0,&j\leq n,\\
 j e_j,&j>n.
 \end{cases}                                          \tag{91.16}
\]
Each \(H_n\geq0\) has compact resolvent, but
\[
 \operatorname{Tr}e^{-tH_n}
 =n+\sum_{j>n}e^{-tj}\longrightarrow\infty.            \tag{91.17}
\]
Thus a uniform scalar lower bound contains no state-counting
information.
\end{counterexample}

\begin{counterexample}[Pointwise form and strong-resolvent loss]
\label{cnt:v15v91-drifting}
Let \(H_*e_j=je_j\), and let \(H_n\) agree with \(H_*\) except that
\(H_ne_n=0\).  The common form domain is
\[
 {\cal V}=\left\{u:\sum_{j\geq1}j|u_j|^2<\infty\right\},
\]
and, for every \(u\in{\cal V}\),
\[
 q_n[u]-q_*[u]=-n|u_n|^2\longrightarrow0.             \tag{91.18}
\]
The resolvents converge strongly because the resolvent difference
is supported in the drifting coordinate \(n\).  Nevertheless,
\[
 \|e^{-tH_n}-e^{-tH_*}\|_1=1-e^{-tn}\longrightarrow1. \tag{91.19}
\]
Pointwise form convergence and strong-resolvent convergence
therefore do not replace uniform form domination plus
norm-resolvent alignment.
\end{counterexample}

\begin{counterexample}[Loss of the coercive fraction]
\label{cnt:v15v91-alpha}
With \(H_*e_j=je_j\), take \(H_n=n^{-1}H_*\).  Every heat operator
is trace class and \(H_n\geq0\), but
\[
 \operatorname{Tr}e^{-tH_n}
 =\sum_{j\geq1}e^{-tj/n}
 ={1\over e^{t/n}-1}\sim{n\over t}.                   \tag{91.20}
\]
This is the degeneration of (91.3) when its positive coefficient
tends to zero.
\end{counterexample}

\section{Standard Model--Einstein interaction ledger}

V91 reduces the interacting boundary problem to two estimates:
\begin{enumerate}
\item a cutoff-independent relative-form constant
      \(\alpha<1\) for the renormalized gauge, Higgs, Yukawa,
      fermionic, and gravitational boundary interactions;
\item common-form convergence strong enough to imply norm-resolvent
      convergence after gauge quotient and anomaly-consistent
      determinant-line transport.
\end{enumerate}
The theorem permits a negative interaction and does not assume an
interacting gap.  It still needs the bosonic free gap used to make
the unperturbed bosonic Fock heat factor nuclear.  A separate
infrared regulator can be used only if its removal preserves the
uniform estimates.

\begin{assessment}[Physical status]
\label{ass:v15v91-status}
V91 proves a uniform interacting Gibbs theorem and trace-norm
continuum criterion from relative form domination and
norm-resolvent convergence.  It transfers the free Fock work of
V89 into a mathematically complete conditional interaction
compiler.

No cutoff-independent relative-form estimate has been proved for
the full four-dimensional Standard Model--Einstein interaction,
and no removal of a bosonic infrared regulator has been justified.
\end{assessment}

\begin{decision}[V92 creation--annihilation form bounds]
\label{dec:v15v91-next}
Derive the elementary Fock estimates needed to test (91.1):
number-operator bounds for creation and annihilation maps, smeared
field bounds relative to \(d\Gamma(K)\), and Wick-monomial bounds
on fixed ultraviolet and infrared cutoffs.  Identify exactly which
constants diverge when either cutoff is removed, rather than
assuming the uniform relative bound.
\end{decision}

\begin{firewall}[Claim boundary after V91]
\label{fire:v15v91}
V91 is conditional on uniform relative-form and convergence
hypotheses.  It does not establish those hypotheses for the
renormalized Standard Model--Einstein interaction, prove a bosonic
mass gap, or construct the full continuum.
\end{firewall}

\section*{Version V91 append-only record}
\addcontentsline{toc}{section}{Version V91 append-only record}
The complete V90 source was retained before this continuation.  It
had \(890431\) bytes, \(26365\) lines, and SHA--256
\[
 \texttt{D5C4E3CED26C2D0A3A841DD5BF1BCC5539A56EA0F62E15C2B2A54822F2AE1B99}.
\]
V91 proved uniform interacting heat trace bounds from a strict
relative-form estimate, upgraded norm-resolvent convergence to
trace-norm heat convergence using free Fock tails, and established
convergence of normalized boundary Gibbs states.
\chapter{Fifteenth-cycle continuation V92: creation--annihilation bounds and cutoff losses}
\label{sec:v15v92}

V91 reduced interacting heat control to a uniform relative-form
estimate.  This note derives the elementary Fock inequalities behind
that estimate and tracks their ultraviolet and infrared constants.
The outcome is deliberately asymmetric: linear and quadratic
smeared fields admit useful energy bounds, while higher bosonic
Wick words require a stronger reference energy or a positivity
argument.

\section{Bosonic energy estimates}

Let \(K\geq0\) be self-adjoint on the one-particle space
\({\cal h}\), and let
\[
 H_0=d\Gamma_+(K),\qquad N=d\Gamma_+(I)                \tag{92.1}
\]
on bosonic Fock space.  Creation and annihilation operators are
denoted by \(a^\dagger(f)\) and \(a(f)\).

\begin{theorem}[One-particle energy bounds]
\label{thm:v15v92-energy}
If \(f\in{\cal D}(K^{-1/2})\), with the inverse equal to zero only
after an explicitly chosen kernel projection, then for every
\(\Psi\in{\cal D}(H_0^{1/2})\),
\[
 \|a(f)\Psi\|
 \leq \|K^{-1/2}f\|\,\|H_0^{1/2}\Psi\|,                \tag{92.2}
\]
\[
 \|a^\dagger(f)\Psi\|^2
 \leq \|K^{-1/2}f\|^2\|H_0^{1/2}\Psi\|^2
      \|f\|^2\|\Psi\|^2.                              \tag{92.3}
\]
If \(K\geq\mu I\), then also
\[
 \|a(f)\Psi\|\leq\|f\|\,\|N^{1/2}\Psi\|,
 \qquad
 \|a^\dagger(f)\Psi\|
 \leq\|f\|\,\|(N+1)^{1/2}\Psi\|,                      \tag{92.4}
\]
and \(N\leq\mu^{-1}H_0\).
\end{theorem}

\begin{proof}
On the \(n\)-particle sector, write the annihilation integral in
the spectral representation of \(K\).  Cauchy--Schwarz with weights
\(K^{-1/2}\) and \(K^{1/2}\), followed by integration over the
remaining variables, gives
\[
 \|a(f)\Psi\|^2
 \leq\|K^{-1/2}f\|^2
 \langle\Psi,d\Gamma(K)\Psi\rangle.                   \tag{92.5}
\]
The canonical commutation relation gives
\[
 \|a^\dagger(f)\Psi\|^2
 =\|a(f)\Psi\|^2+\|f\|^2\|\Psi\|^2,                   \tag{92.6}
\]
which proves (92.3).  Repeating the argument with \(K=I\) gives
(92.4).  Finally \(K\geq\mu I\) implies
\(d\Gamma(K)\geq\mu d\Gamma(I)\) sector by sector.
\end{proof}

For the real smeared field
\[
 \Phi(f)=a(f)+a^\dagger(f),                            \tag{92.7}
\]
the energy estimate has an immediate form consequence.

\begin{corollary}[Linear fields are infinitesimally form bounded]
\label{cor:v15v92-linear}
For \(f\in{\cal D}(K^{-1/2})\), every \(\epsilon>0\), and
\(\Psi\in{\cal D}(H_0^{1/2})\),
\[
 |\langle\Psi,\Phi(f)\Psi\rangle|
 \leq \epsilon\langle\Psi,H_0\Psi\rangle
 +{\|K^{-1/2}f\|^2\over\epsilon}\|\Psi\|^2.            \tag{92.8}
\]
Moreover,
\[
 \|\Phi(f)\Psi\|^2
 \leq4\|K^{-1/2}f\|^2\langle\Psi,H_0\Psi\rangle
      +2\|f\|^2\|\Psi\|^2.                            \tag{92.9}
\]
\end{corollary}

\begin{proof}
The expectation of \(a^\dagger(f)\) is the conjugate of the
expectation of \(a(f)\).  Thus (92.2) and
\(2xy\leq\epsilon x^2+\epsilon^{-1}y^2\) give (92.8).
The inequality
\(\|x+y\|^2\leq2\|x\|^2+2\|y\|^2\), together with
(92.2)--(92.3), gives (92.9).
\end{proof}

\begin{corollary}[A quadratic-field test]
\label{cor:v15v92-quadratic}
For a negative quadratic perturbation
\(-g\Phi(f)^2\), \(g\geq0\), one has
\[
 -g\|\Phi(f)\Psi\|^2
 \geq-4g\|K^{-1/2}f\|^2
        \langle\Psi,H_0\Psi\rangle
      -2g\|f\|^2\|\Psi\|^2.                           \tag{92.10}
\]
It satisfies the strict V91 form condition whenever
\[
 4g\|K^{-1/2}f\|^2<1.                                 \tag{92.11}
\]
A positive quadratic perturbation needs no such smallness for the
lower bound.
\end{corollary}

\section{Fermionic fields and mixed vertices}

On fermionic Fock space the canonical anticommutation relations give
a different estimate.

\begin{proposition}[Fermionic boundedness]
\label{prop:v15v92-car}
For every \(g\in{\cal h}\),
\[
 \|a(g)\|_{\rm op}\leq\|g\|,
 \qquad
 \|a^\dagger(g)\|_{\rm op}\leq\|g\|.                  \tag{92.12}
\]
Consequently a product of \(r\) smeared fermionic creation or
annihilation maps has norm at most \(\prod_{\ell=1}^r\|g_\ell\|\).
\end{proposition}

\begin{proof}
The CAR identity
\(a(g)a^\dagger(g)+a^\dagger(g)a(g)=\|g\|^2I\) and
positivity imply both one-factor bounds.  Submultiplicativity gives
the product estimate.
\end{proof}

\begin{corollary}[Fixed-cutoff Yukawa form]
\label{cor:v15v92-yukawa}
Let \(B\) be a bounded fermionic bilinear and let
\[
 Y=\Phi(f)\otimes B
\]
on a boson--fermion tensor-product Fock space.  Then, for every
\(\epsilon>0\),
\[
 |\langle\Psi,Y\Psi\rangle|
 \leq \epsilon\langle\Psi,(H_0\otimes I)\Psi\rangle
 +{\|B\|^2\|K^{-1/2}f\|^2\over\epsilon}\|\Psi\|^2.     \tag{92.13}
\]
\end{corollary}

\begin{proof}
Apply (92.2) to the bosonic factor and the operator bound for \(B\),
then use the same weighted Cauchy inequality as in (92.8).
\end{proof}

Formula (92.13) is useful at fixed cutoff.  It is not uniform unless
both the bosonic energy norm of \(f\) and the fermionic bilinear norm
stay bounded.

\section{The degree barrier for bosonic words}

\begin{lemma}[Particle-number word estimate]
\label{lem:v15v92-words}
Let \(A_\ell\) be either \(a(f_\ell)\) or
\(a^\dagger(f_\ell)\), \(1\leq\ell\leq r\).  On the finite-particle
core,
\[
 \|A_1\cdots A_r\Psi\|
 \leq r^{r/2}\left(\prod_{\ell=1}^r\|f_\ell\|\right)
       \|(N+1)^{r/2}\Psi\|.                           \tag{92.14}
\]
Hence a product of \(r\) fields \(\Phi(f_\ell)\) obeys the same
bound with the additional factor \(2^r\).
\end{lemma}

\begin{proof}
On the initial \(n\)-particle sector, each successive creation or
annihilation map has norm at most
\(\|f_\ell\|\sqrt{n+r}\).  Their product is bounded by
\((n+r)^{r/2}\prod_\ell\|f_\ell\|\).  Since
\(n+r\leq r(n+1)\), summation over sectors gives (92.14).
Expanding \(r\) fields produces \(2^r\) words.
\end{proof}

\begin{assessment}[The degree barrier]
\label{ass:v15v92-degree}
Lemma~\ref{lem:v15v92-words} controls a degree-\(r\) bosonic Wick
word by \((N+1)^{r/2}\).  With a free Hamiltonian comparable only
to \(N\), this yields a direct relative bound only through degree
two.  Cubic and quartic gauge or Higgs words require additional
structure: positivity, spatial estimates, cancellations,
renormalization, or inclusion of a stable polynomial in the
reference Hamiltonian.
\end{assessment}

\section{Explicit cutoff scaling in three spatial dimensions}

For a scalar dispersion
\(\omega(k)=(|k|^2+m^2)^{1/2}\), \(m>0\), define the point-field
cutoff vector
\[
 f_{\Lambda,x}(k)
 ={ {\bf1}_{\{|k|\leq\Lambda\}}e^{-ik\cdot x}
    \over(2\omega(k))^{1/2}}
\quad\hbox{in }L^2(\mathbb R^3,dk/(2\pi)^3).           \tag{92.15}
\]

\begin{proposition}[Ultraviolet growth of the elementary constants]
\label{prop:v15v92-uv}
As \(\Lambda\to\infty\),
\[
 \|f_{\Lambda,x}\|^2
 ={1\over2}\int_{|k|\leq\Lambda}{dk\over(2\pi)^3\omega(k)}
 =\Theta(\Lambda^2),                                  \tag{92.16}
\]
\[
 \|\omega^{-1/2}f_{\Lambda,x}\|^2
 ={1\over2}\int_{|k|\leq\Lambda}{dk\over(2\pi)^3\omega(k)^2}
 =\Theta(\Lambda).                                    \tag{92.17}
\]
The constants are independent of \(x\), but neither is uniform in
the ultraviolet cutoff.
\end{proposition}

\begin{proof}
The phase has modulus one.  In spherical coordinates the two radial
integrands are, up to the fixed factor \(4\pi/(2\pi)^3\),
\[
 {r^2\over2(r^2+m^2)^{1/2}},
 \qquad
 {r^2\over2(r^2+m^2)}.                                \tag{92.18}
\]
For \(r\geq m\) these are bounded above and below by fixed positive
multiples of \(r\) and \(1\), respectively.  Integration from
\(m\) to \(\Lambda\) proves (92.16)--(92.17); the bounded interval
below \(m\) contributes only a constant.
\end{proof}

\begin{corollary}[Why the fixed-cutoff form estimate is not uniform]
\label{cor:v15v92-uv-loss}
For the point-field family (92.15), the coefficient of \(H_0\) in
(92.10) grows as \(\Theta(g\Lambda)\), while the scalar remainder
grows as \(\Theta(g\Lambda^2)\).  For a Yukawa term (92.13), even
before tracking its fermionic cutoff vectors, the scalar remainder
contains a factor \(\Theta(\Lambda)\).
\end{corollary}

\begin{proof}
Insert (92.16)--(92.17) into (92.10) and (92.13).
\end{proof}

Spatial smearing changes the result.  If a test function has Fourier
transform decaying fast enough that
\[
 \int_{\mathbb R^3}{|\widehat g(k)|^2\over\omega(k)^2}\,dk<\infty,
                                                               \tag{92.19}
\]
then the energy-weighted constant converges as
\(\Lambda\to\infty\).  A continuum construction must prove estimates
for its actual local interaction density, not replace it silently by
a fixed smooth smearing.

\section{Infrared failure and a scalar counterexample}

\begin{counterexample}[Zero-mode energy norm]
\label{cnt:v15v92-ir}
If \(Ke_0=0\) and \(\langle e_0,f\rangle\ne0\), then
\[
 \|K^{-1/2}f\|=\infty.                                \tag{92.20}
\]
Thus (92.2) and the V91 relative-form route fail before second
quantization.  Projecting out \(e_0\) is valid only when the
constraint quotient identifies that projection coherently.
\end{counterexample}

\begin{counterexample}[A diverging linear source]
\label{cnt:v15v92-linear-source}
For one oscillator with \(H_0=\omega N\), let
\(V_n=c_n(a+a^\dagger)\), where \(c_n\to\infty\).  Completing the
square gives
\[
 H_0+V_n
 =\omega\left(a^\dagger+{c_n\over\omega}\right)
          \left(a+{c_n\over\omega}\right)
   -{c_n^2\over\omega}.                               \tag{92.21}
\]
Each \(V_n\) is infinitesimally form bounded, but every uniform
lower-bound constant \(\beta\) in V91 is at least
\(c_n^2/\omega\) and therefore diverges.
\end{counterexample}

\section{Interaction ledger}

At a fixed ultraviolet cutoff and after a controlled zero-mode
quotient, V92 supplies:
\begin{enumerate}
\item infinitesimal form bounds for linear bosonic sources;
\item strict relative bounds for sufficiently small negative
      quadratic pieces;
\item bounded fermionic words and a mixed Yukawa estimate;
\item number-operator bounds for all finite Wick words.
\end{enumerate}
It also proves why these facts alone do not supply the uniform
\(\alpha,\beta\) of V91 when the cutoff is removed.  In particular,
bare point-field estimates lose powers of \(\Lambda\), and a free
number reference does not dominate cubic or quartic bosonic words.

\begin{assessment}[Physical status]
\label{ass:v15v92-status}
V92 rigorously verifies the V91 hypothesis for elementary smeared
linear and some quadratic cutoff interactions, and identifies the
exact ultraviolet, infrared, and polynomial-degree losses for the
remaining vertices.  No cancellation between Standard Model
counterterms has been assumed.
\end{assessment}

\begin{decision}[V93 stable polynomial reference]
\label{dec:v15v92-next}
Go upstream of the degree barrier.  Replace the purely free
reference by a stable cutoff Higgs polynomial, prove lower and form
comparisons for Wick-shifted quartics, and determine which divergent
vacuum and quadratic counterterms can be separated from normalized
states.  Do not treat a cutoff-dependent scalar energy shift as a
uniform partition-function bound.
\end{decision}

\begin{firewall}[Claim boundary after V92]
\label{fire:v15v92}
The estimates are fixed-cutoff operator inequalities with explicit
losses.  They do not prove uniform renormalized form domination,
construct a local four-dimensional interaction, or establish the
full Standard Model--Einstein continuum.
\end{firewall}

\section*{Version V92 append-only record}
\addcontentsline{toc}{section}{Version V92 append-only record}
The complete V91 source was retained before this continuation.  It
had \(901702\) bytes, \(26680\) lines, and SHA--256
\[
 \texttt{B88519F3EE13208BA86BA53D3A42F6C4099513FB33AA23849EF8767D14EC8015}.
\]
V92 proved energy and number bounds for bosonic fields, boundedness
of fermionic words, fixed-cutoff Yukawa estimates, and the explicit
ultraviolet and zero-mode losses that obstruct a uniform V91 form
bound.
\chapter{Fifteenth-cycle continuation V93: stable Higgs references and vacuum shifts}
\label{sec:v15v93}

The free number Hamiltonian in V92 does not dominate quartic
bosonic words.  The correct upstream move is to include the stable
quartic part in the reference form.  At finite cutoff this converts
lower-degree interactions into relative-form perturbations.  The
renormalized continuum problem remains in the convergence of these
new reference Hamiltonians.

\section{Exact finite-cutoff Wick algebra}

Let \(X=(X_1,\ldots,X_r)\) be commuting real field coordinates with
centered Gaussian contraction
\[
 {\mathbb E}(X_aX_b)=C\delta_{ab},\qquad C\geq0.        \tag{93.1}
\]
Wick ordering with respect to this contraction gives
\[
 :|X|^2:_C=|X|^2-rC,                                  \tag{93.2}
\]
\[
 :|X|^4:_C
 =|X|^4-2(r+2)C|X|^2+r(r+2)C^2.                      \tag{93.3}
\]

\begin{lemma}[Derivation of the radial Wick quartic]
\label{lem:v15v93-wick}
Equations (93.2)--(93.3) hold as polynomial identities.
\end{lemma}

\begin{proof}
The quadratic identity subtracts the \(r\) self-contractions.
Expand
\[
 |X|^4=\sum_{a,b}X_a^2X_b^2.                           \tag{93.4}
\]
For \(a=b\),
\(:X_a^4:=X_a^4-6CX_a^2+3C^2\).  For \(a\ne b\),
\[
 :X_a^2X_b^2:
 =X_a^2X_b^2-CX_a^2-CX_b^2+C^2.                      \tag{93.5}
\]
Summing the ordered pairs gives coefficient
\(-6C-2(r-1)C=-2(r+2)C\) for each \(X_a^2\), and constant
\(3rC^2+r(r-1)C^2=r(r+2)C^2\).
\end{proof}

\begin{theorem}[Exact stable counterterm choice]
\label{thm:v15v93-stable-wick}
Let \(\lambda>0\), \(m_R^2\geq0\), and define
\[
 \nu_C=2\lambda(r+2)C+m_R^2,\qquad
 E_C=\lambda r(r+2)C^2+r m_R^2C.                     \tag{93.6}
\]
Then the Wick polynomial
\[
 W_C(X)
 =\lambda:|X|^4:_C+\nu_C:|X|^2:_C+E_C               \tag{93.7}
\]
satisfies the exact identity
\[
 W_C(X)=\lambda|X|^4+m_R^2|X|^2\geq0.                \tag{93.8}
\]
\end{theorem}

\begin{proof}
Insert (93.2)--(93.3) into (93.7).  The coefficient of
\(|X|^2\) is
\(\nu_C-2\lambda(r+2)C=m_R^2\).  The remaining scalar is
\[
 \lambda r(r+2)C^2-r\nu_CC+E_C=0.                    \tag{93.9}
\]
This proves (93.8).
\end{proof}

For the Standard Model Higgs doublet there are \(r=4\) real
components.  The algebraic coefficients in this convention become
\[
 \nu_C=12\lambda C+m_R^2,\qquad
 E_C=24\lambda C^2+4m_R^2C.                          \tag{93.10}
\]
These identities do not assert that the physical counterterms have
already been constructed; they state exactly what a Gaussian Wick
reference contributes at a fixed cutoff.

\section{Stable polynomial reference form}

Let \(H_{G,\Lambda}\geq0\) be a finite-cutoff Gaussian Hamiltonian
with compact resolvent on \(L^2(\mathbb R^{M_\Lambda})\).  Let
\(X_\ell\) denote a finite family of commuting cutoff field
coordinates with positive weights \(w_\ell\), and put
\[
 U_\Lambda
 =\sum_\ell w_\ell
   \bigl(\lambda|X_\ell|^4+m_R^2|X_\ell|^2\bigr),
 \qquad
 H_{{\rm ref},\Lambda}
 =H_{G,\Lambda}\dotplus U_\Lambda.                    \tag{93.11}
\]

\begin{proposition}[Nuclear stable reference]
\label{prop:v15v93-reference}
The form (93.11) is closed and nonnegative on
\({\cal D}(H_{G,\Lambda}^{1/2})\cap{\cal D}(U_\Lambda^{1/2})\),
has compact resolvent, and
\[
 \operatorname{Tr}e^{-tH_{{\rm ref},\Lambda}}
 \leq\operatorname{Tr}e^{-tH_{G,\Lambda}}
 \qquad(t>0).                                         \tag{93.12}
\]
\end{proposition}

\begin{proof}
The sum of two nonnegative closed forms is closed on the
intersection of their form domains.  Its form norm dominates the
Gaussian form norm, so its unit form ball embeds compactly whenever
the Gaussian one does.  Thus the resolvent is compact.  The min--max
principle gives
\(\lambda_j(H_{{\rm ref},\Lambda})
\geq\lambda_j(H_{G,\Lambda})\); exponentiation and summation prove
(93.12).
\end{proof}

\begin{lemma}[Absorption of lower polynomial degrees]
\label{lem:v15v93-young}
For \(0<p<4\), \(\delta>0\), and \(x\geq0\),
\[
 x^p\leq\delta x^4+
 \left(1-{p\over4}\right)
 \left({p\over4\delta}\right)^{p/(4-p)}.              \tag{93.13}
\]
\end{lemma}

\begin{proof}
The maximum of \(x^p-\delta x^4\) occurs at
\(x^{4-p}=p/(4\delta)\).  Substitution gives the displayed
constant.  The endpoints give no larger value.
\end{proof}

\begin{theorem}[Finite-cutoff polynomial form compiler]
\label{thm:v15v93-polynomial}
Suppose a real multiplication perturbation has the pointwise bound
\[
 |P_\Lambda(X)|
 \leq\sum_\ell w_\ell\sum_{0\leq p<4}
       c_{p,\Lambda}|X_\ell|^p.                       \tag{93.14}
\]
For every prescribed \(0<\alpha<1\), there is an explicit
\(\beta_\Lambda<\infty\) such that
\[
 |\langle\Psi,P_\Lambda\Psi\rangle|
 \leq\alpha\langle\Psi,U_\Lambda\Psi\rangle
      +\beta_\Lambda\|\Psi\|^2.                       \tag{93.15}
\]
One may take \(\beta_\Lambda\) to be the weighted sum of the
constants in (93.13), after assigning positive
\(\delta_{p,\Lambda}\) with
\[
 \sum_{p<4}c_{p,\Lambda}\delta_{p,\Lambda}
 \leq\alpha\lambda.                                   \tag{93.16}
\]
\end{theorem}

\begin{proof}
Apply Lemma~\ref{lem:v15v93-young} at each field coordinate, with
the chosen \(\delta_{p,\Lambda}\), multiply by
\(w_\ell c_{p,\Lambda}\), and sum.  The quartic terms are bounded
by \(\alpha U_\Lambda\); discarding its nonnegative quadratic part
only weakens the bound.  The remaining constants give
\(\beta_\Lambda\).  Integration against \(|\Psi|^2\) proves
(93.15).
\end{proof}

The theorem solves the polynomial degree barrier at each finite
cutoff.  Uniformity requires bounded coefficient combinations and
a controlled weighted volume \(\sum_\ell w_\ell\).  If Wick
contractions make \(c_{p,\Lambda}\) diverge, the theorem honestly
returns a divergent \(\beta_\Lambda\).

\section{Scalar energy shifts and what they preserve}

The vacuum term in (93.6) is a scalar multiple of the identity after
spatial integration.  Such a term cancels from normalized states
but not from the boundary amplitude.

\begin{proposition}[Vacuum-shift dichotomy]
\label{prop:v15v93-vacuum-shift}
Let \(H\) be self-adjoint with \(e^{-tH}\) trace class and let
\(\mathcal E\in\mathbb R\).  Then
\[
 e^{-t(H+\mathcal EI)}
 =e^{-t\mathcal E}e^{-tH},\qquad
 Z_{H+\mathcal E}(t)=e^{-t\mathcal E}Z_H(t),           \tag{93.17}
\]
whereas
\[
 {e^{-t(H+\mathcal EI)}\over Z_{H+\mathcal E}(t)}
 ={e^{-tH}\over Z_H(t)}.                              \tag{93.18}
\]
\end{proposition}

\begin{proof}
The scalar commutes with \(H\), giving the functional-calculus
identity and its trace.  The factors cancel in the quotient.
\end{proof}

\begin{counterexample}[Normalized convergence without amplitude
convergence]
\label{cnt:v15v93-shift}
Fix one Gibbs Hamiltonian \(H\) and take
\(H_n=H+\mathcal E_nI\).  All normalized Gibbs states are identical.
If \(\mathcal E_n\to+\infty\), then
\[
 \|e^{-tH_n}\|_1\to0;
\]
if \(\mathcal E_n\to-\infty\), then the trace norm diverges.  Thus
normalized-state convergence does not determine a sewing
normalization.
\end{counterexample}

\begin{proof}
Apply Proposition~\ref{prop:v15v93-vacuum-shift}.
\end{proof}

For a prescribed counterterm \(\mathcal E_\Lambda\), define the
renormalized amplitude
\[
 \widehat T_\Lambda(t)
 =e^{t\mathcal E_\Lambda}e^{-t(H_\Lambda+\mathcal E_\Lambda I)}
 =e^{-tH_\Lambda}.                                    \tag{93.19}
\]
Equation (93.19) is an identity only after the counterterm has been
fixed.  It cannot be used to choose an arbitrary normalization at
each cutoff, because V31--V38 and V81--V84 require determinant,
stress, sewing, and phase compatibilities.

\section{Exact Higgs separation at varying covariance}

\begin{corollary}[Wick covariance ledger]
\label{cor:v15v93-covariance}
At each cutoff \(\Lambda\), let the local Higgs covariance be
\(C_\Lambda\).  Choosing \(\nu_{C_\Lambda}\) and
\(E_{C_\Lambda}\) as in (93.6) turns the Wick-ordered local
polynomial exactly into the nonnegative reference (93.8).
If \(C_\Lambda\to\infty\), then
\[
 \nu_{C_\Lambda}=\Theta(C_\Lambda),\qquad
 E_{C_\Lambda}=\Theta(C_\Lambda^2)                    \tag{93.20}
\]
for fixed \(\lambda>0\).
\end{corollary}

\begin{proof}
The identity is Theorem~\ref{thm:v15v93-stable-wick}.  The leading
terms in (93.6) give (93.20).
\end{proof}

The quadratic divergence changes the operator and must be included
in the stable reference or canceled by a proved renormalization
condition.  The scalar divergence can be removed from normalized
states by (93.18), but its value remains part of the gravitational
source and the gluing amplitude.

\section{Continuum handoff and remaining obstruction}

\begin{theorem}[Conditional V91 handoff with a stable reference]
\label{thm:v15v93-handoff}
Suppose the stable cutoff references
\(H_{{\rm ref},\Lambda}\) satisfy
\[
 e^{-sH_{{\rm ref},\Lambda}}
 \longrightarrow e^{-sH_{\rm ref}}
 \quad\hbox{in trace norm for every }s>0,              \tag{93.21}
\]
after common-carrier transport.  Suppose the remaining interaction
forms satisfy (93.15) with one \(\alpha<1\), one
\(\sup_\Lambda\beta_\Lambda<\infty\), and their full form sums
converge in norm resolvent.  Then the interacting heat factors and
normalized states converge in trace norm.
\end{theorem}

\begin{proof}
Use \(H_{{\rm ref},\Lambda}\) as \(H_{0,n}\) in
Theorem~\ref{thm:v15v91-gibbs}, then apply
Corollary~\ref{cor:v15v91-states}.
\end{proof}

Condition (93.21) is a constructive interacting reference problem,
not a consequence of the finite-cutoff polynomial algebra.  It
contains the non-Gaussian compactness and renormalization work for
the Higgs sector.

\begin{assessment}[Physical status]
\label{ass:v15v93-status}
V93 proves an exact Wick counterterm identity for an
\(r\)-component Higgs field, a stable finite-cutoff reference
theorem, and explicit absorption of all lower polynomial degrees.
It also proves that vacuum shifts cancel in normalized states while
remaining essential to partition functions, gluing, and
gravitational response.

The cutoff-uniform convergence (93.21), the physical counterterm
flow, and compatibility with gauge and chiral sectors remain
unproved.
\end{assessment}

\begin{decision}[V94 signed vertices over the stable reference]
\label{dec:v15v93-next}
Test gauge--Higgs cubic and quartic cross terms, Yukawa terms, and
bounded fermionic factors against the stable reference using
Young and Cauchy inequalities.  Separate estimates that are uniform
at fixed spatial smearing from constants that inherit the
ultraviolet losses of V92.  Treat the gravitational conformal
direction separately rather than hiding it in a positive Higgs
reference.
\end{decision}

\begin{firewall}[Claim boundary after V93]
\label{fire:v15v93}
V93 is a finite-cutoff stability and normalization theorem.  It
does not construct the renormalized Higgs measure, determine the
Standard Model counterterms, control dynamical gravity, or prove
the full continuum.
\end{firewall}

\section*{Version V93 append-only record}
\addcontentsline{toc}{section}{Version V93 append-only record}
The complete V92 source was retained before this continuation.  It
had \(913639\) bytes, \(27029\) lines, and SHA--256
\[
 \texttt{BBC60B0B48106989ECEA978B42DD6B64E8359E427DA7AD7ABB89022FF5F39E74}.
\]
V93 rebased the bosonic interaction on a stable quartic reference,
computed its exact Wick mass and vacuum counterterms, proved
lower-degree form absorption, and separated normalized Gibbs-state
invariance from unnormalized amplitude convergence.
\chapter{Fifteenth-cycle continuation V94: signed vertices, gauge flats, and gravitational separation}
\label{sec:v15v94}

V93 supplies a stable Higgs quartic reference.  This note tests the
remaining schematic Standard Model vertices against stable
reservoirs.  Covariant derivative and curvature cross terms are
best controlled as parts of positive squares.  Yukawa terms are
form bounded at fixed cutoff.  Two independent obstructions remain:
commuting Yang--Mills directions evade the commutator quartic, and
matter stability cannot repair the Euclidean gravitational
conformal direction.

\section{Covariant squares}

\begin{lemma}[Two-reservoir cross estimate]
\label{lem:v15v94-square}
For vectors \(X,Y\) in a Hilbert space and every \(\epsilon>0\),
\[
 2|\operatorname{Re}\langle X,Y\rangle|
 \leq\epsilon\|X\|^2+\epsilon^{-1}\|Y\|^2.            \tag{94.1}
\]
In particular,
\[
 \|X+Y\|^2\geq0                                       \tag{94.2}
\]
and its signed cross term is relatively form bounded by the two
diagonal reservoirs.
\end{lemma}

\begin{proof}
Positivity of
\(\|\epsilon^{1/2}X-\epsilon^{-1/2}e^{i\theta}Y\|^2\),
with \(\theta\) chosen to align the inner product, gives (94.1).
Equation (94.2) is immediate.
\end{proof}

At a finite geometric and ultraviolet cutoff, write schematically
\[
 F_A=dA+g[A,A],\qquad D_A\phi=d\phi+gA\phi.            \tag{94.3}
\]
Lemma~\ref{lem:v15v94-square} gives
\[
 2g|\langle dA,[A,A]\rangle|
 \leq\epsilon\|dA\|^2+\epsilon^{-1}g^2\|[A,A]\|^2,    \tag{94.4}
\]
\[
 2g|\langle d\phi,A\phi\rangle|
 \leq\epsilon\|d\phi\|^2+\epsilon^{-1}g^2\|A\phi\|^2. \tag{94.5}
\]

\begin{proposition}[Do not split a positive covariant square]
\label{prop:v15v94-covariant}
If the cutoff form contains the complete terms
\(\|F_A\|^2\) and \(\|D_A\phi\|^2\) with positive coefficients,
they contribute no negative part to the V91 lower bound.  If a
comparison argument splits them, (94.4)--(94.5) control the signed
cubic terms by the corresponding quadratic and quartic pieces with
arbitrarily assigned reciprocal weights.
\end{proposition}

\begin{proof}
The unsplit assertion is positivity.  The split assertion is
Lemma~\ref{lem:v15v94-square} integrated over the cutoff boundary.
\end{proof}

This elementary reorganization matters: estimating the cubic terms
by particle-number words, as in the degree barrier of V92, discards
the square structure and produces a needlessly strong power of
\(N\).

\section{Mixed quartics and finite-dimensional Yukawa absorption}

\begin{lemma}[Mixed quartic absorption]
\label{lem:v15v94-mixed}
For \(a,b\geq0\) and every \(\eta>0\),
\[
 a^2b^2\leq{\eta\over2}a^4+{1\over2\eta}b^4.          \tag{94.6}
\]
\end{lemma}

\begin{proof}
Apply \(2xy\leq x^2+y^2\) to
\(x=\eta^{1/2}a^2\) and \(y=\eta^{-1/2}b^2\).
\end{proof}

\begin{proposition}[Two stable scalar sectors]
\label{prop:v15v94-two-sector}
Suppose a cutoff reference contains
\[
 U(A,\phi)=\lambda_A|A|^4+\lambda_H|\phi|^4,
 \qquad \lambda_A,\lambda_H>0.                         \tag{94.7}
\]
Then multiplication by \(c|A|^2|\phi|^2\) is form bounded relative
to \(U\).  Its relative coefficient can be made strictly below one
provided
\[
 |c|<2(\lambda_A\lambda_H)^{1/2}.                     \tag{94.8}
\]
\end{proposition}

\begin{proof}
Choose \(\eta=(\lambda_A/\lambda_H)^{1/2}\) in (94.6).  Then
\[
 |c||A|^2|\phi|^2
 \leq {|c|\over2(\lambda_A\lambda_H)^{1/2}}
       \bigl(\lambda_A|A|^4+\lambda_H|\phi|^4\bigr).   \tag{94.9}
\]
The coefficient is below one exactly under (94.8).
\end{proof}

For a cutoff Yukawa vertex, the fermionic factor is bounded by
Proposition~\ref{prop:v15v92-car}.  Quartic Higgs stability yields
another useful estimate.

\begin{proposition}[Yukawa absorption into the Higgs quartic]
\label{prop:v15v94-yukawa}
Let \(B_\Lambda\) be bounded and self-adjoint on a cutoff fermionic
Fock space, let \(y\in\mathbb R\), and let \(\phi\) be a commuting
Higgs coordinate.  For every \(\delta>0\),
\[
 \bigl|\langle\Psi,y\,\phi\otimes B_\Lambda\Psi\rangle\bigr|
 \leq\delta\lambda\langle\Psi,|\phi|^4\Psi\rangle
 +C_{\delta,\Lambda}\|\Psi\|^2,                       \tag{94.10}
\]
where
\[
 C_{\delta,\Lambda}
 ={3\over4}
 (|y|\|B_\Lambda\|)^{4/3}(4\delta\lambda)^{-1/3}.      \tag{94.11}
\]
\end{proposition}

\begin{proof}
In the direct-integral representation of the commuting Higgs
coordinate,
\[
 |\langle\Psi(\phi),B_\Lambda\Psi(\phi)\rangle|
 \leq\|B_\Lambda\|\,\|\Psi(\phi)\|^2.                 \tag{94.12}
\]
The maximum over \(x\geq0\) of
\(|y|\|B_\Lambda\|x-\delta\lambda x^4\) is (94.11).
Integrate the resulting pointwise inequality.
\end{proof}

At fixed cutoff, choose \(\delta\) below the unused fraction of the
Higgs quartic reservoir.  This proves a V91 form bound.  Uniformity
still requires \(\sup_\Lambda\|B_\Lambda\|<\infty\).
For local cutoff fermion fields, Proposition~\ref{prop:v15v92-car}
expresses \(\|B_\Lambda\|\) through one-particle smearing norms,
which generally diverge as the smearing is removed.

\section{The Yang--Mills quartic is not radially coercive}

The positive term \(\|[A,A]\|^2\) in (94.4) is a reservoir only
transverse to commuting configurations.

\begin{proposition}[Cartan flat-direction obstruction]
\label{prop:v15v94-cartan}
Let \(\mathfrak g\) be a compact Lie algebra of positive rank.  There
is no \(c>0\) such that, for every finite family \(A_i\in\mathfrak g\),
\[
 \sum_{i<j}\|[A_i,A_j]\|^2
 \geq c\left(\sum_i\|A_i\|^2\right)^2.                \tag{94.13}
\]
\end{proposition}

\begin{proof}
Choose a nonzero \(H\) in a Cartan subalgebra and take
\(A_1=tH\), with all other \(A_i=0\).  Every bracket vanishes while
the right-hand side equals \(ct^4\|H\|^4>0\) for \(t\ne0\).
\end{proof}

\begin{corollary}[Conditional nature of gauge quartic absorption]
\label{cor:v15v94-gauge-flat}
The replacement of \(\lambda_A|A|^4\) in
Proposition~\ref{prop:v15v94-two-sector} by the Yang--Mills
commutator quartic is invalid unless a separate gauge-slice,
derivative, holonomy, or Higgs mechanism controls the commuting
directions.
\end{corollary}

\begin{proof}
On the sequence in Proposition~\ref{prop:v15v94-cartan}, the
commutator reservoir is zero while \(|A|^4\) is unbounded.
\end{proof}

The local Wilson-symbol floor in YM V41 is not such a nonlinear
coercivity theorem: V90 already showed that it covers only one
momentum cell and has not been lifted to the common quotient
Hamiltonian.

\section{Matter stability cannot repair the conformal direction}

Let \(q_{\rm grav}\) be a gravitational form and
\(q_{\rm mat}\geq0\) a stable matter form.  The following simple
restriction principle prevents a false combined lower bound.

\begin{lemma}[Tensor-factor restriction]
\label{lem:v15v94-tensor}
Suppose there are normalized gravitational vectors \(u_n\) with
\[
 q_{\rm grav}[u_n]\longrightarrow-\infty.             \tag{94.14}
\]
Let \(v\) be one normalized matter vector with
\(q_{\rm mat}[v]<\infty\), and suppose the cross interaction is
uniformly bounded on \(u_n\otimes v\).  Then the combined form on
\(u_n\otimes v\) is unbounded below.  In particular, no positive
matter quartic can by itself provide a V91 lower bound for the
combined theory.
\end{lemma}

\begin{proof}
On product vectors the sum of the uncoupled forms is
\[
 q_{\rm grav}[u_n]+q_{\rm mat}[v].                    \tag{94.15}
\]
The second term is fixed and the cross term is bounded by
hypothesis, whereas the first tends to \(-\infty\).
\end{proof}

V66 exhibited precisely such negative conformal behavior for the
Euclidean Einstein action.  The Higgs field may be fixed at a
finite-energy configuration, including its vacuum sector, along
that sequence.  Therefore the gravitational gate must be handled
by an independently justified contour, constraint reduction, or
reflection-compatible stabilizer.  It cannot be hidden inside the
Higgs estimate of V93.

\section{Conditional finite-cutoff assembly}

\begin{theorem}[Matter-sector form assembly]
\label{thm:v15v94-assembly}
At a fixed cutoff, assume:
\begin{enumerate}
\item the Higgs polynomial is stabilized as in V93;
\item Yang--Mills curvature and gauge--Higgs derivatives are kept
      as their complete positive covariant squares;
\item every additional signed mixed quartic has stable reservoirs
      satisfying the coefficient condition (94.8);
\item the cutoff fermionic bilinears in the Yukawa terms are bounded;
\item gauge-fixing and ghost contributions have a separately proved
      self-adjoint lower-form realization.
\end{enumerate}
Then the signed matter interaction, excluding gravity, admits a
relative-form bound of the V91 type.  Its scalar constant is an
explicit sum of (93.13), (94.11), and the cutoff smearing constants
of V92.
\end{theorem}

\begin{proof}
Keep the covariant squares as nonnegative pieces by
Proposition~\ref{prop:v15v94-covariant}.  Absorb lower Higgs
polynomials by Theorem~\ref{thm:v15v93-polynomial}, admissible mixed
quartics by Proposition~\ref{prop:v15v94-two-sector}, and Yukawa
terms by Proposition~\ref{prop:v15v94-yukawa}.  Choose the finitely
many assigned relative coefficients so their sum is below one.
The separately assumed gauge-fixing and ghost bound completes the
form estimate.
\end{proof}

The theorem is intentionally finite-cutoff and conditional in its
gauge-fixing row.  Proposition~\ref{prop:v15v94-cartan} and
Lemma~\ref{lem:v15v94-tensor} identify two directions that the
polynomial algebra does not control.

\begin{assessment}[Physical status]
\label{ass:v15v94-status}
V94 proves that the main signed matter vertices can be controlled
without a particle-number estimate when their positive square and
quartic structure is retained.  It isolates the exact coefficient
condition for mixed quartics and gives a fixed-cutoff Yukawa bound.

Uniform fermionic smearing, nonlinear gauge flat directions,
ghost/reflection positivity, and the gravitational conformal
direction remain open.  No combined continuum form has been
constructed.
\end{assessment}

\begin{decision}[V95 companion audit]
\label{dec:v15v94-next}
Perform the compulsory YM/NS audit.  In particular, test whether the
newest YM note has advanced from a local Wilson-symbol cell to a
global quotient coercivity statement relevant to
Proposition~\ref{prop:v15v94-cartan}.  Then continue in V96 with the
gauge-slice and holonomy control actually available.
\end{decision}

\begin{firewall}[Claim boundary after V94]
\label{fire:v15v94}
V94 proves finite-cutoff form estimates and two obstruction
theorems.  It does not prove gauge-flat coercivity, construct the
ghost Hilbert space, stabilize Euclidean gravity, or establish the
full Standard Model--Einstein continuum.
\end{firewall}

\section*{Version V94 append-only record}
\addcontentsline{toc}{section}{Version V94 append-only record}
The complete V93 source was retained before this continuation.  It
had \(925111\) bytes, \(27366\) lines, and SHA--256
\[
 \texttt{E39A19EE6CE3C98AA04A45B6DD94F208B6211DF2635E91E07D9603DC6C41841C}.
\]
V94 controlled covariant-square cross terms, mixed quartics, and
fixed-cutoff Yukawa vertices over the stable Higgs reference, and
proved independent Cartan-flat and gravitational-conformal
obstructions.
\chapter{Fifteenth-cycle continuation V95: companion audit at the gauge-coercivity gate}
\label{sec:v15v95}

This is the compulsory audit after V90.  It occurs immediately after
V94 isolated commuting gauge directions as a missing coercive row.
The audit therefore tests the companion notes specifically for a
global gauge-slice, holonomy, or quotient-Hamiltonian estimate.

\section{Current companion files}

The highest-numbered Yang--Mills filename is
\[
 \texttt{Renewal\_Yang\_Mills\_Mass\_Gap\_Research\_Notes\_V42.tex}.
                                                               \tag{95.1}
\]
It has \(356260\) bytes, \(9752\) lines, and SHA--256
\[
 \texttt{0C111A921ABD8DC8970DB8E19412172D72469898CC94670F92D12CC9BB8DD498}.
                                                               \tag{95.2}
\]
Its bytes, length, and digest are identical to the V41 file audited
in V90.  Its final mathematical section is still labeled V41, and
its final decision still assigns the missing orientation test to
V42.  The V42 filename therefore contains no new mathematical
continuation to import.

The newest Navier--Stokes file remains
\[
 \texttt{Renewal\_NS\_Stability\_Research\_Notes\_V26.tex},       \tag{95.3}
\]
with \(278283\) bytes, \(5319\) lines, and SHA--256
\[
 \texttt{82C887F8C2C69E4F28FAD7C3DC8DE5B9099CB5A4BF431EBA1FFF3E91935978AD}.
                                                               \tag{95.4}
\]
It remains the rank-one Oseen-kernel norm computation read at V90.

\section{Import decision}

\begin{theorem}[V95 companion non-import]
\label{thm:v15v95-nonimport}
The companion checkpoint adds no premise to the conditional
matter-sector assembly of V94.
\end{theorem}

\begin{proof}
The Yang--Mills bytes are unchanged from the V41 theorem already
analyzed in V90.  That theorem gives a Wilson-pencil floor only on
one selected momentum quarter cell and explicitly withholds
full-torus, covariant-operator, continuum-measure, and Hamiltonian
mass-gap conclusions.  It does not control the Cartan family in
Proposition~\ref{prop:v15v94-cartan}.

The Navier--Stokes bytes and result are unchanged.  Its restricted
rank-one kernel norm has no invariant-input theorem for the quantum
gauge form, as proved in Lemma~\ref{lem:v15v90-cone}.  Hence neither
checkpoint supplies a new hypothesis for
Theorem~\ref{thm:v15v94-assembly}.
\end{proof}

\begin{assessment}[Direction after the audit]
\label{ass:v15v95-direction}
The correct next step is not to extrapolate the YM local symbol
atlas.  The missing statement is geometric: after a genuine gauge
choice, curvature must control the connection transverse to flat
and harmonic modes.  The commuting modes then require a separate
holonomy analysis.
\end{assessment}

\begin{decision}[V96 local gauge-slice coercivity]
\label{dec:v15v95-next}
On a fixed compact boundary, prove the linear Coulomb--Hodge
estimate away from harmonic one-forms and its small-connection
nonlinear curvature consequence.  Identify the flat holonomy
moduli and show why compactness of that moduli space alone does not
give a uniform transverse spectral gap across reducible strata.
Keep gauge-slice existence as an explicit hypothesis unless it is
proved.
\end{decision}

\begin{firewall}[Claim boundary after V95]
\label{fire:v15v95}
V95 is a read-only companion audit.  It imports no mass gap,
global Yang--Mills coercivity, Navier--Stokes theorem, or continuum
construction.
\end{firewall}

\section*{Version V95 append-only record}
\addcontentsline{toc}{section}{Version V95 append-only record}
The complete V94 source was retained before this continuation.  It
had \(936304\) bytes, \(27672\) lines, and SHA--256
\[
 \texttt{D7B7E319798E1E929978F8CA01B0DCBB87D907D915F92CFC03642D97E8AE7770}.
\]
V95 rechecked the current YM and NS companions, recorded that the
YM V42 filename is an exact V41 duplicate, imported no theorem, and
selected local gauge-slice and holonomy coercivity for V96.
\chapter{Fifteenth-cycle continuation V96: Coulomb coercivity and flat holonomy strata}
\label{sec:v15v96}

V94 showed that the Yang--Mills commutator quartic vanishes on
commuting configurations.  V95 found no companion theorem closing
that gap.  This note proves the local replacement: after a Coulomb
condition and projection away from the Hodge kernel, small
connections are controlled by their curvature.  It then identifies
the precise obstruction to making that estimate uniform over flat
holonomy moduli.

\section{Linear Hodge estimate}

Let \(\Sigma\) be a closed Riemannian manifold of dimension
\(d\leq4\), and let \(\mathfrak g\) be a compact Lie algebra with an
invariant inner product.  Write
\({\cal H}^1=\ker d\cap\ker d^*\) for the harmonic
\(\mathfrak g\)-valued one-forms.

\begin{proposition}[Coulomb--Hodge coercivity]
\label{prop:v15v96-hodge}
There is \(C_H<\infty\) such that every
\(a\in H^1(\Sigma;T^*\Sigma\otimes\mathfrak g)\) obeys
\[
 \|a\|_{H^1}
 \leq C_H\bigl(
       \|da\|_2+\|d^*a\|_2+\|\Pi_{\cal H}a\|_2\bigr).  \tag{96.1}
\]
Consequently, on the slice
\[
 d^*a=0,\qquad a\perp{\cal H}^1,                      \tag{96.2}
\]
one has
\[
 \|a\|_{H^1}\leq C_H\|da\|_2.                         \tag{96.3}
\]
\end{proposition}

\begin{proof}
The first-order operator \(d+d^*\), augmented by the projection
onto its finite-dimensional kernel, is elliptic and injective.
Its graph norm is therefore equivalent to the \(H^1\) norm.  This
gives (96.1).  Conditions (96.2) remove the last two terms.
\end{proof}

Let \(C_S\) be the Sobolev constant for
\[
 \|a\|_4\leq C_S\|a\|_{H^1};                          \tag{96.4}
\]
this embedding holds for \(d\leq4\).  Choose \(c_{\mathfrak g}\)
so that pointwise
\[
 |[a\wedge a]|\leq c_{\mathfrak g}|a|^2.              \tag{96.5}
\]

\section{Nonlinear curvature control near the trivial connection}

Use the convention
\[
 F_a=da+g[a\wedge a],\qquad g\geq0.                   \tag{96.6}
\]

\begin{theorem}[Local Coulomb curvature coercivity]
\label{thm:v15v96-local}
Assume (96.2) and
\[
 \|a\|_{H^1}\leq\delta,\qquad
 C_Hg c_{\mathfrak g}C_S^2\delta\leq{1\over2}.        \tag{96.7}
\]
Then
\[
 \|a\|_{H^1}\leq2C_H\|F_a\|_2.                       \tag{96.8}
\]
In particular,
\[
 \|F_a\|_2^2\geq{1\over4C_H^2}\|a\|_{H^1}^2          \tag{96.9}
\]
throughout this slice ball.
\end{theorem}

\begin{proof}
From (96.3), (96.6), Hölder, (96.4), and (96.5),
\[
\begin{split}
 \|a\|_{H^1}
 &\leq C_H\|da\|_2\\
 &\leq C_H\|F_a\|_2+
 C_Hg c_{\mathfrak g}\|a\|_4^2\\
 &\leq C_H\|F_a\|_2+
 C_Hg c_{\mathfrak g}C_S^2\|a\|_{H^1}^2.             \tag{96.10}
\end{split}
\]
Use \(\|a\|_{H^1}^2\leq\delta\|a\|_{H^1}\) and (96.7) to absorb
the last term.  Squaring (96.8) gives (96.9).
\end{proof}

The smallness in (96.7) is a chart radius, not a mass gap.  The
theorem assumes that the gauge orbit has already been put in the
slice (96.2).

\section{Slices about nontrivial flat connections}

Let \(A_0\) be flat and let \(d_{A_0}\) be its covariant
differential.  Put
\[
 \Delta_{A_0}^{(1)}
 =d_{A_0}d_{A_0}^*+d_{A_0}^*d_{A_0}.                 \tag{96.11}
\]
Let \(\Pi_{A_0}\) project onto its kernel.

\begin{proposition}[Flat-background local coercivity]
\label{prop:v15v96-flat-local}
There are constants \(C(A_0)\) and \(\delta(A_0)>0\) such that, if
\[
 d_{A_0}^*a=0,\qquad \Pi_{A_0}a=0,\qquad
 \|a\|_{H^1}\leq\delta(A_0),                           \tag{96.12}
\]
then
\[
 \|a\|_{H^1}
 \leq2C(A_0)\|F_{A_0+a}\|_2.                          \tag{96.13}
\]
\end{proposition}

\begin{proof}
Ellipticity of \(d_{A_0}+d_{A_0}^*\), restricted to the orthogonal
complement of its kernel, gives
\[
 \|a\|_{H^1}\leq C(A_0)\|d_{A_0}a\|_2                \tag{96.14}
\]
under (96.12).  Flatness gives
\[
 F_{A_0+a}=d_{A_0}a+g[a\wedge a].                    \tag{96.15}
\]
Repeat (96.10) and choose
\(C(A_0)g c_{\mathfrak g}C_S^2\delta(A_0)\leq1/2\).
\end{proof}

\section{Compact holonomy moduli and the kernel-rank gate}

If \(\pi_1(\Sigma)\) is finitely presented and \(G\) is compact,
the representation variety
\[
 {\rm Hom}(\pi_1(\Sigma),G)                           \tag{96.16}
\]
is a closed subset of a finite Cartesian power of \(G\), hence
compact.  Its quotient by conjugation is compact.  This describes
flat holonomies at the topological level.  It does not automatically
make \(C(A_0)\) uniform.

\begin{theorem}[Uniform transverse gap under constant kernel rank]
\label{thm:v15v96-uniform-gap}
Let \(X\) be compact and suppose
\(\Delta_x\geq0\), \(x\in X\), is a norm-resolvent continuous
family of self-adjoint elliptic operators with compact resolvent on
one common carrier.  If
\[
 \dim\ker\Delta_x=m\quad\hbox{for every }x\in X,       \tag{96.17}
\]
then there is \(\gamma>0\) such that
\[
 \sigma(\Delta_x)\cap(0,\gamma)=\varnothing
 \qquad(x\in X).                                      \tag{96.18}
\]
The corresponding Hodge constants on the kernel complements are
uniform.
\end{theorem}

\begin{proof}
Norm-resolvent continuity and compact resolvent imply continuity of
each ordered eigenvalue, with multiplicity.  Under (96.17),
\[
 \lambda_1(x)=\cdots=\lambda_m(x)=0,\qquad
 \lambda_{m+1}(x)>0.                                  \tag{96.19}
\]
The continuous function \(\lambda_{m+1}\) has a positive minimum
on compact \(X\).  Take this minimum as \(\gamma\).  The elliptic
estimate on the spectral complement then has a uniform
\(L^2\)-coercive term, while the principal ellipticity constants
are uniform on the compact family.
\end{proof}

\begin{counterexample}[Kernel jumps destroy the uniform gap]
\label{cnt:v15v96-kernel-jump}
For \(s\in[0,1]\), let
\[
 \Delta_s=
 \begin{pmatrix}s^2&0\\0&1\end{pmatrix}.              \tag{96.20}
\]
This is a norm-continuous nonnegative family on a compact parameter
space.  For \(s>0\) its smallest positive eigenvalue is \(s^2\);
at \(s=0\) the kernel has dimension one and the remaining positive
eigenvalue is one.  After projecting away the exact kernel at every
\(s\), the infimum of the positive spectrum over the family is
still zero.
\end{counterexample}

\begin{proof}
The eigenvalues are read directly from (96.20).
\end{proof}

Reducible flat connections are precisely where stabilizers and
cohomology dimensions can change.  Compactness of the holonomy
quotient therefore does not replace the constant-rank hypothesis
in Theorem~\ref{thm:v15v96-uniform-gap}.

\section{A conditional global small-curvature compiler}

\begin{theorem}[From slice coverage to quotient coercivity]
\label{thm:v15v96-global-conditional}
Let \({\cal M}_{\rm flat}\) be a compact family of flat connections.
Assume:
\begin{enumerate}
\item every connection \(A\) with
      \(\|F_A\|_2\leq\varepsilon_0\) is gauge equivalent to
      \(A_0+a\), where \(A_0\in{\cal M}_{\rm flat}\) and \(a\)
      satisfies (96.12) with one radius \(\delta\);
\item the Hodge constants \(C(A_0)\) are bounded by one \(C_*\);
\item \(C_*g c_{\mathfrak g}C_S^2\delta\leq1/2\).
\end{enumerate}
Then
\[
 \operatorname{dist}_{H^1}([A],{\cal M}_{\rm flat})
 \leq2C_*\|F_A\|_2                                    \tag{96.21}
\]
for all such \(A\).
\end{theorem}

\begin{proof}
Use the gauge and flat representative supplied by the first
hypothesis.  Proposition~\ref{prop:v15v96-flat-local}, with the
uniform constants in the remaining hypotheses, gives
\(\|a\|_{H^1}\leq2C_*\|F_A\|_2\).  Taking the quotient distance
proves (96.21).
\end{proof}

This theorem isolates the unproved global input instead of calling
the local estimate a global gauge theorem.  Uhlenbeck-type gauge
compactness is the relevant source of the first row, while
Theorem~\ref{thm:v15v96-uniform-gap} explains the constant-rank
requirement in the second.

\section{Handoff to the boundary interaction programme}

Inside a constant-rank slice chart, (96.9) or (96.13) makes the
positive curvature square a quadratic reservoir for transverse
connection fluctuations.  Flat holonomy coordinates remain
finite-dimensional quotient variables.  They can be integrated
separately only after proving:
\begin{enumerate}
\item a finite slice atlas with controlled overlap maps;
\item treatment of reducible stabilizers and Faddeev--Popov zero
      modes;
\item a measure and determinant density on each holonomy stratum;
\item compatibility with the common boundary carrier of V86--V89.
\end{enumerate}

\begin{assessment}[Physical status]
\label{ass:v15v96-status}
V96 proves local nonlinear curvature coercivity in Coulomb gauge,
extends it to flat backgrounds, and proves a uniform transverse
gap theorem over compact constant-kernel families.  It also proves
that kernel jumps invalidate the naive compactness argument.

The existence of a global uniform slice atlas, control across
reducible strata, and the associated ghost determinant remain
unproved.
\end{assessment}

\begin{decision}[V97 stratified slice and determinant ledger]
\label{dec:v15v96-next}
Construct the finite-dimensional change-of-variables theorem inside
one constant-rank slice chart: orbit directions, Faddeev--Popov
operator, stabilizer kernel, reduced determinant, and overlap
Jacobian.  Prove what is invariant under a change of slice and show
how an eigenvalue approaching zero destroys uniform determinant
control near a reducible stratum.
\end{decision}

\begin{firewall}[Claim boundary after V96]
\label{fire:v15v96}
V96 assumes a Coulomb representative and, in its global compiler,
uniform slice coverage.  It does not prove a global gauge fixing,
a Yang--Mills mass gap, reflection positivity of ghosts, or the
full Standard Model--Einstein continuum.
\end{firewall}

\section*{Version V96 append-only record}
\addcontentsline{toc}{section}{Version V96 append-only record}
The complete V95 source was retained before this continuation.  It
had \(940251\) bytes, \(27769\) lines, and SHA--256
\[
 \texttt{D08BEB6A95472C95575D759F4140B5C6F3C31D1465581A0B2948C0085495E712}.
\]
V96 proved local Coulomb curvature coercivity, its flat-background
extension, a uniform constant-rank Hodge-gap theorem, and the
kernel-jump obstruction at reducible holonomy strata.
\chapter{Fifteenth-cycle continuation V97: reduced gauge determinants across slice strata}
\label{sec:v15v97}

V96 reduced small-curvature gauge coercivity to constant-kernel
slice charts.  The same kernel controls the finite-cutoff
Faddeev--Popov determinant.  This note develops that statement in a
finite-dimensional model, where every Jacobian is an ordinary
determinant and the rank-jump obstruction can be proved exactly.

\section{Orbit map and reduced Gram determinant}

Let a compact Lie group \(G\) act orthogonally on a Euclidean
configuration space \(X\).  Give its Lie algebra \(\mathfrak g\) an
invariant inner product.  At \(x\in X\), define the infinitesimal
orbit map
\[
 R_x:\mathfrak g\longrightarrow X,\qquad
 R_x\xi={d\over dt}\bigg|_{t=0}\exp(t\xi)x,            \tag{97.1}
\]
and the positive orbit Gram operator
\[
 M_x=R_x^*R_x\geq0.                                   \tag{97.2}
\]
Its kernel is the stabilizer algebra
\[
 \mathfrak h_x=\ker R_x=\ker M_x.                     \tag{97.3}
\]

\begin{proposition}[Reduced orbit Jacobian]
\label{prop:v15v97-orbit}
On \(\mathfrak h_x^\perp\), \(M_x\) is positive definite and
\[
 J_{\rm orb}(x)
 =\sqrt{\det{}'M_x}
 =\prod_{j=1}^{r_x}s_j(R_x),                           \tag{97.4}
\]
where \(r_x=\operatorname{rank}R_x\), the prime omits the stabilizer
kernel, and \(s_j\) are the nonzero singular values.  This is the
Euclidean volume expansion of the orbit map.
\end{proposition}

\begin{proof}
Choose orthonormal bases adapted to
\(\mathfrak h_x\oplus\mathfrak h_x^\perp\) and
\(\ker R_x^*\oplus\operatorname{ran}R_x\).  Singular-value
decomposition makes the restriction of \(R_x\) diagonal with
entries \(s_1,\ldots,s_{r_x}>0\).  Its Gram matrix has eigenvalues
\(s_j^2\), so its volume expansion is their product, which equals
\(\sqrt{\det{}'M_x}\).
\end{proof}

\section{One orthogonal slice chart}

Let \(S\subset X\) be a submanifold through \(x\) satisfying
\[
 T_xX=T_x(Gx)\oplus T_xS
 \quad\hbox{orthogonally}.                             \tag{97.5}
\]
After quotienting the stabilizer, suppose
\[
 \Phi:S\times(G/G_x)\longrightarrow X,\qquad
 \Phi(s,[g])=gs                                      \tag{97.6}
\]
is one-to-one on a neighborhood.

\begin{theorem}[Local quotient integration]
\label{thm:v15v97-coarea}
For every nonnegative \(G\)-invariant measurable function \(F\)
supported in the chart image,
\[
 \int_XF(y)\,dy
 =\int_S F(s)\,J_{\rm orb}(s)\,
       {\rm vol}(G/G_s)\,d\sigma(s),                  \tag{97.7}
\]
with Haar and slice measures normalized consistently.  Equivalently,
after division by orbit volume, the local quotient density is
\[
 d\mu_{\rm quot}(s)=J_{\rm orb}(s)\,d\sigma(s).        \tag{97.8}
\]
\end{theorem}

\begin{proof}
Apply the ordinary change-of-variables formula to (97.6).  In the
orthogonal decomposition (97.5), its differential has a slice
block of determinant one relative to \(d\sigma\), and an orbit
block with volume expansion (97.4).  Integration over the orbit
gives its Haar volume.  Gauge invariance removes the orbit variable
from \(F\), giving (97.7)--(97.8).
\end{proof}

For the orthogonal gauge condition, the linearized gauge map is
\[
 C_x=R_x^*,\qquad C_xR_x=M_x.                         \tag{97.9}
\]
Thus the reduced Faddeev--Popov determinant is
\(\det{}'M_x\).  The square root in (97.8) and the full determinant
in a delta-function gauge formula are consistent: the delta
distribution for the normal coordinate contributes the reciprocal
square root.

\section{Independence of the slice coordinate}

\begin{proposition}[Overlap invariance]
\label{prop:v15v97-overlap}
Let \(S_\alpha,S_\beta\) be two slice charts meeting the same regular
orbit stratum, and let
\[
 \tau_{\beta\alpha}:S_\alpha/G_\alpha
 \longrightarrow S_\beta/G_\beta                    \tag{97.10}
\]
be the induced quotient-coordinate change.  Their quotient
densities satisfy
\[
 d\mu_{{\rm quot},\alpha}
 =\tau_{\beta\alpha}^*
   d\mu_{{\rm quot},\beta}.                           \tag{97.11}
\]
On triple overlaps, the transition Jacobians obey the usual
cocycle identity.
\end{proposition}

\begin{proof}
Both sides are obtained by disintegrating the same Euclidean
measure on \(X\) along the same \(G\)-orbits, as in
Theorem~\ref{thm:v15v97-coarea}.  The uniqueness of the ordinary
change-of-variables density gives (97.11).  Applying the chain rule
to
\(\tau_{\gamma\alpha}
=\tau_{\gamma\beta}\circ\tau_{\beta\alpha}\)
gives the determinant cocycle.
\end{proof}

This proves coordinate invariance on one regular stratum.  It does
not identify determinant phases for chiral fermions; that separate
line-bundle problem was treated conditionally in V81--V84.

\section{Uniform determinant control on a constant-rank stratum}

\begin{theorem}[Reduced determinant stability]
\label{thm:v15v97-det-stability}
Let \(M_n,M\) be positive operators on an \(r\)-dimensional
Euclidean space with one common kernel \(H\) of dimension \(m\).
Assume
\[
 M_n|_{H^\perp}\geq\gamma I,\qquad
 M|_{H^\perp}\geq\gamma I,\qquad \gamma>0,             \tag{97.12}
\]
and \(\|M_n-M\|_{\rm op}\to0\).  Then
\[
 \det{}'M_n\longrightarrow\det{}'M,                   \tag{97.13}
\]
and
\[
 |\log\det{}'M_n-\log\det{}'M|
 \leq {r-m\over\gamma}\|M_n-M\|_{\rm op}.              \tag{97.14}
\]
\end{theorem}

\begin{proof}
Restrict to \(H^\perp\) and join \(M\) to \(M_n\) by the line
segment \(M_s\).  The lower bound \(\gamma I\) persists.  The
finite-dimensional derivative formula gives
\[
 {d\over ds}\log\det M_s
 =\operatorname{Tr}\bigl(M_s^{-1}(M_n-M)\bigr).        \tag{97.15}
\]
Its absolute value is at most
\((r-m)\gamma^{-1}\|M_n-M\|_{\rm op}\).  Integration proves
(97.14), hence (97.13).
\end{proof}

If the kernels vary continuously with constant rank, the Kato
unitaries from V33 and V86 transport them to one fixed kernel.
Theorem~\ref{thm:v15v97-det-stability} then applies on the common
complement.

\begin{counterexample}[Reduced determinant discontinuity at a
stabilizer jump]
\label{cnt:v15v97-jump}
For
\[
 M_\epsilon=
 \begin{pmatrix}\epsilon&0\\0&1\end{pmatrix},
 \qquad \epsilon\geq0,                                \tag{97.16}
\]
one has
\[
 \det{}'M_\epsilon=
 \begin{cases}
 \epsilon,&\epsilon>0,\\
 1,&\epsilon=0.
 \end{cases}                                          \tag{97.17}
\]
Thus the reduced determinant is discontinuous when the kernel rank
jumps.  For \(\epsilon>0\),
\[
 {d\over d\epsilon}\log\det{}'M_\epsilon={1\over\epsilon},          \tag{97.18}
\]
so logarithmic stresses also diverge.
\end{counterexample}

\begin{proof}
The prime removes only exact zero eigenvalues.  The eigenvalues of
(97.16) give (97.17)--(97.18) directly.
\end{proof}

\section{Faddeev--Popov signs and the Gribov boundary}

The positivity above used the orthogonal choice \(C_x=R_x^*\).  For
a general gauge condition \(C_x\), the Faddeev--Popov operator
\[
 {\cal M}_x=C_xR_x                                    \tag{97.19}
\]
need not be positive or even self-adjoint.  A sign or phase choice
for \(\det{}'{\cal M}_x\) is then additional data.

\begin{proposition}[Uniform logarithmic control requires a singular
value floor]
\label{prop:v15v97-singular}
Let \({\cal M}_x\) be an invertible \(q\times q\) matrix family.
If
\[
 s_{\min}({\cal M}_x)\geq\gamma>0                     \tag{97.20}
\]
and \({\cal M}_x\) is differentiable, then
\[
 |d\log\det{\cal M}_x|
 \leq q\gamma^{-1}\|d{\cal M}_x\|_{\rm op}.            \tag{97.21}
\]
When the least singular value tends to zero, no such uniform bound
follows.
\end{proposition}

\begin{proof}
Use
\(d\log\det{\cal M}=\operatorname{Tr}({\cal M}^{-1}d{\cal M})\)
and \(\|{\cal M}^{-1}\|_{\rm op}\leq\gamma^{-1}\).
Counterexample~\ref{cnt:v15v97-jump} supplies the sharp failure.
\end{proof}

The zero set of the determinant is the local Gribov or stabilizer
boundary.  Representing the determinant by ghost variables does
not remove (97.20); it rewrites the same Jacobian.  Reflection
positivity of that representation is a separate requirement.

\section{Stratified quotient ledger}

On each constant-stabilizer stratum, V97 provides:
\begin{enumerate}
\item a reduced orbit Jacobian and local quotient density;
\item invariance under changes of slice coordinates;
\item determinant and logarithmic-derivative control from a
      uniform transverse singular-value floor;
\item common-complement transport when kernel projections converge.
\end{enumerate}
Across strata, Counterexample~\ref{cnt:v15v97-jump} forbids treating
the prime determinant as a continuous scalar.  The correct object
must include stabilizer volume, normal coordinates to the stratum,
and compatible measure disintegration.  A single regular-stratum
ghost determinant is incomplete there.

\begin{assessment}[Physical status]
\label{ass:v15v97-status}
V97 proves the finite-cutoff quotient Jacobian and its invariance on
regular slice overlaps, together with sharp determinant control on
constant-rank strata.  It proves discontinuity and divergent
logarithmic response at a stabilizer-rank jump.

No continuum stratified measure, Gribov-global slice, or
reflection-positive ghost construction has been proved.
\end{assessment}

\begin{decision}[V98 holonomy collective-coordinate heat]
\label{dec:v15v97-next}
Treat the compact flat-holonomy quotient as a genuine collective
coordinate with an elliptic kinetic operator.  Prove heat
nuclearity on a compact regular quotient or orbifold, show why a
mere direct integral over a continuous holonomy parameter is not
trace class, and state the extra interface conditions needed at
singular strata.
\end{decision}

\begin{firewall}[Claim boundary after V97]
\label{fire:v15v97}
V97 is a finite-dimensional slice theorem.  It does not prove a
global gauge fixing, continuity across reducible strata, a
Yang--Mills mass gap, or the full Standard Model--Einstein
continuum.
\end{firewall}

\section*{Version V97 append-only record}
\addcontentsline{toc}{section}{Version V97 append-only record}
The complete V96 source was retained before this continuation.  It
had \(950339\) bytes, \(28066\) lines, and SHA--256
\[
 \texttt{73B5469D99FB7D6112C1AF8EFD42BCA533A2DD9BC31CBBDDF10B065A2011AEEA}.
\]
V97 constructed the reduced orbit and Faddeev--Popov determinant
ledger on regular slice strata, proved overlap invariance and
uniform determinant control, and isolated the exact rank-jump
failure at reducible configurations.
\chapter{Fifteenth-cycle continuation V98: holonomy collective-coordinate heat}
\label{sec:v15v98}

V97 produced a quotient density on each regular gauge stratum.
This density alone does not define a nuclear transfer operator over
a continuous holonomy parameter.  The missing ingredient is kinetic
propagation in the collective coordinate.  This note proves the
compact regular case and isolates the singular-stratum interface
problem.

\section{Weighted quotient Laplacian}

Let \(Q\) be a compact connected \(r\)-dimensional Riemannian
manifold without boundary, representing one regular flat-holonomy
quotient.  Let \(w\) be a quotient density satisfying
\[
 0<w_-\leq w(q)\leq w_+<\infty,\qquad
 w\in W^{1,\infty}(Q).                                \tag{98.1}
\]
On \(L^2(Q,w\,d{\rm vol})\), define the form
\[
 \ell[u]
 =\int_Q\bigl(\langle A(q)\nabla u,\nabla u\rangle
                  +V(q)|u|^2\bigr)w(q)\,d{\rm vol},  \tag{98.2}
\]
where
\[
 a_-|\xi|^2\leq\langle A(q)\xi,\xi\rangle
 \leq a_+|\xi|^2,\qquad V\geq-b.                      \tag{98.3}
\]
Let \(L_Q\) be the associated self-adjoint operator.

\begin{theorem}[Regular holonomy heat nuclearity]
\label{thm:v15v98-holonomy}
Under (98.1)--(98.3), \(L_Q\) is bounded below, has compact
resolvent, and satisfies
\[
 N_{L_Q}(\Lambda)
 \leq C(1+\Lambda+b)^{r/2}                            \tag{98.4}
\]
for \(\Lambda\geq-b\).  Consequently
\[
 e^{-tL_Q}\in{\cal L}^1,\qquad
 \operatorname{Tr}e^{-tL_Q}<\infty                   \tag{98.5}
\]
for every \(t>0\).
\end{theorem}

\begin{proof}
The bounds on \(w\) and \(A\) make the form norm in (98.2), after
adding \(b+1\), equivalent to the \(H^1(Q)\) norm.  Rellich
compactness gives compact resolvent.  The min--max principle compares
\(L_Q+b\) above and below with fixed positive multiples of the
scalar Laplacian plus one.  The standard Laplace counting estimate
on a compact \(r\)-manifold gives (98.4).  Summation by parts, as in
V88, turns (98.4) into an exponentially weighted summable
eigenvalue sequence, proving (98.5).
\end{proof}

The same conclusion holds on a compact orbifold with finite local
stabilizers: a finite orbifold atlas and invariant partition of
unity reduce the form and counting estimates to finite-group
quotients of ordinary elliptic charts.

\section{Why a bare holonomy direct integral fails}

\begin{theorem}[Continuous multiplicity obstruction]
\label{thm:v15v98-direct-integral}
Let \((Q,\mu)\) be non-atomic with positive finite measure and let
\(0\ne T\) be a compact operator on a separable Hilbert space
\({\cal K}\).  The constant decomposable operator
\[
 {\cal T}=I_{L^2(Q,\mu)}\otimes T                     \tag{98.6}
\]
on \(L^2(Q,\mu)\otimes{\cal K}\) is not compact and hence is not
trace class.
\end{theorem}

\begin{proof}
Choose \(v\in{\cal K}\) with \(Tv\ne0\) and an orthonormal sequence
\((f_n)\) in \(L^2(Q,\mu)\).  Then \(f_n\otimes v\) is bounded and
\[
 {\cal T}(f_n\otimes v)=f_n\otimes Tv                 \tag{98.7}
\]
is an orthogonal sequence of one fixed nonzero norm.  It has no
convergent subsequence, so \({\cal T}\) is not compact.
\end{proof}

\begin{corollary}[Fiber traces cannot simply be integrated]
\label{cor:v15v98-fiber}
Even if \(T=e^{-tH_{\rm fib}}\) is trace class,
\[
 \int_Q\operatorname{Tr}T\,d\mu(q)<\infty             \tag{98.8}
\]
does not make the decomposable transfer operator trace class.
Equation (98.8) is a scalar integral, not its Hilbert-space trace.
\end{corollary}

\begin{proof}
Theorem~\ref{thm:v15v98-direct-integral} applies because a heat
operator is nonzero.
\end{proof}

Thus a continuous flat-holonomy label cannot be treated as an
uncounted degeneracy.  The kinetic operator \(L_Q\) converts that
label into a discrete collective-coordinate spectrum.

\section{Product with the transverse Fock sector}

Let \(H_{\rm tr}\) be a transverse, gauge-fixed Fock Hamiltonian
whose heat factor is trace class.  On
\[
 {\cal H}_{\rm coll}=L^2(Q,w\,d{\rm vol})\otimes{\cal F}_{\rm tr},
                                                               \tag{98.9}
\]
put
\[
 H_{\rm coll}=L_Q\otimes I+I\otimes H_{\rm tr}.        \tag{98.10}
\]

\begin{proposition}[Collective-coordinate product trace]
\label{prop:v15v98-product}
For every \(t>0\),
\[
 e^{-tH_{\rm coll}}
 =e^{-tL_Q}\otimes e^{-tH_{\rm tr}},                  \tag{98.11}
\]
\[
 \operatorname{Tr}e^{-tH_{\rm coll}}
 =\operatorname{Tr}e^{-tL_Q}\,
  \operatorname{Tr}e^{-tH_{\rm tr}}<\infty.           \tag{98.12}
\]
\end{proposition}

\begin{proof}
The two self-adjoint summands strongly commute, so the spectral
calculus gives (98.11).  The tensor product of two trace-class
operators is trace class and its trace is the product of traces.
\end{proof}

\begin{proposition}[Trace-norm stability of the product]
\label{prop:v15v98-product-continuity}
If \(A_n\to A\) and \(B_n\to B\) in trace norm, then
\[
 A_n\otimes B_n\longrightarrow A\otimes B
 \quad\hbox{in trace norm}.                            \tag{98.13}
\]
More precisely,
\[
 \|A_n\otimes B_n-A\otimes B\|_1
 \leq\|A_n-A\|_1\|B_n\|_1
      +\|A\|_1\|B_n-B\|_1.                            \tag{98.14}
\]
\end{proposition}

\begin{proof}
Add and subtract \(A\otimes B_n\), use the triangle inequality and
\(\|C\otimes D\|_1=\|C\|_1\|D\|_1\).
\end{proof}

V88--V94 supply conditional trace-norm control for the transverse
factor.  Theorem~\ref{thm:v15v98-holonomy} supplies nuclearity of
the collective factor.  If the metrics, weights, and potentials in
(98.2) converge in a common uniformly elliptic form norm, the V88
form-resolvent and heat-tail argument gives trace-norm convergence
of \(e^{-tL_{Q,n}}\); Proposition~\ref{prop:v15v98-product-continuity}
then joins the two limits.

\section{Rank jumps produce a stratified interface problem}

At a reducible flat connection, the stabilizer changes and the
weight derived in V97 may vanish or change dimension.  Assumption
(98.1) then fails.  The regular quotient may acquire a boundary,
cone point, or a higher-codimension singular stratum.

\begin{proposition}[Interface data are operator data]
\label{prop:v15v98-interface}
Let \(Q_{\rm reg}\) be a regular stratum whose metric completion has
a singular boundary.  The differential expression in (98.2) does
not by itself determine a unique heat semigroup unless its minimal
operator is essentially self-adjoint.  Otherwise a self-adjoint
extension, equivalently admissible interface or boundary data, is
part of the definition of \(L_Q\).
\end{proposition}

\begin{proof}
A closed semibounded quadratic form determines one self-adjoint
operator, while distinct closed extensions of the same
compactly-supported differential expression give distinct
self-adjoint operators and heat semigroups.  If the minimal
operator is not essentially self-adjoint, at least two such
extensions exist by definition.  Hence the expression alone cannot
select the transfer operator.
\end{proof}

The physically admissible extension must agree with measure
disintegration across strata, gauge invariance, and reflection
sewing.  Choosing the Friedrichs extension only because it exists
would add a boundary condition not derived from the cutoff theory.

\section{Conditional holonomy handoff}

\begin{theorem}[Regular-stratum handoff]
\label{thm:v15v98-handoff}
Suppose a finite family of compact regular quotient components
\(Q_\alpha\) covers the physical flat-holonomy sector up to
measure-zero finite-orbifold identifications, each satisfies
(98.1)--(98.3), and their overlap conditions define one
self-adjoint elliptic operator \(L_{\rm hol}\).  Then
\(e^{-tL_{\rm hol}}\) is trace class.  Tensoring it with any
transverse Fock heat factor satisfying V91 preserves nuclearity and
trace-norm convergence.
\end{theorem}

\begin{proof}
A finite elliptic atlas gives the global compact-resolvent and
counting argument of Theorem~\ref{thm:v15v98-holonomy}.  Apply
Propositions~\ref{prop:v15v98-product} and
\ref{prop:v15v98-product-continuity}.
\end{proof}

\begin{assessment}[Physical status]
\label{ass:v15v98-status}
V98 proves nuclear collective-coordinate propagation on compact
regular holonomy quotients and proves that a bare continuous direct
integral is not trace class.  It also identifies self-adjoint
interface data as indispensable at singular strata.

No stratified holonomy operator, cutoff-derived interface condition,
or reflection-positive ghost/collective measure has been
constructed.
\end{assessment}

\begin{decision}[V99 conditional sector assembly]
\label{dec:v15v98-next}
Assemble the proved compilers into one conditional boundary transfer
theorem: finite tensor products of collective, transverse bosonic,
fermionic, Higgs, and anomaly-line factors, followed by a
relative-form interaction.  State every uniform hypothesis in one
ledger and prove trace-norm convergence.  Keep gravity as a
separate required factor with an explicit conformal-stability row.
\end{decision}

\begin{firewall}[Claim boundary after V98]
\label{fire:v15v98}
V98 treats compact regular quotient geometry conditionally.  It
does not resolve reducible strata, choose physical interface data,
prove a gauge mass gap, or construct the full Standard
Model--Einstein continuum.
\end{firewall}

\section*{Version V98 append-only record}
\addcontentsline{toc}{section}{Version V98 append-only record}
The complete V97 source was retained before this continuation.  It
had \(960740\) bytes, \(28367\) lines, and SHA--256
\[
 \texttt{AE997B239235176035BC1FF846BBE1D3A45BDEB29B387FFEDC9C354D7A7F5F77}.
\]
V98 proved heat nuclearity for elliptic holonomy collective
coordinates, the continuous-multiplicity obstruction for bare
direct integrals, tensor-product stability, and the necessity of
singular-stratum interface data.
\chapter{Fifteenth-cycle continuation V99: conditional Standard Model--Einstein sector assembly}
\label{sec:v15v99}

V86--V98 now provide separate compilers for common carriers,
one-particle heat smoothing, Fock lifts, stable Higgs interactions,
local gauge quotients, and holonomy collective coordinates.  This
note proves that these pieces assemble once their remaining uniform
hypotheses are supplied.  The result is conditional, but the
conditions are now stated at the exact operator level needed for a
continuum boundary transfer.

\section{Finite tensor products of nuclear sectors}

\begin{theorem}[Finite-sector tensor convergence]
\label{thm:v15v99-tensor}
Let \(m<\infty\).  For \(1\leq j\leq m\), suppose
\[
 A_{j,n}\longrightarrow A_j\quad\hbox{in trace norm}. \tag{99.1}
\]
Then
\[
 {\cal A}_n=\bigotimes_{j=1}^m A_{j,n}
 \longrightarrow
 {\cal A}=\bigotimes_{j=1}^m A_j
 \quad\hbox{in trace norm},                            \tag{99.2}
\]
and
\[
 \operatorname{Tr}{\cal A}_n
 =\prod_{j=1}^m\operatorname{Tr}A_{j,n}
 \longrightarrow
 \prod_{j=1}^m\operatorname{Tr}A_j.                   \tag{99.3}
\]
\end{theorem}

\begin{proof}
Insert the factors one at a time:
\[
 {\cal A}_n-{\cal A}
 =\sum_{k=1}^m
 \left(\bigotimes_{j<k}A_j\right)
 \otimes(A_{k,n}-A_k)
 \otimes\left(\bigotimes_{j>k}A_{j,n}\right).          \tag{99.4}
\]
The trace norm of a tensor product is the product of trace norms.
Every sequence \(\|A_{j,n}\|_1\) is bounded by (99.1), so the sum
of the resulting bounds tends to zero.  The trace factorization
and its limit give (99.3).
\end{proof}

For self-adjoint sector Hamiltonians \(H_{j,n}\) on common carriers,
put
\[
 H_{0,n}
 =\sum_{j=1}^m
 I\otimes\cdots\otimes H_{j,n}\otimes\cdots\otimes I. \tag{99.5}
\]
The summands strongly commute, and therefore
\[
 e^{-tH_{0,n}}
 =\bigotimes_{j=1}^m e^{-tH_{j,n}}.                   \tag{99.6}
\]

\begin{corollary}[Free assembled heat]
\label{cor:v15v99-free}
If
\[
 e^{-tH_{j,n}}\longrightarrow e^{-tH_j}
 \quad\hbox{in trace norm for every }t>0              \tag{99.7}
\]
in every sector, then the assembled free heat factors (99.6)
converge in trace norm for every \(t>0\).
\end{corollary}

\begin{proof}
Apply Theorem~\ref{thm:v15v99-tensor} at each \(t\).
\end{proof}

\section{Anomaly-line scalar and positivity}

After determinant-line transport, a residual finite-dimensional
factor may act by a scalar \(z_n\).

\begin{proposition}[Scalar-line handoff]
\label{prop:v15v99-line}
If \(z_n\to z\) and \({\cal A}_n\to{\cal A}\) in trace norm, then
\[
 z_n{\cal A}_n\longrightarrow z{\cal A}
 \quad\hbox{in trace norm}.                            \tag{99.8}
\]
If \({\cal A}\) is a nonzero positive operator, then
\(z{\cal A}\) is positive exactly when \(z\) is real and
nonnegative.
\end{proposition}

\begin{proof}
Use
\[
 \|z_n{\cal A}_n-z{\cal A}\|_1
 \leq|z_n-z|\,\|{\cal A}_n\|_1
      +|z|\,\|{\cal A}_n-{\cal A}\|_1.                \tag{99.9}
\]
For the last assertion, evaluate the quadratic form on a vector
where \(\langle v,{\cal A}v\rangle>0\).
\end{proof}

\begin{counterexample}[Uncancelled phase oscillation]
\label{cnt:v15v99-phase}
Let \({\cal A}\ne0\) be trace class and \(z_n=(-1)^n\).  Every
\(\|z_n{\cal A}\|_1\) is the same, but the amplitudes do not
converge.  Nuclear magnitude control therefore does not replace
anomaly-line trivialization and phase convergence.
\end{counterexample}

\section{Interacting assembled transfer theorem}

\begin{theorem}[Conditional assembled boundary transfer]
\label{thm:v15v99-assembled}
Assume (99.7) for a finite family of sector Hamiltonians and let
\(q_{0,n}\) be the form of (99.5).  Let \(v_n\) be a symmetric
interaction form satisfying, with cutoff-independent constants,
\[
 |v_n[\Psi]|
 \leq\alpha q_{0,n}[\Psi]+\beta\|\Psi\|^2,
 \qquad 0\leq\alpha<1,\quad\beta<\infty.               \tag{99.10}
\]
Let \(H_n=H_{0,n}\dotplus V_n\), and suppose
\[
 H_n\longrightarrow H
 \quad\hbox{in norm resolvent}.                        \tag{99.11}
\]
Then, for every \(t>0\),
\[
 e^{-tH_n}\longrightarrow e^{-tH}
 \quad\hbox{in trace norm},                            \tag{99.12}
\]
\[
 Z_n(t)\longrightarrow Z(t)>0,\qquad
 {e^{-tH_n}\over Z_n(t)}
 \longrightarrow {e^{-tH}\over Z(t)}
 \quad\hbox{in trace norm}.                            \tag{99.13}
\]
\end{theorem}

\begin{proof}
Corollary~\ref{cor:v15v99-free} gives the free heat hypothesis
(91.4).  Apply Theorem~\ref{thm:v15v91-gibbs} to obtain (99.12),
then Corollary~\ref{cor:v15v91-states} to obtain (99.13).
\end{proof}

\section{Exact hypothesis ledger}

Theorem~\ref{thm:v15v99-assembled} applies to the intended boundary
theory only after the following rows are supplied.
\[
\begin{array}{p{0.23\textwidth}|p{0.68\textwidth}}
\hbox{sector}&\hbox{required uniform statement}\\ \hline
\hbox{common carrier}&
isometric transport of cutoff boundary spaces, convergent low
spectral projections, and common heat tails as in V86--V88\\
\hbox{fermions}&
one-particle trace-norm heat convergence, uniformly bounded
zero-mode multiplicity, and the anomaly/bordism trivialization of
V81--V84\\
\hbox{bosons}&
one-particle trace-norm heat convergence and a uniform positive
gap after the physical zero-mode quotient, as in V89\\
\hbox{Higgs}&
trace-norm heat convergence of the stable non-Gaussian reference
in V93, with a fixed renormalization prescription\\
\hbox{gauge transverse}&
uniform slice coverage and curvature coercivity beyond the local
theorem of V96\\
\hbox{holonomy}&
a compact elliptic collective-coordinate operator with
cutoff-derived interface conditions across the V97 strata\\
\hbox{ghost/Jacobian}&
a globally compatible reduced determinant, controlled phase, and
reflection-compatible realization across rank jumps\\
\hbox{gravity}&
a lower-bounded, nuclear, reflection-compatible gravitational heat
factor with the conformal direction treated independently\\
\hbox{interaction}&
the cutoff-independent relative-form estimate (99.10), including
Yukawa and all gauge--Higgs cross terms\\
\hbox{limit}&
full norm-resolvent convergence (99.11), cofinal independence,
sewing, Ward identities, and local projective compatibility.
\end{array}                                            \tag{99.14}
\]

\begin{proposition}[No compensation for a missing sector trace]
\label{prop:v15v99-no-compensation}
If one positive sector factor \(A_{k,n}\) has
\(\operatorname{Tr}A_{k,n}\to\infty\), while every other positive
factor has trace bounded below by a positive constant, then the
free product trace in (99.3) diverges.
\end{proposition}

\begin{proof}
The product formula (99.3) bounds the full trace below by a fixed
positive multiple of \(\operatorname{Tr}A_{k,n}\).
\end{proof}

Thus a divergent bosonic zero mode or an unstabilized gravitational
sector cannot be canceled by multiplying it with a well-behaved
matter trace.  A signed formal cancellation would also forfeit the
positive heat-operator hypothesis.

\section{What the theorem does and does not assemble}

The theorem proves nuclearity, trace-norm convergence, partition
convergence, and bounded-observable convergence once every ledger
row holds.  It also permits interactions through V91 rather than
requiring exact tensor factorization of the full theory.

It does not derive Osterwalder--Schrader reflection positivity from
operator positivity alone.  The Euclidean measure, interface
conditions, gauge quotient, and anomaly phase must be shown to
produce the positive transfer operator in the first place.  It also
does not prove locality or cofinal independence merely from
trace-norm convergence on one exhaustion.

\begin{assessment}[Physical status]
\label{ass:v15v99-status}
V99 proves the finite-sector assembly theorem for the conditional
Standard Model--Einstein boundary transfer and consolidates all
remaining hypotheses in one operator ledger.  The theorem is ready
to accept any future YM full-quotient estimate or gravitational
stabilization without changing its proof.

Several ledger rows remain open, most sharply the bosonic infrared
gap, stratified gauge interface, stable Higgs reference limit,
uniform renormalized interaction form, and gravitational conformal
stability.
\end{assessment}

\begin{decision}[V100 fifteenth-cycle audit]
\label{dec:v15v99-next}
Audit V86--V99 as one boundary nuclearity and interaction block.
Record its proved implication chain, its independent counterexamples,
and the smallest unresolved inputs.  Freeze V100 as the fifteenth
legacy file, then restart a compact V1 with context and a summary
of the major programme results before continuing upstream.
\end{decision}

\begin{firewall}[Claim boundary after V99]
\label{fire:v15v99}
The assembled theorem is conditional on every row of (99.14).  It
does not prove those rows, reflection positivity, the physical mass
gap, or the full Standard Model--Einstein continuum.
\end{firewall}

\section*{Version V99 append-only record}
\addcontentsline{toc}{section}{Version V99 append-only record}
The complete V98 source was retained before this continuation.  It
had \(970477\) bytes, \(28629\) lines, and SHA--256
\[
 \texttt{ECEF99528584613F2D98FBF8716DD8DF7909C92FC7A644AE39AD205162A9162C}.
\]
V99 proved finite-sector tensor and interacting assembly, isolated
anomaly-phase convergence, and placed every remaining
Standard Model--Einstein boundary hypothesis in one exact ledger.
\chapter{Fifteenth-cycle V100 audit: boundary nuclearity, interactions, and quotient heat}
\label{sec:v15v100}

This audit closes the fifteenth cycle.  Its purpose is to state the
logical product of V86--V99 without converting any conditional
compiler into a physical existence theorem.  The cycle began with
the common boundary carrier and ends with a single assembled
transfer theorem.

\section{Proved implication chain}

\begin{theorem}[Fifteenth-cycle boundary transfer chain]
\label{thm:v15v100-chain}
The following implications were proved.
\begin{enumerate}
\item Common finite-rank approximation or one common compact
      smoother turns weak boundary amplitudes into strong limits
      (V86--V87).
\item Uniform elliptic eigenvalue counting plus norm-resolvent
      convergence gives trace-norm convergence of one-particle heat
      smoothers (V88).
\item One-particle trace class lifts to fermionic Fock trace class.
      It lifts to bosonic Fock trace class exactly under a strict
      post-quotient one-particle gap.  Under the corresponding
      uniform hypotheses, the full Fock heat operators converge in
      trace norm (V89).
\item A cutoff-independent relative-form estimate with coefficient
      below one transfers free Fock heat tails to interacting heat
      tails.  Norm-resolvent convergence then gives trace-norm
      convergence of interacting Gibbs operators and normalized
      states (V91).
\item Creation--annihilation estimates verify the form mechanism for
      fixed-cutoff linear, quadratic, and Yukawa vertices, while
      exposing ultraviolet, zero-mode, and polynomial-degree losses
      (V92).
\item A stable Wick-ordered Higgs quartic absorbs lower polynomial
      degrees at finite cutoff.  Vacuum scalar shifts cancel from
      normalized states but remain in partition functions, sewing,
      and gravitational response (V93).
\item Curvature and covariant-derivative cross terms are controlled
      by retaining their positive-square structure.  Yang--Mills
      Cartan directions and the gravitational conformal factor are
      independent obstructions (V94).
\item In a Coulomb slice and off the Hodge kernel, sufficiently
      small connections are controlled in \(H^1\) by curvature.
      This is uniform over compact flat families only with a
      uniform transverse gap; constant kernel rank plus
      norm-resolvent continuity is sufficient (V96).
\item On a regular finite-cutoff gauge stratum, the reduced orbit
      and Faddeev--Popov determinants define a slice-invariant
      quotient density.  Kernel-rank jumps make the prime
      determinant discontinuous and its logarithmic response
      singular (V97).
\item A compact elliptic flat-holonomy collective coordinate has a
      nuclear heat operator.  A bare direct integral over a
      continuous holonomy label is not compact.  Singular strata
      require cutoff-derived self-adjoint interface data (V98).
\item Finite nuclear sectors, followed by a uniformly form-bounded
      interaction, assemble into the conditional trace-norm
      Standard Model--Einstein boundary transfer theorem of V99.
\end{enumerate}
\end{theorem}

\begin{proof}
Each row is the cited theorem or its stated corollary.  The
composition is valid because V99 uses exactly the trace-norm free
heat hypothesis delivered by the earlier sector compilers and then
applies V91.
\end{proof}

\section{Independent failure certificates}

\begin{audit}[Counterexamples retained]
\label{aud:v15v100-counterexamples}
The cycle proved that none of the following substitutions is valid:
\begin{enumerate}
\item uniform Schatten norms without common-carrier alignment;
\item local symbol positivity on a proper momentum cell for a
      global operator gap;
\item one-particle heat-trace control for a bosonic Fock trace when
      the spectral gap closes;
\item a scalar lower bound for heat-trace multiplicity control;
\item pointwise form or strong-resolvent convergence for
      trace-norm heat convergence;
\item fixed-cutoff infinitesimal form boundedness for a uniform
      renormalized lower bound;
\item a stable matter quartic for the gravitational conformal
      direction;
\item the Yang--Mills commutator quartic for Cartan-valued gauge
      configurations;
\item compactness of flat-holonomy moduli for a uniform transverse
      gap when stabilizer rank jumps;
\item a prime determinant as a continuous density across such a
      rank jump;
\item a finite integral of fiber traces for trace class of a
      continuous decomposable transfer operator;
\item trace-class magnitudes for convergence or positivity of an
      uncancelled anomaly phase.
\end{enumerate}
\end{audit}

\begin{proof}
These are respectively the drifting-carrier examples of V86--V87,
Proposition~\ref{prop:v15v90-symbol}, the zero-mode examples of V89,
Counterexamples~\ref{cnt:v15v91-multiplicity} and
\ref{cnt:v15v91-drifting}, V92,
Lemma~\ref{lem:v15v94-tensor},
Proposition~\ref{prop:v15v94-cartan},
Counterexample~\ref{cnt:v15v96-kernel-jump},
Counterexample~\ref{cnt:v15v97-jump},
Theorem~\ref{thm:v15v98-direct-integral}, and
Counterexample~\ref{cnt:v15v99-phase}.
\end{proof}

\section{Companion audit record}

The required audits at V90 and V95 inspected the current YM and NS
files.  YM V41 proved exact Wilson-pencil positivity on one selected
transverse quarter cell but withheld full-torus and Hamiltonian
claims.  The apparent YM V42 file was byte-identical to V41 and
contained no V42 continuation.  NS V26 remained a classical
rank-one Oseen-kernel norm refinement.  Neither supplied a theorem
used in the assembled transfer.

\section{Open acquisition ledger}

\begin{ledger}[Smallest unresolved inputs after V100]
\label{led:v15v100-open}
The assembled theorem still requires:
\begin{enumerate}
\item a lower-bounded, nuclear, reflection-compatible gravitational
      reference with a justified treatment of the conformal
      direction;
\item a uniform bosonic infrared mechanism after the physical gauge
      quotient, or a regulator-removal theorem that replaces it;
\item a global gauge atlas with controlled reducible strata,
      reduced determinant, and self-adjoint holonomy interfaces;
\item trace-norm convergence of the stable Higgs reference and
      cutoff-independent renormalized interaction constants;
\item anomaly-line, ghost, and reflection-sewing compatibility for
      the combined transfer;
\item cofinal independence, local projective consistency, Ward/BRST
      passage, and OS reconstruction for the resulting joint state;
\item identification of the limiting stress tensor and support on
      the ordered Einstein equation.
\end{enumerate}
\end{ledger}

These rows are not linearly ordered.  The gravitational first row
is upstream of the combined Gibbs theorem because
Lemma~\ref{lem:v15v94-tensor} prevents matter from compensating for
an unbounded conformal direction.  It is therefore the next
direction.

\section{Cycle conclusion}

\begin{assessment}[Position after the fifteenth cycle]
\label{ass:v15v100-position}
The cycle converts a collection of boundary spectral assumptions
into a rigorous conditional interacting transfer theorem and
pinpoints the exact failures at bosonic zero modes, reducible gauge
strata, continuous holonomy multiplicity, ultraviolet point fields,
and the gravitational conformal factor.

The full Standard Model--Einstein continuum has not been
constructed.  In particular, none of the seven rows in
Ledger~\ref{led:v15v100-open} may be inferred from the assembly
theorem.
\end{assessment}

\begin{decision}[Sixteenth-cycle start]
\label{dec:v15v100-next}
Freeze this V100 with suffix \(\texttt{\_f15}\), restart at a compact
V1 containing context and the major programme results, and then go
upstream to the gravitational reference.  The first substantive
continuation should compare three honest branches: a real coercive
stabilizer, a complex conformal contour, and a constraint-reduced
Lorentzian transfer.  For each branch, require lower boundedness or
sectorial control, heat nuclearity, reflection reality, and
compatibility with the Einstein source derivative before it can
enter V99's ledger.
\end{decision}

\begin{firewall}[Final fifteenth-cycle claim boundary]
\label{fire:v15v100}
V100 audits a conditional boundary transfer architecture.  It does
not prove the physical hypotheses in the open ledger, a Yang--Mills
or Higgs mass gap, a reflection-positive gravitational theory, or
the full Standard Model--Einstein continuum.
\end{firewall}

\section*{Version V100 append-only record}
\addcontentsline{toc}{section}{Version V100 append-only record}
The complete V99 source was retained before this continuation.  It
had \(979839\) bytes, \(28892\) lines, and SHA--256
\[
 \texttt{22E920B196E828B15E977F2EB9958CD227D398167C83849D964394778AF0DB20}.
\]
V100 audited the entire fifteenth-cycle boundary nuclearity,
interaction, gauge-quotient, holonomy, and assembly block and fixed
the gravitational reference as the next upstream acquisition.
\end{document}
